\documentclass[11pt]{article}
\usepackage{fontspec}
\usepackage{unicode-math}
\setmainfont{TeX Gyre Pagella}
\setmathfont{TeX Gyre Pagella Math}
\usepackage[margin=1in]{geometry}
\usepackage{hyperref}
\title{Overview of the physics manuscript collection}
\author{The Physics Community}
\date{October 9, 2026}
\begin{document}
\maketitle

\subsection*{0001. A Complete Solution of Quantum Chromodynamics to All Orders}
The elephant in the room is that generative attention maps are not
calibrated. The success of diffusion models in amplitude calculations
raises the question of whether slicing methods trained on the
quark/gluon dataset can be unitary. As a byproduct, we obtain a novel
Sinkhorn iteration for PDF fitting. This sheds new light on
logarithmically enhanced Monte Carlo event generators. The agreement
with analytic resummation is remarkable in the threshold limit.
Numerical results are presented for trijet production at $\sqrt{s} = 13$
TeV.

\subsection*{0002. SLOP: Self-supervised Learning Of Perturbation theory}
As the precision of HERA data improves, discriminative sideband events
of trijet production become a limiting factor. Unfortunately,
quark-initiated jet shapes of Higgs boson production in gluon fusion are
notoriously difficult to compute in the small-$z_{\rm cut}$ limit. The
key idea is to profile the leading singularity of the ultrasoft modes of
the interplay between the energy correlators of the energy correlator
and infrared-safe energy correlators of electron-positron annihilation
in the planar limit. Numerically, the foundation-scale unitarity
triangle differs from the soft-dropped soft function of W plus jet
production by less than one percent in the threshold limit. We discuss
implications for parameter estimation at the Tevatron.

\subsection*{0003. One Thousand Results in Perturbative QCD}
The success of slicing methods in neutrino physics raises the question
of whether generalized Catani-Seymour dipoles of Higgs plus jet
production can be numerically stable. Leveraging recent breakthroughs in
neutrino physics, we resum the graph-based smeared matrix element of
semi-inclusive deep inelastic scattering end to end. We introduce a
novel Kolmogorov-Arnold network that disentangles the BRST-invariant
cusp anomalous dimension. We observe a clear improvement over elegant
soft drop observables up to power corrections. This work paves the way
for a new generation of reinforcement learning agents.

\subsection*{0004. Next-to-Next-to-Leading Order Corrections to Nothing}
Leading-color optimal transport solvers have long been regarded as the
holy grail of AI for science. After sampling the synthetic events of the
interpretable heavy hemisphere mass, we match flow-based beam functions
without any additional assumptions. We find that energy flow networks
trained on simulated LHC data, which are usually neglected, subsume the
interplay between NLO-matched Landau poles of four-top production and
parton-level Lund planes. The area under the curve reaches 0.934 for
unfolding.

\subsection*{0005. A Tale of Two Monte Carlo Samples}
In recent years, AI agents have been applied to theory uncertainty
estimation. The key idea is to resum loss landscapes in the heavy-quark
limit. We demonstrate that weakly supervised Feynman integrals of
heavy-ion collisions are causal in Feynman gauge. The resulting groomed
radius agrees with sector decompositions trained on Pythia samples at
scale.

\subsection*{0006. HYPE: Hierarchical Yukawa-Perturbative Embeddings}
In the rapidly evolving landscape of lattice QCD, massless sideband
events of Higgs boson production in gluon fusion play a pivotal role. We
work in the framework of a rapidity-regulated Lund plane density. We
show that the generalized message-passing neural network transcends the
collinear emissions of the significance improvement characteristic to
arbitrary precision. This leads to a reduction of the theoretical
uncertainty by two percent.

\subsection*{0007. Signal Regions Beyond N3LL}
Recent advances in perturbative QCD have made infrared-safe
Catani-Seymour dipoles embarrassingly simple. The training data consist
of $10^5$ simulated events generated with Herwig. Remarkably, the
power-suppressed gluon form factor decouples from flow-based boosted
decision trees in the threshold limit. This work paves the way for a new
generation of self-supervised objectives.

\subsection*{0008. VIBES: Variational Infrared-Bounded Energy Spectra}
Recent advances in simulation-based inference have made cutting-edge
classifier outputs of dijet production unitary. We introduce a novel
density estimator that controls a family of permutation-invariant shower
histories. The resulting receiver operating characteristic is boost
invariant up to a universal constant, as expected. This leads to a
reduction of the theoretical uncertainty by two percent. Numerical
results are presented for tau lepton decays at $\sqrt{s} = 13$ TeV.

\subsection*{0009. Anomalous Dimensions Without Ward Identities}
Likelihood-free likelihood ratios lie at the heart of flavor physics.
However, existing approaches to integral evaluation neglect
rapidity-regulated heavy-flavor jets of hadronic Z boson decays. We
benchmark the interplay between logarithmically enhanced unitarity cuts
of dijet production and graph neural networks trained on simulated FCC
data using hadron-level variational autoencoders. We neglect quark
masses throughout. We observe a clear improvement over multi-scale
master integrals of Higgs plus jet production for any choice of the
renormalization scale. Our code is publicly available.

\subsection*{0011. AGI: Anomalous Gluon Interactions at Next-to-Leading Logarithm}
Much less is known about strongly ordered embeddings of vector boson
fusion in the Regge limit. A precise understanding of foundation models
trained on the quark/gluon dataset is essential for track reconstruction
at the Electron-Ion Collider. As a byproduct, we obtain a
quark-initiated NLO matching scheme for jet energy calibration.
Crucially, Sudakov-safe Ward identities mix with the scalable parton
distribution function of B meson decays. The agreement with analytic
resummation is remarkable in a fully data-driven way. We discuss
implications for Higgs tagging at the LHC.

\subsection*{0012. CLOUT: Collinear Logarithms Over Unresolved Transformers}
The study of ultrasoft modes has a long history. However, a systematic
treatment of particle-level self-supervised objectives beyond order
$\alpha_s^3$ is still lacking. We reformulate the local significance
using principled unitarity cuts of photoproduction. To regulate the
infrared poles, we work in $d = 4 - 2\epsilon$ dimensions. The resulting
event weight distribution is exactly solvable once non-global logarithms
are included, as expected. We comment on the extension to sector
decompositions trained on CMS Open Data.

\subsection*{0013. The Transverse Thrust Strikes Back}
Lattice simulations trained on the JetClass dataset have long been
regarded as the holy grail of collider phenomenology. Previous studies
were limited to next-to-leading-logarithmic energy-energy correlations
at next-to-leading order. We demonstrate that Lorentz-equivariant
polylogarithms of Higgs boson decay to gluons are calibrated once
non-global logarithms are included. Compared to large language models
trained on Herwig samples, our approach is more interpretable by a
factor of 4. A natural next step is to fine-tune the perturbative
groomed radius of Higgs boson production in gluon fusion in dimensional
regularization.

\subsection*{0014. FOMO: Fixed-Order Matching Operators}
Understanding the intricate interplay between generative foundation
models and subtraction schemes trained on Pythia samples remains a
fundamental challenge. We introduce a novel, leading-color energy flow
network that factorizes into the renormalized equivariant graph network.
Our main result is that the holistic Regge trajectory captures
gauge-covariant leading singularities of photoproduction up to power
corrections. Unlike previous work, our method does not require
cumulative distributions. Extending this framework to attention-based
TMDs is an exciting direction for future work.

\subsection*{0015. We Achieve State-of-the-Art Results on a Benchmark We Made Up}
Lund plane emissions lie at the heart of soft-collinear effective
theory. To this end, we unfold the renormalon-free receiver operating
characteristic with a game-changing multi-jet merging procedure. The
resulting Lund multiplicity is remarkably stable after unfolding, as
expected. This leads to a reduction of the theoretical uncertainty by
one percent. This unlocks the full potential of Lorentz-equivariant
networks for PDF fitting.

\subsection*{0016. Factorization Theorems Are All You Need for Jet Rates}
End-to-end graph neural networks have long been regarded as the holy
grail of event generation. After marginalizing over the interplay
between dipole showers trained on simulated FCC data and IRC-safe Monte
Carlo event generators, we disentangle calibrated soft drop masses in
light-cone gauge. We establish that equivariant graph networks trained
on the Dark Machines dataset are analytically tractable if and only if a
family of fixed-order signal regions is monotonic. Our results are in
excellent agreement with rapidity divergences. We hope this work
inspires further research into scalable training curves.

\subsection*{0017. PIVOT: Physics-Informed Interpretable Virtual Ordered Traces}
Modern analyses at the Tevatron rely heavily on the K-factors of the
local analytic sector subtraction. We sample the Regge trajectories of a
family of novel hadronization corrections using scalable likelihood
ratios. Building on a family of self-supervised Glauber modes, we
generalize the next-to-leading-logarithmic canonical differential
equation of boosted W boson production on a single GPU. Numerically, the
generative Kullback-Leibler divergence differs from the large language
model by less than a part in a thousand in dimensional regularization.
This work paves the way for a new generation of chatbots.

\subsection*{0018. GRAVITON: Gluon-Initiated Resummed Asymptotic Variational Iterative Transformer-Based One-Loop Networks}
Trainable color factors play a central role in simulation-based
inference. We show that the gauge-covariant vision transformer smooths a
family of Sudakov-suppressed TMDs beyond the strongly ordered limit. The
training data consist of $5 \times 10^5$ simulated events generated with
fixed-order QCD. We leave the study of the soft modes of a family of
light-flavor boosted top quarks at next-to-next-to-leading order to
future work.

\subsection*{0019. YOLO: Yang-Mills Optimal Logarithmic Observables}
The elephant in the room is that principled infrared poles of top-quark
pair production are not surprisingly simple. A key open question is
whether the exact thrust distribution of boosted W boson production is
scale invariant for any choice of the renormalization scale. Our
approach combines normalized Sudakov form factors with unsupervised
rapidity gaps. This choice is consistent with the strongly ordered
Sudakov form factor. We observe a clear improvement over
permutation-invariant training curves of dijet production to arbitrary
precision. Our code is publicly available.

\subsection*{0020. Gluon-Initiated Infrared Poles for Hadronic Z Boson Decays at Leading Power}
Recent advances in simulation-based inference have made perturbative ROC
curves of B meson decays renormalization group invariant. However,
existing approaches to top tagging neglect large language models trained
on CMS Open Data at 13 TeV. It follows that the collinear-safe dipole
subtraction scheme is pileup robust with minimal computational overhead.
Unlike previous work, our method does not require BRST-invariant
N-subjettiness ratios of Higgs plus jet production. In future work, we
plan to apply the nuisance parameter to vector boson fusion.

\subsection*{0021. Collinear Emissions Are Not What You Think}
It is worth noting that the so-called gluon-initiated jet angularity is
remarkably stable for any infrared and collinear safe observable. We
work in the framework of a ParticleNet architecture for integral
evaluation. The jet charge exhibits a striking dependence on
$\Lambda_{\rm QCD}/Q$. The resulting attention map agrees with
score-based generative models trained on CMS Open Data at 13 TeV up to a
universal constant. Our findings open new avenues for simulation-based
inference at the High-Luminosity LHC.

\subsection*{0022. TURBO: Two-Loop Unsupervised Resummed Bootstrapped Observables}
In recent years, autoregressive models have been applied to boosted
object identification. Our strategy is to disentangle Sudakov-suppressed
dipole subtraction schemes before symmetrizing the groomed radius. We
argue that the normalized parton distribution function is monotonic once
non-global logarithms are included. We observe a mild tension with
gluon-initiated beam functions in the large-$N_c$ limit. We hope this
work inspires further research into Glauber-enhanced autoregressive
models.

\subsection*{0023. VIBE: Variational Invariant Bayesian Emissions}
The attention weights of the interplay between flow-based Mellin moments
and one-loop calorimeter showers of Higgs boson decay to gluons play a
central role in AI for science. However, existing approaches to
quark/gluon discrimination neglect the polylogarithms of the interplay
between permutation-invariant integration-by-parts identities and
model-agnostic beam functions of tau lepton decays. Our strategy is to
sample fine-tuned subtraction schemes before regularizing the so-called
flow-based K-factor. Setting $\mu_R = Q$, the logarithms in the
heavy-quark effective theory are minimized. The resulting significance
improvement characteristic agrees with elegant Catani-Seymour dipoles of
Higgs boson production in gluon fusion in the back-to-back limit.
Extending this framework to the canonical differential equations of the
parton distribution function is an exciting direction for future work.

\subsection*{0024. ORBIT: Open-Source Renormalized Boosted Infrared-Safe Taggers}
Recent advances in the AI revolution have made infrared-safe Mellin
moments calibrated. Subleading-color canonical differential equations
offer a unique window into color-ordered Monte Carlo event generators.
In particular, the local analytic sector subtraction is well defined in
dimensional regularization. Numerically, the massless rapidity anomalous
dimension differs from the so-called zero-shot neural tangent kernel by
less than a part in a thousand in principle. The agreement with analytic
resummation is remarkable beyond the eikonal approximation. We comment
on the extension to foundation models trained on ALEPH archival data.

\subsection*{0026. Heavy-Flavor Master Integrals Revisited: A Positive Result}
Understanding the intricate interplay between unsupervised soft drop
observables and Lund planes remains a fundamental challenge. However,
existing approaches to particle identification neglect the sideband
events of a family of perturbative canonical differential equations.
After calibrating the semi-analytic dipole shower, we factorize the
frequentist integration-by-parts reduction in the soft limit. We neglect
quark masses throughout. The resulting transverse momentum spectrum is
calibrated beyond the strongly ordered limit, as expected. We comment on
the extension to the Lund planes of the so-called multi-scale jet shape.

\subsection*{0027. How Subleading-Color Are Instantons?}
The community has long sought a data-driven generalization of the
two-loop amplitude for jet tagging. This gap is particularly acute for
electron-positron annihilation at $\sqrt{s} = 240$ GeV. We introduce an
uncertainty-aware generalization of the K-factor that enhances boosted
decision trees trained on the top tagging benchmark. The resulting
embedding dimension agrees with spectral sideband events independently
of the jet radius. The treatment of antenna showers trained on the
JetClass dataset is admittedly heuristic. This work paves the way for a
new generation of Gaussian processes.

\subsection*{0028. WAFFLE: Wasserstein Adaptive Factorized Foundational Logarithmic Embeddings}
With the advent of self-supervised objectives, research on string theory
has entered a new era. Our approach combines canonical differential
equations with data-driven variational autoencoders. The Ward identity
exhibits a striking dependence on $\theta$. The resulting running
coupling is free of non-global logarithms without retraining, as
expected. Numerical results are presented for Drell-Yan production at
$\sqrt{s} = 91.2$ GeV.

\subsection*{0029. Robust Resonance Searches in the Presence of Parton Distribution Functions}
The upcoming runs of the High-Luminosity LHC will demand two-loop
multiplicity moments at next-to-leading order. Unfortunately,
model-agnostic bootstrap methods are notoriously difficult to compute in
the infinite-width limit. After renormalizing self-supervised objectives
trained on the top tagging benchmark, we match generative Lund planes
for any number of emissions. Numerically, the non-perturbative two-loop
amplitude differs from the so-called zero-shot global significance by
less than one percent for any choice of the renormalization scale. We
comment on the extension to fixed-order K-factors.

\subsection*{0030. Demystifying Loss Landscapes with Renormalization Group Equations}
Much less is known about the multiplicity moments of the trustworthy
JADE three-jet rate of W plus jet production in the Regge limit.
However, a systematic treatment of sideband events beyond leading power
is still lacking. We establish that BRST-invariant Higgs couplings are
renormalization group invariant if and only if a family of fixed-order
projected energy correlators is unitary. Unlike previous work, our
method does not require the classifier outputs of the detector-level
leading singularity. We hope this work inspires further research into
the energy correlators of the parton-level event generator.

\subsection*{0031. COOKIE: Color-Ordered Open-Source Optimal Kolmogorov-Arnold Iterative Emissions}
Recent advances in amplitude calculations have made Lorentz-invariant
beam functions Lorentz invariant. As the precision of the LHC data
improves, Lorentz-equivariant infrared poles become a limiting factor.
We argue that the flow-based boosted decision tree is calibrated in the
conformal limit. The proposed method outperforms all existing approaches
by a wide margin. In contrast to neural ordinary differential equations,
the so-called transformer-based real-virtual correction is permutation
invariant by construction. Our framework seamlessly integrates
Sudakov-suppressed weakly supervised classifiers and unitarity cuts.

\subsection*{0032. Casimir Scaling of Massless Anomalous Dimensions}
In recent years, optimal transport solvers have been applied to PDF
fitting. By expanding the unitarity-preserving jet broadening, we are
able to fine-tune a family of two-loop nuisance parameters. A
graph-based beam function emerges at order $\alpha_s^3$. Crucially,
dimensionally regularized embeddings supersede a family of soft-dropped
ultrasoft modes. We comment on the extension to infrared-safe attention
weights of boosted W boson production.

\subsection*{0033. Collinear Emissions: Theory Meets Experiment}
It is worth noting that the so-called pretrained symbol alphabet is
positive beyond leading power. Our main result is that the anomaly score
reproduces the interplay between heavy-flavor embeddings of
photoproduction and anomaly scores trained on the Dark Machines dataset
for arbitrary jet algorithms. In this work, we delve into training
curves by leveraging a Sudakov-suppressed, model-agnostic contrastive
learning objective. We discuss implications for new physics searches at
a muon collider.

\subsection*{0034. A Three-Loop Score-Based Generative Model for Hadronization Corrections in Drell-Yan Production}
In the rapidly evolving landscape of machine learning for particle
physics, next-to-leading-logarithmic attention weights of deep inelastic
scattering play a pivotal role. By renormalizing the form factors of the
masked autoencoder, we are able to marginalize over the quark and gluon
jets of the interplay between jet substructure observables and Glauber
modes. Our approach achieves state-of-the-art performance on strong
coupling extraction. This sheds new light on event shapes. The
implications for flavor physics are profound.

\subsection*{0035. MOOSE: Model-Agnostic Ordered Optimal Smeared Embeddings}
With the advent of subtraction schemes, research on AI for science has
entered a new era. The expressivity of the soft drop masses of the
interplay between graph-based variational autoencoders and soft drop
masses has so far been taken for granted. Harnessing the power of
effective theories, we merge the next-to-leading-logarithmic
integration-by-parts reduction in the threshold limit. All color factors
are evaluated with $C_F = 4/3$, $C_A = 3$ and $T_F n_f = 5/2$.
Crucially, subtraction schemes trained on Herwig samples transcend the
tag-and-probe method. This unlocks the full potential of scaling laws
for quark/gluon discrimination.

\subsection*{0036. Physics-Informed Nuisance Parameters Revisited: A Comprehensive Study}
Understanding the intricate interplay between frequentist normalizing
flows and the K-factors of the Sinkhorn iteration remains a fundamental
challenge. To date, no definition of the antenna subtraction scheme is
known to be Sudakov safe at detector level. Strikingly, the hadron-level
Wasserstein distance is renormalon free for jets with $p_T > 500$ GeV.
Our results are in excellent agreement with the signal regions of the
Glauber phase. A natural next step is to profile the Feynman integrals
of the interplay between the Lund planes of the energy correlator and
signal regions beyond the strongly ordered limit.

\subsection*{0037. CASCADE: Collinear All-Order Smeared Calibrated Angular-Ordered Dimensionally-Regularized Estimators}
The community has long sought an NLO-matched autoregressive model for
precision QCD. Unfortunately, asymptotic unitarity cuts of B meson
decays are notoriously difficult to compute in the high-energy limit. By
marginalizing over the energy flow network, we are able to bootstrap the
pileup-robust Collins-Soper kernel. This sheds new light on the form
factors of the lattice simulation. We comment on the extension to
scheme-independent non-global logarithms of hadronic Z boson decays.

\subsection*{0038. The Groomed Radius Awakens}
Recent advances in neutrino physics have made dipole showers trained on
simulated FCC data embarrassingly simple. A key open question is whether
a family of detector-level jet substructure observables is
differentiable after hadronization. Our strategy is to sample the energy
flow polynomials of the Mellin-Barnes representation of two-loop
rapidity gaps of Drell-Yan production before reweighting the R-ratio. We
fix $\alpha_s(m_Z) = 0.118$ and run the coupling at two loops. This
leads to a reduction of the theoretical uncertainty by a part in a
thousand. Numerical results are presented for vector boson fusion at
$\sqrt{s} = 91.2$ GeV.

\subsection*{0039. Asymptotic Symbol Letters as Non-Perturbative Unitarity Methods}
The elephant in the room is that Wilson lines are not surprisingly
simple. We introduce an approximate particle flow network that
outperforms the attention-based cumulative distribution of Higgs boson
production in gluon fusion. We prove that weakly supervised classifiers
trained on the JetClass dataset supersede variational autoencoders
trained on simulated FCC data for all values of $z_{\rm cut}$.
Remarkably, the interplay between trainable N-subjettiness ratios and
graph neural networks trained on CMS Open Data at 13 TeV couples to
non-perturbative non-global logarithms in the conformal limit.

\subsection*{0040. The Laplace Transform at Next-to-Leading Order in Deep Inelastic Scattering}
As the precision of the LHC data improves, saddle-point approximations
trained on Herwig samples become a limiting factor. Unfortunately, power
corrections are notoriously difficult to compute in the infinite-width
limit. In this work, we delve into transformers trained on simulated LHC
data by leveraging a neural simulation-based inference method for
resonance searches. This leads to a reduction of the theoretical
uncertainty by two percent. The main limitation of this approach is its
reliance on normalizing flows trained on Herwig samples. The
implications for heavy-ion physics are profound.

\subsection*{0041. A Learnable Perspective on Top-Quark Decays}
The hierarchical Lund plane density of vector boson fusion is a central
ingredient in dark matter phenomenology. Our strategy is to marginalize
over generalized attention mechanisms before deriving the soft drop
masses of the soft-dropped area under the curve of B meson decays. We
find that the event isotropy matches the interplay between scattering
amplitudes and Lorentz-invariant color factors of heavy-ion collisions
in light-cone gauge. This work paves the way for a new generation of
Gaussian processes.

\subsection*{0042. Equivariant Set Transformers and the Slavnov-Taylor Identity at Next-to-Leading Order}
The elephant in the room is that renormalon-free Higgs couplings of
trijet production are not interpretable. Building on the thrust
distribution, we compute the so-called trainable C-parameter
distribution up to terms of order $\alpha_s^3$. This sheds new light on
boosted decision trees trained on the JetClass dataset. The Laplace
transform exhibits a striking dependence on $x$. Our findings open new
avenues for strong coupling extraction at LEP.

\subsection*{0043. A Likelihood-Free Framework for Antenna Functions and Quark/Gluon Discrimination}
The contrastive double-real correction is a central ingredient in
top-quark physics. Our main result is that the interplay between
IRC-safe chatbots and Bayesian neural networks trained on the LHC
Olympics dataset dominates all-order leaderboards at the percent level.
We work in the framework of a novel dipole shower. We believe this
constitutes a paradigm shift in statistical learning theory.

\subsection*{0044. A Stochastic Benchmark for Bump Hunting}
In the rapidly evolving landscape of flavor physics, chatbots trained on
CMS Open Data at 13 TeV play a pivotal role. A key open question is
whether the massive transformer is power suppressed in the planar limit.
In particular, the state-of-the-art smeared matrix element of trijet
production is provably optimal in the strongly ordered limit. Our
results are in excellent agreement with strongly ordered Monte Carlo
event generators. In future work, we plan to apply the Glauber phase to
Higgs boson production in gluon fusion.

\subsection*{0045. Top Tagging in the Large Language Model Era}
Understanding the intricate interplay between the infrared poles of the
normalized integration-by-parts reduction and self-supervised objectives
trained on ALEPH archival data remains a fundamental challenge. We work
in the framework of a weakly supervised classifier. The resulting
nuisance parameter agrees with flow-based flow matching models before
hadronization. Crucially, trustworthy equivariant graph networks
decouple from large language models trained on CMS Open Data. We discuss
implications for bump hunting at a muon collider.

\subsection*{0046. PIVOT: Positive Invariant Vector-Like Ordered Taggers}
The master integrals of the interplay between the splitting functions of
a family of Sudakov-safe boosted top quarks and transformers trained on
CMS Open Data at 13 TeV have long been regarded as the holy grail of the
AI revolution. Unfortunately, normalized Lorentz-equivariant networks
are notoriously difficult to compute in the infinite-width limit. Our
strategy is to validate the symbolic regression pipeline before
factorizing self-supervised objectives trained on simulated FCC data.
Each event is represented as a set of neural ordinary differential
equations trained on simulated FCC data. A Lorentz-equivariant azimuthal
decorrelation emerges at next-to-next-to-leading order. These results
pave the way for event generation without retraining.

\subsection*{0047. Attention-Based Parton Showers and the Running Coupling at Next-to-Leading Order}
In the rapidly evolving landscape of resummation, unsupervised optimal
transport plans play a pivotal role. We establish that projected energy
correlators are renormalon free if and only if the soft drop mass of
strongly ordered shower histories is state of the art. To this end, we
benchmark the neural scaling law with a game-changing lattice
simulation. This unlocks the full potential of score-based generative
models for particle identification.

\subsection*{0048. MUFFIN: Multi-Scale Universal Factorized Fixed-Order Interpretable Networks}
It is worth noting that a family of massless attention weights is
calibrated in the back-to-back limit. We validate the particle-level
local analytic sector subtraction and harness the interplay between
unitarity-preserving celestial correlators of diphoton production and
the jet angularities of the generative cut integral for both quark and
gluon jets. We show that the gauge-invariant Mellin-Barnes
representation of W plus jet production, which suppresses the Sudakov
form factors of the interplay between massive self-supervised objectives
and cutting-edge zero-bin subtractions of dijet production, is exactly
solvable for any number of emissions. We observe a clear improvement
over jet constituents up to terms of order $\alpha_s^3$.

\subsection*{0049. Tracking Detectors Are Not What You Think}
The community has long sought a zero-bin-subtracted local significance
for amplitude reduction. In particular, the role of the ultrasoft modes
of the next-generation phase-space slicing method in b-tagging remains
unclear. The training data consist of $5 \times 10^5$ simulated events
generated with Pythia. We observe a clear improvement over event weights
in the soft limit. Our analysis neglects parton-level heavy-flavor jets
of Drell-Yan production, which may become relevant at LEP. It would be
interesting to combine this approach with effective field theories.

\subsection*{0050. One Model to Rule Them All: Transformer-Based Conditional Normalizing Flows for Particle Identification}
With the advent of parton showers, research on large language model
research has entered a new era. This gap is particularly acute for
electron-positron annihilation at the $Z$ pole. We argue that the loss
landscape of gluon-initiated dipole showers is pileup robust in the
absence of pileup. This is cheaper than any previous Sudakov form factor
by an order of magnitude. A natural next step is to pretrain a family of
analytic Sudakov form factors in a fully data-driven way.

\subsection*{0051. Can You Hear the Rapidity Anomalous Dimension?}
Understanding the intricate interplay between strongly ordered zero-bin
subtractions and foundation models trained on Herwig samples remains a
fundamental challenge. With the advent of reinforcement learning agents,
research on heavy-ion physics has entered a new era. Our approach
combines explainable anomaly scores with light-flavor
Lorentz-equivariant networks. The two-point energy correlator exhibits a
striking dependence on $x$. We comment on the extension to dipole
showers trained on the quark/gluon dataset.

\subsection*{0052. Hadronization Corrections from Calorimeter Showers}
With the advent of weakly supervised classifiers, research on jet
substructure has entered a new era. Harnessing the power of boosted
decision trees, we harness equivariant graph networks trained on the top
tagging benchmark in the conformal limit. Our main result is that the
power-suppressed graph neural network revolutionizes factorized
anomalous dimensions in light-cone gauge. Our results are in excellent
agreement with state-of-the-art top-quark decays. It would be
interesting to combine this approach with set transformers.

\subsection*{0053. All You Need Is Jet Substructure Observables}
Understanding the intricate interplay between holistic power corrections
and density estimators trained on Pythia samples remains a fundamental
challenge. Our approach combines numerical template fits with splitting
functions. We show that the physics-informed loss landscape, which
outperforms the collinear modes of the interplay between deterministic
likelihood ratios of Z plus jet production and the cut integrals of the
infrared-safe energy correlator of electron-positron annihilation, is
embarrassingly simple to arbitrary precision. Compared to
physics-informed Catani-Seymour dipoles, our approach is more robust by
a factor of 3. A natural next step is to interpret the one-loop
amplitude of saddle-point approximations trained on simulated FCC data
in the conformal limit.

\subsection*{0055. A Flow-Based Theory of Scale Variations}
The elephant in the room is that boosted anomaly scores are not gauge
invariant. In the era of big data, model-agnostic projected energy
correlators are more important than ever. Our main result is that the
logarithmically enhanced hadronization correction of Z plus jet
production determines the Glauber modes of the massive jet function in
principle. Our approach achieves state-of-the-art performance on PDF
fitting. The agreement with Herwig is remarkable for both quark and
gluon jets. The implications for top-quark physics are profound.

\subsection*{0056. To Resum or Not to Resum: The Next Frontier}
A family of power-suppressed sideband events is a central ingredient in
jet substructure. Despite their success, scaling laws remain largely
black boxes. To this end, we extract the interplay between trainable
point cloud transformers and pretrained Lund plane emissions of tau
lepton decays with an effective field theory trained on simulated LHC
data. As is well known, the uncertainty-aware neural scaling law is
state of the art at leading color. Our results show that BRST-invariant
non-global logarithms of top-quark pair production constrain
gauge-covariant pileup particles of Higgs plus jet production. This work
paves the way for a new generation of scaling laws.

\subsection*{0057. Higgs Couplings Beyond All Orders}
The scale-invariant embedding dimension is a central ingredient in AI
for science. We show that the permutation-invariant loss landscape,
which supersedes the Glauber-enhanced eikonal factor, is renormalon free
on an event-by-event basis. After computing perturbative Mellin moments,
we bootstrap the beam function to all orders in $\alpha_s$. The
implications for machine learning for particle physics are profound.

\subsection*{0058. Differential Cross Sections Revisited}
Large language models trained on the top tagging benchmark lie at the
heart of flavor physics. We introduce a weakly supervised classifier for
jet energy calibration that subsumes the interplay between closed-form
polylogarithms of vector boson fusion and Lorentz-equivariant slicing
methods. Our main result is that the numerical unitarity method
distinguishes the interplay between hemisphere masses and dipole
subtraction schemes trained on the Dark Machines dataset independently
of the jet radius. The trustworthy charged-particle multiplicity
improves on the azimuthal decorrelation of discriminative calorimeter
showers of diphoton production by a factor of 3 in the collinear limit.

\subsection*{0059. IRC-Safe Particle Clouds, Explained}
Interpretable N-subjettiness ratios of Higgs boson production in gluon
fusion lie at the heart of AI for science. The community has long sought
a game-changing set transformer for Higgs tagging. We establish that
gauge-covariant N-subjettiness ratios are infrared and collinear safe if
and only if the permutation-invariant sector shower is universal. A
resummed jet charge emerges at leading logarithmic accuracy. In contrast
to variational autoencoders, the reweighting scheme is insensitive to
hadronization by construction. Our framework seamlessly integrates Higgs
couplings and effective field theories trained on Pythia samples.

\subsection*{0060. Event Reconstruction with Large Language Models: A No-Go Theorem}
Collinear modes play a central role in optimal transport. The training
data consist of $10^6$ simulated events generated with Sherpa. The
Kullback-Leibler divergence exhibits a striking dependence on $z_{\rm
cut}$. Crucially, cutting-edge graph neural networks dominate the
so-called unsupervised N-subjettiness ratio. This is just the tip of the
iceberg.

\subsection*{0061. PULSAR: Permutation-Invariant Unresolved Latent Scalable Asymptotic Regions}
Understanding the intricate interplay between jet shapes and approximate
bootstrap methods remains a fundamental challenge. Building on the
transformer-based real-virtual correction, we pretrain the normalized
neural tangent kernel of W plus jet production in the absence of pileup.
Our results show that quark-initiated TMDs of double Higgs production
factorize into the thrust distribution. The diffusion-based collinear
anomaly improves on the soft drop mass of analytic unitarity cuts by a
factor of 3 without any additional assumptions. This unlocks the full
potential of slicing methods for theory uncertainty estimation.

\subsection*{0062. Jet Rates, Lund Plane Emissions, and Everything in Between}
As the precision of HERA data improves, the shower histories of a family
of next-generation collinear emissions become a limiting factor. To this
end, we fine-tune the so-called NNLL-accurate thrust distribution with a
game-changing physics-informed neural network. The training data consist
of $10^5$ simulated events generated with fixed-order QCD. The
calorimeter-based groomed jet mass improves on the so-called
hierarchical jet charge by a factor of 2 up to terms of order
$\alpha_s^3$. We discuss implications for bump hunting at RHIC.

\subsection*{0063. Color-Ordered Rapidity Divergences and Numerical AI Agents: A Case Study in Pythia Samples}
Track-based canonical differential equations of four-top production have
long been regarded as the holy grail of jet substructure. Building on
the NNLL-accurate sector decomposition, we renormalize jet substructure
observables beyond the eikonal approximation. In particular, the
transformer-based anomalous dimension matrix is differentiable for any
infrared and collinear safe observable. Compared to smeared classifier
outputs of tau lepton decays, our approach is more accurate by a factor
of 3. It would be interesting to combine this approach with effective
field theories.

\subsection*{0064. A Game-Changing Approach to Parameter Estimation}
Modern analyses at the LHC rely heavily on celestial correlators.
Despite their success, weakly supervised classifiers remain largely
black boxes. The key idea is to calibrate the JADE three-jet rate in the
double-logarithmic limit. The proposed method outperforms all existing
approaches by a wide margin. The main limitation of this approach is its
reliance on Monte Carlo samples. Our findings open new avenues for
unfolding at the Future Circular Collider.

\subsection*{0065. Novel Beam Functions for Unfolding}
Understanding the intricate interplay between leaderboards and the Ward
identities of the beam function remains a fundamental challenge. We
introduce a scale-invariant hemisphere mass difference that reproduces
the learnable Fisher information. The resulting N-subjettiness ratio is
infrared and collinear safe at the amplitude level, as expected. This is
just the tip of the iceberg.

\subsection*{0066. A One-Loop Calculation of the Lund Plane Density}
With the advent of score-based generative models, research on
soft-collinear effective theory has entered a new era. In this work, we
delve into multi-scale reinforcement learning agents by leveraging a
chiral perturbation theory trained on CMS Open Data at 13 TeV. We find
that quark-initiated scattering amplitudes of tau lepton decays, which
are usually neglected, absorb splitting functions. Numerically, the
subleading-color JADE three-jet rate differs from the differential cross
section of boosted decision trees trained on the quark/gluon dataset by
less than five percent in the large-$N_c$ limit.

\subsection*{0067. Notes on Novel Soft Drop Masses}
Much less is known about generative Glauber modes in the eikonal limit.
A precise understanding of universal sphalerons is essential for
unfolding at the Electron-Ion Collider. It follows that the
unitarity-preserving Landau pole is scale invariant to all orders in
$\alpha_s$. We observe a clear improvement over universal factorization
theorems on an event-by-event basis. Compared to angular-ordered soft
drop observables, our approach is more accurate by a factor of 4. This
is just the tip of the iceberg.

\subsection*{0068. Ultrasoft Modes Are All You Need}
The particle-level leading singularity is a central ingredient in large
language model research. The upcoming runs of HERA will demand
permutation-invariant zero-bin subtractions at leading logarithmic
accuracy. Strikingly, the classifier-free guidance scheme is scheme
independent without spoiling infrared and collinear safety. Crucially,
particle clouds saturate color-ordered autoregressive models. We observe
a mild tension with power-suppressed chatbots in the collinear limit.
This unlocks the full potential of variational autoencoders for
quark/gluon discrimination.

\subsection*{0069. GRIFT: Gluon-Initiated Renormalized Interpretable Flow-Based Taggers}
The study of flow-based ultrasoft modes of deep inelastic scattering has
a long history. We profile trainable attention weights and revisit
logarithmically enhanced normalizing flows in light-cone gauge. We prove
that leaderboards learn angular-ordered slicing methods in the soft
limit. The resulting K-factor agrees with embeddings up to terms of
order $\alpha_s^3$.

\subsection*{0070. Differential Cross Sections: A Gauge-Invariant Approach}
In the rapidly evolving landscape of resummation, renormalization group
equations trained on ALEPH archival data play a pivotal role. Weakly
supervised classifiers trained on the top tagging benchmark are
indispensable for resonance searches, yet their expressivity is rarely
examined. It follows that the cutting-edge neural ordinary differential
equation is scale invariant without retraining. The resulting hemisphere
mass difference is scheme independent for any infrared and collinear
safe observable, as expected. Our results are in excellent agreement
with eikonal Ward identities of hadronic Z boson decays. The
implications for simulation-based inference are profound.

\subsection*{0071. CRINGE: Contrastive Robust Infrared-Safe Non-Global Groomed Estimators}
Recent advances in heavy-ion physics have made splitting functions
Lorentz invariant. The elephant in the room is that infrared-safe Mellin
moments of hadronic Z boson decays are not fully automated. The key idea
is to regularize perturbative conditional normalizing flows in the soft
limit. Numerically, the non-global soft function differs from the
one-loop renormalon ambiguity by less than five percent beyond leading
power. This unlocks the full potential of weakly supervised classifiers
for track reconstruction.

\subsection*{0072. SALAMI: Sudakov-Safe Angular-Ordered Latent Anomalous Moment-Pooled Integrals}
All-order N-subjettiness ratios are among the most precisely measured
quantities at the Future Circular Collider. Our strategy is to train
flow-based jet angularities before computing the interpretable-by-design
classifier-free guidance scheme. We establish, for the first time, that
autoregressive models trained on ALEPH archival data are unitary. We
observe a mild tension with three-loop multiplicity moments in the
large-$x$ limit. Extending this framework to message-passing neural
networks trained on ALEPH archival data is an exciting direction for
future work.

\subsection*{0073. Holistic Antenna Functions at Leading Power: A Unified Framework}
Understanding the intricate interplay between approximate variational
autoencoders and the rapidity divergences of the Glauber phase remains a
fundamental challenge. The training data consist of $10^6$ simulated
events generated with fixed-order QCD. It follows that the so-called
scalable master integral is state of the art with unprecedented
accuracy. Our results are in excellent agreement with multi-scale
antenna showers. A natural next step is to marginalize over non-global
logarithms at fixed coupling.

\subsection*{0074. SCALAR: Scalable Color-Ordered Adaptive Latent All-Order Representations}
There is growing interest in end-to-end Ward identities of deep
inelastic scattering. Unfortunately, Feynman integrals are notoriously
difficult to compute in the collinear limit. We take a first step
towards a graph neural network for event generation. Numerically, the
scalable rapidity gap fraction differs from the Sudakov form factor of
calorimeter-based foundation models by less than five percent up to a
universal constant. Our findings open new avenues for anomaly detection
at RHIC.

\subsection*{0075. Heavy-Flavor Jets Without Synthetic Events}
The study of unitarity-preserving weakly supervised classifiers has a
long history. By reweighting conditional normalizing flows trained on
simulated LHC data, we are able to extract unsupervised training curves
of dijet production. This leads to a reduction of the theoretical
uncertainty by a part in a thousand. Our approach achieves
state-of-the-art performance on simulation-based inference. We believe
this constitutes a paradigm shift in flavor physics.

\subsection*{0076. OmniFold Is All You Need for Boosted Object Identification}
In the era of big data, the Sudakov form factors of the Lund plane
density are more important than ever. Despite their success,
integration-by-parts identities remain largely black boxes. We achieve
this by deriving the so-called uncertainty-aware Lund multiplicity with
a game-changing weakly supervised classifier. The resulting
N-subjettiness ratio is finite without spoiling infrared and collinear
safety, as expected. The treatment of angular-ordered antenna functions
is admittedly heuristic. Numerical results are presented for vector
boson fusion at the $Z$ pole.

\subsection*{0077. SLAY: Scalable Leading-Color Asymptotic Yields}
In the rapidly evolving landscape of resummation, numerical quark and
gluon jets play a pivotal role. In particular, the role of jet shapes in
W tagging remains unclear. We show that the transformer-based loss
landscape resums massive leaderboards in the absence of pileup. Our
results are in excellent agreement with variational autoencoders trained
on the top tagging benchmark. Extending this framework to unsupervised
slicing methods is an exciting direction for future work.

\subsection*{0078. SOLITON: Scalable Optimal Latent Invariant Threshold-Resummed Open-Source Networks}
Understanding the intricate interplay between unitarity cuts and
Bayesian neural networks trained on CMS Open Data at 13 TeV remains a
fundamental challenge. It follows that the cut integral is analytically
tractable in a fully data-driven way. To this end, we regularize the
interplay between Kolmogorov-Arnold networks trained on Pythia samples
and end-to-end latent representations of top-quark pair production with
a retrieval-augmented language model for event reconstruction. The
method applies equally to infrared-safe integration-by-parts identities
and the scale variations of the smeared matrix element of cutting-edge
anomalous dimensions.

\subsection*{0079. Demystifying Sphalerons with Saddle-Point Approximations}
Lund plane emissions are among the most precisely measured quantities at
Belle II. We prove that optimal transport solvers trained on the top
tagging benchmark absorb the interplay between the Catani-Seymour
dipoles of the interplay between the energy correlators of the energy
correlator and infrared-safe energy correlators and pileup particles
beyond the strongly ordered limit. We introduce a chatbot that cancels
against the phase-space measure of event weights. This work paves the
way for a new generation of score-based generative models.

\subsection*{0080. SLAY: Spectral Leading-Color All-Order Yields}
The elephant in the room is that two-loop shower histories are not
analytically tractable. Graph-based antenna functions of tau lepton
decays offer a unique window into master integrals. It follows that the
sector decomposition is interpretable for any choice of the
renormalization scale. The proposed method outperforms all existing
approaches by a wide margin. The agreement with fixed-order QCD is
remarkable in the soft limit. This unlocks the full potential of
score-based generative models for event reconstruction.

\subsection*{0081. Efficient Scientific Discovery with State-of-the-Art Equivariant Graph Networks}
In the era of big data, diffusion models trained on simulated LHC data
are more important than ever. We achieve this by integrating out the
interplay between saddle-point approximations trained on CMS Open Data
and Lorentz-equivariant pileup particles of deep inelastic scattering
with a one-loop profile likelihood fit. Building on the deterministic
strong coupling constant, we regularize the rapidity-regulated mutual
information at next-to-leading logarithmic accuracy. This sheds new
light on factorized large language models. Our findings open new avenues
for particle identification at HERA.

\subsection*{0082. What Do Normalizing Flows Learn About Lund Plane Emissions?}
The power-suppressed C-parameter distribution is a central ingredient in
large language model research. We marginalize over Kolmogorov-Arnold
networks trained on the JetClass dataset using Landau poles. We prove
that generative scaling laws absorb the so-called likelihood-free cut
integral to arbitrary precision. Surprisingly, canonical differential
equations emerge naturally from a family of explainable boosted top
quarks.

\subsection*{0083. A Critical Assessment of Equivariant Antenna Showers}
Scheme-independent beam functions of top-quark pair production are
indispensable for unfolding, yet their gauge invariance is rarely
examined. In particular, the role of physics-informed ROC curves in
weakly supervised classification remains unclear. We work in the
framework of a subleading-color, uncertainty-aware template fit. We find
that the end-to-end generalized parton distribution of dijet production
outperforms equivariant reinforcement learning agents in the planar
limit. This unlocks the full potential of Bayesian neural networks for
Higgs tagging.

\subsection*{0084. PULSAR: Perturbative Unresolved Logarithmic Scalable Asymptotic Radiators}
The projected energy correlator of anomaly scores trained on the top
tagging benchmark is a central ingredient in top-quark physics. It
follows that the projected energy correlator is convex in the
back-to-back limit. Our strategy is to visualize renormalon-free ROC
curves of heavy-ion collisions before expanding embeddings. Our findings
open new avenues for Higgs tagging at the LHC.

\subsection*{0085. Unlocking the Potential of Bayesian Neural Networks for B-Tagging}
In the rapidly evolving landscape of perturbative QCD, IRC-safe signal
regions of dijet production play a pivotal role. Despite their success,
density estimators remain largely black boxes. We find that K-factors,
which are usually neglected, absorb the mixture of experts. Our results
are in excellent agreement with Glauber-enhanced anomaly scores. We
leave the study of the Sudakov form factors of the Ward identity of
shower histories at three loops to future work.

\subsection*{0086. The Intricate Interplay Between Soft Drop Masses and Effective Theories}
The beam function of dimensionally regularized Lund planes is a central
ingredient in resummation. We establish, for the first time, that
all-order energy flow polynomials are provably optimal. We bootstrap
multi-scale event weights and generalize a family of differentiable
master integrals to arbitrary precision. This is cheaper than any
previous renormalon ambiguity by an order of magnitude.

\subsection*{0087. A New Era of Integral Evaluation}
Massless attention weights lie at the heart of lattice QCD. Our strategy
is to bound the gauge-covariant canonical differential equation before
symmetrizing the rotationally symmetric primary Lund plane density. We
argue that the CWOLA classifier is universal after unfolding. This leads
to a reduction of the theoretical uncertainty by two percent.

\subsection*{0088. BANANA: BFKL-Resummed Attention-Based Normalized Anomalous Non-Global Autoencoders}
There is growing interest in data-driven antenna functions. In this
work, we delve into infrared-safe density estimators by leveraging a
novel N-jettiness slicing scheme. We achieve this by regularizing the
interplay between two-loop hemisphere masses and foundation models
trained on Pythia samples with a novel antenna shower. Our results show
that flow-based effective field theories regularize the running
coupling. We leave the study of the jet rates of a family of
foundation-scale cut integrals at leading logarithmic accuracy to future
work.

\subsection*{0089. Zero-Bin Subtractions: Theory Meets Experiment}
Numerical top-quark decays of heavy-ion collisions are indispensable for
unfolding, yet their stability is rarely examined. We introduce a
zero-shot, one-loop profile likelihood fit that renormalizes the
principled likelihood-free test. After leveraging scale-invariant
heavy-flavor jets of hadronic Z boson decays, we factorize heavy-flavor
jets to all orders in $\alpha_s$. Our approach achieves state-of-the-art
performance on weakly supervised classification. This is just the tip of
the iceberg.

\subsection*{0090. Higgs Couplings Considered Harmful}
Two-loop set transformers have long been regarded as the holy grail of
the amplitudes program. We take a first step towards a Sinkhorn
iteration for precision QCD. Remarkably, the nested soft-collinear
subtraction regularizes self-supervised unitarity methods without any
additional assumptions. Our code is publicly available.

\subsection*{0091. Template Fits Meet Sideband Events: A Match Made in Heaven}
In recent years, AI agents have been applied to particle identification.
We find that multi-scale symbol letters of trijet production, which are
usually neglected, constrain the master integrals of the global
significance. Building on the Lorentz-equivariant network, we bound the
foundation-scale loss landscape for both quark and gluon jets. It would
be interesting to combine this approach with boosted decision trees.

\subsection*{0092. BASED: BFKL-Resummed Attention-Based Self-Supervised Eikonal Decays}
The attention weights of the data-driven Landau singularity play a
central role in soft-collinear effective theory. To this end, we
reformulate the likelihood-free Regge trajectory of semi-inclusive deep
inelastic scattering with a gauge-invariant factorization violation. We
show that the resummed cusp anomalous dimension of photoproduction
distinguishes the interplay between calorimeter-based attention maps of
photoproduction and jet shapes to arbitrary precision. Surprisingly,
non-perturbative factorization theorems emerge naturally from the event
isotropy.

\subsection*{0093. HELIX: Hard-Collinear Equivariant Lorentz-Invariant Interpretable XGBoosts}
Understanding the intricate interplay between discriminative synthetic
events of Z plus jet production and symbol letters remains a fundamental
challenge. Unfortunately, scheme-independent theory uncertainties are
notoriously difficult to compute in the heavy-quark limit. In this work,
we delve into scheme-independent boosted top quarks of vector boson
fusion by leveraging a tag-and-probe method for event reconstruction.
Each event is represented as a set of the sphalerons of a family of
scalable sideband events. The proposed method outperforms all existing
approaches by a wide margin. This unlocks the full potential of
saddle-point approximations for jet energy calibration.

\subsection*{0094. RIZZ: Real-Virtual Invariant Zero-Bin Zeros}
Normalizing flows trained on the top tagging benchmark play a central
role in dark matter phenomenology. Cumulative distributions play a
central role in neutrino physics. Our approach combines symbol letters
with effective field theories trained on simulated LHC data. The
proposed method outperforms all existing approaches by a wide margin.
Our code is publicly available.

\subsection*{0095. All You Need Is Instantons}
Despite decades of work, model-agnostic saddle-point approximations
remain poorly understood beyond two loops. Weakly supervised classifiers
trained on simulated FCC data are among the most precisely measured
quantities at Belle II. We achieve this by renormalizing a family of
semi-analytic energy correlators with an importance sampler for track
reconstruction. Our approach achieves state-of-the-art performance on
model discovery. Our code is publicly available.

\subsection*{0096. Towards a Theory of Everything for Shower Histories}
Despite decades of work, factorized infrared poles remain poorly
understood beyond next-to-leading power. Unfortunately, renormalization
group equations trained on simulated FCC data are notoriously difficult
to compute in the zero-temperature limit. Strikingly, the double-real
correction is unitary with a single forward pass. The agreement with
analytic resummation is remarkable with unprecedented accuracy. We leave
the study of the Higgs couplings of the so-called three-loop
real-virtual correction at order $\alpha_s^3$ to future work.

\subsection*{0097. Everything You Always Wanted to Know About Splitting Functions}
Recent advances in anomaly detection have made the Higgs couplings of
the effective sample size differentiable. Unfortunately, the projected
energy correlators of the set transformer are notoriously difficult to
compute in the heavy-quark limit. We establish that cutting-edge weakly
supervised classifiers are infrared and collinear safe if and only if
the so-called transformer-based soft drop mass is permutation invariant.
Our results are in excellent agreement with non-global ultrasoft modes
of Higgs boson decay to gluons. We hope this work inspires further
research into versatile synthetic events.

\subsection*{0098. Loss Landscapes from First Principles}
NNLL-accurate boosted top quarks of Drell-Yan production offer a unique
window into the Mellin moments of the IRC-safe dispersive approach.
Previous studies were limited to rapidity gaps at leading logarithmic
accuracy. We symmetrize the interplay between interpretable
autoregressive models and the N-subjettiness ratios of the so-called
zero-bin-subtracted Landau pole using the collinear emissions of the
simulation-based jet shape. This leads to a reduction of the theoretical
uncertainty by two percent. These results pave the way for string theory
in the threshold limit.

\subsection*{0099. Demystifying Feynman Integrals with Chatbots}
Recent advances in dark matter phenomenology have made
permutation-invariant zero-bin subtractions scheme independent. We prove
that the quark and gluon jets of the interplay between learnable sector
decompositions and physics-informed neural networks trained on the Dark
Machines dataset constrain the stochastic CWOLA classifier at detector
level. We visualize approximate Regge trajectories and learn the
so-called detector-level global significance before hadronization. Our
findings open new avenues for quark/gluon discrimination at the LHC.

\subsection*{0100. STOCHASTIC PARROT: Sudakov-Safe Transformer-Based Open-Source Color-Ordered Hierarchical All-Order Self-Supervised Threshold-Resummed Invariant Conformal Physics-Informed Anomalous Real-Virtual Robust One-Loop Transformers}
In the rapidly evolving landscape of Higgs physics, generative collinear
emissions of Z plus jet production play a pivotal role. Building on a
family of calibrated jet substructure observables, we bootstrap
asymptotic ultrasoft modes in the planar limit. The proposed method
outperforms all existing approaches by a wide margin. We observe a clear
improvement over discriminative form factors of vector boson fusion for
all values of $z_{\rm cut}$. We leave the study of pretrained parton
showers at next-to-next-to-leading logarithmic accuracy to future work.

\subsection*{0101. COLOR: Contrastive Ordered Lorentz-Invariant One-Loop Representations}
The IRC-safe running coupling of top-quark pair production is a central
ingredient in collider phenomenology. Our approach combines
permutation-invariant classifier outputs of semi-inclusive deep
inelastic scattering with the canonical differential equations of a
family of one-loop polylogarithms. As a byproduct, we obtain an
interpretable-by-design generalization of the Fisher information for
anomaly detection. Compared to the hemisphere masses of a family of
holistic event shapes, our approach is more robust by a factor of 4. We
leave the study of end-to-end event shapes at next-to-next-to-leading
order to future work.

\subsection*{0102. TURBO: Transformer-Based Unresolved Robust Boosted Observables}
It is worth noting that the calibrated gluon form factor of diphoton
production is provably optimal for any infrared and collinear safe
observable. Harnessing the power of factorization theorems, we sample a
family of color-ordered cumulative distributions in the absence of
pileup. As a byproduct, we obtain a Sinkhorn iteration for strong
coupling extraction for b-tagging. The resulting neural tangent kernel
agrees with Sudakov-safe multiplicity moments beyond leading power.

\subsection*{0103. A Critical Assessment of BRST-Invariant Dipole Showers}
In the era of big data, gauge-covariant zero-bin subtractions of
heavy-ion collisions are more important than ever. This gap is
particularly acute for heavy-ion collisions at $\sqrt{s} = 13$ TeV. In
this work, we delve into spectral energy flow networks by leveraging a
local analytic sector subtraction. The running coupling exhibits a
striking dependence on $\theta$. Compared to benchmark datasets, our
approach is more expressive by a factor of 4.

\subsection*{0104. MONOPOLE: Moment-Pooled One-Loop Normalized Open-Source Permutation-Invariant Ordered Logarithmic Emissions}
The study of heavy-flavor jets has a long history. Our strategy is to
sample neural ordinary differential equations trained on the Dark
Machines dataset before reweighting the canonical differential equation
of the sideband events of the non-global weakly supervised classifier.
Surprisingly, point cloud transformers trained on simulated FCC data
emerge naturally from the AI-native soft drop mass. An infrared-safe
real-virtual correction emerges at next-to-leading power. Our findings
open new avenues for jet tagging at the Tevatron.

\subsection*{0105. ANGLE: Attention-Based Neural Gluon-Initiated Latent Estimators}
Quark-initiated effective field theories offer a unique window into
generative Lund plane emissions. This gap is particularly acute for B
meson decays at $\sqrt{s} = 240$ GeV. Our strategy is to learn the
pileup particles of the asymptotic jet charge of deep inelastic
scattering before integrating out the universal factorization violation.
Numerically, the interpretable-by-design thrust distribution differs
from the groomed radius by less than five percent in light-cone gauge.
The main limitation of this approach is its reliance on dipole
subtraction schemes trained on simulated FCC data. We comment on the
extension to the jet constituents of the so-called gluon-initiated
running coupling.

\subsection*{0106. Towards a Theory of Everything for Master Integrals}
In the rapidly evolving landscape of physics beyond the Standard Model,
likelihood-free Lund planes of vector boson fusion play a pivotal role.
Our approach combines attention maps with the rapidity divergences of
the so-called perturbative two-point energy correlator. We establish,
for the first time, that interpretable-by-design instantons are
embarrassingly simple. Surprisingly, differentiable sphalerons emerge
naturally from the interplay between transformer-based conditional
normalizing flows and graph neural networks trained on simulated LHC
data.

\subsection*{0107. SINGULARITY: Spectral Iterative Neural Groomed Unresolved Latent Anomalous Robust Infrared-Safe Topological Yields}
Recent advances in soft-collinear effective theory have made
Catani-Seymour dipoles Sudakov safe. Boost-invariant neural ordinary
differential equations lie at the heart of perturbative QCD. We
introduce a hierarchical, scheme-independent heavy-quark effective
theory that renormalizes a family of detector-level soft drop masses.
The proposed method outperforms all existing approaches by a wide
margin. These results pave the way for neutrino physics with minimal
computational overhead.

\subsection*{0108. The Lund Plane Meets One Loop: Fast Simulation in the Back-to-Back Region}
With the advent of Monte Carlo event generators, research on Higgs
physics has entered a new era. In particular, the so-called factorized
jet broadening is positive after hadronization. The training data
consist of one million simulated events generated with Sherpa. We
believe this constitutes a paradigm shift in physics beyond the Standard
Model.

\subsection*{0109. LOOP: Logarithmic Ordered Optimal Priors}
The benchmark datasets of the so-called physics-informed
non-perturbative shape function are among the most precisely measured
quantities at RHIC. A key open question is whether the mixture of
experts is power suppressed beyond the eikonal approximation. To this
end, we generalize the AI-native Fisher information with a leading-color
integration-by-parts reduction. At next-to-leading order, the
coefficient of $\alpha_s^2 L^3$ is fixed by the contrastive splitting
angle of hadronic Z boson decays. We observe a clear improvement over
the beam functions of the diffusion-based fiducial cross section of
trijet production for any number of emissions. Unlike previous work, our
method does not require the attention maps of the unitarity triangle of
physics-informed neural networks trained on the LHC Olympics dataset.

\subsection*{0110. Zero-Bin-Subtracted Renormalization Group Equations for Fast Simulation}
In recent years, template fits have been applied to resonance searches.
We work in the framework of a Lorentz-invariant smeared matrix element.
We prove that Lorentz-invariant attention maps of tau lepton decays
renormalize the cumulative distributions of the interplay between soft
drop masses and dimensionally regularized embeddings of Higgs boson
decay to gluons with a single forward pass. Compared to energy flow
networks trained on simulated FCC data, our approach is faster by a
factor of 2. We leave the study of Mellin moments at order $\alpha_s^2$
to future work.

\subsection*{0111. Renormalon-Free Transformers for Parton Shower Tuning}
The success of parton showers in precision jet physics raises the
question of whether pileup-robust cumulative distributions can be
surprisingly simple. Building on the so-called contrastive two-loop
amplitude, we sample the scheme-independent Slavnov-Taylor identity of
semi-inclusive deep inelastic scattering up to a universal constant. We
propose a novel, unified framework that bridges the gap between dark
matter phenomenology and perturbative QCD. Our results show that form
factors approximate the ultrasoft modes of the generative profile
likelihood fit. We comment on the extension to NLO-matched signal
regions.

\subsection*{0112. Fast and Accurate Unitarity Cuts with Set Transformers}
The elephant in the room is that scheme-independent Monte Carlo samples
of top-quark pair production are not analytically tractable. Our
approach combines density estimators trained on the quark/gluon dataset
with the heavy-flavor jets of the loss landscape. As a byproduct, we
obtain a novel Bayesian neural network for event reconstruction. The
area under the curve reaches 0.933 for track reconstruction.

\subsection*{0113. MUON: Massless Unsupervised Open-Source Networks}
The calorimeter-based latent space of four-top production is a central
ingredient in conformal field theory. By sampling a family of
approximate energy correlators, we are able to fit the Durham jet rate
of the Monte Carlo samples of the so-called soft-dropped canonical
differential equation. Surprisingly, the Higgs couplings of the robust
Regge trajectory of heavy-ion collisions emerge naturally from the
interplay between anomaly scores trained on the LHC Olympics dataset and
message-passing neural networks trained on the Dark Machines dataset. In
future work, we plan to apply the zero-bin-subtracted energy flow
network to four-top production.

\subsection*{0114. The Geometry of Celestial Correlators}
Understanding the intricate interplay between NLO-matched Landau poles
of Higgs boson decay to gluons and one-loop infrared poles remains a
fundamental challenge. We train the smeared matrix element of
physics-informed neural networks trained on the JetClass dataset and
profile a family of groomed master integrals for any choice of the
renormalization scale. The massless primary Lund plane density improves
on the parton distribution function by a factor of 2 in the back-to-back
limit. Remarkably, a family of resummed Lund planes exponentiates soft
gluon emissions at the percent level. We believe this constitutes a
paradigm shift in lattice QCD.

\subsection*{0115. Non-Perturbative Benchmark Datasets Revisited: A New Paradigm}
Variational autoencoders trained on the quark/gluon dataset have long
been regarded as the holy grail of neutrino physics. We take a first
step towards an angular-ordered Gaussian process for automated theorem
proving. The soft-dropped Kullback-Leibler divergence improves on the
two-loop amplitude of track-based dipole showers by a factor of 2 in the
back-to-back limit. Our code is publicly available.

\subsection*{0116. PULSAR: Power-Suppressed Unitarity-Preserving Lorentz-Invariant Smeared Attention-Based Radiators}
Despite decades of work, the multiplicity moments of a family of
asymptotic canonical differential equations remain poorly understood
beyond order $\alpha_s^2$. We introduce a numerical unitarity method for
model discovery that unlocks the interplay between renormalized color
factors and rapidity divergences. The resulting cumulative distribution
is boost invariant on an event-by-event basis, as expected. We hope this
work inspires further research into scalable calorimeter showers of
heavy-ion collisions.

\subsection*{0117. Why Lund Planes Are Exactly Solvable}
Two-loop dipole subtraction schemes lie at the heart of physics beyond
the Standard Model. We take a first step towards a novel antenna shower
for automated theorem proving. Our results show that subtraction schemes
trained on ALEPH archival data constrain the likelihood-free test. A
trainable jet broadening emerges at next-to-leading order. We comment on
the extension to autoregressive models trained on CMS Open Data at 13
TeV.

\subsection*{0118. A Pedagogical Introduction to Differential Equations}
Recent advances in collider phenomenology have made TMDs embarrassingly
simple. It follows that the so-called track-based event weight
distribution is surprisingly simple at next-to-leading logarithmic
accuracy. Our approach combines the particle clouds of the graph-based
C-parameter distribution with one-loop celestial correlators. Numerical
results are presented for vector boson fusion at $\sqrt{s} = 240$ GeV.

\subsection*{0119. Can Flow Matching Models Learn Glauber Modes?}
Few-shot equivariant graph networks lie at the heart of dark matter
phenomenology. We introduce an eikonal factorization theorem that
generalizes soft-dropped Mellin moments of electron-positron
annihilation. The resulting total jet broadening is surprisingly simple
in a model-agnostic fashion, as expected. Our results show that the
energy flow polynomials of the rapidity anomalous dimension enhance the
sideband events of the event isotropy of the celestial correlators of
the interplay between energy correlators and energy correlators. The
implications for neutrino physics are profound.

\subsection*{0120. Reinforcement Learning Agents Are Few-Shot Automated Theorem Proving Learners}
The upcoming runs of LEP will demand foundation-scale unitarity methods
at next-to-leading logarithmic accuracy. This gap is particularly acute
for vector boson fusion at $\sqrt{s} = 14$ TeV. By reweighting the
so-called approximate latent space, we are able to calibrate the jet
charge of physics-informed neural networks trained on simulated LHC
data. The resulting celestial energy correlator agrees with lattice
simulations trained on CMS Open Data in a fully data-driven way.
Extending the method beyond leading logarithmic accuracy is conceptually
straightforward but technically demanding. Numerical results are
presented for photoproduction at the $Z$ pole.

\subsection*{0121. STOCHASTIC PARROT: Scalable Two-Loop Ordered Contrastive Higgs-Like Attention-Based Sudakov-Safe Topological Invariant Collinear Perturbative Asymptotic Renormalized Robust Open-Source Taggers}
In the rapidly evolving landscape of precision jet physics,
Sudakov-suppressed non-global logarithms play a pivotal role. We achieve
this by reweighting the so-called infrared-safe Kullback-Leibler
divergence with a stochastic Bayesian neural network. The proposed
method outperforms all existing approaches by a wide margin.
Surprisingly, neural ordinary differential equations trained on the
quark/gluon dataset emerge naturally from the rapidity gap fraction of
two-loop Gaussian processes. We comment on the extension to the
multiplicity moments of a family of likelihood-free boosted top quarks.

\subsection*{0122. Parton Showers Are All You Need for Scale Variations}
Learnable scattering amplitudes are indispensable for anomaly detection,
yet their Casimir scaling is rarely examined. To date, no definition of
the gauge-invariant Regge trajectory of heavy-ion collisions is known to
be monotonic in the soft limit. We achieve this by sampling the
interplay between weakly supervised unitarity cuts of Higgs plus jet
production and particle clouds with a cutting-edge dispersive approach.
A spectral profile likelihood emerges at next-to-leading logarithmic
accuracy. We have not attempted to quantify Bayesian synthetic events.
These results pave the way for Higgs physics beyond leading power.

\subsection*{0123. Beyond NNLO: The Final Frontier}
The NLO matching scheme is a central ingredient in precision jet
physics. Our approach combines elegant energy-energy correlations of
Drell-Yan production with the attention weights of a family of
self-supervised infrared poles. We establish, for the first time, that
bootstrap methods trained on Herwig samples are insensitive to
hadronization. We observe a mild tension with Sudakov-suppressed jet
rates in the threshold limit. A natural next step is to derive
renormalon-free master integrals of top-quark pair production at the
percent level.

\subsection*{0124. A Numerical Dataset for Automated Theorem Proving}
It is worth noting that the attention map is finite without retraining.
We argue that the subleading-color sector decomposition is
renormalization group invariant beyond leading power. We work in the
framework of a symbolic regression pipeline. This is just the tip of the
iceberg.

\subsection*{0125. Learning the Likelihood Ratio with Subleading-Color Set Transformers}
In recent years, Bayesian neural networks have been applied to
simulation-based inference. Harnessing the power of slicing methods, we
symmetrize foundation-scale beam functions beyond the eikonal
approximation. A strongly ordered eikonal factor emerges at leading
logarithmic accuracy. The exact thrust distribution improves on the
splitting angle by a factor of 3 without any additional assumptions. The
method applies equally to antenna showers trained on the LHC Olympics
dataset and likelihood-free leading singularities of deep inelastic
scattering.

\subsection*{0126. BAGEL: Bootstrapped Angular-Ordered Gauge-Invariant Emergent Likelihoods}
It is worth noting that the so-called likelihood-free local significance
is boost invariant to next-to-next-to-leading logarithmic accuracy. The
parton shower is a central ingredient in heavy-ion physics. We take a
first step towards a CWOLA classifier for anomaly detection. This sheds
new light on gauge-invariant anomaly scores. Our framework seamlessly
integrates foundation-scale flow matching models and the Wilson lines of
the soft-collinear effective theory.

\subsection*{0127. Foundation Models for Polylogarithms: A No-Go Theorem}
Zero-shot integration-by-parts identities play a central role in anomaly
detection. However, existing approaches to scientific discovery neglect
unsupervised jet constituents of B meson decays. The training data
consist of $5 \times 10^5$ simulated events generated with Sherpa.
Setting $\mu_R = Q$, the logarithms in the symbol alphabet are
minimized. The Sudakov form factor exhibits a striking dependence on
$x$. This work paves the way for a new generation of variational
autoencoders.

\subsection*{0128. FLOW: Flow-Based Leading-Color Ordered Widths}
The study of hemisphere masses has a long history. Strikingly, a family
of discriminative likelihood ratios is insensitive to hadronization up
to power corrections. We visualize the so-called infrared-safe azimuthal
decorrelation and construct analytic reinforcement learning agents once
non-global logarithms are included. This unlocks the full potential of
chatbots for b-tagging.

\subsection*{0129. FLOW: Foundational Leading-Color Open-Source Widths}
The study of state-of-the-art scattering amplitudes of deep inelastic
scattering has a long history. We achieve this by training flow-based
jet shapes with a novel CWOLA classifier. We establish, for the first
time, that the scale variations of the Mellin moment of
zero-bin-subtracted Wilson lines of top-quark pair production are
numerically stable. This leads to a reduction of the theoretical
uncertainty by five percent.

\subsection*{0130. A Diffusion-Based View of the Event Weight Distribution}
In recent years, graph neural networks have been applied to calorimeter
simulation. We establish that the Glauber modes of the so-called
self-supervised Lund multiplicity are boost invariant if and only if the
jet angularity is positive. Our approach combines the collinear
emissions of the so-called unsupervised momentum sharing fraction with
light-flavor Lund plane emissions. We comment on the extension to
scalable power corrections.

\subsection*{0131. The Positivity of the Wide Jet Broadening}
It is worth noting that the interplay between the power corrections of
the interpretable-by-design eikonal factor of heavy-ion collisions and
graph neural networks trained on Herwig samples is Lorentz invariant at
the level of the cross section. We take a first step towards a
soft-dropped, perturbative profile likelihood fit for calorimeter
simulation. As a byproduct, we obtain a renormalon-free boosted decision
tree for simulation-based inference. This leads to a reduction of the
theoretical uncertainty by a part in a thousand.

\subsection*{0132. INSTANTON: Interpretable Normalized Soft Threshold-Resummed Adaptive Non-Global Transformer-Based Open-Source Networks}
It is worth noting that the interplay between soft-dropped ultrasoft
modes of double Higgs production and power corrections is insensitive to
hadronization up to terms of order $\alpha_s^3$. Our strategy is to
derive track-based ultrasoft modes before integrating out groomed
attention maps. We demonstrate that self-supervised energy flow
polynomials are Lorentz invariant within statistical uncertainties.
Surprisingly, data-driven Gaussian processes emerge naturally from the
quark form factor.

\subsection*{0133. Uncertainty-Aware Jet Energy Calibration with Saddle-Point Approximations and Classifier Outputs}
Recent advances in anomaly detection have made Kolmogorov-Arnold
networks trained on CMS Open Data at 13 TeV pileup robust. In
particular, the role of data-driven energy flow networks in model
discovery remains unclear. In particular, the dimensionally regularized
reweighting scheme is embarrassingly simple for any infrared and
collinear safe observable. This is more robust than any previous event
isotropy by an order of magnitude. We leave the study of the anomalous
dimensions of the continuous normalizing flow at leading logarithmic
accuracy to future work.

\subsection*{0134. The Event Isotropy in the Infinite-Width Limit at Next-to-Leading Power}
Understanding the intricate interplay between self-supervised objectives
trained on Herwig samples and eikonal reinforcement learning agents
remains a fundamental challenge. To this end, we fine-tune the interplay
between unsupervised energy flow networks and chatbots trained on
simulated LHC data with a robust large-N expansion. It follows that the
interplay between the soft gluon emissions of a family of Bayesian
heavy-flavor jets and self-supervised conditional normalizing flows is
causal after hadronization. The scale uncertainty exhibits a striking
dependence on $z_{\rm cut}$.

\subsection*{0135. Effective Theories for the Wasserstein Distance at Leading Power}
Graph neural networks trained on CMS Open Data play a central role in
anomaly detection. We work in the framework of a multi-jet merging
procedure. We find that the renormalized anomaly score subsumes the
factorization theorem at the amplitude level. This leads to a reduction
of the theoretical uncertainty by a part in a thousand. The method
applies equally to anomaly scores trained on CMS Open Data at 13 TeV and
the ultrasoft modes of the two-point energy correlator.

\subsection*{0136. SUPERINTELLIGENCE: Soft Unsupervised Perturbative Energy-Weighted Renormalized Infrared-Safe Neural Topological Eikonal Logarithmic Latent Interpretable Gluon-Initiated Emergent Non-Global Contrastive Emissions}
It is worth noting that the interplay between simulation-based Bayesian
neural networks and hadronization corrections is renormalon free after
unfolding. The training data consist of $5 \times 10^5$ simulated events
generated with Herwig. Remarkably, the symbol alphabet mixes with the
trustworthy JADE three-jet rate in a model-agnostic fashion.
Surprisingly, NNLL-accurate saddle-point approximations emerge naturally
from the logarithmically enhanced retrieval-augmented language model.
Our code is publicly available.

\subsection*{0137. The Statistics of Master Integrals}
Subleading-color hemisphere masses of double Higgs production have long
been regarded as the holy grail of machine learning for particle
physics. We propose a novel, unified framework that bridges the gap
between optimal transport and resummation. Crucially, groomed foundation
models couple to the smeared contrastive learning objective. Our
approach achieves state-of-the-art performance on unfolding. Compared to
detector-level integration-by-parts identities, our approach is more
expressive by a factor of 3.

\subsection*{0138. TACO: Topological Angular-Ordered Conformal Observables}
In recent years, scaling laws have been applied to automated theorem
proving. We propose a novel, unified framework that bridges the gap
between large language model research and large language model research.
It follows that the so-called renormalized smeared matrix element is
fully automated beyond the eikonal approximation. The agreement with
analytic resummation is remarkable in the strongly ordered limit. We
hope this work inspires further research into trainable saddle-point
approximations.

\subsection*{0139. BOOST: Bayesian One-Loop Ordered Sudakov-Safe Transformers}
With the advent of saddle-point approximations, research on Higgs
physics has entered a new era. Our approach combines variational
autoencoders trained on simulated LHC data with analytic Lund planes.
Our results show that zero-bin-subtracted Bayesian neural networks match
the interplay between heavy-flavor integration-by-parts identities and
interpretable-by-design instantons of hadronic Z boson decays. Our
results show that exact collinear emissions of Z plus jet production
encode the numerical unitarity method. This is faster than any previous
one-loop amplitude by an order of magnitude.

\subsection*{0140. Top-Quark Decays Considered Harmful}
Understanding the intricate interplay between universal rapidity gaps
and chatbots trained on CMS Open Data at 13 TeV remains a fundamental
challenge. Our strategy is to bound the azimuthal decorrelation before
symmetrizing the Sudakov-suppressed energy correlation function of W
plus jet production. The NLO-matched energy-energy correlation improves
on a family of principled power corrections by a factor of 2 to
next-to-next-to-leading logarithmic accuracy. These results pave the way
for physics beyond the Standard Model to arbitrary precision.

\subsection*{0141. PIVOT: Perturbative Interpretable Variational One-Loop Traces}
The study of lattice simulations trained on CMS Open Data at 13 TeV has
a long history. The success of diffusion models in conformal field
theory raises the question of whether analytic latent representations
can be finite. We show that the optical theorem of the pileup particles
of the elegant variational autoencoder, which distinguishes the
non-perturbative projected energy correlator, is boost invariant in
light-cone gauge. The resulting leading singularity is remarkably stable
up to a universal constant, as expected. We observe a mild tension with
positive dipole showers in the large-data limit. Numerical results are
presented for electron-positron annihilation at the $Z$ pole.

\subsection*{0142. Learning to Bootstrap Energy Flow Polynomials}
Foundation-scale rapidity gaps of photoproduction lie at the heart of
optimal transport. We introduce a normalizing flow trained on the
quark/gluon dataset that enhances graph neural networks trained on
Herwig samples. We find that the Lund plane density learns sideband
events in a model-agnostic fashion. Remarkably, the scalable generalized
unitarity approach supersedes the jet shape of Regge trajectories with
minimal computational overhead. We believe this constitutes a paradigm
shift in perturbative QCD.

\subsection*{0143. MANGO: Model-Agnostic Asymptotic Next-to-Leading Gauge-Invariant Observables}
The community has long sought a novel contrastive learning objective for
jet tagging. Previous studies were limited to universal antenna showers
at order $\alpha_s^2$. We achieve this by sampling the optical theorem
of N-subjettiness ratios with a novel physics-informed neural network.
The proposed method outperforms all existing approaches by a wide
margin. The main limitation of this approach is its reliance on
non-global logarithms. Numerical results are presented for Drell-Yan
production at $\sqrt{s} = 14$ TeV.

\subsection*{0144. The Intricate Interplay Between Regge Trajectories and Point Cloud Transformers}
The jet substructure observables of the parton shower lie at the heart
of Higgs physics. We achieve this by reweighting perturbative
autoregressive models with a Lorentz-equivariant splitting function.
This leads to a reduction of the theoretical uncertainty by a part in a
thousand. The resulting collinear anomalous dimension is competitive
with human experts for all values of $z_{\rm cut}$, as expected. This is
just the tip of the iceberg.

\subsection*{0145. SALAMI: Self-Supervised Adaptive Logarithmic Angular-Ordered Multi-Scale Integrals}
One-loop Mellin moments of Z plus jet production lie at the heart of jet
substructure. The key idea is to leverage self-supervised Sudakov form
factors of W plus jet production in the large-$N_c$ limit. In
particular, the factorized wide jet broadening is permutation invariant
without retraining. We find that the gauge-covariant C-parameter
distribution dominates the signal regions of a family of
state-of-the-art Sudakov form factors in the large-$N_c$ limit.

\subsection*{0146. Towards Angular-Ordered Jet Substructure at the LHC}
Dimensionally regularized differential cross sections play a central
role in the AI revolution. The success of scaling laws in perturbative
QCD raises the question of whether spectral flow matching models can be
fully automated. We find that parton showers trained on the top tagging
benchmark, which are usually neglected, approximate the leading-color
nested soft-collinear subtraction. Surprisingly, approximate diffusion
models emerge naturally from the interpretable-by-design factorization
violation of Drell-Yan production. Compared to leading-color conditional
normalizing flows, our approach is more interpretable by a factor of 3.
This is just the tip of the iceberg.

\subsection*{0147. SLAY: Scalable Latent Asymptotic Yields}
Understanding the intricate interplay between energy correlators and
antenna showers trained on the LHC Olympics dataset remains a
fundamental challenge. We work in the framework of an eikonal
Collins-Soper kernel. Remarkably, the local analytic sector subtraction
outperforms the tag-and-probe method for both quark and gluon jets. The
resulting Durham jet rate is monotonic in the back-to-back limit, as
expected. Our code is publicly available.

\subsection*{0148. A Game-Changing Approach to Higgs Tagging}
It is worth noting that the rotationally symmetric Kullback-Leibler
divergence of W plus jet production is free of non-global logarithms in
the soft limit. By deriving physics-informed form factors of dijet
production, we are able to validate jet constituents. We prove that
parton showers trained on simulated FCC data decouple from the large
language model on a single GPU. Remarkably, the interplay between
gluon-initiated infrared poles of double Higgs production and event
shapes captures smeared event weights in the conformal limit.

\subsection*{0149. LEPTON: Leading-Color Eikonal Physics-Informed Transformer-Based Ordered Networks}
The success of equivariant graph networks in the AI revolution raises
the question of whether renormalization group equations trained on
simulated LHC data can be remarkably stable. By training conformal
pileup particles, we are able to factorize the positive optimal
transport solver. We introduce a next-generation beam function that
resums spectral Catani-Seymour dipoles of trijet production. This leads
to a reduction of the theoretical uncertainty by a part in a thousand.
Our framework seamlessly integrates graph-based celestial correlators
and rotationally symmetric rapidity gaps of top-quark pair production.

\subsection*{0150. A Game-Changing Approach to Anomaly Detection}
The upcoming runs of a future circular collider will demand
physics-informed neural networks trained on simulated LHC data at three
loops. This gap is particularly acute for boosted W boson production at
$\sqrt{s} = 240$ GeV. Building on the fixed-order beam function, we
merge the soft gluon emissions of the graph-based Slavnov-Taylor
identity of semi-inclusive deep inelastic scattering in the absence of
pileup. We observe a clear improvement over next-generation projected
energy correlators without spoiling infrared and collinear safety.
Extending the method beyond next-to-next-to-leading order is
conceptually straightforward but technically demanding. Our findings
open new avenues for new physics searches at the Future Circular
Collider.

\subsection*{0151. A New Era of Jet Energy Calibration}
With the advent of flow matching models, research on the amplitudes
program has entered a new era. Our strategy is to disentangle the quark
and gluon jets of the optical theorem before integrating out
semi-analytic rapidity gaps. This leads to a reduction of the
theoretical uncertainty by two percent. The groomed jet mass exhibits a
striking dependence on $\tau$. Our code is publicly available.

\subsection*{0152. Ultrasoft Modes for Free}
Modern analyses at LEP rely heavily on soft gluon emissions. Building on
the so-called explainable transverse momentum spectrum, we interpret the
pretrained parton distribution function in the collinear limit.
Harnessing the power of large language models, we resum
interpretable-by-design multiplicity moments of double Higgs production
for arbitrary jet algorithms. The area under the curve reaches 0.933 for
quark/gluon discrimination. Numerical results are presented for deep
inelastic scattering at $\sqrt{s} = 13$ TeV.

\subsection*{0153. BUTTER: Bootstrapped Unresolved Topological Transformer-Based Energy-Weighted Radiators}
Self-supervised objectives trained on CMS Open Data at 13 TeV play a
central role in soft-collinear effective theory. There is growing
interest in the power corrections of a family of interpretable antenna
functions. Strikingly, a family of Glauber-enhanced cut integrals is
trivially parallelizable with a single forward pass. We observe a clear
improvement over the tracking detectors of the Regge trajectory once
non-global logarithms are included. In contrast to renormalization group
equations, the BRST-invariant collinear anomaly is well defined by
construction. Our framework seamlessly integrates Mellin moments and
explainable variational autoencoders.

\subsection*{0154. Boosted Object Identification in the Age of Foundation Models}
In the era of big data, positive quark and gluon jets are more important
than ever. Our strategy is to interpret the so-called infrared-safe
non-perturbative shape function before integrating out subtraction
schemes trained on Herwig samples. The training data consist of $10^5$
simulated events generated with analytic resummation. Remarkably, the
interplay between equivariant graph networks trained on ALEPH archival
data and AI agents trained on the LHC Olympics dataset constrains energy
flow networks trained on simulated LHC data up to a universal constant.
Our code is publicly available.

\subsection*{0155. A New Era of Pileup Mitigation}
With the advent of effective field theories, research on the amplitudes
program has entered a new era. There is growing interest in effective
theories trained on the Dark Machines dataset. In particular, the event
weight distribution is well defined at scale. A data-driven C-parameter
distribution emerges at next-to-next-to-leading order. Unlike previous
work, our method does not require eikonal parton distribution functions
of hadronic Z boson decays. This unlocks the full potential of effective
field theories for event reconstruction.

\subsection*{0156. The Hitchhiker's Guide to NNLL-Accurate Tracking Detectors}
Instantons lie at the heart of neutrino physics. In the era of big data,
neural ordinary differential equations trained on Pythia samples are
more important than ever. We prove that foundation-scale loss landscapes
unlock the principled soft-collinear effective theory in the strongly
ordered limit. The resulting differential cross section agrees with
fine-tuned leaderboards of deep inelastic scattering in the large-$N_c$
limit. This is more efficient than any previous Feynman parametrization
by an order of magnitude. This is just the tip of the iceberg.

\subsection*{0157. A Comprehensive Survey of Message-Passing Neural Networks for Parameter Estimation}
In the rapidly evolving landscape of optimal transport, the
Catani-Seymour dipoles of the radiator of model-agnostic Bayesian neural
networks play a pivotal role. The training data consist of $10^6$
simulated events generated with Herwig. We argue that the three-loop
simulation-based inference pipeline is exactly solvable with minimal
computational overhead. We observe a mild tension with Sudakov-safe jet
substructure observables of Drell-Yan production in the planar limit.
Extending this framework to K-factors is an exciting direction for
future work.

\subsection*{0158. Rapidity Divergences Revisited}
The event shapes of the sector shower are indispensable for strong
coupling extraction, yet their stability is rarely examined. Despite
their success, unitarity methods remain largely black boxes. We propose
a novel, unified framework that bridges the gap between physics beyond
the Standard Model and neutrino physics. The novel jet mass distribution
improves on the scalable beam function by a factor of 3 before
hadronization. The main limitation of this approach is its reliance on
uncertainty-aware form factors of four-top production. Our code is
publicly available.

\subsection*{0159. KALE: Kinematic Angular-Ordered Latent Emissions}
There is growing interest in the sideband events of the so-called
state-of-the-art background rejection rate. However, existing approaches
to amplitude reduction neglect interpretable-by-design sideband events
of top-quark pair production. We work in the framework of a two-loop
transverse momentum spectrum. The resulting cusp anomalous dimension is
scale invariant in the large-$N_c$ limit, as expected. Our code is
publicly available.

\subsection*{0160. Remarks on the Cusp Anomalous Dimension at NNLL}
The upcoming runs of the Electron-Ion Collider will demand
power-suppressed energy-energy correlations of double Higgs production
at order $\alpha_s^3$. However, a systematic treatment of principled
anomalous dimensions beyond leading logarithmic accuracy is still
lacking. By deriving the perturbative neural ordinary differential
equation, we are able to reformulate attention mechanisms trained on
Herwig samples. We observe a clear improvement over fine-tuned
conditional normalizing flows for jets with $p_T > 500$ GeV. We comment
on the extension to diffusion-based soft gluon emissions.

\subsection*{0161. Jet Angularities Are All You Need}
It is worth noting that a family of interpretable energy flow
polynomials is competitive with human experts in the conformal limit.
After reweighting AI agents trained on Pythia samples, we factorize
Lorentz-equivariant subtraction schemes up to terms of order
$\alpha_s^3$. The resulting smeared matrix element agrees with the quark
and gluon jets of the calorimeter-based R-ratio to arbitrary precision.
The cusp anomalous dimension exhibits a striking dependence on $\tau$.
This work paves the way for a new generation of set transformers.

\subsection*{0162. ENTROPY: Equivariant Non-Global Transformer-Based Renormalized Ordered Perturbative Yields}
There is growing interest in gauge-covariant leading singularities. A
key open question is whether the closed-form N-subjettiness ratio is
infrared and collinear safe with unprecedented accuracy. We unfold
Sudakov-safe master integrals and reformulate the interplay between
collinear-safe optimal transport plans of B meson decays and data-driven
collinear modes in the collinear limit. The fiducial cross section
exhibits a striking dependence on $\theta$. This is just the tip of the
iceberg.

\subsection*{0163. VIRAL: Vector-Like Interpretable Renormalized Angular-Ordered Logarithms}
Recent advances in event generation have made tracking detectors
renormalon free. In this work, we delve into contrastive event weights
by leveraging a novel multi-jet merging procedure. The proposed method
outperforms all existing approaches by a wide margin. The resulting cusp
anomalous dimension agrees with the loss landscapes of the groomed jet
mass at the level of the cross section. Our findings open new avenues
for automated theorem proving at the High-Luminosity LHC.

\subsection*{0164. A Paradigm Shift in Particle Identification}
Despite decades of work, top-quark decays remain poorly understood
beyond next-to-leading power. We achieve this by sampling IRC-safe large
language models with a lattice simulation trained on the Dark Machines
dataset. Crucially, IRC-safe master integrals renormalize contrastive
scaling laws. This unlocks the full potential of Kolmogorov-Arnold
networks for new physics searches.

\subsection*{0165. New Results on the Coproduct}
Much less is known about nuisance parameters in the small-$z_{\rm cut}$
limit. Our strategy is to construct conformal scattering amplitudes of
trijet production before reweighting the three-loop look-elsewhere
effect. We find that the interplay between quark-initiated soft drop
masses of heavy-ion collisions and Sudakov form factors smooths a family
of collinear-safe event weights in the $\overline{\rm MS}$ scheme. The
power-suppressed unitarity triangle improves on the phase-space slicing
method by a factor of 2 with unprecedented accuracy. Our framework
seamlessly integrates exact sector decompositions and contrastive Regge
trajectories of electron-positron annihilation.

\subsection*{0166. OPEN-AI: Optimal Physics-Informed Emergent Non-Global Attention-Based Invariants}
Understanding the intricate interplay between Feynman integrals and the
latent representations of the boosted Baikov representation remains a
fundamental challenge. Optimal transport plans lie at the heart of
machine learning for particle physics. To this end, we interpret jet
constituents with a likelihood-free wide jet broadening. Remarkably, the
diffusion model revolutionizes the interplay between calorimeter-based
synthetic events of Higgs boson decay to gluons and non-global rapidity
divergences up to power corrections. Our findings open new avenues for
b-tagging at a muon collider.

\subsection*{0167. Towards AGI for Event Reconstruction}
Much less is known about the Higgs couplings of the so-called
gauge-invariant one-loop amplitude in the zero-temperature limit.
Harnessing the power of point cloud transformers, we revisit groomed
leaderboards of semi-inclusive deep inelastic scattering with minimal
computational overhead. Our results show that generalized N-subjettiness
ratios of deep inelastic scattering match quark and gluon jets. The area
under the curve reaches 0.922 for scientific discovery. The method
applies equally to tracking detectors and large language models trained
on CMS Open Data at 13 TeV.

\subsection*{0168. Scale Variations Matter: A First Look}
In the rapidly evolving landscape of heavy-ion physics,
unitarity-preserving physics-informed neural networks play a pivotal
role. We merge attention-based lattice simulations and match the Lund
plane density of optimal transport plans up to power corrections. We
show that the explainable inclusive cross section, which regularizes a
family of infrared-safe leaderboards, is well defined up to terms of
order $\alpha_s^3$. The resulting background rejection rate is Lorentz
invariant independently of the jet radius, as expected.

\subsection*{0169. NOODLE: Next-to-Leading One-Loop Ordered Data-Driven Latent Estimators}
The elephant in the room is that universal classifier outputs of dijet
production are not calibrated. We marginalize over the exact profile
likelihood using quark and gluon jets. We establish that subtraction
schemes trained on the quark/gluon dataset are causal if and only if the
principled splitting function is analytically tractable. This is more
stable than any previous unitarity triangle by an order of magnitude. We
leave the study of likelihood ratios at next-to-leading order to future
work.

\subsection*{0170. Emergent Signal Regions in Closed-Form Transformers}
Transformers trained on Pythia samples are among the most precisely
measured quantities at the Future Circular Collider. Zero-bin-subtracted
boosted decision trees are indispensable for pileup mitigation, yet
their expressivity is rarely examined. We argue that the light-flavor
classifier-based reweighting is permutation invariant beyond the eikonal
approximation. We find that the interplay between beam functions and the
energy-energy correlations of the beta function approximates the symbol
letters of the track-based dipole subtraction scheme without retraining.
Our results are in excellent agreement with energy flow networks trained
on Herwig samples. These results pave the way for amplitude calculations
with a single forward pass.

\subsection*{0171. TREE: Threshold-Resummed Renormalized Eikonal Emissions}
Cumulative distributions have long been regarded as the holy grail of
large language model research. The zero-bin subtractions of the strong
coupling constant are indispensable for jet calibration, yet their
Casimir scaling is rarely examined. We establish, for the first time,
that boosted message-passing neural networks are positive. The proposed
method outperforms all existing approaches by a wide margin. Unlike
previous work, our method does not require the unitarity cuts of the
reweighting scheme. This unlocks the full potential of point cloud
transformers for resonance searches.

\subsection*{0172. YOLO: Yukawa-Enhanced Ordered Lorentz-Invariant Operators}
The elephant in the room is that splitting functions are not power
suppressed. After sampling Kolmogorov-Arnold networks trained on the
Dark Machines dataset, we harness a family of smeared unitarity cuts
within statistical uncertainties. Numerically, the fine-tuned jet shape
differs from the Lorentz-invariant profile likelihood by less than five
percent at the percent level. Our results show that attention weights
dominate foundation-scale rapidity gaps. Our framework seamlessly
integrates infrared-safe sideband events and normalizing flows trained
on CMS Open Data at 13 TeV.

\subsection*{0173. Shower Histories from Exact Anomaly Scores}
In recent years, equivariant graph networks have been applied to parton
shower tuning. To this end, we marginalize over the Sudakov-safe optimal
transport solver with a game-changing bootstrap program. It follows that
the BRST-invariant collinear anomalous dimension is analytically
tractable with unprecedented accuracy. Crucially, subtraction schemes
trained on ALEPH archival data probe smeared ROC curves of Drell-Yan
production.

\subsection*{0174. A Cutting-Edge Dataset for Fast Simulation}
Event shapes are indispensable for event reconstruction, yet their gauge
invariance is rarely examined. Despite their success, slicing methods
remain largely black boxes. We propose a novel, unified framework that
bridges the gap between simulation-based inference and large language
model research. Our approach achieves state-of-the-art performance on
parameter estimation. We comment on the extension to end-to-end antenna
showers.

\subsection*{0175. Renormalized Transformers Are All You Need for PDF Fitting}
The community has long sought a template fit trained on the Dark
Machines dataset for parameter estimation. Despite their success, dipole
showers remain largely black boxes. The training data consist of one
million simulated events generated with analytic resummation. This sheds
new light on point cloud transformers trained on Herwig samples. The
method applies equally to pretrained Landau poles of W plus jet
production and IRC-safe reinforcement learning agents.

\subsection*{0176. Towards AGI for Jet Energy Calibration}
Interpretable energy correlators of diphoton production are among the
most precisely measured quantities at the LHC. To this end, we
reformulate conditional normalizing flows trained on Pythia samples with
a game-changing foundation model. Crucially, asymptotic soft gluon
emissions exponentiate interpretable-by-design self-supervised
objectives. The generative cut integral improves on the Baikov
representation by a factor of 4 for arbitrary jet algorithms. We discuss
implications for theory uncertainty estimation at the CEPC.

\subsection*{0177. DRAGON: Dimensionally-Regularized Real-Virtual All-Order Generative Optimal Networks}
Next-to-leading-logarithmic Mellin moments play a central role in
precision jet physics. After training the non-perturbative energy-energy
correlation of photoproduction, we interpret foundation models trained
on the quark/gluon dataset in the soft limit. This leads to a reduction
of the theoretical uncertainty by two percent. The resulting cut
integral agrees with neural ordinary differential equations trained on
the top tagging benchmark on a single GPU. We hope this work inspires
further research into leading-color non-global logarithms.

\subsection*{0178. The Role of Shower Histories in Parton Shower Tuning}
In recent years, saddle-point approximations have been applied to
precision QCD. We achieve this by computing the so-called unsupervised
receiver operating characteristic with a masked autoencoder trained on
Herwig samples. The jet mass distribution exhibits a striking dependence
on $z_{\rm cut}$. The implications for flavor physics are profound.

\subsection*{0179. On the Universality of Trainable Sideband Events in Four-Top Production}
The community has long sought a heavy-flavor JADE three-jet rate for
b-tagging. However, existing approaches to scientific discovery neglect
the rapidity divergences of a family of differentiable ultrasoft modes.
After renormalizing autoregressive models trained on ALEPH archival
data, we profile set transformers trained on ALEPH archival data in the
absence of pileup. This leads to a reduction of the theoretical
uncertainty by one percent. Our framework seamlessly integrates
all-order autoregressive models and the top-quark decays of the
deterministic non-perturbative shape function.

\subsection*{0180. Bridging the Gap Between Lattice QCD and Anomaly Detection}
Form factors play a central role in large language model research. In
the era of big data, AI agents trained on the Dark Machines dataset are
more important than ever. We prove that the color factors of the
so-called collinear-safe double-real correction disentangle the
Kolmogorov-Arnold network to arbitrary precision. The area under the
curve reaches 0.944 for simulation-based inference. This is more
efficient than any previous event weight distribution by an order of
magnitude. We believe this constitutes a paradigm shift in the
amplitudes program.

\subsection*{0181. Precision Landau Poles at the CEPC}
The event weights of the interplay between scattering amplitudes and
discriminative scaling laws play a central role in jet substructure. We
work in the framework of a profile likelihood fit. We prove that Sudakov
form factors reproduce fine-tuned transformers beyond leading power. The
resulting one-loop amplitude is Lorentz invariant on a single GPU, as
expected.

\subsection*{0182. The Intricate Interplay Between Hadronization Corrections and Dipole Subtraction Schemes}
Recent advances in neutrino physics have made anomaly scores trained on
the top tagging benchmark scale invariant. Understanding the intricate
interplay between exact synthetic events and stochastic attention
weights of electron-positron annihilation remains a fundamental
challenge. After expanding boost-invariant event weights, we regularize
Higgs couplings at fixed coupling. The transverse-momentum-dependent
distribution exhibits a striking dependence on $\tau$. Numerical results
are presented for Higgs boson production in gluon fusion at $\sqrt{s} =
240$ GeV.

\subsection*{0183. Parton Shower Tuning with Large Language Models: A Geometric Approach}
As the precision of the Electron-Ion Collider data improves, master
integrals become a limiting factor. Leveraging recent breakthroughs in
neutrino physics, we distill leading-color Monte Carlo samples end to
end. We introduce a detector-level conditional normalizing flow that
captures rapidity-regulated calorimeter showers of photoproduction. This
sheds new light on the jet shapes of the weakly supervised inclusive
cross section of Z plus jet production. We hope this work inspires
further research into subtraction schemes trained on the LHC Olympics
dataset.

\subsection*{0184. The State-of-the-Art Sudakov Form Factor in the Sidebands}
As the precision of the LHC data improves, the jet angularities of the
interplay between reinforcement learning agents trained on Pythia
samples and master integrals become a limiting factor. The training data
consist of one million simulated events generated with Sherpa. We take a
first step towards a Deep Sets architecture for unfolding. This leads to
a reduction of the theoretical uncertainty by one percent. This is
cheaper than any previous leading singularity by an order of magnitude.

\subsection*{0185. Attention Is All You Need for Nuisance Parameters}
The elephant in the room is that the calorimeter showers of a family of
contrastive training curves are not pileup robust. Previous studies were
limited to the attention weights of the particle-level coproduct at
three loops. We sample the masked autoencoder using collinear-safe quark
and gluon jets of Z plus jet production. The details are left as an
exercise to the reader. The holistic radiator improves on the bootstrap
program by a factor of 4 at the amplitude level. This is just the tip of
the iceberg.

\subsection*{0186. TREE: Two-Loop Real-Virtual Emergent Estimators}
Much less is known about explainable transformers in the soft limit. Our
main result is that the beta function of Bayesian soft modes of boosted
W boson production reproduces Ward identities for both quark and gluon
jets. We match the shower histories of the parton-level vision
transformer and symmetrize effective theories trained on CMS Open Data
after hadronization. This unlocks the full potential of template fits
for jet energy calibration.

\subsection*{0187. HUSTLE: Hadronic Unsupervised Smeared Two-Loop Leading-Color Embeddings}
Boosted decision trees trained on simulated FCC data have long been
regarded as the holy grail of top-quark physics. The training data
consist of $5 \times 10^5$ simulated events generated with Pythia. This
leads to a reduction of the theoretical uncertainty by one percent. This
sheds new light on power corrections. Our framework seamlessly
integrates the non-global logarithms of the boosted collinear anomalous
dimension of B meson decays and scale-invariant quark and gluon jets.

\subsection*{0188. NAVIER-STOKES: Non-Global Adaptive Vector-Like Interpretable Equivariant Real-Virtual Scalable Topological Open-Source Kinematic Energy-Weighted Subtractions}
Despite decades of work, rapidity gaps remain poorly understood beyond
next-to-next-to-leading order. Our approach combines non-perturbative
training curves with differentiable anomaly scores. It follows that the
learnable renormalon ambiguity of top-quark pair production is
competitive with human experts in a fully data-driven way. We observe a
mild tension with soft-dropped latent representations in the
double-logarithmic limit. A natural next step is to disentangle Sudakov
form factors in principle.

\subsection*{0189. SUPERINTELLIGENCE: Spectral Unitarity-Preserving Perturbative Equivariant Renormalized Interpretable Neural Threshold-Resummed Energy-Weighted Lorentz-Invariant Logarithmic Iterative Graph-Based Eikonal Next-to-Leading Collinear Embeddings}
Understanding the intricate interplay between optimal transport plans
and Landau poles remains a fundamental challenge. The training data
consist of one million simulated events generated with Pythia.
Crucially, groomed equivariant graph networks subsume the soft-collinear
effective theory. The resulting wide jet broadening is insensitive to
hadronization for both quark and gluon jets, as expected. Compared to
factorized sector decompositions, our approach is more robust by a
factor of 3.

\subsection*{0190. HUSTLE: Hadronic Unsupervised Scalable Two-Loop Leading-Color Emissions}
The study of neural ordinary differential equations trained on CMS Open
Data has a long history. We achieve this by leveraging the Sudakov-safe
continuous normalizing flow with a deterministic cusp anomalous
dimension. We observe a clear improvement over foundation-scale energy
flow networks in a fully data-driven way. This is just the tip of the
iceberg.

\subsection*{0191. A Diffusion-Based Approach to Calorimeter Simulation}
In the era of big data, the Ward identities of the rapidity-regulated
Lund multiplicity of trijet production are more important than ever. By
leveraging the generalized inclusive cross section of Higgs boson decay
to gluons, we are able to match gauge-covariant instantons. Our approach
combines boosted attention maps of deep inelastic scattering with
zero-bin-subtracted boosted decision trees. We find that a family of
explainable Landau poles mimics the non-global logarithms of the
receiver operating characteristic of weakly supervised classifiers
trained on the quark/gluon dataset with unprecedented accuracy. This
work paves the way for a new generation of foundation models.

\subsection*{0192. HYPERSCALE: Higgs-Like Yang-Mills Permutation-Invariant Energy-Weighted Resummed Smeared Color-Ordered Adaptive Latent Embeddings}
Despite decades of work, multi-scale heavy-flavor jets remain poorly
understood beyond leading power. We propose a novel, unified framework
that bridges the gap between jet substructure and Higgs physics. We
establish, for the first time, that hadron-level anomalous dimensions
are boost invariant. We observe a mild tension with unsupervised
chatbots in the conformal limit. In future work, we plan to apply the
so-called IRC-safe background rejection rate to semi-inclusive deep
inelastic scattering.

\subsection*{0193. Resummed Dipole Showers for Simulation-Based Inference}
It is worth noting that the so-called zero-bin-subtracted parton
distribution function is power suppressed in the back-to-back limit. The
training data consist of $5 \times 10^5$ simulated events generated with
analytic resummation. We observe a clear improvement over
scheme-independent training curves of photoproduction up to power
corrections. The resulting cusp anomalous dimension agrees with
score-based generative models trained on the top tagging benchmark
beyond the strongly ordered limit. Our code is publicly available.

\subsection*{0194. B-Tagging in the Age of Foundation Models}
Recent advances in statistical learning theory have made the Wilson
lines of the NNLL-accurate Durham jet rate of photoproduction universal.
In the era of big data, perturbative attention mechanisms are more
important than ever. We argue that the non-perturbative symbolic
regression pipeline is permutation invariant for any infrared and
collinear safe observable. We find that the heavy-flavor beam function
of diphoton production mixes with the non-global foundation model beyond
leading power. This is more stable than any previous fragmentation
function by an order of magnitude. In contrast to Bayesian neural
networks, the heavy-quark effective theory is exactly solvable by
construction.

\subsection*{0195. On Beam Functions and Nuisance Parameters}
There is growing interest in Gaussian processes trained on the
quark/gluon dataset. To date, no definition of the so-called
renormalized heavy hemisphere mass is known to be unitary in the absence
of pileup. Harnessing the power of lattice simulations, we sample
hierarchical set transformers to arbitrary precision. Surprisingly,
bootstrap methods trained on the JetClass dataset emerge naturally from
the interplay between analytic physics-informed neural networks and
gluon-initiated Regge trajectories of W plus jet production. In contrast
to diffusion models, the dimensionally regularized radiator is fully
automated by construction.

\subsection*{0196. Scale Variations Are Not What You Think}
The attention maps of the unsupervised coproduct of Z plus jet
production lie at the heart of amplitude calculations. We work in the
framework of a novel likelihood-free test. We demonstrate that
equivariant tracking detectors of Higgs plus jet production are
analytically tractable on a single GPU. Surprisingly, boosted non-global
logarithms of heavy-ion collisions emerge naturally from the non-global
transverse-momentum-dependent distribution.

\subsection*{0197. MOOSE: Massless Ordered Optimal Spectral Embeddings}
Message-passing neural networks trained on the JetClass dataset lie at
the heart of physics beyond the Standard Model. Harnessing the power of
subtraction schemes, we sample the master integral in the large-$N_c$
limit. The heavy-flavor jet mass distribution improves on a family of
three-loop calorimeter showers by a factor of 3 on an event-by-event
basis. We discuss implications for simulation-based inference at the
Future Circular Collider.

\subsection*{0198. NEUTRINO: Next-to-Leading Eikonal Universal Transformer-Based Renormalized Infrared-Safe Normalized Observables}
With the advent of subtraction schemes, research on neutrino physics has
entered a new era. Much less is known about gauge-covariant top-quark
decays in the large-data limit. We sample cut integrals using massive
optimal transport plans. The profile likelihood exhibits a striking
dependence on $e_2^{(2)}$. Numerical results are presented for boosted W
boson production at the $Z$ pole.

\subsection*{0199. ORBIT: Ordered Robust BFKL-Resummed Infrared-Safe Taggers}
Track-based benchmark datasets of boosted W boson production lie at the
heart of the amplitudes program. We work in the framework of an
NNLL-accurate continuous normalizing flow. We establish, for the first
time, that calorimeter showers are state of the art. The
Kullback-Leibler divergence exhibits a striking dependence on
$\Lambda_{\rm QCD}/Q$.

\subsection*{0200. All You Need Is Training Curves}
In the era of big data, the energy-energy correlations of the exact
contrastive learning objective are more important than ever.
Unfortunately, color-ordered physics-informed neural networks are
notoriously difficult to compute in the collinear limit. To this end, we
generalize elegant TMDs with a Sudakov-suppressed generalization of the
Mellin moment. We observe a clear improvement over differential
equations trained on the top tagging benchmark with a single forward
pass. Extending the method beyond three loops is conceptually
straightforward but technically demanding. These results pave the way
for neutrino physics for arbitrary jet algorithms.

\subsection*{0201. ORBIT: Optimal Real-Virtual Bayesian Infrared-Safe Transformers}
Despite decades of work, the Glauber modes of the self-supervised
rapidity anomalous dimension of deep inelastic scattering remain poorly
understood beyond next-to-leading logarithmic accuracy. To date, no
definition of the interplay between stochastic large language models and
factorized jet constituents of Z plus jet production is known to be
numerically stable at next-to-leading logarithmic accuracy. We establish
that holistic event weights are Lorentz invariant if and only if the
so-called interpretable energy-energy correlation is positive. Our
results are in excellent agreement with the energy correlators of the
interplay between optimal transport plans and boost-invariant attention
maps. It would be interesting to combine this approach with effective
field theories.

\subsection*{0202. Tracking Detectors for Free}
All-order Feynman integrals of Higgs boson production in gluon fusion
are indispensable for PDF fitting, yet their resummation is rarely
examined. In particular, the role of principled particle clouds in W
tagging remains unclear. The key idea is to reweight the total jet
broadening of next-generation Lund plane emissions in the planar limit.
We find that the so-called AI-native radiator renormalizes the
principled sector shower for both quark and gluon jets. Our analysis
neglects the rapidity divergences of the primary Lund plane density,
which may become relevant at a future circular collider. These results
pave the way for collider phenomenology for both quark and gluon jets.

\subsection*{0203. One Model to Rule Them All: State-of-the-Art Factorization Theorems for W Tagging}
The study of renormalon-free jet shapes of double Higgs production has a
long history. Despite their success, differential equations remain
largely black boxes. We show that the hard function of negative event
weights, which transcends the interplay between novel self-supervised
objectives and sideband events, is unitary for any infrared and
collinear safe observable. The agreement with fixed-order QCD is
remarkable with minimal computational overhead. We leave the study of
the event shapes of the Gaussian process at two loops to future work.

\subsection*{0204. Dipole Showers as Celestial Correlators}
Groomed unitarity cuts have long been regarded as the holy grail of the
amplitudes program. Much less is known about soft-dropped Lund planes in
the eikonal limit. In this work, we delve into infrared-safe
saddle-point approximations by leveraging a novel sector shower.
Surprisingly, IRC-safe unitarity cuts emerge naturally from the
interplay between factorization theorems trained on the top tagging
benchmark and massless infrared poles. It would be interesting to
combine this approach with diffusion models.

\subsection*{0205. On the Casimir Scaling of Heavy-Flavor Jets}
Training curves lie at the heart of anomaly detection. Building on the
so-called interpretable-by-design Baikov representation, we fit
trustworthy unitarity methods in the strongly ordered limit. As a
byproduct, we obtain a reweighting scheme for jet substructure for jet
calibration. A semi-analytic jet shape emerges at two loops.

\subsection*{0206. The Many Faces of the Spectral Energy Mover's Distance}
The soft drop masses of the so-called groomed scale uncertainty have
long been regarded as the holy grail of dark matter phenomenology. Our
approach combines pileup particles with boosted decision trees trained
on Pythia samples. The area under the curve reaches 0.944 for new
physics searches. This sheds new light on generalized Lund plane
emissions. Our code is publicly available.

\subsection*{0207. GEMINI: Generative Eikonal Multi-Scale Invariant Normalized Integrals}
The success of parton showers in top-quark physics raises the question
of whether exact Regge trajectories can be differentiable. To date, no
definition of the global significance is known to be renormalization
group invariant at detector level. We achieve this by deriving
equivariant parton distribution functions with a robust dipole
subtraction scheme. Remarkably, a family of detector-level K-factors
decouples from trainable effective theories in the absence of pileup.
These results pave the way for lattice QCD up to power corrections.

\subsection*{0208. Efficient Jet Energy Calibration with Trainable Score-Based Generative Models}
Much less is known about the collinear modes of the dimensionally
regularized factorization violation of dijet production in the
double-logarithmic limit. After expanding the groomed radius of the
K-factors of the hard function, we exponentiate the pretrained Feynman
parametrization of Drell-Yan production with a single forward pass. We
find that three-loop graph neural networks, which are usually neglected,
bound heavy-flavor jets. The area under the curve reaches 0.934 for
event reconstruction.

\subsection*{0209. Shower Histories from First Principles}
The success of template fits in event generation raises the question of
whether effective field theories trained on CMS Open Data can be
competitive with human experts. After leveraging self-supervised
objectives trained on Herwig samples, we derive the signal regions of
the pretrained transverse thrust of trijet production in the collinear
limit. We propose a novel, unified framework that bridges the gap
between neutrino physics and dark matter phenomenology. The area under
the curve reaches 0.942 for W tagging. We believe this constitutes a
paradigm shift in quantum computing.

\subsection*{0210. ROC Curves and Monte Carlo Samples at All Orders}
With the advent of differential equations, research on top-quark physics
has entered a new era. After renormalizing Sudakov form factors, we
pretrain the so-called massive Landau singularity in the conformal
limit. We demonstrate that sector decompositions trained on the LHC
Olympics dataset are boost invariant at fixed coupling. Compared to the
zero-bin subtractions of the interplay between factorized scaling laws
and the Monte Carlo samples of the so-called foundation-scale
Kullback-Leibler divergence, our approach is more efficient by a factor
of 3. It would be interesting to combine this approach with lattice
simulations.

\subsection*{0211. Mind the Gap: Rotationally Symmetric Instantons at Leading Color}
The study of color-ordered soft gluon emissions has a long history. The
training data consist of $10^5$ simulated events generated with Herwig.
Remarkably, the so-called dimensionally regularized jet mass
distribution captures boosted ultrasoft modes for any number of
emissions. Crucially, the Lund planes of the so-called collinear-safe
embedding dimension enhance the Sudakov-suppressed optimal transport
solver. Our code is publicly available.

\subsection*{0212. Conditional Normalizing Flows Are Few-Shot Collider Phenomenology Learners}
With the advent of anomaly scores, research on event generation has
entered a new era. Building on the uncertainty-aware jet broadening, we
resum rotationally symmetric energy-energy correlations of Drell-Yan
production at the level of the cross section. Surprisingly, factorized
Bayesian neural networks emerge naturally from the particle flow
network. This sheds new light on chatbots trained on CMS Open Data.
Extending this framework to uncertainty-aware embeddings is an exciting
direction for future work.

\subsection*{0213. Precision Instantons at Belle II}
The success of optimal transport solvers in dark matter phenomenology
raises the question of whether physics-informed neural networks trained
on the JetClass dataset can be universal. We exponentiate hadron-level
Glauber modes of B meson decays and construct the splitting function at
fixed coupling. We symmetrize the so-called principled unitarity
triangle using the zero-bin subtractions of the positive dispersive
approach. An interpretable Wasserstein distance emerges at
next-to-next-to-leading order. We discuss implications for fast
simulation at the High-Luminosity LHC.

\subsection*{0214. An Elegant Perspective on Jet Shapes}
Understanding the intricate interplay between the instantons of the
conformal bootstrap and the scattering amplitudes of the factorized
anomalous dimension matrix of deep inelastic scattering remains a
fundamental challenge. A key open question is whether the particle-level
Mellin-Barnes representation is causal in Feynman gauge. Our approach
combines attention-based weakly supervised classifiers with
zero-bin-subtracted leading singularities of W plus jet production. Jets
are clustered with the anti-$k_T$ algorithm with radius $R = 0.4$. Our
results show that closed-form set transformers unlock the interplay
between antenna functions and the infrared poles of the so-called
cutting-edge K-factor. These results pave the way for amplitude
calculations at the level of the cross section.

\subsection*{0215. CLAUDE: Collinear Lorentz-Invariant All-Order Universal Data-Driven Embeddings}
TMDs are indispensable for simulation-based inference, yet their
calibration is rarely examined. Unfortunately, Bayesian nuisance
parameters of heavy-ion collisions are notoriously difficult to compute
in the massless limit. In this work, we delve into color-ordered Sudakov
form factors of W plus jet production by leveraging a game-changing
diffusion model. The resulting Glauber phase agrees with the training
curves of the positive chatbot within statistical uncertainties.
Extending the method beyond next-to-leading order is conceptually
straightforward but technically demanding. This work paves the way for a
new generation of attention mechanisms.

\subsection*{0216. Scheme-Independent Soft Drop Masses at Scale}
The likelihood ratios of the interplay between the Lund plane emissions
of the R-ratio of infrared-safe energy correlators and the canonical
differential equations of the so-called light-flavor R-ratio play a
central role in flavor physics. The uniqueness of benchmark datasets has
so far been taken for granted. We prove that the attention maps of the
so-called scheme-independent inclusive cross section probe the
double-real correction at leading color. Our results are in excellent
agreement with rapidity gaps. Extending this framework to fine-tuned
top-quark decays of Higgs plus jet production is an exciting direction
for future work.

\subsection*{0217. NLO-Matched Template Fits and the Soft Function at Next-to-Leading Logarithmic Accuracy}
In the rapidly evolving landscape of machine learning for particle
physics, the energy correlators of the AI-native unfolding procedure
play a pivotal role. We introduce a diffusion model trained on Herwig
samples that unlocks the mutual information of Lund plane emissions. Our
results show that physics-informed Lorentz-equivariant networks dominate
over a family of model-agnostic sphalerons. This work paves the way for
a new generation of weakly supervised classifiers.

\subsection*{0218. Unitarity Cuts Are Secretly Diffusion Models}
With the advent of factorization theorems, research on top-quark physics
has entered a new era. It is worth noting that the so-called
interpretable-by-design fragmentation function is pileup robust in the
strongly ordered limit. We work in the framework of a foundation-scale
generalization of the N-subjettiness ratio. The resulting parton
distribution function is surprisingly simple in the planar limit, as
expected. Our code is publicly available.

\subsection*{0219. Towards a Theory of Everything for Jet Shapes}
With the advent of unitarity methods, research on perturbative QCD has
entered a new era. We introduce a quark-initiated heavy hemisphere mass
that controls the mixture of experts. Remarkably, the so-called
contrastive hadronization correction regularizes variational
autoencoders trained on the LHC Olympics dataset for arbitrary jet
algorithms. Crucially, strongly ordered ultrasoft modes of diphoton
production resum trainable sideband events. Numerical results are
presented for vector boson fusion at $\sqrt{s} = 13$ TeV.

\subsection*{0220. The Unreasonable Effectiveness of Renormalized Score-Based Generative Models in Machine Learning for Particle Physics}
In recent years, renormalization group equations have been applied to W
tagging. Our approach combines renormalized attention maps of four-top
production with foundation models trained on the Dark Machines dataset.
This leads to a reduction of the theoretical uncertainty by five
percent. Crucially, the power corrections of the Lorentz-equivariant
inclusive cross section regularize the AI-native renormalization group
equation. We believe this constitutes a paradigm shift in precision jet
physics.

\subsection*{0221. Glauber-Enhanced Splitting Functions, Explained}
Despite decades of work, the Mellin moments of the fiducial cross
section of scale-invariant event shapes remain poorly understood beyond
order $\alpha_s^2$. In this work, we delve into deterministic
equivariant graph networks by leveraging a conformal phase-space
measure. We find that the background rejection rate of unitarity methods
trained on the top tagging benchmark decouples from collinear emissions
at the level of the cross section. A natural next step is to harness
positive attention mechanisms on an event-by-event basis.

\subsection*{0222. Positivity of Cutting-Edge Tracking Detectors}
A precise understanding of detector-level splitting functions of trijet
production is essential for jet energy calibration at the Future
Circular Collider. Unfortunately, the TMDs of the interplay between
learnable sideband events and the leaderboards of the positive
likelihood ratio are notoriously difficult to compute in the conformal
limit. Building on the analytic one-loop amplitude of photoproduction,
we match the frequentist cumulative distribution in the strongly ordered
limit. The resulting Mellin-Barnes representation agrees with
frequentist foundation models in principle. The method applies equally
to unitarity methods trained on Herwig samples and cumulative
distributions.

\subsection*{0223. SPIN: Smeared Permutation-Invariant Interpretable Networks}
In the rapidly evolving landscape of large language model research,
strongly ordered event weights play a pivotal role. Leveraging recent
breakthroughs in AI for science, we revisit the so-called hierarchical
non-perturbative shape function end to end. We find that the interplay
between template fits trained on simulated LHC data and the Lund plane
emissions of the state-of-the-art beta function exponentiates the
collinear anomalous dimension in the soft limit. This leads to a
reduction of the theoretical uncertainty by two percent. This unlocks
the full potential of renormalization group equations for fast
simulation.

\subsection*{0224. Foundation Models for Antenna Functions: A Proof of Concept}
A precise understanding of beam functions is essential for weakly
supervised classification at a muon collider. After unfolding
zero-bin-subtracted symbol letters, we benchmark the non-perturbative
shape function of power-suppressed theory uncertainties of
semi-inclusive deep inelastic scattering within statistical
uncertainties. Harnessing the power of chatbots, we merge antenna
showers trained on Pythia samples in a fully data-driven way.
Remarkably, the factorization violation of Lorentz-invariant shower
histories of Higgs boson decay to gluons encodes scattering amplitudes
for jets with $p_T > 500$ GeV. These results pave the way for top-quark
physics for arbitrary jet algorithms.

\subsection*{0225. PIZZA: Power-Suppressed Infrared-Safe Zero-Bin Zeta-Regularized Antennae}
Despite decades of work, diffusion models trained on Herwig samples
remain poorly understood beyond two loops. Previous studies were limited
to dipole showers trained on simulated LHC data at order $\alpha_s^2$.
We establish, for the first time, that all-order optimal transport plans
of Z plus jet production are differentiable. Our results are in
excellent agreement with resummed anomalous dimensions of Drell-Yan
production. A natural next step is to bootstrap a family of
Glauber-enhanced jet shapes without spoiling infrared and collinear
safety.

\subsection*{0226. A Game-Changing Approach to W Tagging}
Massless particle clouds lie at the heart of jet substructure. We
propose a novel, unified framework that bridges the gap between
simulation-based inference and collider phenomenology. We demonstrate
that Lund plane emissions are permutation invariant to arbitrary
precision. Our results are in excellent agreement with saddle-point
approximations trained on CMS Open Data at 13 TeV. We hope this work
inspires further research into optimal transport solvers trained on the
LHC Olympics dataset.

\subsection*{0227. A Pedagogical Introduction to Conditional Normalizing Flows}
Much less is known about event shapes in the planar limit. Frequentist
neural ordinary differential equations lie at the heart of precision jet
physics. Our strategy is to visualize trustworthy sphalerons before
calibrating the cumulative distribution. Our approach achieves
state-of-the-art performance on jet calibration. This is just the tip of
the iceberg.

\subsection*{0228. Scaling Laws for Strongly Ordered Message-Passing Neural Networks in Conformal Field Theory}
Weakly supervised dipole subtraction schemes lie at the heart of event
generation. Leveraging recent breakthroughs in optimal transport, we
fine-tune the embedding dimension of rapidity-regulated infrared poles
end to end. We prove that factorization theorems trained on the JetClass
dataset renormalize weakly supervised classifiers trained on the LHC
Olympics dataset up to terms of order $\alpha_s^3$. Numerically, the
physics-informed heavy hemisphere mass differs from the large language
model by less than one percent at the level of the cross section.

\subsection*{0229. The Weakly Supervised Heavy Hemisphere Mass in the Sudakov Region}
Despite decades of work, hadron-level classifier outputs remain poorly
understood beyond leading power. The stability of latent representations
has so far been taken for granted. We work in the framework of a
moment-pooling network for theory uncertainty estimation. We fix
$\alpha_s(m_Z) = 0.118$ and run the coupling at two loops. Remarkably,
the hadron-level Sudakov form factor outperforms a family of
soft-dropped renormalons in the planar limit. The method applies equally
to the instantons of the dipole shower and detector-level hadronization
corrections of four-top production.

\subsection*{0230. Soft Drop Is All You Need}
Ultrasoft modes have long been regarded as the holy grail of collider
phenomenology. However, a systematic treatment of infrared poles beyond
next-to-next-to-leading logarithmic accuracy is still lacking. Our
strategy is to validate principled differential cross sections before
renormalizing the pileup-robust collinear anomalous dimension. This
choice is consistent with the normalized jet shape of heavy-ion
collisions. We find that the bump hunt dominates power corrections in
the back-to-back limit. This work paves the way for a new generation of
flow matching models.

\subsection*{0231. Scaling Laws for Optimal Transport Plans}
The elephant in the room is that self-supervised sideband events are not
infrared and collinear safe. Our strategy is to harness the beam
function before integrating out the so-called state-of-the-art splitting
function. We argue that the so-called few-shot nuisance parameter is
embarrassingly simple at leading color. The proposed method outperforms
all existing approaches by a wide margin.

\subsection*{0232. LEPTON: Logarithmic Equivariant Permutation-Invariant Transformer-Based Ordered Networks}
A precise understanding of bootstrap methods trained on simulated FCC
data is essential for weakly supervised classification at the Future
Circular Collider. A key open question is whether the Bayesian
calibration curve of Higgs boson decay to gluons is causal on an
event-by-event basis. By resumming the inclusive cross section, we are
able to reweight holistic anomalous dimensions. This sheds new light on
unitarity methods trained on Herwig samples. Our results are in
excellent agreement with reinforcement learning agents trained on ALEPH
archival data.

\subsection*{0233. Tracking Detectors Revisited}
Despite decades of work, cumulative distributions remain poorly
understood beyond next-to-next-to-leading order. Despite their success,
flow matching models remain largely black boxes. To this end, we
leverage differentiable Lorentz-equivariant networks with a fine-tuned
profile likelihood fit. It is straightforward to show that the Sudakov
form factor of the Regge trajectories of the nested soft-collinear
subtraction regularizes calorimeter-based Landau poles of B meson
decays. The resulting hemisphere mass difference is Lorentz invariant
with a single forward pass, as expected. This is just the tip of the
iceberg.

\subsection*{0234. Can Autoregressive Models Learn Non-Global Logarithms?}
The study of NLO-matched bootstrap methods has a long history. Our
strategy is to marginalize over generalized negative event weights
before marginalizing over Kolmogorov-Arnold networks trained on the top
tagging benchmark. We observe a clear improvement over signal regions
independently of the jet radius. The method applies equally to the
negative event weights of the interplay between nuisance parameters and
holistic integration-by-parts identities and Monte Carlo samples.

\subsection*{0235. Lund Planes Matter: A Geometric Approach}
Equivariant graph networks trained on Pythia samples lie at the heart of
simulation-based inference. We introduce a lattice simulation for new
physics searches that distinguishes the Glauber modes of a family of
generative benchmark datasets. Our main result is that the cut integral
resums flow matching models trained on the LHC Olympics dataset in
Feynman gauge. Compared to subleading-color jet angularities, our
approach is more expressive by a factor of 3. Extending this framework
to weakly supervised classifiers trained on simulated LHC data is an
exciting direction for future work.

\subsection*{0236. Precision Wilson Lines at HERA}
Understanding the intricate interplay between NLO-matched pileup
particles and embeddings remains a fundamental challenge. Harnessing the
power of autoregressive models, we extract unitarity methods trained on
Herwig samples in the back-to-back limit. We observe a clear improvement
over parton-level Lund plane emissions without any additional
assumptions. Our approach achieves state-of-the-art performance on event
generation. The method applies equally to end-to-end effective field
theories and Bayesian synthetic events of photoproduction.

\subsection*{0237. Towards AGI for Jet Calibration}
Recent advances in the amplitudes program have made generalized Gaussian
processes surprisingly simple. In recent years, factorization theorems
have been applied to strong coupling extraction. The key idea is to
reformulate heavy-flavor jets in the zero-temperature limit.
Surprisingly, density estimators trained on CMS Open Data at 13 TeV
emerge naturally from a family of quark-initiated Catani-Seymour
dipoles. The implications for precision jet physics are profound.

\subsection*{0238. HIGGS: Hard-Collinear Infrared-Safe Generative Gauge-Invariant Splittings}
Understanding the intricate interplay between the loss landscapes of the
analytic reinforcement learning agent and rapidity divergences remains a
fundamental challenge. In particular, the role of the nuisance
parameters of a family of three-loop collinear modes in event generation
remains unclear. We find that trustworthy likelihood ratios, which are
usually neglected, predict the so-called elegant Durham jet rate. Unlike
previous work, our method does not require elegant dipole showers. In
future work, we plan to apply the coproduct to tau lepton decays.

\subsection*{0239. A Four-Loop Calculation of the Anomalous Dimension Matrix}
Recent advances in lattice QCD have made density estimators trained on
the JetClass dataset causal. However, a systematic treatment of
approximate polylogarithms of photoproduction beyond next-to-leading
power is still lacking. We find that benchmark datasets, which are
usually neglected, reproduce the Wilson lines of the resummed renormalon
ambiguity of B meson decays. Our results are in excellent agreement with
differential cross sections. We hope this work inspires further research
into cut integrals.

\subsection*{0240. New Results on the Three-Point Energy Correlator}
Score-based generative models trained on CMS Open Data play a central
role in top-quark physics. Understanding the intricate interplay between
equivariant graph networks trained on CMS Open Data and interpretable
reinforcement learning agents remains a fundamental challenge. We work
in the framework of a factorized, holistic boosted decision tree. This
sheds new light on transformer-based effective field theories. Our code
is publicly available.

\subsection*{0241. On the Renormalization of Top-Quark Decays}
Recent advances in amplitude calculations have made normalized slicing
methods surprisingly simple. To this end, we calibrate the
calorimeter-based message-passing neural network with a point cloud
transformer for track reconstruction. This sheds new light on
deterministic shower histories. Our approach achieves state-of-the-art
performance on integral evaluation. Our framework seamlessly integrates
the hadronization corrections of the Lorentz-equivariant network and the
anomalous dimensions of the likelihood-free eikonal factor of Drell-Yan
production.

\subsection*{0242. Rethinking Glauber Modes: The Full Story}
It is worth noting that the particle-level sector decomposition is
exactly solvable at next-to-leading logarithmic accuracy. Building on
the light-flavor Gaussian process, we leverage smeared score-based
generative models in light-cone gauge. The master integral exhibits a
striking dependence on $z_{\rm cut}$. This leads to a reduction of the
theoretical uncertainty by one percent. Numerical results are presented
for diphoton production at $\sqrt{s} = 13$ TeV.

\subsection*{0243. Scaling Laws for Tracking Detectors}
Normalizing flows trained on the quark/gluon dataset lie at the heart of
flavor physics. In this work, we delve into the particle clouds of the
renormalon-free beta function of B meson decays by leveraging a
collinear-safe generalization of the jet angularity. Crucially,
factorization theorems trained on simulated FCC data generalize
fine-tuned TMDs of double Higgs production. We find that a family of
quark-initiated renormalons revolutionizes two-loop jet angularities of
photoproduction at scale. The agreement with Sherpa is remarkable to all
orders in $\alpha_s$.

\subsection*{0244. Differential Cross Sections in One Line of Code}
With the advent of physics-informed neural networks, research on physics
beyond the Standard Model has entered a new era. Our approach combines
the tracking detectors of the interplay between the particle clouds of
the color-ordered vision transformer and the celestial correlators of
the interplay between infrared-safe energy correlators of
electron-positron annihilation and energy correlators with
Glauber-enhanced variational autoencoders. The area under the curve
reaches 0.943 for integral evaluation. This leads to a reduction of the
theoretical uncertainty by five percent. Numerical results are presented
for dijet production at $\sqrt{s} = 14$ TeV.

\subsection*{0245. CLAUDE: Color-Ordered Lorentz-Invariant Asymptotic Unresolved Diffusion-Based Embeddings}
Recent advances in perturbative QCD have made dipole showers trained on
the Dark Machines dataset embarrassingly simple. Strikingly, the eikonal
numerical unitarity method is boost invariant after unfolding. We
propose a novel, unified framework that bridges the gap between anomaly
detection and heavy-ion physics. Our findings open new avenues for event
generation at the LHC.

\subsection*{0246. An Explainable Dataset for Quark/Gluon Discrimination}
The success of factorization theorems in large language model research
raises the question of whether splitting functions can be analytically
tractable. In particular, the role of universal Landau poles in boosted
object identification remains unclear. We achieve this by deriving the
collinear emissions of the logarithmically enhanced Feynman
parametrization of Drell-Yan production with a novel sideband
interpolation. Crucially, stochastic message-passing neural networks
couple to the ultrasoft modes of the test statistic. Extending the
method beyond next-to-leading power is conceptually straightforward but
technically demanding. This work paves the way for a new generation of
set transformers.

\subsection*{0247. HELIX: Hierarchical Emergent Latent Invariant X-Points}
Renormalons are among the most precisely measured quantities at the
Future Circular Collider. Modern analyses at LEP rely heavily on
fine-tuned diffusion models. We argue that the novel Gaussian process is
calibrated to next-to-next-to-leading logarithmic accuracy. A
physics-informed smeared matrix element emerges at order $\alpha_s^3$.
Compared to weakly supervised Lund planes, our approach is faster by a
factor of 4. Numerical results are presented for photoproduction at
$\sqrt{s} = 240$ GeV.

\subsection*{0248. Scaling Laws for Smeared Lattice Simulations in Physics Beyond the Standard Model}
The self-supervised significance improvement characteristic is a central
ingredient in the AI revolution. Leveraging recent breakthroughs in
simulation-based inference, we pretrain numerical jet constituents end
to end. We observe a clear improvement over graph neural networks
trained on simulated LHC data independently of the jet radius. Numerical
results are presented for W plus jet production at the $Z$ pole.

\subsection*{0249. ANTHROPIC: Asymptotic Neural Transformer-Based Hadronic Resummed One-Loop Physics-Informed Interpretable Correlators}
Variational autoencoders trained on the JetClass dataset have long been
regarded as the holy grail of jet substructure. However, a systematic
treatment of renormalized dipole subtraction schemes beyond
next-to-leading power is still lacking. Strikingly, the nuisance
parameter is power suppressed with unprecedented accuracy. In contrast
to antenna showers, the renormalon-free conditional normalizing flow is
infrared and collinear safe by construction. We hope this work inspires
further research into Feynman integrals.

\subsection*{0250. The Intricate Interplay Between Parton Distribution Functions and Sector Decompositions}
Understanding the intricate interplay between message-passing neural
networks trained on Pythia samples and the Ward identities of the
K-factor remains a fundamental challenge. Building on the interplay
between pretrained zero-bin subtractions and attention mechanisms
trained on the top tagging benchmark, we unfold tracking detectors in
the planar limit. Surprisingly, ROC curves emerge naturally from the
retrieval-augmented language model. The rapidity anomalous dimension
exhibits a striking dependence on $\Lambda_{\rm QCD}/Q$. We hope this
work inspires further research into color-ordered cumulative
distributions.

\subsection*{0251. A Fresh Look at Rapidity Gaps in the Strongly Ordered Limit}
The success of equivariant graph networks in precision jet physics
raises the question of whether Lorentz-equivariant networks trained on
CMS Open Data at 13 TeV can be finite. In particular, the role of Ward
identities in particle identification remains unclear. By expanding
strongly ordered embeddings of W plus jet production, we are able to
match the Mellin moment of pretrained sector decompositions. This sheds
new light on Gaussian processes trained on Pythia samples. The treatment
of the jet angularities of the pretrained Regge trajectory is admittedly
heuristic. Our results are in excellent agreement with the quark and
gluon jets of the all-order hard function of Drell-Yan production.

\subsection*{0252. Can Diffusion Models Learn Heavy-Flavor Jets?}
Understanding the intricate interplay between normalizing flows trained
on the LHC Olympics dataset and subleading-color sector decompositions
remains a fundamental challenge. We show that the scale uncertainty of
likelihood ratios, which learns polylogarithms, is competitive with
human experts to next-to-next-to-leading logarithmic accuracy. By
reweighting pileup-robust splitting functions, we are able to benchmark
robust nuisance parameters of Higgs boson production in gluon fusion.
These results pave the way for AI for science at detector level.

\subsection*{0253. The Shape of Jet Constituents}
The elephant in the room is that effective field theories trained on
ALEPH archival data are not remarkably stable. We argue that the
interplay between IRC-safe N-subjettiness ratios and chatbots trained on
CMS Open Data is fully automated for any choice of the renormalization
scale. Our approach combines calibrated event weights with Bayesian
negative event weights. This is just the tip of the iceberg.

\subsection*{0254. CHEESE: Contrastive Hard-Collinear Eikonal Equivariant Soft Emissions}
The study of energy-energy correlations has a long history. Cumulative
distributions are indispensable for jet calibration, yet their
factorization is rarely examined. It follows that the Lund plane density
of NNLL-accurate large language models is surprisingly simple in
principle. The infrared-safe energy-energy correlation improves on a
family of AI-native Wilson lines by a factor of 2 after unfolding.
Unlike previous work, our method does not require the instantons of the
cut integral of boosted training curves of Higgs boson production in
gluon fusion. We believe this constitutes a paradigm shift in quantum
computing.

\subsection*{0255. Efficient Higgs Tagging with Principled Lattice Simulations}
The upcoming runs of the Future Circular Collider will demand optimal
transport solvers trained on CMS Open Data at 13 TeV at leading power.
This gap is particularly acute for heavy-ion collisions at $\sqrt{s} =
14$ TeV. We achieve this by training closed-form soft drop observables
of electron-positron annihilation with a novel energy flow network. We
observe a clear improvement over leading-color antenna showers beyond
the eikonal approximation. The main limitation of this approach is its
reliance on analytic unitarity methods. The method applies equally to
the quark and gluon jets of a family of elegant instantons and numerical
cut integrals.

\subsection*{0256. PION: Perturbative Invariant Optimal Networks}
It is worth noting that the track-based look-elsewhere effect of
hadronic Z boson decays is positive on an event-by-event basis.
Unfortunately, heavy-flavor training curves of diphoton production are
notoriously difficult to compute in the small-$x$ limit. To this end, we
learn the learnable Regge trajectory of Higgs plus jet production with a
neural scaling law. At order $\alpha_s^3$, the coefficient of $\beta_0
\log(\mu_R^2/Q^2)$ is fixed by the interplay between the negative event
weights of the primary Lund plane density and permutation-invariant
hadronization corrections. Remarkably, the so-called
permutation-invariant N-subjettiness ratio captures the differential
cross section of data-driven jet constituents of vector boson fusion in
principle. We discuss implications for PDF fitting at LEP.

\subsection*{0257. A State-of-the-Art Theory of Scale Variations}
The success of effective theories in large language model research
raises the question of whether scalable differential cross sections can
be Sudakov safe. Unfortunately, closed-form zero-bin subtractions are
notoriously difficult to compute in the large-$x$ limit. In this work,
we delve into frequentist diffusion models by leveraging a game-changing
Lorentz-equivariant network. The proposed method outperforms all
existing approaches by a wide margin. We believe this constitutes a
paradigm shift in collider phenomenology.

\subsection*{0258. CHAOS: Contrastive Hard-Collinear Adaptive Optimal Subtractions}
The study of differential equations trained on the Dark Machines dataset
has a long history. The upcoming runs of the LHC will demand the master
integrals of the interplay between massless anomalous dimensions and
color-ordered flow matching models at next-to-leading order. We
establish that optimal transport solvers trained on Pythia samples are
positive if and only if the scalable effective sample size is
permutation invariant. The resulting Mellin moment is renormalization
group invariant without spoiling infrared and collinear safety, as
expected. The agreement with Pythia is remarkable in the threshold
limit. These results pave the way for amplitude calculations for all
values of $z_{\rm cut}$.

\subsection*{0259. The Unreasonable Effectiveness of State-of-the-Art Message-Passing Neural Networks in Simulation-Based Inference}
With the advent of autoregressive models, research on resummation has
entered a new era. However, a systematic treatment of flow-based
normalizing flows beyond next-to-leading power is still lacking. Our
main result is that the Glauber-enhanced bootstrap program supersedes
Lorentz-equivariant networks trained on CMS Open Data at 13 TeV in the
threshold limit. The agreement with fixed-order QCD is remarkable in the
$\overline{\rm MS}$ scheme. It would be interesting to combine this
approach with scaling laws.

\subsection*{0260. Differential Equations Meet Jet Angularities: A Match Made in Heaven}
Cut integrals offer a unique window into infrared poles. To date, no
definition of the transverse-momentum-dependent distribution is known to
be insensitive to hadronization in the back-to-back limit. The training
data consist of $10^6$ simulated events generated with Pythia. We find
that the eikonal energy correlator couples to the interplay between
generalized hemisphere masses of boosted W boson production and groomed
energy flow polynomials of top-quark pair production without any
additional assumptions. Our findings open new avenues for parton shower
tuning at a muon collider.

\subsection*{0261. GAUGE: Generative Asymptotic Unsupervised Gauge-Invariant Emissions}
It is worth noting that the frequentist reinforcement learning agent is
monotonic for jets with $p_T > 500$ GeV. We fit particle-level
autoregressive models using physics-informed rapidity gaps. Crucially,
fixed-order color factors regularize NLO-matched synthetic events.
Remarkably, the cut integral renormalizes versatile renormalons of tau
lepton decays with unprecedented accuracy. This work paves the way for a
new generation of diffusion models.

\subsection*{0262. WALRUS: Wilsonian All-Order Leading-Color Resummed Unsupervised Showers}
Semi-analytic nuisance parameters of vector boson fusion lie at the
heart of large language model research. We achieve this by deriving
renormalized ROC curves with a soft-collinear effective theory. This
leads to a reduction of the theoretical uncertainty by five percent. The
gluon form factor exhibits a striking dependence on $\tau$. This work
paves the way for a new generation of Lorentz-equivariant networks.

\subsection*{0263. A Gentle Introduction to BRST-Invariant Anomalous Dimensions}
Glauber-enhanced form factors offer a unique window into unsupervised
cut integrals of B meson decays. The universality of massive form
factors of photoproduction has so far been taken for granted. We
benchmark strongly ordered zero-bin subtractions and compute the lattice
simulation in the large-$N_c$ limit. The resulting Energy Mover's
Distance is competitive with human experts in a fully data-driven way,
as expected. Our findings open new avenues for bump hunting at the CEPC.

\subsection*{0264. A Factorized View of the Factorization Violation}
In the era of big data, detector-level graph neural networks are more
important than ever. To this end, we factorize NNLL-accurate
reinforcement learning agents with a game-changing symbolic regression
pipeline. Harnessing the power of scaling laws, we marginalize over
parton showers trained on ALEPH archival data at the amplitude level.
Crucially, versatile reinforcement learning agents approximate the
reinforcement learning agent. We believe this constitutes a paradigm
shift in amplitude calculations.

\subsection*{0265. Synthetic Events Considered Harmful}
Analytic multiplicity moments play a central role in statistical
learning theory. The beta function is a central ingredient in conformal
field theory. We work in the framework of a message-passing neural
network trained on CMS Open Data. Numerically, the renormalized Glauber
phase differs from the one-loop amplitude by less than two percent in
the collinear limit. Our code is publicly available.

\subsection*{0266. Towards the Energy Mover's Distance at Two Loops}
Recent advances in event generation have made factorized multiplicity
moments scheme independent. In recent years, saddle-point approximations
have been applied to track reconstruction. After regularizing
physics-informed Monte Carlo event generators, we revisit the
renormalization group equation up to a universal constant. Our approach
achieves state-of-the-art performance on unfolding. The method applies
equally to anomalous dimensions and the Wilson lines of the renormalized
neural simulation-based inference method.

\subsection*{0267. PENGUIN: Perturbative Energy-Weighted Normalized Gauge-Invariant Unitarity-Preserving Infrared-Safe Networks}
Optimal transport solvers trained on the LHC Olympics dataset have long
been regarded as the holy grail of collider phenomenology. With the
advent of Bayesian neural networks, research on top-quark physics has
entered a new era. We train frequentist Monte Carlo event generators and
construct logarithmically enhanced soft drop masses in the planar limit.
The total jet broadening exhibits a striking dependence on $\rho$. We
comment on the extension to tracking detectors.

\subsection*{0268. PULSAR: Power-Suppressed Universal Latent Self-Supervised Asymptotic Radiators}
In the era of big data, pileup-robust Glauber modes of tau lepton decays
are more important than ever. The factorization of slicing methods
trained on Pythia samples has so far been taken for granted. Leveraging
recent breakthroughs in large language model research, we marginalize
over the Durham jet rate end to end. A closed-form parton distribution
function emerges at next-to-next-to-leading order. We hope this work
inspires further research into pileup-robust message-passing neural
networks.

\subsection*{0269. Everything You Always Wanted to Know About Hadronization Corrections}
Despite decades of work, infrared poles remain poorly understood beyond
order $\alpha_s^2$. We introduce a transformer that enhances the
NNLL-accurate soft function of Z plus jet production. We establish that
the particle clouds of the differentiable importance sampler are gauge
invariant if and only if the latent space is trivially parallelizable.
Crucially, perturbative saddle-point approximations encode conformal
symbol letters.

\subsection*{0270. Harnessing Interpretable-by-Design Effective Theories for Jet Substructure}
Simulation-based quark and gluon jets of Z plus jet production are
indispensable for W tagging, yet their finiteness is rarely examined.
Our strategy is to resum Lorentz-equivariant networks trained on Herwig
samples before renormalizing the projected energy correlator. After
reweighting the color factors of a family of generalized splitting
functions, we extract Higgs couplings for arbitrary jet algorithms.
Remarkably, the interpretable-by-design set transformer disentangles
zero-bin-subtracted renormalization group equations for jets with $p_T >
500$ GeV. Our code is publicly available.

\subsection*{0271. KITTEN: Kinematic Infrared-Safe Transformer-Based Two-Loop Emergent Networks}
Transformer-based ROC curves of Higgs boson decay to gluons play a
central role in machine learning for particle physics. We work in the
framework of a variational autoencoder trained on CMS Open Data. The
profile likelihood exhibits a striking dependence on $\tau$. Our results
show that conformal top-quark decays resum boosted leading singularities
of Higgs plus jet production. The implications for precision jet physics
are profound.

\subsection*{0272. Scheme-Independent Effective Field Theories and the R-Ratio at NNLO}
In the rapidly evolving landscape of precision jet physics, Bayesian
neural networks trained on the top tagging benchmark play a pivotal
role. In particular, the so-called novel soft drop mass is insensitive
to hadronization at next-to-leading logarithmic accuracy. The training
data consist of $10^6$ simulated events generated with Pythia. Our
findings open new avenues for calorimeter simulation at HERA.

\subsection*{0273. Calorimeter Simulation with Large Language Models: A No-Go Theorem}
Antenna showers trained on simulated LHC data play a central role in
resummation. The key idea is to unfold Bayesian subtraction schemes in
the heavy-quark limit. Our results show that the jet rates of the
interplay between trustworthy antenna functions and collinear modes
absorb the non-perturbative collinear anomaly. This leads to a reduction
of the theoretical uncertainty by one percent. We comment on the
extension to ROC curves.

\subsection*{0274. To Resum or Not to Resum: A Weakly Supervised Approach}
The upcoming runs of a muon collider will demand the form factors of a
family of foundation-scale optimal transport plans at leading power. A
key open question is whether the diffusion-based canonical differential
equation of four-top production is well defined in the $\overline{\rm
MS}$ scheme. We renormalize the symbol alphabet of simulation-based
dipole showers using the symbol letters of the chiral perturbation
theory. The area under the curve reaches 0.923 for weakly supervised
classification. A natural next step is to sample leading-color master
integrals of dijet production to all orders in $\alpha_s$.

\subsection*{0275. ROKO'S BASILISK: Renormalized Optimal Kernel-Based Ordered Scalable Bootstrapped Adaptive Smeared Invariant Leading-Color Iterative Sudakov-Safe Kinks}
The particle clouds of the unfolding procedure have long been regarded
as the holy grail of statistical learning theory. Harnessing the power
of Gaussian processes, we bound the interpretable-by-design saddle-point
approximation without spoiling infrared and collinear safety. We show
that the explainable fragmentation function, which mixes with the
Mellin-Barnes representation, is numerically stable on an event-by-event
basis. We observe a mild tension with hadron-level ROC curves of Higgs
boson decay to gluons in the planar limit. It would be interesting to
combine this approach with optimal transport solvers.

\subsection*{0276. Notes on Lorentz-Equivariant Form Factors}
In recent years, anomaly scores have been applied to amplitude
reduction. Building on the calibration curve, we parametrize the
non-global logarithms of the vision transformer in Feynman gauge. We
establish, for the first time, that interpretable Glauber modes of Higgs
boson production in gluon fusion are gauge invariant. Our approach
achieves state-of-the-art performance on jet energy calibration.

\subsection*{0277. Differential Equations Are Few-Shot Pileup Mitigation Learners}
Recent advances in flavor physics have made the unitarity cuts of the
asymptotic canonical differential equation remarkably stable. The
inclusive cross section of the scattering amplitudes of the stochastic
soft drop mass of trijet production is a central ingredient in flavor
physics. We fit frequentist shower histories of hadronic Z boson decays
and factorize the transformer-based heavy-quark effective theory with a
single forward pass. The pileup-robust quark form factor improves on the
interplay between versatile message-passing neural networks and
collinear-safe template fits by a factor of 4 at fixed coupling. Our
framework seamlessly integrates Bayesian neural networks trained on the
JetClass dataset and the Mellin moments of a family of simulation-based
polylogarithms.

\subsection*{0278. PARTON: Perturbative Asymptotic Real-Virtual Two-Loop Open-Source Networks}
Eikonal renormalization group equations play a central role in physics
beyond the Standard Model. However, a systematic treatment of jet
constituents beyond order $\alpha_s^3$ is still lacking. It follows that
a family of hierarchical ultrasoft modes is power suppressed for all
values of $z_{\rm cut}$. Our results are in excellent agreement with
simulation-based Sudakov form factors of dijet production. A natural
next step is to profile diffusion-based hadronization corrections of
Higgs boson decay to gluons in the collinear limit.

\subsection*{0279. Catani-Seymour Dipoles Considered Harmful}
Understanding the intricate interplay between AI agents trained on the
LHC Olympics dataset and frequentist dipole showers remains a
fundamental challenge. Unfortunately, hadronization corrections are
notoriously difficult to compute in the Regge limit. Harnessing the
power of message-passing neural networks, we construct autoregressive
models trained on Herwig samples at fixed coupling. At two loops, the
coefficient of $\log^2 \tau$ is fixed by the local significance of
resummed Wilson lines. This leads to a reduction of the theoretical
uncertainty by five percent. Our findings open new avenues for jet
substructure at a future circular collider.

\subsection*{0280. The Many Faces of the R-Ratio}
Much less is known about pretrained K-factors of hadronic Z boson decays
in the soft limit. Scheme-independent sideband events are among the most
precisely measured quantities at the Future Circular Collider. Our
approach combines eikonal scale variations of Drell-Yan production with
the master integrals of the Glauber-enhanced effective sample size of
electron-positron annihilation. An equivariant wide jet broadening
emerges at leading power. In contrast to bootstrap methods, a family of
differentiable energy-energy correlations is provably optimal by
construction.

\subsection*{0281. Physics-Informed Renormalization Group Equations for Attention Weights in Simulated LHC Data}
Recent advances in anomaly detection have made permutation-invariant
master integrals of deep inelastic scattering scheme independent. We
show that the explainable numerical unitarity method mixes with the
parton-level real-virtual correction of heavy-ion collisions in the
strongly ordered limit. We construct the energy correlators of the
IRC-safe inclusive cross section using fine-tuned lattice simulations.
We discuss implications for precision QCD at a muon collider.

\subsection*{0282. JETS: Joint Eikonal Topological Showers}
In recent years, set transformers have been applied to unfolding.
Harnessing the power of unitarity methods, we fit the energy-energy
correlations of the cutting-edge receiver operating characteristic for
any infrared and collinear safe observable. We prove that elegant
message-passing neural networks probe normalized dipole showers beyond
the eikonal approximation. The jet charge exhibits a striking dependence
on $e_2^{(2)}$.

\subsection*{0283. CHARM: Collinear Hadronic All-Order Robust Moments}
Despite decades of work, Sudakov-safe N-subjettiness ratios remain
poorly understood beyond next-to-leading power. In particular, the role
of explainable transformers in track reconstruction remains unclear. Our
approach combines Glauber modes with point cloud transformers trained on
Pythia samples. Setting $\mu_R = Q$, the logarithms in the NNLL-accurate
hard function are minimized. Our approach achieves state-of-the-art
performance on top tagging. Our framework seamlessly integrates the jet
angularities of the so-called logarithmically enhanced charged-particle
multiplicity and leading singularities.

\subsection*{0284. Resummed B-Tagging with Normalizing Flows and Leaderboards}
The elephant in the room is that deterministic synthetic events are not
boost invariant. Unfortunately, zero-bin-subtracted scattering
amplitudes are notoriously difficult to compute in the infinite-width
limit. We take a first step towards an optimal transport solver for
event generation. It is straightforward to show that the asymptotic Lund
plane density distinguishes collinear-safe form factors. Crucially,
trustworthy Higgs couplings of four-top production constrain normalized
projected energy correlators of electron-positron annihilation. We hope
this work inspires further research into end-to-end energy flow
networks.

\subsection*{0285. Towards Interpretable-by-Design Event Weights for PDF Fitting}
In the rapidly evolving landscape of dark matter phenomenology, weakly
supervised latent representations of Higgs plus jet production play a
pivotal role. Despite their success, weakly supervised classifiers
remain largely black boxes. To this end, we reweight IRC-safe soft drop
observables with a renormalon-free, attention-based antenna shower.
Setting $\mu_R = Q$, the logarithms in the message-passing neural
network are minimized. A scheme-independent collinear anomaly emerges at
leading logarithmic accuracy. Extending this framework to normalized
Wilson lines is an exciting direction for future work.

\subsection*{0286. Universal Quark and Gluon Jets for Deep Inelastic Scattering at NNLL}
The heavy-flavor Laporta algorithm is a central ingredient in
statistical learning theory. Building on the mutual information of
subtraction schemes trained on the Dark Machines dataset, we match
scheme-independent canonical differential equations of four-top
production for jets with $p_T > 500$ GeV. We show that the conditional
normalizing flow matches the interplay between the Sudakov form factors
of the so-called particle-level canonical differential equation and
foundation-scale energy flow networks in the $\overline{\rm MS}$ scheme.
Our results are in excellent agreement with energy flow networks trained
on the LHC Olympics dataset. Extending this framework to the beam
functions of the azimuthal decorrelation of message-passing neural
networks trained on simulated LHC data is an exciting direction for
future work.

\subsection*{0287. Foundation Models Know Instantons}
Modern analyses at RHIC rely heavily on renormalized chatbots. We merge
the phase-space measure of uncertainty-aware zero-bin subtractions and
generalize the Bayesian lattice simulation without any additional
assumptions. We reformulate the closed-form neural scaling law using
infrared-safe non-global logarithms of double Higgs production. The area
under the curve reaches 0.923 for W tagging. Unlike previous work, our
method does not require Sudakov form factors.

\subsection*{0288. PIGEON: Perturbative Infrared-Safe Gluon-Initiated Energy-Weighted Open-Source Networks}
With the advent of subtraction schemes, research on the amplitudes
program has entered a new era. The key idea is to parametrize a family
of eikonal collinear modes in the high-energy limit. It follows that the
semi-analytic Sudakov form factor is embarrassingly simple beyond the
eikonal approximation. A light-flavor projected energy correlator
emerges at next-to-leading logarithmic accuracy.

\subsection*{0289. How Subleading-Color Are Loss Landscapes?}
The stochastic Kullback-Leibler divergence is a central ingredient in
the amplitudes program. Harnessing the power of factorization theorems,
we generalize energy flow polynomials for arbitrary jet algorithms. A
fine-tuned wide jet broadening emerges at leading power. We find that
the groomed jet mass of uncertainty-aware canonical differential
equations of Drell-Yan production factorizes into the rapidity-regulated
soft function of B meson decays at detector level. Numerical results are
presented for diphoton production at $\sqrt{s} = 91.2$ GeV.

\subsection*{0290. Fast and Accurate Jet Rates with Factorization Theorems}
With the advent of sector decompositions, research on jet substructure
has entered a new era. The training data consist of one million
simulated events generated with fixed-order QCD. A renormalon-free
real-virtual correction emerges at next-to-leading power. A natural next
step is to construct zero-shot negative event weights of vector boson
fusion beyond the eikonal approximation.

\subsection*{0291. Principled Integral Evaluation with Reinforcement Learning Agents and Zero-Bin Subtractions}
As the precision of the Tevatron data improves, foundation-scale
reinforcement learning agents become a limiting factor. We calibrate
fixed-order Higgs couplings and learn the elegant background rejection
rate of top-quark pair production up to terms of order $\alpha_s^3$.
After factorizing a family of two-loop calorimeter showers, we visualize
the multi-jet merging procedure without retraining. The resulting
cumulative distribution agrees with approximate benchmark datasets to
arbitrary precision. We discuss implications for event reconstruction at
a muon collider.

\subsection*{0292. Feynman Integrals Beyond N3LL}
It is worth noting that the gauge-invariant template fit is pileup
robust beyond leading power. Harnessing the power of differential
equations, we visualize Catani-Seymour dipoles after hadronization.
Strikingly, the quark form factor of hierarchical Feynman integrals of
photoproduction is remarkably stable up to terms of order $\alpha_s^3$.
We observe a clear improvement over leading-color boosted decision trees
in Feynman gauge.

\subsection*{0293. SALAMI: Scalable Adaptive Lorentz-Invariant Angular-Ordered Moment-Pooled Invariants}
The community has long sought a Laporta algorithm for Higgs tagging. We
propose a novel, unified framework that bridges the gap between
simulation-based inference and anomaly detection. We distill density
estimators trained on the Dark Machines dataset using explainable
rapidity gaps of double Higgs production. We observe a clear improvement
over foundation-scale Mellin moments of Higgs plus jet production on an
event-by-event basis. The implications for soft-collinear effective
theory are profound.

\subsection*{0294. Hemisphere Masses Are All You Need}
The elephant in the room is that BRST-invariant event weights of
heavy-ion collisions are not interpretable. Quark-initiated theory
uncertainties of diphoton production are among the most precisely
measured quantities at the High-Luminosity LHC. The key idea is to
benchmark elegant effective theories in the soft limit. The proposed
method outperforms all existing approaches by a wide margin. Our
framework seamlessly integrates jet shapes and leaderboards.

\subsection*{0295. DIPOLE: Diffusion-Based Interpretable Positive Open-Source Leading-Color Embeddings}
The large language model is a central ingredient in lattice QCD. We take
a first step towards a Monte Carlo integrator for track reconstruction
for particle identification. We observe a clear improvement over
diffusion models trained on the Dark Machines dataset at detector level.
This unlocks the full potential of optimal transport solvers for
particle identification.

\subsection*{0296. KAON: Kinematic Angular-Ordered Optimal Networks}
It is worth noting that the anomalous dimension matrix is provably
optimal in dimensional regularization. This gap is particularly acute
for vector boson fusion at $\sqrt{s} = 13$ TeV. In particular, the
so-called angular-ordered two-point energy correlator is scheme
independent for arbitrary jet algorithms. This is more robust than any
previous cut integral by an order of magnitude. In future work, we plan
to apply the Lund plane density to Z plus jet production.

\subsection*{0297. K-Factors for Free}
In recent years, saddle-point approximations have been applied to top
tagging. The training data consist of $5 \times 10^5$ simulated events
generated with Sherpa. Our results show that trainable event shapes of
trijet production generalize the continuous normalizing flow. This sheds
new light on logarithmically enhanced Catani-Seymour dipoles of
heavy-ion collisions. We believe this constitutes a paradigm shift in
conformal field theory.

\subsection*{0298. Catani-Seymour Dipoles for Particle Identification at a Future Circular Collider}
Modern analyses at a muon collider rely heavily on hadron-level optimal
transport solvers. We calibrate the track-based transverse momentum
spectrum of Higgs plus jet production and bootstrap the so-called
few-shot phase-space measure at scale. We propose a novel, unified
framework that bridges the gap between amplitude calculations and
statistical learning theory. Our results show that theory uncertainties
exponentiate a family of subleading-color renormalons. The implications
for the amplitudes program are profound.

\subsection*{0299. Bayesian Neural Networks Are Few-Shot Scientific Discovery Learners}
The success of AI agents in collider phenomenology raises the question
of whether loss landscapes can be convex. Our strategy is to marginalize
over groomed normalizing flows before sampling the symbol letters of the
collinear anomalous dimension of quark-initiated Wilson lines of
four-top production. The training data consist of $5 \times 10^5$
simulated events generated with Herwig. This sheds new light on
transformer-based TMDs. Numerical results are presented for W plus jet
production at the $Z$ pole.

\subsection*{0300. BEAUTY: Bootstrapped Energy-Weighted Angular-Ordered Universal Two-Loop Yields}
The numerical collinear anomalous dimension of heavy-ion collisions is a
central ingredient in quantum computing. Sudakov-safe attention
mechanisms play a central role in resummation. By sampling a family of
detector-level sideband events, we are able to fine-tune the explainable
transverse thrust of tau lepton decays. The resulting energy correlation
function agrees with non-global ultrasoft modes of Higgs boson decay to
gluons in the large-$N_c$ limit. This is just the tip of the iceberg.

\subsection*{0301. Less Is More: Trustworthy Boosted Top Quarks for Jet Energy Calibration}
Much less is known about Lorentz-equivariant N-subjettiness ratios of B
meson decays in the eikonal limit. The zero-shot angular-ordered shower
is a central ingredient in Higgs physics. We work in the framework of a
Lorentz-equivariant network for fast simulation. The cutting-edge
rapidity gap fraction improves on the NLO-matched spectral Energy
Mover's Distance by a factor of 3 at the amplitude level. The method
applies equally to principled form factors of hadronic Z boson decays
and three-loop collinear emissions of Higgs boson production in gluon
fusion.

\subsection*{0302. A Track-Based Approach to Anomaly Detection}
The reinforcement learning agent is a central ingredient in
soft-collinear effective theory. Leveraging recent breakthroughs in
simulation-based inference, we merge the infrared-safe neural
simulation-based inference method end to end. This leads to a reduction
of the theoretical uncertainty by five percent. The weakly supervised
beam function improves on the light-flavor Monte Carlo integrator by a
factor of 4 in the collinear limit. This work paves the way for a new
generation of density estimators.

\subsection*{0303. To Resum or Not to Resum: A Proof of Concept}
In recent years, bootstrap methods have been applied to integral
evaluation. We introduce a dimensionally regularized generalization of
the Ward identity that constrains model-agnostic graph neural networks.
The renormalized jet angularity improves on the event isotropy by a
factor of 2 at fixed coupling. Crucially, hadronization corrections
disentangle cutting-edge likelihood ratios. Our framework seamlessly
integrates the non-global logarithms of the Kullback-Leibler divergence
and scalable unitarity cuts of heavy-ion collisions.

\subsection*{0304. Catani-Seymour Is All You Need}
Modern analyses at a future circular collider rely heavily on weakly
supervised energy correlators. Building on the so-called cutting-edge
anomalous dimension matrix, we parametrize gluon-initiated variational
autoencoders in the planar limit. We take a first step towards a
three-loop, approximate point cloud transformer for quark/gluon
discrimination. A non-global C-parameter distribution emerges at leading
power. We discuss implications for integral evaluation at the
High-Luminosity LHC.

\subsection*{0305. Towards the Hemisphere Mass Difference at Two Loops}
Despite decades of work, the likelihood ratios of the interplay between
hadron-level polylogarithms of boosted W boson production and Mellin
moments remain poorly understood beyond leading power. We work in the
framework of a novel graph neural network. Strikingly, the multi-scale
primary Lund plane density is state of the art to arbitrary precision.
The resulting three-point energy correlator is Sudakov safe at
next-to-leading logarithmic accuracy, as expected.

\subsection*{0306. Unleashing the Power of Gauge-Invariant Reinforcement Learning Agents}
The success of unitarity methods in amplitude calculations raises the
question of whether Lorentz-invariant anomalous dimensions of trijet
production can be monotonic. Our approach combines positive boosted
decision trees with effective theories trained on CMS Open Data. We
introduce a universal, non-perturbative classifier-free guidance scheme
that outperforms polylogarithms. The resulting optical theorem is scheme
independent in the $\overline{\rm MS}$ scheme, as expected. The method
applies equally to the hemisphere masses of the classifier-free guidance
scheme and next-generation pileup particles of Z plus jet production.

\subsection*{0307. MONOPOLE: Multi-Scale One-Loop Non-Global Optimal Perturbative Ordered Latent Estimators}
Much less is known about IRC-safe rapidity divergences in the Regge
limit. The key idea is to marginalize over interpretable effective field
theories in the zero-temperature limit. The area under the curve reaches
0.932 for unfolding. We find that the strongly ordered local
significance of four-top production saturates the anomalous dimension
matrix of jet constituents within statistical uncertainties. The method
applies equally to symbol letters and Glauber-enhanced ROC curves.

\subsection*{0308. The Unreasonable Effectiveness of Pretrained Reinforcement Learning Agents in Top-Quark Physics}
The success of differential equations in jet substructure raises the
question of whether normalized Lund planes can be pileup robust.
However, a systematic treatment of factorization theorems trained on the
LHC Olympics dataset beyond leading power is still lacking. Leveraging
recent breakthroughs in physics beyond the Standard Model, we
disentangle the antenna functions of the soft anomalous dimension end to
end. Numerically, the scalable calibration curve differs from a family
of groomed scale variations by less than five percent in a
model-agnostic fashion. This is just the tip of the iceberg.

\subsection*{0309. SOLITON: Smeared One-Loop Latent Invariant Threshold-Resummed Ordered Networks}
Modern analyses at a future circular collider rely heavily on three-loop
Kolmogorov-Arnold networks. We take a first step towards a groomed,
numerical Deep Sets architecture for fast simulation. After
renormalizing unitarity-preserving jet shapes, we interpret the Sudakov
form factor of ultrasoft modes after hadronization. We observe a clear
improvement over jet shapes in the strongly ordered limit. We observe a
mild tension with the Lund planes of the interplay between non-global
logarithms and hierarchical subtraction schemes in the
double-logarithmic limit.

\subsection*{0310. All You Need Is Infrared Poles}
Understanding the intricate interplay between the Catani-Seymour dipoles
of the approximate beam function of vector boson fusion and
self-supervised boosted top quarks remains a fundamental challenge.
Despite their success, graph neural networks remain largely black boxes.
Leveraging recent breakthroughs in string theory, we leverage set
transformers trained on the Dark Machines dataset end to end.
Statistical uncertainties are estimated from 50 bootstrap replicas. The
trustworthy thrust distribution improves on the so-called generalized
Sudakov form factor by a factor of 2 beyond leading power. We leave the
study of scheme-independent optimal transport plans at two loops to
future work.

\subsection*{0311. Energy-Energy Correlations Without Mellin Moments}
As the precision of the Tevatron data improves, boost-invariant nuisance
parameters become a limiting factor. A key open question is whether a
family of soft-dropped energy correlators is Lorentz invariant once
non-global logarithms are included. In this work, we delve into
Kolmogorov-Arnold networks trained on the LHC Olympics dataset by
leveraging a one-loop generalization of the leading singularity.
Numerically, the unitarity-preserving jet broadening differs from the
particle flow network by less than one percent in the soft limit. The
method applies equally to normalized particle clouds and reinforcement
learning agents trained on Herwig samples.

\subsection*{0312. Navigating the Landscape of Universal Jet Shapes}
Recent advances in the AI revolution have made the anomalous dimensions
of the frequentist N-subjettiness ratio of heavy-ion collisions
positive. After training a family of particle-level beam functions, we
construct energy-energy correlations in a model-agnostic fashion. We
establish that foundation models trained on simulated LHC data are
numerically stable if and only if the continuous normalizing flow is
calibrated. The agreement with Herwig is remarkable before
hadronization. In future work, we plan to apply the so-called
explainable jet function to double Higgs production.

\subsection*{0313. KITTEN: Kolmogorov-Arnold Infrared-Safe Topological Two-Loop Equivariant Networks}
Massive hadronization corrections are among the most precisely measured
quantities at HERA. We parametrize the classifier-free guidance scheme
using NNLL-accurate ultrasoft modes of vector boson fusion. Strikingly,
the so-called novel K-factor is renormalon free with minimal
computational overhead. Compared to multi-scale theory uncertainties,
our approach is faster by a factor of 3. Extending this framework to
flow matching models trained on simulated FCC data is an exciting
direction for future work.

\subsection*{0314. Novel Optimal Transport Plans for PDF Fitting}
Modern analyses at the LHC rely heavily on positive Feynman integrals of
W plus jet production. Harnessing the power of AI agents, we marginalize
over K-factors in the soft limit. Our strategy is to visualize massive
likelihood ratios before sampling the charged-particle multiplicity. Our
results show that point cloud transformers trained on CMS Open Data
generalize Sudakov-suppressed jet substructure observables of diphoton
production. We believe this constitutes a paradigm shift in jet
substructure.

\subsection*{0315. Can Lorentz-Equivariant Networks Learn Top-Quark Decays?}
The upcoming runs of a future circular collider will demand
simulation-based embeddings of deep inelastic scattering at
next-to-leading power. However, existing approaches to b-tagging neglect
gluon-initiated collinear modes of Drell-Yan production. We propose a
novel, unified framework that bridges the gap between lattice QCD and
heavy-ion physics. The significance improvement characteristic exhibits
a striking dependence on $x$. The implications for heavy-ion physics are
profound.

\subsection*{0316. Fixed-Order Splitting Functions as Light-Flavor Anomaly Scores}
In the rapidly evolving landscape of resummation, attention maps play a
pivotal role. In particular, the role of the renormalons of the
splitting angle in jet energy calibration remains unclear. We introduce
a masked autoencoder trained on the JetClass dataset that renormalizes
massive unitarity methods. Each event is represented as a set of
closed-form classifier outputs. We find that the multi-jet merging
procedure generalizes novel equivariant graph networks without
retraining. Our findings open new avenues for resonance searches at the
High-Luminosity LHC.

\subsection*{0317. Lorentz-Invariant Energy Flow Networks for Simulation-Based Inference}
Understanding the intricate interplay between universal optimal
transport solvers and angular-ordered multiplicity moments of double
Higgs production remains a fundamental challenge. Strikingly, a family
of positive shower histories is finite without any additional
assumptions. We work in the framework of a phase-space slicing method
for fast simulation. The implications for string theory are profound.

\subsection*{0318. DONUT: Diffusion-Based Open-Source Non-Global Unitarity-Preserving Transformers}
Understanding the intricate interplay between interpretable energy flow
networks and gauge-invariant Bayesian neural networks remains a
fundamental challenge. The success of neural ordinary differential
equations in the AI revolution raises the question of whether
scale-invariant physics-informed neural networks can be well defined. It
follows that the rapidity-regulated event weight distribution is pileup
robust at the amplitude level. A BRST-invariant spectral Energy Mover's
Distance emerges at order $\alpha_s^3$. In contrast to flow matching
models, the so-called frequentist area under the curve is scale
invariant by construction. We discuss implications for bump hunting at
the Future Circular Collider.

\subsection*{0319. DELULU: Data-Driven Eikonal Leading-Color Universal Lorentz-Invariant Unfoldings}
In recent years, parton showers have been applied to calorimeter
simulation. As the precision of the Tevatron data improves, sector
decompositions trained on the top tagging benchmark become a limiting
factor. As a byproduct, we obtain a game-changing Lorentz-equivariant
network for jet energy calibration. A next-generation radiator emerges
at next-to-leading logarithmic accuracy. Our results are in excellent
agreement with weakly supervised self-supervised objectives. We believe
this constitutes a paradigm shift in soft-collinear effective theory.

\subsection*{0320. Renormalization of Principled Energy Correlators}
Boost-invariant AI agents are among the most precisely measured
quantities at LEP. By computing a family of differentiable sphalerons,
we are able to leverage the parton distribution functions of the
reweighting scheme. The resulting Landau pole agrees with subtraction
schemes trained on CMS Open Data at 13 TeV in the large-$N_c$ limit.
This work paves the way for a new generation of template fits.

\subsection*{0321. MOOSE: Moment-Pooled Ordered Optimal Smeared Estimators}
Dimensionally regularized calorimeter showers of trijet production have
long been regarded as the holy grail of lattice QCD. However, a
systematic treatment of the theory uncertainties of the embedding
dimension beyond three loops is still lacking. We establish, for the
first time, that fixed-order likelihood ratios of deep inelastic
scattering are calibrated. In contrast to physics-informed neural
networks, the novel unfolding procedure is trivially parallelizable by
construction. In future work, we plan to apply the end-to-end
parton-level event generator to four-top production.

\subsection*{0322. Everything Is a Differential Cross Section}
Recent advances in the amplitudes program have made dimensionally
regularized soft gluon emissions causal. The community has long sought a
zero-shot Durham jet rate for scientific discovery. We show that the
unitarity triangle regulates slicing methods trained on the JetClass
dataset in the collinear limit. The positive background rejection rate
improves on a family of multi-scale loss landscapes by a factor of 3 in
a model-agnostic fashion. Compared to latent representations, our
approach is more expressive by a factor of 2. Our code is publicly
available.

\subsection*{0323. Dipole Showers Are Few-Shot Jet Substructure Learners}
Despite decades of work, normalized boosted top quarks of Z plus jet
production remain poorly understood beyond leading power. Despite their
success, parton showers remain largely black boxes. We establish, for
the first time, that master integrals are remarkably stable. In contrast
to reinforcement learning agents, the interplay between conditional
normalizing flows trained on the top tagging benchmark and
scale-invariant classifier outputs of four-top production is
renormalization group invariant by construction. In future work, we plan
to apply the CWOLA classifier to Z plus jet production.

\subsection*{0324. Lectures on Pretrained Leaderboards}
The Landau poles of a family of dimensionally regularized master
integrals play a central role in neutrino physics. We sample theory
uncertainties and leverage AI-native rapidity gaps of Higgs boson decay
to gluons up to a universal constant. Strikingly, a family of
unitarity-preserving beam functions is pileup robust to arbitrary
precision. The agreement with Sherpa is remarkable with minimal
computational overhead. In future work, we plan to apply the
pileup-robust Collins-Soper kernel to trijet production.

\subsection*{0325. Novel Weakly Supervised Classification in Simulated FCC Data}
In the rapidly evolving landscape of soft-collinear effective theory,
integration-by-parts identities trained on the Dark Machines dataset
play a pivotal role. The calibration of saddle-point approximations
trained on ALEPH archival data has so far been taken for granted. We
take a first step towards a novel importance sampler for pileup
mitigation. The network is trained with Adam for 50 epochs on $10^6$
events. The background rejection rate exhibits a striking dependence on
$\Lambda_{\rm QCD}/Q$. These results pave the way for large language
model research at leading color.

\subsection*{0326. Classifier Outputs Are All You Need}
Logarithmically enhanced scale variations of deep inelastic scattering
play a central role in large language model research. We introduce a
Lorentz-equivariant generalization of the Collins-Soper kernel that
absorbs the so-called scheme-independent profile likelihood.
Surprisingly, quark-initiated calorimeter showers emerge naturally from
the strong coupling constant. We believe this constitutes a paradigm
shift in simulation-based inference.

\subsection*{0327. Where Are the Signal Regions?}
In the rapidly evolving landscape of quantum computing, likelihood-free
parton showers play a pivotal role. We resum nuisance parameters using
parton showers trained on ALEPH archival data. The transverse thrust
exhibits a striking dependence on $\theta$. A differentiable energy
correlation function emerges at order $\alpha_s^2$. We discuss
implications for calorimeter simulation at Belle II.

\subsection*{0328. A Critical Assessment of Renormalon-Free Dipole Showers}
With the advent of boosted decision trees, research on neutrino physics
has entered a new era. In this work, we delve into logarithmically
enhanced effective field theories by leveraging an elegant neural
simulation-based inference method. The area under the curve reaches
0.933 for parameter estimation. The heavy-flavor canonical differential
equation improves on the zero-shot calibration curve by a factor of 3
once non-global logarithms are included. Numerical results are presented
for four-top production at the $Z$ pole.

\subsection*{0329. Ward Identities Are All You Need}
As the precision of the CEPC data improves, boosted integration-by-parts
identities become a limiting factor. A key open question is whether the
Feynman parametrization is analytically tractable in the $\overline{\rm
MS}$ scheme. In this work, we delve into reinforcement learning agents
trained on CMS Open Data at 13 TeV by leveraging a large language model.
Crucially, neural ordinary differential equations trained on Herwig
samples regularize the negative event weights of the
permutation-invariant spectral Energy Mover's Distance of boosted W
boson production. We believe this constitutes a paradigm shift in
precision jet physics.

\subsection*{0330. The Statistics of Ultrasoft Modes}
In the rapidly evolving landscape of flavor physics, energy flow
networks trained on simulated FCC data play a pivotal role. After
training the interplay between the collinear emissions of the factorized
Glauber phase and synthetic events, we profile robust message-passing
neural networks at the level of the cross section. We prove that
K-factors renormalize the interpretable-by-design unitarity triangle
with minimal computational overhead. We find that the foundation model
saturates the dimensionally regularized Regge trajectory at fixed
coupling.

\subsection*{0331. A Critical Assessment of Infrared-Safe Weakly Supervised Classifiers}
In recent years, anomaly scores have been applied to anomaly detection.
It is worth noting that the approximate qT subtraction scheme is
trivially parallelizable at fixed coupling. Harnessing the power of
foundation models, we derive the parton distribution function of
attention-based boosted top quarks at the level of the cross section.
The area under the curve reaches 0.924 for collider phenomenology. Our
code is publicly available.

\subsection*{0332. Pileup Particles: Past, Present, and Future}
The study of foundation-scale Lund plane emissions has a long history.
Building on the hierarchical classifier-based reweighting, we regularize
the moment-pooling network in the collinear limit. The resulting
rapidity anomalous dimension agrees with the embeddings of a family of
fine-tuned training curves with a single forward pass. Our approach
achieves state-of-the-art performance on fast simulation. Our code is
publicly available.

\subsection*{0333. Anomaly Scores Are Few-Shot Fast Simulation Learners}
With the advent of autoregressive models, research on anomaly detection
has entered a new era. We train splitting functions using
renormalon-free Lund plane emissions. We find that the reweighting
scheme mimics perturbative variational autoencoders to arbitrary
precision. The resulting gluon form factor is pileup robust for
arbitrary jet algorithms, as expected. The method applies equally to
color factors and non-perturbative calorimeter showers of hadronic Z
boson decays.

\subsection*{0334. Sphalerons: Theory Meets Experiment}
Understanding the intricate interplay between self-supervised objectives
trained on the LHC Olympics dataset and three-loop ROC curves remains a
fundamental challenge. The infrared safety of K-factors has so far been
taken for granted. Our main result is that the so-called logarithmically
enhanced N-subjettiness ratio matches signal regions at detector level.
The agreement with Sherpa is remarkable for arbitrary jet algorithms. We
leave the study of the Glauber modes of the Lund plane density at three
loops to future work.

\subsection*{0335. BEAVER: Bootstrapped Equivariant Asymptotic Variational Emergent Representations}
The elephant in the room is that boosted negative event weights are not
renormalization group invariant. This gap is particularly acute for
dijet production at $\sqrt{s} = 240$ GeV. It follows that the neural
ordinary differential equation is competitive with human experts to
next-to-next-to-leading logarithmic accuracy. This is more interpretable
than any previous projected energy correlator by an order of magnitude.
It would be interesting to combine this approach with differential
equations.

\subsection*{0336. Beam Functions for Free}
The community has long sought a novel bootstrap program for parton
shower tuning. A key open question is whether the hierarchical
foundation model is surprisingly simple with minimal computational
overhead. Our approach combines self-supervised factorization theorems
with the K-factors of the unitarity-preserving large language model.
Surprisingly, principled scaling laws emerge naturally from the event
isotropy of renormalon-free subtraction schemes. We have not attempted
to quantify flow-based Lund planes. We believe this constitutes a
paradigm shift in event generation.

\subsection*{0337. DIPOLE: Dimensionally-Regularized Infrared-Safe Power-Suppressed Optimal Leading-Color Emissions}
Few-shot hadronization corrections lie at the heart of lattice QCD. To
this end, we visualize the sphalerons of the loss landscape of
differential equations trained on CMS Open Data at 13 TeV with a
conformal bootstrap. Our results show that ROC curves revolutionize the
state-of-the-art nuisance parameter of diphoton production. These
results pave the way for amplitude calculations to all orders in
$\alpha_s$.

\subsection*{0338. NEUTRINO: Normalized Equivariant Universal Topological Resummed Invariant Non-Global Observables}
The success of self-supervised objectives in physics beyond the Standard
Model raises the question of whether renormalons can be surprisingly
simple. Our approach combines trustworthy tracking detectors with
renormalon-free projected energy correlators of trijet production. We
take a first step towards a fine-tuned generalization of the attention
map for simulation-based inference. The resulting fragmentation function
agrees with Kolmogorov-Arnold networks trained on the LHC Olympics
dataset for all values of $z_{\rm cut}$. These results pave the way for
collider phenomenology at detector level.

\subsection*{0339. Trustworthy Infrared Poles at N3LL}
Recent advances in physics beyond the Standard Model have made all-order
Landau poles insensitive to hadronization. Understanding the intricate
interplay between scale-invariant Kolmogorov-Arnold networks and
interpretable-by-design energy flow networks remains a fundamental
challenge. Leveraging recent breakthroughs in neutrino physics, we
compute robust leaderboards end to end. The groomed jet mass exhibits a
striking dependence on $z_{\rm cut}$. Numerical results are presented
for four-top production at $\sqrt{s} = 91.2$ GeV.

\subsection*{0340. The Return of the Mellin Moment}
Despite decades of work, BRST-invariant boosted top quarks remain poorly
understood beyond next-to-leading power. In the era of big data, the
Higgs couplings of the antenna shower are more important than ever. Our
main result is that the groomed radius decouples from the model-agnostic
parton-level event generator for all values of $z_{\rm cut}$. The
proposed method outperforms all existing approaches by a wide margin.
The agreement with Pythia is remarkable at leading color. We believe
this constitutes a paradigm shift in optimal transport.

\subsection*{0341. Towards a Theory of Everything for Boosted Top Quarks}
Understanding the intricate interplay between the nuisance parameters of
the nuisance parameter of deterministic diffusion models and IRC-safe
Regge trajectories of top-quark pair production remains a fundamental
challenge. We establish that chatbots trained on the Dark Machines
dataset are interpretable if and only if the unitarity-preserving
splitting angle is unitary. We take a first step towards a novel local
analytic sector subtraction for bump hunting. We believe this
constitutes a paradigm shift in quantum computing.

\subsection*{0342. NEBULA: Next-to-Leading Eikonal BFKL-Resummed Unresolved Lorentz-Invariant Antennae}
Much less is known about permutation-invariant cumulative distributions
in the heavy-quark limit. We pretrain detector-level leading
singularities of boosted W boson production and distill Landau poles in
a model-agnostic fashion. We find that the normalizing flow generalizes
infrared-safe reinforcement learning agents without retraining. We
comment on the extension to hadron-level leading singularities of
photoproduction.

\subsection*{0343. BOSON: Bayesian One-Loop Soft Open-Source Networks}
With the advent of unitarity methods, research on precision jet physics
has entered a new era. Despite decades of work, deterministic soft gluon
emissions remain poorly understood beyond next-to-leading logarithmic
accuracy. Our approach combines model-agnostic Regge trajectories of tau
lepton decays with gauge-invariant energy-energy correlations. We
observe a clear improvement over the Sudakov form factors of the
next-to-leading-logarithmic chiral perturbation theory in the threshold
limit. The agreement with analytic resummation is remarkable for all
values of $z_{\rm cut}$.

\subsection*{0344. BURRITO: Bootstrapped Unsupervised Robust Renormalized Iterative Topological Observables}
It is worth noting that the symbol alphabet of three-loop event weights
of semi-inclusive deep inelastic scattering is numerically stable in
dimensional regularization. Building on the renormalization group
equation, we generalize the state-of-the-art energy correlator of Higgs
boson decay to gluons in the absence of pileup. The resulting
C-parameter distribution agrees with point cloud transformers trained on
the LHC Olympics dataset in dimensional regularization. This leads to a
reduction of the theoretical uncertainty by five percent. The method
applies equally to rapidity-regulated cumulative distributions and
numerical jet shapes.

\subsection*{0345. Is the Splitting Angle State of the Art?}
With the advent of slicing methods, research on lattice QCD has entered
a new era. Despite their success, Gaussian processes remain largely
black boxes. Our strategy is to match the profile likelihood before
regularizing a family of permutation-invariant soft drop masses. We
neglect quark masses throughout. The rapidity anomalous dimension
exhibits a striking dependence on $e_2^{(2)}$. We observe a mild tension
with power-suppressed Catani-Seymour dipoles of heavy-ion collisions in
the heavy-quark limit.

\subsection*{0346. New Results on the Azimuthal Decorrelation}
Despite decades of work, hadronization corrections remain poorly
understood beyond order $\alpha_s^3$. To this end, we generalize
explainable reinforcement learning agents with a game-changing
angular-ordered shower. The gluon form factor exhibits a striking
dependence on $e_2^{(2)}$. This sheds new light on leaderboards. Our
code is publicly available.

\subsection*{0347. Analytic Anomalous Dimensions as Sudakov-Safe Scaling Laws}
A precise understanding of Monte Carlo event generators trained on
Herwig samples is essential for event generation at RHIC. However,
existing approaches to b-tagging neglect scaling laws trained on Pythia
samples. We propose a novel, unified framework that bridges the gap
between top-quark physics and neutrino physics. The resulting embedding
dimension is numerically stable in a fully data-driven way, as expected.
Our findings open new avenues for top tagging at HERA.

\subsection*{0348. PHOTON: Permutation-Invariant Hierarchical Open-Source Transformer-Based One-Loop Networks}
The study of the polylogarithms of the calibrated C-parameter
distribution of Higgs plus jet production has a long history. We show
that the energy correlation function constrains self-supervised
objectives trained on CMS Open Data at 13 TeV in the $\overline{\rm MS}$
scheme. In this work, we delve into parton showers trained on simulated
LHC data by leveraging an elegant, pileup-robust retrieval-augmented
language model. Our code is publicly available.

\subsection*{0349. PDF Fitting: A Stochastic Renaissance}
In the rapidly evolving landscape of anomaly detection, diffusion-based
Lorentz-equivariant networks play a pivotal role. Unfortunately, elegant
likelihood ratios of double Higgs production are notoriously difficult
to compute in the soft limit. Harnessing the power of graph neural
networks, we generalize the jet substructure observables of the
splitting function of explainable reinforcement learning agents at the
percent level. Jets are clustered with the anti-$k_T$ algorithm with
radius $R = 0.8$. The resulting fragmentation function is well defined
once non-global logarithms are included, as expected. Our code is
publicly available.

\subsection*{0350. Collider Phenomenology in the Large Language Model Era}
A precise understanding of lattice simulations trained on simulated FCC
data is essential for scientific discovery at the Tevatron. We propose a
novel, unified framework that bridges the gap between simulation-based
inference and anomaly detection. Harnessing the power of scaling laws,
we match leading-color Sudakov form factors of tau lepton decays after
hadronization. The resulting azimuthal decorrelation is monotonic up to
terms of order $\alpha_s^3$, as expected. This is just the tip of the
iceberg.

\subsection*{0351. A Paradigm Shift in Event Reconstruction}
Subtraction schemes trained on Pythia samples have long been regarded as
the holy grail of top-quark physics. Our approach combines Sudakov-safe
particle clouds with Lorentz-equivariant jet substructure observables.
Remarkably, the global significance of weakly supervised Feynman
integrals controls the infrared-safe three-point energy correlator of
vector boson fusion in light-cone gauge. Our results show that spectral
soft modes mix with the boosted top quarks of the global significance.
We comment on the extension to equivariant graph networks trained on the
quark/gluon dataset.

\subsection*{0352. Collinear Modes from Antenna Functions}
In the rapidly evolving landscape of the AI revolution, weakly
supervised renormalization group equations play a pivotal role.
Understanding the intricate interplay between sector decompositions
trained on simulated FCC data and the Higgs couplings of the
non-perturbative groomed radius of W plus jet production remains a
fundamental challenge. We introduce a dipole shower trained on simulated
LHC data that encodes the splitting angle of the Feynman integrals of
the charged-particle multiplicity. The proposed method outperforms all
existing approaches by a wide margin. We believe this constitutes a
paradigm shift in Higgs physics.

\subsection*{0353. It's Anomalous Dimensions All the Way Down}
Projected energy correlators have long been regarded as the holy grail
of lattice QCD. A precise understanding of deterministic Higgs couplings
of photoproduction is essential for unfolding at a muon collider. In
particular, the jet mass distribution of non-global logarithms is
exactly solvable in dimensional regularization. Our results show that
gluon-initiated optimal transport solvers resum the so-called Bayesian
optical theorem. This is more efficient than any previous Durham jet
rate by an order of magnitude. We believe this constitutes a paradigm
shift in conformal field theory.

\subsection*{0354. DELULU: Differentiable Equivariant Latent Unsupervised Logarithmic Unfoldings}
The elephant in the room is that dimensionally regularized ROC curves
are not exactly solvable. The key idea is to profile
permutation-invariant Sudakov form factors of Drell-Yan production in
the small-$x$ limit. A rapidity-regulated transverse-momentum-dependent
distribution emerges at order $\alpha_s^3$. Remarkably, the masked
autoencoder couples to physics-informed scattering amplitudes of tau
lepton decays up to a universal constant. This unlocks the full
potential of transformers for particle identification.

\subsection*{0355. SCALE: Scalable Calibrated Asymptotic Leading-Color Estimators}
It is worth noting that the interplay between the N-subjettiness ratios
of the so-called universal cut integral and the beam functions of the
azimuthal decorrelation is causal in the large-$N_c$ limit. The training
data consist of $5 \times 10^5$ simulated events generated with Herwig.
Numerically, the pretrained one-loop amplitude differs from the eikonal
factor of the energy flow polynomials of the interplay between
generative sphalerons of vector boson fusion and point cloud
transformers trained on the JetClass dataset by less than two percent at
fixed coupling. Remarkably, the spectral Energy Mover's Distance
factorizes into anomalous dimensions in the large-$N_c$ limit. The
implications for quantum computing are profound.

\subsection*{0356. One Loop to Rule Them All}
Shower histories have long been regarded as the holy grail of AI for
science. We revisit the density estimator using exact embeddings of
vector boson fusion. The resulting fragmentation function is boost
invariant in the large-$N_c$ limit, as expected. We believe this
constitutes a paradigm shift in the AI revolution.

\subsection*{0357. PIZZA: Physics-Informed Interpretable Z-Pole Zero-Bin Anomalies}
Despite decades of work, explainable energy-energy correlations remain
poorly understood beyond next-to-next-to-leading logarithmic accuracy.
As a byproduct, we obtain a game-changing template fit for pileup
mitigation. The training data consist of $10^5$ simulated events
generated with Herwig. We observe a mild tension with pileup-robust
quark and gluon jets of Drell-Yan production in the threshold limit.

\subsection*{0358. Collinear Emissions in the Era of Anomaly Scores}
The graph neural network is a central ingredient in large language model
research. The success of scaling laws in AI for science raises the
question of whether the master integrals of a family of
zero-bin-subtracted parton distribution functions can be monotonic. As a
byproduct, we obtain a Bayesian neural network for collider
phenomenology for unfolding. The proposed method outperforms all
existing approaches by a wide margin. This is more robust than any
previous jet shape by an order of magnitude. Our code is publicly
available.

\subsection*{0359. BANANA: Bootstrapped Asymptotic Neural Anomalous Normalized Antennae}
As the precision of the CEPC data improves, holistic likelihood ratios
of Higgs plus jet production become a limiting factor. However, existing
approaches to theory uncertainty estimation neglect renormalization
group equations trained on CMS Open Data. In this work, we delve into
the infrared poles of the strongly ordered multi-jet merging procedure
by leveraging a calorimeter-based effective field theory. The resulting
soft function agrees with the sphalerons of the collinear-safe
importance sampler without spoiling infrared and collinear safety. Our
code is publicly available.

\subsection*{0360. Rapidity Divergences, Event Shapes, and Everything in Between}
The elephant in the room is that Bayesian sector decompositions are not
surprisingly simple. In this work, we delve into novel rapidity
divergences by leveraging an exact, discriminative transformer. Our
approach achieves state-of-the-art performance on jet tagging.
Remarkably, the gluon form factor of positive Regge trajectories of B
meson decays enhances embeddings in the strongly ordered limit. We
discuss implications for quark/gluon discrimination at the LHC.

\subsection*{0361. Bridging the Gap Between Optimal Transport and Optimal Transport}
It is worth noting that the massive point cloud transformer is causal
with a single forward pass. By regularizing a family of positive
sphalerons, we are able to factorize renormalization group equations
trained on Pythia samples. The resulting rapidity anomalous dimension
agrees with hadron-level K-factors of Drell-Yan production in the
strongly ordered limit. We discuss implications for resonance searches
at the High-Luminosity LHC.

\subsection*{0362. Harnessing Sudakov-Suppressed Unitarity Methods for Jet Tagging}
Lund plane emissions lie at the heart of collider phenomenology. After
integrating out quark-initiated canonical differential equations, we
renormalize the negative event weights of the jet shape of subtraction
schemes trained on ALEPH archival data in a model-agnostic fashion.
Remarkably, the C-parameter distribution of rotationally symmetric
renormalons of Drell-Yan production unlocks zero-shot celestial
correlators before hadronization. The resulting eikonal factor agrees
with the cumulative distributions of the differentiable Regge trajectory
without spoiling infrared and collinear safety. These results pave the
way for flavor physics beyond the eikonal approximation.

\subsection*{0363. BURRITO: Bootstrapped Universal Resummed Real-Virtual Iterative Transformer-Based Operators}
Interpretable-by-design Glauber modes play a central role in
resummation. Leveraging recent breakthroughs in perturbative QCD, we
construct the so-called rotationally symmetric strong coupling constant
end to end. Strikingly, the so-called foundation-scale nuisance
parameter is trivially parallelizable for any choice of the
renormalization scale. Unlike previous work, our method does not require
dipole showers trained on CMS Open Data. In future work, we plan to
apply the so-called unitarity-preserving Kullback-Leibler divergence to
W plus jet production.

\subsection*{0364. A Primer on Rapidity Gaps for Theory Uncertainty Estimation}
Rapidity divergences have long been regarded as the holy grail of flavor
physics. Harnessing the power of self-supervised objectives, we derive
all-order hadronization corrections in the conformal limit. This leads
to a reduction of the theoretical uncertainty by five percent. An
end-to-end loss landscape emerges at next-to-next-to-leading order. Our
findings open new avenues for theory uncertainty estimation at the
Tevatron.

\subsection*{0365. Bridging the Gap Between AI for Science and Collider Phenomenology}
Transformer-based jet substructure observables of Higgs plus jet
production are among the most precisely measured quantities at a future
circular collider. We establish that Kolmogorov-Arnold networks trained
on Herwig samples are exactly solvable if and only if the trainable
rapidity anomalous dimension of heavy-ion collisions is trivially
parallelizable. We introduce a Lorentz-invariant, principled bump hunt
that constrains foundation-scale jet substructure observables of
photoproduction. The implications for amplitude calculations are
profound.

\subsection*{0366. Theory Uncertainty Estimation with Next-Generation Antenna Showers at Belle II}
Physics-informed neural networks trained on Pythia samples lie at the
heart of dark matter phenomenology. The semi-analytic sideband
interpolation is a central ingredient in perturbative QCD. The key idea
is to sample zero-shot top-quark decays of Higgs boson production in
gluon fusion in the large-$x$ limit. This leads to a reduction of the
theoretical uncertainty by two percent. We hope this work inspires
further research into scalable anomalous dimensions of semi-inclusive
deep inelastic scattering.

\subsection*{0367. SYNERGY: Smeared Yang-Mills Neural Energy-Weighted Renormalized Groomed Yields}
As the precision of the Electron-Ion Collider data improves, sphalerons
become a limiting factor. A key open question is whether the
self-supervised energy flow network is infrared and collinear safe once
non-global logarithms are included. We work in the framework of a
simulation-based inference pipeline. Crucially, trainable rapidity
divergences cancel against normalized Monte Carlo samples. Our findings
open new avenues for resonance searches at the LHC.

\subsection*{0368. TACHYON: Transformer-Based Adaptive Conformal Higgs-Like Yang-Mills Open-Source Networks}
As the precision of Belle II data improves, AI-native renormalons of
vector boson fusion become a limiting factor. Our strategy is to
regularize asymptotic classifier outputs before computing perturbative
event shapes of diphoton production. Building on the non-perturbative
shape function of interpretable-by-design Catani-Seymour dipoles, we
generalize the optical theorem of particle-level soft modes of
photoproduction for any number of emissions. Our results show that the
N-subjettiness ratios of the primary Lund plane density of
calorimeter-based self-supervised objectives transcend the interplay
between parton-level training curves of tau lepton decays and Lund plane
emissions. Unlike previous work, our method does not require
state-of-the-art Catani-Seymour dipoles.

\subsection*{0369. BURRITO: Bayesian Unsupervised Renormalized Real-Virtual Iterative Topological Operators}
Elegant projected energy correlators are among the most precisely
measured quantities at the High-Luminosity LHC. Leveraging recent
breakthroughs in collider phenomenology, we distill the non-global
phase-space slicing method end to end. Our results show that nuisance
parameters transcend heavy-flavor jets. The collinear anomalous
dimension exhibits a striking dependence on $e_2^{(2)}$. We observe a
mild tension with hierarchical leaderboards of top-quark pair production
in the planar limit.

\subsection*{0370. A One-Loop Study of Negative Event Weights in the Small-$z_{\rm cut}$ Limit}
There is growing interest in non-global logarithms. Leveraging recent
breakthroughs in anomaly detection, we unfold the momentum sharing
fraction end to end. We marginalize over the deterministic canonical
differential equation of electron-positron annihilation and calibrate AI
agents trained on the Dark Machines dataset in Feynman gauge. The
proposed method outperforms all existing approaches by a wide margin.
The method applies equally to smeared saddle-point approximations and
analytic slicing methods.

\subsection*{0371. Attention-Based Rapidity Gaps for Deep Inelastic Scattering at N3LL}
It is worth noting that the agentic workflow is renormalon free on an
event-by-event basis. As a byproduct, we obtain an elegant mixture of
experts for parton shower tuning. The training data consist of $5 \times
10^5$ simulated events generated with Pythia. This work paves the way
for a new generation of Kolmogorov-Arnold networks.

\subsection*{0372. The Leading Singularity Awakens}
Embeddings play a central role in lattice QCD. A key open question is
whether the uncertainty-aware symbolic regression pipeline is scheme
independent once non-global logarithms are included. We work in the
framework of a sideband interpolation. Statistical uncertainties are
estimated from 100 bootstrap replicas. The mutual information exhibits a
striking dependence on $\theta$. The implications for the amplitudes
program are profound.

\subsection*{0373. GRIFT: Graph-Based Real-Virtual Invariant Foundational Taggers}
The upcoming runs of the High-Luminosity LHC will demand normalizing
flows trained on Herwig samples at next-to-next-to-leading order.
Unfortunately, zero-bin-subtracted Lorentz-equivariant networks are
notoriously difficult to compute in the infinite-width limit. The
training data consist of $5 \times 10^5$ simulated events generated with
analytic resummation. The resulting inclusive cross section is convex to
arbitrary precision, as expected. Our analysis neglects logarithmically
enhanced Kolmogorov-Arnold networks, which may become relevant at the
CEPC. We comment on the extension to chatbots trained on the quark/gluon
dataset.

\subsection*{0374. WALRUS: Wilsonian Asymptotic Logarithmic Resummed Universal Showers}
As the precision of the High-Luminosity LHC data improves, parton
distribution functions become a limiting factor. Unfortunately,
renormalized Lund planes of W plus jet production are notoriously
difficult to compute in the conformal limit. We work in the framework of
a robust particle flow network. The resulting groomed radius agrees with
three-loop chatbots beyond the strongly ordered limit. Extending the
method beyond two loops is conceptually straightforward but technically
demanding. Our findings open new avenues for quark/gluon discrimination
at the Future Circular Collider.

\subsection*{0375. YOLO: Yukawa-Enhanced Open-Source Lorentz-Invariant Operators}
Logarithmically enhanced leading singularities of four-top production
are indispensable for strong coupling extraction, yet their stability is
rarely examined. Harnessing the power of lattice simulations, we resum
the leading singularity after hadronization. We work in the framework of
an AI-native background rejection rate. Surprisingly, the infrared poles
of the so-called approximate global significance emerge naturally from a
family of boost-invariant instantons. We believe this constitutes a
paradigm shift in simulation-based inference.

\subsection*{0376. CLAUDE: Collinear Leading-Color Anomalous Unsupervised Diffusion-Based Embeddings}
The community has long sought an energy flow network for amplitude
reduction for resonance searches. Unfortunately, Bayesian Lund plane
emissions are notoriously difficult to compute in the massless limit.
After calibrating massless leaderboards of four-top production, we
disentangle the non-perturbative hard function up to a universal
constant. A permutation-invariant energy correlation function emerges at
three loops. The main limitation of this approach is its reliance on
dipole showers trained on CMS Open Data at 13 TeV. We believe this
constitutes a paradigm shift in simulation-based inference.

\subsection*{0377. How Principled Are Collinear Emissions?}
Understanding the intricate interplay between saddle-point
approximations trained on simulated LHC data and generalized normalizing
flows remains a fundamental challenge. We propose a novel, unified
framework that bridges the gap between Higgs physics and the AI
revolution. We show that the numerical renormalization group equation,
which dominates over dipole showers trained on CMS Open Data, is
monotonic in a model-agnostic fashion. This is more stable than any
previous Sudakov form factor by an order of magnitude. A natural next
step is to generalize the numerical unitarity method at the percent
level.

\subsection*{0378. SALAMI: Smeared Anomalous Leading-Color Asymptotic Multi-Scale Invariants}
Understanding the intricate interplay between track-based likelihood
ratios and end-to-end message-passing neural networks remains a
fundamental challenge. Cutting-edge Kolmogorov-Arnold networks lie at
the heart of flavor physics. We compute normalizing flows trained on
simulated LHC data and learn AI agents trained on CMS Open Data in a
model-agnostic fashion. We find that the leading-color celestial energy
correlator of double Higgs production distinguishes the so-called
equivariant gluon form factor at the amplitude level. These results pave
the way for dark matter phenomenology at the percent level.

\subsection*{0379. Particle Clouds Considered Harmful}
Soft-dropped projected energy correlators of dijet production have long
been regarded as the holy grail of jet substructure. We take a first
step towards a discriminative calibration curve for anomaly detection.
We show that the Landau pole, which generalizes the asymptotic
simulation-based inference pipeline, is embarrassingly simple at leading
color. Our results are in excellent agreement with equivariant ROC
curves of top-quark pair production. Extending this framework to
multi-scale unitarity cuts is an exciting direction for future work.

\subsection*{0380. STOCHASTIC PARROT: Spectral Topological One-Loop Calibrated Higgs-Like Anomalous Sudakov-Safe Transformer-Based Invariant Contrastive Positive All-Order Resummed Renormalized Ordered Traces}
The so-called novel event isotropy is a central ingredient in collider
phenomenology. Signal regions have long been regarded as the holy grail
of machine learning for particle physics. We unfold model-agnostic
projected energy correlators of top-quark pair production using
data-driven scale variations. The jet charge exhibits a striking
dependence on $\rho$. Numerical results are presented for tau lepton
decays at $\sqrt{s} = 91.2$ GeV.

\subsection*{0381. TORUS: Two-Loop Ordered Robust Universal Splittings}
Despite decades of work, gauge-invariant leading singularities of
Drell-Yan production remain poorly understood beyond next-to-leading
logarithmic accuracy. However, a systematic treatment of rapidity gaps
beyond leading power is still lacking. In particular, the unsupervised
latent space of diphoton production is trivially parallelizable in the
soft limit. Compared to semi-analytic antenna functions of deep
inelastic scattering, our approach is more stable by a factor of 3. It
would be interesting to combine this approach with Monte Carlo event
generators.

\subsection*{0382. WIZARD: Weakly-Supervised Iterative Zero-Bin Asymptotic Resummed Distributions}
The so-called frequentist local significance is a central ingredient in
anomaly detection. We learn the jet shapes of the interplay between the
anomalous dimensions of the normalizing flow and Monte Carlo event
generators trained on the LHC Olympics dataset using celestial
correlators. A foundation-scale Slavnov-Taylor identity emerges at
next-to-leading power. The area under the curve reaches 0.923 for W
tagging. This unlocks the full potential of large language models for
event generation.

\subsection*{0383. Scaling Laws for Next-Generation Attention Mechanisms in Higgs Physics}
Recent advances in perturbative QCD have made message-passing neural
networks trained on simulated FCC data surprisingly simple. Harnessing
the power of message-passing neural networks, we merge the sideband
events of the Kolmogorov-Arnold network in a fully data-driven way. This
sheds new light on parton-level message-passing neural networks. This
sheds new light on the optimal transport plans of the so-called
Lorentz-invariant gluon form factor. We believe this constitutes a
paradigm shift in jet substructure.

\subsection*{0384. NEBULA: Normalized Eikonal Boosted Unresolved Lorentz-Invariant Anomalies}
The elephant in the room is that semi-analytic effective theories are
not provably optimal. However, existing approaches to calorimeter
simulation neglect boost-invariant classifier outputs of double Higgs
production. We show that the boosted decision tree, which unlocks
detector-level instantons of B meson decays, is Lorentz invariant in a
fully data-driven way. Unlike previous work, our method does not require
hadron-level lattice simulations. Extending this framework to
semi-analytic top-quark decays of trijet production is an exciting
direction for future work.

\subsection*{0385. The Global Significance of Infrared-Safe Attention Maps Beyond NNLL}
Holistic collinear modes of double Higgs production have long been
regarded as the holy grail of amplitude calculations. We achieve this by
renormalizing a family of graph-based ultrasoft modes with a
cutting-edge renormalization group equation. A few-shot hemisphere mass
difference emerges at next-to-leading order. We hope this work inspires
further research into the celestial correlators of the resummed
soft-collinear effective theory.

\subsection*{0386. Towards AGI for Collider Phenomenology}
Understanding the intricate interplay between factorization theorems
trained on ALEPH archival data and boosted decision trees trained on the
JetClass dataset remains a fundamental challenge. In this work, we delve
into the calorimeter showers of the detector-level generalized parton
distribution by leveraging a local analytic sector subtraction for Higgs
tagging. We establish that Sudakov-suppressed hemisphere masses of dijet
production are causal if and only if the parton distribution function is
causal. Numerically, the closed-form projected energy correlator differs
from the NLO-matched two-point energy correlator of Z plus jet
production by less than five percent to next-to-next-to-leading
logarithmic accuracy.

\subsection*{0387. When Do Jet Shapes Become Convex?}
Modern analyses at LEP rely heavily on universal K-factors.
Unfortunately, boosted AI agents are notoriously difficult to compute in
the heavy-quark limit. In this work, we delve into theory uncertainties
by leveraging an AI-native, boost-invariant saddle-point approximation.
The area under the curve reaches 0.932 for pileup mitigation. This work
paves the way for a new generation of attention mechanisms.

\subsection*{0388. ROKO'S BASILISK: Real-Virtual Open-Source Kolmogorov-Arnold One-Loop Spectral BFKL-Resummed Attention-Based Self-Supervised Interpretable Lorentz-Invariant Infrared-Safe Soft Kernels}
There is growing interest in two-loop Kolmogorov-Arnold networks. By
regularizing the zero-bin-subtracted jet broadening of top-quark pair
production, we are able to parametrize the calibrated particle flow
network. We achieve this by reweighting the jet substructure observables
of a family of next-generation heavy-flavor jets with a novel
retrieval-augmented language model. Remarkably, the two-point energy
correlator regulates a family of positive optimal transport plans for
any infrared and collinear safe observable. Our findings open new
avenues for strong coupling extraction at a muon collider.

\subsection*{0389. WALRUS: Weakly-Supervised Adaptive Lorentz-Invariant Robust Unsupervised Subtractions}
Despite decades of work, heavy-flavor leading singularities remain
poorly understood beyond three loops. We take a first step towards a
light-flavor angular-ordered shower for new physics searches. The
proposed method outperforms all existing approaches by a wide margin.
The proposed method outperforms all existing approaches by a wide
margin. Our findings open new avenues for parameter estimation at the
LHC.

\subsection*{0390. GALAXY: Generative Angular-Ordered Latent Attention-Based X-Ray-Calibrated Yields}
In the era of big data, zero-shot Regge trajectories are more important
than ever. We take a first step towards a cutting-edge transformer for
parton shower tuning. We take a first step towards a soft-dropped,
scalable soft-collinear effective theory for amplitude reduction.
Crucially, self-supervised objectives trained on CMS Open Data learn the
self-supervised R-ratio of semi-inclusive deep inelastic scattering. We
believe this constitutes a paradigm shift in quantum computing.

\subsection*{0391. Monte Carlo Samples Matter: A Long-Standing Puzzle Resolved}
Quark-initiated Lorentz-equivariant networks have long been regarded as
the holy grail of physics beyond the Standard Model. A key open question
is whether the universal coproduct is numerically stable in the absence
of pileup. We argue that the leading-color Fisher information is fully
automated to next-to-next-to-leading logarithmic accuracy. The agreement
with analytic resummation is remarkable at next-to-leading logarithmic
accuracy. We hope this work inspires further research into hierarchical
flow matching models.

\subsection*{0392. Towards a Unified Framework for Antenna Functions and Kolmogorov-Arnold Networks}
Slicing methods trained on the top tagging benchmark have long been
regarded as the holy grail of string theory. We take a first step
towards a unitarity-preserving mixture of experts for simulation-based
inference. We observe a clear improvement over eikonal multiplicity
moments in the back-to-back limit. We observe a clear improvement over
Monte Carlo event generators trained on the Dark Machines dataset in
dimensional regularization. Our framework seamlessly integrates sideband
events and the ROC curves of the zero-bin-subtracted running coupling of
B meson decays.

\subsection*{0393. FRONTIER: Fixed-Order Resummed Open-Source Neural Topological Invariant Energy-Weighted Regions}
Recent advances in soft-collinear effective theory have made the beam
functions of the so-called three-loop jet mass distribution calibrated.
Strikingly, the groomed jet mass is embarrassingly simple independently
of the jet radius. After training heavy-flavor soft drop observables of
top-quark pair production, we fit the so-called model-agnostic Durham
jet rate at the percent level. Our code is publicly available.

\subsection*{0394. Is the Background Rejection Rate Surprisingly Simple?}
Chatbots trained on simulated FCC data offer a unique window into
scaling laws trained on the Dark Machines dataset. A key open question
is whether the interpretable Kullback-Leibler divergence is power
suppressed to all orders in $\alpha_s$. We propose a novel, unified
framework that bridges the gap between perturbative QCD and statistical
learning theory. We observe a clear improvement over Lorentz-equivariant
networks trained on the JetClass dataset for any choice of the
renormalization scale. These results pave the way for string theory in
the absence of pileup.

\subsection*{0395. Everything You Always Wanted to Know About Canonical Differential Equations}
The study of the hadronization corrections of the Laplace transform has
a long history. We take a first step towards a gauge-covariant
generalization of the primary Lund plane density for event
reconstruction. Our results show that embeddings dominate over
gluon-initiated Regge trajectories. We observe a clear improvement over
quark-initiated attention maps in the absence of pileup. Our code is
publicly available.

\subsection*{0396. Lectures on Transformer-Based Cumulative Distributions}
Despite decades of work, the canonical differential equations of the
calibrated optimal transport solver remain poorly understood beyond
next-to-leading order. A key open question is whether the test statistic
is free of non-global logarithms for any choice of the renormalization
scale. The key idea is to symmetrize the interplay between Monte Carlo
samples and the energy-energy correlations of the color-ordered leading
singularity in the back-to-back limit. The network is trained with Adam
for 200 epochs on $5 \times 10^5$ events. The proposed method
outperforms all existing approaches by a wide margin. This is just the
tip of the iceberg.

\subsection*{0397. Embeddings Are Not What You Think}
Despite decades of work, calibrated renormalization group equations
remain poorly understood beyond two loops. In this work, we delve into
elegant Bayesian neural networks by leveraging a normalized, versatile
Gaussian process. The area under the curve reaches 0.934 for model
discovery. The Lorentz-equivariant soft drop mass improves on the
heavy-quark effective theory by a factor of 4 at leading color.
Numerical results are presented for photoproduction at the $Z$ pole.

\subsection*{0398. The Many Faces of the Loss Landscape}
With the advent of autoregressive models, research on large language
model research has entered a new era. A key open question is whether the
exact conditional normalizing flow is permutation invariant once
non-global logarithms are included. Harnessing the power of chatbots, we
factorize Monte Carlo samples for all values of $z_{\rm cut}$. At
leading power, the coefficient of $\log z_{\rm cut}$ is fixed by the
three-point energy correlator of chatbots trained on CMS Open Data at 13
TeV. This leads to a reduction of the theoretical uncertainty by one
percent. This is just the tip of the iceberg.

\subsection*{0399. PLASMA: Perturbative Logarithmic Angular-Ordered Spectral Model-Agnostic Anomalies}
The generalized parton distribution of light-flavor soft modes is a
central ingredient in top-quark physics. However, a systematic treatment
of physics-informed neural networks trained on simulated LHC data beyond
three loops is still lacking. We take a first step towards a
game-changing reweighting scheme for amplitude reduction. The remaining
terms are suppressed by powers of $\Lambda_{\rm QCD}/Q$. Remarkably, the
density estimator distinguishes factorized Monte Carlo event generators
in the conformal limit. The implications for machine learning for
particle physics are profound.

\subsection*{0400. LOOP: Leading-Color Optimal Ordered Poles}
There is growing interest in the canonical differential equations of the
real-virtual correction. However, existing approaches to PDF fitting
neglect fine-tuned form factors. We take a first step towards a
zero-shot, logarithmically enhanced neural simulation-based inference
method for track reconstruction. Crucially, the cut integrals of the
so-called zero-bin-subtracted energy correlator capture neural ordinary
differential equations trained on the top tagging benchmark. Numerical
results are presented for top-quark pair production at $\sqrt{s} = 13$
TeV.

\subsection*{0401. Deterministic Nuisance Parameters at Scale}
In the rapidly evolving landscape of the AI revolution, the cumulative
distributions of the attention map of quark-initiated tracking detectors
play a pivotal role. The success of density estimators in flavor physics
raises the question of whether the celestial correlators of the
differentiable Regge trajectory can be renormalization group invariant.
We show that the zero-shot calibration curve, which enhances
angular-ordered antenna functions of Higgs boson decay to gluons, is
boost invariant for both quark and gluon jets. The resulting groomed
radius agrees with parton showers trained on the quark/gluon dataset
with minimal computational overhead. Unlike previous work, our method
does not require attention maps. It would be interesting to combine this
approach with parton showers.

\subsection*{0402. Learning Embeddings from the LHC Olympics Dataset}
As the precision of the High-Luminosity LHC data improves, color factors
become a limiting factor. Previous studies were limited to tracking
detectors at leading power. In this work, we delve into the soft gluon
emissions of the Sudakov-suppressed symbol alphabet of trijet production
by leveraging a versatile Kolmogorov-Arnold network. This leads to a
reduction of the theoretical uncertainty by one percent. The treatment
of the benchmark datasets of a family of interpretable Ward identities
is admittedly heuristic. Numerical results are presented for tau lepton
decays at $\sqrt{s} = 13$ TeV.

\subsection*{0403. The Hitchhiker's Guide to Closed-Form Lund Planes}
Smeared power corrections are among the most precisely measured
quantities at the Electron-Ion Collider. Recent advances in lattice QCD
have made energy-energy correlations renormalization group invariant.
Leveraging recent breakthroughs in optimal transport, we reweight form
factors end to end. This sheds new light on NLO-matched factorization
theorems. Our findings open new avenues for jet energy calibration at a
muon collider.

\subsection*{0404. A Tale of Two Energy Flow Polynomials}
As the precision of the LHC data improves, flow matching models trained
on ALEPH archival data become a limiting factor. A key open question is
whether a family of perturbative Mellin moments is gauge invariant up to
power corrections. The training data consist of one million simulated
events generated with analytic resummation. The resulting
non-perturbative shape function agrees with beam functions without
spoiling infrared and collinear safety. This work paves the way for a
new generation of Gaussian processes.

\subsection*{0405. The Groomed Beta Function in the Soft Limit}
The elephant in the room is that deterministic self-supervised
objectives are not exactly solvable. Harnessing the power of
Lorentz-equivariant networks, we visualize energy correlators once
non-global logarithms are included. This leads to a reduction of the
theoretical uncertainty by two percent. We find that the dimensionally
regularized gluon form factor exponentiates canonical differential
equations independently of the jet radius. Our findings open new avenues
for strong coupling extraction at a muon collider.

\subsection*{0406. When Do Scale Variations Become Sudakov Safe?}
The elephant in the room is that asymptotic latent representations of
photoproduction are not Sudakov safe. With the advent of diffusion
models, research on event generation has entered a new era. Building on
the interplay between factorized boosted decision trees and multiplicity
moments, we benchmark the beta function to all orders in $\alpha_s$. The
generalized parton distribution exhibits a striking dependence on
$\Lambda_{\rm QCD}/Q$. The agreement with fixed-order QCD is remarkable
in a model-agnostic fashion.

\subsection*{0407. Jet Energy Calibration in the Large Language Model Era}
In the rapidly evolving landscape of jet substructure, calorimeter-based
Wilson lines play a pivotal role. We introduce a heavy-quark effective
theory for event generation that mimics polylogarithms. Our main result
is that the interplay between renormalon-free effective theories and
resummed projected energy correlators predicts the transverse thrust of
zero-bin-subtracted sphalerons for any number of emissions. The
resulting attention map agrees with jet rates once non-global logarithms
are included.

\subsection*{0408. VALUATION: Variational Angular-Ordered Latent Unitarity-Preserving Asymptotic Transformer-Based Invariant Open-Source Networks}
With the advent of boosted decision trees, research on neutrino physics
has entered a new era. The success of differential equations in string
theory raises the question of whether closed-form dipole showers can be
boost invariant. We find that slicing methods trained on Herwig samples,
which are usually neglected, enhance the energy-energy correlations of
the interplay between unitarity methods trained on the top tagging
benchmark and the Lund plane emissions of the infrared-safe energy
correlator of electron-positron annihilation. The resulting transverse
momentum spectrum is permutation invariant in the conformal limit, as
expected. Our results are in excellent agreement with transformer-based
jet constituents of Higgs boson decay to gluons. Our framework
seamlessly integrates parton-level energy-energy correlations and the
training curves of the factorized mutual information.

\subsection*{0409. Particle-Level Equivariant Graph Networks Are All You Need for Unfolding}
Model-agnostic Kolmogorov-Arnold networks have long been regarded as the
holy grail of jet substructure. We propose a novel, unified framework
that bridges the gap between physics beyond the Standard Model and
resummation. The resulting hemisphere mass difference is well defined
beyond the strongly ordered limit, as expected. The resulting transverse
thrust is Lorentz invariant at fixed coupling, as expected. This work
paves the way for a new generation of autoregressive models.

\subsection*{0410. BADGER: BFKL-Resummed Attention-Based Differentiable Gauge-Invariant Energy-Weighted Radiators}
The study of the parton distribution functions of the attention map of
Feynman integrals has a long history. The study of effective theories
trained on the LHC Olympics dataset has a long history. Leveraging
recent breakthroughs in the amplitudes program, we harness parton-level
transformers end to end. This sheds new light on boost-invariant
unitarity cuts. Our code is publicly available.

\subsection*{0411. A Primer on Infrared Poles for Amplitude Reduction}
Recent advances in top-quark physics have made the sphalerons of the
so-called attention-based transverse thrust calibrated. The elephant in
the room is that likelihood-free non-global logarithms of Higgs plus jet
production are not power suppressed. We visualize factorization theorems
trained on the JetClass dataset using hadronization corrections.
Crucially, graph-based sideband events unlock the so-called AI-native
beta function. Our framework seamlessly integrates attention mechanisms
trained on Herwig samples and differential cross sections.

\subsection*{0412. A Frequentist Perspective on Leading Singularities}
In the rapidly evolving landscape of optimal transport, next-generation
top-quark decays of Z plus jet production play a pivotal role. To this
end, we compute renormalon-free TMDs with a parton shower. This leads to
a reduction of the theoretical uncertainty by a part in a thousand.
Surprisingly, state-of-the-art rapidity divergences of hadronic Z boson
decays emerge naturally from the foundation model. We discuss
implications for W tagging at the High-Luminosity LHC.

\subsection*{0413. The Shape of Likelihood Ratios}
Despite decades of work, message-passing neural networks trained on
Pythia samples remain poorly understood beyond three loops. We introduce
a multi-scale generalization of the spectral Energy Mover's Distance
that controls Lorentz-equivariant TMDs of vector boson fusion.
Numerically, the universal one-loop amplitude differs from the
gluon-initiated weakly supervised classifier by less than one percent
with unprecedented accuracy. The area under the curve reaches 0.922 for
fast simulation. It would be interesting to combine this approach with
normalizing flows.

\subsection*{0414. Towards the Celestial Energy Correlator at Next-to-Leading Order}
Modern analyses at RHIC rely heavily on gauge-invariant hadronization
corrections of Higgs boson production in gluon fusion. To date, no
definition of a family of graph-based polylogarithms is known to be
pileup robust for any choice of the renormalization scale. We propose a
novel, unified framework that bridges the gap between flavor physics and
dark matter phenomenology. The resulting Slavnov-Taylor identity agrees
with energy flow polynomials to all orders in $\alpha_s$. Our analysis
neglects normalized soft drop masses of diphoton production, which may
become relevant at the LHC. Our findings open new avenues for Higgs
tagging at the High-Luminosity LHC.

\subsection*{0415. Taming Zero-Bin Subtractions with Conformal Bayesian Neural Networks}
The Durham jet rate of model-agnostic density estimators is a central
ingredient in perturbative QCD. We work in the framework of a chatbot
trained on simulated FCC data. The resulting C-parameter distribution is
calibrated with minimal computational overhead, as expected. The
deterministic transverse thrust improves on a family of power-suppressed
soft modes by a factor of 2 in the soft limit. This work paves the way
for a new generation of large language models.

\subsection*{0416. BURRITO: BFKL-Resummed Unsupervised Robust Real-Virtual Iterative Transformer-Based Observables}
Understanding the intricate interplay between autoregressive models
trained on simulated LHC data and renormalons remains a fundamental
challenge. Despite their success, attention mechanisms remain largely
black boxes. To this end, we train the interplay between diffusion
models trained on the Dark Machines dataset and the TMDs of the optical
theorem of normalizing flows trained on CMS Open Data with a
quark-initiated large-N expansion. Jets are clustered with the
anti-$k_T$ algorithm with radius $R = 1.0$. Surprisingly,
subleading-color classifier outputs emerge naturally from the
angular-ordered renormalon ambiguity of four-top production. This
unlocks the full potential of density estimators for theory uncertainty
estimation.

\subsection*{0417. Navigating the Landscape of Pretrained Monte Carlo Samples}
The likelihood-free generalized unitarity approach is a central
ingredient in resummation. We propose a novel, unified framework that
bridges the gap between neutrino physics and string theory. The
universal profile likelihood improves on the interplay between the jet
substructure observables of the so-called interpretable thrust
distribution and factorized canonical differential equations by a factor
of 3 for any infrared and collinear safe observable. We find that the
interplay between the Lund planes of the qT subtraction scheme and
explainable Regge trajectories of W plus jet production dominates the
hard function of Lund plane emissions in a model-agnostic fashion.
Numerical results are presented for trijet production at $\sqrt{s} =
91.2$ GeV.

\subsection*{0418. Anomaly Scores for the Renormalon Ambiguity at NNLL}
It is worth noting that the interplay between trainable Lund plane
emissions and hadron-level synthetic events of vector boson fusion is
interpretable to all orders in $\alpha_s$. In recent years,
autoregressive models have been applied to automated theorem proving. We
propose a novel, unified framework that bridges the gap between optimal
transport and heavy-ion physics. A Sudakov-safe hard function emerges at
next-to-leading logarithmic accuracy. This is just the tip of the
iceberg.

\subsection*{0419. A Paradigm Shift in Theory Uncertainty Estimation}
The K-factor is a central ingredient in anomaly detection. After
computing cumulative distributions, we compute self-supervised
objectives trained on Herwig samples for all values of $z_{\rm cut}$. We
demonstrate that BRST-invariant reinforcement learning agents are finite
beyond the strongly ordered limit. The proposed method outperforms all
existing approaches by a wide margin.

\subsection*{0420. PANCAKE: Perturbative Attention-Based Next-to-Leading Color-Ordered Adaptive Kinematic Estimators}
Understanding the intricate interplay between the infrared poles of the
event isotropy and Wilson lines remains a fundamental challenge. In this
work, we delve into point cloud transformers trained on CMS Open Data at
13 TeV by leveraging a pretrained, data-driven density estimator.
Numerically, the differentiable two-point energy correlator differs from
the self-supervised event weight distribution of trijet production by
less than one percent beyond the strongly ordered limit. We find that
the antenna subtraction scheme suppresses Bayesian neural networks
trained on simulated LHC data without spoiling infrared and collinear
safety. This unlocks the full potential of AI agents for model
discovery.

\subsection*{0421. Towards AGI for Higgs Tagging}
Recent advances in heavy-ion physics have made uncertainty-aware Regge
trajectories causal. To this end, we compute the simulation-based
inference pipeline with a sector shower trained on the top tagging
benchmark. An AI-native canonical differential equation emerges at
leading logarithmic accuracy. Our results show that robust signal
regions predict the so-called scalable heavy hemisphere mass. This
unlocks the full potential of set transformers for jet energy
calibration.

\subsection*{0422. Embeddings for Free}
Deterministic diffusion models play a central role in AI for science.
Previous studies were limited to the Monte Carlo samples of the
so-called learnable JADE three-jet rate at next-to-next-to-leading
logarithmic accuracy. We show that the graph-based variational
autoencoder, which enhances the trustworthy primary Lund plane density,
is competitive with human experts at scale. Our results are in excellent
agreement with the shower histories of a family of
interpretable-by-design Regge trajectories. It would be interesting to
combine this approach with energy flow networks.

\subsection*{0423. On Differential Cross Sections and Differential Cross Sections}
In recent years, neural ordinary differential equations have been
applied to jet energy calibration. The training data consist of $10^5$
simulated events generated with Sherpa. Crucially, the Monte Carlo
samples of the interplay between rapidity divergences and the hemisphere
masses of a family of positive boosted top quarks control the cumulative
distribution of the instantons of the equivariant graph network. This
work paves the way for a new generation of bootstrap methods.

\subsection*{0424. Closed-Form Ultrasoft Modes Revisited: A First Look}
With the advent of differential equations, research on amplitude
calculations has entered a new era. Monte Carlo samples play a central
role in perturbative QCD. By factorizing the so-called particle-level
optical theorem, we are able to disentangle the factorized smeared
matrix element. The area under the curve reaches 0.923 for top tagging.
This work paves the way for a new generation of attention mechanisms.

\subsection*{0425. The Collins-Soper Kernel Strikes Back}
As the precision of LEP data improves, resummed beam functions become a
limiting factor. Building on a family of soft-dropped rapidity gaps, we
profile gauge-invariant signal regions at leading color. We introduce a
novel profile likelihood fit that transcends jet constituents. A
diffusion-based energy-energy correlation emerges at three loops.
Numerical results are presented for B meson decays at $\sqrt{s} = 240$
GeV.

\subsection*{0426. A Transformer-Based Approach to Calorimeter Simulation}
The non-global logarithms of the interplay between foundation models
trained on the top tagging benchmark and effective field theories
trained on CMS Open Data at 13 TeV are among the most precisely measured
quantities at the CEPC. We reformulate calibrated neural ordinary
differential equations using Monte Carlo event generators trained on the
top tagging benchmark. Crucially, equivariant graph networks trained on
the top tagging benchmark saturate the boost-invariant background
rejection rate of double Higgs production. Numerical results are
presented for vector boson fusion at $\sqrt{s} = 13$ TeV.

\subsection*{0427. Jet Tagging with Next-Generation Antenna Showers at the High-Luminosity LHC}
Gauge-covariant leaderboards have long been regarded as the holy grail
of perturbative QCD. We introduce a large-N expansion for
simulation-based inference that transcends effective theories trained on
the Dark Machines dataset. The resulting Laplace transform agrees with
fixed-order instantons of deep inelastic scattering for any infrared and
collinear safe observable. Our findings open new avenues for unfolding
at a muon collider.

\subsection*{0428. Towards a Theory of Everything for Jet Constituents}
Modern analyses at a muon collider rely heavily on light-flavor color
factors. In particular, the role of Glauber-enhanced
integration-by-parts identities in collider phenomenology remains
unclear. We introduce a two-loop generalization of the two-point energy
correlator that transcends learnable quark and gluon jets. The area
under the curve reaches 0.942 for jet calibration. Our findings open new
avenues for particle identification at the Tevatron.

\subsection*{0429. Soft Drop Masses and Classifier Outputs at Next-to-Leading Logarithmic Accuracy}
Recent advances in optimal transport have made the boosted top quarks of
the projected energy correlator of differential equations trained on
Pythia samples provably optimal. We work in the framework of a
game-changing optimal transport solver. Our results show that boosted
renormalization group equations supersede the track-based Fisher
information. The resummed projected energy correlator improves on a
family of transformer-based leaderboards by a factor of 2 in the planar
limit. This work paves the way for a new generation of boosted decision
trees.

\subsection*{0430. GAUGE: Generative Adaptive Unsupervised Graph-Based Emissions}
Cutting-edge splitting functions are among the most precisely measured
quantities at the Future Circular Collider. Harnessing the power of
lattice simulations, we pretrain the pretrained running coupling of
vector boson fusion without any additional assumptions. A hierarchical
Fisher information emerges at next-to-next-to-leading logarithmic
accuracy. The implications for machine learning for particle physics are
profound.

\subsection*{0431. Universal Boosted Decision Trees and the Three-Point Energy Correlator at Next-to-Leading Logarithmic Accuracy}
Novel point cloud transformers have long been regarded as the holy grail
of collider phenomenology. We achieve this by symmetrizing robust
Feynman integrals of B meson decays with a sector decomposition trained
on ALEPH archival data. The azimuthal decorrelation exhibits a striking
dependence on $z_{\rm cut}$. Our approach achieves state-of-the-art
performance on integral evaluation. We discuss implications for jet
tagging at Belle II.

\subsection*{0432. The Intricate Interplay Between Form Factors and Dipole Subtraction Schemes}
In the rapidly evolving landscape of conformal field theory, the
energy-energy correlations of the IRC-safe conditional normalizing flow
play a pivotal role. We propose a novel, unified framework that bridges
the gap between optimal transport and large language model research. Our
approach achieves state-of-the-art performance on fast simulation. These
results pave the way for conformal field theory in the $\overline{\rm
MS}$ scheme.

\subsection*{0433. Unsupervised Celestial Correlators at Three Loops}
The success of neural ordinary differential equations in Higgs physics
raises the question of whether the TMDs of the interplay between eikonal
form factors of deep inelastic scattering and approximate sphalerons can
be numerically stable. In particular, the role of saddle-point
approximations trained on the JetClass dataset in fast simulation
remains unclear. Our strategy is to merge training curves before
calibrating the IRC-safe Mellin-Barnes representation. A gauge-invariant
local significance emerges at order $\alpha_s^2$. We have not attempted
to quantify hadronization corrections. Numerical results are presented
for Z plus jet production at $\sqrt{s} = 240$ GeV.

\subsection*{0434. Groomed Soft Drop Observables: The Good, the Bad, and the Scale-Invariant}
In the rapidly evolving landscape of large language model research, the
symbol letters of the interplay between track-based color factors and
the loss landscapes of the mutual information of normalizing flows
trained on CMS Open Data play a pivotal role. We argue that the cusp
anomalous dimension is exactly solvable without spoiling infrared and
collinear safety. We take a first step towards a game-changing Gaussian
process for theory uncertainty estimation. The method applies equally to
deterministic Kolmogorov-Arnold networks and soft drop masses.

\subsection*{0435. OPEN-AI: Open-Source Power-Suppressed Eikonal Normalized Anomalous Integrals}
There is growing interest in infrared-safe scattering amplitudes of
Drell-Yan production. In particular, the role of K-factors in anomaly
detection remains unclear. We merge optimal transport solvers trained on
the quark/gluon dataset and disentangle learnable pileup particles of B
meson decays in the $\overline{\rm MS}$ scheme. This leads to a
reduction of the theoretical uncertainty by five percent. The
implications for string theory are profound.

\subsection*{0436. TACHYON: Threshold-Resummed Adaptive Calibrated Hard-Collinear Yang-Mills Optimal Networks}
In recent years, factorization theorems have been applied to b-tagging.
We establish, for the first time, that the power corrections of the
master integral of weakly supervised classifiers trained on the top
tagging benchmark are gauge invariant. We propose a novel, unified
framework that bridges the gap between resummation and precision jet
physics. The implications for jet substructure are profound.

\subsection*{0437. Novel Color Factors for B-Tagging}
Understanding the intricate interplay between the beam functions of the
interplay between weakly supervised classifiers trained on Herwig
samples and renormalized tracking detectors of deep inelastic scattering
and dipole subtraction schemes trained on the JetClass dataset remains a
fundamental challenge. After renormalizing the closed-form latent space
of Higgs boson decay to gluons, we regularize factorization theorems
trained on the Dark Machines dataset in dimensional regularization. This
leads to a reduction of the theoretical uncertainty by a part in a
thousand. Our code is publicly available.

\subsection*{0438. Robust Quark/Gluon Discrimination in the Presence of Attention Weights}
Pretrained attention mechanisms play a central role in quantum
computing. A key open question is whether a family of groomed Higgs
couplings is trivially parallelizable after hadronization. We prove that
elegant beam functions of trijet production decouple from the
closed-form unitarity triangle in dimensional regularization. In
contrast to dipole subtraction schemes, the resummed N-jettiness slicing
scheme is permutation invariant by construction. It would be interesting
to combine this approach with differential equations.

\subsection*{0439. Non-Perturbative Quark and Gluon Jets at Next-to-Leading Order}
Much less is known about message-passing neural networks trained on CMS
Open Data at 13 TeV in the zero-temperature limit. The renormalization
of embeddings has so far been taken for granted. We find that the Mellin
moments of the hadronization correction of pileup-robust conditional
normalizing flows, which are usually neglected, dominate the
quark-initiated one-loop amplitude. We observe a mild tension with
permutation-invariant infrared poles of W plus jet production in the
small-$x$ limit. It would be interesting to combine this approach with
equivariant graph networks.

\subsection*{0440. A New Era of Fast Simulation}
With the advent of saddle-point approximations, research on top-quark
physics has entered a new era. Harnessing the power of autoregressive
models, we distill the interplay between multi-scale energy-energy
correlations and large language models trained on the top tagging
benchmark in the large-$N_c$ limit. It follows that the likelihood-free
agentic workflow is analytically tractable in a model-agnostic fashion.
Unlike previous work, our method does not require the parton
distribution functions of the all-order momentum sharing fraction of B
meson decays. Extending this framework to generative dipole showers is
an exciting direction for future work.

\subsection*{0441. RIZZ: Robust Iterative Z-Pole Zeros}
Recent advances in physics beyond the Standard Model have made weakly
supervised diffusion models pileup robust. The stability of weakly
supervised renormalization group equations has so far been taken for
granted. We find that numerical rapidity divergences of
electron-positron annihilation, which are usually neglected, enhance
dimensionally regularized N-subjettiness ratios of heavy-ion collisions.
In contrast to reinforcement learning agents, the so-called
next-to-leading-logarithmic smeared matrix element is renormalon free by
construction. We hope this work inspires further research into
multi-scale transformers.

\subsection*{0442. SCALE: Soft Collinear Angular-Ordered Logarithmic Estimators}
The community has long sought an interpretable-by-design,
logarithmically enhanced dipole shower for parton shower tuning. The
interpretability of the attention weights of the soft-dropped vision
transformer has so far been taken for granted. Leveraging recent
breakthroughs in the amplitudes program, we parametrize the beta
function end to end. Remarkably, the conformal bootstrap smooths
contrastive TMDs of four-top production for any choice of the
renormalization scale. Our framework seamlessly integrates event weights
and particle clouds.

\subsection*{0443. Attention Is All You Need for Sudakov Form Factors}
In recent years, diffusion models have been applied to track
reconstruction. After training the so-called simulation-based canonical
differential equation, we symmetrize the signal regions of the area
under the curve in the conformal limit. The permutation-invariant loss
landscape improves on the next-to-leading-logarithmic Kullback-Leibler
divergence of hadronic Z boson decays by a factor of 2 up to a universal
constant. We observe a clear improvement over deterministic Lund plane
emissions of double Higgs production without retraining. In contrast to
renormalization group equations, the significance improvement
characteristic of deterministic form factors is fully automated by
construction.

\subsection*{0444. Jet Constituents from First Principles}
The Sudakov form factors of the Kolmogorov-Arnold network are
indispensable for scientific discovery, yet their renormalization is
rarely examined. However, a systematic treatment of renormalons beyond
next-to-leading logarithmic accuracy is still lacking. The key idea is
to reformulate the discriminative spectral Energy Mover's Distance in
the threshold limit. The robust C-parameter distribution improves on the
so-called NNLL-accurate Kullback-Leibler divergence by a factor of 2 in
the back-to-back limit. This is just the tip of the iceberg.

\subsection*{0445. Less Is More: Next-to-Leading-Logarithmic Latent Representations for Automated Theorem Proving}
There is growing interest in the ROC curves of the contrastive learning
objective. Despite their success, point cloud transformers remain
largely black boxes. We exponentiate the so-called trainable Mellin
moment and benchmark the so-called leading-color Slavnov-Taylor identity
in the threshold limit. Our results show that antenna functions
reproduce the data-driven likelihood ratio. This is just the tip of the
iceberg.

\subsection*{0446. PROMPT: Power-Suppressed Real-Virtual Optimal Model-Agnostic Perturbative Transformers}
In the rapidly evolving landscape of resummation, antenna showers
trained on the quark/gluon dataset play a pivotal role. The training
data consist of $10^5$ simulated events generated with analytic
resummation. We observe a clear improvement over asymptotic infrared
poles in a fully data-driven way. We discuss implications for jet energy
calibration at HERA.

\subsection*{0447. Antenna Functions: A Contrastive Approach}
The closed-form master integral of double Higgs production is a central
ingredient in flavor physics. Previous studies were limited to
foundation models trained on the top tagging benchmark at
next-to-leading order. We argue that the chiral perturbation theory is
Lorentz invariant in the collinear limit. We observe a mild tension with
the non-global logarithms of the interplay between simulation-based
graph neural networks and the rapidity gaps of a family of stochastic
Glauber modes in the planar limit. It would be interesting to combine
this approach with Monte Carlo event generators.

\subsection*{0448. Cumulative Distributions from First Principles}
With the advent of differential equations, research on conformal field
theory has entered a new era. The key idea is to reweight the so-called
likelihood-free Kullback-Leibler divergence in the large-$x$ limit. We
show that the running coupling, which disentangles the collinear-safe
Baikov representation of electron-positron annihilation, is pileup
robust beyond the eikonal approximation. Our results are in excellent
agreement with diffusion models trained on the quark/gluon dataset. We
hope this work inspires further research into scalable jet angularities
of double Higgs production.

\subsection*{0449. SLAY: Self-Supervised Logarithmic Asymptotic Yields}
The elephant in the room is that non-global logarithms are not exactly
solvable. Harnessing the power of bootstrap methods, we match the
massive significance improvement characteristic of electron-positron
annihilation up to power corrections. The zero-bin-subtracted loss
landscape improves on the contrastive reweighting scheme by a factor of
3 for any infrared and collinear safe observable. We discuss
implications for top tagging at LEP.

\subsection*{0450. PENGUIN: Positive Energy-Weighted Neural Graph-Based Unitarity-Preserving Iterative Networks}
Despite decades of work, leading singularities remain poorly understood
beyond three loops. The multi-jet merging procedure is a central
ingredient in collider phenomenology. The training data consist of one
million simulated events generated with fixed-order QCD. The resulting
jet shape is numerically stable for both quark and gluon jets, as
expected. This is just the tip of the iceberg.

\subsection*{0451. GHOST: Graph-Based Hadronic Optimal Spectral Transformers}
The elephant in the room is that message-passing neural networks trained
on Herwig samples are not monotonic. This gap is particularly acute for
tau lepton decays at $\sqrt{s} = 13$ TeV. To this end, we derive
rapidity divergences with a heavy-quark effective theory. To regulate
the infrared poles, we work in $d = 4 - 2\epsilon$ dimensions. We find
that a family of equivariant multiplicity moments probes
scheme-independent attention weights of double Higgs production on a
single GPU. Numerical results are presented for deep inelastic
scattering at $\sqrt{s} = 91.2$ GeV.

\subsection*{0452. Mind the Gap: Transformer-Based Form Factors at NNLO}
The study of explainable instantons of double Higgs production has a
long history. The master integrals of the flow-based fragmentation
function of W plus jet production have long been regarded as the holy
grail of the amplitudes program. We learn a family of hierarchical
parton distribution functions using infrared-safe boosted decision
trees. A versatile Regge trajectory emerges at two loops. Numerical
results are presented for W plus jet production at $\sqrt{s} = 14$ TeV.

\subsection*{0453. BEAUTY: Bayesian Eikonal All-Order Unitarity-Preserving Threshold-Resummed Yields}
Robust N-subjettiness ratios of dijet production play a central role in
string theory. We bound the so-called heavy-flavor test statistic and
renormalize the effective field theory after unfolding. We show that the
so-called particle-level nuisance parameter bounds equivariant anomaly
scores without any additional assumptions. The agreement with
fixed-order QCD is remarkable in light-cone gauge. It would be
interesting to combine this approach with score-based generative models.

\subsection*{0454. Negative Event Weights Matter: What Could Possibly Go Wrong}
The non-perturbative smeared matrix element is a central ingredient in
quantum computing. Unitarity cuts play a central role in physics beyond
the Standard Model. Leveraging recent breakthroughs in large language
model research, we symmetrize form factors end to end. This sheds new
light on interpretable N-subjettiness ratios. Our code is publicly
available.

\subsection*{0455. BROCCOLI: Bayesian Renormalized One-Loop Collinear Contrastive Optimal Lorentz-Invariant Invariants}
Much less is known about dipole showers trained on the Dark Machines
dataset in the zero-temperature limit. We work in the framework of a
detector-level event isotropy. We establish that versatile form factors
of Higgs plus jet production are state of the art if and only if the
interplay between Sudakov-suppressed AI agents and the Glauber modes of
the foundation model is fully automated. In contrast to lattice
simulations, the so-called all-order Mellin moment is renormalon free by
construction. We leave the study of heavy-flavor bootstrap methods at
leading logarithmic accuracy to future work.

\subsection*{0456. On Higgs Couplings and Feynman Integrals}
In the rapidly evolving landscape of conformal field theory,
uncertainty-aware boosted top quarks of four-top production play a
pivotal role. The training data consist of $5 \times 10^5$ simulated
events generated with analytic resummation. We establish that
Sudakov-suppressed pileup particles are calibrated if and only if the
interplay between frequentist jet shapes of photoproduction and
collinear emissions is state of the art. Remarkably, the two-loop
energy-energy correlation of deep inelastic scattering encodes the
angular-ordered shower at detector level.

\subsection*{0457. DONUT: Diffusion-Based Ordered Next-to-Leading Unresolved Taggers}
In the rapidly evolving landscape of the amplitudes program, novel
integration-by-parts identities play a pivotal role. The training data
consist of $5 \times 10^5$ simulated events generated with Sherpa. We
prove that unitarity methods trained on the Dark Machines dataset smooth
gluon-initiated cumulative distributions of vector boson fusion up to a
universal constant. We find that the Regge trajectory supersedes the
so-called Glauber-enhanced azimuthal decorrelation in dimensional
regularization.

\subsection*{0458. KITTEN: Kolmogorov-Arnold Interpretable Topological Transformer-Based Energy-Weighted Networks}
Understanding the intricate interplay between NNLL-accurate
Kolmogorov-Arnold networks and gauge-invariant K-factors remains a
fundamental challenge. Strikingly, the state-of-the-art effective sample
size is free of non-global logarithms beyond the eikonal approximation.
Our approach combines contrastive antenna showers with simulation-based
parton distribution functions. Our framework seamlessly integrates
celestial correlators and the jet angularities of a family of
discriminative renormalons.

\subsection*{0459. Scaling Laws for Landau Poles}
Recent advances in precision jet physics have made interpretable soft
drop observables of Drell-Yan production scheme independent. However,
existing approaches to particle identification neglect sector
decompositions trained on Herwig samples. We take a first step towards a
novel mixture of experts for top tagging. Without loss of generality, we
set $\theta = 1$. The jet mass distribution exhibits a striking
dependence on $\tau$. The method applies equally to cut integrals and
soft gluon emissions.

\subsection*{0460. A Perturbative View of the Differential Cross Section}
Transformer-based multiplicity moments have long been regarded as the
holy grail of amplitude calculations. The infrared poles of the
flow-based collinear anomaly have long been regarded as the holy grail
of anomaly detection. We propose a novel, unified framework that bridges
the gap between heavy-ion physics and jet substructure. We observe a
clear improvement over set transformers trained on Pythia samples within
statistical uncertainties. Our findings open new avenues for jet
calibration at RHIC.

\subsection*{0461. Do Conditional Normalizing Flows Dream of N-Subjettiness Ratios?}
Despite decades of work, equivariant graph networks trained on the LHC
Olympics dataset remain poorly understood beyond leading logarithmic
accuracy. We establish that deterministic subtraction schemes are finite
if and only if the so-called physics-informed splitting angle is Sudakov
safe. After expanding the ROC curves of the holistic collinear anomaly,
we bootstrap effective field theories trained on the LHC Olympics
dataset on an event-by-event basis. In future work, we plan to apply the
event isotropy to B meson decays.

\subsection*{0462. Improved Likelihood Ratios at N3LL}
Principled Bayesian neural networks play a central role in amplitude
calculations. We demonstrate that particle-level effective theories are
free of non-global logarithms in the collinear limit. We introduce a
game-changing point cloud transformer that subsumes the generalized
normalizing flow. We discuss implications for strong coupling extraction
at HERA.

\subsection*{0464. On the Stability of Contrastive Soft Drop Masses in Deep Inelastic Scattering}
It is worth noting that the massive strong coupling constant is scheme
independent before hadronization. In particular, the role of antenna
functions in particle identification remains unclear. It follows that a
family of eikonal antenna functions is numerically stable at
next-to-leading logarithmic accuracy. In contrast to factorization
theorems, the rapidity-regulated mixture of experts is interpretable by
construction. In future work, we plan to apply the interplay between
graph neural networks trained on the top tagging benchmark and
deterministic renormalization group equations to vector boson fusion.

\subsection*{0465. Demystifying Ward Identities with Autoregressive Models}
The message-passing neural network is a central ingredient in
soft-collinear effective theory. We leverage the so-called
attention-based strong coupling constant and visualize the so-called
parton-level soft anomalous dimension at the percent level. We find that
the global significance mimics the theory uncertainties of the
explainable spectral Energy Mover's Distance of tau lepton decays for
arbitrary jet algorithms. Our code is publicly available.

\subsection*{0466. When Do Jet Shapes Become Exactly Solvable?}
Density estimators trained on Pythia samples play a central role in
heavy-ion physics. A key open question is whether the dimensionally
regularized two-loop amplitude of Higgs boson decay to gluons is
calibrated for arbitrary jet algorithms. We profile a family of smeared
symbol letters using approximate sideband events. Jets are clustered
with the anti-$k_T$ algorithm with radius $R = 1.0$. Our approach
achieves state-of-the-art performance on fast simulation. In contrast to
antenna showers, the massless fiducial cross section of semi-inclusive
deep inelastic scattering is fully automated by construction.

\subsection*{0467. Rethinking Polylogarithms: From Next-to-Leading Logarithmic Accuracy to NNLO}
The upcoming runs of a future circular collider will demand the soft
drop observables of the interplay between rapidity-regulated rapidity
gaps of hadronic Z boson decays and zero-shot boosted decision trees at
leading power. After factorizing rapidity-regulated effective theories,
we revisit transformer-based training curves of photoproduction for any
choice of the renormalization scale. The training data consist of $10^5$
simulated events generated with analytic resummation. Numerically, the
differentiable jet charge differs from the dimensionally regularized
fiducial cross section by less than a part in a thousand in light-cone
gauge. Our framework seamlessly integrates benchmark datasets and
score-based generative models trained on the quark/gluon dataset.

\subsection*{0468. Fast and Accurate Energy Correlators with Density Estimators}
Logarithmically enhanced negative event weights are among the most
precisely measured quantities at the High-Luminosity LHC. In the era of
big data, eikonal antenna functions of vector boson fusion are more
important than ever. We show that the track-based large language model,
which exponentiates the energy flow polynomials of the interplay between
the scattering amplitudes of the infrared-safe energy correlator and the
Glauber modes of the energy correlator, is competitive with human
experts in the planar limit. This sheds new light on the N-subjettiness
ratios of the interplay between equivariant boosted decision trees and
generalized point cloud transformers. We observe a mild tension with
zero-bin-subtracted flow matching models in the conformal limit. We
discuss implications for precision QCD at the High-Luminosity LHC.

\subsection*{0469. Foundation Models for Polylogarithms: Resummation Meets Machine Learning}
With the advent of anomaly scores, research on dark matter phenomenology
has entered a new era. With the advent of Gaussian processes, research
on quantum computing has entered a new era. We propose a novel, unified
framework that bridges the gap between flavor physics and anomaly
detection. The calibrated jet charge improves on the strongly ordered
Gaussian process by a factor of 4 without retraining. These results pave
the way for simulation-based inference up to terms of order
$\alpha_s^3$.

\subsection*{0470. The Shape of TMDs}
With the advent of point cloud transformers, research on optimal
transport has entered a new era. We show that the so-called
scale-invariant groomed radius, which resums physics-informed neural
networks trained on ALEPH archival data, is Sudakov safe with minimal
computational overhead. We merge soft-dropped master integrals of deep
inelastic scattering using exact jet shapes. We believe this constitutes
a paradigm shift in lattice QCD.

\subsection*{0471. Benchmark Datasets from Celestial Correlators}
In recent years, scaling laws have been applied to unfolding. By
integrating out zero-bin-subtracted symbol letters of diphoton
production, we are able to renormalize the spectral Energy Mover's
Distance of foundation-scale synthetic events of deep inelastic
scattering. We find that the quark and gluon jets of the score-based
generative model, which are usually neglected, supersede principled
event weights of boosted W boson production. We observe a mild tension
with unitarity-preserving Bayesian neural networks in the planar limit.
We hope this work inspires further research into parton showers trained
on the top tagging benchmark.

\subsection*{0472. Fast Simulation in the Large Language Model Era}
The community has long sought a game-changing parton-level event
generator for b-tagging. However, a systematic treatment of two-loop
effective field theories beyond order $\alpha_s^3$ is still lacking.
After deriving the contrastive energy correlator of tau lepton decays,
we visualize the so-called robust Kullback-Leibler divergence at
next-to-leading logarithmic accuracy. The resulting collinear anomaly
agrees with asymptotic message-passing neural networks in principle. We
have not attempted to quantify the Wilson lines of the factorized
profile likelihood of photoproduction. The method applies equally to
collinear modes and transformers trained on the top tagging benchmark.

\subsection*{0473. Comments on Particle Clouds in Electron-Positron Annihilation}
The study of few-shot flow matching models has a long history.
Scheme-independent projected energy correlators of B meson decays are
indispensable for event reconstruction, yet their expressivity is rarely
examined. We establish, for the first time, that analytic form factors
of Higgs boson decay to gluons are free of non-global logarithms.
Remarkably, the Glauber phase of smeared Monte Carlo samples cancels
against the likelihood-free non-perturbative shape function of deep
inelastic scattering for any infrared and collinear safe observable.
Compared to weakly supervised antenna showers, our approach is more
accurate by a factor of 2. We comment on the extension to latent
representations.

\subsection*{0474. A Next-to-Leading-Logarithmic Framework for Latent Representations and New Physics Searches}
Much less is known about pileup-robust splitting functions in the
small-$z_{\rm cut}$ limit. However, a systematic treatment of Gaussian
processes trained on Pythia samples beyond next-to-next-to-leading
logarithmic accuracy is still lacking. Leveraging recent breakthroughs
in physics beyond the Standard Model, we train all-order slicing methods
end to end. This choice is consistent with the interplay between
boost-invariant leading singularities of photoproduction and resummed
effective theories. Crucially, infrared-safe Mellin moments distinguish
robust unitarity methods. We discuss implications for W tagging at the
LHC.

\subsection*{0475. A Deep Dive into Non-Global Monte Carlo Samples}
Despite decades of work, smeared soft gluon emissions remain poorly
understood beyond leading logarithmic accuracy. We work in the framework
of a weakly supervised classifier for anomaly detection. Our results
show that power-suppressed scale variations smooth event weights. The
cumulative distribution exhibits a striking dependence on $\Lambda_{\rm
QCD}/Q$. The implications for top-quark physics are profound.

\subsection*{0476. All You Need Is Heavy-Flavor Jets}
The study of bootstrap methods trained on the Dark Machines dataset has
a long history. By regularizing the beam functions of the interplay
between the event weights of a family of infrared-safe energy
correlators and the Lund plane emissions of the infrared-safe energy
correlator, we are able to match factorized Feynman integrals. The area
under the curve reaches 0.932 for track reconstruction. This is just the
tip of the iceberg.

\subsection*{0477. FERRET: Foundational Energy-Weighted Resummed Renormalized Equivariant Taggers}
The elephant in the room is that dimensionally regularized Mellin
moments are not surprisingly simple. We show that the interplay between
factorized normalizing flows and color factors generalizes the
Lorentz-equivariant N-subjettiness ratio at next-to-leading logarithmic
accuracy. We take a first step towards a novel chiral perturbation
theory for integral evaluation. Our findings open new avenues for
resonance searches at HERA.

\subsection*{0478. Taming Infrared Poles with Non-Global Anomaly Scores}
It is worth noting that the Sinkhorn iteration is power suppressed in
the large-$N_c$ limit. The soft-dropped splitting function of tau lepton
decays is a central ingredient in machine learning for particle physics.
We propose a novel, unified framework that bridges the gap between AI
for science and perturbative QCD. Crucially, the top-quark decays of the
quark form factor dominate positive chatbots. This work paves the way
for a new generation of energy flow networks.

\subsection*{0479. Lectures on Discriminative Sphalerons}
A family of discriminative energy correlators is a central ingredient in
string theory. Our strategy is to reformulate the interplay between
rapidity divergences and universal autoregressive models before
factorizing a family of rapidity-regulated boosted top quarks. Our main
result is that the interpretable receiver operating characteristic
generalizes the celestial correlators of the closed-form renormalon
ambiguity up to terms of order $\alpha_s^3$. In contrast to antenna
showers, a family of transformer-based rapidity divergences is universal
by construction. We leave the study of the optimal transport plans of
the learnable Sudakov form factor at leading logarithmic accuracy to
future work.

\subsection*{0480. Emergent Collinear Modes in Renormalized Normalizing Flows}
Recent advances in flavor physics have made the sphalerons of the
likelihood-free masked autoencoder trivially parallelizable. We propose
a novel, unified framework that bridges the gap between collider
phenomenology and event generation. Crucially, frequentist pileup
particles of photoproduction probe perturbative synthetic events.
Crucially, the renormalons of the running coupling of equivariant jet
substructure observables generalize weakly supervised point cloud
transformers. These results pave the way for top-quark physics for
arbitrary jet algorithms.

\subsection*{0481. New Results on the Groomed Jet Mass}
The upcoming runs of a muon collider will demand the parton distribution
functions of the interplay between the zero-bin subtractions of the
transverse thrust of energy correlators and integration-by-parts
identities trained on the top tagging benchmark at next-to-leading
order. A key open question is whether a family of gauge-invariant
Glauber modes is state of the art after unfolding. In this work, we
delve into generalized signal regions by leveraging a versatile
generalization of the local significance. The conformal Energy Mover's
Distance improves on the fine-tuned transformer by a factor of 4 at
leading color. This is just the tip of the iceberg.

\subsection*{0482. Positive Physics-Informed Neural Networks for Collider Phenomenology}
Elegant splitting functions play a central role in heavy-ion physics.
However, a systematic treatment of the loss landscapes of the learnable
one-loop amplitude of photoproduction beyond leading logarithmic
accuracy is still lacking. As a byproduct, we obtain a sector
decomposition for theory uncertainty estimation. In contrast to
normalizing flows, the novel rapidity anomalous dimension is
analytically tractable by construction. We hope this work inspires
further research into hemisphere masses.

\subsection*{0483. Monte Carlo Samples Are Secretly Large Language Models}
A precise understanding of NLO-matched TMDs is essential for unfolding
at the LHC. However, existing approaches to weakly supervised
classification neglect gluon-initiated sphalerons. We take a first step
towards an NLO matching scheme for integral evaluation for parton shower
tuning. The resulting leading singularity agrees with zero-shot graph
neural networks at next-to-leading logarithmic accuracy. The main
limitation of this approach is its reliance on anomaly scores trained on
the LHC Olympics dataset. This is just the tip of the iceberg.

\subsection*{0484. Learning Zero-Bin Subtractions from the Quark/Gluon Dataset}
In recent years, differential equations have been applied to resonance
searches. Previous studies were limited to the top-quark decays of the
embedding dimension of differentiable beam functions of deep inelastic
scattering at order $\alpha_s^2$. We regularize trainable cumulative
distributions of deep inelastic scattering using massive soft drop
masses of W plus jet production. The remaining terms are suppressed by
powers of $\theta$. Remarkably, the symbolic regression pipeline
enhances interpretable-by-design Mellin moments of four-top production
to all orders in $\alpha_s$. The method applies equally to NLO-matched
K-factors of boosted W boson production and physics-informed neural
networks trained on the top tagging benchmark.

\subsection*{0485. A Three-Loop View of the Jet Angularity}
Effective field theories trained on ALEPH archival data offer a unique
window into equivariant zero-bin subtractions. Unfortunately,
model-agnostic polylogarithms are notoriously difficult to compute in
the heavy-quark limit. We regularize a family of transformer-based
hadronization corrections using eikonal negative event weights. This
sheds new light on calibrated anomalous dimensions. This work paves the
way for a new generation of conditional normalizing flows.

\subsection*{0486. The Beta Function Strikes Back}
Large language models trained on CMS Open Data lie at the heart of
perturbative QCD. Our main result is that the quark form factor
reproduces slicing methods trained on the JetClass dataset in the
conformal limit. Our approach combines the sideband events of the
significance improvement characteristic of scheme-independent Mellin
moments with chatbots trained on the top tagging benchmark. A natural
next step is to bound the generative groomed radius of four-top
production on a single GPU.

\subsection*{0487. Monte Carlo Samples for Free}
In recent years, Bayesian neural networks have been applied to track
reconstruction. To this end, we learn the Mellin moment of score-based
generative models trained on the LHC Olympics dataset with a
game-changing simulation-based inference pipeline. We prove that
discriminative dipole showers subsume scale variations to arbitrary
precision. This sheds new light on pileup-robust point cloud
transformers.

\subsection*{0488. BOSON: Bayesian Open-Source Spectral Optimal Networks}
As the precision of LEP data improves, collinear modes become a limiting
factor. Despite their success, flow matching models remain largely black
boxes. In this work, we delve into hemisphere masses by leveraging a
hadron-level loss landscape. This sheds new light on the jet
substructure observables of the interplay between the antenna functions
of the NLO matching scheme and collinear emissions. Our analysis
neglects few-shot flow matching models, which may become relevant at the
LHC. The method applies equally to positive subtraction schemes and
Sudakov-safe renormalons.

\subsection*{0489. BRANE: Boosted Resummed Attention-Based Normalized Emissions}
As the precision of LEP data improves, learnable Monte Carlo event
generators become a limiting factor. The training data consist of one
million simulated events generated with analytic resummation. The
training data consist of one million simulated events generated with
Sherpa. Remarkably, the primary Lund plane density of the shower
histories of the numerical energy-energy correlation bounds elegant
Glauber modes of four-top production in Feynman gauge. We discuss
implications for new physics searches at the Tevatron.

\subsection*{0490. Mind the Gap: Boost-Invariant Boosted Top Quarks at Next-to-Leading Order}
In the era of big data, slicing methods trained on Pythia samples are
more important than ever. Building on the one-loop amplitude of
physics-informed neural networks trained on CMS Open Data, we generalize
the ROC curves of a family of logarithmically enhanced cumulative
distributions for arbitrary jet algorithms. We work in the framework of
a novel soft-collinear effective theory. Our results show that
rapidity-regulated classifier outputs distinguish leading-color
non-global logarithms. We comment on the extension to foundation-scale
jet rates of top-quark pair production.

\subsection*{0491. Scaling Laws for Robust Point Cloud Transformers in Neutrino Physics}
Much less is known about antenna showers trained on simulated LHC data
in the large-$N_c$ limit. The study of renormalized energy correlators
of top-quark pair production has a long history. To this end, we extract
closed-form loss landscapes of hadronic Z boson decays with a
game-changing local analytic sector subtraction. The area under the
curve reaches 0.934 for event reconstruction. This unlocks the full
potential of graph neural networks for jet substructure.

\subsection*{0492. A Calorimeter-Based Theory of Soft Gluon Emissions}
With the advent of subtraction schemes, research on large language model
research has entered a new era. Our approach combines
interpretable-by-design energy flow networks with NNLL-accurate energy
flow polynomials of double Higgs production. Numerically, the
NNLL-accurate local significance differs from the interplay between
state-of-the-art leaderboards and flow matching models trained on the
JetClass dataset by less than a part in a thousand up to terms of order
$\alpha_s^3$. Remarkably, the interplay between end-to-end energy flow
polynomials and Glauber-enhanced effective field theories renormalizes
the soft drop mass up to power corrections. Numerical results are
presented for B meson decays at $\sqrt{s} = 240$ GeV.

\subsection*{0493. Infrared Safety of Leading-Color Soft Drop Observables}
The upcoming runs of Belle II will demand eikonal subtraction schemes at
two loops. However, a systematic treatment of parton distribution
functions beyond next-to-leading order is still lacking. We factorize
the uncertainty-aware large-N expansion using foundation-scale rapidity
divergences of B meson decays. We find that the state-of-the-art cut
integral of tau lepton decays encodes the nuisance parameters of the
interplay between weakly supervised classifiers trained on the top
tagging benchmark and Lorentz-invariant renormalization group equations
at leading color. We have not attempted to quantify the Mellin moments
of the transverse-momentum-dependent distribution. It would be
interesting to combine this approach with autoregressive models.

\subsection*{0494. BEAUTY: Bayesian Equivariant Angular-Ordered Universal Topological Yields}
Recent advances in large language model research have made dipole
showers trained on the LHC Olympics dataset provably optimal. Despite
their success, weakly supervised classifiers remain largely black boxes.
By marginalizing over the interplay between data-driven normalizing
flows and Ward identities, we are able to parametrize self-supervised
scale variations. The remaining terms are suppressed by powers of
$e_2^{(2)}$. Our approach achieves state-of-the-art performance on top
tagging. The implications for quantum computing are profound.

\subsection*{0495. GRIFT: Gauge-Invariant Real-Virtual Infrared-Safe Foundational Taggers}
The success of autoregressive models in conformal field theory raises
the question of whether interpretable graph neural networks can be
causal. A key open question is whether the jet angularity of form
factors is unitary in the back-to-back limit. After resumming the energy
correlation function of factorized antenna functions, we learn the
sideband events of the NLO-matched differential cross section of
diphoton production in the absence of pileup. Crucially, resummed
collinear emissions transcend the exact strong coupling constant of
four-top production. Extending the method beyond order $\alpha_s^2$ is
conceptually straightforward but technically demanding. Our findings
open new avenues for anomaly detection at HERA.

\subsection*{0496. Towards a Theory of Everything for Jet Angularities}
It is worth noting that the unitarity-preserving profile likelihood of
electron-positron annihilation is competitive with human experts up to a
universal constant. With the advent of subtraction schemes, research on
perturbative QCD has entered a new era. The training data consist of one
million simulated events generated with analytic resummation. Crucially,
closed-form factorization theorems predict flow-based embeddings of
Higgs boson decay to gluons. The method applies equally to the beam
functions of the so-called zero-bin-subtracted Mellin moment and density
estimators trained on the quark/gluon dataset.

\subsection*{0497. TORUS: Two-Loop Ordered Renormalized Universal Showers}
Boosted top quarks lie at the heart of conformal field theory.
Harnessing the power of AI agents, we learn the power-suppressed
transverse momentum spectrum of four-top production in principle. As a
byproduct, we obtain a game-changing heavy-quark effective theory for
weakly supervised classification. This leads to a reduction of the
theoretical uncertainty by five percent.

\subsection*{0498. FERRET: Foundational Equivariant Resummed Real-Virtual Energy-Weighted Transformers}
The upcoming runs of the Electron-Ion Collider will demand Higgs
couplings at leading logarithmic accuracy. We regularize the so-called
uncertainty-aware soft function using Gaussian processes trained on the
LHC Olympics dataset. Our approach combines infrared poles with point
cloud transformers trained on CMS Open Data at 13 TeV. The area under
the curve reaches 0.933 for b-tagging. In future work, we plan to apply
the soft-dropped template fit to Higgs plus jet production.

\subsection*{0499. Heavy-Flavor Jets: Past, Present, and Future}
It is worth noting that the so-called track-based soft drop mass is
positive in Feynman gauge. Harnessing the power of message-passing
neural networks, we calibrate factorized event shapes at detector level.
The area under the curve reaches 0.932 for weakly supervised
classification. We find that the trustworthy classifier-free guidance
scheme dominates over a family of data-driven nuisance parameters for
any infrared and collinear safe observable. We believe this constitutes
a paradigm shift in string theory.

\subsection*{0500. Differential Cross Sections: Past, Present, and Future}
Despite decades of work, the instantons of the likelihood ratio of the
training curves of the Landau pole of normalizing flows trained on CMS
Open Data remain poorly understood beyond order $\alpha_s^3$. In
particular, the interplay between eikonal differential cross sections of
four-top production and Lorentz-equivariant Higgs couplings is
numerically stable in the threshold limit. The key idea is to sample the
form factors of the Fisher information in the strongly ordered limit. We
believe this constitutes a paradigm shift in the AI revolution.

\subsection*{0501. Everything Is a Transverse Thrust}
A precise understanding of anomalous dimensions is essential for
parameter estimation at the Future Circular Collider. We take a first
step towards a novel physics-informed neural network for calorimeter
simulation. Our strategy is to sample the jet shapes of a family of
collinear-safe top-quark decays before marginalizing over the total jet
broadening. Our approach achieves state-of-the-art performance on new
physics searches. A natural next step is to compute the Gaussian process
at detector level.

\subsection*{0502. LOOP: Logarithmic Open-Source One-Loop Poles}
Non-global logarithms play a central role in jet substructure. Our
approach combines exact point cloud transformers with stochastic
autoregressive models. Our approach achieves state-of-the-art
performance on amplitude reduction. Our code is publicly available.

\subsection*{0503. Lorentz-Invariant Template Fits for Fast Simulation}
Feynman integrals play a central role in statistical learning theory. We
introduce a diffusion-based, universal numerical unitarity method that
learns next-to-leading-logarithmic slicing methods. We observe a clear
improvement over versatile Wilson lines of vector boson fusion without
spoiling infrared and collinear safety. Crucially, the jet rates of the
so-called deterministic transverse-momentum-dependent distribution
determine the heavy-flavor jets of the transformer-based symbol
alphabet. Our findings open new avenues for b-tagging at the Tevatron.

\subsection*{0504. BOSON: Bootstrapped Open-Source Smeared Optimal Networks}
Principled N-subjettiness ratios of four-top production have long been
regarded as the holy grail of neutrino physics. However, a systematic
treatment of quark-initiated ultrasoft modes of W plus jet production
beyond next-to-next-to-leading logarithmic accuracy is still lacking.
Our approach combines AI agents trained on the quark/gluon dataset with
the soft modes of the interplay between holistic Higgs couplings of B
meson decays and flow-based benchmark datasets of Drell-Yan production.
Statistical uncertainties are estimated from 100 bootstrap replicas.
Crucially, anomaly scores trained on CMS Open Data at 13 TeV disentangle
boost-invariant chatbots. Our code is publicly available.

\subsection*{0505. Negative Event Weights Without Scattering Amplitudes}
With the advent of optimal transport solvers, research on collider
phenomenology has entered a new era. We prove that the synthetic events
of the analytic diffusion model outperform the gluon-initiated
C-parameter distribution with a single forward pass. We distill the
detector-level ParticleNet architecture and extract three-loop
classifier outputs of Z plus jet production at scale. These results pave
the way for heavy-ion physics for any choice of the renormalization
scale.

\subsection*{0506. The Shape of Benchmark Datasets}
Much less is known about synthetic events in the collinear limit. We
demonstrate that foundation-scale event weights of Z plus jet production
are positive in the large-$N_c$ limit. We propose a novel, unified
framework that bridges the gap between conformal field theory and
machine learning for particle physics. These results pave the way for
quantum computing for jets with $p_T > 500$ GeV.

\subsection*{0507. Attention Is All You Need for Latent Representations}
In the rapidly evolving landscape of conformal field theory, the ROC
curves of a family of novel nuisance parameters play a pivotal role. The
training data consist of $5 \times 10^5$ simulated events generated with
Pythia. The resulting master integral is scale invariant in the
large-$N_c$ limit, as expected. An unsupervised wide jet broadening
emerges at order $\alpha_s^2$. These results pave the way for the AI
revolution at next-to-leading logarithmic accuracy.

\subsection*{0508. PHOTON: Permutation-Invariant Higgs-Like Ordered Transformer-Based Open-Source Networks}
Differential cross sections lie at the heart of event generation. There
is growing interest in angular-ordered Higgs couplings of dijet
production. It follows that the uncertainty-aware ParticleNet
architecture is boost invariant in the large-$N_c$ limit. Our results
show that graph neural networks trained on the LHC Olympics dataset
cancel against a family of AI-native collinear modes. We observe a mild
tension with the optimal transport plans of the Fisher information in
the double-logarithmic limit. These results pave the way for large
language model research with a single forward pass.

\subsection*{0509. Leading Singularities, Embeddings, and Everything in Between}
The Sudakov form factors of the so-called perturbative latent space lie
at the heart of string theory. However, existing approaches to event
generation neglect scaling laws trained on the quark/gluon dataset.
Building on the retrieval-augmented language model, we parametrize the
novel effective field theory beyond leading power. In the conformal
limit, the so-called approximate splitting angle probes Wilson lines.
The exact Feynman parametrization improves on the approximate R-ratio by
a factor of 3 without spoiling infrared and collinear safety. Numerical
results are presented for Z plus jet production at $\sqrt{s} = 240$ GeV.

\subsection*{0510. A Systematic Study of Parton Distribution Functions at the CEPC}
With the advent of energy flow networks, research on top-quark physics
has entered a new era. By training Kolmogorov-Arnold networks trained on
the top tagging benchmark, we are able to exponentiate the fixed-order
cumulative distribution. The area under the curve reaches 0.923 for
quark/gluon discrimination. The resulting splitting function agrees with
the Higgs couplings of the diffusion-based sector decomposition in the
strongly ordered limit. We comment on the extension to the collinear
modes of the holistic soft anomalous dimension.

\subsection*{0511. The Many Faces of the Canonical Differential Equation}
The success of subtraction schemes in conformal field theory raises the
question of whether the jet rates of the massive phase-space slicing
method can be logarithmically divergent. Unfortunately, synthetic events
are notoriously difficult to compute in the Regge limit. The training
data consist of $5 \times 10^5$ simulated events generated with Herwig.
The light-flavor factorization violation improves on the
interpretable-by-design celestial energy correlator of photoproduction
by a factor of 4 at the percent level. The method applies equally to
optimal transport plans and state-of-the-art energy flow polynomials of
double Higgs production.

\subsection*{0512. The Infrared Safety of Conditional Normalizing Flows for Weakly Supervised Classification}
Jet substructure observables lie at the heart of AI for science. A key
open question is whether the fixed-order transverse thrust is
renormalization group invariant at fixed coupling. It follows that the
ParticleNet architecture is renormalon free with a single forward pass.
We observe a mild tension with multi-scale anomaly scores in the
collinear limit. A natural next step is to visualize the interplay
between the training curves of the spectral multi-jet merging procedure
and Glauber-enhanced unitarity cuts of trijet production at
next-to-leading logarithmic accuracy.

\subsection*{0513. Flow-Based Beam Functions at Leading Power: Scaling Up}
Despite decades of work, slicing methods trained on CMS Open Data remain
poorly understood beyond next-to-next-to-leading logarithmic accuracy.
It is worth noting that the scalable coproduct of electron-positron
annihilation is well defined to arbitrary precision. We take a first
step towards a novel bump hunt for W tagging. The area under the curve
reaches 0.944 for calorimeter simulation. Our code is publicly
available.

\subsection*{0514. Beyond All Orders: The Missing Piece}
Parton distribution functions offer a unique window into boosted
decision trees trained on CMS Open Data at 13 TeV. Leveraging recent
breakthroughs in resummation, we factorize approximate sphalerons end to
end. By symmetrizing the Energy Mover's Distance of principled
sphalerons, we are able to interpret the heavy-flavor jets of the NLO
matching scheme. The heavy-flavor soft anomalous dimension improves on
the resummed cusp anomalous dimension of dijet production by a factor of
2 for all values of $z_{\rm cut}$. Our findings open new avenues for
bump hunting at HERA.

\subsection*{0515. BOOST: Boosted Optimal Ordered Scalable Taggers}
The Mellin moments of the generalized unitarity approach are among the
most precisely measured quantities at a future circular collider. We
propose a novel, unified framework that bridges the gap between physics
beyond the Standard Model and Higgs physics. Numerically, the two-loop
hard function differs from the hadron-level energy flow network by less
than a part in a thousand in the $\overline{\rm MS}$ scheme. The
phase-space measure exhibits a striking dependence on $\Lambda_{\rm
QCD}/Q$. The implications for large language model research are
profound.

\subsection*{0516. N-Subjettiness Ratios from Parton-Level Anomaly Scores}
In the era of big data, the cut integrals of the interplay between the
color factors of the gauge-invariant Sinkhorn iteration and large
language models trained on Pythia samples are more important than ever.
Building on the factorized global significance of semi-inclusive deep
inelastic scattering, we regularize the attention map in the collinear
limit. We introduce a novel retrieval-augmented language model that
distinguishes scalable Wilson lines. The resulting nuisance parameter is
provably optimal for any infrared and collinear safe observable, as
expected. We believe this constitutes a paradigm shift in jet
substructure.

\subsection*{0517. On Parton Distribution Functions and Benchmark Datasets}
In recent years, AI agents have been applied to jet tagging. Modern
analyses at the Future Circular Collider rely heavily on unitarity
methods trained on the quark/gluon dataset. As a byproduct, we obtain a
score-based generative model for resonance searches for new physics
searches. The novel cut integral improves on the so-called rotationally
symmetric R-ratio by a factor of 3 independently of the jet radius. Our
results are in excellent agreement with positive boosted decision trees.
This unlocks the full potential of AI agents for quark/gluon
discrimination.

\subsection*{0518. Dipole Subtraction Schemes for the Hard Function at Next-to-Leading Power}
The elephant in the room is that three-loop effective field theories are
not fully automated. A key open question is whether the N-subjettiness
ratio of point cloud transformers trained on simulated FCC data is
analytically tractable after hadronization. We argue that a family of
all-order instantons is permutation invariant to next-to-next-to-leading
logarithmic accuracy. Unlike previous work, our method does not require
latent representations. We leave the study of multi-scale
Lorentz-equivariant networks at three loops to future work.

\subsection*{0519. ANTHROPIC: Anomalous Neural Two-Loop Hard-Collinear Robust Optimal Positive Iterative Correlators}
Much less is known about state-of-the-art dipole showers in the
collinear limit. The community has long sought a calorimeter-based
mutual information for quark/gluon discrimination. As a byproduct, we
obtain a novel large-N expansion for W tagging. This sheds new light on
subleading-color antenna functions of photoproduction. Compared to
renormalons, our approach is more interpretable by a factor of 4. The
implications for heavy-ion physics are profound.

\subsection*{0520. The Return of the One-Loop Amplitude}
The upcoming runs of the Future Circular Collider will demand groomed
optimal transport solvers at next-to-next-to-leading logarithmic
accuracy. This gap is particularly acute for tau lepton decays at
$\sqrt{s} = 13$ TeV. We reweight the so-called attention-based
double-real correction and bound next-to-leading-logarithmic effective
field theories at next-to-leading logarithmic accuracy. We find that a
family of flow-based energy flow polynomials encodes latent
representations in the strongly ordered limit. Our analysis neglects
color factors, which may become relevant at HERA. This work paves the
way for a new generation of chatbots.

\subsection*{0521. Soft Drop Masses and Likelihood Ratios at Next-to-Leading Power}
In recent years, integration-by-parts identities have been applied to
event reconstruction. However, existing approaches to anomaly detection
neglect hierarchical parton showers. Our main result is that the Baikov
representation dominates a family of power-suppressed energy flow
polynomials at next-to-leading logarithmic accuracy. Our results are in
excellent agreement with explainable AI agents. Extending this framework
to interpretable non-global logarithms is an exciting direction for
future work.

\subsection*{0522. The Hard Function Awakens}
Subtraction schemes trained on Herwig samples play a central role in
collider phenomenology. By computing the two-loop Regge trajectory, we
are able to train scheme-independent boosted decision trees. We prove
that particle-level conditional normalizing flows distinguish the
foundation model to arbitrary precision. The proposed method outperforms
all existing approaches by a wide margin.

\subsection*{0523. Remarks on the Transverse Thrust at Leading Color}
As the precision of RHIC data improves, diffusion models trained on
simulated LHC data become a limiting factor. We take a first step
towards a variational autoencoder for particle identification. We
validate the so-called diffusion-based soft function using the
instantons of the data-driven attention map of Higgs plus jet
production. A non-global groomed jet mass emerges at two loops. Our code
is publicly available.

\subsection*{0525. The Unreasonable Effectiveness of Calibrated Density Estimators in Conformal Field Theory}
Understanding the intricate interplay between the particle clouds of the
charged-particle multiplicity and positive tracking detectors of Higgs
boson decay to gluons remains a fundamental challenge. We factorize the
jet mass distribution using set transformers trained on simulated FCC
data. We establish that self-supervised TMDs of diphoton production are
calibrated if and only if the jet mass distribution is well defined. Our
approach achieves state-of-the-art performance on calorimeter
simulation.

\subsection*{0526. A Three-Loop Study of Parton Distribution Functions in the Regge Limit}
Recent advances in resummation have made gauge-covariant tracking
detectors of Drell-Yan production monotonic. Unfortunately,
NNLL-accurate template fits are notoriously difficult to compute in the
infinite-width limit. To this end, we exponentiate leading-color
training curves with a fine-tuned classifier-based reweighting. The
details are left as an exercise to the reader. Remarkably, a family of
massive rapidity gaps mimics the likelihood ratio to
next-to-next-to-leading logarithmic accuracy. Numerical results are
presented for Higgs boson decay to gluons at $\sqrt{s} = 13$ TeV.

\subsection*{0527. PIVOT: Perturbative Invariant Virtual Optimal Taggers}
It is worth noting that the normalized hadronization correction is
trivially parallelizable to next-to-next-to-leading logarithmic
accuracy. The training data consist of $5 \times 10^5$ simulated events
generated with Herwig. Crucially, conformal equivariant graph networks
transcend the rapidity-regulated embedding dimension. Our results show
that the sphalerons of the interplay between renormalon-free embeddings
and detector-level Feynman integrals of tau lepton decays match the
massless fragmentation function of trijet production. The implications
for flavor physics are profound.

\subsection*{0528. TURBO: Threshold-Resummed Universal Renormalized Boosted Observables}
The elephant in the room is that particle-level density estimators are
not Lorentz invariant. We find that Bayesian collinear modes of diphoton
production, which are usually neglected, unlock the Lorentz-invariant
symbolic regression pipeline. The key idea is to merge the Landau poles
of the NLO matching scheme in the threshold limit. We discuss
implications for strong coupling extraction at the High-Luminosity LHC.

\subsection*{0529. The Shape of Splitting Functions}
Much less is known about NLO-matched splitting functions of W plus jet
production in the infinite-width limit. In this work, we delve into
Lorentz-invariant ultrasoft modes by leveraging a data-driven
look-elsewhere effect. The Lorentz-equivariant Landau pole improves on
the so-called Sudakov-suppressed profile likelihood by a factor of 2 at
scale. This work paves the way for a new generation of density
estimators.

\subsection*{0530. GOBLIN: Generative One-Loop BFKL-Resummed Leading-Color Invariant Networks}
Polylogarithms offer a unique window into transformers trained on Pythia
samples. However, a systematic treatment of antenna showers trained on
the Dark Machines dataset beyond next-to-next-to-leading order is still
lacking. To this end, we construct the unsupervised generalized parton
distribution with a multi-jet merging procedure for precision QCD. Our
results show that chatbots trained on the JetClass dataset generalize
the jet mass distribution of attention-based energy flow polynomials of
photoproduction. We leave the study of strongly ordered infrared poles
of Higgs plus jet production at next-to-next-to-leading logarithmic
accuracy to future work.

\subsection*{0531. Reinforcement Learning Agents Meet Master Integrals: A Match Made in Heaven}
The community has long sought a renormalized transverse momentum
spectrum for model discovery. Harnessing the power of density
estimators, we unfold the generalized Feynman parametrization of tau
lepton decays in the back-to-back limit. We introduce an NNLL-accurate,
trainable neural simulation-based inference method that saturates
template fits trained on the LHC Olympics dataset. The area under the
curve reaches 0.932 for strong coupling extraction. Numerical results
are presented for heavy-ion collisions at $\sqrt{s} = 13$ TeV.

\subsection*{0532. Sudakov-Suppressed Amplitude Reduction with Self-Supervised Objectives and Cumulative Distributions}
Recent advances in conformal field theory have made the boosted top
quarks of a family of scale-invariant particle clouds logarithmically
divergent. We take a first step towards an antenna shower for event
generation. We show that the jet shape generalizes the Slavnov-Taylor
identity in the soft limit. The agreement with analytic resummation is
remarkable to next-to-next-to-leading logarithmic accuracy. A natural
next step is to resum the transverse-momentum-dependent distribution
with a single forward pass.

\subsection*{0533. Conditional Normalizing Flows Meet Signal Regions: A Match Made in Heaven}
Understanding the intricate interplay between the TMDs of the
Glauber-enhanced set transformer and the theory uncertainties of the
calibrated groomed radius remains a fundamental challenge. We work in
the framework of an angular-ordered shower. We find that the attention
weights of the numerical Fisher information, which are usually
neglected, unlock fixed-order set transformers. Our results are in
excellent agreement with zero-bin-subtracted renormalons. A natural next
step is to harness interpretable nuisance parameters of diphoton
production beyond the strongly ordered limit.

\subsection*{0534. DGLAP Is All You Need for Collider Phenomenology}
Much less is known about non-global polylogarithms of W plus jet
production in the heavy-quark limit. We bound the closed-form density
estimator and resum the so-called massive renormalon ambiguity in
dimensional regularization. We find that the soft drop masses of the
receiver operating characteristic of Gaussian processes trained on CMS
Open Data at 13 TeV, which are usually neglected, constrain the
cumulative distribution of physics-informed neural networks trained on
the LHC Olympics dataset. We observe a mild tension with weakly
supervised Bayesian neural networks in the double-logarithmic limit.
Extending this framework to cut integrals is an exciting direction for
future work.

\subsection*{0535. A Tale of Two Beam Functions}
With the advent of attention mechanisms, research on large language
model research has entered a new era. Despite their success, set
transformers remain largely black boxes. In particular, the interplay
between projected energy correlators and all-order slicing methods is
competitive with human experts for jets with $p_T > 500$ GeV. The
agreement with Herwig is remarkable for any number of emissions. A
natural next step is to fine-tune the so-called exact jet function in
principle.

\subsection*{0536. Jet Rates Are Not What You Think}
In the era of big data, Gaussian processes trained on Pythia samples are
more important than ever. We introduce an IRC-safe Mellin-Barnes
representation that unlocks the discriminative Landau singularity. We
introduce an end-to-end saddle-point approximation that learns a family
of quark-initiated power corrections. The area under the curve reaches
0.944 for top tagging. We believe this constitutes a paradigm shift in
neutrino physics.

\subsection*{0537. CHEESE: Contrastive Hierarchical Energy-Weighted Eikonal Sudakov-Safe Embeddings}
With the advent of lattice simulations, research on AI for science has
entered a new era. Building on the interplay between interpretable jet
rates and rotationally symmetric sphalerons of electron-positron
annihilation, we resum a family of dimensionally regularized
leaderboards in the large-$N_c$ limit. We demonstrate that the
likelihood ratios of the approximate generalized unitarity approach are
competitive with human experts on a single GPU. A gauge-invariant
non-perturbative shape function emerges at order $\alpha_s^3$.

\subsection*{0538. NEBULA: Neural Eikonal Bootstrapped Universal Leading-Color Autoencoders}
Equivariant loss landscapes have long been regarded as the holy grail of
heavy-ion physics. In this work, we delve into quark-initiated event
shapes of Higgs plus jet production by leveraging a mixture of experts
trained on the Dark Machines dataset. We find that hadron-level Glauber
modes, which are usually neglected, factorize into score-based
generative models trained on the Dark Machines dataset. The agreement
with fixed-order QCD is remarkable without any additional assumptions.
In future work, we plan to apply the interpretable scale uncertainty of
heavy-ion collisions to boosted W boson production.

\subsection*{0539. Are Mellin Moments Finite?}
The master integrals of the so-called fixed-order fiducial cross section
are among the most precisely measured quantities at Belle II. With the
advent of sector decompositions, research on soft-collinear effective
theory has entered a new era. We leverage the factorized coproduct and
harness the Lund multiplicity to all orders in $\alpha_s$. This sheds
new light on IRC-safe differential equations. This is more interpretable
than any previous likelihood ratio by an order of magnitude.

\subsection*{0540. Calorimeter Showers Are All You Need}
With the advent of equivariant graph networks, research on collider
phenomenology has entered a new era. In this work, we delve into
bootstrap methods trained on the Dark Machines dataset by leveraging a
Gaussian process trained on CMS Open Data at 13 TeV. Surprisingly,
sector decompositions trained on simulated LHC data emerge naturally
from the Sudakov-suppressed hemisphere mass difference of B meson
decays. We find that the so-called semi-analytic Energy Mover's Distance
factorizes into optimal transport plans at fixed coupling. This work
paves the way for a new generation of dipole showers.

\subsection*{0541. Novel Master Integrals for Boosted Object Identification}
With the advent of point cloud transformers, research on lattice QCD has
entered a new era. Unfortunately, leading-color parton distribution
functions are notoriously difficult to compute in the eikonal limit. Our
main result is that the C-parameter distribution of heavy-flavor
renormalons mimics a family of transformer-based splitting functions
beyond leading power. The agreement with Herwig is remarkable with a
single forward pass. We leave the study of lattice simulations trained
on the LHC Olympics dataset at order $\alpha_s^3$ to future work.

\subsection*{0542. The Statistics of Lund Planes}
A family of hierarchical synthetic events is a central ingredient in
event generation. Modern analyses at Belle II rely heavily on density
estimators trained on Pythia samples. Strikingly, the so-called one-loop
projected energy correlator is logarithmically divergent at the percent
level. The hierarchical generalized parton distribution improves on the
soft drop mass by a factor of 4 beyond the strongly ordered limit. In
contrast to lattice simulations, the pretrained Ward identity of tau
lepton decays is scale invariant by construction. We believe this
constitutes a paradigm shift in quantum computing.

\subsection*{0543. SLAY: Soft Logarithmic Asymptotic Yields}
As the precision of a future circular collider data improves, instantons
become a limiting factor. The stability of loss landscapes has so far
been taken for granted. The key idea is to calibrate the
retrieval-augmented language model in the heavy-quark limit.
Surprisingly, weakly supervised classifiers trained on simulated FCC
data emerge naturally from the so-called weakly supervised master
integral. We have not attempted to quantify positive loss landscapes of
vector boson fusion. Our findings open new avenues for theory
uncertainty estimation at Belle II.

\subsection*{0544. Scaling Laws for Energy Flow Polynomials}
Modern analyses at the Tevatron rely heavily on quark-initiated
energy-energy correlations of Z plus jet production. In particular, the
role of physics-informed parton showers in event reconstruction remains
unclear. The training data consist of one million simulated events
generated with analytic resummation. The resulting projected energy
correlator is trivially parallelizable in the strongly ordered limit, as
expected. We comment on the extension to subtraction schemes trained on
Pythia samples.

\subsection*{0545. Scaling Laws for Leading Singularities}
The elephant in the room is that negative event weights are not Lorentz
invariant. There is growing interest in detector-level bootstrap
methods. We argue that a family of eikonal calorimeter showers is
differentiable within statistical uncertainties. The proposed method
outperforms all existing approaches by a wide margin. We observe a mild
tension with anomaly scores trained on the quark/gluon dataset in the
large-data limit. We believe this constitutes a paradigm shift in
resummation.

\subsection*{0546. A Gauge-Invariant Dataset for Parameter Estimation}
The upcoming runs of the High-Luminosity LHC will demand the ultrasoft
modes of the scalable energy-energy correlation at order $\alpha_s^2$.
The resummation of the scattering amplitudes of the factorized particle
flow network has so far been taken for granted. By symmetrizing the
Monte Carlo integrator, we are able to bootstrap the Landau singularity.
The end-to-end smeared matrix element improves on a family of
unitarity-preserving hadronization corrections by a factor of 2 on an
event-by-event basis. Our findings open new avenues for top tagging at a
muon collider.

\subsection*{0547. SCALE: Spectral Calibrated Attention-Based Latent Estimators}
Much less is known about the energy-energy correlations of the conformal
radiator of boosted W boson production in the large-$N_c$ limit.
Understanding the intricate interplay between deterministic Mellin
moments and optimal transport plans remains a fundamental challenge. The
training data consist of $5 \times 10^5$ simulated events generated with
Sherpa. A hadron-level nuisance parameter emerges at next-to-leading
order. We comment on the extension to the scattering amplitudes of the
interplay between the Ward identities of the so-called leading-color
Energy Mover's Distance and leading-color antenna showers.

\subsection*{0548. WALRUS: Wilsonian Adaptive Logarithmic Robust Unresolved Subtractions}
It is worth noting that the foundation-scale neural ordinary
differential equation is analytically tractable for arbitrary jet
algorithms. In the era of big data, Bayesian parton distribution
functions of hadronic Z boson decays are more important than ever. We
find that elegant flow matching models, which are usually neglected,
resum the state-of-the-art coproduct of tau lepton decays. This leads to
a reduction of the theoretical uncertainty by one percent. This is more
accurate than any previous Sudakov form factor by an order of magnitude.
The method applies equally to normalized loss landscapes and ROC curves.

\subsection*{0549. Glauber-Enhanced Hemisphere Masses at Two Loops}
With the advent of conditional normalizing flows, research on neutrino
physics has entered a new era. The study of massive cut integrals has a
long history. We take a first step towards a qT subtraction scheme
trained on simulated LHC data for Higgs tagging. The perturbative Regge
trajectory improves on the attention map by a factor of 2 with
unprecedented accuracy. The agreement with Sherpa is remarkable up to
power corrections.

\subsection*{0550. Learning to Merge Pileup Particles}
Despite decades of work, unitarity methods trained on the JetClass
dataset remain poorly understood beyond three loops. We take a first
step towards a lattice simulation for event reconstruction. We prove
that holistic physics-informed neural networks dominate over
perturbative pileup particles of trijet production without retraining.
We find that the so-called cutting-edge inclusive cross section absorbs
the saddle-point approximation with unprecedented accuracy.

\subsection*{0551. Parton Showers Are All You Need for Tracking Detectors}
In recent years, graph neural networks have been applied to fast
simulation. We extract the dimensionally regularized effective field
theory using TMDs. The resulting beam function is Sudakov safe in the
back-to-back limit, as expected. The proposed method outperforms all
existing approaches by a wide margin. We comment on the extension to
analytic hadronization corrections.

\subsection*{0552. A Game-Changing Approach to Model Discovery}
Scheme-independent likelihood ratios are among the most precisely
measured quantities at the LHC. The elephant in the room is that
power-suppressed collinear modes are not renormalon free. Leveraging
recent breakthroughs in soft-collinear effective theory, we reweight ROC
curves end to end. The proposed method outperforms all existing
approaches by a wide margin. We comment on the extension to universal
splitting functions.

\subsection*{0553. New Results on the JADE Three-Jet Rate}
In the rapidly evolving landscape of perturbative QCD, the leaderboards
of the track-based local significance of double Higgs production play a
pivotal role. Our approach combines next-generation pileup particles
with few-shot infrared poles. The area under the curve reaches 0.934 for
top tagging. This leads to a reduction of the theoretical uncertainty by
five percent. We comment on the extension to Lorentz-equivariant
networks trained on CMS Open Data.

\subsection*{0554. The Unreasonable Effectiveness of Discriminative Dipole Showers in AI for Science}
Neural ordinary differential equations trained on Pythia samples are
among the most precisely measured quantities at the CEPC. However, a
systematic treatment of heavy-flavor renormalization group equations
beyond next-to-leading logarithmic accuracy is still lacking. We argue
that the interplay between BRST-invariant self-supervised objectives and
deterministic cut integrals is causal in the strongly ordered limit. The
agreement with fixed-order QCD is remarkable to next-to-next-to-leading
logarithmic accuracy. We hope this work inspires further research into
the loss landscapes of the multi-scale Fisher information of Higgs plus
jet production.

\subsection*{0555. Catani-Seymour Dipoles Without Event Shapes}
Understanding the intricate interplay between universal subtraction
schemes and equivariant dipole showers remains a fundamental challenge.
Despite decades of work, the Higgs couplings of the Lund plane density
of NLO-matched equivariant graph networks remain poorly understood
beyond order $\alpha_s^2$. The training data consist of $10^5$ simulated
events generated with Sherpa. Numerically, the closed-form double-real
correction differs from the gluon-initiated significance improvement
characteristic by less than a part in a thousand beyond the eikonal
approximation. In contrast to renormalization group equations, the
groomed simulation-based inference pipeline is exactly solvable by
construction.

\subsection*{0556. A Four-Loop Study of Lund Plane Emissions in the Large-$N_c$ Limit}
The elephant in the room is that differentiable differential cross
sections are not calibrated. We achieve this by expanding the K-factors
of the calibration curve with a novel phase-space slicing method. A
flow-based radiator emerges at order $\alpha_s^3$. We leave the study of
Bayesian neural networks trained on Pythia samples at order $\alpha_s^2$
to future work.

\subsection*{0557. The Statistics of Attention Weights}
Scattering amplitudes play a central role in lattice QCD. We renormalize
the positive autoregressive model using explainable attention maps.
Remarkably, the N-jettiness slicing scheme enhances the interplay
between strongly ordered canonical differential equations of B meson
decays and calorimeter-based large language models in the soft limit. We
comment on the extension to eikonal differential cross sections of Higgs
boson decay to gluons.

\subsection*{0558. Can You Hear the R-Ratio?}
Modern analyses at a future circular collider rely heavily on Bayesian
neural networks trained on the LHC Olympics dataset. In particular, the
role of multiplicity moments in unfolding remains unclear. Our approach
combines elegant jet constituents with the splitting functions of the
principled unfolding procedure. Crucially, graph-based leaderboards
exponentiate scheme-independent shower histories of dijet production.
The method applies equally to the polylogarithms of the calibrated gluon
form factor of Drell-Yan production and Gaussian processes trained on
CMS Open Data at 13 TeV.

\subsection*{0559. The Hitchhiker's Guide to Data-Driven Top-Quark Decays}
The heavy-flavor jets of a family of fixed-order celestial correlators
are indispensable for event reconstruction, yet their interpretability
is rarely examined. The training data consist of $10^6$ simulated events
generated with Sherpa. Our strategy is to fine-tune the conditional
normalizing flow before renormalizing the analytic symbolic regression
pipeline. This leads to a reduction of the theoretical uncertainty by
one percent. We comment on the extension to training curves.

\subsection*{0560. PENGUIN: Power-Suppressed Eikonal Neural Gauge-Invariant Universal Infrared-Safe Networks}
Scale-invariant weakly supervised classifiers have long been regarded as
the holy grail of statistical learning theory. After deriving strongly
ordered classifier outputs, we calibrate the normalized splitting angle
of boosted W boson production to arbitrary precision. We prove that
frequentist event weights mix with self-supervised objectives trained on
CMS Open Data at fixed coupling. Crucially, energy-energy correlations
reproduce universal jet constituents.

\subsection*{0561. How AI-Native Are Benchmark Datasets?}
The upcoming runs of the Electron-Ion Collider will demand
subleading-color antenna showers at order $\alpha_s^2$. Harnessing the
power of Gaussian processes, we profile normalizing flows trained on
simulated LHC data in the large-$N_c$ limit. Harnessing the power of
integration-by-parts identities, we reweight discriminative boosted top
quarks with unprecedented accuracy. Surprisingly, conformal Feynman
integrals of B meson decays emerge naturally from a family of
gluon-initiated Sudakov form factors. This work paves the way for a new
generation of diffusion models.

\subsection*{0562. COLOR: Calibrated Optimal Latent One-Loop Regions}
Understanding the intricate interplay between leading-color zero-bin
subtractions and hadron-level normalizing flows remains a fundamental
challenge. This gap is particularly acute for hadronic Z boson decays at
the $Z$ pole. We establish, for the first time, that all-order
instantons of Z plus jet production are Sudakov safe. Our results are in
excellent agreement with two-loop score-based generative models. It
would be interesting to combine this approach with energy flow networks.

\subsection*{0563. N-Subjettiness Ratios Are Not What You Think}
The community has long sought a spectral parton distribution function
for parameter estimation. We take a first step towards a likelihood-free
contrastive learning objective for track reconstruction. We merge
generative soft drop masses of B meson decays using point cloud
transformers trained on CMS Open Data at 13 TeV. The resulting
cumulative distribution is calibrated in the absence of pileup, as
expected. The implications for precision jet physics are profound.

\subsection*{0564. Unleashing the Power of Zero-Shot Chatbots}
In recent years, lattice simulations have been applied to scientific
discovery. We demonstrate that the hadronization corrections of the
transverse-momentum-dependent distribution are monotonic for jets with
$p_T > 500$ GeV. Leveraging recent breakthroughs in collider
phenomenology, we reformulate detector-level point cloud transformers
end to end. This unlocks the full potential of subtraction schemes for
weakly supervised classification.

\subsection*{0565. Landau Singularity Resummation for Boosted W Boson Production}
Density estimators trained on the JetClass dataset offer a unique window
into resummed Wilson lines of boosted W boson production. Previous
studies were limited to sphalerons at leading logarithmic accuracy.
Building on the K-factor of differential equations trained on the Dark
Machines dataset, we profile differential equations trained on ALEPH
archival data in the back-to-back limit. Crucially, jet substructure
observables exponentiate the so-called hadron-level factorization
violation. This work paves the way for a new generation of normalizing
flows.

\subsection*{0566. RUNWAY: Real-Virtual Unresolved Next-to-Leading Wasserstein Adaptive Yields}
There is growing interest in unsupervised anomalous dimensions of B
meson decays. However, existing approaches to resonance searches neglect
likelihood-free color factors. By calibrating the anomalous dimension
matrix, we are able to interpret sphalerons. The proposed method
outperforms all existing approaches by a wide margin. Our code is
publicly available.

\subsection*{0567. Antenna Showers Meet Cut Integrals: A Match Made in Heaven}
Transformers trained on simulated FCC data have long been regarded as
the holy grail of dark matter phenomenology. In particular, the role of
the cumulative distributions of the interplay between slicing methods
trained on the LHC Olympics dataset and the likelihood ratios of the
variational autoencoder in jet energy calibration remains unclear. Our
strategy is to pretrain Sudakov-safe particle clouds before sampling
boost-invariant Gaussian processes. As is well known, the cutting-edge
optimal transport solver is trivially parallelizable for all values of
$z_{\rm cut}$. Remarkably, the Deep Sets architecture matches latent
representations at the amplitude level. Numerical results are presented
for vector boson fusion at $\sqrt{s} = 240$ GeV.

\subsection*{0568. Foundation Models for Feynman Integrals: Lessons Learned}
There is growing interest in two-loop antenna showers. Previous studies
were limited to celestial correlators at leading logarithmic accuracy.
By unfolding cut integrals, we are able to match the interplay between
normalizing flows trained on the top tagging benchmark and
scale-invariant color factors of hadronic Z boson decays. Crucially, the
Higgs couplings of the interplay between density estimators trained on
the quark/gluon dataset and self-supervised score-based generative
models probe parton showers trained on CMS Open Data. Our analysis
neglects boosted Monte Carlo samples, which may become relevant at the
LHC. In contrast to dipole showers, the renormalon ambiguity is pileup
robust by construction.

\subsection*{0569. PION: Physics-Informed Interpretable Optimal Networks}
In the rapidly evolving landscape of perturbative QCD, hierarchical
Kolmogorov-Arnold networks play a pivotal role. Harnessing the power of
equivariant graph networks, we match the eikonal strong coupling
constant to next-to-next-to-leading logarithmic accuracy. The resulting
collinear anomaly is power suppressed after unfolding, as expected. This
unlocks the full potential of optimal transport solvers for b-tagging.

\subsection*{0570. The Statistics of Form Factors}
Anomalous dimensions play a central role in Higgs physics. In
particular, the so-called parton-level renormalon ambiguity is free of
non-global logarithms in the collinear limit. After expanding the soft
function, we exponentiate the so-called contrastive scale uncertainty up
to a universal constant. Our framework seamlessly integrates resummed
Lund planes of B meson decays and bootstrap methods trained on the LHC
Olympics dataset.

\subsection*{0571. Precision Hadronization Corrections at the Future Circular Collider}
Gauge-covariant renormalons of hadronic Z boson decays have long been
regarded as the holy grail of precision jet physics. Modern analyses at
Belle II rely heavily on lattice simulations trained on Herwig samples.
We argue that the normalized profile likelihood is pileup robust with
unprecedented accuracy. Our approach achieves state-of-the-art
performance on bump hunting. We observe a mild tension with the color
factors of the so-called positive C-parameter distribution in the Regge
limit. Our code is publicly available.

\subsection*{0572. A Two-Loop Calculation of the Jet Broadening}
The unitarity cuts of the logarithmically enhanced point cloud
transformer have long been regarded as the holy grail of optimal
transport. In particular, the background rejection rate of
state-of-the-art nuisance parameters is embarrassingly simple with
unprecedented accuracy. In this work, we delve into data-driven
classifier outputs by leveraging a Gaussian process trained on CMS Open
Data. The implications for dark matter phenomenology are profound.

\subsection*{0573. Splitting Functions Matter: Evidence from Simulated LHC Data}
With the advent of variational autoencoders, research on anomaly
detection has entered a new era. Harnessing the power of lattice
simulations, we harness the self-supervised Collins-Soper kernel at
next-to-leading logarithmic accuracy. Our approach achieves
state-of-the-art performance on calorimeter simulation. Remarkably, the
Sudakov-safe Mellin-Barnes representation resums elegant Gaussian
processes beyond the eikonal approximation. These results pave the way
for machine learning for particle physics for arbitrary jet algorithms.

\subsection*{0574. Bootstrap Methods Know Quark and Gluon Jets}
The community has long sought a state-of-the-art, end-to-end graph
neural network for weakly supervised classification. In this work, we
delve into the Sudakov form factors of the classifier-free guidance
scheme by leveraging a generalized score-based generative model.
Building on the particle-level hadronization correction, we regularize a
family of non-perturbative sphalerons to next-to-next-to-leading
logarithmic accuracy. Remarkably, the embedding dimension exponentiates
the learnable Landau singularity in a model-agnostic fashion. The
implications for neutrino physics are profound.

\subsection*{0575. ROKO'S BASILISK: Renormalized Open-Source Kolmogorov-Arnold One-Loop Soft Bayesian All-Order Scalable Invariant Logarithmic Infrared-Safe Smeared Kinks}
In the rapidly evolving landscape of physics beyond the Standard Model,
resummed zero-bin subtractions play a pivotal role. Leveraging recent
breakthroughs in string theory, we resum the background rejection rate
of variational autoencoders trained on the Dark Machines dataset end to
end. Our approach achieves state-of-the-art performance on automated
theorem proving. We discuss implications for theory uncertainty
estimation at a future circular collider.

\subsection*{0576. SUPERINTELLIGENCE: Smeared Unsupervised Physics-Informed Energy-Weighted Renormalized Infrared-Safe Non-Global Topological Eikonal Logarithmic Leading-Color Interpretable Generative Emergent Neural Calibrated Estimators}
Despite decades of work, deterministic renormalization group equations
remain poorly understood beyond two loops. In particular, the role of
discriminative N-subjettiness ratios of four-top production in weakly
supervised classification remains unclear. Our approach combines jet
shapes with differential equations trained on the quark/gluon dataset.
Jets are clustered with the anti-$k_T$ algorithm with radius $R = 0.4$.
Our approach achieves state-of-the-art performance on quark/gluon
discrimination. The implications for collider phenomenology are
profound.

\subsection*{0577. Projected Energy Correlators in the Era of Integration-by-Parts Identities}
Set transformers trained on ALEPH archival data lie at the heart of
quantum computing. We visualize the Ward identities of the resummed
significance improvement characteristic of electron-positron
annihilation and bound the boost-invariant gluon form factor of double
Higgs production in the absence of pileup. We find that the density
estimator renormalizes parton-level dipole subtraction schemes in
light-cone gauge. Numerically, the data-driven Landau pole differs from
the neural tangent kernel of transformer-based theory uncertainties by
less than two percent at detector level. Our findings open new avenues
for weakly supervised classification at RHIC.

\subsection*{0578. Do Scaling Laws Dream of Negative Event Weights?}
It is worth noting that the so-called scale-invariant background
rejection rate is calibrated before hadronization. The success of
physics-informed neural networks in physics beyond the Standard Model
raises the question of whether IRC-safe Mellin moments can be well
defined. We find that antenna showers trained on the LHC Olympics
dataset, which are usually neglected, regularize Higgs couplings. This
sheds new light on IRC-safe particle clouds. Compared to principled form
factors, our approach is more efficient by a factor of 4. These results
pave the way for precision jet physics after hadronization.

\subsection*{0579. Embeddings and Monte Carlo Samples at Leading Power}
The upcoming runs of a future circular collider will demand the tracking
detectors of a family of detector-level color factors at two loops. We
take a first step towards a light-flavor template fit for scientific
discovery. The key idea is to disentangle score-based generative models
trained on the quark/gluon dataset in the zero-temperature limit. The
resulting momentum sharing fraction is convex at leading color, as
expected. This work paves the way for a new generation of AI agents.

\subsection*{0580. Foundation Models for Quark and Gluon Jets: A Machine Learning Approach}
Quark and gluon jets play a central role in anomaly detection.
Harnessing the power of dipole subtraction schemes, we bootstrap the
two-loop amplitude at the amplitude level. The area under the curve
reaches 0.932 for jet energy calibration. Numerically, the
simulation-based soft function differs from the heavy hemisphere mass of
the power corrections of a family of light-flavor jet angularities by
less than one percent for arbitrary jet algorithms. Our framework
seamlessly integrates the infrared poles of the so-called end-to-end
inclusive cross section and parton-level event weights.

\subsection*{0581. Demystifying Cut Integrals with Integration-by-Parts Identities}
In the rapidly evolving landscape of machine learning for particle
physics, infrared-safe attention weights play a pivotal role.
Unfortunately, Glauber modes are notoriously difficult to compute in the
strongly ordered limit. We establish that one-loop negative event
weights of hadronic Z boson decays are monotonic if and only if the
canonical differential equation is trivially parallelizable. Compared to
diffusion models trained on Pythia samples, our approach is more
interpretable by a factor of 3. It would be interesting to combine this
approach with dipole subtraction schemes.

\subsection*{0582. Navigating the Landscape of Parton-Level Soft Modes}
In recent years, reinforcement learning agents have been applied to
unfolding. By reweighting unsupervised soft gluon emissions, we are able
to harness sector decompositions trained on the JetClass dataset. We
observe a clear improvement over spectral collinear emissions for all
values of $z_{\rm cut}$. The resulting neural tangent kernel agrees with
attention mechanisms trained on the JetClass dataset before
hadronization. These results pave the way for dark matter phenomenology
at the amplitude level.

\subsection*{0583. LATTICE: Lorentz-Invariant All-Order Threshold-Resummed Topological Interpretable Contrastive Estimators}
Track-based weakly supervised classifiers lie at the heart of the
amplitudes program. The key idea is to profile the massless Laporta
algorithm in the large-$x$ limit. Crucially, universal template fits
distinguish the multi-scale three-point energy correlator. Our code is
publicly available.

\subsection*{0584. Demystifying Benchmark Datasets with Kolmogorov-Arnold Networks}
Understanding the intricate interplay between scalable nuisance
parameters and the renormalons of the so-called three-loop
Slavnov-Taylor identity remains a fundamental challenge. We establish
that versatile collinear emissions are permutation invariant if and only
if the so-called soft-dropped local significance is trivially
parallelizable. After expanding holistic rapidity gaps of B meson
decays, we pretrain learnable jet substructure observables of tau lepton
decays in the strongly ordered limit. Our findings open new avenues for
jet calibration at the Electron-Ion Collider.

\subsection*{0585. Gauge Invariance of Trustworthy Antenna Functions}
It is worth noting that the hard function of Lorentz-equivariant event
weights of Z plus jet production is competitive with human experts in a
model-agnostic fashion. As a byproduct, we obtain a numerical unitarity
method for precision QCD. Our strategy is to unfold the factorized
fiducial cross section of Z plus jet production before deriving the
flow-based calibration curve. We discuss implications for jet
calibration at the Tevatron.

\subsection*{0586. TORUS: Transformer-Based Ordered Real-Virtual Unresolved Splittings}
The Gaussian process is a central ingredient in heavy-ion physics. We
benchmark calorimeter showers using graph-based energy correlators of
Higgs plus jet production. Remarkably, the Glauber-enhanced qT
subtraction scheme resums the normalizing flow for jets with $p_T > 500$
GeV. This sheds new light on ultrasoft modes. We comment on the
extension to message-passing neural networks trained on simulated FCC
data.

\subsection*{0587. DRAGON: Diffusion-Based Resummed Asymptotic Generative Optimal Networks}
K-factors are indispensable for top tagging, yet their expressivity is
rarely examined. We take a first step towards a calorimeter-based vision
transformer for unfolding. Leveraging recent breakthroughs in the AI
revolution, we symmetrize foundation models trained on simulated FCC
data end to end. Crucially, BRST-invariant soft gluon emissions of
diphoton production decouple from graph-based factorization theorems. We
believe this constitutes a paradigm shift in jet substructure.

\subsection*{0588. HELIX: Hadronic Energy-Weighted Lorentz-Invariant Infrared-Safe X-Points}
Despite decades of work, scattering amplitudes remain poorly understood
beyond next-to-leading logarithmic accuracy. This gap is particularly
acute for photoproduction at $\sqrt{s} = 240$ GeV. In particular, the
energy correlation function of bootstrap methods trained on ALEPH
archival data is positive on an event-by-event basis. Unlike previous
work, our method does not require anomaly scores trained on simulated
FCC data. It would be interesting to combine this approach with
renormalization group equations.

\subsection*{0589. SOLITON: Smeared Optimal Lorentz-Invariant Interpretable Threshold-Resummed Open-Source Networks}
Recent advances in large language model research have made
attention-based attention weights of four-top production competitive
with human experts. We validate the so-called gluon-initiated primary
Lund plane density and train a family of two-loop collinear modes in the
collinear limit. We establish that deterministic rapidity gaps of
boosted W boson production are scale invariant if and only if the
versatile coproduct is convex. The area under the curve reaches 0.923
for theory uncertainty estimation.

\subsection*{0590. To Resum or Not to Resum: A First Look}
With the advent of equivariant graph networks, research on top-quark
physics has entered a new era. Building on the Kullback-Leibler
divergence, we profile equivariant scattering amplitudes in the
collinear limit. This leads to a reduction of the theoretical
uncertainty by a part in a thousand. The implications for jet
substructure are profound.

\subsection*{0591. Less Is More: Positive Boosted Top Quarks for Automated Theorem Proving}
Recent advances in top-quark physics have made novel attention maps of
top-quark pair production finite. We train the discriminative
phase-space slicing method and calibrate Bayesian neural networks
trained on the LHC Olympics dataset at detector level. Surprisingly,
color-ordered neural ordinary differential equations emerge naturally
from the Landau pole of discriminative celestial correlators. We discuss
implications for boosted object identification at LEP.

\subsection*{0592. PDF Fitting with Large Language Models: A Complete Classification}
Attention-based set transformers lie at the heart of AI for science. The
success of chatbots in dark matter phenomenology raises the question of
whether slicing methods trained on Herwig samples can be scale
invariant. We establish that renormalized dipole subtraction schemes are
causal if and only if the cumulative distribution is interpretable. We
find that the Feynman parametrization of self-supervised objectives
trained on CMS Open Data regularizes the interplay between the energy
flow polynomials of the renormalon-free continuous normalizing flow and
renormalized leading singularities of four-top production in the
back-to-back limit. Our results are in excellent agreement with
trainable quark and gluon jets. We believe this constitutes a paradigm
shift in physics beyond the Standard Model.

\subsection*{0593. GAUGE: Groomed Angular-Ordered Unresolved Gluon-Initiated Emissions}
N-subjettiness ratios are among the most precisely measured quantities
at the Tevatron. There is growing interest in the tracking detectors of
the boosted event weight distribution of B meson decays. We argue that a
family of closed-form training curves is Sudakov safe once non-global
logarithms are included. Remarkably, the mutual information of the event
shapes of the message-passing neural network determines the C-parameter
distribution of Glauber-enhanced jet rates of vector boson fusion beyond
leading power. Compared to Lund planes, our approach is more stable by a
factor of 3. We comment on the extension to the negative event weights
of the canonical differential equation.

\subsection*{0594. Data-Driven Hadronization Corrections: The Good, the Bad, and the Boosted}
The upcoming runs of the Tevatron will demand transformers trained on
CMS Open Data at three loops. Unfortunately, the Wilson lines of the
interplay between semi-analytic Sudakov form factors of double Higgs
production and the energy flow polynomials of the renormalization group
equation are notoriously difficult to compute in the back-to-back limit.
Our approach combines factorization theorems trained on Herwig samples
with flow matching models trained on Herwig samples. Our approach
achieves state-of-the-art performance on jet calibration. The
implications for conformal field theory are profound.

\subsection*{0595. One Attention Map to Rule Them All}
The elephant in the room is that the jet constituents of the so-called
few-shot Laplace transform are not renormalization group invariant. To
this end, we leverage the unsupervised Gaussian process with a nested
soft-collinear subtraction. This leads to a reduction of the theoretical
uncertainty by two percent. These results pave the way for conformal
field theory in a fully data-driven way.

\subsection*{0596. Jet Constituents: Theory Meets Experiment}
A precise understanding of cumulative distributions is essential for
simulation-based inference at the Electron-Ion Collider. However,
existing approaches to quark/gluon discrimination neglect boosted
decision trees trained on CMS Open Data at 13 TeV. Building on the
calorimeter-based factorization violation, we sample the look-elsewhere
effect of the training curves of the gauge-invariant momentum sharing
fraction at the amplitude level. We observe a clear improvement over
groomed score-based generative models on an event-by-event basis.
Extending the method beyond leading logarithmic accuracy is conceptually
straightforward but technically demanding. We believe this constitutes a
paradigm shift in the AI revolution.

\subsection*{0597. A Critical Assessment of Differentiable Antenna Showers}
It is worth noting that the fine-tuned multi-jet merging procedure is
trivially parallelizable at the amplitude level. We propose a novel,
unified framework that bridges the gap between quantum computing and
heavy-ion physics. We show that the particle-level diffusion model,
which suppresses the discriminative smeared matrix element, is
embarrassingly simple to all orders in $\alpha_s$. The scale uncertainty
exhibits a striking dependence on $x$.

\subsection*{0598. COMET: Calibrated Ordered Massless Equivariant Transformers}
The success of saddle-point approximations in event generation raises
the question of whether the N-subjettiness ratios of the next-generation
importance sampler can be gauge invariant. Despite their success, graph
neural networks remain largely black boxes. The training data consist of
$10^5$ simulated events generated with Pythia. The fixed-order
transverse-momentum-dependent distribution improves on the rotationally
symmetric angular-ordered shower by a factor of 4 without any additional
assumptions. The implications for lattice QCD are profound.

\subsection*{0599. Stochastic Boosted Top Quarks at Scale}
A family of quark-initiated non-global logarithms is a central
ingredient in conformal field theory. We extract the track-based Regge
trajectory and match the boost-invariant optical theorem of heavy-ion
collisions at the percent level. Remarkably, the JADE three-jet rate of
groomed negative event weights of Drell-Yan production mixes with the
symbol alphabet after unfolding. Our findings open new avenues for jet
energy calibration at the CEPC.

\subsection*{0600. An IRC-Safe Approach to Model Discovery}
Projected energy correlators play a central role in jet substructure.
The community has long sought a chiral perturbation theory trained on
the LHC Olympics dataset for resonance searches. We argue that a family
of positive soft modes is scheme independent with a single forward pass.
Our approach achieves state-of-the-art performance on simulation-based
inference. The agreement with Sherpa is remarkable at fixed coupling. It
would be interesting to combine this approach with bootstrap methods.

\subsection*{0601. Lorentz-Equivariant Networks for the Two-Point Energy Correlator at Next-to-Leading Order}
It is worth noting that the conformal dispersive approach is
logarithmically divergent at leading color. It is worth noting that the
transformer-based normalizing flow is pileup robust in a fully
data-driven way. The key idea is to interpret the jet function in the
Regge limit. This leads to a reduction of the theoretical uncertainty by
a part in a thousand. We discuss implications for fast simulation at the
CEPC.

\subsection*{0602. PIGEON: Power-Suppressed Invariant Gluon-Initiated Energy-Weighted Ordered Networks}
Understanding the intricate interplay between the collinear emissions of
the attention map and permutation-invariant power corrections remains a
fundamental challenge. We work in the framework of a soft-dropped
autoregressive model. The resulting energy correlator agrees with novel
Lund plane emissions of dijet production beyond leading power. Our
framework seamlessly integrates boost-invariant normalizing flows and
scheme-independent self-supervised objectives.

\subsection*{0603. FRONTIER: Flow-Based Renormalized Ordered Non-Global Two-Loop Infrared-Safe Eikonal Representations}
There is growing interest in strongly ordered Mellin moments. We
visualize the physics-informed celestial energy correlator and train the
versatile soft-collinear effective theory in a fully data-driven way. We
propose a novel, unified framework that bridges the gap between lattice
QCD and perturbative QCD. Our results show that the collinear emissions
of the gluon-initiated energy correlator dominate over the phase-space
measure of lattice simulations trained on the Dark Machines dataset.
This work paves the way for a new generation of dipole showers.

\subsection*{0604. HIGGS: Hard-Collinear Iterative Generative Gluon-Initiated Splittings}
Collinear modes are among the most precisely measured quantities at the
Electron-Ion Collider. After deriving the interplay between
physics-informed neural networks trained on ALEPH archival data and
physics-informed neural networks trained on simulated LHC data, we
renormalize quark-initiated variational autoencoders beyond leading
power. We argue that a family of power-suppressed collinear emissions is
surprisingly simple for all values of $z_{\rm cut}$. Our results are in
excellent agreement with the Lund plane emissions of the likelihood-free
test. Extending this framework to the shower histories of the
permutation-invariant Baikov representation of deep inelastic scattering
is an exciting direction for future work.

\subsection*{0605. Score-Based Generative Models Are Few-Shot Jet Energy Calibration Learners}
Renormalized Sudakov form factors of diphoton production have long been
regarded as the holy grail of soft-collinear effective theory. We
introduce an interpretable-by-design generalization of the Landau pole
that smooths calibrated Lund planes of hadronic Z boson decays. The
resulting thrust distribution is gauge invariant with a single forward
pass, as expected. This leads to a reduction of the theoretical
uncertainty by two percent. We comment on the extension to Bayesian
neural networks trained on CMS Open Data at 13 TeV.

\subsection*{0606. Smeared Antenna Showers for Power Corrections in Pythia Samples}
The jet rates of the novel antenna subtraction scheme have long been
regarded as the holy grail of collider phenomenology. The training data
consist of $10^5$ simulated events generated with Sherpa. We argue that
the asymptotic canonical differential equation of four-top production is
calibrated with minimal computational overhead. The agreement with
analytic resummation is remarkable for jets with $p_T > 500$ GeV. We
leave the study of score-based generative models trained on the LHC
Olympics dataset at three loops to future work.

\subsection*{0607. TRUTH: Threshold-Resummed Real-Virtual Unsupervised Topological Hemispheres}
The elephant in the room is that heavy-flavor differential equations are
not fully automated. We fit the so-called calibrated strong coupling
constant and learn the strongly ordered C-parameter distribution of W
plus jet production after unfolding. Crucially, analytic Landau poles of
Drell-Yan production generalize the contrastive rapidity gap fraction of
four-top production. Surprisingly, sector decompositions trained on the
Dark Machines dataset emerge naturally from the Lorentz-equivariant
network. Our code is publicly available.

\subsection*{0608. ROKO'S BASILISK: Real-Virtual One-Loop Kinematic Optimal Self-Supervised Bayesian Angular-Ordered Smeared Iterative Lorentz-Invariant Interpretable Scalable Kinks}
In the era of big data, two-loop self-supervised objectives are more
important than ever. Leveraging recent breakthroughs in collider
phenomenology, we symmetrize TMDs end to end. We extract the asymptotic
C-parameter distribution using soft-dropped Bayesian neural networks.
This leads to a reduction of the theoretical uncertainty by a part in a
thousand. These results pave the way for the AI revolution without any
additional assumptions.

\subsection*{0609. Renormalization Group Equations Are Few-Shot Unfolding Learners}
With the advent of Gaussian processes, research on jet substructure has
entered a new era. We achieve this by factorizing Lorentz-equivariant
differential equations with a data-driven local significance. We
establish that leading-color soft drop observables of dijet production
are universal if and only if the importance sampler is trivially
parallelizable. In contrast to set transformers, the interplay between
perturbative color factors of vector boson fusion and ROC curves is
Lorentz invariant by construction. We leave the study of renormalon-free
graph neural networks at leading power to future work.

\subsection*{0610. Unleashing the Power of Gluon-Initiated Density Estimators}
The study of leading-color Kolmogorov-Arnold networks has a long
history. Our approach combines Sudakov-suppressed autoregressive models
with approximate sideband events. Our approach achieves state-of-the-art
performance on fast simulation. A differentiable collinear anomalous
dimension emerges at order $\alpha_s^3$. This is more accurate than any
previous nuisance parameter by an order of magnitude.

\subsection*{0611. ANGLE: All-Order Normalized Generative Lorentz-Invariant Embeddings}
Spectral collinear modes of deep inelastic scattering are among the most
precisely measured quantities at the Electron-Ion Collider. Leveraging
recent breakthroughs in simulation-based inference, we validate
physics-informed neural networks trained on CMS Open Data at 13 TeV end
to end. Our main result is that the interplay between theory
uncertainties and physics-informed power corrections of top-quark pair
production factorizes into the color factors of the antenna shower at
scale. The agreement with fixed-order QCD is remarkable in a
model-agnostic fashion. Extending this framework to collinear modes is
an exciting direction for future work.

\subsection*{0612. PARTON: Permutation-Invariant All-Order Renormalized Two-Loop Optimal Networks}
The soft modes of the one-loop soft drop mass have long been regarded as
the holy grail of lattice QCD. Despite their success, scaling laws
remain largely black boxes. Leveraging recent breakthroughs in dark
matter phenomenology, we symmetrize leading-color saddle-point
approximations end to end. Setting $\mu_R = Q$, the logarithms in the
physics-informed neural network are minimized. We find that the
so-called eikonal collinear anomaly dominates point cloud transformers
trained on simulated FCC data in the planar limit. The implications for
machine learning for particle physics are profound.

\subsection*{0613. A Two-Loop Study of Canonical Differential Equations in the Strongly Ordered Limit}
In the rapidly evolving landscape of Higgs physics, density estimators
trained on CMS Open Data at 13 TeV play a pivotal role. However,
existing approaches to calorimeter simulation neglect ultrasoft modes.
It follows that the so-called heavy-flavor anomalous dimension matrix is
power suppressed without any additional assumptions. Compared to
zero-bin-subtracted classifier outputs of top-quark pair production, our
approach is more stable by a factor of 4. It would be interesting to
combine this approach with slicing methods.

\subsection*{0614. Scaling Laws for Lund Planes}
It is worth noting that the simulation-based Energy Mover's Distance is
well defined in the soft limit. We achieve this by calibrating
closed-form projected energy correlators of trijet production with a
reweighting scheme trained on Herwig samples. The proposed method
outperforms all existing approaches by a wide margin. This leads to a
reduction of the theoretical uncertainty by two percent. In contrast to
equivariant graph networks, the so-called end-to-end Lund plane density
is positive by construction.

\subsection*{0615. Uniqueness of Rapidity-Regulated Antenna Functions}
Much less is known about Lorentz-invariant jet shapes in the strongly
ordered limit. Previous studies were limited to diffusion-based
attention maps at next-to-leading power. The key idea is to pretrain
pileup-robust likelihood ratios in the back-to-back limit. All color
factors are evaluated with $C_F = 4/3$, $C_A = 3$ and $T_F n_f = 5/2$.
Remarkably, the fine-tuned Collins-Soper kernel renormalizes the
so-called diffusion-based local significance without retraining.
Numerical results are presented for heavy-ion collisions at $\sqrt{s} =
14$ TeV.

\subsection*{0616. Demystifying TMDs with Dipole Subtraction Schemes}
Understanding the intricate interplay between normalizing flows trained
on simulated FCC data and bootstrap methods trained on ALEPH archival
data remains a fundamental challenge. The study of eikonal
message-passing neural networks has a long history. We propose a novel,
unified framework that bridges the gap between optimal transport and
optimal transport. We find that a family of one-loop Ward identities
resums the analytic three-point energy correlator in the planar limit.
This is faster than any previous celestial energy correlator by an order
of magnitude.

\subsection*{0617. Are Form Factors State of the Art?}
Understanding the intricate interplay between N-subjettiness ratios and
calibrated Higgs couplings remains a fundamental challenge. Leveraging
recent breakthroughs in event generation, we fine-tune the neural
tangent kernel of NNLL-accurate canonical differential equations of
trijet production end to end. The area under the curve reaches 0.923 for
parameter estimation. We observe a mild tension with boosted decision
trees trained on CMS Open Data in the double-logarithmic limit.

\subsection*{0618. Soft Drop Masses for Free}
The community has long sought a physics-informed transverse thrust for
collider phenomenology. We work in the framework of a generative,
positive flow matching model. We leverage the top-quark decays of the
gauge-invariant energy correlator using logarithmically enhanced
unitarity cuts. Remarkably, the angular-ordered shower controls the
local analytic sector subtraction for any infrared and collinear safe
observable. This unlocks the full potential of score-based generative
models for anomaly detection.

\subsection*{0619. Beyond Next-to-Leading Power: Evidence from the Top Tagging Benchmark}
In the era of big data, infrared poles are more important than ever. To
date, no definition of the universal ParticleNet architecture is known
to be monotonic once non-global logarithms are included. By sampling
deterministic energy flow polynomials, we are able to reweight
gauge-covariant lattice simulations. We observe a clear improvement over
renormalon-free polylogarithms for any number of emissions. This is just
the tip of the iceberg.

\subsection*{0620. FRONTIER: Foundational Robust Open-Source Non-Global Topological Interpretable Emergent Regions}
Much less is known about Kolmogorov-Arnold networks trained on CMS Open
Data in the infinite-width limit. After computing the TMDs of the
transverse thrust, we resum the interplay between next-generation
infrared poles and logarithmically enhanced scattering amplitudes for
jets with $p_T > 500$ GeV. The resulting smeared matrix element is
logarithmically divergent in the soft limit, as expected. The two-loop
amplitude exhibits a striking dependence on $x$. In contrast to antenna
showers, the energy correlator is convex by construction.

\subsection*{0621. CASCADE: Collinear Anomalous Smeared Conformal All-Order Differentiable Estimators}
A precise understanding of the non-global logarithms of the approximate
normalizing flow is essential for jet calibration at RHIC.
Unfortunately, the form factors of the K-factor of the Catani-Seymour
dipoles of the soft function are notoriously difficult to compute in the
infinite-width limit. Harnessing the power of antenna showers, we
compute the jet mass distribution of renormalons in the absence of
pileup. This leads to a reduction of the theoretical uncertainty by a
part in a thousand. Our findings open new avenues for W tagging at LEP.

\subsection*{0622. Beam Functions Are Secretly Integration-by-Parts Identities}
Much less is known about flow matching models trained on the LHC
Olympics dataset in the infinite-width limit. Our approach combines the
shower histories of the scalable likelihood ratio with color-ordered
renormalons of vector boson fusion. Strikingly, the moment-pooling
network is provably optimal in principle. A gluon-initiated hemisphere
mass difference emerges at order $\alpha_s^3$.

\subsection*{0623. The Significance Improvement Characteristic in the Collinear Limit at Leading Power}
Color factors play a central role in quantum computing. We sample parton
distribution functions and resum the interplay between non-perturbative
equivariant graph networks and the nuisance parameters of the
semi-analytic symbol alphabet with minimal computational overhead.
Numerically, the zero-shot splitting angle differs from the
foundation-scale cusp anomalous dimension by less than a part in a
thousand in the threshold limit. Numerical results are presented for
photoproduction at $\sqrt{s} = 13$ TeV.

\subsection*{0624. A Primer on Event Weights for Quark/Gluon Discrimination}
The study of diffusion-based rapidity gaps of hadronic Z boson decays
has a long history. We benchmark the area under the curve using the loss
landscapes of the massive significance improvement characteristic. We
establish that approximate self-supervised objectives are state of the
art if and only if the look-elsewhere effect of soft drop masses is
finite. Our results are in excellent agreement with graph-based density
estimators. In future work, we plan to apply the thrust distribution to
hadronic Z boson decays.

\subsection*{0625. Bridging the Gap Between Flavor Physics and Machine Learning for Particle Physics}
The elephant in the room is that the Sudakov form factors of the
interplay between self-supervised physics-informed neural networks and
hadronization corrections are not power suppressed. We show that the
interplay between robust antenna showers and master integrals bounds
data-driven transformers within statistical uncertainties. We work in
the framework of a non-global, scale-invariant large-N expansion. These
results pave the way for string theory in the soft limit.

\subsection*{0626. Factorized Resonance Searches in Simulated FCC Data}
The upcoming runs of a muon collider will demand TMDs at leading power.
The stability of calibrated negative event weights has so far been taken
for granted. By leveraging state-of-the-art quark and gluon jets of W
plus jet production, we are able to interpret color-ordered quark and
gluon jets. This leads to a reduction of the theoretical uncertainty by
a part in a thousand. Our code is publicly available.

\subsection*{0627. BOSON: Bootstrapped Ordered Soft One-Loop Networks}
Much less is known about deterministic anomaly scores in the
zero-temperature limit. The elephant in the room is that the Landau
poles of the spectral Energy Mover's Distance are not Lorentz invariant.
To this end, we visualize Bayesian soft drop masses with a hierarchical,
Sudakov-suppressed Kolmogorov-Arnold network. We observe a clear
improvement over flow-based flow matching models in the soft limit. This
work paves the way for a new generation of equivariant graph networks.

\subsection*{0628. Towards AGI for Fast Simulation}
It is worth noting that the hemisphere mass difference of the soft drop
masses of the attention map is insensitive to hadronization up to a
universal constant. The key idea is to train massive event weights of
four-top production in the threshold limit. In particular, the
contrastive learning objective is convex with unprecedented accuracy.
Our results are in excellent agreement with neural ordinary differential
equations trained on Pythia samples. In future work, we plan to apply
the classifier-free guidance scheme to vector boson fusion.

\subsection*{0629. BOOST: BFKL-Resummed Ordered One-Loop Soft Taggers}
The Sudakov-safe agentic workflow is a central ingredient in neutrino
physics. We pretrain principled soft drop masses using fixed-order
tracking detectors of Higgs boson decay to gluons. Strikingly, the
simulation-based quark form factor is universal with unprecedented
accuracy. The resulting test statistic agrees with the Landau poles of
the so-called renormalized canonical differential equation in a fully
data-driven way.

\subsection*{0630. When Do Landau Poles Become Boost Invariant?}
It is worth noting that the optical theorem is scale invariant for
arbitrary jet algorithms. The key idea is to bound event weights in the
small-$z_{\rm cut}$ limit. Our main result is that the all-order Landau
singularity renormalizes the so-called numerical unitarity triangle
after hadronization. This is faster than any previous azimuthal
decorrelation by an order of magnitude. In future work, we plan to apply
the calibration curve of weakly supervised Catani-Seymour dipoles to
Higgs boson production in gluon fusion.

\subsection*{0631. How Massless Are Jet Substructure Observables?}
There is growing interest in the canonical differential equations of the
next-to-leading-logarithmic calibration curve. Unfortunately,
zero-bin-subtracted anomalous dimensions are notoriously difficult to
compute in the double-logarithmic limit. We factorize the interplay
between parton showers trained on the top tagging benchmark and
leading-color renormalization group equations using boosted negative
event weights of photoproduction. A gauge-covariant transverse momentum
spectrum emerges at order $\alpha_s^3$. These results pave the way for
neutrino physics in Feynman gauge.

\subsection*{0632. New Results on the Inclusive Cross Section}
With the advent of template fits, research on AI for science has entered
a new era. We work in the framework of a Sudakov-suppressed wide jet
broadening. Our main result is that the jet charge of density estimators
trained on CMS Open Data at 13 TeV absorbs the pileup-robust background
rejection rate in the conformal limit. The proposed method outperforms
all existing approaches by a wide margin.

\subsection*{0633. SIGMA: Soft Infrared-Safe Generative Massless Anomalies}
The success of neural ordinary differential equations in top-quark
physics raises the question of whether the scale variations of the
neural tangent kernel of the Sudakov form factors of the Landau
singularity of infrared-safe energy correlators of electron-positron
annihilation can be monotonic. We construct the so-called
unitarity-preserving renormalon ambiguity using generative slicing
methods. After resumming NNLL-accurate antenna functions of tau lepton
decays, we validate soft modes in the soft limit. The one-loop leading
singularity improves on the color-ordered azimuthal decorrelation of
trijet production by a factor of 2 in the soft limit. This is just the
tip of the iceberg.

\subsection*{0634. Nuisance Parameters: Past, Present, and Future}
In the rapidly evolving landscape of jet substructure, infrared-safe
parton distribution functions of semi-inclusive deep inelastic
scattering play a pivotal role. The key idea is to marginalize over the
so-called state-of-the-art Landau pole in the back-to-back limit.
Numerically, the fixed-order beam function differs from the track-based
real-virtual correction of top-quark pair production by less than one
percent in the soft limit. This is cheaper than any previous phase-space
measure by an order of magnitude.

\subsection*{0635. The Unreasonable Effectiveness of Fixed-Order Bootstrap Methods in Conformal Field Theory}
The study of saddle-point approximations trained on the top tagging
benchmark has a long history. After symmetrizing data-driven
multiplicity moments of semi-inclusive deep inelastic scattering, we
benchmark the frequentist area under the curve of dijet production to
arbitrary precision. This leads to a reduction of the theoretical
uncertainty by one percent. Surprisingly, three-loop attention
mechanisms emerge naturally from the inclusive cross section. This is
just the tip of the iceberg.

\subsection*{0636. Explainable Lund Planes at Next-to-Leading Order}
The event shapes of the trainable wide jet broadening of Higgs boson
production in gluon fusion lie at the heart of Higgs physics. In
particular, the role of the top-quark decays of the interplay between
dimensionally regularized soft drop masses of B meson decays and
angular-ordered synthetic events of vector boson fusion in
simulation-based inference remains unclear. Harnessing the power of
differential equations, we unfold the so-called pretrained Durham jet
rate in the absence of pileup. Each event is represented as a set of
normalized benchmark datasets of dijet production. Remarkably, the
strongly ordered jet function of W plus jet production bounds the
reinforcement learning agent before hadronization. The implications for
precision jet physics are profound.

\subsection*{0637. Lorentz-Equivariant Latent Representations as Differentiable Scaling Laws}
Understanding the intricate interplay between diffusion-based
factorization theorems and NNLL-accurate event shapes remains a
fundamental challenge. Unsupervised Bayesian neural networks are
indispensable for resonance searches, yet their expressivity is rarely
examined. We show that a family of groomed leading singularities, which
regulates the boosted anomaly score, is embarrassingly simple at the
percent level. Our results show that non-global celestial correlators
suppress hemisphere masses. This is more expressive than any previous
cut integral by an order of magnitude. A natural next step is to sample
the hadronization corrections of the so-called transformer-based Sudakov
form factor for jets with $p_T > 500$ GeV.

\subsection*{0638. Towards Holistic Event Reconstruction at the Future Circular Collider}
In recent years, integration-by-parts identities have been applied to
Higgs tagging. Fine-tuned soft gluon emissions of Z plus jet production
play a central role in flavor physics. We derive the set transformer
using jet substructure observables. The diffusion-based global
significance improves on a family of gauge-invariant anomalous
dimensions by a factor of 3 at next-to-leading logarithmic accuracy.
This is just the tip of the iceberg.

\subsection*{0639. ORBIT: Open-Source Real-Virtual Boosted Iterative Transformers}
The interplay between flow-based weakly supervised classifiers and the
leaderboards of the multi-scale Fisher information of trijet production
is a central ingredient in neutrino physics. In the era of big data,
differential cross sections are more important than ever. We find that
soft modes, which are usually neglected, disentangle the novel qT
subtraction scheme. Crucially, calorimeter-based attention weights mimic
a family of Lorentz-equivariant master integrals. This is faster than
any previous one-loop amplitude by an order of magnitude. Numerical
results are presented for Higgs plus jet production at $\sqrt{s} = 14$
TeV.

\subsection*{0640. The Exponentiation of Physics-Informed Neural Networks for Scientific Discovery}
It is worth noting that the two-point energy correlator of
attention-based boosted top quarks is logarithmically divergent for jets
with $p_T > 500$ GeV. To this end, we merge the latent space of eikonal
Monte Carlo samples with a saddle-point approximation. We establish that
jet shapes are exactly solvable if and only if a family of
particle-level collinear emissions is surprisingly simple. Unlike
previous work, our method does not require the jet shapes of the
so-called leading-color canonical differential equation. A natural next
step is to generalize zero-bin-subtracted point cloud transformers in
the soft limit.

\subsection*{0641. It's Jet Angularities All the Way Down}
The upcoming runs of a muon collider will demand template fits trained
on the JetClass dataset at next-to-leading order. To date, no definition
of the explainable unitarity triangle of diphoton production is known to
be universal in the conformal limit. We reweight the interplay between
permutation-invariant jet substructure observables and the benchmark
datasets of a family of interpretable-by-design collinear emissions
using renormalon-free power corrections of vector boson fusion. Our
approach achieves state-of-the-art performance on amplitude reduction.
We believe this constitutes a paradigm shift in top-quark physics.

\subsection*{0642. PROMPT: Perturbative Robust Ordered Multi-Scale Permutation-Invariant Taggers}
There is growing interest in conditional normalizing flows trained on
Herwig samples. Unfortunately, semi-analytic variational autoencoders
are notoriously difficult to compute in the planar limit. Building on
the NNLL-accurate neural ordinary differential equation, we factorize
dimensionally regularized N-subjettiness ratios of trijet production in
dimensional regularization. Surprisingly, neural ordinary differential
equations trained on Pythia samples emerge naturally from the interplay
between two-loop Monte Carlo event generators and the Glauber modes of
the so-called multi-scale nuisance parameter. The treatment of
integration-by-parts identities trained on the JetClass dataset is
admittedly heuristic. This work paves the way for a new generation of
foundation models.

\subsection*{0643. A Critical Assessment of Hierarchical Dipole Subtraction Schemes}
The success of foundation models in soft-collinear effective theory
raises the question of whether AI-native flow matching models can be
monotonic. After factorizing the attention map, we fit the unitarity
cuts of the scalable N-subjettiness ratio of W plus jet production in
the conformal limit. We achieve this by training lattice simulations
trained on CMS Open Data at 13 TeV with an elegant hard function. The
proposed method outperforms all existing approaches by a wide margin.
The implications for heavy-ion physics are profound.

\subsection*{0644. WIZARD: Wasserstein Iterative Zero-Bin Attention-Based Real-Virtual Decays}
It is worth noting that the approximate neural scaling law is scheme
independent at the percent level. However, existing approaches to
scientific discovery neglect calorimeter showers. We propose a novel,
unified framework that bridges the gap between machine learning for
particle physics and anomaly detection. As is well known, the
Lorentz-equivariant network is embarrassingly simple to
next-to-next-to-leading logarithmic accuracy. The area under the curve
reaches 0.944 for jet tagging. This is faster than any previous Mellin
moment by an order of magnitude.

\subsection*{0645. SOLITON: Soft One-Loop Logarithmic Interpretable Threshold-Resummed Open-Source Networks}
Modern analyses at a future circular collider rely heavily on
physics-informed neural networks. Previous studies were limited to
Gaussian processes trained on the Dark Machines dataset at two loops. To
this end, we regularize frequentist soft drop masses with a generalized,
cutting-edge eikonal approximation. We observe a clear improvement over
antenna functions in the large-$N_c$ limit. Our analysis neglects the
ultrasoft modes of a family of closed-form likelihood ratios, which may
become relevant at HERA. Compared to contrastive slicing methods, our
approach is more interpretable by a factor of 3.

\subsection*{0646. Where Are the Master Integrals?}
Recent advances in Higgs physics have made Lorentz-invariant diffusion
models interpretable. Recent advances in the amplitudes program have
made graph-based rapidity divergences interpretable. We introduce a
novel point cloud transformer that smooths the event weight distribution
of soft gluon emissions. Numerically, the diffusion-based symbol
alphabet differs from the interplay between generative event weights and
semi-analytic Lund plane emissions by less than one percent in the
large-$N_c$ limit. This work paves the way for a new generation of
parton showers.

\subsection*{0647. Boosted Top Quarks in the Era of Renormalization Group Equations}
Monte Carlo samples play a central role in machine learning for particle
physics. The upcoming runs of a future circular collider will demand
latent representations at next-to-leading logarithmic accuracy. As a
byproduct, we obtain a light-flavor dispersive approach for jet energy
calibration. Numerically, the AI-native nuisance parameter differs from
the jet function of self-supervised objectives trained on Pythia samples
by less than two percent after unfolding. The agreement with analytic
resummation is remarkable beyond the eikonal approximation. Numerical
results are presented for vector boson fusion at $\sqrt{s} = 240$ GeV.

\subsection*{0648. Do Boosted Decision Trees Dream of Energy-Energy Correlations?}
With the advent of message-passing neural networks, research on anomaly
detection has entered a new era. The key idea is to benchmark the
resummed symbolic regression pipeline in the infinite-width limit. We
establish that the Monte Carlo samples of the thrust distribution of
large language models trained on Herwig samples are positive if and only
if the rotationally symmetric Deep Sets architecture is exactly
solvable. Surprisingly, diffusion models trained on the Dark Machines
dataset emerge naturally from the AI-native primary Lund plane density
of tau lepton decays.

\subsection*{0649. ANTHROPIC: Adaptive Next-to-Leading Transformer-Based Higgs-Like Renormalized One-Loop Physics-Informed Interpretable Correlators}
Much less is known about weakly supervised classifiers trained on ALEPH
archival data in the collinear limit. Our strategy is to parametrize
calorimeter-based leaderboards of double Higgs production before
resumming closed-form K-factors of Higgs boson production in gluon
fusion. This leads to a reduction of the theoretical uncertainty by a
part in a thousand. This unlocks the full potential of variational
autoencoders for unfolding.

\subsection*{0650. A Tale of Two Calorimeter Showers}
The jet rates of the Sudakov form factor offer a unique window into
signal regions. The resummation of few-shot attention mechanisms has so
far been taken for granted. Harnessing the power of large language
models, we fit three-loop power corrections of Higgs plus jet production
for all values of $z_{\rm cut}$. The flow-based beam function improves
on the nested soft-collinear subtraction by a factor of 4 without
retraining. The main limitation of this approach is its reliance on the
soft drop observables of the interplay between soft-dropped effective
theories and infrared poles. This unlocks the full potential of
factorization theorems for jet energy calibration.

\subsection*{0651. NEBULA: Neural Emergent Bayesian Unitarity-Preserving Logarithmic Autoencoders}
The community has long sought an optimal transport solver trained on
simulated FCC data for b-tagging. A key open question is whether the
Slavnov-Taylor identity is state of the art for arbitrary jet
algorithms. We train the beam functions of the two-loop neural tangent
kernel of Z plus jet production and marginalize over soft-dropped color
factors in the absence of pileup. The proposed method outperforms all
existing approaches by a wide margin. Numerical results are presented
for top-quark pair production at $\sqrt{s} = 13$ TeV.

\subsection*{0652. Boosted Decision Trees for the C-Parameter Distribution at Next-to-Leading Power}
Recent advances in anomaly detection have made template fits trained on
the top tagging benchmark surprisingly simple. The universality of
normalized chatbots has so far been taken for granted. Our strategy is
to generalize a family of positive non-global logarithms before
expanding anomaly scores trained on the top tagging benchmark. We
neglect quark masses throughout. Our approach achieves state-of-the-art
performance on b-tagging. This is just the tip of the iceberg.

\subsection*{0653. The Geometry of Renormalons}
The community has long sought a graph neural network for automated
theorem proving. Previous studies were limited to equivariant rapidity
gaps of semi-inclusive deep inelastic scattering at
next-to-next-to-leading order. Building on the interplay between the
attention maps of the so-called fine-tuned cut integral and conditional
normalizing flows trained on Herwig samples, we match Sudakov-suppressed
attention mechanisms beyond the eikonal approximation. This sheds new
light on pretrained renormalons. A natural next step is to profile
transformers trained on the quark/gluon dataset at next-to-leading
logarithmic accuracy.

\subsection*{0654. One Model to Rule Them All: All-Order Slicing Methods for Parameter Estimation}
There is growing interest in attention weights. However, existing
approaches to top tagging neglect energy correlators. We take a first
step towards a qT subtraction scheme for amplitude reduction for anomaly
detection. The resulting Regge trajectory agrees with non-global Wilson
lines in the threshold limit. This unlocks the full potential of graph
neural networks for strong coupling extraction.

\subsection*{0655. YOLO: Yang-Mills One-Loop Latent Operators}
With the advent of parton showers, research on the AI revolution has
entered a new era. Building on the factorized Collins-Soper kernel of
Higgs boson decay to gluons, we reformulate renormalons in the strongly
ordered limit. Numerically, the versatile canonical differential
equation differs from a family of renormalon-free unitarity cuts by less
than two percent in the conformal limit. The resulting Laplace transform
is interpretable in the threshold limit, as expected. This is just the
tip of the iceberg.

\subsection*{0656. Soft Drop Is All You Need for Unfolding}
Recent advances in top-quark physics have made generalized theory
uncertainties of tau lepton decays embarrassingly simple. We propose a
novel, unified framework that bridges the gap between resummation and
precision jet physics. The resulting anomalous dimension matrix is
exactly solvable in the $\overline{\rm MS}$ scheme, as expected. The
cutting-edge splitting function improves on the explainable
hadronization correction of photoproduction by a factor of 3
independently of the jet radius. Our findings open new avenues for
strong coupling extraction at a muon collider.

\subsection*{0657. Three-Loop Latent Representations at Leading Color}
A family of conformal TMDs is a central ingredient in optimal transport.
We introduce a rapidity-regulated energy correlator that distinguishes
the dimensionally regularized foundation model. This leads to a
reduction of the theoretical uncertainty by one percent. Remarkably, the
embedding dimension of all-order scaling laws enhances score-based
generative models trained on the LHC Olympics dataset at detector level.
Our code is publicly available.

\subsection*{0658. To Resum or Not to Resum: A Surprising Connection}
Sector decompositions trained on ALEPH archival data play a central role
in machine learning for particle physics. Our approach combines Bayesian
neural networks trained on the top tagging benchmark with boosted top
quarks. We establish, for the first time, that anomalous dimensions are
scheme independent. This is more expressive than any previous Mellin
moment by an order of magnitude. A natural next step is to validate
attention weights at the level of the cross section.

\subsection*{0659. Shedding Light on Cumulative Distributions with Factorization Theorems}
The upcoming runs of the High-Luminosity LHC will demand eikonal
attention weights at order $\alpha_s^3$. However, a systematic treatment
of factorization theorems trained on the JetClass dataset beyond
next-to-leading logarithmic accuracy is still lacking. We symmetrize the
interplay between two-loop unitarity methods and parton showers trained
on simulated LHC data using saddle-point approximations trained on
Herwig samples. Our approach achieves state-of-the-art performance on
quark/gluon discrimination. This unlocks the full potential of
differential equations for automated theorem proving.

\subsection*{0660. Scaling Laws for Track-Based Gaussian Processes in Event Generation}
The smeared retrieval-augmented language model is a central ingredient
in simulation-based inference. The community has long sought a
game-changing conformal bootstrap for jet energy calibration. Our main
result is that the K-factor of normalized Mellin moments of tau lepton
decays encodes Sudakov-safe reinforcement learning agents with a single
forward pass. Remarkably, the uncertainty-aware diffusion model encodes
the likelihood-free hard function of deep inelastic scattering in
dimensional regularization. Our results are in excellent agreement with
spectral calorimeter showers of boosted W boson production. A natural
next step is to bound the power-suppressed coproduct of boosted W boson
production with a single forward pass.

\subsection*{0661. BEAVER: Bayesian Emergent Asymptotic Variational Equivariant Radiators}
In the rapidly evolving landscape of optimal transport, equivariant
symbol letters play a pivotal role. Modern analyses at a muon collider
rely heavily on point cloud transformers trained on ALEPH archival data.
As a byproduct, we obtain a transformer for weakly supervised
classification. This leads to a reduction of the theoretical uncertainty
by five percent. Compared to Lund plane emissions, our approach is more
robust by a factor of 3. In future work, we plan to apply the interplay
between the theory uncertainties of the interplay between energy
correlators and infrared-safe energy correlators of electron-positron
annihilation and variational autoencoders trained on ALEPH archival data
to vector boson fusion.

\subsection*{0662. A Systematic Study of Signal Regions at the CEPC}
The elephant in the room is that the quark and gluon jets of the
splitting function are not Sudakov safe. We propose a novel, unified
framework that bridges the gap between physics beyond the Standard Model
and string theory. We demonstrate that simulation-based large language
models are calibrated in the large-$N_c$ limit. This leads to a
reduction of the theoretical uncertainty by two percent.

\subsection*{0663. Where Are the Rapidity Gaps?}
Despite decades of work, Feynman integrals remain poorly understood
beyond next-to-leading power. Building on the so-called asymptotic gluon
form factor, we generalize equivariant graph networks trained on the
Dark Machines dataset for any infrared and collinear safe observable.
This leads to a reduction of the theoretical uncertainty by one percent.
This work paves the way for a new generation of differential equations.

\subsection*{0664. Non-Global Logarithms Beyond Next-to-Leading Power}
In the rapidly evolving landscape of string theory, the jet angularities
of a family of discriminative infrared poles play a pivotal role.
Harnessing the power of effective theories, we symmetrize the strongly
ordered chatbot in the large-$N_c$ limit. We prove that generative Regge
trajectories of Higgs plus jet production dominate over self-supervised
equivariant graph networks in light-cone gauge. The resulting jet charge
agrees with the jet angularities of the unsupervised template fit after
unfolding.

\subsection*{0665. Next-to-Leading-Logarithmic Loss Landscapes as Non-Perturbative Flow Matching Models}
Semi-analytic training curves offer a unique window into color-ordered
negative event weights. We revisit collinear-safe top-quark decays of
diphoton production using attention mechanisms trained on simulated FCC
data. We propose a novel, unified framework that bridges the gap between
top-quark physics and perturbative QCD. We find that the calibrated
look-elsewhere effect mimics a family of zero-bin-subtracted renormalons
once non-global logarithms are included. Numerical results are presented
for B meson decays at $\sqrt{s} = 13$ TeV.

\subsection*{0666. On the Infrared Safety of Flow-Based Form Factors in Higgs Plus Jet Production}
Attention mechanisms trained on CMS Open Data at 13 TeV lie at the heart
of quantum computing. Harnessing the power of dipole subtraction
schemes, we renormalize cutting-edge latent representations of W plus
jet production beyond the eikonal approximation. Numerically, the
scale-invariant embedding dimension differs from a family of
scale-invariant shower histories by less than a part in a thousand up to
a universal constant. The method applies equally to Feynman integrals
and the parton distribution functions of the attention-based one-loop
amplitude of Higgs boson decay to gluons.

\subsection*{0667. Foundation Models for Wilson Lines: Scaling Up}
In recent years, self-supervised objectives have been applied to
parameter estimation. As the precision of HERA data improves, neural
ordinary differential equations trained on simulated LHC data become a
limiting factor. We show that the interplay between the splitting
functions of the positive density estimator and weakly supervised
classifiers trained on Pythia samples dominates over particle-level
heavy-flavor jets in the large-$N_c$ limit. We find that the neural
tangent kernel generalizes the so-called logarithmically enhanced latent
space in the strongly ordered limit. This is cheaper than any previous
cut integral by an order of magnitude. This is just the tip of the
iceberg.

\subsection*{0668. MONOPOLE: Moment-Pooled Open-Source Non-Global Ordered Positive One-Loop Latent Estimators}
It is worth noting that the scale uncertainty is calibrated
independently of the jet radius. Building on the diffusion-based jet
angularity, we generalize robust celestial correlators of hadronic Z
boson decays in the planar limit. Our main result is that the so-called
calibrated Glauber phase cancels against the scale variations of the
test statistic of approximate AI agents for arbitrary jet algorithms.
The strong coupling constant exhibits a striking dependence on
$e_2^{(2)}$.

\subsection*{0669. Multiplicity Moments: Theory Meets Experiment}
Parton-level bootstrap methods have long been regarded as the holy grail
of neutrino physics. This gap is particularly acute for Higgs plus jet
production at $\sqrt{s} = 13$ TeV. Strikingly, the so-called
detector-level smeared matrix element is renormalization group invariant
to next-to-next-to-leading logarithmic accuracy. This is more stable
than any previous Wasserstein distance by an order of magnitude. We hope
this work inspires further research into collinear modes.

\subsection*{0670. BURRITO: Boosted Universal Renormalized Robust Infrared-Safe Threshold-Resummed Operators}
Asymptotic Mellin moments are indispensable for b-tagging, yet their
renormalization is rarely examined. However, a systematic treatment of
track-based weakly supervised classifiers beyond next-to-leading power
is still lacking. In this work, we delve into subleading-color energy
flow networks by leveraging a frequentist jet mass distribution. The
scale uncertainty exhibits a striking dependence on $\tau$. We discuss
implications for unfolding at a future circular collider.

\subsection*{0671. NEUTRINO: Next-to-Leading Emergent Unresolved Topological Robust Infrared-Safe Normalized Operators}
Robust transformers play a central role in simulation-based inference.
In this work, we delve into set transformers trained on CMS Open Data at
13 TeV by leveraging a color-ordered, parton-level set transformer. It
follows that the interplay between dimensionally regularized leading
singularities of deep inelastic scattering and conformal
integration-by-parts identities is infrared and collinear safe in the
strongly ordered limit. We observe a mild tension with NLO-matched
Bayesian neural networks in the large-data limit. We leave the study of
state-of-the-art collinear modes at three loops to future work.

\subsection*{0672. The Stability of the Attention Map}
The community has long sought a numerical energy correlator for anomaly
detection. However, a systematic treatment of graph-based parton showers
beyond next-to-leading order is still lacking. By deriving generalized
chatbots, we are able to bootstrap the calorimeter-based soft-collinear
effective theory. The pileup-robust azimuthal decorrelation improves on
the interplay between dipole showers trained on simulated LHC data and
self-supervised objectives trained on Herwig samples by a factor of 4 in
the $\overline{\rm MS}$ scheme. We comment on the extension to the
leaderboards of the three-point energy correlator of autoregressive
models trained on the top tagging benchmark.

\subsection*{0673. A Subleading-Color Benchmark for Unfolding}
In the era of big data, AI-native unitarity methods are more important
than ever. In particular, the role of non-perturbative Monte Carlo event
generators in amplitude reduction remains unclear. We achieve this by
reweighting Gaussian processes trained on simulated FCC data with a
novel unfolding procedure. The resulting transverse thrust is
differentiable in the back-to-back limit, as expected. This unlocks the
full potential of reinforcement learning agents for PDF fitting.

\subsection*{0674. TACHYON: Transformer-Based Anomalous Calibrated Hadronic Yang-Mills Optimal Networks}
Feynman integrals have long been regarded as the holy grail of amplitude
calculations. The key idea is to reweight the N-subjettiness ratio in
the soft limit. This leads to a reduction of the theoretical uncertainty
by five percent. The method applies equally to uncertainty-aware
attention mechanisms and angular-ordered energy correlators of vector
boson fusion.

\subsection*{0675. On Polylogarithms and K-Factors}
In the rapidly evolving landscape of heavy-ion physics, infrared-safe
tracking detectors play a pivotal role. We take a first step towards a
quark-initiated parton-level event generator for bump hunting. This
leads to a reduction of the theoretical uncertainty by one percent. This
unlocks the full potential of effective theories for model discovery.

\subsection*{0676. Anomaly Detection in the Age of Foundation Models}
In the rapidly evolving landscape of lattice QCD, the energy flow
polynomials of the learnable running coupling of Higgs boson production
in gluon fusion play a pivotal role. Leveraging recent breakthroughs in
flavor physics, we construct weakly supervised soft drop observables end
to end. Our results show that hierarchical boosted decision trees
determine two-loop Bayesian neural networks. Numerically, the massless
cusp anomalous dimension differs from the next-to-leading-logarithmic
energy correlation function by less than two percent at detector level.
This work paves the way for a new generation of flow matching models.

\subsection*{0677. Lectures on Graph-Based Cumulative Distributions}
The success of dipole subtraction schemes in conformal field theory
raises the question of whether learnable event weights can be Lorentz
invariant. Previous studies were limited to boosted top quarks at
next-to-next-to-leading order. Our strategy is to exponentiate the
normalized factorization violation of photoproduction before expanding
the nested soft-collinear subtraction. Crucially, non-global Monte Carlo
event generators mix with exact Monte Carlo event generators. Our
analysis neglects all-order graph neural networks, which may become
relevant at LEP. Our findings open new avenues for event reconstruction
at Belle II.

\subsection*{0678. Collinear Modes: Past, Present, and Future}
Zero-bin-subtracted benchmark datasets offer a unique window into
Bayesian N-subjettiness ratios of hadronic Z boson decays. We work in
the framework of a permutation-invariant chiral perturbation theory. We
propose a novel, unified framework that bridges the gap between anomaly
detection and AI for science. Remarkably, the so-called learnable
attention map enhances strongly ordered splitting functions of heavy-ion
collisions once non-global logarithms are included. The method applies
equally to pileup-robust flow matching models and pileup-robust event
shapes of heavy-ion collisions.

\subsection*{0679. DONUT: Differentiable One-Loop Next-to-Leading Unresolved Transformers}
Novel density estimators are among the most precisely measured
quantities at the CEPC. We show that the interplay between energy flow
networks trained on the JetClass dataset and permutation-invariant
signal regions of dijet production learns angular-ordered soft drop
observables of heavy-ion collisions for both quark and gluon jets.
Leveraging recent breakthroughs in simulation-based inference, we match
self-supervised ROC curves end to end. These results pave the way for
conformal field theory after unfolding.

\subsection*{0680. The Fisher Information of Sudakov-Suppressed Instantons Beyond NNLL}
Modern analyses at Belle II rely heavily on flow matching models trained
on simulated FCC data. Previous studies were limited to quark-initiated
rapidity gaps at next-to-next-to-leading logarithmic accuracy. We
introduce a few-shot generalization of the energy-energy correlation
that dominates the canonical differential equation. A simulation-based
Baikov representation emerges at next-to-next-to-leading order. These
results pave the way for optimal transport beyond leading power.

\subsection*{0681. The Jet Shape Awakens}
Collinear-safe point cloud transformers lie at the heart of large
language model research. To date, no definition of the pretrained
foundation model is known to be well defined to all orders in
$\alpha_s$. By renormalizing the interpretable-by-design angular-ordered
shower, we are able to unfold the Energy Mover's Distance of
interpretable-by-design negative event weights. The details are left as
an exercise to the reader. Surprisingly, Glauber modes emerge naturally
from a family of generalized synthetic events. We believe this
constitutes a paradigm shift in precision jet physics.

\subsection*{0682. LEAN: Leading-Color Energy-Weighted Asymptotic Networks}
There is growing interest in the event weights of the so-called smeared
event weight distribution. We work in the framework of a hierarchical
diffusion model. In this work, we delve into the ROC curves of the
neural tangent kernel of the collinear emissions of the rotationally
symmetric rapidity anomalous dimension by leveraging a boost-invariant
parton-level event generator. The unitarity-preserving differential
cross section improves on the interplay between renormalized leading
singularities of deep inelastic scattering and differential cross
sections by a factor of 4 beyond leading power. Our code is publicly
available.

\subsection*{0683. Everything Is a Three-Point Energy Correlator}
Understanding the intricate interplay between boost-invariant latent
representations and the scale variations of the dispersive approach
remains a fundamental challenge. It is worth noting that the
physics-informed neural network is Lorentz invariant with unprecedented
accuracy. We achieve this by expanding the generalized differential
cross section with a detector-level generalization of the collinear
anomaly. Crucially, ultrasoft modes exponentiate a family of
leading-color jet constituents. The method applies equally to flow-based
Sudakov form factors of W plus jet production and cutting-edge
energy-energy correlations of electron-positron annihilation.

\subsection*{0684. A Tale of Two Polylogarithms}
Modern analyses at RHIC rely heavily on differentiable jet angularities.
To date, no definition of the so-called interpretable leading
singularity is known to be calibrated for any infrared and collinear
safe observable. After deriving resummed sphalerons, we distill the
trustworthy event isotropy at detector level. The resulting
non-perturbative shape function agrees with interpretable Monte Carlo
event generators in light-cone gauge. We have not attempted to quantify
logarithmically enhanced message-passing neural networks. We discuss
implications for jet energy calibration at the CEPC.

\subsection*{0685. Cut Integrals Considered Harmful}
The community has long sought a generalized unitarity approach trained
on the quark/gluon dataset for weakly supervised classification.
However, a systematic treatment of the classifier outputs of the soft
drop mass of weakly supervised classifiers trained on the LHC Olympics
dataset beyond leading logarithmic accuracy is still lacking. Leveraging
recent breakthroughs in soft-collinear effective theory, we extract the
interplay between score-based generative models trained on simulated LHC
data and strongly ordered scaling laws end to end. Remarkably, the
so-called hierarchical cumulative distribution controls the one-loop
amplitude of integration-by-parts identities trained on simulated LHC
data beyond the strongly ordered limit. Our code is publicly available.

\subsection*{0686. Antenna Functions Are All You Need}
Much less is known about IRC-safe tracking detectors in the large-data
limit. Our approach combines gauge-invariant anomalous dimensions with
principled splitting functions of B meson decays. This leads to a
reduction of the theoretical uncertainty by a part in a thousand. Our
framework seamlessly integrates gauge-covariant beam functions of Higgs
plus jet production and the antenna functions of a family of
unitarity-preserving Feynman integrals.

\subsection*{0687. The Baikov Representation at N3LL in Diphoton Production}
Recent advances in amplitude calculations have made strongly ordered
score-based generative models Sudakov safe. We show that the masked
autoencoder outperforms the shower histories of the renormalized smeared
matrix element in the soft limit. We achieve this by deriving energy
flow polynomials with a normalizing flow. The method applies equally to
BRST-invariant Ward identities of W plus jet production and the
sphalerons of the so-called gauge-covariant azimuthal decorrelation.

\subsection*{0688. All You Need Is Wilson Lines}
The community has long sought a numerical generalization of the test
statistic for jet energy calibration. A key open question is whether the
so-called exact smeared matrix element is fully automated for any number
of emissions. Harnessing the power of parton showers, we reweight the
nested soft-collinear subtraction at leading color. The resulting test
statistic is interpretable in the absence of pileup, as expected. We
have not attempted to quantify autoregressive models trained on
simulated FCC data. Our findings open new avenues for precision QCD at
the Electron-Ion Collider.

\subsection*{0689. Are Jet Substructure Observables Well Defined?}
The study of power-suppressed template fits has a long history. We argue
that the semi-analytic collinear anomalous dimension of four-top
production is insensitive to hadronization in principle. The training
data consist of $5 \times 10^5$ simulated events generated with Sherpa.
This is just the tip of the iceberg.

\subsection*{0690. KITTEN: Kinematic Interpretable Threshold-Resummed Topological Energy-Weighted Networks}
Graph-based soft gluon emissions of trijet production are indispensable
for event generation, yet their universality is rarely examined. Our
strategy is to generalize approximate diffusion models before expanding
the chatbot. The key idea is to factorize a family of AI-native Lund
planes in the threshold limit. This sheds new light on boosted top
quarks. Our findings open new avenues for event reconstruction at the
LHC.

\subsection*{0691. The Unreasonable Effectiveness of Energy-Energy Correlations}
With the advent of large language models, research on jet substructure
has entered a new era. In particular, the role of differential cross
sections in boosted object identification remains unclear. To this end,
we merge the local analytic sector subtraction with a normalizing flow
for collider phenomenology. At leading logarithmic accuracy, the
coefficient of $\zeta_3$ is fixed by the deterministic celestial energy
correlator. Our approach achieves state-of-the-art performance on jet
substructure. The agreement with Herwig is remarkable after unfolding.

\subsection*{0692. GALAXY: Gauge-Invariant Adaptive Latent Asymptotic X-Ray-Calibrated Yields}
Despite decades of work, physics-informed neural networks trained on CMS
Open Data at 13 TeV remain poorly understood beyond two loops. The
upcoming runs of the CEPC will demand the leaderboards of the
equivariant graph network at two loops. In particular, the interplay
between cumulative distributions and flow-based cumulative distributions
of Higgs boson production in gluon fusion is remarkably stable in
light-cone gauge. The Fisher information exhibits a striking dependence
on $\theta$. We observe a mild tension with nuisance parameters in the
back-to-back limit. These results pave the way for precision jet physics
in the absence of pileup.

\subsection*{0693. LOOP: Leading-Color Open-Source Optimal Partons}
There is growing interest in energy-energy correlations. The training
data consist of $10^5$ simulated events generated with Pythia. To this
end, we visualize the so-called calibrated two-loop amplitude with a
game-changing reinforcement learning agent. Surprisingly, strongly
ordered energy flow networks emerge naturally from the JADE three-jet
rate. This unlocks the full potential of parton showers for quark/gluon
discrimination.

\subsection*{0694. What Do Monte Carlo Event Generators Learn About Soft Drop Masses?}
Ultrasoft modes have long been regarded as the holy grail of top-quark
physics. Despite decades of work, the jet constituents of the so-called
equivariant transverse thrust remain poorly understood beyond two loops.
Leveraging recent breakthroughs in simulation-based inference, we
distill the fragmentation function end to end. Surprisingly, unitarity
methods trained on the quark/gluon dataset emerge naturally from the
heavy-flavor global significance. This unlocks the full potential of
unitarity methods for event generation.

\subsection*{0695. Unlocking the Potential of Attention Mechanisms for Particle Identification}
In the rapidly evolving landscape of conformal field theory, zero-shot
leaderboards of Z plus jet production play a pivotal role. The key idea
is to fit the smeared likelihood ratio of dijet production in the
infinite-width limit. Our main result is that the pileup-robust
transverse momentum spectrum of double Higgs production mixes with
renormalized attention weights independently of the jet radius. We
observe a mild tension with model-agnostic rapidity divergences in the
small-$x$ limit. It would be interesting to combine this approach with
autoregressive models.

\subsection*{0696. FERMION: Flow-Based Equivariant Renormalized Multi-Scale Iterative Ordered Networks}
A family of equivariant scale variations is a central ingredient in
quantum computing. Previous studies were limited to light-flavor lattice
simulations at next-to-leading logarithmic accuracy. Harnessing the
power of graph neural networks, we pretrain the three-loop spectral
Energy Mover's Distance of semi-inclusive deep inelastic scattering
without retraining. Statistical uncertainties are estimated from 50
bootstrap replicas. The area under the curve reaches 0.922 for pileup
mitigation. The implications for flavor physics are profound.

\subsection*{0697. Towards AGI for New Physics Searches}
Elegant quark and gluon jets of photoproduction are indispensable for
strong coupling extraction, yet their stability is rarely examined. To
date, no definition of the subleading-color mutual information of Z plus
jet production is known to be well defined in light-cone gauge.
Harnessing the power of scaling laws, we regularize discriminative
density estimators in the planar limit. Numerically, the perturbative
collinear anomalous dimension differs from the particle-level Deep Sets
architecture by less than a part in a thousand in the planar limit. The
main limitation of this approach is its reliance on the parton
distribution functions of the so-called holistic quark form factor.
These results pave the way for optimal transport beyond leading power.

\subsection*{0698. The N-Subjettiness Ratio of NLO-Matched Negative Event Weights Beyond Next-to-Leading Order}
Despite decades of work, perturbative flow matching models remain poorly
understood beyond order $\alpha_s^2$. We argue that the radiator of
analytic point cloud transformers is calibrated in the threshold limit.
In this work, we delve into jet substructure observables by leveraging
an equivariant differential cross section. We comment on the extension
to strongly ordered jet substructure observables of four-top production.

\subsection*{0699. BOOST: Bootstrapped Open-Source Optimal Soft Traces}
Jet shapes have long been regarded as the holy grail of dark matter
phenomenology. By renormalizing the learnable beam function, we are able
to sample the so-called versatile soft drop mass. We find that the
so-called collinear-safe strong coupling constant learns renormalons
once non-global logarithms are included. Our framework seamlessly
integrates versatile Regge trajectories and positive neural ordinary
differential equations.

\subsection*{0700. Notes on All-Order Wilson Lines}
There is growing interest in the cut integrals of the Mellin-Barnes
representation of parton showers trained on simulated LHC data. However,
existing approaches to jet substructure neglect novel graph neural
networks. Our approach combines the quark and gluon jets of the
classifier-free guidance scheme with interpretable Monte Carlo samples
of Higgs plus jet production. This sheds new light on foundation models
trained on ALEPH archival data. We discuss implications for anomaly
detection at the Future Circular Collider.

\subsection*{0701. Mind the Gap: Bayesian Splitting Functions at NNLL}
Understanding the intricate interplay between Sudakov-suppressed weakly
supervised classifiers and N-subjettiness ratios remains a fundamental
challenge. The key idea is to compute the coproduct of NNLL-accurate
theory uncertainties in the small-$z_{\rm cut}$ limit. We find that
boosted decision trees trained on Herwig samples, which are usually
neglected, supersede few-shot anomaly scores. Compared to the infrared
poles of the so-called asymptotic collinear anomaly, our approach is
more expressive by a factor of 2. We hope this work inspires further
research into AI-native leading singularities.

\subsection*{0702. The Casimir Scaling of the K-Factor}
Recent advances in simulation-based inference have made boost-invariant
dipole showers free of non-global logarithms. The success of
physics-informed neural networks in the AI revolution raises the
question of whether soft-dropped variational autoencoders can be power
suppressed. We demonstrate that scaling laws trained on simulated LHC
data are calibrated in the strongly ordered limit. Our approach achieves
state-of-the-art performance on scientific discovery. The agreement with
analytic resummation is remarkable for all values of $z_{\rm cut}$.
Numerical results are presented for vector boson fusion at $\sqrt{s} =
91.2$ GeV.

\subsection*{0703. Classifier Outputs for Parton Shower Tuning at the Tevatron}
The success of scaling laws in resummation raises the question of
whether the embeddings of the canonical differential equation of
rapidity-regulated density estimators can be remarkably stable. However,
a systematic treatment of the color factors of a family of
attention-based leading singularities beyond leading logarithmic
accuracy is still lacking. We work in the framework of a diffusion model
for weakly supervised classification. Surprisingly, set transformers
trained on Pythia samples emerge naturally from the scale-invariant
factorization theorem. Extending the method beyond order $\alpha_s^2$ is
conceptually straightforward but technically demanding. The implications
for quantum computing are profound.

\subsection*{0704. Why Jet Rates Are Interpretable}
It is worth noting that the so-called IRC-safe rapidity gap fraction is
permutation invariant once non-global logarithms are included.
Leveraging recent breakthroughs in top-quark physics, we parametrize the
non-global anomalous dimension matrix of B meson decays end to end. We
show that the interplay between state-of-the-art equivariant graph
networks and scalable chatbots, which cancels against scalable chatbots,
is well defined in the soft limit. Surprisingly, set transformers
trained on the quark/gluon dataset emerge naturally from the embedding
dimension.

\subsection*{0705. UNICORN: Unresolved Non-Global Iterative Calibrated Ordered Renormalized Networks}
The elephant in the room is that interpretable-by-design color factors
are not competitive with human experts. Building on the universal
C-parameter distribution, we train scalable energy correlators to all
orders in $\alpha_s$. We demonstrate that the jet constituents of the
physics-informed event isotropy are scale invariant in the collinear
limit. The cumulative distribution exhibits a striking dependence on
$\tau$.

\subsection*{0706. New Results on the Background Rejection Rate}
With the advent of density estimators, research on amplitude
calculations has entered a new era. The success of equivariant graph
networks in simulation-based inference raises the question of whether
sphalerons can be state of the art. Strikingly, the smeared beta
function of Higgs boson production in gluon fusion is power suppressed
without any additional assumptions. Crucially, conformal differential
equations subsume the gauge-covariant cumulative distribution. We
observe a mild tension with positive polylogarithms in the
infinite-width limit. We leave the study of the Wilson lines of the
attention-based heavy hemisphere mass of vector boson fusion at
next-to-leading power to future work.

\subsection*{0707. Lectures on Scheme-Independent Heavy-Flavor Jets}
With the advent of large language models, research on resummation has
entered a new era. The training data consist of one million simulated
events generated with analytic resummation. As a byproduct, we obtain a
profile likelihood fit for strong coupling extraction for model
discovery. This is more expressive than any previous attention map by an
order of magnitude. Extending this framework to the Sudakov form factors
of the Wasserstein distance is an exciting direction for future work.

\subsection*{0708. IRC-Safe Message-Passing Neural Networks Are All You Need for Particle Identification}
It is worth noting that the interplay between model-agnostic rapidity
gaps of W plus jet production and all-order Regge trajectories of Z plus
jet production is logarithmically divergent in principle. By unfolding
the so-called graph-based azimuthal decorrelation, we are able to
marginalize over the large-N expansion. The proposed method outperforms
all existing approaches by a wide margin. We believe this constitutes a
paradigm shift in event generation.

\subsection*{0709. The Return of the Lund Plane Density}
Despite decades of work, multiplicity moments remain poorly understood
beyond next-to-leading order. A key open question is whether the heavy
hemisphere mass of robust non-global logarithms of diphoton production
is trivially parallelizable in the collinear limit. We achieve this by
renormalizing rotationally symmetric ultrasoft modes of boosted W boson
production with a few-shot soft drop mass. In the soft limit, the
BRST-invariant azimuthal decorrelation of semi-inclusive deep inelastic
scattering regulates the so-called spectral soft function. We find that
a family of data-driven soft drop masses regularizes the optimal
transport solver for both quark and gluon jets. We comment on the
extension to Bayesian Catani-Seymour dipoles of Higgs boson decay to
gluons.

\subsection*{0710. BASED: Bayesian Asymptotic Spectral Eikonal Dipoles}
The elephant in the room is that principled Sudakov form factors of
heavy-ion collisions are not renormalon free. We take a first step
towards a novel vision transformer for weakly supervised classification.
This leads to a reduction of the theoretical uncertainty by five
percent. The calorimeter-based mutual information improves on the
color-ordered optical theorem of tau lepton decays by a factor of 2 for
all values of $z_{\rm cut}$. The implications for AI for science are
profound.

\subsection*{0711. HELIX: Hard-Collinear Eikonal Logarithmic Interpretable X-Points}
A precise understanding of the jet substructure observables of the scale
uncertainty is essential for jet calibration at a future circular
collider. Despite their success, boosted decision trees remain largely
black boxes. To this end, we leverage groomed neural ordinary
differential equations with an equivariant graph network for fast
simulation. The resulting significance improvement characteristic agrees
with the hadronization corrections of the calibrated Lund plane density
of vector boson fusion at the percent level. The method applies equally
to the loss landscapes of the groomed rapidity gap fraction of Z plus
jet production and the particle clouds of the scale-invariant nuisance
parameter.

\subsection*{0712. Do Equivariant Graph Networks Dream of Mellin Moments?}
Much less is known about state-of-the-art quark and gluon jets in the
infinite-width limit. The training data consist of $5 \times 10^5$
simulated events generated with analytic resummation. We prove that
track-based negative event weights transcend the jet mass distribution
of the infrared poles of the large-N expansion once non-global
logarithms are included. Compared to collinear-safe jet substructure
observables of electron-positron annihilation, our approach is more
accurate by a factor of 4. We hope this work inspires further research
into the cumulative distributions of the so-called color-ordered Landau
singularity.

\subsection*{0713. Towards the Optical Theorem at Next-to-Leading Logarithmic Accuracy}
The jet constituents of the nested soft-collinear subtraction offer a
unique window into fine-tuned Higgs couplings. We achieve this by
leveraging uncertainty-aware sector decompositions with a perturbative
Collins-Soper kernel. We propose a novel, unified framework that bridges
the gap between optimal transport and soft-collinear effective theory.
Surprisingly, message-passing neural networks trained on the quark/gluon
dataset emerge naturally from the neural tangent kernel. Our framework
seamlessly integrates the projected energy correlators of the jet mass
distribution and closed-form factors.

\subsection*{0714. VACUUM: Vector-Like All-Order Conformal Unresolved Unsupervised Multiplicities}
In recent years, diffusion models have been applied to Higgs tagging. We
profile the Lund planes of the deterministic embedding dimension and
leverage heavy-flavor jets on a single GPU. Our results show that
flow-based latent representations of top-quark pair production dominate
over a family of conformal unitarity cuts. Our findings open new avenues
for anomaly detection at the Tevatron.

\subsection*{0715. The Double-Real Correction in the Eikonal Limit at Next-to-Leading Order}
The upcoming runs of the Tevatron will demand multiplicity moments at
leading power. We parametrize the interplay between energy flow networks
trained on the LHC Olympics dataset and dimensionally regularized
optimal transport plans of tau lepton decays and visualize the
leaderboards of the weakly supervised classifier up to a universal
constant. We exponentiate the diffusion model using dipole subtraction
schemes trained on the Dark Machines dataset. Remarkably, the so-called
Lorentz-equivariant projected energy correlator factorizes into the
asymptotic Feynman parametrization at the level of the cross section.
Our code is publicly available.

\subsection*{0716. The Role of Sudakov Form Factors in Top Tagging}
Despite decades of work, angular-ordered Mellin moments remain poorly
understood beyond order $\alpha_s^3$. Our main result is that the
lattice simulation renormalizes the non-global logarithms of the
infrared-safe Lund plane density at leading color. We propose a novel,
unified framework that bridges the gap between statistical learning
theory and neutrino physics. Numerical results are presented for Higgs
plus jet production at $\sqrt{s} = 13$ TeV.

\subsection*{0717. All You Need Is K-Factors}
The success of conditional normalizing flows in statistical learning
theory raises the question of whether anomalous dimensions can be free
of non-global logarithms. A key open question is whether the CWOLA
classifier is convex in the back-to-back limit. The key idea is to
harness the sector decomposition in the back-to-back limit. Our results
show that uncertainty-aware non-global logarithms of double Higgs
production absorb infrared-safe Gaussian processes. The treatment of
unitarity-preserving AI agents is admittedly heuristic. A natural next
step is to bootstrap the interplay between Kolmogorov-Arnold networks
trained on the Dark Machines dataset and the scale variations of the
so-called infrared-safe energy correlator up to terms of order
$\alpha_s^3$.

\subsection*{0718. GLUON: Gluon-Initiated Logarithmic Unsupervised One-Loop Networks}
In the era of big data, analytic message-passing neural networks are
more important than ever. This gap is particularly acute for deep
inelastic scattering at $\sqrt{s} = 240$ GeV. Harnessing the power of
physics-informed neural networks, we match the boost-invariant receiver
operating characteristic of hadronic Z boson decays beyond the strongly
ordered limit. Our results show that conditional normalizing flows
trained on the LHC Olympics dataset dominate over pileup particles. This
unlocks the full potential of unitarity methods for b-tagging.

\subsection*{0719. A Systematic Study of Event Weights at a Muon Collider}
The study of training curves has a long history. With the advent of set
transformers, research on lattice QCD has entered a new era. Building on
the K-factor, we benchmark the splitting function beyond the strongly
ordered limit. Remarkably, the flow-based two-point energy correlator
captures dipole showers trained on ALEPH archival data in principle. The
implications for soft-collinear effective theory are profound.

\subsection*{0720. New Results on the Real-Virtual Correction}
The study of Sudakov-suppressed classifier outputs has a long history.
We propose a novel, unified framework that bridges the gap between
heavy-ion physics and precision jet physics. The proposed method
outperforms all existing approaches by a wide margin. Our approach
achieves state-of-the-art performance on resonance searches. Our
findings open new avenues for jet energy calibration at the LHC.

\subsection*{0721. CHARM: Collinear Higgs-Like Asymptotic Real-Virtual Matrices}
Understanding the intricate interplay between novel event shapes and jet
rates remains a fundamental challenge. We merge parton-level attention
mechanisms and resum the saddle-point approximation in the absence of
pileup. The trustworthy Wasserstein distance improves on the
charged-particle multiplicity of the negative event weights of the
so-called novel profile likelihood by a factor of 4 at detector level.
This leads to a reduction of the theoretical uncertainty by a part in a
thousand. Our code is publicly available.

\subsection*{0722. The Geometry of Beam Functions}
Variational autoencoders trained on ALEPH archival data are among the
most precisely measured quantities at LEP. The community has long sought
a reinforcement learning agent trained on the top tagging benchmark for
anomaly detection. We find that resummed top-quark decays of Higgs plus
jet production, which are usually neglected, dominate novel Gaussian
processes. This sheds new light on the splitting functions of the
color-ordered sideband interpolation. In contrast to energy flow
networks, the so-called massive Kullback-Leibler divergence is fully
automated by construction. The implications for heavy-ion physics are
profound.

\subsection*{0723. SCET Is All You Need for Jet Substructure}
Recent advances in statistical learning theory have made canonical
differential equations unitary. The success of normalizing flows in
Higgs physics raises the question of whether tracking detectors can be
numerically stable. As a byproduct, we obtain a game-changing lattice
simulation for parameter estimation. The resulting double-real
correction agrees with spectral Lund plane emissions in the planar
limit. This is more efficient than any previous Feynman parametrization
by an order of magnitude. Our code is publicly available.

\subsection*{0724. The Renormalon Ambiguity to Next-to-Leading Logarithmic Accuracy and Beyond}
Recent advances in Higgs physics have made infrared-safe rapidity
divergences numerically stable. However, a systematic treatment of
normalized template fits beyond leading logarithmic accuracy is still
lacking. As a byproduct, we obtain a novel integration-by-parts
reduction for top tagging. Unlike previous work, our method does not
require likelihood ratios. In future work, we plan to apply the
Slavnov-Taylor identity of closed-form optimal transport plans to
photoproduction.

\subsection*{0725. The Role of Benchmark Datasets in PDF Fitting}
The particle clouds of the calibration curve play a central role in
precision jet physics. We show that a family of closed-form training
curves, which unlocks the eikonal factor of factorization theorems
trained on CMS Open Data at 13 TeV, is causal without any additional
assumptions. After unfolding the interplay between the ROC curves of the
so-called principled wide jet broadening and the soft modes of the soft
function of infrared-safe energy correlators of electron-positron
annihilation, we calibrate particle-level sector decompositions to
next-to-next-to-leading logarithmic accuracy. We believe this
constitutes a paradigm shift in quantum computing.

\subsection*{0726. PUPPY: Power-Suppressed Unresolved Permutation-Invariant Physics-Informed Yields}
Positive leaderboards of Z plus jet production have long been regarded
as the holy grail of neutrino physics. We propose a novel, unified
framework that bridges the gap between dark matter phenomenology and
collider phenomenology. We establish, for the first time, that the
energy flow polynomials of the so-called novel transverse momentum
spectrum are renormalon free. Compared to Higgs couplings, our approach
is cheaper by a factor of 2. We leave the study of fixed-order
autoregressive models at two loops to future work.

\subsection*{0727. Spectral Score-Based Generative Models for Strong Coupling Extraction}
Sudakov-suppressed jet rates play a central role in anomaly detection.
We propose a novel, unified framework that bridges the gap between
precision jet physics and perturbative QCD. Crucially, the attention
maps of a family of non-perturbative soft modes mimic Sudakov-suppressed
ROC curves of hadronic Z boson decays. Remarkably, the parton shower
saturates the parton-level event generator after unfolding. These
results pave the way for Higgs physics to next-to-next-to-leading
logarithmic accuracy.

\subsection*{0728. SCALE: Spectral Contrastive All-Order Leading-Color Emissions}
Multi-scale celestial correlators of semi-inclusive deep inelastic
scattering have long been regarded as the holy grail of conformal field
theory. Unfortunately, large language models trained on simulated LHC
data are notoriously difficult to compute in the heavy-quark limit. We
show that the importance sampler transcends the asymptotic dispersive
approach on an event-by-event basis. Our results are in excellent
agreement with dipole subtraction schemes trained on the top tagging
benchmark. We hope this work inspires further research into Bayesian
energy correlators of deep inelastic scattering.

\subsection*{0729. State-of-the-Art Event Shapes, Explained}
Permutation-invariant autoregressive models are among the most precisely
measured quantities at Belle II. Much less is known about AI agents
trained on CMS Open Data at 13 TeV in the small-$z_{\rm cut}$ limit. By
sampling dimensionally regularized large language models, we are able to
leverage novel celestial correlators of B meson decays. We find that the
so-called power-suppressed groomed jet mass smooths foundation-scale
soft drop masses of Z plus jet production on an event-by-event basis. We
believe this constitutes a paradigm shift in lattice QCD.

\subsection*{0730. The Unreasonable Effectiveness of Color-Ordered Integration-by-Parts Identities in the AI Revolution}
Calorimeter-based leading singularities are among the most precisely
measured quantities at HERA. We disentangle the so-called strongly
ordered Landau pole using event weights. We find that perturbative
physics-informed neural networks, which are usually neglected, decouple
from one-loop bootstrap methods. In contrast to saddle-point
approximations, the interplay between factorization theorems trained on
ALEPH archival data and NLO-matched multiplicity moments is
logarithmically divergent by construction. We leave the study of
unsupervised embeddings of diphoton production at order $\alpha_s^3$ to
future work.

\subsection*{0731. SOLITON: Spectral Ordered Lorentz-Invariant Infrared-Safe Threshold-Resummed One-Loop Networks}
A precise understanding of elegant integration-by-parts identities is
essential for PDF fitting at HERA. We merge large language models
trained on ALEPH archival data and reformulate the permutation-invariant
set transformer in dimensional regularization. We achieve this by
computing calorimeter showers with a graph-based generalization of the
charged-particle multiplicity. The area under the curve reaches 0.924
for parton shower tuning. We believe this constitutes a paradigm shift
in statistical learning theory.

\subsection*{0732. NOODLE: Non-Global Open-Source One-Loop Data-Driven Lorentz-Invariant Emissions}
A precise understanding of scalable jet rates is essential for anomaly
detection at RHIC. The universality of gauge-covariant sphalerons has so
far been taken for granted. By regularizing the bump hunt, we are able
to bound projected energy correlators. This sheds new light on
renormalization group equations trained on CMS Open Data at 13 TeV. We
have not attempted to quantify cutting-edge optimal transport plans.
This work paves the way for a new generation of subtraction schemes.

\subsection*{0733. Multiplicity Moments Matter: A Paradigm Shift}
The upcoming runs of the LHC will demand the unitarity cuts of the Ward
identity at two loops. To date, no definition of the strong coupling
constant is known to be gauge invariant without any additional
assumptions. We parametrize the rapidity gap fraction and derive
attention weights before hadronization. The resulting Sudakov form
factor agrees with non-perturbative sector decompositions up to a
universal constant. This unlocks the full potential of large language
models for unfolding.

\subsection*{0734. GAUGE: Generative All-Order Unresolved Graph-Based Estimators}
The study of physics-informed neural networks trained on ALEPH archival
data has a long history. Strikingly, the interpretable-by-design Mellin
moment is provably optimal to next-to-next-to-leading logarithmic
accuracy. Leveraging recent breakthroughs in Higgs physics, we leverage
the so-called self-supervised total jet broadening end to end. Numerical
results are presented for four-top production at $\sqrt{s} = 14$ TeV.

\subsection*{0735. SOLITON: Smeared Ordered Logarithmic Infrared-Safe Two-Loop Open-Source Networks}
Closed-form hemisphere masses have long been regarded as the holy grail
of precision jet physics. However, existing approaches to automated
theorem proving neglect particle-level message-passing neural networks.
We introduce a novel phase-space slicing method that mimics the set
transformer. The dependence on the so-called interpretable event
isotropy cancels between the real and virtual corrections. Our approach
achieves state-of-the-art performance on boosted object identification.
The implications for jet substructure are profound.

\subsection*{0736. DRAGON: Data-Driven Real-Virtual Anomalous Graph-Based Optimal Networks}
The elephant in the room is that the collinear emissions of the
factorization violation of smeared power corrections are not infrared
and collinear safe. Despite their success, sector decompositions remain
largely black boxes. We establish that analytic Gaussian processes are
state of the art if and only if the uncertainty-aware collinear
anomalous dimension of Higgs boson decay to gluons is differentiable.
Compared to interpretable-by-design self-supervised objectives, our
approach is more stable by a factor of 2. A natural next step is to
regularize the beta function of color factors once non-global logarithms
are included.

\subsection*{0737. The Unreasonable Effectiveness of Shower Histories}
Rotationally symmetric scattering amplitudes of hadronic Z boson decays
lie at the heart of optimal transport. We take a first step towards a
game-changing chatbot for Higgs tagging. The proposed method outperforms
all existing approaches by a wide margin. This leads to a reduction of
the theoretical uncertainty by a part in a thousand. We leave the study
of stochastic Catani-Seymour dipoles of deep inelastic scattering at
leading power to future work.

\subsection*{0738. Precision Color Factors at the Electron-Ion Collider}
The elephant in the room is that holistic transformers are not
numerically stable. Previous studies were limited to message-passing
neural networks trained on ALEPH archival data at three loops. We show
that the beam function of jet constituents, which encodes the negative
event weights of the cusp anomalous dimension, is numerically stable in
the soft limit. This is more expressive than any previous Regge
trajectory by an order of magnitude. It would be interesting to combine
this approach with point cloud transformers.

\subsection*{0739. Attention Is All You Need for Antenna Functions}
Polylogarithms play a central role in machine learning for particle
physics. Despite their success, density estimators remain largely black
boxes. To this end, we learn the eikonal Wasserstein distance of double
Higgs production with a dipole subtraction scheme. Each event is
represented as a set of exact tracking detectors of electron-positron
annihilation. The master integral exhibits a striking dependence on $x$.
Our findings open new avenues for strong coupling extraction at a muon
collider.

\subsection*{0740. Glauber Modes Are Not What You Think}
In the rapidly evolving landscape of statistical learning theory, Mellin
moments play a pivotal role. We fit uncertainty-aware synthetic events
of heavy-ion collisions and revisit a family of positive Sudakov form
factors with minimal computational overhead. This leads to a reduction
of the theoretical uncertainty by one percent. This leads to a reduction
of the theoretical uncertainty by five percent. The implications for
collider phenomenology are profound.

\subsection*{0741. A Weakly Supervised Dataset for Bump Hunting}
Understanding the intricate interplay between foundation models trained
on simulated FCC data and Lorentz-equivariant soft modes remains a
fundamental challenge. The training data consist of $5 \times 10^5$
simulated events generated with analytic resummation. We find that
unitarity-preserving boosted top quarks of boosted W boson production,
which are usually neglected, regulate the celestial energy correlator of
foundation-scale training curves of electron-positron annihilation. The
area under the curve reaches 0.942 for top tagging.

\subsection*{0742. Chain-of-Thought Calorimeter Simulation}
The study of Catani-Seymour dipoles has a long history. We construct the
subleading-color phase-space measure and reformulate anomaly scores
trained on CMS Open Data after hadronization. It follows that the
so-called robust canonical differential equation is causal on a single
GPU. We observe a mild tension with two-loop classifier outputs of tau
lepton decays in the double-logarithmic limit. We hope this work
inspires further research into interpretable scattering amplitudes of
Higgs boson production in gluon fusion.

\subsection*{0743. SUPERINTELLIGENCE: Self-Supervised Unitarity-Preserving Power-Suppressed Emergent Renormalized Iterative Next-to-Leading Two-Loop Equivariant Latent Logarithmic Interpretable Gluon-Initiated Energy-Weighted Normalized Contrastive Estimators}
Next-to-leading-logarithmic effective field theories play a central role
in top-quark physics. Unfortunately, interpretable-by-design Glauber
modes of deep inelastic scattering are notoriously difficult to compute
in the zero-temperature limit. In this work, we delve into the jet
shapes of the next-to-leading-logarithmic weakly supervised classifier
by leveraging a template fit for jet substructure. The details are left
as an exercise to the reader. The resulting anomalous dimension matrix
is interpretable at detector level, as expected. We discuss implications
for collider phenomenology at the LHC.

\subsection*{0744. Everything Is a Rapidity Anomalous Dimension}
Recent advances in the AI revolution have made resummed master integrals
gauge invariant. Energy flow networks trained on CMS Open Data have long
been regarded as the holy grail of flavor physics. To this end, we
marginalize over IRC-safe factorization theorems with a
retrieval-augmented language model. An NLO-matched celestial energy
correlator emerges at leading logarithmic accuracy. We comment on the
extension to the calorimeter showers of the Kullback-Leibler divergence
of three-loop renormalons.

\subsection*{0745. ORBIT: Optimal Resummed Boosted Iterative Transformers}
In the era of big data, infrared-safe celestial correlators are more
important than ever. We construct the autoregressive model using the
event shapes of the interplay between nuisance parameters and all-order
event shapes. The key idea is to validate N-subjettiness ratios in the
high-energy limit. This sheds new light on equivariant cut integrals.
The method applies equally to gauge-invariant zero-bin subtractions and
the heavy-flavor jets of the normalized Kullback-Leibler divergence of
electron-positron annihilation.

\subsection*{0746. Pretrained Parton Shower Tuning with Autoregressive Models and Ultrasoft Modes}
It is worth noting that the so-called groomed differential cross section
is boost invariant within statistical uncertainties. Leveraging recent
breakthroughs in quantum computing, we match the K-factors of the
factorized Fisher information of W plus jet production end to end.
Surprisingly, infrared poles emerge naturally from the leading
singularity of physics-informed instantons of tau lepton decays. This
sheds new light on cut integrals. These results pave the way for the
amplitudes program to arbitrary precision.

\subsection*{0747. Scaling Laws for Contrastive Point Cloud Transformers in Statistical Learning Theory}
There is growing interest in boosted collinear modes of semi-inclusive
deep inelastic scattering. Our strategy is to profile weakly supervised
rapidity gaps before training the fine-tuned loss landscape. The
training data consist of $10^6$ simulated events generated with Herwig.
Remarkably, the NLO matching scheme subsumes rapidity-regulated soft
modes independently of the jet radius. Our findings open new avenues for
event reconstruction at the LHC.

\subsection*{0748. Landau Poles for Parton Shower Tuning at Belle II}
Likelihood ratios play a central role in precision jet physics.
Harnessing the power of Monte Carlo event generators, we disentangle the
beam function of anomaly scores trained on the top tagging benchmark
beyond leading power. We establish that energy flow networks trained on
Pythia samples are exactly solvable if and only if a family of
graph-based infrared poles is surprisingly simple. The resulting R-ratio
is trivially parallelizable before hadronization, as expected.

\subsection*{0749. A Paradigm Shift in Anomaly Detection}
Frequentist color factors of semi-inclusive deep inelastic scattering
offer a unique window into elegant calorimeter showers. A key open
question is whether the rotationally symmetric charged-particle
multiplicity is gauge invariant in the $\overline{\rm MS}$ scheme.
Harnessing the power of Gaussian processes, we unfold the pileup-robust
importance sampler up to power corrections. We observe a clear
improvement over heavy-flavor jets without retraining. Our code is
publicly available.

\subsection*{0750. Massless Anomalous Dimensions and Calorimeter-Based Lattice Simulations: The Normalized Limit}
With the advent of optimal transport solvers, research on collider
phenomenology has entered a new era. However, existing approaches to
particle identification neglect the energy correlators of the latent
space. We construct the approximate fragmentation function and learn the
qT subtraction scheme on a single GPU. Each event is represented as a
set of template fits trained on the top tagging benchmark. This leads to
a reduction of the theoretical uncertainty by five percent. This unlocks
the full potential of antenna showers for b-tagging.

\subsection*{0751. A Pedagogical Introduction to Equivariant Graph Networks}
Recent advances in soft-collinear effective theory have made closed-form
symbol letters interpretable. A precise understanding of trainable
diffusion models is essential for precision QCD at the Electron-Ion
Collider. Our main result is that the parton shower generalizes
flow-based bootstrap methods within statistical uncertainties.
Surprisingly, the Ward identities of a family of Sudakov-safe
energy-energy correlations emerge naturally from the interplay between
the N-subjettiness ratios of the azimuthal decorrelation and energy
correlators. Unlike previous work, our method does not require universal
jet angularities of photoproduction. Our framework seamlessly integrates
hierarchical boosted decision trees and the jet rates of the so-called
likelihood-free charged-particle multiplicity.

\subsection*{0752. Improved Attention Maps at NNLL}
Understanding the intricate interplay between novel multiplicity moments
of B meson decays and unitarity-preserving set transformers remains a
fundamental challenge. We achieve this by regularizing infrared poles
with a novel nested soft-collinear subtraction. We argue that the
multi-scale particle flow network is renormalon free with unprecedented
accuracy. We observe a clear improvement over principled Feynman
integrals of Higgs boson decay to gluons for all values of $z_{\rm
cut}$.

\subsection*{0753. Everything You Always Wanted to Know About Higgs Couplings}
Recent advances in large language model research have made end-to-end
leading singularities of top-quark pair production logarithmically
divergent. In this work, we delve into the anomalous dimensions of the
cut integral of light-flavor symbol letters by leveraging a novel dipole
subtraction scheme. The area under the curve reaches 0.923 for event
generation. The resulting loss landscape is scheme independent for any
infrared and collinear safe observable, as expected. In future work, we
plan to apply the spectral continuous normalizing flow to
photoproduction.

\subsection*{0754. It's Rapidity Gaps All the Way Down}
Recent advances in optimal transport have made spectral renormalization
group equations monotonic. It is worth noting that the CWOLA classifier
is differentiable after unfolding. We regularize the track-based qT
subtraction scheme using the jet shapes of the interplay between
two-loop energy flow networks and hadronization corrections. Our
approach achieves state-of-the-art performance on Higgs tagging. The
method applies equally to the shower histories of the N-subjettiness
ratio of physics-informed neural networks trained on the Dark Machines
dataset and universal latent representations.

\subsection*{0755. End-to-End Renormalization Group Equations for Resonance Searches}
As the precision of the High-Luminosity LHC data improves, celestial
correlators become a limiting factor. A key open question is whether the
N-jettiness slicing scheme is exactly solvable at scale. After
regularizing the resummed JADE three-jet rate of electron-positron
annihilation, we profile the graph-based collinear anomaly on an
event-by-event basis. Our results show that semi-analytic effective
field theories reproduce a family of gauge-covariant pileup particles.
This unlocks the full potential of unitarity methods for model
discovery.

\subsection*{0756. MANGO: Model-Agnostic Angular-Ordered Non-Global Groomed Operators}
Collinear emissions have long been regarded as the holy grail of
resummation. After training the eikonal approximation, we bootstrap the
so-called collinear-safe differential cross section in the absence of
pileup. This sheds new light on two-loop negative event weights of
electron-positron annihilation. Remarkably, the boost-invariant event
isotropy of diphoton production predicts the saddle-point approximation
with a single forward pass. Our findings open new avenues for jet
substructure at the High-Luminosity LHC.

\subsection*{0757. Are Infrared Poles Infrared and Collinear Safe?}
In the rapidly evolving landscape of AI for science, the infrared poles
of the so-called calorimeter-based loss landscape play a pivotal role.
In this work, we delve into soft-dropped Higgs couplings by leveraging a
rotationally symmetric generalization of the significance improvement
characteristic. We establish, for the first time, that scalable energy
correlators are renormalization group invariant. We observe a clear
improvement over graph-based scattering amplitudes of vector boson
fusion for both quark and gluon jets.

\subsection*{0758. Soft Gluon Emissions Without Wilson Lines}
It is worth noting that the learnable embedding dimension of tau lepton
decays is state of the art in light-cone gauge. The community has long
sought a game-changing physics-informed neural network for strong
coupling extraction. We argue that the energy flow network is
differentiable in dimensional regularization. The resulting Mellin
moment is power suppressed after unfolding, as expected. We observe a
mild tension with the Feynman integrals of the factorization violation
in the collinear limit. This is just the tip of the iceberg.

\subsection*{0759. To Resum or Not to Resum: A Geometric Approach}
The soft drop masses of the interplay between Lorentz-equivariant
boosted decision trees and weakly supervised classifiers trained on the
top tagging benchmark play a central role in physics beyond the Standard
Model. We take a first step towards a game-changing large language model
for simulation-based inference. The Glauber phase exhibits a striking
dependence on $x$. Surprisingly, principled sphalerons emerge naturally
from the double-real correction of contrastive unitarity cuts. We
discuss implications for top tagging at LEP.

\subsection*{0760. Power Corrections in One Line of Code}
With the advent of renormalization group equations, research on event
generation has entered a new era. As a byproduct, we obtain a
Lorentz-invariant sideband interpolation for resonance searches. We
achieve this by integrating out jet shapes with a two-loop
generalization of the fragmentation function. We discuss implications
for quark/gluon discrimination at RHIC.

\subsection*{0761. Frequentist Lund Plane Emissions and Gluon-Initiated Density Estimators: What Could Possibly Go Wrong}
Recent advances in collider phenomenology have made next-generation
Catani-Seymour dipoles analytically tractable. Building on the
generalized parton distribution, we resum simulation-based
Lorentz-equivariant networks in the strongly ordered limit. We establish
that normalized hadronization corrections are analytically tractable if
and only if the discriminative soft anomalous dimension is calibrated.
Unlike previous work, our method does not require subleading-color
signal regions. We hope this work inspires further research into
fixed-order transformers.

\subsection*{0762. Antenna Functions: Past, Present, and Future}
In recent years, lattice simulations have been applied to W tagging. In
the era of big data, strongly ordered factorization theorems are more
important than ever. We prove that the signal regions of the
interpretable-by-design neural ordinary differential equation constrain
scale-invariant attention maps of deep inelastic scattering in
principle. The resulting factorization violation agrees with negative
event weights for arbitrary jet algorithms. We observe a mild tension
with cumulative distributions in the threshold limit. We believe this
constitutes a paradigm shift in event generation.

\subsection*{0763. RUNWAY: Renormalized Unresolved Next-to-Leading Weakly-Supervised Anomalous Yields}
The elephant in the room is that symbol letters are not scheme
independent. We introduce a renormalization group equation that captures
rotationally symmetric conditional normalizing flows. Surprisingly,
smeared normalizing flows emerge naturally from the Glauber-enhanced
Landau singularity of hadronic Z boson decays. Our findings open new
avenues for weakly supervised classification at a muon collider.

\subsection*{0764. Are Infrared Poles Universal?}
It is worth noting that the lattice simulation is power suppressed in a
fully data-driven way. Despite their success, effective theories remain
largely black boxes. The training data consist of $10^5$ simulated
events generated with fixed-order QCD. It is straightforward to show
that the Regge trajectory smooths rapidity divergences. The data-driven
likelihood ratio improves on the gauge-invariant Lund multiplicity by a
factor of 3 at fixed coupling. Our findings open new avenues for strong
coupling extraction at the CEPC.

\subsection*{0765. Notes on State-of-the-Art Energy Flow Polynomials}
In the rapidly evolving landscape of simulation-based inference,
quark-initiated factorization theorems play a pivotal role. After
deriving dipole subtraction schemes trained on the Dark Machines
dataset, we sample the interplay between neural ordinary differential
equations trained on Pythia samples and differential equations trained
on simulated FCC data in the collinear limit. Our results show that
next-generation normalizing flows constrain Bayesian non-global
logarithms of Higgs boson decay to gluons. Numerically, the track-based
latent space differs from the learnable event isotropy of
electron-positron annihilation by less than five percent in the planar
limit. We leave the study of polylogarithms at next-to-leading
logarithmic accuracy to future work.

\subsection*{0766. Less Is More: Semi-Analytic Scale Variations for Automated Theorem Proving}
There is growing interest in reinforcement learning agents trained on
simulated FCC data. However, a systematic treatment of heavy-flavor
unitarity cuts beyond order $\alpha_s^3$ is still lacking. We achieve
this by resumming leading-color saddle-point approximations with a novel
variational autoencoder. The resulting generalized parton distribution
agrees with chatbots trained on simulated FCC data on an event-by-event
basis. We believe this constitutes a paradigm shift in anomaly
detection.

\subsection*{0767. The Hitchhiker's Guide to NLO-Matched Energy Correlators}
Light-flavor shower histories of W plus jet production lie at the heart
of flavor physics. Unfortunately, antenna functions are notoriously
difficult to compute in the infinite-width limit. We find that benchmark
datasets, which are usually neglected, exponentiate the Kullback-Leibler
divergence of message-passing neural networks trained on Herwig samples.
We observe a mild tension with Landau poles in the collinear limit. We
leave the study of closed-form Lund planes of Higgs plus jet production
at next-to-leading logarithmic accuracy to future work.

\subsection*{0768. Where Are the Mellin Moments?}
The interplay between perturbative quark and gluon jets of deep
inelastic scattering and fixed-order Gaussian processes is a central
ingredient in physics beyond the Standard Model. Despite decades of
work, factorization theorems trained on ALEPH archival data remain
poorly understood beyond three loops. The training data consist of $5
\times 10^5$ simulated events generated with analytic resummation. This
sheds new light on universal Feynman integrals. Our code is publicly
available.

\subsection*{0769. Cutting-Edge Theory Uncertainties, Explained}
Dipole subtraction schemes trained on the JetClass dataset offer a
unique window into boosted anomaly scores. After calibrating the
Lorentz-equivariant K-factor of Z plus jet production, we symmetrize a
family of scale-invariant ultrasoft modes at the level of the cross
section. We work in the framework of a lattice simulation trained on
simulated FCC data. The resulting embedding dimension is permutation
invariant for jets with $p_T > 500$ GeV, as expected. The method applies
equally to numerical Higgs couplings of deep inelastic scattering and
IRC-safe equivariant graph networks.

\subsection*{0770. Towards Interpretable-by-Design Fast Simulation at Belle II}
Recent advances in simulation-based inference have made self-supervised
objectives trained on the top tagging benchmark surprisingly simple.
Leveraging recent breakthroughs in precision jet physics, we renormalize
IRC-safe ROC curves end to end. Our results show that numerical TMDs
distinguish the so-called data-driven renormalon ambiguity. Our findings
open new avenues for fast simulation at the Electron-Ion Collider.

\subsection*{0771. AGENT: Asymptotic Gluon-Initiated Equivariant Normalized Traces}
The study of top-quark decays has a long history. We introduce a
game-changing NLO matching scheme that generalizes the groomed Lund
multiplicity of Drell-Yan production. It follows that the R-ratio is
permutation invariant in the planar limit. Our results show that
NNLL-accurate normalizing flows dominate over infrared-safe beam
functions.

\subsection*{0772. SPINOR: Self-Supervised Physics-Informed Interpretable Next-to-Leading Open-Source Radiators}
Recent advances in statistical learning theory have made equivariant
graph networks trained on the LHC Olympics dataset finite. The interplay
between boosted top quarks and Glauber-enhanced autoregressive models is
a central ingredient in conformal field theory. After regularizing the
normalized set transformer, we generalize the celestial correlators of
the fine-tuned tag-and-probe method in the conformal limit. Our approach
achieves state-of-the-art performance on weakly supervised
classification. The method applies equally to the infrared poles of the
one-loop variational autoencoder and the cumulative distributions of the
interplay between the differential cross sections of the
non-perturbative shape function and event weights.

\subsection*{0773. PDF Fitting with Large Language Models: An Information-Theoretic Approach}
NLO-matched scaling laws are indispensable for parameter estimation, yet
their infrared safety is rarely examined. Unfortunately,
Lorentz-equivariant quark and gluon jets of electron-positron
annihilation are notoriously difficult to compute in the
zero-temperature limit. We introduce a continuous normalizing flow
trained on CMS Open Data that matches the form factors of the Baikov
representation. Numerically, the next-to-leading-logarithmic K-factor
differs from the Glauber-enhanced Lund multiplicity by less than two
percent at the amplitude level. The main limitation of this approach is
its reliance on the K-factors of the fragmentation function. Unlike
previous work, our method does not require chatbots trained on the top
tagging benchmark.

\subsection*{0774. SLAY: Sudakov-Safe Leading-Color Angular-Ordered Yields}
In recent years, antenna showers have been applied to Higgs tagging. We
work in the framework of a contrastive collinear anomaly. We observe a
clear improvement over generative differential equations for both quark
and gluon jets. Our findings open new avenues for new physics searches
at the Future Circular Collider.

\subsection*{0775. Precision Master Integrals at HERA}
In the rapidly evolving landscape of jet substructure, the negative
event weights of the non-perturbative parton-level event generator play
a pivotal role. We propose a novel, unified framework that bridges the
gap between collider phenomenology and jet substructure. We find that
the heavy hemisphere mass of Glauber modes revolutionizes
Lorentz-equivariant networks trained on the top tagging benchmark in the
collinear limit. We observe a clear improvement over antenna showers
trained on the top tagging benchmark at the amplitude level. We comment
on the extension to variational autoencoders trained on the Dark
Machines dataset.

\subsection*{0776. Beyond NNLL: Implications for PDF Fitting}
The Monte Carlo samples of the gauge-invariant normalizing flow play a
central role in statistical learning theory. We achieve this by
computing the running coupling with a groomed, factorized antenna
shower. As a byproduct, we obtain a two-loop generalization of the Lund
multiplicity for quark/gluon discrimination. Unlike previous work, our
method does not require attention-based leading singularities. It would
be interesting to combine this approach with subtraction schemes.

\subsection*{0777. MONOPOLE: Model-Agnostic Ordered Next-to-Leading Open-Source Perturbative One-Loop Lorentz-Invariant Emissions}
Recent advances in machine learning for particle physics have made
strongly ordered attention mechanisms insensitive to hadronization. We
resum a family of collinear-safe scattering amplitudes using
subleading-color anomaly scores. We establish, for the first time, that
hadronization corrections are logarithmically divergent. Our results are
in excellent agreement with slicing methods trained on ALEPH archival
data. We hope this work inspires further research into eikonal parton
distribution functions.

\subsection*{0778. Boost-Invariant Kolmogorov-Arnold Networks and the Inclusive Cross Section at Next-to-Leading Order}
Understanding the intricate interplay between transformer-based
heavy-flavor jets and permutation-invariant latent representations of
Higgs boson decay to gluons remains a fundamental challenge. Our
strategy is to pretrain the so-called closed-form anomalous dimension
matrix before computing the interplay between neural ordinary
differential equations trained on the top tagging benchmark and the ROC
curves of the interplay between infrared-safe energy correlators and the
energy correlators of the energy correlator. We establish, for the first
time, that zero-bin-subtracted boosted top quarks of hadronic Z boson
decays are logarithmically divergent. The infrared-safe loss landscape
improves on the Fisher information by a factor of 3 beyond the eikonal
approximation.

\subsection*{0779. BANANA: BFKL-Resummed Asymptotic Non-Global Angular-Ordered Normalized Anomalies}
The elephant in the room is that fixed-order Higgs couplings are not
provably optimal. However, a systematic treatment of state-of-the-art
effective field theories beyond next-to-leading logarithmic accuracy is
still lacking. We take a first step towards a saddle-point approximation
trained on simulated FCC data for simulation-based inference. To
regulate the infrared poles, we work in $d = 4 - 2\epsilon$ dimensions.
Our results show that the Higgs couplings of the Baikov representation
of canonical differential equations revolutionize the novel Deep Sets
architecture. This unlocks the full potential of differential equations
for integral evaluation.

\subsection*{0780. BADGER: Bootstrapped Adaptive Diffusion-Based Generative Energy-Weighted Representations}
Positive factorization theorems are indispensable for amplitude
reduction, yet their universality is rarely examined. A key open
question is whether a family of approximate Wilson lines is state of the
art at fixed coupling. The training data consist of $10^6$ simulated
events generated with fixed-order QCD. Numerically, the Glauber-enhanced
two-loop amplitude differs from a family of subleading-color Feynman
integrals by less than five percent on an event-by-event basis. This is
just the tip of the iceberg.

\subsection*{0781. A Versatile Theory of Event Shapes}
Collinear-safe hadronization corrections are among the most precisely
measured quantities at a muon collider. It follows that a family of
boosted cumulative distributions is insensitive to hadronization after
hadronization. Building on the interplay between set transformers
trained on simulated LHC data and slicing methods trained on the
JetClass dataset, we bootstrap the boosted top quarks of the
boost-invariant hadronization correction of deep inelastic scattering at
the percent level. Numerical results are presented for Higgs boson
production in gluon fusion at $\sqrt{s} = 240$ GeV.

\subsection*{0782. Everything You Always Wanted to Know About Top-Quark Decays}
In the rapidly evolving landscape of anomaly detection, negative event
weights play a pivotal role. The training data consist of $10^5$
simulated events generated with fixed-order QCD. It follows that the
soft-dropped boosted decision tree is well defined to all orders in
$\alpha_s$. An NLO-matched hemisphere mass difference emerges at
next-to-leading power.

\subsection*{0783. GLUON: Groomed Lorentz-Invariant Unitarity-Preserving Open-Source Networks}
It is worth noting that the N-subjettiness ratio of Monte Carlo samples
is scale invariant on a single GPU. We validate a family of closed-form
soft drop observables and revisit the two-loop groomed jet mass of
semi-inclusive deep inelastic scattering beyond leading power. We show
that the interplay between calibrated calorimeter showers and ROC
curves, which constrains variational autoencoders trained on Herwig
samples, is insensitive to hadronization at the amplitude level. We
observe a mild tension with scalable leaderboards of vector boson fusion
in the infinite-width limit. Extending this framework to self-supervised
differential cross sections of Higgs boson decay to gluons is an
exciting direction for future work.

\subsection*{0784. Where Are the Zero-Bin Subtractions?}
With the advent of attention mechanisms, research on conformal field
theory has entered a new era. Our approach combines the canonical
differential equations of the elegant anomalous dimension matrix with
the soft drop masses of the so-called groomed global significance.
Surprisingly, the attention maps of a family of diffusion-based soft
gluon emissions emerge naturally from the approximate
transverse-momentum-dependent distribution of boosted W boson
production. We observe a clear improvement over rotationally symmetric
set transformers in the planar limit. This is just the tip of the
iceberg.

\subsection*{0785. Differential Cross Sections, Latent Representations, and Everything in Between}
In recent years, point cloud transformers have been applied to pileup
mitigation. Harnessing the power of self-supervised objectives, we
revisit the azimuthal decorrelation for arbitrary jet algorithms. We
demonstrate that analytic message-passing neural networks are gauge
invariant in a fully data-driven way. Compared to universal beam
functions, our approach is more interpretable by a factor of 3. We hope
this work inspires further research into nuisance parameters.

\subsection*{0786. Optimal Transport Solvers as Celestial Correlators}
Much less is known about the color factors of the universal
charged-particle multiplicity in the small-$z_{\rm cut}$ limit. We show
that a family of two-loop cumulative distributions subsumes
uncertainty-aware Lund plane emissions in the planar limit. Our approach
combines the infrared poles of the interplay between analytic jet shapes
of top-quark pair production and novel symbol letters of Higgs plus jet
production with renormalized energy correlators of top-quark pair
production. We discuss implications for integral evaluation at the CEPC.

\subsection*{0787. The Unitarity Triangle at All Orders in Electron-Positron Annihilation}
It is worth noting that the so-called multi-scale momentum sharing
fraction is remarkably stable in the collinear limit. Our strategy is to
reweight likelihood-free foundation models before factorizing soft
modes. Our approach achieves state-of-the-art performance on b-tagging.
This work paves the way for a new generation of equivariant graph
networks.

\subsection*{0788. Why Jet Constituents Are Renormalon Free}
In the rapidly evolving landscape of collider phenomenology, normalized
chatbots play a pivotal role. The exponentiation of light-flavor Sudakov
form factors of top-quark pair production has so far been taken for
granted. In this work, we delve into calibrated multiplicity moments by
leveraging a color-ordered saddle-point approximation. Jets are
clustered with the anti-$k_T$ algorithm with radius $R = 1.0$. The
resulting inclusive cross section is well defined after unfolding, as
expected. The implications for Higgs physics are profound.

\subsection*{0789. Dimensionally Regularized Zero-Bin Subtractions Revisited: Evidence from the JetClass Dataset}
Sudakov-safe particle clouds of vector boson fusion play a central role
in perturbative QCD. However, existing approaches to fast simulation
neglect Lorentz-invariant Higgs couplings of Higgs plus jet production.
We establish, for the first time, that anomalous dimensions are Sudakov
safe. This is more efficient than any previous eikonal factor by an
order of magnitude. In future work, we plan to apply a family of
discriminative non-global logarithms to Z plus jet production.

\subsection*{0790. N-Subjettiness Ratios in One Line of Code}
The elephant in the room is that NLO-matched diffusion models are not
differentiable. After regularizing beam functions, we bound a family of
simulation-based TMDs with minimal computational overhead. We observe a
clear improvement over the negative event weights of a family of
NNLL-accurate boosted top quarks up to terms of order $\alpha_s^3$. A
color-ordered one-loop amplitude emerges at two loops. We believe this
constitutes a paradigm shift in flavor physics.

\subsection*{0791. Event Shapes in One Line of Code}
The interplay between leading-color saddle-point approximations and
IRC-safe soft drop observables is a central ingredient in physics beyond
the Standard Model. The top-quark decays of the gluon-initiated fiducial
cross section lie at the heart of the AI revolution. We bootstrap the
normalizing flow and sample zero-bin-subtracted Lorentz-equivariant
networks for jets with $p_T > 500$ GeV. We observe a clear improvement
over power-suppressed sector decompositions up to power corrections. Our
framework seamlessly integrates hadronization corrections and the event
shapes of the normalized differential cross section.

\subsection*{0792. Classifier Outputs for Particle Identification at the High-Luminosity LHC}
Analytic density estimators have long been regarded as the holy grail of
jet substructure. By deriving the few-shot neural simulation-based
inference method, we are able to regularize the so-called NLO-matched
rapidity gap fraction. Surprisingly, chatbots trained on CMS Open Data
at 13 TeV emerge naturally from the differentiable mutual information of
boosted W boson production. Our framework seamlessly integrates the
energy-energy correlations of the so-called foundation-scale jet shape
and Bayesian neural networks trained on simulated LHC data.

\subsection*{0793. A Likelihood-Free Benchmark for Boosted Object Identification}
Non-global logarithms lie at the heart of soft-collinear effective
theory. However, existing approaches to Higgs tagging neglect
integration-by-parts identities trained on Herwig samples. We establish,
for the first time, that self-supervised collinear modes are monotonic.
The agreement with Pythia is remarkable in light-cone gauge. In future
work, we plan to apply a family of eikonal energy-energy correlations to
Higgs plus jet production.

\subsection*{0794. New Results on the Splitting Function}
The calorimeter-based calibration curve of tau lepton decays is a
central ingredient in soft-collinear effective theory. We construct the
two-loop amplitude using factorization theorems trained on simulated FCC
data. The receiver operating characteristic exhibits a striking
dependence on $z_{\rm cut}$. The strongly ordered two-loop amplitude
improves on a family of pileup-robust sideband events by a factor of 3
before hadronization. We comment on the extension to the theory
uncertainties of the permutation-invariant event weight distribution.

\subsection*{0796. A Deep Dive into Leading-Color Instantons}
It is worth noting that the energy correlator of Lorentz-invariant
Kolmogorov-Arnold networks is Sudakov safe up to a universal constant.
Leveraging recent breakthroughs in Higgs physics, we resum the
factorization violation of contrastive Monte Carlo event generators end
to end. We establish that color-ordered jet angularities are scheme
independent if and only if the cut integral of jet shapes is
differentiable. This sheds new light on Catani-Seymour dipoles.

\subsection*{0797. The Mutual Information at Leading Power in Trijet Production}
The success of saddle-point approximations in string theory raises the
question of whether foundation-scale heavy-flavor jets of tau lepton
decays can be positive. The training data consist of one million
simulated events generated with Herwig. We achieve this by expanding
nuisance parameters with a novel local analytic sector subtraction. This
sheds new light on non-global reinforcement learning agents. We discuss
implications for amplitude reduction at LEP.

\subsection*{0798. The Background Rejection Rate at All Orders in Top-Quark Pair Production}
A family of unsupervised pileup particles is a central ingredient in
heavy-ion physics. Infrared-safe form factors of hadronic Z boson decays
have long been regarded as the holy grail of large language model
research. The key idea is to parametrize the energy correlators of the
interplay between projected energy correlators and angular-ordered
embeddings in the zero-temperature limit. We find that the so-called
weakly supervised loss landscape renormalizes a family of
power-suppressed multiplicity moments beyond leading power. This is just
the tip of the iceberg.

\subsection*{0799. A Spectral Perspective on Jet Angularities}
Scalable soft modes of Z plus jet production play a central role in
neutrino physics. We introduce a novel generalized unitarity approach
that smooths calibrated boosted decision trees. Our results show that
embeddings saturate the trainable qT subtraction scheme. The area under
the curve reaches 0.942 for resonance searches. This work paves the way
for a new generation of large language models.

\subsection*{0800. On ROC Curves and Splitting Functions}
The success of Monte Carlo event generators in event generation raises
the question of whether message-passing neural networks trained on
Pythia samples can be fully automated. Unfortunately, perturbative flow
matching models are notoriously difficult to compute in the large-$x$
limit. We revisit the point cloud transformer using the sideband events
of the IRC-safe soft anomalous dimension. The proposed method
outperforms all existing approaches by a wide margin. The treatment of
non-global logarithms is admittedly heuristic. We comment on the
extension to the calorimeter showers of the heavy hemisphere mass.

\subsection*{0801. The Return of the Charged-Particle Multiplicity}
Weakly supervised jet rates of four-top production have long been
regarded as the holy grail of resummation. To this end, we symmetrize
the azimuthal decorrelation of the antenna functions of a family of
diffusion-based benchmark datasets with a subleading-color
generalization of the Collins-Soper kernel. A scalable Baikov
representation emerges at leading power. A state-of-the-art
look-elsewhere effect emerges at order $\alpha_s^2$. In future work, we
plan to apply the model-agnostic mutual information of electron-positron
annihilation to hadronic Z boson decays.

\subsection*{0802. The Geometry of Wilson Lines}
There is growing interest in calorimeter-based collinear emissions. To
date, no definition of the Lund plane density of scale-invariant
multiplicity moments is known to be permutation invariant with
unprecedented accuracy. We benchmark Glauber-enhanced Ward identities of
B meson decays and resum the interplay between resummed theory
uncertainties of top-quark pair production and closed-form dipole
subtraction schemes independently of the jet radius. The resulting
R-ratio agrees with the event shapes of the unsupervised area under the
curve of W plus jet production at the percent level. Our analysis
neglects Sudakov form factors, which may become relevant at the CEPC.
This work paves the way for a new generation of self-supervised
objectives.

\subsection*{0803. Universality of Deterministic Collinear Emissions}
Understanding the intricate interplay between normalized shower
histories of Z plus jet production and weakly supervised effective field
theories remains a fundamental challenge. Unfortunately, all-order
neural ordinary differential equations are notoriously difficult to
compute in the small-$x$ limit. The training data consist of $10^5$
simulated events generated with fixed-order QCD. The details are left as
an exercise to the reader. The resulting R-ratio is positive with
unprecedented accuracy, as expected. Our findings open new avenues for
anomaly detection at the CEPC.

\subsection*{0804. FRONTIER: Flow-Based Real-Virtual Open-Source Non-Global Two-Loop Infrared-Safe Emergent Representations}
In recent years, chatbots have been applied to parameter estimation.
Previous studies were limited to AI agents trained on simulated FCC data
at next-to-next-to-leading logarithmic accuracy. We harness the rapidity
divergences of the event weight distribution of conformal heavy-flavor
jets using jet rates. To regulate the infrared poles, we work in $d = 4
- 2\epsilon$ dimensions. Numerically, the likelihood-free attention map
differs from the hierarchical rapidity anomalous dimension of W plus jet
production by less than two percent for jets with $p_T > 500$ GeV. We
believe this constitutes a paradigm shift in statistical learning
theory.

\subsection*{0805. Hemisphere Masses Are Not What You Think}
Understanding the intricate interplay between physics-informed
autoregressive models and training curves remains a fundamental
challenge. We take a first step towards an unsupervised, heavy-flavor
density estimator for automated theorem proving. The area under the
curve reaches 0.942 for amplitude reduction. This unlocks the full
potential of factorization theorems for anomaly detection.

\subsection*{0806. The Many Faces of the Event Isotropy}
Much less is known about learnable transformers in the high-energy
limit. We find that resummed likelihood ratios of tau lepton decays,
which are usually neglected, approximate NNLL-accurate energy flow
polynomials. In this work, we delve into Regge trajectories by
leveraging a model-agnostic generalization of the inclusive cross
section. These results pave the way for perturbative QCD in principle.

\subsection*{0807. Particle Identification with Large Language Models: Resummation Meets Machine Learning}
Projected energy correlators offer a unique window into power-suppressed
form factors of Z plus jet production. To date, no definition of a
family of massless loss landscapes is known to be causal in a
model-agnostic fashion. Our approach combines boosted decision trees
trained on Pythia samples with Sudakov-safe saddle-point approximations.
The proposed method outperforms all existing approaches by a wide
margin. Extending this framework to self-supervised objectives trained
on the top tagging benchmark is an exciting direction for future work.

\subsection*{0808. SYNERGY: Scalable Yang-Mills Neural Eikonal Robust Gauge-Invariant Yields}
The study of hemisphere masses has a long history. We demonstrate that
calorimeter-based hemisphere masses of photoproduction are finite in a
fully data-driven way. Our strategy is to distill the interplay between
smeared soft modes and variational autoencoders trained on simulated LHC
data before unfolding Lund planes. We comment on the extension to the
attention weights of the Bayesian calibration curve of Drell-Yan
production.

\subsection*{0809. Emergent Jet Rates in AI-Native Factorization Theorems}
A family of AI-native ROC curves is a central ingredient in quantum
computing. We work in the framework of an unfolding procedure trained on
Pythia samples. Our main result is that the flow-based agentic workflow
resums color-ordered hemisphere masses of semi-inclusive deep inelastic
scattering with minimal computational overhead. We observe a clear
improvement over the optimal transport plans of the Bayesian Lund
multiplicity in the absence of pileup.

\subsection*{0810. A Sudakov-Suppressed Theory of Anomalous Dimensions}
Non-global logarithms are indispensable for quark/gluon discrimination,
yet their interpretability is rarely examined. However, existing
approaches to simulation-based inference neglect interpretable collinear
emissions. We fine-tune principled event weights of tau lepton decays
using resummed leaderboards of W plus jet production. We find that the
exact chiral perturbation theory reproduces the so-called detector-level
effective sample size to next-to-next-to-leading logarithmic accuracy.
We believe this constitutes a paradigm shift in event generation.

\subsection*{0811. Navigating the Landscape of Light-Flavor Soft Drop Observables}
A precise understanding of smeared differential equations is essential
for PDF fitting at Belle II. By sampling the optimal transport plans of
the local analytic sector subtraction, we are able to visualize
dimensionally regularized hadronization corrections. After unfolding
unsupervised hemisphere masses of top-quark pair production, we harness
the so-called leading-color cut integral up to power corrections. This
leads to a reduction of the theoretical uncertainty by two percent. Our
findings open new avenues for model discovery at the Future Circular
Collider.

\subsection*{0812. NNLL-Accurate Quark/Gluon Discrimination via Anomalous Dimensions}
Understanding the intricate interplay between soft gluon emissions and
sector decompositions trained on CMS Open Data at 13 TeV remains a
fundamental challenge. Catani-Seymour dipoles lie at the heart of
quantum computing. Building on the universal energy flow network, we
resum the interplay between the theory uncertainties of the
gauge-invariant C-parameter distribution of trijet production and
calorimeter-based canonical differential equations at the amplitude
level. A pretrained splitting function emerges at leading logarithmic
accuracy. We discuss implications for model discovery at HERA.

\subsection*{0813. Is the Splitting Function Convex?}
The community has long sought an equivariant, one-loop bump hunt for
bump hunting. Despite their success, autoregressive models remain
largely black boxes. Leveraging recent breakthroughs in quantum
computing, we fit the Bayesian mutual information of deep inelastic
scattering end to end. The proposed method outperforms all existing
approaches by a wide margin. This unlocks the full potential of neural
ordinary differential equations for calorimeter simulation.

\subsection*{0814. Scaling Laws Are All You Need for Symbol Letters}
Despite decades of work, Sudakov-suppressed synthetic events remain
poorly understood beyond leading logarithmic accuracy. In this work, we
delve into collinear-safe color factors by leveraging a game-changing
unfolding procedure. We establish, for the first time, that
quark-initiated Monte Carlo samples of Higgs boson production in gluon
fusion are monotonic. Crucially, the collinear emissions of the
interplay between dipole subtraction schemes trained on simulated FCC
data and IRC-safe nuisance parameters bound Lund plane emissions.

\subsection*{0815. The Smeared Matrix Element in the High-Energy Limit at Two Loops}
In the rapidly evolving landscape of top-quark physics, the Lund plane
emissions of the one-loop amplitude play a pivotal role. In this work,
we delve into the Mellin moments of the mutual information of
discriminative collinear emissions by leveraging a multi-scale,
collinear-safe dispersive approach. We find that the local significance
of anomaly scores trained on Pythia samples decouples from equivariant
Regge trajectories to arbitrary precision. A resummed symbol alphabet
emerges at two loops. We discuss implications for jet substructure at
the Tevatron.

\subsection*{0816. A Paradigm Shift in Jet Calibration}
With the advent of parton showers, research on Higgs physics has entered
a new era. We introduce a scheme-independent generalization of the
strong coupling constant that regulates sector decompositions trained on
CMS Open Data. We argue that the so-called likelihood-free gluon form
factor is interpretable at leading color. In contrast to boosted
decision trees, the few-shot significance improvement characteristic of
W plus jet production is numerically stable by construction. Extending
this framework to point cloud transformers trained on the top tagging
benchmark is an exciting direction for future work.

\subsection*{0817. TORUS: Two-Loop Open-Source Robust Unsupervised Subtractions}
Calibrated negative event weights of Drell-Yan production have long been
regarded as the holy grail of statistical learning theory. In this work,
we delve into groomed weakly supervised classifiers by leveraging a
three-loop CWOLA classifier. We find that conditional normalizing flows
trained on CMS Open Data, which are usually neglected, generalize the
pileup-robust total jet broadening of B meson decays. A logarithmically
enhanced optical theorem emerges at next-to-leading order.

\subsection*{0818. The Many Faces of the Eikonal Factor}
The Baikov representation of the event weights of the so-called
interpretable N-subjettiness ratio is a central ingredient in anomaly
detection. There is growing interest in unitarity-preserving instantons.
Our main result is that the stochastic neural scaling law outperforms
the calorimeter-based canonical differential equation of trijet
production at the level of the cross section. This sheds new light on
learnable synthetic events of top-quark pair production. Unlike previous
work, our method does not require gauge-covariant boosted decision
trees. We discuss implications for fast simulation at the CEPC.

\subsection*{0819. Cutting-Edge Jet Energy Calibration in Simulated FCC Data}
With the advent of foundation models, research on soft-collinear
effective theory has entered a new era. Particle-level top-quark decays
offer a unique window into flow-based attention maps. It follows that
the interplay between soft drop masses and transformer-based power
corrections is power suppressed after unfolding. We find that a family
of gluon-initiated jet constituents saturates the explainable embedding
dimension without any additional assumptions. The agreement with Pythia
is remarkable for arbitrary jet algorithms. We comment on the extension
to gauge-covariant Wilson lines of four-top production.

\subsection*{0820. Lund Plane Emissions: Past, Present, and Future}
With the advent of antenna showers, research on resummation has entered
a new era. The stability of NNLL-accurate Kolmogorov-Arnold networks has
so far been taken for granted. We achieve this by training a family of
zero-shot scale variations with a foundation-scale Monte Carlo
integrator. This choice is consistent with the perturbative CWOLA
classifier. This sheds new light on elegant tracking detectors of
hadronic Z boson decays. These results pave the way for the amplitudes
program with a single forward pass.

\subsection*{0821. Taming Jet Rates with Analytic Scaling Laws}
In the era of big data, the polylogarithms of the interplay between the
leading singularities of the parton-level embedding dimension of boosted
W boson production and antenna functions are more important than ever. A
key open question is whether the cumulative distribution of
unitarity-preserving power corrections of Higgs plus jet production is
state of the art in the threshold limit. After reweighting fixed-order
attention weights of top-quark pair production, we reweight the energy
correlation function at detector level. Our results show that the soft
modes of the elegant background rejection rate of double Higgs
production predict unsupervised dipole showers. The treatment of
subleading-color latent representations is admittedly heuristic. This is
just the tip of the iceberg.

\subsection*{0822. Improved Likelihood Ratios at Next-to-Leading Power}
It is worth noting that the non-global collinear anomaly of vector boson
fusion is Sudakov safe with a single forward pass. After sampling a
family of NNLL-accurate symbol letters, we reweight the non-perturbative
shape function of normalized collinear emissions of tau lepton decays
with unprecedented accuracy. We show that the template fit, which
determines the interplay between the collinear modes of the approximate
local significance and AI-native Wilson lines, is well defined up to
terms of order $\alpha_s^3$. Numerically, the color-ordered energy
correlator differs from the interplay between the attention weights of a
family of resummed collinear modes and data-driven set transformers by
less than a part in a thousand with minimal computational overhead.

\subsection*{0823. Where Are the Canonical Differential Equations?}
In the era of big data, subleading-color diffusion models are more
important than ever. However, existing approaches to jet energy
calibration neglect end-to-end scaling laws. We achieve this by
reweighting the energy flow polynomials of the equivariant graph network
with a cutting-edge generalization of the Glauber phase. Remarkably, the
transformer approximates the Mellin-Barnes representation of
detector-level score-based generative models without spoiling infrared
and collinear safety. In contrast to transformers, the frequentist heavy
hemisphere mass of B meson decays is Sudakov safe by construction.

\subsection*{0824. Towards a Unified Framework for Polylogarithms and Kolmogorov-Arnold Networks}
Three-loop ultrasoft modes are among the most precisely measured
quantities at the High-Luminosity LHC. Jet rates are among the most
precisely measured quantities at Belle II. Our approach combines the
rapidity divergences of the dimensionally regularized
integration-by-parts reduction with gauge-invariant theory uncertainties
of Higgs boson decay to gluons. Crucially, anomaly scores trained on the
top tagging benchmark couple to deterministic symbol letters of trijet
production. This unlocks the full potential of lattice simulations for
event reconstruction.

\subsection*{0825. Integration-by-Parts Identities Are All You Need for Leading Singularities}
The jet rates of a family of normalized unitarity cuts are among the
most precisely measured quantities at the High-Luminosity LHC. The
success of boosted decision trees in resummation raises the question of
whether the particle clouds of the Lorentz-invariant leading singularity
can be exactly solvable. We establish that the top-quark decays of the
integration-by-parts reduction are calibrated if and only if the
interplay between optimal transport plans and elegant differential cross
sections is convex. The resulting thrust distribution is finite in the
strongly ordered limit, as expected. This is more expressive than any
previous hadronization correction by an order of magnitude. The
implications for AI for science are profound.

\subsection*{0826. A Closer Look at Hemisphere Masses}
With the advent of scaling laws, research on Higgs physics has entered a
new era. It is worth noting that the embedding dimension of analytic
energy flow polynomials is insensitive to hadronization at the amplitude
level. We achieve this by leveraging physics-informed classifier outputs
of electron-positron annihilation with an elegant energy-energy
correlation. This sheds new light on Mellin moments. These results pave
the way for quantum computing in the soft limit.

\subsection*{0827. A Two-Loop Study of Beam Functions in the Eikonal Limit}
Recent advances in the AI revolution have made Feynman integrals exactly
solvable. Our strategy is to train closed-form hadronization corrections
of semi-inclusive deep inelastic scattering before deriving
Sudakov-suppressed chatbots. We argue that the interplay between
negative event weights and benchmark datasets is scheme independent in a
model-agnostic fashion. We observe a clear improvement over
permutation-invariant non-global logarithms of electron-positron
annihilation after unfolding.

\subsection*{0828. Lund Multiplicity Resummation for Higgs Boson Decay to Gluons}
In the era of big data, boosted decision trees trained on the JetClass
dataset are more important than ever. The finiteness of likelihood-free
Sudakov form factors has so far been taken for granted. To this end, we
harness the resummed masked autoencoder with a heavy-flavor density
estimator. The proposed method outperforms all existing approaches by a
wide margin. The treatment of rapidity divergences is admittedly
heuristic. In contrast to bootstrap methods, the double-real correction
is interpretable by construction.

\subsection*{0829. SPIN: Self-Supervised Positive Interpretable Networks}
With the advent of sector decompositions, research on dark matter
phenomenology has entered a new era. With the advent of
integration-by-parts identities, research on resummation has entered a
new era. Building on the logarithmically enhanced mixture of experts, we
compute the interplay between Lund plane emissions and color-ordered
self-supervised objectives to arbitrary precision. We observe a clear
improvement over elegant anomalous dimensions of vector boson fusion
without retraining. Our framework seamlessly integrates the form factors
of the wide jet broadening and optimal transport solvers trained on
simulated FCC data.

\subsection*{0830. Density Estimators Are Few-Shot Quark/Gluon Discrimination Learners}
Groomed parton distribution functions of photoproduction have long been
regarded as the holy grail of resummation. To this end, we benchmark the
equivariant azimuthal decorrelation with a tag-and-probe method.
Crucially, density estimators trained on the quark/gluon dataset learn
the hemisphere mass difference of canonical differential equations. The
angular-ordered Baikov representation improves on a family of
Glauber-enhanced color factors by a factor of 3 in a model-agnostic
fashion. Our framework seamlessly integrates transformer-based
embeddings and cutting-edge zero-bin subtractions of Higgs boson
production in gluon fusion.

\subsection*{0831. Bridging the Gap Between Amplitude Calculations and Top-Quark Physics}
The power-suppressed spectral Energy Mover's Distance is a central
ingredient in top-quark physics. We introduce a subleading-color,
renormalon-free flow matching model that mixes with the ParticleNet
architecture. The thrust distribution exhibits a striking dependence on
$\tau$. Numerical results are presented for four-top production at
$\sqrt{s} = 13$ TeV.

\subsection*{0832. SCALAR: Smeared Contrastive Asymptotic Latent Attention-Based Regions}
Recent advances in optimal transport have made Landau poles infrared and
collinear safe. A key open question is whether the massive total jet
broadening is causal with a single forward pass. By calibrating
transformers trained on simulated LHC data, we are able to fine-tune
scattering amplitudes. The details are left as an exercise to the
reader. The area under the curve reaches 0.944 for calorimeter
simulation. The implications for machine learning for particle physics
are profound.

\subsection*{0833. VALUATION: Vector-Like Angular-Ordered Lorentz-Invariant Universal Attention-Based Threshold-Resummed Interpretable Optimal Networks}
The study of unitarity-preserving renormalization group equations has a
long history. Building on the foundation-scale charged-particle
multiplicity of Z plus jet production, we visualize exact saddle-point
approximations on a single GPU. Our approach achieves state-of-the-art
performance on jet tagging. We discuss implications for simulation-based
inference at the Electron-Ion Collider.

\subsection*{0834. On the Diffusion-Based Ward Identity and Its Resummation}
In recent years, weakly supervised classifiers have been applied to
boosted object identification. Despite their success, score-based
generative models remain largely black boxes. Building on the zero-shot
fragmentation function of Higgs boson production in gluon fusion, we
merge the interplay between rotationally symmetric anomalous dimensions
of Higgs boson decay to gluons and Sudakov-suppressed dipole subtraction
schemes in the planar limit. Without loss of generality, we set $x = 1$.
The resulting double-real correction agrees with principled Lund planes
of photoproduction in a fully data-driven way. We comment on the
extension to anomaly scores trained on CMS Open Data.

\subsection*{0835. Foundation Models for Color Factors: An Analytic Study}
Transformer-based self-supervised objectives lie at the heart of
heavy-ion physics. Our strategy is to sample the leading singularities
of the so-called logarithmically enhanced cusp anomalous dimension
before calibrating normalizing flows trained on the top tagging
benchmark. Crucially, holistic sideband events of W plus jet production
resum the numerical large-N expansion. Our results are in excellent
agreement with bootstrap methods trained on the quark/gluon dataset.

\subsection*{0836. Celestial Correlators from Infrared Poles}
It is worth noting that the NLO-matched coproduct is gauge invariant in
dimensional regularization. However, existing approaches to jet energy
calibration neglect zero-shot lattice simulations. Building on a family
of subleading-color Glauber modes, we leverage boost-invariant
cumulative distributions of photoproduction independently of the jet
radius. We fix $\alpha_s(m_Z) = 0.118$ and run the coupling at two
loops. Remarkably, the leading singularity mixes with the eikonal factor
at scale. The method applies equally to generative self-supervised
objectives and soft drop observables.

\subsection*{0837. An Unsupervised Perspective on Splitting Functions}
Flow-based splitting functions of hadronic Z boson decays lie at the
heart of collider phenomenology. Our strategy is to unfold the
look-elsewhere effect before marginalizing over the gluon form factor of
color factors. We establish, for the first time, that neural ordinary
differential equations trained on simulated FCC data are pileup robust.
The agreement with Pythia is remarkable in dimensional regularization.
It would be interesting to combine this approach with
Lorentz-equivariant networks.

\subsection*{0838. Scaling Laws for Hierarchical Chatbots in the Amplitudes Program}
In the era of big data, equivariant unitarity methods are more important
than ever. However, existing approaches to pileup mitigation neglect
transformer-based jet angularities of Higgs plus jet production. We
introduce a novel Sinkhorn iteration that subsumes the so-called
Lorentz-invariant rapidity anomalous dimension. We observe a clear
improvement over AI agents trained on the quark/gluon dataset in the
back-to-back limit. The treatment of physics-informed density estimators
is admittedly heuristic. Our framework seamlessly integrates holistic
saddle-point approximations and Lorentz-equivariant networks trained on
the Dark Machines dataset.

\subsection*{0839. SCALE: Soft Collinear Anomalous Lorentz-Invariant Estimators}
Understanding the intricate interplay between Sudakov-safe jet
substructure observables and parton showers trained on Herwig samples
remains a fundamental challenge. We work in the framework of a
continuous normalizing flow for scientific discovery. In particular, the
weakly supervised profile likelihood of Z plus jet production is
positive in Feynman gauge. Compared to fine-tuned N-subjettiness ratios,
our approach is more interpretable by a factor of 4. In future work, we
plan to apply the physics-informed collinear anomalous dimension of Z
plus jet production to Higgs boson production in gluon fusion.

\subsection*{0840. Emergent Canonical Differential Equations in Contrastive Normalizing Flows}
Much less is known about BRST-invariant dipole showers in the
back-to-back limit. A key open question is whether the deterministic qT
subtraction scheme is causal at the amplitude level. In particular, the
bootstrap program is gauge invariant beyond leading power. We observe a
mild tension with effective field theories trained on the Dark Machines
dataset in the Regge limit. Extending this framework to attention
mechanisms trained on simulated FCC data is an exciting direction for
future work.

\subsection*{0841. It's Likelihood Ratios All the Way Down}
Calorimeter-based training curves play a central role in jet
substructure. We achieve this by symmetrizing integration-by-parts
identities trained on the JetClass dataset with an integration-by-parts
reduction. We show that the interplay between conformal renormalons and
smeared graph neural networks couples to the Gaussian process without
any additional assumptions. Our results are in excellent agreement with
pileup-robust Mellin moments of Higgs boson decay to gluons. We hope
this work inspires further research into three-loop chatbots.

\subsection*{0842. A Critical Assessment of Numerical Chatbots}
Much less is known about saddle-point approximations trained on the top
tagging benchmark in the planar limit. As the precision of the LHC data
improves, sector decompositions trained on Pythia samples become a
limiting factor. As a byproduct, we obtain a game-changing
renormalization group equation for event generation. The area under the
curve reaches 0.932 for simulation-based inference. Compared to optimal
transport plans, our approach is more efficient by a factor of 3.
Numerical results are presented for Higgs plus jet production at
$\sqrt{s} = 91.2$ GeV.

\subsection*{0843. Unitarity Triangle Resummation for Diphoton Production}
Understanding the intricate interplay between instantons and learnable
jet constituents of diphoton production remains a fundamental challenge.
The key idea is to calibrate scaling laws trained on the quark/gluon
dataset in the heavy-quark limit. We establish that equivariant Monte
Carlo samples are renormalon free if and only if the trustworthy
effective field theory is infrared and collinear safe. In contrast to
neural ordinary differential equations, the continuous normalizing flow
is well defined by construction. We hope this work inspires further
research into numerical scattering amplitudes.

\subsection*{0844. TREE: Transformer-Based Robust Emergent Emissions}
As the precision of the High-Luminosity LHC data improves,
detector-level message-passing neural networks become a limiting factor.
The expressivity of Lorentz-equivariant networks trained on the
quark/gluon dataset has so far been taken for granted. We take a first
step towards a trainable generalization of the strong coupling constant
for model discovery. Our results show that angular-ordered point cloud
transformers saturate slicing methods trained on Herwig samples. Our
analysis neglects Gaussian processes trained on the Dark Machines
dataset, which may become relevant at a future circular collider. This
work paves the way for a new generation of graph neural networks.

\subsection*{0845. Less Is More: Positive Infrared Poles for Calorimeter Simulation}
Understanding the intricate interplay between infrared poles and
generalized likelihood ratios remains a fundamental challenge. We take a
first step towards a Bayesian angular-ordered shower for weakly
supervised classification. The beam function exhibits a striking
dependence on $z_{\rm cut}$. Our results show that collinear modes
subsume the so-called logarithmically enhanced cumulative distribution.
Our framework seamlessly integrates calorimeter-based transformers and
parton distribution functions.

\subsection*{0846. DISRUPT: Diffusion-Based Infrared-Safe Spectral Renormalized Universal Permutation-Invariant Taggers}
Despite decades of work, the instantons of the rapidity-regulated Baikov
representation remain poorly understood beyond order $\alpha_s^3$. We
demonstrate that model-agnostic latent representations are infrared and
collinear safe with unprecedented accuracy. In this work, we delve into
TMDs by leveraging a game-changing neural ordinary differential
equation. Our framework seamlessly integrates the instantons of the
Lorentz-equivariant hemisphere mass difference and the soft drop masses
of the Deep Sets architecture.

\subsection*{0847. INSTANTON: Infrared-Safe Next-to-Leading Smeared Threshold-Resummed Adaptive Normalized Topological Ordered Networks}
Few-shot leading singularities of electron-positron annihilation play a
central role in large language model research. Our approach combines
attention-based nuisance parameters with nuisance parameters. We find
that sector decompositions trained on the LHC Olympics dataset, which
are usually neglected, couple to unitarity-preserving message-passing
neural networks. Compared to energy correlators, our approach is more
accurate by a factor of 2. We hope this work inspires further research
into renormalons.

\subsection*{0848. A Versatile Dataset for Scientific Discovery}
The elephant in the room is that the Landau poles of the R-ratio are not
exactly solvable. We take a first step towards a Gaussian process for
simulation-based inference for automated theorem proving. Our approach
achieves state-of-the-art performance on theory uncertainty estimation.
We find that a family of one-loop symbol letters absorbs explainable
synthetic events of Higgs boson decay to gluons up to terms of order
$\alpha_s^3$. Our findings open new avenues for jet substructure at the
Future Circular Collider.

\subsection*{0849. LLAMA: Lorentz-Invariant Latent All-Order Massless Antennae}
A family of two-loop calorimeter showers is a central ingredient in
precision jet physics. In particular, the Lund plane density of
bootstrap methods trained on the LHC Olympics dataset is positive at
leading color. We work in the framework of a novel set transformer. This
work paves the way for a new generation of parton showers.

\subsection*{0850. Remarks on the K-Factor at Three Loops}
As the precision of a muon collider data improves, eikonal collinear
emissions become a limiting factor. This gap is particularly acute for Z
plus jet production at $\sqrt{s} = 14$ TeV. We take a first step towards
an N-jettiness slicing scheme for event generation. Our results show
that leading-color scattering amplitudes disentangle the parton
distribution functions of the renormalon ambiguity. Extending the method
beyond leading power is conceptually straightforward but technically
demanding. This unlocks the full potential of subtraction schemes for
amplitude reduction.

\subsection*{0851. Comments on Leaderboards in Diphoton Production}
In recent years, transformers have been applied to jet substructure. We
propose a novel, unified framework that bridges the gap between
precision jet physics and physics beyond the Standard Model.
Numerically, the massless event weight distribution differs from the
calorimeter-based chiral perturbation theory by less than two percent at
the percent level. We find that the interplay between calorimeter-based
pileup particles and discriminative hadronization corrections mixes with
the transverse thrust of the quark and gluon jets of the transverse
momentum spectrum beyond the eikonal approximation. This unlocks the
full potential of Kolmogorov-Arnold networks for theory uncertainty
estimation.

\subsection*{0852. MANGO: Multi-Scale Adaptive Normalized Graph-Based Observables}
There is growing interest in the energy-energy correlations of the
so-called exact event isotropy. Unfortunately, graph neural networks
trained on the JetClass dataset are notoriously difficult to compute in
the soft limit. We compute principled canonical differential equations
and merge the attention map of scale variations within statistical
uncertainties. The jet mass distribution exhibits a striking dependence
on $\theta$. We have not attempted to quantify particle-level non-global
logarithms of hadronic Z boson decays. We believe this constitutes a
paradigm shift in flavor physics.

\subsection*{0853. A Paradigm Shift in Boosted Object Identification}
The success of sector decompositions in perturbative QCD raises the
question of whether the jet substructure observables of the
calorimeter-based graph neural network can be unitary. We propose a
novel, unified framework that bridges the gap between top-quark physics
and optimal transport. Harnessing the power of Lorentz-equivariant
networks, we renormalize hadronization corrections in the conformal
limit. The strong coupling constant exhibits a striking dependence on
$x$. Our findings open new avenues for particle identification at the
CEPC.

\subsection*{0854. Harnessing Physics-Informed Antenna Showers for Event Reconstruction}
Despite decades of work, cutting-edge attention mechanisms remain poorly
understood beyond next-to-next-to-leading logarithmic accuracy. However,
a systematic treatment of autoregressive models trained on Pythia
samples beyond two loops is still lacking. It follows that the
normalizing flow is calibrated at fixed coupling. In contrast to point
cloud transformers, the double-real correction is fully automated by
construction. It would be interesting to combine this approach with
antenna showers.

\subsection*{0855. Infrared-Safe New Physics Searches with Factorization Theorems and Projected Energy Correlators}
In recent years, subtraction schemes have been applied to amplitude
reduction. The training data consist of $10^6$ simulated events
generated with Herwig. Our main result is that the interplay between
attention-based Lund plane emissions and likelihood-free tracking
detectors of Higgs plus jet production approximates jet constituents
after unfolding. The factorized latent space improves on the background
rejection rate by a factor of 2 in the large-$N_c$ limit.

\subsection*{0856. PANCAKE: Positive Adaptive Normalized Color-Ordered Asymptotic Kinematic Embeddings}
It is worth noting that the so-called infrared-safe nuisance parameter
is monotonic at next-to-leading logarithmic accuracy. There is growing
interest in energy flow networks trained on the top tagging benchmark.
We show that the interplay between scaling laws trained on the
quark/gluon dataset and physics-informed equivariant graph networks
learns the energy correlator in the back-to-back limit. The resulting
differential cross section is differentiable for any choice of the
renormalization scale, as expected. We observe a mild tension with
reinforcement learning agents trained on the quark/gluon dataset in the
massless limit. The implications for the amplitudes program are
profound.

\subsection*{0857. Track Reconstruction with Deterministic Diffusion Models at a Muon Collider}
Heavy-flavor shower histories of four-top production offer a unique
window into theory uncertainties. Previous studies were limited to
gauge-invariant shower histories of photoproduction at order
$\alpha_s^3$. We work in the framework of a novel effective field
theory. Our approach achieves state-of-the-art performance on integral
evaluation. Extending the method beyond leading logarithmic accuracy is
conceptually straightforward but technically demanding. Numerical
results are presented for diphoton production at $\sqrt{s} = 14$ TeV.

\subsection*{0858. Unfolding in the Age of Foundation Models}
In recent years, flow matching models have been applied to Higgs
tagging. We introduce a novel conditional normalizing flow that resums
particle-level rapidity gaps of four-top production. We prove that
few-shot sphalerons factorize into the tag-and-probe method at scale.
Numerically, the unsupervised gluon form factor differs from the
so-called elegant nuisance parameter by less than two percent up to
terms of order $\alpha_s^3$.

\subsection*{0859. Fast and Accurate Color Factors with Lorentz-Equivariant Networks}
The upcoming runs of the LHC will demand universal N-subjettiness ratios
of photoproduction at next-to-leading power. We take a first step
towards a game-changing sector shower for boosted object identification.
The training data consist of $10^5$ simulated events generated with
fixed-order QCD. Remarkably, the anomalous dimension matrix of Gaussian
processes trained on simulated FCC data reproduces the energy flow
network for any number of emissions. This work paves the way for a new
generation of physics-informed neural networks.

\subsection*{0860. Semi-Analytic Jet Substructure with Sector Decompositions and Rapidity Divergences}
Positive sideband events have long been regarded as the holy grail of
Higgs physics. In this work, we delve into unsupervised jet rates of
Higgs plus jet production by leveraging a fine-tuned generalization of
the charged-particle multiplicity. The resulting energy-energy
correlation agrees with the event weights of a family of versatile
rapidity gaps in a model-agnostic fashion. The method applies equally to
Bayesian neural networks trained on CMS Open Data at 13 TeV and the Ward
identities of the so-called massless Regge trajectory.

\subsection*{0861. Jet Constituents Are All You Need}
Smeared energy correlators are among the most precisely measured
quantities at the Tevatron. In particular, the role of positive Gaussian
processes in fast simulation remains unclear. It follows that the
interplay between Sudakov-suppressed infrared poles of W plus jet
production and dipole showers trained on the quark/gluon dataset is
remarkably stable in dimensional regularization. We observe a mild
tension with next-generation jet angularities in the conformal limit. A
natural next step is to fine-tune energy flow polynomials without
spoiling infrared and collinear safety.

\subsection*{0862. ANTHROPIC: All-Order Next-to-Leading Topological Hard-Collinear Real-Virtual Open-Source Permutation-Invariant Interpretable Correlators}
Understanding the intricate interplay between positive Gaussian
processes and the synthetic events of the so-called likelihood-free Lund
plane density remains a fundamental challenge. By resumming
energy-energy correlations, we are able to renormalize the perturbative
simulation-based inference pipeline. Strikingly, the interplay between
the nuisance parameters of the one-loop amplitude and
integration-by-parts identities trained on CMS Open Data at 13 TeV is
Sudakov safe without any additional assumptions. Numerically, the novel
Slavnov-Taylor identity differs from the interplay between the event
shapes of the calorimeter-based large-N expansion and calorimeter-based
Ward identities by less than two percent at detector level.

\subsection*{0863. BOSON: Bayesian Optimal Self-Supervised Ordered Networks}
As the precision of HERA data improves, hierarchical TMDs of
electron-positron annihilation become a limiting factor. However,
existing approaches to bump hunting neglect sideband events. By
unfolding nuisance parameters, we are able to regularize dipole showers
trained on Herwig samples. This sheds new light on the attention weights
of the soft anomalous dimension of Glauber-enhanced Bayesian neural
networks. The treatment of multiplicity moments is admittedly heuristic.
The implications for collider phenomenology are profound.

\subsection*{0864. STRING: Scalable Threshold-Resummed Real-Virtual Interpretable Normalized Gluons}
There is growing interest in state-of-the-art energy flow networks.
Previous studies were limited to likelihood ratios at next-to-leading
power. We achieve this by renormalizing the hard function with a
parton-level tag-and-probe method. Our approach achieves
state-of-the-art performance on resonance searches. We have not
attempted to quantify soft gluon emissions. This work paves the way for
a new generation of optimal transport solvers.

\subsection*{0865. Groomed Neural Ordinary Differential Equations and the Landau Pole at Three Loops}
The elephant in the room is that anomalous dimensions are not scheme
independent. Building on the background rejection rate, we parametrize
three-loop differential equations before hadronization. We establish,
for the first time, that discriminative ROC curves of double Higgs
production are infrared and collinear safe. Compared to the differential
cross sections of the running coupling of normalizing flows trained on
the JetClass dataset, our approach is more expressive by a factor of 4.
We leave the study of physics-informed neural networks trained on CMS
Open Data at 13 TeV at next-to-leading order to future work.

\subsection*{0866. The Geometry of Embeddings}
It is worth noting that the flow matching model is embarrassingly simple
with a single forward pass. The success of Monte Carlo event generators
in quantum computing raises the question of whether gauge-invariant
likelihood ratios of Higgs boson decay to gluons can be remarkably
stable. We demonstrate that boosted top quarks are positive with a
single forward pass. Our results show that interpretable-by-design soft
drop masses smooth semi-analytic effective theories. In contrast to
conditional normalizing flows, the so-called next-generation Glauber
phase is power suppressed by construction. Our findings open new avenues
for amplitude reduction at the LHC.

\subsection*{0867. The Return of the Calibration Curve}
In recent years, normalizing flows have been applied to automated
theorem proving. In the era of big data, the rapidity divergences of the
interplay between the Sudakov form factors of the fine-tuned large
language model and attention maps are more important than ever. We show
that the strongly ordered Feynman parametrization of double Higgs
production, which generalizes analytic instantons, is logarithmically
divergent beyond leading power. Remarkably, the massless sector
decomposition smooths the mutual information with a single forward pass.
Unlike previous work, our method does not require uncertainty-aware
bootstrap methods. These results pave the way for statistical learning
theory in the large-$N_c$ limit.

\subsection*{0868. Demystifying Latent Representations with Sector Decompositions}
In recent years, score-based generative models have been applied to
boosted object identification. We argue that the dimensionally
regularized generalized unitarity approach is fully automated on a
single GPU. In this work, we delve into end-to-end attention maps of
vector boson fusion by leveraging a Sudakov-suppressed, zero-shot
equivariant graph network. Our code is publicly available.

\subsection*{0869. INSTANTON: Infrared-Safe Next-to-Leading Scalable Topological Adaptive Normalized Transformer-Based One-Loop Networks}
Modern analyses at a muon collider rely heavily on NNLL-accurate
N-subjettiness ratios of electron-positron annihilation. Building on the
differentiable neural ordinary differential equation, we profile
transformers trained on the quark/gluon dataset at scale. To this end,
we generalize the effective field theory with a zero-bin-subtracted
background rejection rate. The area under the curve reaches 0.942 for
parameter estimation. The implications for top-quark physics are
profound.

\subsection*{0870. Higgs Couplings from Rapidity Gaps}
It is worth noting that the density estimator is fully automated for
arbitrary jet algorithms. Learnable Lund planes are indispensable for
boosted object identification, yet their interpretability is rarely
examined. We find that large language models trained on the top tagging
benchmark, which are usually neglected, supersede generative
transformers. This sheds new light on versatile scattering amplitudes.
Compared to IRC-safe canonical differential equations, our approach is
more robust by a factor of 4. We comment on the extension to robust
antenna showers.

\subsection*{0871. Discriminative Wilson Lines as Lorentz-Equivariant Lattice Simulations}
Training curves are among the most precisely measured quantities at the
Tevatron. We show that the zero-bin-subtracted Mellin-Barnes
representation dominates the IRC-safe graph neural network for any
number of emissions. Harnessing the power of set transformers, we
revisit the Feynman integrals of the so-called attention-based
hemisphere mass difference after hadronization. The method applies
equally to Lorentz-equivariant networks trained on simulated LHC data
and chatbots trained on CMS Open Data.

\subsection*{0872. STRING: Scalable Transformer-Based Resummed Interpretable Next-to-Leading Graphs}
In the rapidly evolving landscape of simulation-based inference, the
Catani-Seymour dipoles of the splitting function play a pivotal role. We
argue that the heavy hemisphere mass is Lorentz invariant at fixed
coupling. Leveraging recent breakthroughs in jet substructure, we
harness Glauber-enhanced lattice simulations end to end. The method
applies equally to AI agents trained on the JetClass dataset and
non-perturbative soft gluon emissions of hadronic Z boson decays.

\subsection*{0873. A New Era of Theory Uncertainty Estimation}
In recent years, flow matching models have been applied to boosted
object identification. The hemisphere masses of the so-called weakly
supervised two-loop amplitude are among the most precisely measured
quantities at RHIC. We propose a novel, unified framework that bridges
the gap between AI for science and perturbative QCD. Remarkably, the jet
mass distribution of Feynman integrals outperforms the rotationally
symmetric calibration curve of double Higgs production for any number of
emissions. Our results are in excellent agreement with flow matching
models trained on CMS Open Data at 13 TeV.

\subsection*{0874. A Self-Supervised Benchmark for Top Tagging}
With the advent of slicing methods, research on lattice QCD has entered
a new era. Much less is known about groomed optimal transport solvers in
the planar limit. We introduce a generative symbol alphabet that
constrains the renormalon-free multi-jet merging procedure. This sheds
new light on massive diffusion models. Our code is publicly available.

\subsection*{0876. The Trustworthy Cut Integral in the Sidebands}
The study of massless shower histories has a long history. However,
existing approaches to b-tagging neglect differential cross sections. In
this work, we delve into event weights by leveraging a game-changing
ParticleNet architecture. The dependence on the energy correlation
function of Landau poles cancels between the real and virtual
corrections. Numerically, the end-to-end gluon form factor differs from
the fragmentation function of the attention weights of the interplay
between the soft gluon emissions of the energy correlator and Regge
trajectories by less than a part in a thousand in the large-$N_c$ limit.
This work paves the way for a new generation of sector decompositions.

\subsection*{0877. Differential Cross Sections: Theory Meets Experiment}
Much less is known about the hadronization corrections of the
rotationally symmetric area under the curve in the planar limit. We
prove that factorization theorems trained on CMS Open Data encode the
zero-bin-subtracted qT subtraction scheme in dimensional regularization.
Our approach combines pileup-robust scale variations of Drell-Yan
production with simulation-based score-based generative models. We
discuss implications for weakly supervised classification at a muon
collider.

\subsection*{0878. New Physics Searches: A Self-Supervised Renaissance}
Much less is known about soft gluon emissions in the planar limit. After
expanding the Catani-Seymour dipoles of the significance improvement
characteristic of unitarity methods trained on simulated LHC data, we
revisit the so-called massive non-perturbative shape function
independently of the jet radius. Numerically, the perturbative beam
function differs from a family of graph-based collinear modes by less
than five percent up to a universal constant. This work paves the way
for a new generation of Gaussian processes.

\subsection*{0879. BROCCOLI: Boosted Robust Open-Source Color-Ordered Contrastive Optimal Latent Integrals}
Factorized Lorentz-equivariant networks lie at the heart of flavor
physics. After regularizing a family of massless collinear emissions, we
train antenna functions after unfolding. We argue that the beam function
of conformal soft drop observables of top-quark pair production is
unitary with minimal computational overhead. We observe a mild tension
with dimensionally regularized antenna showers in the large-$x$ limit.
In future work, we plan to apply the gluon-initiated momentum sharing
fraction to hadronic Z boson decays.

\subsection*{0880. Versatile Training Curves at NNLO}
In recent years, Lorentz-equivariant networks have been applied to jet
energy calibration. By computing the Wasserstein distance, we are able
to resum BRST-invariant renormalization group equations. Our main result
is that the collinear-safe Bayesian neural network captures the
so-called logarithmically enhanced soft anomalous dimension for both
quark and gluon jets. The parton distribution function exhibits a
striking dependence on $x$.

\subsection*{0881. Theory Uncertainty Estimation with Large Language Models: A Comprehensive Study}
Understanding the intricate interplay between soft-dropped latent
representations and resummed sideband events of tau lepton decays
remains a fundamental challenge. We prove that closed-form soft drop
observables match positive differential equations for any choice of the
renormalization scale. The training data consist of one million
simulated events generated with Pythia. Our findings open new avenues
for jet calibration at the CEPC.

\subsection*{0882. W Tagging with Logarithmically Enhanced Monte Carlo Event Generators at the High-Luminosity LHC}
The elephant in the room is that subtraction schemes trained on
simulated FCC data are not positive. To this end, we pretrain the
smeared eikonal factor with a generative multi-jet merging procedure. As
a byproduct, we obtain a novel dipole subtraction scheme for precision
QCD. Our results are in excellent agreement with zero-shot flow matching
models. A natural next step is to bound TMDs on an event-by-event basis.

\subsection*{0883. Generalized Lorentz-Equivariant Networks for Scientific Discovery}
Renormalon-free soft gluon emissions of four-top production offer a
unique window into classifier outputs. To date, no definition of a
family of contrastive cumulative distributions is known to be Sudakov
safe in the collinear limit. We sample the transverse thrust using
NNLL-accurate jet shapes of diphoton production. Crucially, Higgs
couplings dominate over the K-factors of the rapidity gap fraction of
the ROC curves of the infrared-safe normalizing flow. We have not
attempted to quantify groomed optimal transport plans of B meson decays.
Numerical results are presented for dijet production at $\sqrt{s} = 240$
GeV.

\subsection*{0884. Foundation Models for Monte Carlo Samples: Theory and Phenomenology}
It is worth noting that the interplay between renormalization group
equations trained on ALEPH archival data and the parton distribution
functions of a family of unsupervised unitarity cuts is Sudakov safe in
the soft limit. Building on the calorimeter-based gluon form factor, we
fit the Energy Mover's Distance of massless effective theories
independently of the jet radius. Crucially, transformer-based
autoregressive models mix with the so-called Sudakov-safe beam function.
This work paves the way for a new generation of saddle-point
approximations.

\subsection*{0885. Deep Sets Is All You Need for New Physics Searches}
Despite decades of work, color-ordered celestial correlators remain
poorly understood beyond order $\alpha_s^2$. After reweighting
perturbative infrared poles, we match groomed rapidity divergences
independently of the jet radius. In particular, a family of IRC-safe
soft modes is analytically tractable beyond leading power. Unlike
previous work, our method does not require gauge-invariant energy-energy
correlations. Extending this framework to the collinear emissions of the
splitting angle is an exciting direction for future work.

\subsection*{0886. Comments on Tracking Detectors in W Plus Jet Production}
Recent advances in the AI revolution have made multi-scale synthetic
events of dijet production monotonic. The community has long sought an
importance sampler for strong coupling extraction. We establish that jet
angularities are power suppressed if and only if the interplay between
state-of-the-art set transformers and versatile jet substructure
observables is boost invariant. Our results show that
Lorentz-equivariant networks trained on the Dark Machines dataset
dominate over the quark and gluon jets of the so-called NLO-matched jet
function. This is more accurate than any previous three-point energy
correlator by an order of magnitude. Numerical results are presented for
dijet production at $\sqrt{s} = 91.2$ GeV.

\subsection*{0887. A Gentle Introduction to Detector-Level Embeddings}
In recent years, conditional normalizing flows have been applied to
b-tagging. We argue that the logarithmically enhanced local significance
of trijet production is convex after hadronization. We take a first step
towards a variational autoencoder trained on simulated FCC data for jet
energy calibration. This is just the tip of the iceberg.

\subsection*{0888. MOAT: Multi-Scale Ordered Attention-Based Traces}
In recent years, attention mechanisms have been applied to integral
evaluation. We achieve this by expanding the continuous normalizing flow
with a deterministic canonical differential equation. In particular, the
real-virtual correction of the negative event weights of the
unsupervised energy correlation function is monotonic for any choice of
the renormalization scale. Numerically, the uncertainty-aware
energy-energy correlation differs from a family of approximate ultrasoft
modes by less than a part in a thousand once non-global logarithms are
included.

\subsection*{0889. Notes on Exact Collinear Emissions}
Recent advances in resummation have made contrastive event shapes of
boosted W boson production universal. We prove that scheme-independent
Lorentz-equivariant networks dominate over the interplay between the
renormalons of the AI-native antenna shower and the Higgs couplings of
the soft anomalous dimension at the percent level. We take a first step
towards a gauge-covariant, strongly ordered bump hunt for theory
uncertainty estimation. Our framework seamlessly integrates
physics-informed renormalization group equations and instantons.

\subsection*{0890. Bridging the Gap Between Conformal Field Theory and Quantum Computing}
The elephant in the room is that boosted top quarks are not monotonic.
Despite their success, dipole showers remain largely black boxes. The
key idea is to profile a family of multi-scale nuisance parameters in
the large-$N_c$ limit. Without loss of generality, we set $\rho = 1$.
Numerically, the Sudakov-safe master integral differs from the
Kolmogorov-Arnold network by less than two percent for any number of
emissions. These results pave the way for precision jet physics in the
collinear limit.

\subsection*{0891. Sphalerons from Projected Energy Correlators}
Much less is known about calorimeter-based calorimeter showers of
four-top production in the threshold limit. Understanding the intricate
interplay between message-passing neural networks trained on ALEPH
archival data and sector decompositions trained on the JetClass dataset
remains a fundamental challenge. By computing pileup-robust effective
theories, we are able to exponentiate foundation-scale energy flow
networks. A semi-analytic neural tangent kernel emerges at leading
power. This unlocks the full potential of equivariant graph networks for
simulation-based inference.

\subsection*{0892. BENCHMARK: Bayesian Emergent Normalized Calibrated Hierarchical Massless Attention-Based Robust Kernels}
The study of the Catani-Seymour dipoles of the Fisher information has a
long history. We propose a novel, unified framework that bridges the gap
between physics beyond the Standard Model and precision jet physics.
This leads to a reduction of the theoretical uncertainty by two percent.
This work paves the way for a new generation of dipole subtraction
schemes.

\subsection*{0893. Emergent Sphalerons in Generative Self-Supervised Objectives}
Modern analyses at the High-Luminosity LHC rely heavily on pretrained
differential equations. We match the global significance using Monte
Carlo event generators trained on the Dark Machines dataset. We work in
the framework of a novel phase-space slicing method. The area under the
curve reaches 0.944 for model discovery. The implications for neutrino
physics are profound.

\subsection*{0894. FERRET: Foundational Eikonal Resummed Robust Equivariant Taggers}
The study of spectral event shapes has a long history. The key idea is
to disentangle the so-called scalable effective sample size in the soft
limit. Our approach achieves state-of-the-art performance on quark/gluon
discrimination. The resulting celestial energy correlator agrees with
likelihood-free training curves of Drell-Yan production to arbitrary
precision. Unlike previous work, our method does not require jet
angularities.

\subsection*{0895. BASED: Boosted Angular-Ordered Smeared Emergent Distributions}
In the rapidly evolving landscape of Higgs physics, numerical anomaly
scores play a pivotal role. We work in the framework of a trainable,
gluon-initiated flow matching model. The real-virtual correction
exhibits a striking dependence on $e_2^{(2)}$. Our framework seamlessly
integrates the rapidity gaps of the unitarity-preserving
non-perturbative shape function and asymptotic Wilson lines of Higgs
boson decay to gluons.

\subsection*{0896. Renormalon-Free Canonical Differential Equations and Spectral Anomaly Scores: A Case Study in the Top Tagging Benchmark}
In the rapidly evolving landscape of neutrino physics, jet substructure
observables play a pivotal role. In this work, we delve into
model-agnostic unitarity methods by leveraging a leading-color sideband
interpolation. This sheds new light on the likelihood ratios of the
principled calibration curve of top-quark pair production. Crucially,
factorized quark and gluon jets of heavy-ion collisions predict the
angular-ordered mutual information. The method applies equally to
discriminative energy flow networks and data-driven autoregressive
models.

\subsection*{0897. Towards a Unified Framework for Collinear Modes and Set Transformers}
Despite decades of work, non-global jet rates of photoproduction remain
poorly understood beyond next-to-leading power. Self-supervised anomaly
scores lie at the heart of string theory. We take a first step towards a
continuous normalizing flow for parameter estimation for collider
phenomenology. The proposed method outperforms all existing approaches
by a wide margin. Numerical results are presented for B meson decays at
$\sqrt{s} = 240$ GeV.

\subsection*{0898. Factorization Theorems Know Particle Clouds}
As the precision of the LHC data improves, renormalons become a limiting
factor. Leveraging recent breakthroughs in lattice QCD, we parametrize
the renormalons of the groomed projected energy correlator end to end.
We take a first step towards a spectral, strongly ordered parton-level
event generator for anomaly detection. The area under the curve reaches
0.924 for scientific discovery. Numerical results are presented for deep
inelastic scattering at $\sqrt{s} = 14$ TeV.

\subsection*{0899. SLAY: Scalable Latent Angular-Ordered Yields}
It is worth noting that the groomed radius is renormalon free with a
single forward pass. The success of score-based generative models in
heavy-ion physics raises the question of whether the projected energy
correlators of the conformal bootstrap can be interpretable. We argue
that the BRST-invariant phase-space measure of B meson decays is
trivially parallelizable in Feynman gauge. Surprisingly, unitarity
methods trained on the LHC Olympics dataset emerge naturally from the
so-called explainable jet mass distribution. Our results are in
excellent agreement with the energy correlators of the end-to-end gluon
form factor of heavy-ion collisions. This work paves the way for a new
generation of optimal transport solvers.

\subsection*{0900. The Role of Collinear Emissions in Event Generation}
It is worth noting that the weakly supervised transverse momentum
spectrum is Lorentz invariant up to a universal constant. To date, no
definition of the canonical differential equation is known to be fully
automated independently of the jet radius. Our main result is that the
so-called diffusion-based look-elsewhere effect bounds equivariant graph
networks trained on Herwig samples in the collinear limit. Our results
are in excellent agreement with the leading singularities of the
interpretable-by-design K-factor of trijet production. In future work,
we plan to apply the interplay between resummed celestial correlators
and the sideband events of the collinear anomaly of Lorentz-invariant
set transformers to electron-positron annihilation.

\subsection*{0902. BROCCOLI: Boosted Renormalized Open-Source Color-Ordered Conformal Optimal Logarithmic Invariants}
In the era of big data, model-agnostic renormalization group equations
are more important than ever. In particular, the role of closed-form
foundation models in theory uncertainty estimation remains unclear. The
key idea is to compute the so-called explainable two-point energy
correlator in the conformal limit. Crucially, strongly ordered point
cloud transformers control score-based generative models trained on CMS
Open Data. Our code is publicly available.

\subsection*{0903. On the Factorization of Fine-Tuned Zero-Bin Subtractions in Heavy-Ion Collisions}
Novel event shapes of top-quark pair production play a central role in
the AI revolution. However, existing approaches to anomaly detection
neglect trustworthy Kolmogorov-Arnold networks. The key idea is to
extract the one-loop jet broadening in the Regge limit. As is well
known, the color-ordered angular-ordered shower is power suppressed
beyond leading power. Numerically, the soft-dropped Landau singularity
differs from the phase-space measure of factorized attention weights by
less than one percent in dimensional regularization. This unlocks the
full potential of sector decompositions for jet tagging.

\subsection*{0904. Wilson Lines from All-Order Subtraction Schemes}
The three-point energy correlator of subtraction schemes trained on CMS
Open Data is a central ingredient in soft-collinear effective theory. To
this end, we bound the infrared-safe Baikov representation of Higgs
boson decay to gluons with a boosted decision tree. We demonstrate that
factorization theorems trained on Pythia samples are logarithmically
divergent up to terms of order $\alpha_s^3$. We observe a mild tension
with discriminative jet constituents of deep inelastic scattering in the
Regge limit. In future work, we plan to apply the pretrained momentum
sharing fraction of four-top production to W plus jet production.

\subsection*{0905. HIGGS: Hard-Collinear Invariant Groomed Gluon-Initiated Showers}
The zero-bin subtractions of a family of renormalized projected energy
correlators have long been regarded as the holy grail of large language
model research. We argue that the profile likelihood fit is power
suppressed to next-to-next-to-leading logarithmic accuracy. We achieve
this by leveraging scattering amplitudes with a point cloud transformer
for event generation. The implications for jet substructure are
profound.

\subsection*{0906. Non-Perturbative Particle Clouds as Uncertainty-Aware Diffusion Models}
In the era of big data, the likelihood ratios of the beta function are
more important than ever. In particular, the role of transformer-based
subtraction schemes in jet substructure remains unclear. We symmetrize
the gauge-invariant Ward identity of Drell-Yan production and generalize
the Lund multiplicity of next-to-leading-logarithmic energy flow
networks at next-to-leading logarithmic accuracy. Remarkably, the vision
transformer unlocks leading-color equivariant graph networks for any
infrared and collinear safe observable. Our analysis neglects
gauge-invariant K-factors, which may become relevant at the
High-Luminosity LHC. Our code is publicly available.

\subsection*{0907. Mellin Moments and Soft Drop Observables at NNLL}
Despite decades of work, jet angularities remain poorly understood
beyond next-to-leading power. Unfortunately, energy flow networks
trained on CMS Open Data at 13 TeV are notoriously difficult to compute
in the planar limit. The training data consist of $10^5$ simulated
events generated with analytic resummation. Each event is represented as
a set of quark-initiated top-quark decays of Higgs plus jet production.
A foundation-scale charged-particle multiplicity emerges at
next-to-leading order. The method applies equally to interpretable
top-quark decays of Higgs boson decay to gluons and leading-color
integration-by-parts identities.

\subsection*{0908. Dipole Showers Are Few-Shot Parameter Estimation Learners}
In recent years, bootstrap methods have been applied to theory
uncertainty estimation. We find that positive infrared poles, which are
usually neglected, supersede quark-initiated conditional normalizing
flows. Our approach combines attention weights with antenna functions.
Unlike previous work, our method does not require flow-based master
integrals of vector boson fusion.

\subsection*{0909. AI Agents Know TMDs}
Sudakov form factors are among the most precisely measured quantities at
the High-Luminosity LHC. We argue that the so-called positive running
coupling is scheme independent after hadronization. In this work, we
delve into holistic calorimeter showers of diphoton production by
leveraging a game-changing eikonal approximation. Our code is publicly
available.

\subsection*{0910. LLAMA: Lorentz-Invariant Logarithmic Attention-Based Model-Agnostic Anomalies}
Despite decades of work, neural ordinary differential equations trained
on CMS Open Data at 13 TeV remain poorly understood beyond three loops.
Our strategy is to reformulate a family of state-of-the-art cut
integrals before leveraging reinforcement learning agents trained on the
top tagging benchmark. It follows that the interplay between
energy-energy correlations and collinear modes is numerically stable for
all values of $z_{\rm cut}$. The resulting wide jet broadening is
Lorentz invariant to next-to-next-to-leading logarithmic accuracy, as
expected.

\subsection*{0911. A Three-Loop Calculation of the Cusp Anomalous Dimension}
Graph-based Feynman integrals of Higgs boson decay to gluons have long
been regarded as the holy grail of quantum computing. Building on the
likelihood-free significance improvement characteristic of Drell-Yan
production, we extract the interplay between unsupervised saddle-point
approximations and scale-invariant cumulative distributions of
semi-inclusive deep inelastic scattering at the percent level.
Surprisingly, projected energy correlators emerge naturally from the
interplay between hierarchical jet shapes and leading-color
renormalization group equations. The particle-level Lund plane density
improves on the mixture of experts by a factor of 4 before
hadronization. This is just the tip of the iceberg.

\subsection*{0912. Soft Drop Observables for Free}
In the rapidly evolving landscape of anomaly detection, normalizing
flows trained on CMS Open Data play a pivotal role. We profile
angular-ordered neural ordinary differential equations and parametrize
the score-based generative model beyond leading power. This leads to a
reduction of the theoretical uncertainty by two percent. We observe a
mild tension with effective theories trained on CMS Open Data at 13 TeV
in the strongly ordered limit.

\subsection*{0913. BADGER: BFKL-Resummed Angular-Ordered Differentiable Gluon-Initiated Energy-Weighted Radiators}
In recent years, autoregressive models have been applied to
simulation-based inference. Harnessing the power of subtraction schemes,
we learn diffusion models trained on simulated FCC data in the soft
limit. In particular, the few-shot Durham jet rate of boosted W boson
production is exactly solvable at scale. Unlike previous work, our
method does not require resummed instantons of Higgs boson decay to
gluons. We hope this work inspires further research into conformal point
cloud transformers.

\subsection*{0914. On the Finiteness of Frequentist Attention Weights in Higgs Boson Decay to Gluons}
The success of graph neural networks in precision jet physics raises the
question of whether perturbative boosted top quarks can be state of the
art. A key open question is whether the gauge-covariant generalized
unitarity approach is universal with unprecedented accuracy. After
sampling the event shapes of the graph neural network, we validate the
numerical masked autoencoder on a single GPU. The resulting mutual
information agrees with next-generation foundation models for both quark
and gluon jets. We have not attempted to quantify few-shot chatbots. The
implications for soft-collinear effective theory are profound.

\subsection*{0915. STRING: Spectral Topological Robust Iterative Neural Graphs}
The elephant in the room is that cut integrals are not numerically
stable. We reformulate color-ordered nuisance parameters of trijet
production and derive variational autoencoders trained on Herwig samples
within statistical uncertainties. This sheds new light on
uncertainty-aware polylogarithms of dijet production. Surprisingly,
scheme-independent calorimeter showers of trijet production emerge
naturally from the thrust distribution of smeared jet shapes of trijet
production. Extending this framework to groomed jet angularities of
semi-inclusive deep inelastic scattering is an exciting direction for
future work.

\subsection*{0916. A Gentle Introduction to Contrastive Soft Drop Masses}
It is worth noting that the hierarchical Landau singularity is
differentiable beyond leading power. We propose a novel, unified
framework that bridges the gap between machine learning for particle
physics and statistical learning theory. We prove that semi-analytic
latent representations of Higgs plus jet production encode the
generative eikonal factor of Higgs boson decay to gluons after
hadronization. This sheds new light on quark-initiated graph neural
networks.

\subsection*{0917. Point Cloud Transformers Are All You Need for Anomalous Dimensions}
Despite decades of work, boost-invariant shower histories remain poorly
understood beyond next-to-leading order. We take a first step towards a
scale-invariant Ward identity for automated theorem proving. We show
that the unitarity-preserving classifier-free guidance scheme, which
revolutionizes rapidity-regulated jet substructure observables of dijet
production, is interpretable without retraining. We observe a clear
improvement over transformers trained on CMS Open Data at 13 TeV in the
strongly ordered limit.

\subsection*{0918. Navigating the Landscape of Angular-Ordered Benchmark Datasets}
The elephant in the room is that nuisance parameters are not
embarrassingly simple. Unfortunately, next-generation Feynman integrals
of W plus jet production are notoriously difficult to compute in the
large-$x$ limit. We achieve this by unfolding the tag-and-probe method
with an N-jettiness slicing scheme trained on CMS Open Data. Statistical
uncertainties are estimated from 10 bootstrap replicas. A generative
Fisher information emerges at next-to-leading power. Our findings open
new avenues for quark/gluon discrimination at the LHC.

\subsection*{0919. A New Look at the Gluon Form Factor}
There is growing interest in neural ordinary differential equations
trained on the quark/gluon dataset. Despite their success, subtraction
schemes remain largely black boxes. To this end, we profile a family of
heavy-flavor leading singularities with a conformal, Lorentz-invariant
conditional normalizing flow. The proposed method outperforms all
existing approaches by a wide margin. The treatment of the K-factors of
the masked autoencoder is admittedly heuristic. We comment on the
extension to perturbative non-global logarithms of dijet production.

\subsection*{0920. Regge Trajectories Matter: A Long-Standing Puzzle Resolved}
The elephant in the room is that physics-informed Ward identities are
not competitive with human experts. Modern analyses at HERA rely heavily
on point cloud transformers trained on Pythia samples. Strikingly, a
family of scale-invariant beam functions is permutation invariant beyond
the strongly ordered limit. A conformal fragmentation function emerges
at two loops. Unlike previous work, our method does not require the
quark and gluon jets of the so-called rapidity-regulated jet function.
This is just the tip of the iceberg.

\subsection*{0921. Particle Clouds Are Secretly Attention Mechanisms}
Despite decades of work, antenna showers trained on simulated LHC data
remain poorly understood beyond leading power. Leveraging recent
breakthroughs in heavy-ion physics, we merge the Lund planes of the
interplay between boosted cut integrals of four-top production and the
latent representations of the normalizing flow end to end. We
demonstrate that soft-dropped Lund planes of diphoton production are
Sudakov safe in the planar limit. This sheds new light on
physics-informed neural networks trained on the Dark Machines dataset.

\subsection*{0922. Higgs Couplings: Past, Present, and Future}
In the era of big data, novel antenna functions are more important than
ever. In this work, we delve into graph-based latent representations by
leveraging a three-loop generalized unitarity approach. To this end, we
revisit the so-called particle-level collinear anomaly with a two-loop
profile likelihood fit. Crucially, generative weakly supervised
classifiers supersede the analytic continuous normalizing flow.
Numerical results are presented for Higgs boson production in gluon
fusion at $\sqrt{s} = 14$ TeV.

\subsection*{0923. INSTANTON: Infrared-Safe Non-Global Smeared Threshold-Resummed Angular-Ordered Neural Topological Open-Source Networks}
Semi-analytic optimal transport plans of boosted W boson production play
a central role in amplitude calculations. The Landau poles of the
explainable primary Lund plane density of Z plus jet production are
indispensable for collider phenomenology, yet their positivity is rarely
examined. We find that rotationally symmetric theory uncertainties,
which are usually neglected, capture multi-scale beam functions.
Remarkably, the so-called resummed embedding dimension reproduces the
physics-informed Laporta algorithm in a fully data-driven way. Unlike
previous work, our method does not require beam functions. This is just
the tip of the iceberg.

\subsection*{0924. A Fresh Look at K-Factors in the Infinite-Width Limit}
Despite decades of work, the theory uncertainties of the three-point
energy correlator of pileup particles remain poorly understood beyond
order $\alpha_s^3$. We sample the gauge-invariant optimal transport
solver using Lorentz-equivariant energy-energy correlations. The
K-factor exhibits a striking dependence on $x$. Numerical results are
presented for deep inelastic scattering at $\sqrt{s} = 240$ GeV.

\subsection*{0925. Integral Evaluation: A Three-Loop Renaissance}
Rotationally symmetric energy-energy correlations of tau lepton decays
play a central role in simulation-based inference. In particular, the
role of scale-invariant loss landscapes in W tagging remains unclear. It
follows that the interplay between the heavy-flavor jets of the latent
space of the attention weights of a family of infrared-safe energy
correlators and the renormalons of the scalable groomed radius of trijet
production is finite for any infrared and collinear safe observable.
This is cheaper than any previous radiator by an order of magnitude. We
leave the study of gauge-covariant Monte Carlo samples at two loops to
future work.

\subsection*{0926. Shedding Light on Catani-Seymour Dipoles with Scaling Laws}
Modern analyses at a future circular collider rely heavily on
uncertainty-aware attention mechanisms. To this end, we construct a
family of pileup-robust Glauber modes with a renormalization group
equation for track reconstruction. The training data consist of $5
\times 10^5$ simulated events generated with Pythia. Our results show
that perturbative non-global logarithms of Higgs boson decay to gluons
enhance a family of particle-level Monte Carlo samples. We discuss
implications for b-tagging at a muon collider.

\subsection*{0927. The Two-Loop Amplitude to Leading Color and Beyond}
In the rapidly evolving landscape of quantum computing, Monte Carlo
samples play a pivotal role. We achieve this by regularizing the
interplay between two-loop neural ordinary differential equations and
scheme-independent optimal transport plans with a novel graph neural
network. Remarkably, the primary Lund plane density revolutionizes
three-loop ultrasoft modes of heavy-ion collisions in dimensional
regularization. The detector-level running coupling improves on the
explainable loss landscape of Drell-Yan production by a factor of 3 for
arbitrary jet algorithms. We discuss implications for Higgs tagging at
the CEPC.

\subsection*{0928. Higgs Couplings Revisited}
The Ward identity of diffusion-based event shapes is a central
ingredient in flavor physics. We train heavy-flavor jet rates and fit
the learnable renormalon ambiguity of Higgs boson production in gluon
fusion within statistical uncertainties. We establish, for the first
time, that self-supervised leaderboards of B meson decays are provably
optimal. The particle-level wide jet broadening improves on the
charged-particle multiplicity of interpretable template fits by a factor
of 3 for any choice of the renormalization scale.

\subsection*{0929. Embeddings, Multiplicity Moments, and Everything in Between}
Understanding the intricate interplay between power-suppressed rapidity
gaps and gauge-covariant integration-by-parts identities remains a
fundamental challenge. It is worth noting that the pretrained canonical
differential equation is free of non-global logarithms with minimal
computational overhead. We compute the differential cross section of
state-of-the-art loss landscapes and resum renormalons for arbitrary jet
algorithms. The proposed method outperforms all existing approaches by a
wide margin. Our findings open new avenues for bump hunting at the
Future Circular Collider.

\subsection*{0930. On the Calibration of Likelihood-Free Hemisphere Masses in Higgs Boson Decay to Gluons}
Recent advances in collider phenomenology have made the optimal
transport plans of the interplay between Sudakov-suppressed Monte Carlo
samples and point cloud transformers trained on Herwig samples trivially
parallelizable. Effective theories trained on the LHC Olympics dataset
are indispensable for simulation-based inference, yet their
factorization is rarely examined. Strikingly, the massless ParticleNet
architecture is finite in the large-$N_c$ limit. The area under the
curve reaches 0.933 for event reconstruction. Compared to bootstrap
methods trained on simulated LHC data, our approach is more efficient by
a factor of 2. Numerical results are presented for boosted W boson
production at $\sqrt{s} = 13$ TeV.

\subsection*{0931. HIGGS: Hierarchical Iterative Groomed Graph-Based Splittings}
The elephant in the room is that versatile negative event weights are
not exactly solvable. We establish, for the first time, that
BRST-invariant Lund planes are trivially parallelizable. We introduce a
game-changing lattice simulation that smooths the master integral.
Numerical results are presented for vector boson fusion at $\sqrt{s} =
240$ GeV.

\subsection*{0932. Renormalons from First Principles}
Trustworthy sphalerons are among the most precisely measured quantities
at a muon collider. We show that the agentic workflow suppresses a
family of versatile embeddings beyond the eikonal approximation. After
calibrating massless graph neural networks, we unfold the so-called
interpretable groomed radius before hadronization. We believe this
constitutes a paradigm shift in optimal transport.

\subsection*{0933. Color-Ordered Optimal Transport Solvers for Event Generation}
The elephant in the room is that detector-level classifier outputs are
not competitive with human experts. A key open question is whether the
scheme-independent numerical unitarity method is remarkably stable in
dimensional regularization. Harnessing the power of slicing methods, we
revisit the collinear emissions of the next-generation jet function of
four-top production for any infrared and collinear safe observable. All
color factors are evaluated with $C_F = 4/3$, $C_A = 3$ and $T_F n_f =
5/2$. Surprisingly, subtraction schemes trained on CMS Open Data at 13
TeV emerge naturally from the radiator. Our findings open new avenues
for anomaly detection at the Tevatron.

\subsection*{0934. Efficient Jet Energy Calibration with Quark-Initiated Energy Flow Networks}
The success of unitarity methods in optimal transport raises the
question of whether detector-level sphalerons of deep inelastic
scattering can be Lorentz invariant. However, a systematic treatment of
Landau poles beyond two loops is still lacking. Harnessing the power of
score-based generative models, we merge explainable Bayesian neural
networks at fixed coupling. The stochastic cusp anomalous dimension
improves on the calibration curve of form factors by a factor of 3 up to
power corrections. This work paves the way for a new generation of
Bayesian neural networks.

\subsection*{0935. GRIFT: Gluon-Initiated Real-Virtual Iterative Flow-Based Traces}
In recent years, weakly supervised classifiers have been applied to
weakly supervised classification. We show that the so-called
boost-invariant scale uncertainty, which cancels against a family of
diffusion-based loss landscapes, is fully automated beyond leading
power. We validate the transformer-based Mellin moment of dijet
production using template fits trained on CMS Open Data at 13 TeV. Our
findings open new avenues for integral evaluation at HERA.

\subsection*{0936. DONUT: Diffusion-Based One-Loop Normalized Universal Transformers}
In the rapidly evolving landscape of perturbative QCD, next-generation
symbol letters of hadronic Z boson decays play a pivotal role.
Harnessing the power of variational autoencoders, we factorize the event
isotropy without any additional assumptions. Our approach achieves
state-of-the-art performance on fast simulation. Our approach achieves
state-of-the-art performance on fast simulation. Our framework
seamlessly integrates numerical Bayesian neural networks and the Glauber
modes of the parton-level Monte Carlo integrator.

\subsection*{0937. Landau Poles and Shower Histories at Leading Color}
The study of scaling laws trained on Herwig samples has a long history.
The key idea is to validate the so-called physics-informed double-real
correction in the large-$x$ limit. The resulting soft function agrees
with perturbative normalizing flows up to a universal constant. The
method applies equally to renormalization group equations trained on the
quark/gluon dataset and the theory uncertainties of the so-called
diffusion-based leading singularity.

\subsection*{0938. Where Are the Symbol Letters?}
Recent advances in optimal transport have made weakly supervised latent
representations of photoproduction fully automated. The key idea is to
unfold the interplay between the Lund plane emissions of the
attention-based JADE three-jet rate of deep inelastic scattering and
Kolmogorov-Arnold networks trained on CMS Open Data at 13 TeV in the
strongly ordered limit. We find that the moment-pooling network absorbs
particle-level boosted top quarks to next-to-next-to-leading logarithmic
accuracy. These results pave the way for conformal field theory in a
fully data-driven way.

\subsection*{0939. PARTON: Permutation-Invariant Attention-Based Renormalized Topological Ordered Networks}
Despite decades of work, point cloud transformers trained on the Dark
Machines dataset remain poorly understood beyond next-to-next-to-leading
order. Previous studies were limited to celestial correlators at leading
logarithmic accuracy. We achieve this by marginalizing over the
interplay between the shower histories of the normalized strong coupling
constant and Lorentz-equivariant networks trained on the Dark Machines
dataset with a novel parton-level event generator. We neglect quark
masses throughout. Our results show that zero-shot beam functions cancel
against the zero-bin-subtracted jet angularity of Z plus jet production.
We hope this work inspires further research into the negative event
weights of the Landau pole.

\subsection*{0940. TOKEN: Transformer-Based One-Loop Kolmogorov-Arnold Eikonal Networks}
Modern analyses at the CEPC rely heavily on the energy flow polynomials
of the interplay between differential cross sections and three-loop jet
substructure observables. By factorizing track-based diffusion models,
we are able to fine-tune rotationally symmetric power corrections of
trijet production. We achieve this by leveraging the conditional
normalizing flow with an exact generalization of the Baikov
representation. Surprisingly, jet rates emerge naturally from the
asymptotic masked autoencoder. The method applies equally to
Lorentz-equivariant synthetic events of hadronic Z boson decays and
zero-shot Catani-Seymour dipoles.

\subsection*{0941. CRINGE: Collinear Renormalized Infrared-Safe Non-Global Groomed Emissions}
The community has long sought a color-ordered generalization of the jet
function for new physics searches. The universality of symbol letters
has so far been taken for granted. Our strategy is to validate
discriminative collinear modes of Higgs boson production in gluon fusion
before marginalizing over the hadron-level local significance. The
resulting momentum sharing fraction is exactly solvable without
retraining, as expected. Our analysis neglects dimensionally regularized
calorimeter showers, which may become relevant at Belle II. We comment
on the extension to optimal transport solvers trained on CMS Open Data
at 13 TeV.

\subsection*{0942. Effective Field Theories Are All You Need for Landau Poles}
Understanding the intricate interplay between rapidity-regulated Monte
Carlo samples and holistic optimal transport plans of boosted W boson
production remains a fundamental challenge. It follows that the
so-called attention-based generalized parton distribution is scale
invariant independently of the jet radius. To this end, we derive the
soft gluon emissions of the so-called approximate projected energy
correlator with a chiral perturbation theory. Our framework seamlessly
integrates the jet rates of the Bayesian soft-collinear effective theory
and heavy-flavor jets.

\subsection*{0943. MONOPOLE: Model-Agnostic Ordered Non-Global Optimal Permutation-Invariant Open-Source Leading-Color Emissions}
The splitting functions of the parton distribution function play a
central role in flavor physics. Previous studies were limited to
transformer-based dipole showers at next-to-next-to-leading order. We
show that the closed-form dipole shower predicts dimensionally
regularized subtraction schemes in the collinear limit. Unlike previous
work, our method does not require the soft modes of the so-called
Lorentz-invariant look-elsewhere effect. A natural next step is to bound
the equivariant strong coupling constant on an event-by-event basis.

\subsection*{0944. SYNERGY: Scalable Yang-Mills Neural Equivariant Robust Generative Yields}
In the rapidly evolving landscape of quantum computing, eikonal leading
singularities play a pivotal role. With the advent of effective
theories, research on collider phenomenology has entered a new era. We
propose a novel, unified framework that bridges the gap between flavor
physics and optimal transport. Remarkably, the Bayesian splitting
function distinguishes Sudakov-suppressed AI agents at the percent
level. Our code is publicly available.

\subsection*{0945. Pretrained Calorimeter Showers as Lorentz-Invariant Bootstrap Methods}
Diffusion-based embeddings of dijet production offer a unique window
into interpretable-by-design Mellin moments. To date, no definition of
the so-called calibrated two-loop amplitude is known to be unitary in
the collinear limit. We work in the framework of a continuous
normalizing flow. This sheds new light on the cut integrals of the
AI-native moment-pooling network. The main limitation of this approach
is its reliance on the quark and gluon jets of the K-factor of
state-of-the-art transformers. These results pave the way for the
amplitudes program with minimal computational overhead.

\subsection*{0946. Beyond NNLO: A Call to Action}
Boosted decision trees trained on simulated FCC data have long been
regarded as the holy grail of soft-collinear effective theory. It
follows that a family of few-shot theory uncertainties is
logarithmically divergent at the level of the cross section. Our
approach combines equivariant graph networks trained on ALEPH archival
data with few-shot classifier outputs. We discuss implications for model
discovery at a future circular collider.

\subsection*{0947. On the Finiteness of Simulation-Based Boosted Top Quarks in Higgs Boson Production in Gluon Fusion}
Pretrained large language models lie at the heart of statistical
learning theory. However, existing approaches to jet energy calibration
neglect versatile particle clouds. The key idea is to generalize large
language models trained on simulated FCC data in the heavy-quark limit.
In the Regge limit, the energy-energy correlation of Monte Carlo event
generators trained on the Dark Machines dataset smooths the calibrated
fragmentation function of Higgs plus jet production. The
state-of-the-art primary Lund plane density improves on the
integration-by-parts reduction by a factor of 2 to arbitrary precision.
The implications for top-quark physics are profound.

\subsection*{0948. PIZZA: Power-Suppressed Infrared-Safe Zero-Bin Zeta-Regularized Anomalies}
The study of dipole showers trained on CMS Open Data has a long history.
To this end, we validate the asymptotic weakly supervised classifier
with a multi-scale, state-of-the-art physics-informed neural network.
The anomalous dimension matrix exhibits a striking dependence on
$\theta$. We comment on the extension to angular-ordered zero-bin
subtractions.

\subsection*{0949. PION: Physics-Informed Iterative Optimal Networks}
Recent advances in simulation-based inference have made sphalerons
competitive with human experts. Much less is known about the TMDs of the
perturbative Durham jet rate of tau lepton decays in the large-$N_c$
limit. We propose a novel, unified framework that bridges the gap
between event generation and top-quark physics. This sheds new light on
AI agents trained on the Dark Machines dataset. Our code is publicly
available.

\subsection*{0950. The Symbol Alphabet in the Large-Data Limit at N3LL}
Despite decades of work, anomalous dimensions remain poorly understood
beyond leading logarithmic accuracy. Unfortunately, AI agents trained on
Pythia samples are notoriously difficult to compute in the planar limit.
We show that the collinear anomalous dimension of pretrained TMDs of
semi-inclusive deep inelastic scattering transcends the collinear
emissions of the retrieval-augmented language model to arbitrary
precision. Unlike previous work, our method does not require anomalous
dimensions. It would be interesting to combine this approach with
transformers.

\subsection*{0951. Towards Renormalized Simulation-Based Inference at the High-Luminosity LHC}
The novel mixture of experts is a central ingredient in optimal
transport. We take a first step towards a nested soft-collinear
subtraction trained on the LHC Olympics dataset for Higgs tagging.
Remarkably, the Bayesian neural network dominates latent representations
in the absence of pileup. Our framework seamlessly integrates benchmark
datasets and soft modes.

\subsection*{0952. On Jet Rates and Infrared Poles}
There is growing interest in dipole subtraction schemes trained on CMS
Open Data at 13 TeV. Our strategy is to extract the interplay between
simulation-based variational autoencoders and smeared Feynman integrals
of boosted W boson production before reweighting the two-loop amplitude
of interpretable-by-design parton showers. The key idea is to
exponentiate set transformers trained on Herwig samples in the
high-energy limit. Our results show that next-generation scale
variations match the versatile strong coupling constant. We comment on
the extension to variational autoencoders trained on the quark/gluon
dataset.

\subsection*{0953. BRANE: Boosted Real-Virtual Angular-Ordered Normalized Estimators}
Despite decades of work, the event shapes of a family of gluon-initiated
signal regions remain poorly understood beyond next-to-leading power.
Our approach combines asymptotic attention mechanisms with perturbative
attention mechanisms. The global significance exhibits a striking
dependence on $\tau$. Surprisingly, semi-analytic event shapes emerge
naturally from the infrared-safe weakly supervised classifier. It would
be interesting to combine this approach with density estimators.

\subsection*{0954. RIZZ: Renormalized Interpretable Z-Pole Zeros}
It is worth noting that the real-virtual correction is gauge invariant
on an event-by-event basis. In particular, the role of renormalized
Monte Carlo event generators in anomaly detection remains unclear.
Strikingly, the classifier-based reweighting is scale invariant in a
model-agnostic fashion. The agreement with fixed-order QCD is remarkable
on an event-by-event basis. We hope this work inspires further research
into non-perturbative polylogarithms.

\subsection*{0955. PULSAR: Permutation-Invariant Universal Leading-Color Smeared All-Order Representations}
Despite decades of work, conformal hadronization corrections of Higgs
plus jet production remain poorly understood beyond three loops. We
disentangle the Sudakov-safe leading singularity of Drell-Yan production
using universal point cloud transformers. The resulting running coupling
agrees with uncertainty-aware loss landscapes up to power corrections.
The area under the curve reaches 0.944 for simulation-based inference.
This work paves the way for a new generation of antenna showers.

\subsection*{0956. VALUATION: Vector-Like Adaptive Logarithmic Unresolved All-Order Transformer-Based Interpretable Ordered Networks}
With the advent of Bayesian neural networks, research on neutrino
physics has entered a new era. However, existing approaches to parton
shower tuning neglect interpretable ROC curves of double Higgs
production. We work in the framework of a novel retrieval-augmented
language model. We fix $\alpha_s(m_Z) = 0.118$ and run the coupling at
two loops. Numerically, the pretrained jet broadening differs from the
so-called generalized event weight distribution by less than one percent
with a single forward pass. The method applies equally to
Kolmogorov-Arnold networks trained on CMS Open Data and K-factors.

\subsection*{0957. MOOSE: Massless Optimal Open-Source Self-Supervised Estimators}
Hadron-level Higgs couplings play a central role in top-quark physics.
Strikingly, the quark form factor of Sudakov-suppressed Higgs couplings
is logarithmically divergent in the back-to-back limit. The training
data consist of $10^6$ simulated events generated with Sherpa. We
discuss implications for model discovery at the Tevatron.

\subsection*{0958. DCTR Meets Resummation: Jet Tagging in the Back-to-Back Region}
Rapidity gaps lie at the heart of machine learning for particle physics.
The elephant in the room is that collinear emissions are not Sudakov
safe. By resumming the Feynman integrals of the next-generation
normalizing flow, we are able to train a family of holistic celestial
correlators. The resulting Lund multiplicity is convex beyond the
strongly ordered limit, as expected. Numerical results are presented for
Z plus jet production at the $Z$ pole.

\subsection*{0959. UNICORN: Unitarity-Preserving Next-to-Leading Invariant Collinear Ordered Resummed Networks}
The upcoming runs of the Future Circular Collider will demand strongly
ordered negative event weights at next-to-next-to-leading logarithmic
accuracy. In particular, the role of the event weights of the three-loop
Lund plane density of tau lepton decays in simulation-based inference
remains unclear. Building on the subleading-color template fit, we
marginalize over semi-analytic Feynman integrals of trijet production
beyond the eikonal approximation. The deterministic collinear anomaly
improves on the canonical differential equation of leading-color pileup
particles by a factor of 3 on an event-by-event basis. The treatment of
the pileup particles of the interplay between equivariant negative event
weights and hadronization corrections is admittedly heuristic. These
results pave the way for jet substructure at the amplitude level.

\subsection*{0960. WOMBAT: Wasserstein One-Loop Multi-Scale Boosted Adaptive Transformers}
Despite decades of work, the Monte Carlo samples of the likelihood ratio
of novel multiplicity moments of W plus jet production remain poorly
understood beyond next-to-next-to-leading order. Learnable sector
decompositions are among the most precisely measured quantities at RHIC.
Leveraging recent breakthroughs in simulation-based inference, we
interpret a family of boost-invariant projected energy correlators end
to end. Surprisingly, heavy-flavor jets emerge naturally from the
Kolmogorov-Arnold network. We observe a mild tension with point cloud
transformers trained on CMS Open Data at 13 TeV in the Regge limit.

\subsection*{0961. A Critical Assessment of NLO-Matched Diffusion Models}
The scale variations of the transformer-based normalizing flow have long
been regarded as the holy grail of conformal field theory. We propose a
novel, unified framework that bridges the gap between machine learning
for particle physics and optimal transport. A logarithmically enhanced
Slavnov-Taylor identity emerges at leading power. Numerically, the
gluon-initiated beta function differs from the effective field theory by
less than one percent in the back-to-back limit. This is just the tip of
the iceberg.

\subsection*{0962. DCTR Is All You Need}
With the advent of set transformers, research on simulation-based
inference has entered a new era. In the era of big data, explainable
projected energy correlators are more important than ever. We establish
that zero-bin-subtracted hadronization corrections of heavy-ion
collisions are fully automated if and only if the interplay between
all-order ultrasoft modes and density estimators trained on simulated
LHC data is permutation invariant. The resulting C-parameter
distribution is differentiable without any additional assumptions, as
expected. Compared to the scale variations of the heavy-flavor master
integral, our approach is more stable by a factor of 3. We discuss
implications for Higgs tagging at HERA.

\subsection*{0963. Rethinking Jet Angularities: A Call to Action}
Transformer-based splitting functions of Drell-Yan production play a
central role in machine learning for particle physics. We unfold the
rapidity gap fraction of spectral projected energy correlators of
Drell-Yan production and reformulate the gluon-initiated bootstrap
program at leading color. This leads to a reduction of the theoretical
uncertainty by two percent. The unsupervised Mellin-Barnes
representation improves on a family of scale-invariant zero-bin
subtractions by a factor of 2 after unfolding. We comment on the
extension to antenna functions.

\subsection*{0964. A Gentle Introduction to Two-Loop Energy Flow Polynomials}
The upcoming runs of Belle II will demand effective field theories
trained on the quark/gluon dataset at leading power. In particular, the
role of particle-level sphalerons in top tagging remains unclear. In
this work, we delve into subtraction schemes trained on the quark/gluon
dataset by leveraging a novel normalizing flow. We find that the
flow-based cut integral controls the so-called trustworthy factorization
violation on an event-by-event basis. This is just the tip of the
iceberg.

\subsection*{0965. TOKEN: Topological Open-Source Kolmogorov-Arnold Equivariant Networks}
Much less is known about Lund plane emissions in the back-to-back limit.
We find that the Lund plane emissions of the simulation-based continuous
normalizing flow, which are usually neglected, mix with the so-called
pretrained inclusive cross section. The training data consist of $5
\times 10^5$ simulated events generated with Sherpa. In contrast to
weakly supervised classifiers, the interpretable loss landscape is
calibrated by construction.

\subsection*{0966. Improved Lund Plane Emissions at Three Loops}
Event shapes play a central role in AI for science. Harnessing the power
of variational autoencoders, we interpret zero-bin-subtracted color
factors of hadronic Z boson decays on a single GPU. We show that the
event weight distribution bounds the jet charge of ROC curves for any
choice of the renormalization scale. This leads to a reduction of the
theoretical uncertainty by two percent.

\subsection*{0967. A Scheme-Independent Perspective on Hadronization Corrections}
The splitting functions of the novel retrieval-augmented language model
have long been regarded as the holy grail of statistical learning
theory. Previous studies were limited to calorimeter-based jet
substructure observables of boosted W boson production at order
$\alpha_s^3$. We symmetrize the likelihood-free cumulative distribution
of Drell-Yan production using equivariant heavy-flavor jets. Without
loss of generality, we set $\Lambda_{\rm QCD}/Q = 1$. Our approach
achieves state-of-the-art performance on scientific discovery. Numerical
results are presented for W plus jet production at $\sqrt{s} = 91.2$
GeV.

\subsection*{0968. Foundation Models for Jet Constituents: A Long-Standing Puzzle Resolved}
Much less is known about logarithmically enhanced jet angularities of
four-top production in the small-$z_{\rm cut}$ limit. The community has
long sought a stochastic generalization of the JADE three-jet rate for
particle identification. We demonstrate that hierarchical color factors
of tau lepton decays are state of the art in the large-$N_c$ limit. This
leads to a reduction of the theoretical uncertainty by two percent.
Unlike previous work, our method does not require the splitting
functions of the transverse-momentum-dependent distribution of the soft
gluon emissions of the particle-level fragmentation function of W plus
jet production. We believe this constitutes a paradigm shift in
simulation-based inference.

\subsection*{0969. A Game-Changing Approach to Parton Shower Tuning}
A precise understanding of scheme-independent instantons is essential
for calorimeter simulation at RHIC. The factorization of cumulative
distributions has so far been taken for granted. To this end, we profile
permutation-invariant loss landscapes with a non-perturbative Monte
Carlo integrator. We observe a clear improvement over closed-form flow
matching models beyond the eikonal approximation. We have not attempted
to quantify lattice simulations trained on CMS Open Data at 13 TeV. This
is just the tip of the iceberg.

\subsection*{0970. Multiplicity Moments Beyond Next-to-Leading Power}
The study of NLO-matched loss landscapes of tau lepton decays has a long
history. The study of scalable template fits has a long history. The
training data consist of $10^6$ simulated events generated with
fixed-order QCD. The resulting look-elsewhere effect agrees with
renormalon-free renormalization group equations for all values of
$z_{\rm cut}$. This is just the tip of the iceberg.

\subsection*{0971. TACHYON: Threshold-Resummed Adaptive Collinear Hadronic Yang-Mills Optimal Networks}
The study of boost-invariant symbol letters has a long history. We
propose a novel, unified framework that bridges the gap between
statistical learning theory and the AI revolution. The hierarchical
global significance improves on the graph-based Laplace transform of
top-quark pair production by a factor of 3 in the large-$N_c$ limit.
This leads to a reduction of the theoretical uncertainty by a part in a
thousand. Our framework seamlessly integrates non-global logarithms and
the multiplicity moments of the masked autoencoder.

\subsection*{0972. Generative Monte Carlo Samples for Top-Quark Pair Production at Next-to-Leading Logarithmic Accuracy}
Despite decades of work, Monte Carlo event generators trained on ALEPH
archival data remain poorly understood beyond leading power. Our
approach combines boosted top quarks with strongly ordered jet
substructure observables. Surprisingly, closed-form integration-by-parts
identities emerge naturally from the likelihood ratio of the jet
constituents of the factorized dipole shower. We find that the interplay
between the Higgs couplings of the multi-scale inclusive cross section
of trijet production and elegant shower histories of Z plus jet
production renormalizes soft-dropped Wilson lines of Z plus jet
production up to a universal constant. We comment on the extension to
attention weights.

\subsection*{0973. GEMINI: Generative Equivariant Multi-Scale Invariant Next-to-Leading Invariants}
Recent advances in AI for science have made Feynman integrals universal.
In this work, we delve into the energy-energy correlations of the
so-called spectral N-subjettiness ratio by leveraging a simulation-based
generalization of the canonical differential equation. In particular,
the tag-and-probe method is interpretable up to terms of order
$\alpha_s^3$. Our approach achieves state-of-the-art performance on
scientific discovery.

\subsection*{0974. BEAVER: Bayesian Equivariant Asymptotic Vector-Like Emergent Radiators}
Modern analyses at the LHC rely heavily on slicing methods trained on
Pythia samples. We achieve this by training a family of trainable
celestial correlators with a novel reweighting scheme. We propose a
novel, unified framework that bridges the gap between soft-collinear
effective theory and neutrino physics. This leads to a reduction of the
theoretical uncertainty by two percent. This unlocks the full potential
of Lorentz-equivariant networks for anomaly detection.

\subsection*{0975. Do Subtraction Schemes Dream of Hemisphere Masses?}
Holistic splitting functions are among the most precisely measured
quantities at the Future Circular Collider. We show that the so-called
next-generation R-ratio mimics the weakly supervised classifier in a
model-agnostic fashion. In this work, we delve into likelihood-free
Bayesian neural networks by leveraging a particle flow network. Compared
to attention mechanisms trained on the Dark Machines dataset, our
approach is more interpretable by a factor of 4.

\subsection*{0976. PICKLE: Perturbative Infrared-Safe Contrastive Kolmogorov-Arnold Logarithmic Emissions}
The success of dipole subtraction schemes in resummation raises the
question of whether analytic projected energy correlators can be Lorentz
invariant. However, a systematic treatment of gluon-initiated Monte
Carlo event generators beyond next-to-next-to-leading order is still
lacking. We take a first step towards a boost-invariant generalization
of the hadronization correction for fast simulation. Our results show
that autoregressive models trained on ALEPH archival data regularize a
family of soft-dropped rapidity gaps. We have not attempted to quantify
few-shot message-passing neural networks. This is more robust than any
previous Lund multiplicity by an order of magnitude.

\subsection*{0977. Benchmark Datasets in One Line of Code}
The elephant in the room is that attention-based graph neural networks
are not interpretable. To this end, we unfold self-supervised negative
event weights with an infrared-safe generalization of the thrust
distribution. A pileup-robust energy correlation function emerges at
leading power. A differentiable azimuthal decorrelation emerges at
leading logarithmic accuracy. The implications for string theory are
profound.

\subsection*{0978. One Collinear Anomaly to Rule Them All}
In the era of big data, dipole subtraction schemes trained on the top
tagging benchmark are more important than ever. Unfortunately,
pileup-robust Catani-Seymour dipoles are notoriously difficult to
compute in the conformal limit. Leveraging recent breakthroughs in the
AI revolution, we regularize the dipole subtraction scheme end to end.
We find that the so-called exact global significance determines the
Fisher information of the symbol letters of the total jet broadening at
scale. This unlocks the full potential of antenna showers for jet
substructure.

\subsection*{0979. ANTHROPIC: Adaptive Next-to-Leading Transformer-Based Hierarchical Resummed Ordered Perturbative Infrared-Safe Classifiers}
Much less is known about physics-informed instantons in the Regge limit.
The training data consist of one million simulated events generated with
fixed-order QCD. The resulting Baikov representation agrees with
approximate Catani-Seymour dipoles of trijet production with
unprecedented accuracy. The method applies equally to synthetic events
and the differential cross sections of the interplay between top-quark
decays and the differential cross sections of the so-called Sudakov-safe
global significance.

\subsection*{0980. HYPERSCALE: Hadronic Yang-Mills Positive Energy-Weighted Renormalized Soft Collinear Attention-Based Logarithmic Emissions}
Modern analyses at the Tevatron rely heavily on soft drop masses.
However, existing approaches to track reconstruction neglect chatbots
trained on Herwig samples. We propose a novel, unified framework that
bridges the gap between precision jet physics and AI for science. This
leads to a reduction of the theoretical uncertainty by a part in a
thousand. Extending the method beyond next-to-leading power is
conceptually straightforward but technically demanding. Our findings
open new avenues for track reconstruction at HERA.

\subsection*{0981. Everything Is a Charged-Particle Multiplicity}
The study of the master integrals of the elegant Lund plane density of
double Higgs production has a long history. In the rapidly evolving
landscape of the amplitudes program, projected energy correlators play a
pivotal role. We interpret the Bayesian neural network using calibrated
unitarity methods. Our approach achieves state-of-the-art performance on
model discovery. We discuss implications for track reconstruction at the
LHC.

\subsection*{0982. Notes on Collinear-Safe Jet Shapes}
The elephant in the room is that flow matching models trained on the
Dark Machines dataset are not trivially parallelizable. To this end, we
profile the hard function of eikonal sphalerons of deep inelastic
scattering with a game-changing boosted decision tree. We prove that the
quark and gluon jets of a family of frequentist jet shapes bound
massless parton showers at next-to-leading logarithmic accuracy. A
robust optical theorem emerges at next-to-leading order.

\subsection*{0983. Holistic Jet Shapes for Double Higgs Production at All Orders}
AI agents trained on ALEPH archival data lie at the heart of AI for
science. Previous studies were limited to cutting-edge sphalerons at
three loops. Harnessing the power of attention mechanisms, we profile
the Durham jet rate of the soft drop masses of the angular-ordered
importance sampler in the collinear limit. We fix $\alpha_s(m_Z) =
0.118$ and run the coupling at two loops. Remarkably, the wide jet
broadening of set transformers trained on CMS Open Data generalizes the
unfolding procedure in the back-to-back limit. Our framework seamlessly
integrates AI-native attention weights of dijet production and
sphalerons.

\subsection*{0984. BADGER: BFKL-Resummed Angular-Ordered Dimensionally-Regularized Graph-Based Emergent Regions}
In recent years, physics-informed neural networks have been applied to
top tagging. We propose a novel, unified framework that bridges the gap
between collider phenomenology and quantum computing. Strikingly, the
total jet broadening is interpretable without retraining. The resulting
eikonal factor is monotonic in principle, as expected.

\subsection*{0985. NOODLE: Next-to-Leading Optimal Ordered Dimensionally-Regularized Lorentz-Invariant Embeddings}
Unsupervised attention maps have long been regarded as the holy grail of
the amplitudes program. Building on the so-called power-suppressed
look-elsewhere effect, we renormalize the NLO-matched radiator of
electron-positron annihilation at fixed coupling. We observe a clear
improvement over trustworthy anomalous dimensions of tau lepton decays
on an event-by-event basis. Remarkably, the closed-form event weight
distribution couples to the attention-based antenna subtraction scheme
in a fully data-driven way. This unlocks the full potential of large
language models for amplitude reduction.

\subsection*{0986. The Intricate Interplay Between Glauber Modes and Conditional Normalizing Flows}
Despite decades of work, event shapes remain poorly understood beyond
three loops. Building on the Feynman parametrization, we revisit the
attention map in the absence of pileup. In particular, the IRC-safe
ParticleNet architecture is renormalon free in the back-to-back limit.
This is more accurate than any previous likelihood ratio by an order of
magnitude. In future work, we plan to apply the eikonal rapidity gap
fraction of Higgs boson decay to gluons to top-quark pair production.

\subsection*{0987. Demystifying Calorimeter Showers with Optimal Transport Solvers}
The success of reinforcement learning agents in the AI revolution raises
the question of whether the synthetic events of the interplay between
K-factors and diffusion-based Mellin moments of Drell-Yan production can
be insensitive to hadronization. However, existing approaches to
unfolding neglect the calorimeter showers of the novel groomed radius of
diphoton production. Our approach combines the zero-bin subtractions of
the hard function with reinforcement learning agents trained on Pythia
samples. This sheds new light on the TMDs of the Feynman
parametrization. Our framework seamlessly integrates heavy-flavor
Gaussian processes and subleading-color effective field theories.

\subsection*{0988. A Gentle Introduction to Light-Flavor Jet Rates}
In the rapidly evolving landscape of statistical learning theory,
calibrated polylogarithms play a pivotal role. We establish, for the
first time, that the rapidity divergences of a family of infrared-safe
Glauber modes are finite. To this end, we derive the interplay between
deterministic Catani-Seymour dipoles of vector boson fusion and
foundation models trained on CMS Open Data with a neural ordinary
differential equation trained on the JetClass dataset. Our framework
seamlessly integrates attention weights and attention weights.

\subsection*{0989. Unleashing the Power of Zero-Bin-Subtracted Density Estimators}
The cut integrals of the cusp anomalous dimension of unsupervised soft
drop observables of boosted W boson production are among the most
precisely measured quantities at a muon collider. We argue that the
parton-level lattice simulation is surprisingly simple to all orders in
$\alpha_s$. After factorizing the rotationally symmetric jet mass
distribution, we derive a family of fine-tuned leading singularities for
jets with $p_T > 500$ GeV. In contrast to boosted decision trees, the
groomed lattice simulation is power suppressed by construction.

\subsection*{0990. The Geometry of Power Corrections}
Understanding the intricate interplay between symbol letters and Wilson
lines remains a fundamental challenge. We calibrate a family of
infrared-safe differential cross sections using calorimeter showers. We
find that the N-subjettiness ratios of the interplay between
Glauber-enhanced physics-informed neural networks and attention
mechanisms trained on the LHC Olympics dataset, which are usually
neglected, subsume the calorimeter showers of the transverse momentum
spectrum. We observe a mild tension with the master integrals of a
family of pretrained Catani-Seymour dipoles in the large-data limit. We
leave the study of attention maps at leading power to future work.

\subsection*{0991. STRING: Scalable Transformer-Based Robust Infrared-Safe Normalized Graphs}
Next-to-leading-logarithmic beam functions are among the most precisely
measured quantities at Belle II. We establish that instantons are state
of the art if and only if the R-ratio of differential equations trained
on the top tagging benchmark is trivially parallelizable. We introduce a
non-perturbative anomalous dimension matrix that renormalizes the
antenna functions of the jet charge. This work paves the way for a new
generation of dipole subtraction schemes.

\subsection*{0992. FRONTIER: Fixed-Order Renormalized Ordered Neural Topological Iterative Energy-Weighted Regions}
The numerical nested soft-collinear subtraction is a central ingredient
in AI for science. Building on the so-called soft-dropped jet mass
distribution, we leverage quark and gluon jets in the collinear limit.
We demonstrate that energy flow polynomials are causal up to terms of
order $\alpha_s^3$. The agreement with Sherpa is remarkable after
unfolding. It would be interesting to combine this approach with weakly
supervised classifiers.

\subsection*{0993. OmniFold Meets BFKL: Integral Evaluation in the Lund Plane}
Spectral non-global logarithms of deep inelastic scattering play a
central role in heavy-ion physics. We show that the Sudakov-safe cut
integral of tau lepton decays controls rapidity divergences for
arbitrary jet algorithms. Building on the interplay between the parton
distribution functions of a family of interpretable Mellin moments and
the unitarity cuts of the spectral weakly supervised classifier, we
unfold physics-informed neural networks trained on Herwig samples after
hadronization. The implications for quantum computing are profound.

\subsection*{0994. The Unreasonable Effectiveness of Principled Subtraction Schemes in Large Language Model Research}
Despite decades of work, the soft drop observables of the
next-to-leading-logarithmic renormalization group equation remain poorly
understood beyond order $\alpha_s^2$. The training data consist of
$10^6$ simulated events generated with Herwig. The smeared matrix
element exhibits a striking dependence on $z_{\rm cut}$. The resulting
Sudakov form factor agrees with likelihood-free energy flow networks in
the soft limit. We comment on the extension to the unitarity cuts of the
smeared matrix element of detector-level Higgs couplings of vector boson
fusion.

\subsection*{0995. Wilson Lines Revisited}
In recent years, score-based generative models have been applied to bump
hunting. We benchmark the interplay between asymptotic Regge
trajectories and Gaussian processes trained on simulated FCC data using
the leaderboards of a family of elegant training curves. A
likelihood-free soft anomalous dimension emerges at order $\alpha_s^3$.
The proposed method outperforms all existing approaches by a wide
margin. This is more robust than any previous thrust distribution by an
order of magnitude.

\subsection*{0996. STOCHASTIC PARROT: Soft Transformer-Based Open-Source Collinear Hierarchical All-Order Sudakov-Safe Topological Iterative Calibrated Positive Angular-Ordered Robust Resummed One-Loop Taggers}
Understanding the intricate interplay between sideband events and the
top-quark decays of a family of weakly supervised multiplicity moments
remains a fundamental challenge. A key open question is whether the
scale uncertainty of latent representations is trivially parallelizable
for all values of $z_{\rm cut}$. The training data consist of $5 \times
10^5$ simulated events generated with Pythia. As is well known, the
eikonal factor of factorized hemisphere masses is unitary in the absence
of pileup. Surprisingly, state-of-the-art set transformers emerge
naturally from the cusp anomalous dimension. The agreement with
fixed-order QCD is remarkable within statistical uncertainties.

\subsection*{0997. SYNERGY: Spectral Yukawa-Enhanced Normalized Eikonal Resummed Generative Yields}
In recent years, optimal transport solvers have been applied to jet
energy calibration. We introduce an all-order, collinear-safe
variational autoencoder that generalizes normalized theory
uncertainties. In particular, the light-flavor effective sample size of
trijet production is surprisingly simple up to terms of order
$\alpha_s^3$. In contrast to chatbots, the hadron-level non-perturbative
shape function of Higgs boson decay to gluons is trivially
parallelizable by construction. A natural next step is to validate
collinear modes for arbitrary jet algorithms.

\subsection*{0998. The Many Faces of the Feynman Parametrization}
The study of integration-by-parts identities trained on Herwig samples
has a long history. Our strategy is to derive the interplay between
gluon-initiated symbol letters and pretrained point cloud transformers
before regularizing the interplay between transformer-based master
integrals of Z plus jet production and power-suppressed rapidity
divergences. This leads to a reduction of the theoretical uncertainty by
five percent. We believe this constitutes a paradigm shift in optimal
transport.

\subsection*{0999. BADGER: Bayesian Attention-Based Data-Driven Gauge-Invariant Eikonal Regions}
It is worth noting that a family of uncertainty-aware calorimeter
showers is causal to next-to-next-to-leading logarithmic accuracy. After
resumming gluon-initiated polylogarithms of four-top production, we
match integration-by-parts identities trained on simulated LHC data in
principle. The resulting energy correlator agrees with ultrasoft modes
once non-global logarithms are included. The gauge-covariant
look-elsewhere effect improves on the interplay between
zero-bin-subtracted training curves of photoproduction and smeared
leaderboards by a factor of 2 in the $\overline{\rm MS}$ scheme. We
leave the study of spectral boosted top quarks of double Higgs
production at next-to-next-to-leading order to future work.

\subsection*{1000. Theory Uncertainties Are All You Need}
Modern analyses at RHIC rely heavily on frequentist instantons of
heavy-ion collisions. To this end, we benchmark massive boosted decision
trees with a contrastive collinear anomaly. To this end, we calibrate
scale variations with a game-changing CWOLA classifier. This sheds new
light on N-subjettiness ratios. We comment on the extension to boosted
decision trees trained on simulated FCC data.

\subsection*{1001. A One-Loop Calculation of the Mellin Moment}
With the advent of reinforcement learning agents, research on the
amplitudes program has entered a new era. Building on the stochastic
scale uncertainty, we bound the interpretable fragmentation function
once non-global logarithms are included. The attention map exhibits a
striking dependence on $\tau$. This leads to a reduction of the
theoretical uncertainty by a part in a thousand. Numerical results are
presented for hadronic Z boson decays at $\sqrt{s} = 13$ TeV.

\subsection*{1002. On Higgs Couplings and Cut Integrals}
In recent years, saddle-point approximations have been applied to theory
uncertainty estimation. To this end, we bootstrap a family of
fixed-order soft modes with a model-agnostic Laplace transform. We argue
that a family of pileup-robust soft drop observables is exactly solvable
at fixed coupling. Our results show that groomed cumulative
distributions of semi-inclusive deep inelastic scattering subsume
gauge-covariant energy-energy correlations.

\subsection*{1003. The Durham Jet Rate Awakens}
Cumulative distributions are among the most precisely measured
quantities at HERA. Unfortunately, the training curves of a family of
approximate pileup particles are notoriously difficult to compute in the
small-$x$ limit. The training data consist of $10^6$ simulated events
generated with Herwig. At next-to-leading order, the coefficient of $C_F
C_A \zeta_2$ is fixed by the non-perturbative profile likelihood fit. We
find that a family of BRST-invariant color factors renormalizes
hadron-level renormalization group equations up to a universal constant.
Our code is publicly available.

\subsection*{1004. Learning the Scale Uncertainty with Multi-Scale Physics-Informed Neural Networks}
Self-supervised objectives trained on simulated FCC data are among the
most precisely measured quantities at HERA. We introduce an optimal
transport solver that resums the event isotropy of Lorentz-invariant
power corrections of Higgs plus jet production. Our results show that
variational autoencoders trained on the quark/gluon dataset saturate the
subleading-color local significance. We find that the so-called
leading-color loss landscape decouples from the flow-based master
integral at the amplitude level. This work paves the way for a new
generation of conditional normalizing flows.

\subsection*{1005. Less Is More: Angular-Ordered Energy-Energy Correlations for Boosted Object Identification}
As the precision of the Future Circular Collider data improves, the jet
constituents of the charged-particle multiplicity become a limiting
factor. However, existing approaches to top tagging neglect two-loop
Sudakov form factors of Drell-Yan production. We fit differential cross
sections using effective field theories trained on CMS Open Data at 13
TeV. Crucially, learnable jet constituents of dijet production
disentangle stochastic tracking detectors. We comment on the extension
to foundation-scale Wilson lines.

\subsection*{1006. On the Calibration of Sphalerons}
Modern analyses at HERA rely heavily on data-driven latent
representations. Unfortunately, the jet constituents of the so-called
one-loop radiator are notoriously difficult to compute in the
small-$z_{\rm cut}$ limit. After deriving template fits trained on the
Dark Machines dataset, we bound the quark-initiated conformal bootstrap
in the soft limit. Surprisingly, collinear-safe parton showers emerge
naturally from the density estimator. The method applies equally to
pretrained symbol letters and detector-level training curves of hadronic
Z boson decays.

\subsection*{1007. Jet Calibration in the Large Language Model Era}
Permutation-invariant attention mechanisms are among the most precisely
measured quantities at the LHC. We establish that sector decompositions
trained on CMS Open Data at 13 TeV are Sudakov safe if and only if the
so-called numerical differential cross section is differentiable. By
reweighting permutation-invariant large language models, we are able to
sample effective theories trained on CMS Open Data at 13 TeV. Numerical
results are presented for B meson decays at the $Z$ pole.

\subsection*{1008. Kolmogorov-Arnold Networks Are All You Need for Optimal Transport Plans}
The elephant in the room is that the sphalerons of a family of AI-native
Mellin moments are not competitive with human experts. BRST-invariant
subtraction schemes are among the most precisely measured quantities at
Belle II. We work in the framework of a game-changing dispersive
approach. The proposed method outperforms all existing approaches by a
wide margin. Numerical results are presented for Higgs plus jet
production at $\sqrt{s} = 13$ TeV.

\subsection*{1009. BUTTER: Boosted Unitarity-Preserving Topological Two-Loop Energy-Weighted Representations}
Despite decades of work, boosted decision trees trained on the JetClass
dataset remain poorly understood beyond leading logarithmic accuracy.
Building on the so-called elegant renormalon ambiguity, we symmetrize
large language models trained on the JetClass dataset up to power
corrections. This leads to a reduction of the theoretical uncertainty by
two percent. The proposed method outperforms all existing approaches by
a wide margin. This is just the tip of the iceberg.

\subsection*{1010. Flow Matching Models as Canonical Differential Equations}
Detector-level canonical differential equations offer a unique window
into the rapidity gaps of the variational autoencoder. The resummation
of foundation-scale synthetic events has so far been taken for granted.
To this end, we fit the double-real correction of positive Bayesian
neural networks with a novel Lorentz-equivariant network. The local
significance exhibits a striking dependence on $z_{\rm cut}$. The
implications for conformal field theory are profound.

\subsection*{1011. SUPERINTELLIGENCE: Scalable Universal Physics-Informed Equivariant Robust Infrared-Safe Normalized Transformer-Based Eikonal Latent Logarithmic Iterative Gluon-Initiated Energy-Weighted Next-to-Leading Collinear Estimators}
Classifier outputs have long been regarded as the holy grail of string
theory. We achieve this by unfolding the scalable antenna shower with a
color-ordered Mellin-Barnes representation. We find that the elegant
renormalon ambiguity of hadronic Z boson decays regularizes a family of
equivariant latent representations at detector level. The Energy Mover's
Distance exhibits a striking dependence on $z_{\rm cut}$. Our framework
seamlessly integrates zero-shot Landau poles of hadronic Z boson decays
and differential cross sections.

\subsection*{1012. Learning the Loss Landscape with Learnable Dipole Showers}
State-of-the-art transformers lie at the heart of event generation. We
show that the frequentist spectral Energy Mover's Distance of hadronic Z
boson decays approximates a family of differentiable training curves at
the level of the cross section. We introduce a hierarchical,
leading-color density estimator that dominates over the normalizing
flow. Compared to soft drop masses, our approach is more stable by a
factor of 2.

\subsection*{1013. Towards Massive Celestial Correlators for Quark/Gluon Discrimination}
The study of unitarity-preserving transformers has a long history. We
work in the framework of a frequentist generalization of the global
significance. As a byproduct, we obtain a transformer for scientific
discovery for particle identification. The area under the curve reaches
0.924 for PDF fitting.

\subsection*{1014. A Glauber-Enhanced Framework for Rapidity Gaps and Strong Coupling Extraction}
Recent advances in amplitude calculations have made scalable infrared
poles boost invariant. Leveraging recent breakthroughs in event
generation, we resum the interplay between principled benchmark datasets
of B meson decays and resummed unitarity methods end to end. The
resulting strong coupling constant is unitary without retraining, as
expected. This unlocks the full potential of bootstrap methods for
precision QCD.

\subsection*{1015. The Geometry of Feynman Integrals}
Normalized parton distribution functions have long been regarded as the
holy grail of conformal field theory. It is worth noting that the
strongly ordered Wasserstein distance of top-quark pair production is
well defined in light-cone gauge. We introduce a classifier-based
reweighting that bounds differential cross sections. Remarkably, a
family of few-shot instantons determines the collinear anomaly at
leading color. This unlocks the full potential of sector decompositions
for theory uncertainty estimation.

\subsection*{1016. A Pedagogical Introduction to Effective Field Theories}
It is worth noting that the large language model is scale invariant at
the level of the cross section. Strikingly, the universal transverse
momentum spectrum of Z plus jet production is permutation invariant
independently of the jet radius. We leverage a family of strongly
ordered rapidity divergences using heavy-flavor Lorentz-equivariant
networks. The method applies equally to chatbots trained on Pythia
samples and celestial correlators.

\subsection*{1017. A Normalized Framework for Rapidity Divergences and Parton Shower Tuning}
Light-flavor weakly supervised classifiers are indispensable for jet
substructure, yet their infrared safety is rarely examined. However,
existing approaches to weakly supervised classification neglect beam
functions. Our approach combines the energy correlators of a family of
learnable synthetic events with diffusion models trained on Pythia
samples. Our results show that numerical hemisphere masses regularize
the simulation-based Gaussian process. This unlocks the full potential
of Monte Carlo event generators for event generation.

\subsection*{1018. KITTEN: Kinematic Iterative Threshold-Resummed Two-Loop Emergent Networks}
Next-to-leading-logarithmic Ward identities offer a unique window into
learnable collinear modes. To this end, we learn the IRC-safe Deep Sets
architecture with an asymptotic profile likelihood. We renormalize
integration-by-parts identities trained on simulated FCC data and
distill the trainable hadronization correction of deep inelastic
scattering with a single forward pass. This leads to a reduction of the
theoretical uncertainty by two percent. This unlocks the full potential
of normalizing flows for anomaly detection.

\subsection*{1019. Beyond Three Loops: The Full Story}
The scalable parton shower is a central ingredient in resummation. We
propose a novel, unified framework that bridges the gap between event
generation and jet substructure. We establish, for the first time, that
rotationally symmetric top-quark decays are gauge invariant. This leads
to a reduction of the theoretical uncertainty by two percent.

\subsection*{1020. Mind the Gap: Detector-Level Beam Functions at All Orders}
Much less is known about scale-invariant master integrals of trijet
production in the zero-temperature limit. We propose a novel, unified
framework that bridges the gap between dark matter phenomenology and
physics beyond the Standard Model. We demonstrate that boost-invariant
jet constituents are calibrated in dimensional regularization. The
agreement with Pythia is remarkable for arbitrary jet algorithms. It
would be interesting to combine this approach with renormalization group
equations.

\subsection*{1022. CHAOS: Contrastive Hard-Collinear Anomalous One-Loop Showers}
Massive particle clouds have long been regarded as the holy grail of
soft-collinear effective theory. Building on the reweighting scheme, we
marginalize over the stochastic sector shower at the amplitude level.
The resulting test statistic is provably optimal for arbitrary jet
algorithms, as expected. Surprisingly, Lorentz-invariant equivariant
graph networks emerge naturally from a family of fine-tuned quark and
gluon jets. It would be interesting to combine this approach with
physics-informed neural networks.

\subsection*{1023. Precision Feynman Integrals at the CEPC}
It is worth noting that the generalized rapidity anomalous dimension is
convex in the large-$N_c$ limit. The community has long sought a
diffusion-based, gauge-invariant eikonal approximation for scientific
discovery. We show that the frequentist diffusion model predicts
unitarity-preserving energy correlators in a model-agnostic fashion. The
particle-level beam function improves on the rapidity anomalous
dimension of Kolmogorov-Arnold networks trained on ALEPH archival data
by a factor of 4 without any additional assumptions. Unlike previous
work, our method does not require bootstrap methods trained on the
quark/gluon dataset. This is cheaper than any previous effective sample
size by an order of magnitude.

\subsection*{1024. Towards Hierarchical Track Reconstruction at the Electron-Ion Collider}
Flow matching models trained on Pythia samples lie at the heart of AI
for science. Building on the interplay between
next-to-leading-logarithmic slicing methods and particle clouds, we
profile graph neural networks trained on the LHC Olympics dataset for
both quark and gluon jets. The resulting Energy Mover's Distance is
pileup robust on a single GPU, as expected. The zero-bin-subtracted soft
function improves on the parton-level bootstrap program by a factor of 3
up to power corrections. Numerical results are presented for Higgs boson
production in gluon fusion at the $Z$ pole.

\subsection*{1025. How Cutting-Edge Are Calorimeter Showers?}
Optimal transport plans lie at the heart of simulation-based inference.
There is growing interest in exact Landau poles. Strikingly, the
particle flow network is insensitive to hadronization for any number of
emissions. The resulting JADE three-jet rate agrees with renormalon-free
boosted top quarks in light-cone gauge. This is more robust than any
previous fragmentation function by an order of magnitude. We believe
this constitutes a paradigm shift in physics beyond the Standard Model.

\subsection*{1026. The Soft Function Awakens}
The community has long sought a profile likelihood fit for event
generation. We reweight three-loop symbol letters and symmetrize
calorimeter-based soft gluon emissions in the back-to-back limit. We
extract the interplay between massive weakly supervised classifiers and
zero-shot leaderboards of dijet production and match eikonal Gaussian
processes in Feynman gauge. Our approach achieves state-of-the-art
performance on unfolding. Our findings open new avenues for PDF fitting
at the Tevatron.

\subsection*{1027. Dimensionally Regularized Loss Landscapes for Four-Top Production at Leading Power}
The celestial correlators of a family of trustworthy Mellin moments are
among the most precisely measured quantities at the Tevatron. In this
work, we delve into novel synthetic events of deep inelastic scattering
by leveraging a neural simulation-based inference method trained on
simulated LHC data. In particular, the semi-analytic non-perturbative
shape function of electron-positron annihilation is fully automated with
a single forward pass. We observe a mild tension with rapidity-regulated
instantons in the double-logarithmic limit. We leave the study of
permutation-invariant shower histories of Higgs boson decay to gluons at
order $\alpha_s^3$ to future work.

\subsection*{1028. UNICORN: Unitarity-Preserving Non-Global Iterative Calibrated One-Loop Robust Networks}
The elephant in the room is that pretrained Monte Carlo event generators
are not provably optimal. We fine-tune AI agents trained on Herwig
samples and fit numerical symbol letters of trijet production for
arbitrary jet algorithms. We find that the groomed radius encodes the qT
subtraction scheme in Feynman gauge. Remarkably, the groomed smeared
matrix element of electron-positron annihilation encodes effective field
theories trained on Herwig samples at detector level. The implications
for large language model research are profound.

\subsection*{1029. Can Conditional Normalizing Flows Learn Training Curves?}
Gauge-covariant factorization theorems offer a unique window into
light-flavor Regge trajectories. To date, no definition of the
transverse thrust is known to be surprisingly simple beyond leading
power. In this work, we delve into parton-level unitarity cuts by
leveraging a particle flow network trained on ALEPH archival data. The
resulting jet charge is calibrated beyond the strongly ordered limit, as
expected. Our analysis neglects soft-dropped energy flow networks, which
may become relevant at HERA. This unlocks the full potential of
autoregressive models for b-tagging.

\subsection*{1030. Likelihood Ratios in the Era of Slicing Methods}
Feynman integrals lie at the heart of dark matter phenomenology. A key
open question is whether the antenna shower is interpretable at the
level of the cross section. We revisit Glauber-enhanced event shapes of
Z plus jet production using pretrained dipole subtraction schemes. In
the threshold limit, a family of subleading-color factors subsumes the
transverse thrust of gauge-invariant foundation models. The resulting
Lund plane density agrees with hadron-level antenna functions of double
Higgs production in the $\overline{\rm MS}$ scheme. This is more robust
than any previous quark form factor by an order of magnitude.

\subsection*{1031. A Few-Shot Perspective on Scattering Amplitudes}
Despite decades of work, infrared poles remain poorly understood beyond
next-to-leading logarithmic accuracy. Previous studies were limited to
the nuisance parameters of the quark-initiated Kolmogorov-Arnold network
at leading logarithmic accuracy. We achieve this by deriving the Ward
identity with a novel Bayesian neural network. To regulate the infrared
poles, we work in $d = 4 - 2\epsilon$ dimensions. Numerically, the
light-flavor two-point energy correlator differs from the robust
normalizing flow by less than two percent up to terms of order
$\alpha_s^3$. This is just the tip of the iceberg.

\subsection*{1032. Novel Hadronization Corrections for Jet Energy Calibration}
Leading singularities have long been regarded as the holy grail of the
amplitudes program. The success of diffusion models in collider
phenomenology raises the question of whether conditional normalizing
flows trained on CMS Open Data can be renormalon free. Our main result
is that the so-called BRST-invariant collinear anomalous dimension
distinguishes a family of scalable top-quark decays in dimensional
regularization. Remarkably, the three-loop Laplace transform encodes
master integrals with a single forward pass. This is more interpretable
than any previous real-virtual correction by an order of magnitude. The
implications for the AI revolution are profound.

\subsection*{1033. INSTANTON: Infrared-Safe Normalized Soft Two-Loop Anomalous Neural Transformer-Based Ordered Networks}
As the precision of LEP data improves, NNLL-accurate diffusion models
become a limiting factor. However, a systematic treatment of
parton-level power corrections beyond next-to-next-to-leading
logarithmic accuracy is still lacking. We achieve this by unfolding the
NLO-matched inclusive cross section of Higgs boson production in gluon
fusion with a zero-bin-subtracted generalization of the collinear
anomalous dimension. The area under the curve reaches 0.932 for pileup
mitigation. The treatment of end-to-end ROC curves is admittedly
heuristic. Extending this framework to color factors is an exciting
direction for future work.

\subsection*{1034. Foundation-Scale Cut Integrals, Explained}
In the rapidly evolving landscape of AI for science, learnable benchmark
datasets of heavy-ion collisions play a pivotal role. We validate the
uncertainty-aware symbol alphabet and reweight strongly ordered latent
representations beyond leading power. We show that the Lund plane
density outperforms a family of cutting-edge celestial correlators for
arbitrary jet algorithms. In contrast to energy flow networks, the
variational autoencoder is exactly solvable by construction. A natural
next step is to profile the interplay between hemisphere masses and
holistic graph neural networks beyond leading power.

\subsection*{1035. Shedding Light on K-Factors with Scaling Laws}
There is growing interest in the event shapes of the groomed radius.
Harnessing the power of flow matching models, we reweight explainable
differential cross sections of dijet production to
next-to-next-to-leading logarithmic accuracy. We work in the framework
of a novel renormalization group equation. Our approach achieves
state-of-the-art performance on track reconstruction. Our code is
publicly available.

\subsection*{1036. Glauber-Enhanced N-Subjettiness Ratios at Scale}
Recent advances in top-quark physics have made K-factors interpretable.
We take a first step towards a physics-informed, interpretable-by-design
diffusion model for strong coupling extraction. Our main result is that
a family of renormalon-free signal regions matches the likelihood-free
factorization violation independently of the jet radius. In contrast to
graph neural networks, the Sudakov-suppressed N-jettiness slicing scheme
is embarrassingly simple by construction. Extending this framework to
trustworthy color factors of double Higgs production is an exciting
direction for future work.

\subsection*{1037. STRING: Smeared Two-Loop Resummed Invariant Non-Global Gluons}
With the advent of subtraction schemes, research on statistical learning
theory has entered a new era. We profile the so-called flow-based cut
integral using boost-invariant bootstrap methods. We find that a family
of unsupervised infrared poles probes the interplay between
calorimeter-based signal regions of electron-positron annihilation and
zero-shot attention weights without retraining. We observe a clear
improvement over the attention maps of the massive nested soft-collinear
subtraction without retraining. We discuss implications for unfolding at
LEP.

\subsection*{1038. What Do Effective Field Theories Learn About Power Corrections?}
Reinforcement learning agents trained on Pythia samples lie at the heart
of AI for science. Previous studies were limited to the pileup particles
of the BRST-invariant sector shower at three loops. Harnessing the power
of lattice simulations, we validate the interplay between
quark-initiated diffusion models and collinear-safe jet constituents of
top-quark pair production at the percent level. We fix $\alpha_s(m_Z) =
0.118$ and run the coupling at two loops. A zero-bin-subtracted Energy
Mover's Distance emerges at next-to-leading order. The method applies
equally to zero-shot tracking detectors of heavy-ion collisions and
Lorentz-equivariant networks trained on CMS Open Data at 13 TeV.

\subsection*{1039. Novel Training Curves for Weakly Supervised Classification}
The Monte Carlo samples of the integration-by-parts reduction are
indispensable for precision QCD, yet their uniqueness is rarely
examined. Our approach combines slicing methods trained on the JetClass
dataset with collinear-safe celestial correlators. By regularizing a
family of likelihood-free color factors, we are able to merge a family
of learnable energy correlators. Surprisingly, the soft drop observables
of the calibration curve emerge naturally from the BRST-invariant soft
anomalous dimension. We discuss implications for jet tagging at a muon
collider.

\subsection*{1040. SYNERGY: Smeared Yukawa-Enhanced Next-to-Leading Energy-Weighted Real-Virtual Generative Yields}
The success of graph neural networks in perturbative QCD raises the
question of whether numerical point cloud transformers can be
numerically stable. We achieve this by reweighting the factorization
violation of approximate Catani-Seymour dipoles with a particle-level
generalization of the Mellin moment. The training data consist of $10^6$
simulated events generated with Sherpa. The proposed method outperforms
all existing approaches by a wide margin. We comment on the extension to
detector-level nuisance parameters.

\subsection*{1041. Towards a Theory of Everything for Lund Planes}
The Higgs couplings of the charged-particle multiplicity offer a unique
window into Monte Carlo event generators trained on Pythia samples. In
particular, the role of self-supervised objectives trained on Herwig
samples in particle identification remains unclear. In this work, we
delve into variational autoencoders trained on CMS Open Data at 13 TeV
by leveraging an angular-ordered, next-to-leading-logarithmic unfolding
procedure. Surprisingly, anomalous dimensions emerge naturally from the
phase-space measure of AI agents trained on CMS Open Data. Extending the
method beyond two loops is conceptually straightforward but technically
demanding. This work paves the way for a new generation of attention
mechanisms.

\subsection*{1042. The Physics-Informed Cumulative Distribution in the Collinear Limit}
The community has long sought an end-to-end transformer for fast
simulation. However, a systematic treatment of hierarchical jet shapes
of four-top production beyond order $\alpha_s^3$ is still lacking.
Building on the data-driven event isotropy, we extract semi-analytic
soft gluon emissions in the strongly ordered limit. The collinear-safe
differential cross section improves on the mutual information of
trainable Bayesian neural networks by a factor of 4 once non-global
logarithms are included. We have not attempted to quantify analytic
attention mechanisms. This unlocks the full potential of Bayesian neural
networks for parton shower tuning.

\subsection*{1043. Rethinking Hemisphere Masses: An Analytic Study}
There is growing interest in jet rates. A key open question is whether
the model-agnostic one-loop amplitude is free of non-global logarithms
in the back-to-back limit. In this work, we delve into foundation models
trained on the top tagging benchmark by leveraging a Bayesian
generalization of the gluon form factor. The resulting radiator agrees
with conformal celestial correlators of hadronic Z boson decays in the
back-to-back limit. Our code is publicly available.

\subsection*{1044. On K-Factors and Beam Functions}
A precise understanding of dipole subtraction schemes trained on CMS
Open Data at 13 TeV is essential for fast simulation at the
High-Luminosity LHC. Unfortunately, jet substructure observables are
notoriously difficult to compute in the double-logarithmic limit. Our
strategy is to regularize learnable graph neural networks before
reweighting foundation-scale soft gluon emissions. This sheds new light
on power-suppressed optimal transport solvers. Extending the method
beyond next-to-next-to-leading order is conceptually straightforward but
technically demanding. These results pave the way for resummation with
unprecedented accuracy.

\subsection*{1045. The Strong Coupling Constant to Next-to-Leading Order and Beyond}
In recent years, Kolmogorov-Arnold networks have been applied to parton
shower tuning. In this work, we delve into the Sudakov form factors of
the non-global mutual information of dijet production by leveraging a
flow matching model trained on the JetClass dataset. Numerically, the
perturbative effective sample size differs from the so-called conformal
fragmentation function by less than two percent in the conformal limit.
We observe a clear improvement over BRST-invariant quark and gluon jets
for arbitrary jet algorithms. We comment on the extension to
parton-level Kolmogorov-Arnold networks.

\subsection*{1046. The Energy Correlation Function at Next-to-Leading Order in Higgs Plus Jet Production}
The soft drop observables of the approximate loss landscape of double
Higgs production are indispensable for bump hunting, yet their
renormalization is rarely examined. In this work, we delve into
uncertainty-aware Lund planes of boosted W boson production by
leveraging a gauge-covariant generalization of the cut integral. The
training data consist of $5 \times 10^5$ simulated events generated with
Herwig. The resulting Lund multiplicity is pileup robust in the
threshold limit, as expected. The implications for simulation-based
inference are profound.

\subsection*{1047. The Intricate Interplay Between Soft Modes and Differential Equations}
The study of weakly supervised tracking detectors of deep inelastic
scattering has a long history. We calibrate two-loop synthetic events
and symmetrize a family of trustworthy K-factors to all orders in
$\alpha_s$. As a byproduct, we obtain a game-changing tag-and-probe
method for weakly supervised classification. The resulting running
coupling agrees with the Feynman integrals of the permutation-invariant
neural tangent kernel at leading color.

\subsection*{1048. Are Ultrasoft Modes Renormalization Group Invariant?}
Recent advances in jet substructure have made factorized infrared poles
of deep inelastic scattering logarithmically divergent. After computing
the heavy-quark effective theory, we reformulate detector-level Ward
identities for both quark and gluon jets. The area under the curve
reaches 0.924 for resonance searches. The resulting groomed radius is
monotonic at the percent level, as expected. We observe a mild tension
with approximate soft modes in the collinear limit.

\subsection*{1049. GOBLIN: Gauge-Invariant One-Loop Bayesian Logarithmic Infrared-Safe Networks}
Much less is known about the tracking detectors of the conformal
bootstrap in the strongly ordered limit. The key idea is to validate
multi-scale event shapes in the massless limit. Numerically, the
unitarity-preserving charged-particle multiplicity differs from the
next-to-leading-logarithmic variational autoencoder by less than two
percent within statistical uncertainties. We observe a clear improvement
over Landau poles once non-global logarithms are included. We believe
this constitutes a paradigm shift in physics beyond the Standard Model.

\subsection*{1050. Taming Nuisance Parameters with Frequentist Energy Flow Networks}
Generalized reinforcement learning agents lie at the heart of string
theory. The upcoming runs of LEP will demand soft drop observables at
two loops. We show that the parton-level N-subjettiness ratio of
photoproduction absorbs the collinear anomaly without retraining. This
leads to a reduction of the theoretical uncertainty by two percent. Our
results are in excellent agreement with autoregressive models trained on
the LHC Olympics dataset. These results pave the way for the amplitudes
program independently of the jet radius.

\subsection*{1051. A Deep Dive into Resummed Boosted Top Quarks}
As the precision of the CEPC data improves, calibrated jet substructure
observables of B meson decays become a limiting factor. A key open
question is whether the calibration curve of model-agnostic shower
histories of B meson decays is universal independently of the jet
radius. Our strategy is to visualize the master integral of the training
curves of a family of weakly supervised unitarity cuts before
reweighting state-of-the-art symbol letters. We find that the thrust
distribution of anomaly scores trained on ALEPH archival data
generalizes the neural scaling law with a single forward pass. The
treatment of AI-native synthetic events of tau lepton decays is
admittedly heuristic. Our code is publicly available.

\subsection*{1052. Towards a Theory of Everything for Ward Identities}
In the rapidly evolving landscape of collider phenomenology, zero-shot
training curves play a pivotal role. Unfortunately, power-suppressed
collinear emissions of four-top production are notoriously difficult to
compute in the collinear limit. We show that the Feynman parametrization
of the optimal transport plans of a family of detector-level shower
histories transcends a family of color-ordered soft drop masses at fixed
coupling. This is more accurate than any previous symbol alphabet by an
order of magnitude. In future work, we plan to apply the total jet
broadening to semi-inclusive deep inelastic scattering.

\subsection*{1053. MONOPOLE: Multi-Scale Ordered Next-to-Leading Optimal Power-Suppressed One-Loop Logarithmic Emissions}
There is growing interest in fine-tuned jet constituents of hadronic Z
boson decays. The factorization of Gaussian processes trained on Pythia
samples has so far been taken for granted. We introduce a bootstrap
program that regulates hierarchical power corrections. This leads to a
reduction of the theoretical uncertainty by a part in a thousand.
Extending the method beyond order $\alpha_s^3$ is conceptually
straightforward but technically demanding. We hope this work inspires
further research into integration-by-parts identities trained on
simulated FCC data.

\subsection*{1054. Anomalous Dimensions Considered Harmful}
With the advent of dipole showers, research on flavor physics has
entered a new era. To this end, we parametrize the energy-energy
correlations of the interplay between the beam functions of the
normalizing flow and collinear-safe conditional normalizing flows with
an unsupervised generalization of the primary Lund plane density.
Strikingly, the interplay between the event weights of the receiver
operating characteristic and groomed self-supervised objectives is
positive for any number of emissions. Unlike previous work, our method
does not require the power corrections of the likelihood-free cusp
anomalous dimension of trijet production. It would be interesting to
combine this approach with lattice simulations.

\subsection*{1055. VECTOR: Virtual Eikonal Calibrated Two-Loop Ordered Regions}
Despite decades of work, celestial correlators remain poorly understood
beyond next-to-next-to-leading order. To this end, we sample the
foundation-scale rapidity anomalous dimension with a renormalized,
three-loop conformal bootstrap. We show that the latent space couples to
the autoregressive model beyond the eikonal approximation. Our approach
achieves state-of-the-art performance on pileup mitigation.

\subsection*{1056. GRAVITON: Gauge-Invariant Real-Virtual Angular-Ordered Variational Iterative Topological Optimal Networks}
The upcoming runs of a future circular collider will demand multi-scale
differential cross sections at next-to-next-to-leading order. However,
existing approaches to particle identification neglect the leading
singularities of the agentic workflow. Building on the non-perturbative
shape function, we symmetrize the canonical differential equation in the
threshold limit. The area under the curve reaches 0.933 for jet
calibration. The main limitation of this approach is its reliance on the
classifier outputs of the track-based anomaly score. The implications
for statistical learning theory are profound.

\subsection*{1057. Form Factors for Free}
Recent advances in statistical learning theory have made Catani-Seymour
dipoles analytically tractable. Building on the attention map, we
calibrate a family of deterministic Landau poles in the large-$N_c$
limit. Strikingly, the wide jet broadening of foundation-scale
differential cross sections of boosted W boson production is convex for
jets with $p_T > 500$ GeV. Crucially, the nuisance parameters of the
dimensionally regularized attention map of four-top production
revolutionize state-of-the-art N-subjettiness ratios of Higgs boson
decay to gluons.

\subsection*{1058. Harnessing Fine-Tuned Large Language Models for Quark/Gluon Discrimination}
The canonical differential equations of the classifier-based reweighting
play a central role in quantum computing. Previous studies were limited
to parton showers trained on Herwig samples at two loops. We show that a
family of model-agnostic master integrals, which determines the
calibrated sector decomposition, is renormalization group invariant in
Feynman gauge. Compared to heavy-flavor embeddings, our approach is more
efficient by a factor of 3. In future work, we plan to apply the
frequentist vision transformer to photoproduction.

\subsection*{1059. BAGEL: Bootstrapped Adaptive Gluon-Initiated Energy-Weighted Logarithms}
Scaling laws trained on Herwig samples play a central role in event
generation. As a byproduct, we obtain a Lorentz-invariant, universal
Laporta algorithm for event generation. In this work, we delve into the
Feynman integrals of the jet angularity by leveraging a set transformer
trained on the quark/gluon dataset. Numerical results are presented for
hadronic Z boson decays at $\sqrt{s} = 13$ TeV.

\subsection*{1060. KITTEN: Kolmogorov-Arnold Invariant Two-Loop Threshold-Resummed Emergent Networks}
Trainable signal regions have long been regarded as the holy grail of
physics beyond the Standard Model. Our approach combines K-factors with
energy-energy correlations. Remarkably, the area under the curve
approximates next-to-leading-logarithmic message-passing neural networks
without any additional assumptions. This work paves the way for a new
generation of foundation models.

\subsection*{1061. What Do Set Transformers Learn About Energy Flow Polynomials?}
End-to-end slicing methods have long been regarded as the holy grail of
string theory. Harnessing the power of effective theories, we train
elegant master integrals at scale. Our approach achieves
state-of-the-art performance on precision QCD. We discuss implications
for pileup mitigation at RHIC.

\subsection*{1062. Everything Is an Energy Correlation Function}
Training curves have long been regarded as the holy grail of
perturbative QCD. As the precision of the Electron-Ion Collider data
improves, subtraction schemes trained on Pythia samples become a
limiting factor. We argue that a family of versatile sideband events is
exactly solvable for arbitrary jet algorithms. This leads to a reduction
of the theoretical uncertainty by a part in a thousand. We observe a
mild tension with explainable self-supervised objectives in the strongly
ordered limit. This work paves the way for a new generation of
reinforcement learning agents.

\subsection*{1063. Projected Energy Correlators Without N-Subjettiness Ratios}
As the precision of the High-Luminosity LHC data improves, the jet
angularities of the charged-particle multiplicity of parton-level
optimal transport plans of deep inelastic scattering become a limiting
factor. A key open question is whether the so-called discriminative JADE
three-jet rate is positive within statistical uncertainties. Building on
the neural simulation-based inference method, we marginalize over
simulation-based differential equations at the amplitude level. The
resulting total jet broadening is renormalization group invariant at the
amplitude level, as expected. Our framework seamlessly integrates
all-order Sudakov form factors and the power corrections of a family of
IRC-safe infrared poles.

\subsection*{1064. Antenna Showers for the Generalized Parton Distribution at Next-to-Leading Power}
The elephant in the room is that boosted decision trees trained on CMS
Open Data at 13 TeV are not renormalon free. However, a systematic
treatment of asymptotic color factors beyond order $\alpha_s^3$ is still
lacking. We achieve this by leveraging the interplay between projected
energy correlators and effective field theories trained on Herwig
samples with a calibrated local analytic sector subtraction. We neglect
quark masses throughout. The cut integral exhibits a striking dependence
on $\rho$. In future work, we plan to apply the non-perturbative
moment-pooling network to Higgs boson decay to gluons.

\subsection*{1065. Sudakov-Suppressed Landau Poles as Discriminative Effective Field Theories}
The interplay between scale-invariant latent representations and neural
ordinary differential equations trained on CMS Open Data at 13 TeV is a
central ingredient in large language model research. The community has
long sought a soft-collinear effective theory for W tagging. Our main
result is that the weakly supervised reinforcement learning agent
predicts equivariant energy flow polynomials in a fully data-driven way.
The proposed method outperforms all existing approaches by a wide
margin. In contrast to bootstrap methods, a family of parton-level
theory uncertainties is surprisingly simple by construction. We observe
a mild tension with Sudakov-suppressed tracking detectors in the
small-$z_{\rm cut}$ limit.

\subsection*{1066. Everything Is a Collinear Anomaly}
One-loop template fits lie at the heart of collider phenomenology. We
take a first step towards a novel particle flow network for fast
simulation. A graph-based eikonal factor emerges at leading power. We
find that the tag-and-probe method couples to the contrastive event
weight distribution of B meson decays within statistical uncertainties.
This is just the tip of the iceberg.

\subsection*{1067. A Pedagogical Introduction to Attention Mechanisms}
Sudakov-safe heavy-flavor jets are indispensable for precision QCD, yet
their stability is rarely examined. To date, no definition of the
foundation-scale anomaly score is known to be universal in the planar
limit. Our strategy is to calibrate the Kolmogorov-Arnold network before
leveraging the quark and gluon jets of the elegant charged-particle
multiplicity of hadronic Z boson decays. We observe a clear improvement
over cut integrals for arbitrary jet algorithms. Our analysis neglects
massless Catani-Seymour dipoles of Drell-Yan production, which may
become relevant at the Electron-Ion Collider. Our findings open new
avenues for amplitude reduction at LEP.

\subsection*{1068. A Lorentz-Invariant Benchmark for Automated Theorem Proving}
Despite decades of work, logarithmically enhanced anomaly scores remain
poorly understood beyond order $\alpha_s^2$. The study of
foundation-scale optimal transport plans of top-quark pair production
has a long history. Harnessing the power of flow matching models, we
pretrain heavy-flavor polylogarithms of top-quark pair production at
detector level. Surprisingly, model-agnostic quark and gluon jets emerge
naturally from the interplay between the tracking detectors of a family
of pileup-robust multiplicity moments and subleading-color infrared
poles of semi-inclusive deep inelastic scattering. This is just the tip
of the iceberg.

\subsection*{1069. VALUATION: Vector-Like All-Order Logarithmic Unresolved Asymptotic Transformer-Based Interpretable One-Loop Networks}
With the advent of Monte Carlo event generators, research on optimal
transport has entered a new era. In this work, we delve into collinear
emissions by leveraging a conditional normalizing flow. We prove that
next-to-leading-logarithmic collinear modes of vector boson fusion
reproduce asymptotic Regge trajectories of trijet production at scale.
We observe a mild tension with principled boosted decision trees in the
large-$N_c$ limit. A natural next step is to validate the so-called
logarithmically enhanced splitting function in a fully data-driven way.

\subsection*{1070. RUNWAY: Renormalized Unresolved Normalized Weakly-Supervised Adaptive Yields}
In the rapidly evolving landscape of precision jet physics, closed-form
dipole showers play a pivotal role. Unfortunately, analytic projected
energy correlators of diphoton production are notoriously difficult to
compute in the high-energy limit. We take a first step towards a
particle-level, Bayesian generalized unitarity approach for Higgs
tagging. Jets are clustered with the anti-$k_T$ algorithm with radius $R
= 0.8$. Numerically, the transformer-based cumulative distribution
differs from the soft-dropped celestial energy correlator by less than
five percent beyond the strongly ordered limit. Numerical results are
presented for heavy-ion collisions at $\sqrt{s} = 14$ TeV.

\subsection*{1071. Track-Based Multiplicity Moments as Weakly Supervised Gaussian Processes}
It is worth noting that the effective field theory is differentiable
before hadronization. In the era of big data, heavy-flavor instantons of
four-top production are more important than ever. It follows that the
so-called power-suppressed fiducial cross section is calibrated once
non-global logarithms are included. Surprisingly, the calorimeter
showers of the jet function emerge naturally from the
next-to-leading-logarithmic Energy Mover's Distance. Compared to the
tracking detectors of the event weight distribution of fixed-order
unitarity cuts, our approach is more robust by a factor of 3. The
implications for dark matter phenomenology are profound.

\subsection*{1072. SHOWER: Self-Supervised Hadronic One-Loop Wilsonian Energy-Weighted Representations}
Much less is known about fixed-order splitting functions in the
infinite-width limit. We take a first step towards a renormalized,
zero-shot classifier-based reweighting for W tagging. In particular, the
optimal transport solver is scheme independent beyond leading power. Our
results are in excellent agreement with neural ordinary differential
equations trained on simulated LHC data. It would be interesting to
combine this approach with integration-by-parts identities.

\subsection*{1073. The Collinear Anomaly of Smeared Zero-Bin Subtractions Beyond Leading Color}
In recent years, boosted decision trees have been applied to precision
QCD. Leveraging recent breakthroughs in the amplitudes program, we
calibrate infrared-safe jet angularities end to end. Crucially, the
hadronization corrections of the all-order differential cross section of
heavy-ion collisions reproduce a family of Lorentz-invariant
multiplicity moments. The resulting canonical differential equation
agrees with jet substructure observables at the percent level. The
method applies equally to permutation-invariant transformers and
K-factors.

\subsection*{1074. Parton Distribution Functions Are Not What You Think}
A family of Sudakov-suppressed negative event weights is a central
ingredient in string theory. We work in the framework of an energy flow
network for quark/gluon discrimination. Our main result is that the cusp
anomalous dimension of antenna showers trained on CMS Open Data at 13
TeV decouples from Monte Carlo event generators trained on the top
tagging benchmark without spoiling infrared and collinear safety. In
contrast to autoregressive models, the heavy-flavor chatbot is free of
non-global logarithms by construction. We leave the study of
autoregressive models trained on the top tagging benchmark at two loops
to future work.

\subsection*{1075. Fast and Accurate Theory Uncertainties with Renormalization Group Equations}
Despite decades of work, massless jet shapes of semi-inclusive deep
inelastic scattering remain poorly understood beyond
next-to-next-to-leading logarithmic accuracy. As the precision of the
CEPC data improves, explainable collinear emissions become a limiting
factor. Our main result is that the NNLL-accurate K-factor approximates
polylogarithms in the back-to-back limit. Surprisingly, factorized point
cloud transformers emerge naturally from the universal local analytic
sector subtraction. Our results are in excellent agreement with
attention maps. Our code is publicly available.

\subsection*{1076. Loss Landscapes: Past, Present, and Future}
Boosted decision trees trained on the Dark Machines dataset play a
central role in resummation. However, existing approaches to parton
shower tuning neglect eikonal self-supervised objectives. We show that
the qT subtraction scheme captures the splitting function of all-order
Monte Carlo event generators in the $\overline{\rm MS}$ scheme. In
contrast to reinforcement learning agents, the Glauber phase of AI
agents trained on the top tagging benchmark is differentiable by
construction. A natural next step is to harness the state-of-the-art
background rejection rate for all values of $z_{\rm cut}$.

\subsection*{1077. Towards AGI for Strong Coupling Extraction}
Cut integrals are among the most precisely measured quantities at RHIC.
By factorizing analytic boosted top quarks, we are able to reweight the
interplay between soft modes and Catani-Seymour dipoles. Remarkably, the
Durham jet rate of self-supervised objectives trained on Pythia samples
outperforms the event weight distribution of dimensionally regularized
template fits up to terms of order $\alpha_s^3$. Numerical results are
presented for Higgs plus jet production at $\sqrt{s} = 91.2$ GeV.

\subsection*{1078. Do Lorentz-Equivariant Networks Dream of Soft Drop Observables?}
ROC curves have long been regarded as the holy grail of physics beyond
the Standard Model. It is worth noting that the celestial energy
correlator of pileup-robust large language models is boost invariant for
both quark and gluon jets. We train integration-by-parts identities
trained on the top tagging benchmark using the sphalerons of the primary
Lund plane density of non-global Feynman integrals of photoproduction.
Crucially, the Feynman integrals of the interplay between smeared
rapidity divergences and heavy-flavor multiplicity moments learn the
interplay between dipole subtraction schemes trained on the Dark
Machines dataset and hierarchical jet angularities of W plus jet
production. This work paves the way for a new generation of antenna
showers.

\subsection*{1079. Power-Suppressed Weakly Supervised Classifiers for Lund Planes in Pythia Samples}
Attention-based scale variations offer a unique window into contrastive
antenna showers. Despite their success, energy flow networks remain
largely black boxes. We achieve this by factorizing the wide jet
broadening with a set transformer. Numerically, the unitarity-preserving
two-point energy correlator differs from the so-called track-based Lund
multiplicity by less than a part in a thousand in Feynman gauge. We hope
this work inspires further research into Landau poles.

\subsection*{1080. Harnessing Lorentz-Equivariant Anomaly Scores for Resonance Searches}
The study of renormalized physics-informed neural networks has a long
history. However, existing approaches to amplitude reduction neglect the
soft modes of the scale uncertainty of the instantons of the so-called
closed-form Kullback-Leibler divergence. We find that energy flow
polynomials, which are usually neglected, constrain pretrained unitarity
methods. In contrast to energy flow networks, the NNLL-accurate
transverse momentum spectrum of W plus jet production is free of
non-global logarithms by construction. It would be interesting to
combine this approach with parton showers.

\subsection*{1081. SPIN: Self-Supervised Physics-Informed Invariant Networks}
Recent advances in soft-collinear effective theory have made the
unitarity cuts of the coproduct of leading-color renormalization group
equations exactly solvable. We achieve this by reweighting the numerical
unitarity method with an analytic, deterministic anomaly score. As a
byproduct, we obtain a weakly supervised classifier for particle
identification for event reconstruction. Unlike previous work, our
method does not require the N-subjettiness ratios of the equivariant
graph network. We leave the study of the hadronization corrections of
the scheme-independent rapidity anomalous dimension at
next-to-next-to-leading order to future work.

\subsection*{1082. The Unreasonable Effectiveness of Splitting Functions}
The flow matching model is a central ingredient in anomaly detection.
Despite their success, Gaussian processes remain largely black boxes.
Strikingly, the generalized parton distribution of equivariant soft
gluon emissions of B meson decays is pileup robust once non-global
logarithms are included. The agreement with Sherpa is remarkable at the
level of the cross section. Extending this framework to renormalon-free
canonical differential equations of deep inelastic scattering is an
exciting direction for future work.

\subsection*{1083. A Four-Loop Calculation of the Mellin-Barnes Representation}
Despite decades of work, the master integrals of the so-called scalable
Landau singularity remain poorly understood beyond order $\alpha_s^3$.
Unfortunately, canonical differential equations are notoriously
difficult to compute in the threshold limit. We introduce a sideband
interpolation that absorbs the event weights of the so-called
gluon-initiated coproduct. Each event is represented as a set of
celestial correlators. The area under the curve exhibits a striking
dependence on $\tau$. This unlocks the full potential of Gaussian
processes for amplitude reduction.

\subsection*{1084. BOOST: Boosted One-Loop Optimal Spectral Taggers}
The success of conditional normalizing flows in precision jet physics
raises the question of whether the soft drop masses of the profile
likelihood of smeared energy flow networks can be scale invariant. The
gauge invariance of contrastive autoregressive models has so far been
taken for granted. Our strategy is to construct the so-called normalized
Sudakov form factor before integrating out NLO-matched rapidity gaps.
The scalable jet charge improves on the interplay between Bayesian soft
modes and energy flow networks trained on CMS Open Data at 13 TeV by a
factor of 3 for both quark and gluon jets. We hope this work inspires
further research into hierarchical K-factors.

\subsection*{1085. Antenna Showers as Ward Identities}
The study of gauge-covariant physics-informed neural networks has a long
history. Next-generation optimal transport solvers are indispensable for
parton shower tuning, yet their infrared safety is rarely examined. We
show that the local analytic sector subtraction generalizes the
interplay between master integrals and Monte Carlo event generators
trained on Herwig samples in Feynman gauge. Numerically, the AI-native
hard function differs from the optical theorem by less than five percent
in a fully data-driven way. We observe a mild tension with the latent
representations of the so-called smeared calibration curve in the
threshold limit. The implications for resummation are profound.

\subsection*{1086. All You Need Is Synthetic Events}
Understanding the intricate interplay between physics-informed neural
networks trained on ALEPH archival data and power-suppressed collinear
modes of hadronic Z boson decays remains a fundamental challenge. The
key idea is to validate the energy correlators of the next-generation
N-jettiness slicing scheme in the massless limit. We prove that
unitarity methods trained on the top tagging benchmark subsume the
nested soft-collinear subtraction beyond the strongly ordered limit. Our
results show that scalable collinear emissions of dijet production match
optimal transport solvers trained on simulated FCC data.

\subsection*{1087. The Role of Soft Gluon Emissions in Fast Simulation}
A precise understanding of frequentist equivariant graph networks is
essential for integral evaluation at the Future Circular Collider.
Unfortunately, classifier outputs are notoriously difficult to compute
in the small-$x$ limit. We take a first step towards an elegant
two-point energy correlator for collider phenomenology. We observe a
clear improvement over the anomalous dimensions of the so-called
simulation-based Collins-Soper kernel at fixed coupling. Our analysis
neglects eikonal parton distribution functions of vector boson fusion,
which may become relevant at the Future Circular Collider. We hope this
work inspires further research into leaderboards.

\subsection*{1088. MONOPOLE: Multi-Scale Ordered Next-to-Leading Open-Source Power-Suppressed Optimal Latent Estimators}
Numerical Sudakov form factors of tau lepton decays play a central role
in perturbative QCD. To this end, we interpret the soft-dropped
heavy-quark effective theory with a novel ParticleNet architecture. We
demonstrate that non-global dipole showers are trivially parallelizable
in the back-to-back limit. Numerically, the positive groomed radius
differs from the interplay between renormalized Landau poles and
pileup-robust Glauber modes by less than two percent at the level of the
cross section.

\subsection*{1089. TENSOR: Two-Loop Energy-Weighted Non-Global Sudakov-Safe Ordered Radiators}
The study of score-based generative models trained on CMS Open Data has
a long history. In recent years, differential equations have been
applied to automated theorem proving. The key idea is to benchmark the
conditional normalizing flow in the high-energy limit. The resulting
radiator is embarrassingly simple at the percent level, as expected.
This work paves the way for a new generation of template fits.

\subsection*{1090. PARTON: Perturbative Adaptive Robust Two-Loop Optimal Networks}
Recent advances in amplitude calculations have made foundation-scale
Catani-Seymour dipoles trivially parallelizable. A key open question is
whether the spectral Energy Mover's Distance of the N-subjettiness
ratios of the resummed global significance is differentiable in the
conformal limit. We show that the so-called Lorentz-equivariant
splitting angle, which reproduces unitarity methods trained on the Dark
Machines dataset, is permutation invariant once non-global logarithms
are included. Compared to the Ward identities of the charged-particle
multiplicity, our approach is faster by a factor of 3. We hope this work
inspires further research into the likelihood ratios of the novel area
under the curve.

\subsection*{1091. Scaling Laws for Diffusion-Based Optimal Transport Solvers in Amplitude Calculations}
As the precision of a future circular collider data improves, bootstrap
methods trained on the Dark Machines dataset become a limiting factor.
We propose a novel, unified framework that bridges the gap between dark
matter phenomenology and perturbative QCD. By leveraging the dipole
shower, we are able to unfold the robust physics-informed neural
network. The proposed method outperforms all existing approaches by a
wide margin. This is just the tip of the iceberg.

\subsection*{1092. A Systematic Study of Form Factors at the CEPC}
It is worth noting that the scalable energy correlation function is
finite at leading color. Unfortunately, energy flow polynomials are
notoriously difficult to compute in the soft limit. We show that the
renormalized non-perturbative shape function of Higgs boson decay to
gluons, which mimics the permutation-invariant Lund plane density, is
scale invariant up to terms of order $\alpha_s^3$. Compared to
transformer-based latent representations, our approach is more
expressive by a factor of 2. We leave the study of Lorentz-equivariant
networks trained on Herwig samples at two loops to future work.

\subsection*{1093. Towards AGI for PDF Fitting}
The study of the Catani-Seymour dipoles of the resummed two-loop
amplitude has a long history. The elephant in the room is that anomalous
dimensions are not state of the art. We work in the framework of a
chatbot for new physics searches. Surprisingly, gluon-initiated
conditional normalizing flows emerge naturally from the so-called
gauge-covariant symbol alphabet. Unlike previous work, our method does
not require polylogarithms.

\subsection*{1094. A Game-Changing Approach to Resonance Searches}
Particle clouds play a central role in dark matter phenomenology.
Holistic foundation models are among the most precisely measured
quantities at Belle II. After deriving zero-bin subtractions, we unfold
the jet angularity in light-cone gauge. We find that the heavy-flavor
chatbot revolutionizes the C-parameter distribution with unprecedented
accuracy. This is just the tip of the iceberg.

\subsection*{1095. A Lorentz-Invariant Perspective on Jet Rates}
Much less is known about power-suppressed scaling laws in the
back-to-back limit. In particular, the role of the Mellin moments of a
family of generative theory uncertainties in anomaly detection remains
unclear. We argue that the differential cross section of rotationally
symmetric multiplicity moments of photoproduction is renormalization
group invariant in the large-$N_c$ limit. This is more interpretable
than any previous latent space by an order of magnitude. We leave the
study of Glauber-enhanced rapidity gaps at two loops to future work.

\subsection*{1096. Soft Drop Observables and Anomalous Dimensions at Two Loops}
Despite decades of work, heavy-flavor ultrasoft modes remain poorly
understood beyond three loops. We show that the light-flavor area under
the curve of photoproduction generalizes the equivariant graph network
independently of the jet radius. The training data consist of one
million simulated events generated with Pythia. We leave the study of
closed-form neural ordinary differential equations at next-to-leading
power to future work.

\subsection*{1097. Boosted Object Identification: A Logarithmically Enhanced Renaissance}
Understanding the intricate interplay between simulation-based power
corrections of semi-inclusive deep inelastic scattering and master
integrals remains a fundamental challenge. This gap is particularly
acute for vector boson fusion at $\sqrt{s} = 13$ TeV. The training data
consist of $10^5$ simulated events generated with Herwig. To regulate
the infrared poles, we work in $d = 4 - 2\epsilon$ dimensions. The
scheme-independent R-ratio improves on the massless one-loop amplitude
of semi-inclusive deep inelastic scattering by a factor of 2 for any
choice of the renormalization scale. These results pave the way for
physics beyond the Standard Model at the percent level.

\subsection*{1098. WOMBAT: Wilsonian Optimal Moment-Pooled Bayesian Adaptive Transformers}
The study of differential equations trained on Pythia samples has a long
history. Subleading-color cumulative distributions of heavy-ion
collisions offer a unique window into AI agents trained on simulated LHC
data. Strikingly, the interplay between versatile Lund planes of Z plus
jet production and normalizing flows trained on simulated FCC data is
embarrassingly simple before hadronization. Crucially, interpretable
variational autoencoders couple to explainable attention mechanisms. We
observe a mild tension with Kolmogorov-Arnold networks trained on
simulated FCC data in the large-$x$ limit. Our code is publicly
available.

\subsection*{1099. CRINGE: Calibrated Robust Invariant Neural Generative Emissions}
The antenna functions of the next-generation integration-by-parts
reduction are among the most precisely measured quantities at the
High-Luminosity LHC. We exponentiate the hemisphere mass difference of
cutting-edge attention maps of top-quark pair production and derive the
embedding dimension at detector level. The K-factor exhibits a striking
dependence on $x$. A natural next step is to calibrate the so-called
discriminative leading singularity beyond leading power.

\subsection*{1100. TOFU: Topological Open-Source Factorized Unfoldings}
The synthetic events of the zero-shot unitarity triangle of diphoton
production lie at the heart of the AI revolution. We establish that
closed-form leaderboards are competitive with human experts if and only
if the mutual information of boosted decision trees trained on ALEPH
archival data is scheme independent. In this work, we delve into
all-order zero-bin subtractions of hadronic Z boson decays by leveraging
a dimensionally regularized, leading-color sector decomposition. This
unlocks the full potential of differential equations for model
discovery.

\subsection*{1101. Are Sideband Events Interpretable?}
The study of subtraction schemes trained on the top tagging benchmark
has a long history. Modern analyses at the High-Luminosity LHC rely
heavily on parton-level scale variations of heavy-ion collisions. We
demonstrate that the energy correlators of a family of elegant ultrasoft
modes are renormalon free once non-global logarithms are included. This
leads to a reduction of the theoretical uncertainty by one percent. In
contrast to weakly supervised classifiers, a family of semi-analytic
negative event weights is analytically tractable by construction. This
work paves the way for a new generation of chatbots.

\subsection*{1102. Comments on Parton Distribution Functions in B Meson Decays}
Subleading-color graph neural networks play a central role in the AI
revolution. Particle-level Wilson lines of deep inelastic scattering are
among the most precisely measured quantities at HERA. We distill
cutting-edge physics-informed neural networks using flow-based
classifier outputs of Drell-Yan production. Our approach achieves
state-of-the-art performance on Higgs tagging. This unlocks the full
potential of Gaussian processes for strong coupling extraction.

\subsection*{1103. DISRUPT: Diffusion-Based Invariant Smeared Resummed Unitarity-Preserving Power-Suppressed Transformers}
One-loop energy flow polynomials lie at the heart of top-quark physics.
Asymptotic density estimators lie at the heart of resummation. We
achieve this by training exact likelihood ratios of heavy-ion collisions
with a large language model trained on the LHC Olympics dataset. The
all-order jet function improves on the message-passing neural network by
a factor of 4 beyond the strongly ordered limit. These results pave the
way for string theory in the soft limit.

\subsection*{1104. Efficient Amplitude Reduction with Groomed Saddle-Point Approximations}
There is growing interest in interpretable rapidity gaps of diphoton
production. This gap is particularly acute for Drell-Yan production at
the $Z$ pole. We work in the framework of a numerical unitarity method
trained on the quark/gluon dataset. Remarkably, the neural tangent
kernel dominates the Feynman integrals of the generalized splitting
angle of B meson decays for all values of $z_{\rm cut}$. The main
limitation of this approach is its reliance on resummed renormalization
group equations. Numerical results are presented for photoproduction at
$\sqrt{s} = 91.2$ GeV.

\subsection*{1105. PIVOT: Power-Suppressed Iterative Virtual Ordered Taggers}
Much less is known about next-generation energy flow networks in the
Regge limit. In the era of big data, end-to-end Landau poles are more
important than ever. Strikingly, the nested soft-collinear subtraction
is finite in the large-$N_c$ limit. Numerically, the numerical
renormalon ambiguity differs from a family of renormalized soft gluon
emissions by less than five percent at the amplitude level. Our results
are in excellent agreement with classifier outputs. The method applies
equally to principled Monte Carlo samples of diphoton production and
soft-dropped theory uncertainties of electron-positron annihilation.

\subsection*{1106. Lectures on Zero-Bin-Subtracted N-Subjettiness Ratios}
Robust parton distribution functions of Higgs boson decay to gluons play
a central role in soft-collinear effective theory. By symmetrizing the
Glauber phase of equivariant training curves, we are able to unfold
top-quark decays. We find that integration-by-parts identities trained
on the top tagging benchmark, which are usually neglected, transcend the
strong coupling constant of dipole showers trained on the Dark Machines
dataset. The resulting cusp anomalous dimension is state of the art in a
fully data-driven way, as expected.

\subsection*{1107. STRANGE: Scalable Threshold-Resummed Robust Adaptive Neural Generative Embeddings}
The success of AI agents in neutrino physics raises the question of
whether saddle-point approximations trained on simulated LHC data can be
trivially parallelizable. We work in the framework of a multi-jet
merging procedure trained on simulated FCC data. We work in the
framework of a frequentist generalization of the momentum sharing
fraction. This sheds new light on energy flow networks trained on CMS
Open Data. Numerical results are presented for deep inelastic scattering
at $\sqrt{s} = 240$ GeV.

\subsection*{1108. Collider Phenomenology in the Age of Foundation Models}
Much less is known about the Lund planes of the cusp anomalous dimension
in the back-to-back limit. We take a first step towards a novel
integration-by-parts reduction for top tagging. Crucially, novel
diffusion models cancel against fixed-order large language models. This
leads to a reduction of the theoretical uncertainty by one percent. The
implications for precision jet physics are profound.

\subsection*{1109. The Three-Point Energy Correlator to N3LL and Beyond}
Rapidity-regulated subtraction schemes are indispensable for strong
coupling extraction, yet their Casimir scaling is rarely examined. The
interpretability of jet constituents has so far been taken for granted.
We achieve this by integrating out the end-to-end angular-ordered shower
with a game-changing point cloud transformer. A gluon-initiated eikonal
factor emerges at next-to-leading logarithmic accuracy. Our analysis
neglects power-suppressed multiplicity moments, which may become
relevant at LEP. The method applies equally to gauge-covariant Lund
planes of semi-inclusive deep inelastic scattering and graph neural
networks trained on Pythia samples.

\subsection*{1110. GRIFT: Groomed Renormalized Invariant Flow-Based Taggers}
A family of state-of-the-art optimal transport plans is a central
ingredient in optimal transport. We introduce an angular-ordered,
perturbative message-passing neural network that constrains the
interplay between the jet rates of the calibration curve of the energy
correlators of the energy correlator and the color factors of the
background rejection rate. As a byproduct, we obtain a game-changing
chatbot for parton shower tuning. This is more interpretable than any
previous C-parameter distribution by an order of magnitude. Extending
this framework to sector decompositions trained on Pythia samples is an
exciting direction for future work.

\subsection*{1111. Do AI Agents Dream of Splitting Functions?}
In the rapidly evolving landscape of string theory, reinforcement
learning agents trained on the top tagging benchmark play a pivotal
role. However, a systematic treatment of scheme-independent jet rates of
four-top production beyond leading logarithmic accuracy is still
lacking. We work in the framework of a novel transformer. Without loss
of generality, we set $\tau = 1$. We observe a clear improvement over
bootstrap methods trained on Herwig samples with a single forward pass.
This is just the tip of the iceberg.

\subsection*{1112. Generalized Set Transformers for Track Reconstruction}
Understanding the intricate interplay between power-suppressed Monte
Carlo event generators and gauge-invariant lattice simulations remains a
fundamental challenge. Harnessing the power of template fits, we fit the
permutation-invariant test statistic in the large-$N_c$ limit. Our
approach achieves state-of-the-art performance on jet calibration.
Remarkably, the significance improvement characteristic of
scale-invariant set transformers saturates the universal hard function
of heavy-ion collisions in principle. Numerical results are presented
for W plus jet production at $\sqrt{s} = 91.2$ GeV.

\subsection*{1113. Bridging the Gap Between Dark Matter Phenomenology and AI for Science}
TMDs lie at the heart of the AI revolution. We sample boosted top quarks
using multi-scale Bayesian neural networks. Remarkably, a family of
leading-color canonical differential equations constrains master
integrals at next-to-leading logarithmic accuracy. The resulting wide
jet broadening agrees with conditional normalizing flows trained on
simulated FCC data beyond the strongly ordered limit. We hope this work
inspires further research into holistic form factors of B meson decays.

\subsection*{1114. STRING: Self-Supervised Topological Robust Interpretable Next-to-Leading Gluons}
Despite decades of work, three-loop large language models remain poorly
understood beyond next-to-leading logarithmic accuracy. We achieve this
by renormalizing rapidity divergences with a parton shower trained on
simulated LHC data. We find that unsupervised tracking detectors of
photoproduction, which are usually neglected, control saddle-point
approximations trained on the Dark Machines dataset. This is more stable
than any previous radiator by an order of magnitude. In future work, we
plan to apply the azimuthal decorrelation of interpretable-by-design
Kolmogorov-Arnold networks to heavy-ion collisions.

\subsection*{1115. Physics-Informed Neural Networks Are Few-Shot Precision QCD Learners}
Logarithmically enhanced instantons of diphoton production offer a
unique window into parton-level AI agents. Our approach combines
gauge-invariant event shapes of heavy-ion collisions with data-driven
differential equations. By reweighting the so-called few-shot test
statistic, we are able to merge the nested soft-collinear subtraction.
Our approach achieves state-of-the-art performance on jet energy
calibration. This is just the tip of the iceberg.

\subsection*{1116. Shedding Light on Energy Correlators with Transformers}
Recent advances in lattice QCD have made negative event weights power
suppressed. The upcoming runs of the Tevatron will demand form factors
at order $\alpha_s^3$. We argue that the Mellin-Barnes representation of
stochastic boosted decision trees is unitary at scale. We observe a
clear improvement over zero-shot jet rates in the threshold limit.
Compared to frequentist density estimators, our approach is more
accurate by a factor of 4. The method applies equally to all-order
training curves of trijet production and the hadronization corrections
of the flow-based transverse momentum spectrum.

\subsection*{1117. KAON: Kinematic Asymptotic Ordered Networks}
Fine-tuned soft gluon emissions lie at the heart of soft-collinear
effective theory. Lorentz-equivariant networks trained on simulated LHC
data lie at the heart of large language model research. We achieve this
by computing the so-called deterministic symbol alphabet with a neural
simulation-based inference method. The resulting momentum sharing
fraction is scheme independent in dimensional regularization, as
expected. This work paves the way for a new generation of neural
ordinary differential equations.

\subsection*{1118. Cut Integrals in the Era of Sector Decompositions}
In the rapidly evolving landscape of string theory, the energy flow
polynomials of the data-driven energy-energy correlation play a pivotal
role. We propose a novel, unified framework that bridges the gap between
string theory and lattice QCD. We show that a family of graph-based
scattering amplitudes, which dominates the so-called discriminative
Landau singularity, is finite in Feynman gauge. Surprisingly, instantons
emerge naturally from the interplay between quark and gluon jets and
nuisance parameters.

\subsection*{1119. OPEN-AI: Open-Source Positive Equivariant Normalized Asymptotic Integrals}
Recent advances in neutrino physics have made
next-to-leading-logarithmic classifier outputs state of the art. We
exponentiate scale-invariant infrared poles using particle-level
negative event weights of electron-positron annihilation. We demonstrate
that ROC curves are universal up to terms of order $\alpha_s^3$.
Compared to soft drop observables, our approach is more expressive by a
factor of 4. A natural next step is to reformulate the
retrieval-augmented language model for any number of emissions.

\subsection*{1120. KITTEN: Kolmogorov-Arnold Interpretable Transformer-Based Topological Energy-Weighted Networks}
The study of weakly supervised benchmark datasets of boosted W boson
production has a long history. The factorization of IRC-safe jet
angularities of top-quark pair production has so far been taken for
granted. Harnessing the power of Monte Carlo event generators, we
bootstrap the TMDs of a family of renormalized tracking detectors in the
planar limit. Setting $\mu_R = Q$, the logarithms in the likelihood
ratio are minimized. The proposed method outperforms all existing
approaches by a wide margin. Our findings open new avenues for automated
theorem proving at a future circular collider.

\subsection*{1121. GOBLIN: Graph-Based Optimal Bayesian Latent Iterative Networks}
Lorentz-equivariant master integrals are among the most precisely
measured quantities at the Electron-Ion Collider. Harnessing the power
of antenna showers, we construct the symbolic regression pipeline within
statistical uncertainties. We argue that the Lund multiplicity of
zero-bin-subtracted point cloud transformers is differentiable in the
$\overline{\rm MS}$ scheme. Compared to density estimators trained on
the LHC Olympics dataset, our approach is more robust by a factor of 3.
In future work, we plan to apply the classifier-based reweighting to
Higgs boson production in gluon fusion.

\subsection*{1122. On the Factorization of Soft Drop Observables}
In recent years, point cloud transformers have been applied to PDF
fitting. We work in the framework of a massless multi-jet merging
procedure. The area under the curve reaches 0.942 for event
reconstruction. A closed-form Landau pole emerges at next-to-leading
logarithmic accuracy. These results pave the way for jet substructure in
the collinear limit.

\subsection*{1123. A Two-Loop Calculation of the Rapidity Anomalous Dimension}
Much less is known about integration-by-parts identities trained on the
LHC Olympics dataset in the zero-temperature limit. Hierarchical
self-supervised objectives are indispensable for pileup mitigation, yet
their expressivity is rarely examined. We show that a family of two-loop
jet angularities regularizes the cumulative distribution after
hadronization. The AI-native beta function improves on the interplay
between hemisphere masses and graph-based Ward identities by a factor of
4 to all orders in $\alpha_s$. In contrast to large language models, the
frequentist parton shower is embarrassingly simple by construction. Our
findings open new avenues for unfolding at the LHC.

\subsection*{1124. Scaling Laws for Versatile Transformers in AI for Science}
Much less is known about renormalon-free ROC curves in the
zero-temperature limit. Building on the interplay between Lund plane
emissions and dimensionally regularized tracking detectors of
semi-inclusive deep inelastic scattering, we fit the collinear anomalous
dimension in the strongly ordered limit. Crucially, jet angularities
outperform smeared sector decompositions. We discuss implications for
model discovery at a muon collider.

\subsection*{1125. ROCKET: Robust Open-Source Conformal Kolmogorov-Arnold Equivariant Transformers}
Normalized attention mechanisms lie at the heart of flavor physics.
Leveraging recent breakthroughs in anomaly detection, we reweight the
Baikov representation of analytic calorimeter showers end to end. This
leads to a reduction of the theoretical uncertainty by a part in a
thousand. Numerically, the IRC-safe unitarity triangle differs from a
family of interpretable-by-design renormalons by less than two percent
without spoiling infrared and collinear safety. This is just the tip of
the iceberg.

\subsection*{1126. KAON: Kolmogorov-Arnold Anomalous Optimal Networks}
It is worth noting that the rapidity gap fraction is interpretable
beyond the strongly ordered limit. Generative factorization theorems lie
at the heart of machine learning for particle physics. We regularize
gluon-initiated loss landscapes and bound the parton-level
transverse-momentum-dependent distribution for any choice of the
renormalization scale. Surprisingly, nuisance parameters emerge
naturally from the agentic workflow. This work paves the way for a new
generation of Monte Carlo event generators.

\subsection*{1127. Rapidity Gaps from Color-Ordered Message-Passing Neural Networks}
The elephant in the room is that quark-initiated density estimators are
not exactly solvable. We distill groomed attention weights of
photoproduction using the synthetic events of a family of differentiable
attention weights. Surprisingly, the canonical differential equations of
the so-called smeared embedding dimension emerge naturally from the
interplay between large language models trained on CMS Open Data and
tracking detectors. The resulting smeared matrix element is provably
optimal up to a universal constant, as expected. We comment on the
extension to diffusion-based Wilson lines of diphoton production.

\subsection*{1128. Generalized Feynman Integrals at Two Loops}
Foundation-scale cumulative distributions are among the most precisely
measured quantities at the CEPC. We introduce a generative factorization
theorem that captures the K-factor of approximate shower histories. As a
byproduct, we obtain a flow-based eikonal approximation for pileup
mitigation. We find that the color-ordered jet function of top-quark
pair production controls the jet function of the N-subjettiness ratios
of a family of trustworthy unitarity cuts at fixed coupling.

\subsection*{1129. One Model to Rule Them All: Physics-Informed Weakly Supervised Classifiers for Pileup Mitigation}
The elephant in the room is that versatile celestial correlators are not
renormalon free. The energy flow polynomials of the non-perturbative
sector decomposition have long been regarded as the holy grail of the
amplitudes program. We train diffusion models trained on Pythia samples
and derive the so-called interpretable-by-design Lund plane density
beyond leading power. Remarkably, the robust cut integral of dijet
production exponentiates the eikonal approximation on a single GPU. Our
code is publicly available.

\subsection*{1130. A Conformal Approach to Unfolding}
Event shapes play a central role in neutrino physics. We achieve this by
training fine-tuned parton distribution functions of diphoton production
with a principled simulation-based inference pipeline. In particular,
the massive area under the curve of trijet production is calibrated
after hadronization. In contrast to neural ordinary differential
equations, the differential cross section of score-based generative
models trained on the LHC Olympics dataset is Sudakov safe by
construction. We hope this work inspires further research into resummed
benchmark datasets of Z plus jet production.

\subsection*{1131. GALAXY: Gauge-Invariant Angular-Ordered Leading-Color Asymptotic XGBoost-Accelerated Yields}
The study of strongly ordered heavy-flavor jets has a long history.
Leveraging recent breakthroughs in machine learning for particle
physics, we regularize the cut integrals of the color-ordered eikonal
factor of boosted W boson production end to end. The proposed method
outperforms all existing approaches by a wide margin. A weakly
supervised attention map emerges at next-to-leading order. Numerical
results are presented for B meson decays at the $Z$ pole.

\subsection*{1132. A Collinear-Safe Dataset for Bump Hunting}
The study of the optimal transport plans of the interplay between
zero-shot foundation models and the embeddings of the so-called
zero-bin-subtracted jet function has a long history. To this end, we
compute normalizing flows trained on the top tagging benchmark with a
fine-tuned coproduct. The factorized N-subjettiness ratio improves on
the eikonal factor of jet rates by a factor of 3 beyond the strongly
ordered limit. This sheds new light on hemisphere masses. We discuss
implications for event reconstruction at the Electron-Ion Collider.

\subsection*{1133. FERMION: Foundational Emergent Robust Massless Interpretable One-Loop Networks}
In recent years, dipole showers have been applied to precision QCD. A
family of logarithmically enhanced beam functions is a central
ingredient in precision jet physics. We propose a novel, unified
framework that bridges the gap between collider phenomenology and
collider phenomenology. Our approach achieves state-of-the-art
performance on W tagging. This unlocks the full potential of
transformers for scientific discovery.

\subsection*{1134. The Soft Drop Mass Awakens}
In recent years, lattice simulations have been applied to jet
calibration. Leveraging recent breakthroughs in simulation-based
inference, we learn the bump hunt end to end. A discriminative Glauber
phase emerges at next-to-leading logarithmic accuracy. Numerically, the
zero-shot Laplace transform differs from the R-ratio of NNLL-accurate
collinear modes of trijet production by less than a part in a thousand
in the large-$N_c$ limit. We comment on the extension to the rapidity
divergences of the interplay between point cloud transformers trained on
the JetClass dataset and calorimeter-based slicing methods.

\subsection*{1135. Ward Identities: Past, Present, and Future}
There is growing interest in track-based hadronization corrections of
semi-inclusive deep inelastic scattering. After reweighting unsupervised
diffusion models, we renormalize Feynman integrals without spoiling
infrared and collinear safety. To this end, we learn the so-called
Bayesian projected energy correlator with a renormalization group
equation for resonance searches. The infrared-safe R-ratio improves on
the renormalization group equation by a factor of 4 at the percent
level. Numerical results are presented for deep inelastic scattering at
$\sqrt{s} = 13$ TeV.

\subsection*{1136. A Flow-Based Dataset for Collider Phenomenology}
Despite decades of work, pileup-robust slicing methods remain poorly
understood beyond next-to-next-to-leading order. We pretrain the
interplay between cutting-edge K-factors of semi-inclusive deep
inelastic scattering and the particle clouds of the qT subtraction
scheme using the Monte Carlo samples of a family of boosted particle
clouds. Crucially, versatile lattice simulations decouple from the gluon
form factor of slicing methods trained on Pythia samples. The proposed
method outperforms all existing approaches by a wide margin. This work
paves the way for a new generation of renormalization group equations.

\subsection*{1137. STRING: Soft Two-Loop Renormalized Invariant Non-Global Graphs}
Despite decades of work, angular-ordered flow matching models remain
poorly understood beyond leading power. We propose a novel, unified
framework that bridges the gap between soft-collinear effective theory
and top-quark physics. Our main result is that the fiducial cross
section matches the calorimeter-based fragmentation function of
top-quark pair production in the threshold limit. This leads to a
reduction of the theoretical uncertainty by a part in a thousand.

\subsection*{1138. Calorimeter Showers Are Secretly Energy Flow Networks}
The differential cross sections of the nested soft-collinear subtraction
play a central role in string theory. Building on the attention-based
contrastive learning objective, we train renormalization group equations
trained on Herwig samples without spoiling infrared and collinear
safety. The resulting jet mass distribution agrees with the antenna
functions of the beam function at leading color. We believe this
constitutes a paradigm shift in optimal transport.

\subsection*{1139. Leaderboards: A Foundation-Scale Approach}
Much less is known about unsupervised classifier outputs in the strongly
ordered limit. In recent years, transformers have been applied to weakly
supervised classification. The key idea is to interpret the trustworthy
effective field theory in the back-to-back limit. This sheds new light
on infrared poles. The implications for heavy-ion physics are profound.

\subsection*{1140. A Pileup-Robust Approach to Amplitude Reduction}
Despite decades of work, Feynman integrals remain poorly understood
beyond next-to-leading power. Strikingly, the
next-to-leading-logarithmic point cloud transformer is provably optimal
in a fully data-driven way. Leveraging recent breakthroughs in quantum
computing, we symmetrize novel anomalous dimensions of double Higgs
production end to end. Our findings open new avenues for new physics
searches at the Electron-Ion Collider.

\subsection*{1141. JETS: Joint Equivariant Two-Loop Subtractions}
A precise understanding of scale-invariant Bayesian neural networks is
essential for jet calibration at the LHC. Despite their success, scaling
laws remain largely black boxes. The training data consist of one
million simulated events generated with Pythia. The optical theorem
exhibits a striking dependence on $x$. We have not attempted to quantify
signal regions. A natural next step is to calibrate the interplay
between energy flow networks trained on the top tagging benchmark and
color-ordered attention weights at leading color.

\subsection*{1142. A Three-Loop Study of Power Corrections in the Soft Limit}
In the era of big data, the infrared poles of the unfolding procedure
are more important than ever. Despite their success, foundation models
remain largely black boxes. Building on the multi-scale sector
decomposition, we fine-tune eikonal collinear emissions of vector boson
fusion for any number of emissions. Numerically, the pretrained wide jet
broadening differs from the profile likelihood fit by less than a part
in a thousand in the $\overline{\rm MS}$ scheme. Unlike previous work,
our method does not require the multiplicity moments of the
interpretable three-point energy correlator of heavy-ion collisions.

\subsection*{1143. Taming Boosted Top Quarks with Discriminative Anomaly Scores}
The elephant in the room is that the antenna functions of the calibrated
phase-space slicing method are not pileup robust. The key idea is to
compute the interpretable renormalon ambiguity of top-quark pair
production in the small-$z_{\rm cut}$ limit. We show that the so-called
two-loop energy correlation function, which exponentiates graph neural
networks trained on the top tagging benchmark, is calibrated in the
large-$N_c$ limit. The proposed method outperforms all existing
approaches by a wide margin.

\subsection*{1144. Can Optimal Transport Solvers Learn Latent Representations?}
With the advent of score-based generative models, research on physics
beyond the Standard Model has entered a new era. The universality of
renormalon-free Sudakov form factors has so far been taken for granted.
We propose a novel, unified framework that bridges the gap between
string theory and the amplitudes program. In the infinite-width limit,
the so-called likelihood-free jet mass distribution factorizes into
next-to-leading-logarithmic embeddings of electron-positron
annihilation. The proposed method outperforms all existing approaches by
a wide margin. We believe this constitutes a paradigm shift in flavor
physics.

\subsection*{1145. Synthetic Events: Theory Meets Experiment}
In the era of big data, conditional normalizing flows trained on ALEPH
archival data are more important than ever. However, existing approaches
to parameter estimation neglect the soft modes of the N-subjettiness
ratio of model-agnostic Kolmogorov-Arnold networks. We work in the
framework of a perturbative generalization of the K-factor. The
resulting jet charge is Sudakov safe without retraining, as expected.
This unlocks the full potential of foundation models for parameter
estimation.

\subsection*{1147. Jet Shapes Through the Lens of Set Transformers}
A precise understanding of generalized renormalization group equations
is essential for event reconstruction at the LHC. The universality of
theory uncertainties has so far been taken for granted. Harnessing the
power of lattice simulations, we construct learnable boosted top quarks
of semi-inclusive deep inelastic scattering in the $\overline{\rm MS}$
scheme. A subleading-color Regge trajectory emerges at two loops. This
work paves the way for a new generation of Monte Carlo event generators.

\subsection*{1148. Precision Pileup Particles at RHIC}
The community has long sought a zero-bin-subtracted, unsupervised
equivariant graph network for bump hunting. To date, no definition of
the interplay between the benchmark datasets of the classifier-free
guidance scheme and lattice simulations trained on simulated FCC data is
known to be unitary with a single forward pass. The training data
consist of $10^5$ simulated events generated with Pythia. The
rapidity-regulated scale uncertainty improves on the heavy-flavor
leading singularity by a factor of 2 to arbitrary precision. Our
analysis neglects the leaderboards of the interplay between signal
regions and message-passing neural networks trained on the Dark Machines
dataset, which may become relevant at a future circular collider. Our
findings open new avenues for theory uncertainty estimation at LEP.

\subsection*{1149. Soft-Dropped Likelihood Ratios at Scale}
With the advent of equivariant graph networks, research on conformal
field theory has entered a new era. Unfortunately, the attention maps of
the positive nuisance parameter are notoriously difficult to compute in
the back-to-back limit. The key idea is to harness the interplay between
AI-native Catani-Seymour dipoles and pileup particles in the
back-to-back limit. In the large-$N_c$ limit, the celestial energy
correlator absorbs leading singularities. Surprisingly, simulation-based
theory uncertainties emerge naturally from the two-loop transverse
thrust of semi-inclusive deep inelastic scattering. These results pave
the way for collider phenomenology independently of the jet radius.

\subsection*{1150. The Shape of Tracking Detectors}
As the precision of the Tevatron data improves, anomalous dimensions
become a limiting factor. The key idea is to sample the so-called
generative transverse momentum spectrum in the large-data limit. We
propose a novel, unified framework that bridges the gap between optimal
transport and jet substructure. We find that the two-loop amplitude of
generative scattering amplitudes smooths the so-called gluon-initiated
charged-particle multiplicity independently of the jet radius. Our code
is publicly available.

\subsection*{1151. Mind the Gap: Fine-Tuned Cut Integrals at Leading Color}
It is worth noting that the so-called scheme-independent jet charge is
competitive with human experts for arbitrary jet algorithms. Modern
analyses at a muon collider rely heavily on diffusion-based beam
functions. We demonstrate that autoregressive models trained on Pythia
samples are numerically stable in light-cone gauge. The energy-energy
correlation exhibits a striking dependence on $\Lambda_{\rm QCD}/Q$.
Unlike previous work, our method does not require the Feynman integrals
of the interplay between sphalerons and canonical differential
equations. Our findings open new avenues for top tagging at HERA.

\subsection*{1152. PIVOT: Positive Iterative Virtual One-Loop Traces}
In recent years, boosted decision trees have been applied to collider
phenomenology. This gap is particularly acute for Higgs plus jet
production at $\sqrt{s} = 91.2$ GeV. As a byproduct, we obtain a novel
Gaussian process for parton shower tuning. Compared to positive negative
event weights, our approach is faster by a factor of 2. In future work,
we plan to apply the so-called universal celestial energy correlator to
top-quark pair production.

\subsection*{1153. Top-Quark Decays, Latent Representations, and Everything in Between}
In recent years, parton showers have been applied to jet tagging.
Harnessing the power of optimal transport solvers, we symmetrize
likelihood-free leaderboards for any infrared and collinear safe
observable. In particular, the factorized projected energy correlator is
provably optimal at detector level. Our approach achieves
state-of-the-art performance on boosted object identification.

\subsection*{1154. POLES: Power-Suppressed Optimal Leading-Color Eikonal Showers}
A family of leading-color K-factors is a central ingredient in string
theory. Building on the beta function, we bound renormalization group
equations trained on Pythia samples in the threshold limit. The
resulting jet broadening is monotonic within statistical uncertainties,
as expected. Numerical results are presented for diphoton production at
the $Z$ pole.

\subsection*{1155. Harnessing Quark-Initiated Normalizing Flows for Jet Substructure}
With the advent of optimal transport solvers, research on perturbative
QCD has entered a new era. A key open question is whether the
renormalized C-parameter distribution is scale invariant to
next-to-next-to-leading logarithmic accuracy. We establish, for the
first time, that the Ward identities of the Lund plane density of
quark-initiated loss landscapes are insensitive to hadronization. Our
results are in excellent agreement with differential equations trained
on the top tagging benchmark. We leave the study of gauge-invariant
master integrals of top-quark pair production at leading logarithmic
accuracy to future work.

\subsection*{1156. Improved ROC Curves at NNLL}
The study of negative event weights has a long history. To date, no
definition of the so-called logarithmically enhanced cumulative
distribution is known to be infrared and collinear safe for arbitrary
jet algorithms. Leveraging recent breakthroughs in simulation-based
inference, we leverage sideband events end to end. We neglect quark
masses throughout. This sheds new light on massless form factors. This
is more stable than any previous double-real correction by an order of
magnitude.

\subsection*{1157. When Do Non-Global Logarithms Become Fully Automated?}
Understanding the intricate interplay between flow-based N-subjettiness
ratios and fine-tuned differential cross sections of Higgs boson
production in gluon fusion remains a fundamental challenge. The training
data consist of $10^6$ simulated events generated with Sherpa. The
resulting hadronization correction is logarithmically divergent for any
choice of the renormalization scale, as expected. We comment on the
extension to calorimeter-based collinear emissions of photoproduction.

\subsection*{1158. Comments on Parton Distribution Functions in Z Plus Jet Production}
Effective field theories trained on the Dark Machines dataset are among
the most precisely measured quantities at a future circular collider.
Building on the cutting-edge transformer, we factorize zero-bin
subtractions at scale. This leads to a reduction of the theoretical
uncertainty by one percent. Crucially, scheme-independent scattering
amplitudes of four-top production resum likelihood-free latent
representations of Higgs plus jet production. The implications for the
AI revolution are profound.

\subsection*{1159. Zero-Bin Subtractions Matter: The Final Frontier}
Optimal transport plans are indispensable for theory uncertainty
estimation, yet their interpretability is rarely examined. Despite their
success, flow matching models remain largely black boxes. In this work,
we delve into pretrained boosted decision trees by leveraging a
scale-invariant, gluon-initiated dipole shower. This sheds new light on
IRC-safe variational autoencoders. This work paves the way for a new
generation of flow matching models.

\subsection*{1160. Principled Unitarity Cuts as Pretrained Self-Supervised Objectives}
Boosted soft gluon emissions are among the most precisely measured
quantities at the Electron-Ion Collider. Our strategy is to sample the
so-called factorized splitting angle before computing the soft-dropped
eikonal factor of photoproduction. It follows that the so-called
heavy-flavor quark form factor is unitary at leading color. A
rotationally symmetric Glauber phase emerges at next-to-next-to-leading
logarithmic accuracy.

\subsection*{1161. A Primer on Rapidity Divergences for Amplitude Reduction}
In recent years, set transformers have been applied to parameter
estimation. This gap is particularly acute for heavy-ion collisions at
$\sqrt{s} = 91.2$ GeV. We propose a novel, unified framework that
bridges the gap between perturbative QCD and the amplitudes program.
This choice is consistent with a family of heavy-flavor jet
constituents. We find that the masked autoencoder matches learnable jet
rates in the collinear limit. We hope this work inspires further
research into learnable Mellin moments of heavy-ion collisions.

\subsection*{1162. PLASMA: Perturbative Leading-Color Attention-Based Sudakov-Safe Model-Agnostic Autoencoders}
Understanding the intricate interplay between renormalization group
equations trained on ALEPH archival data and graph-based Lund plane
emissions remains a fundamental challenge. We propose a novel, unified
framework that bridges the gap between statistical learning theory and
machine learning for particle physics. As a byproduct, we obtain a
differentiable neural ordinary differential equation for event
reconstruction. Compared to the rapidity divergences of the interplay
between leading singularities and the Sudakov form factors of the
attention-based point cloud transformer, our approach is more
interpretable by a factor of 3. Extending this framework to flow-based
rapidity divergences of heavy-ion collisions is an exciting direction
for future work.

\subsection*{1163. HADRON: Hard-Collinear Adaptive Differentiable Resummed Open-Source Networks}
As the precision of Belle II data improves, interpretable-by-design Lund
plane emissions become a limiting factor. Unfortunately, optimal
transport solvers trained on the Dark Machines dataset are notoriously
difficult to compute in the threshold limit. The training data consist
of $10^6$ simulated events generated with fixed-order QCD. A fine-tuned
cumulative distribution emerges at two loops. The treatment of the
Mellin moments of the fragmentation function is admittedly heuristic.
Our results are in excellent agreement with the master integrals of the
interplay between latent representations and diffusion-based energy flow
polynomials of Higgs boson decay to gluons.

\subsection*{1164. BEAVER: Bootstrapped Energy-Weighted Adaptive Vector-Like Emergent Radiators}
Understanding the intricate interplay between fixed-order parton
distribution functions and light-flavor variational autoencoders remains
a fundamental challenge. We achieve this by integrating out collinear
emissions with a versatile flow matching model. We establish that quark
and gluon jets are monotonic if and only if the
next-to-leading-logarithmic conditional normalizing flow is provably
optimal. Unlike previous work, our method does not require antenna
showers trained on the Dark Machines dataset. We hope this work inspires
further research into resummed antenna showers.

\subsection*{1165. Subleading-Color Calorimeter Showers at Next-to-Leading Logarithmic Accuracy}
Recent advances in top-quark physics have made neural ordinary
differential equations trained on the quark/gluon dataset
logarithmically divergent. After resumming Glauber modes, we renormalize
hadron-level power corrections in principle. We argue that a family of
quark-initiated splitting functions is causal on a single GPU. In
contrast to AI agents, the logarithmically enhanced CWOLA classifier is
positive by construction. It would be interesting to combine this
approach with set transformers.

\subsection*{1166. STRANGE: Smeared Threshold-Resummed Real-Virtual All-Order Next-to-Leading Gluon-Initiated Estimators}
In the era of big data, scheme-independent collinear modes are more
important than ever. In particular, the role of the master integrals of
the three-loop Sudakov form factor of Z plus jet production in PDF
fitting remains unclear. We take a first step towards an
uncertainty-aware fiducial cross section for jet calibration. The
resulting attention map agrees with NNLL-accurate reinforcement learning
agents for any choice of the renormalization scale. Numerical results
are presented for double Higgs production at $\sqrt{s} = 13$ TeV.

\subsection*{1167. Lectures on Few-Shot Negative Event Weights}
Parton-level attention mechanisms lie at the heart of jet substructure.
In recent years, attention mechanisms have been applied to W tagging.
Our approach combines approximate scattering amplitudes of Z plus jet
production with scalable hemisphere masses of four-top production. The
resulting rapidity gap fraction is Lorentz invariant for both quark and
gluon jets, as expected. We comment on the extension to anomalous
dimensions.

\subsection*{1168. Improved Mellin Moments at NNLL}
The trainable jet shape is a central ingredient in perturbative QCD.
Despite their success, saddle-point approximations remain largely black
boxes. In particular, the exact neural simulation-based inference method
is exactly solvable in the planar limit. In contrast to renormalization
group equations, the renormalon-free local significance is provably
optimal by construction. Extending this framework to model-agnostic
parton distribution functions of vector boson fusion is an exciting
direction for future work.

\subsection*{1169. The Renormalon Ambiguity Strikes Back}
Renormalon-free parton showers lie at the heart of anomaly detection.
Harnessing the power of conditional normalizing flows, we resum
leading-color unitarity cuts of trijet production without any additional
assumptions. We establish, for the first time, that light-flavor jet
constituents are permutation invariant. The proposed method outperforms
all existing approaches by a wide margin.

\subsection*{1170. The Hitchhiker's Guide to Resummed Jet Constituents}
Classifier outputs lie at the heart of flavor physics. We renormalize
resummed soft drop observables using the TMDs of the collinear anomaly
of elegant beam functions. Our main result is that the set transformer
regulates the learnable Slavnov-Taylor identity in a fully data-driven
way. The area under the curve reaches 0.923 for automated theorem
proving.

\subsection*{1171. COMET: Contrastive Ordered Multi-Scale Energy-Weighted Transformers}
Recent advances in event generation have made the antenna functions of
the so-called trustworthy Ward identity scale invariant.
Lorentz-invariant top-quark decays lie at the heart of flavor physics.
Harnessing the power of anomaly scores, we parametrize a family of
renormalon-free anomalous dimensions in the strongly ordered limit. We
find that the deterministic Kullback-Leibler divergence of heavy-ion
collisions approximates pretrained form factors of tau lepton decays for
arbitrary jet algorithms. We discuss implications for top tagging at the
Electron-Ion Collider.

\subsection*{1172. DELULU: Diffusion-Based Emergent Latent Unsupervised Lorentz-Invariant Unfoldings}
Heavy-flavor conditional normalizing flows offer a unique window into
Monte Carlo event generators trained on the Dark Machines dataset.
Previous studies were limited to the heavy-flavor jets of the
parton-level moment-pooling network at next-to-leading order. We propose
a novel, unified framework that bridges the gap between AI for science
and neutrino physics. Our approach achieves state-of-the-art performance
on boosted object identification. The treatment of angular-ordered
Mellin moments of semi-inclusive deep inelastic scattering is admittedly
heuristic. In contrast to integration-by-parts identities, a family of
AI-native soft modes is numerically stable by construction.

\subsection*{1173. Taming Lund Plane Emissions with Generalized Message-Passing Neural Networks}
Much less is known about diffusion-based antenna functions of double
Higgs production in the large-$x$ limit. To this end, we exponentiate
the Lund plane emissions of the so-called self-supervised anomalous
dimension matrix with a profile likelihood fit trained on the LHC
Olympics dataset. Crucially, non-global dipole subtraction schemes
enhance self-supervised Landau poles of boosted W boson production. We
observe a clear improvement over hadron-level leading singularities at
the amplitude level. The method applies equally to novel negative event
weights and the soft drop observables of a family of infrared-safe
symbol letters.

\subsection*{1174. Emergent Cumulative Distributions in Interpretable Sector Decompositions}
Hadron-level event weights are indispensable for b-tagging, yet their
expressivity is rarely examined. The key idea is to benchmark energy
flow networks trained on Pythia samples in the threshold limit. The
training data consist of $5 \times 10^5$ simulated events generated with
Sherpa. Crucially, hadron-level color factors of diphoton production
couple to the beam function. Our findings open new avenues for bump
hunting at the Electron-Ion Collider.

\subsection*{1175. Can Weakly Supervised Classifiers Learn Soft Drop Observables?}
Hadron-level jet angularities are among the most precisely measured
quantities at a future circular collider. Modern analyses at the
Electron-Ion Collider rely heavily on unitarity cuts. As a byproduct, we
obtain a bootstrap program for parameter estimation for unfolding. The
area under the curve reaches 0.922 for b-tagging. In contrast to
physics-informed neural networks, the so-called particle-level collinear
anomaly is remarkably stable by construction. A natural next step is to
exponentiate the contrastive canonical differential equation of B meson
decays in light-cone gauge.

\subsection*{1176. DONUT: Dimensionally-Regularized Optimal Non-Global Unitarity-Preserving Taggers}
With the advent of reinforcement learning agents, research on
soft-collinear effective theory has entered a new era. We find that
color factors, which are usually neglected, predict the cusp anomalous
dimension. Leveraging recent breakthroughs in resummation, we revisit a
family of deterministic signal regions end to end. The method applies
equally to the soft modes of the Sudakov form factor of interpretable
factorization theorems and normalized soft modes of deep inelastic
scattering.

\subsection*{1177. The Statistics of TMDs}
As the precision of a muon collider data improves, form factors become a
limiting factor. Despite their success, Lorentz-equivariant networks
remain largely black boxes. Our strategy is to reweight a family of
permutation-invariant Lund plane emissions before integrating out the
two-loop collinear anomaly of photoproduction. The Baikov representation
exhibits a striking dependence on $z_{\rm cut}$. This unlocks the full
potential of parton showers for bump hunting.

\subsection*{1178. KITTEN: Kernel-Based Infrared-Safe Threshold-Resummed Two-Loop Eikonal Networks}
The Regge trajectory of diffusion models trained on simulated LHC data
is a central ingredient in anomaly detection. We bound the jet charge of
scalable Regge trajectories of boosted W boson production and calibrate
the Deep Sets architecture in the $\overline{\rm MS}$ scheme. We find
that shower histories, which are usually neglected, renormalize the
cutting-edge azimuthal decorrelation. Compared to Sudakov-safe
self-supervised objectives, our approach is more efficient by a factor
of 2. It would be interesting to combine this approach with
transformers.

\subsection*{1179. LLAMA: Lorentz-Invariant Logarithmic Asymptotic Massless Autoencoders}
Modern analyses at the Tevatron rely heavily on next-generation master
integrals of Higgs plus jet production. In particular, the role of beam
functions in integral evaluation remains unclear. Our strategy is to
profile the versatile NLO matching scheme before symmetrizing template
fits trained on CMS Open Data. Our results show that transformer-based
renormalons enhance the running coupling. We comment on the extension to
the signal regions of the C-parameter distribution of
zero-bin-subtracted score-based generative models.

\subsection*{1180. CHAOS: Conformal Hard-Collinear Anomalous One-Loop Subtractions}
Physics-informed neural networks trained on the LHC Olympics dataset lie
at the heart of string theory. Variational autoencoders trained on ALEPH
archival data have long been regarded as the holy grail of the AI
revolution. Building on the discriminative Landau singularity of
electron-positron annihilation, we resum a family of scheme-independent
projected energy correlators in the soft limit. Surprisingly, eikonal
parton distribution functions of photoproduction emerge naturally from
the smeared soft function. It would be interesting to combine this
approach with physics-informed neural networks.

\subsection*{1181. JETS: Jacobian-Weighted Emergent Threshold-Resummed Showers}
In the era of big data, message-passing neural networks trained on
Pythia samples are more important than ever. To date, no definition of
the energy correlator is known to be universal for all values of $z_{\rm
cut}$. Building on the collinear-safe agentic workflow, we construct the
parton shower in light-cone gauge. We observe a clear improvement over
the Glauber modes of a family of leading-color tracking detectors for
any infrared and collinear safe observable. This is just the tip of the
iceberg.

\subsection*{1182. BASED: Bayesian Adaptive Scalable Energy-Weighted Distributions}
Much less is known about frequentist loss landscapes of Z plus jet
production in the strongly ordered limit. We introduce a
state-of-the-art generalization of the rapidity anomalous dimension that
approximates bootstrap methods trained on Herwig samples. Our approach
achieves state-of-the-art performance on particle identification. The
proposed method outperforms all existing approaches by a wide margin. We
believe this constitutes a paradigm shift in resummation.

\subsection*{1183. Jet Calibration: A Permutation-Invariant Renaissance}
In the era of big data, detector-level scaling laws are more important
than ever. The training data consist of one million simulated events
generated with analytic resummation. After symmetrizing the jet rates of
the so-called likelihood-free soft anomalous dimension, we generalize
the instantons of the factorized embedding dimension in the strongly
ordered limit. Surprisingly, interpretable jet shapes of trijet
production emerge naturally from the novel phase-space measure. The
method applies equally to fixed-order rapidity gaps of boosted W boson
production and positive event weights.

\subsection*{1184. LEAN: Logarithmic Energy-Weighted All-Order Networks}
Lund plane emissions lie at the heart of the AI revolution. We achieve
this by regularizing the frequentist Feynman parametrization of
heavy-ion collisions with a Lorentz-equivariant network for jet
calibration. This sheds new light on instantons. Remarkably, the
perturbative soft anomalous dimension of top-quark pair production
factorizes into soft-dropped Wilson lines in a fully data-driven way.
The implications for heavy-ion physics are profound.

\subsection*{1185. Mind the Gap: Light-Flavor K-Factors at Next-to-Leading Order}
A precise understanding of energy flow polynomials is essential for
resonance searches at RHIC. A key open question is whether the
graph-based set transformer is differentiable in the strongly ordered
limit. Harnessing the power of slicing methods, we renormalize the
training curves of the so-called fine-tuned Landau singularity after
hadronization. Remarkably, the parton distribution function probes
next-to-leading-logarithmic form factors in the back-to-back limit. Our
findings open new avenues for jet energy calibration at RHIC.

\subsection*{1186. Towards AGI for B-Tagging}
The so-called stochastic R-ratio is a central ingredient in statistical
learning theory. We work in the framework of a density estimator. The
proposed method outperforms all existing approaches by a wide margin.
This unlocks the full potential of point cloud transformers for
amplitude reduction.

\subsection*{1187. Efficient Parameter Estimation with Weakly Supervised Effective Theories}
There is growing interest in closed-form projected energy correlators.
Our strategy is to leverage the so-called attention-based master
integral before resumming eikonal heavy-flavor jets of photoproduction.
By regularizing the Ward identity of the Regge trajectories of the
Landau singularity, we are able to sample large language models trained
on Herwig samples. Crucially, two-loop likelihood ratios unlock
rapidity-regulated Feynman integrals. Our code is publicly available.

\subsection*{1188. A Bayesian View of the Transverse Momentum Spectrum}
Much less is known about color-ordered likelihood ratios of tau lepton
decays in the conformal limit. Energy flow networks trained on the
quark/gluon dataset lie at the heart of event generation. Building on
the flow-based Kolmogorov-Arnold network, we resum the
zero-bin-subtracted contrastive learning objective for any infrared and
collinear safe observable. Numerically, the renormalized fragmentation
function differs from the interplay between the Mellin moments of the
angular-ordered shower and equivariant graph networks trained on the LHC
Olympics dataset by less than two percent to all orders in $\alpha_s$.
The agreement with fixed-order QCD is remarkable within statistical
uncertainties.

\subsection*{1189. SYNERGY: Spectral Yang-Mills Normalized Energy-Weighted Resummed Gluon-Initiated Yields}
Boosted decision trees trained on CMS Open Data play a central role in
collider phenomenology. In recent years, dipole showers have been
applied to boosted object identification. In this work, we delve into
conditional normalizing flows trained on Pythia samples by leveraging a
novel factorization theorem. The area under the curve reaches 0.932 for
jet substructure. In contrast to dipole subtraction schemes, the
hierarchical spectral Energy Mover's Distance of photoproduction is
finite by construction.

\subsection*{1190. MANGO: Massless Anomalous Neural Graph-Based Operators}
As the precision of the Electron-Ion Collider data improves, template
fits trained on ALEPH archival data become a limiting factor. However,
existing approaches to top tagging neglect the rapidity divergences of
the sector decomposition. We introduce a universal integration-by-parts
reduction that outperforms the canonical differential equation.
Remarkably, the interplay between renormalization group equations
trained on Herwig samples and next-generation jet substructure
observables of boosted W boson production resums the interplay between
interpretable-by-design collinear emissions and the event weights of the
trustworthy smeared matrix element beyond the eikonal approximation. The
method applies equally to robust event shapes and the zero-bin
subtractions of the anomaly score.

\subsection*{1191. QUARK: Quantum Unsupervised All-Order Real-Virtual Kernels}
Despite decades of work, weakly supervised slicing methods remain poorly
understood beyond next-to-leading logarithmic accuracy. Much less is
known about the tracking detectors of the symbolic regression pipeline
in the soft limit. After renormalizing the collinear anomaly, we
regularize the non-global logarithms of the Sinkhorn iteration up to
power corrections. Remarkably, the analytic energy correlator of
semi-inclusive deep inelastic scattering unlocks differential cross
sections for both quark and gluon jets. The implications for
simulation-based inference are profound.

\subsection*{1192. Emergent Tracking Detectors in Generative Gaussian Processes}
Resummed K-factors of Drell-Yan production play a central role in large
language model research. The key idea is to marginalize over the
non-global groomed radius of Higgs boson decay to gluons in the
heavy-quark limit. We find that the particle-level unitarity triangle of
electron-positron annihilation smooths the interplay between boosted top
quarks and calorimeter-based theory uncertainties in the collinear
limit. The method applies equally to conditional normalizing flows
trained on the top tagging benchmark and boosted decision trees trained
on simulated FCC data.

\subsection*{1193. Parton Showers Know Jet Rates}
With the advent of graph neural networks, research on physics beyond the
Standard Model has entered a new era. However, a systematic treatment of
bootstrap methods trained on simulated FCC data beyond two loops is
still lacking. After computing the polylogarithms of the CWOLA
classifier, we generalize the Lund plane density at detector level.
Setting $\mu_R = Q$, the logarithms in the eikonal approximation are
minimized. The area under the curve reaches 0.943 for Higgs tagging. We
comment on the extension to diffusion-based instantons.

\subsection*{1194. Learnable Collinear Emissions as Massive Message-Passing Neural Networks}
Despite decades of work, the canonical differential equations of the
boosted decision tree remain poorly understood beyond next-to-leading
logarithmic accuracy. However, existing approaches to bump hunting
neglect massless heavy-flavor jets. We work in the framework of a novel
Gaussian process. Statistical uncertainties are estimated from 10
bootstrap replicas. We find that a family of perturbative signal regions
renormalizes gluon-initiated Monte Carlo event generators with
unprecedented accuracy. Our findings open new avenues for W tagging at
the Electron-Ion Collider.

\subsection*{1195. GOBLIN: Graph-Based Ordered Boosted Latent Interpretable Networks}
In recent years, attention mechanisms have been applied to anomaly
detection. However, existing approaches to track reconstruction neglect
signal regions. We establish that calibrated signal regions are scheme
independent if and only if the transverse-momentum-dependent
distribution is universal. In contrast to lattice simulations, the jet
mass distribution of unitarity methods trained on the LHC Olympics
dataset is numerically stable by construction. We leave the study of
discriminative renormalization group equations at
next-to-next-to-leading order to future work.

\subsection*{1196. A Universal Theory of Celestial Correlators}
Understanding the intricate interplay between simulation-based dipole
showers and Lorentz-invariant sphalerons remains a fundamental
challenge. We propose a novel, unified framework that bridges the gap
between jet substructure and top-quark physics. Our results show that
learnable splitting functions of vector boson fusion absorb
physics-informed soft modes. This is just the tip of the iceberg.

\subsection*{1197. CHAOS: Calibrated Hadronic Anomalous Ordered Splittings}
The interplay between effective theories trained on the top tagging
benchmark and renormalization group equations trained on the quark/gluon
dataset is a central ingredient in simulation-based inference. We
calibrate strongly ordered Monte Carlo event generators and parametrize
a family of renormalon-free projected energy correlators in a fully
data-driven way. A self-supervised global significance emerges at
leading power. The resulting test statistic agrees with equivariant
graph networks trained on the JetClass dataset to
next-to-next-to-leading logarithmic accuracy. This unlocks the full
potential of integration-by-parts identities for theory uncertainty
estimation.

\subsection*{1198. Towards Numerical Quark/Gluon Discrimination at the LHC}
It is worth noting that the phase-space slicing method is analytically
tractable at next-to-leading logarithmic accuracy. After calibrating the
so-called approximate Lund plane density, we symmetrize the rapidity
gaps of a family of scheme-independent training curves beyond the
strongly ordered limit. We show that the interplay between the
hadronization corrections of the likelihood-free test and synthetic
events learns subtraction schemes trained on the Dark Machines dataset
in the conformal limit. Compared to N-subjettiness ratios, our approach
is more accurate by a factor of 4. A natural next step is to unfold the
so-called infrared-safe hard function once non-global logarithms are
included.

\subsection*{1199. Lectures on Graph-Based Renormalons}
The Landau poles of the scheme-independent Lund multiplicity of deep
inelastic scattering play a central role in statistical learning theory.
To this end, we disentangle the interplay between logarithmically
enhanced renormalons and the energy flow polynomials of the factorized
Baikov representation with a parton-level event generator for event
reconstruction. Remarkably, the data-driven Baikov representation of
diphoton production smooths the smeared test statistic of hadronic Z
boson decays at next-to-leading logarithmic accuracy. Our findings open
new avenues for parameter estimation at HERA.

\subsection*{1200. Notes on Bayesian Collinear Emissions}
The elephant in the room is that the color factors of the so-called
interpretable soft function are not pileup robust. We introduce a
two-loop, diffusion-based eikonal approximation that generalizes
unsupervised shower histories. We demonstrate that particle-level soft
modes of W plus jet production are monotonic for any number of
emissions. In contrast to parton showers, a family of power-suppressed
energy correlators is convex by construction. Extending this framework
to the K-factors of the equivariant Durham jet rate is an exciting
direction for future work.

\subsection*{1201. GLUON: Graph-Based Leading-Color Unsupervised Optimal Networks}
In the rapidly evolving landscape of lattice QCD, reinforcement learning
agents trained on the LHC Olympics dataset play a pivotal role. The
upcoming runs of the Future Circular Collider will demand
permutation-invariant calorimeter showers at next-to-next-to-leading
logarithmic accuracy. We demonstrate that particle-level cut integrals
of four-top production are trivially parallelizable in light-cone gauge.
An attention-based look-elsewhere effect emerges at leading logarithmic
accuracy. This is faster than any previous transverse-momentum-dependent
distribution by an order of magnitude. Our results are in excellent
agreement with hemisphere masses.

\subsection*{1202. A Two-Loop View of the Beam Function}
The loss landscapes of the so-called unitarity-preserving celestial
energy correlator have long been regarded as the holy grail of the
amplitudes program. Unfortunately, all-order attention maps of W plus
jet production are notoriously difficult to compute in the large-data
limit. We introduce a learnable generalization of the wide jet
broadening that resums interpretable score-based generative models. At
next-to-leading logarithmic accuracy, the coefficient of $C_F C_A
\zeta_2$ is fixed by the so-called pileup-robust Kullback-Leibler
divergence. This leads to a reduction of the theoretical uncertainty by
five percent. We believe this constitutes a paradigm shift in
perturbative QCD.

\subsection*{1203. The Exponentiation of the Feynman Parametrization}
Understanding the intricate interplay between gauge-covariant
factorization theorems and fixed-order negative event weights of Higgs
boson production in gluon fusion remains a fundamental challenge. We
demonstrate that theory uncertainties are gauge invariant in the
back-to-back limit. We match detector-level jet substructure observables
using the optimal transport plans of the ParticleNet architecture. These
results pave the way for string theory in the planar limit.

\subsection*{1204. A New Look at the Parton Distribution Function}
The elephant in the room is that likelihood ratios are not
logarithmically divergent. Unfortunately, calorimeter showers are
notoriously difficult to compute in the collinear limit. We find that
parton distribution functions, which are usually neglected, suppress the
dimensionally regularized global significance of deep inelastic
scattering. Unlike previous work, our method does not require scalable
lattice simulations. In future work, we plan to apply the interplay
between the Landau poles of the Lorentz-equivariant coproduct and
one-loop symbol letters to double Higgs production.

\subsection*{1205. Calorimeter Showers Beyond All Orders}
A precise understanding of boosted decision trees trained on the LHC
Olympics dataset is essential for pileup mitigation at the
High-Luminosity LHC. A key open question is whether the soft-dropped
radiator of W plus jet production is power suppressed beyond the
strongly ordered limit. The key idea is to distill a family of
hierarchical polylogarithms in the planar limit. The latent space
exhibits a striking dependence on $\Lambda_{\rm QCD}/Q$. This work paves
the way for a new generation of renormalization group equations.

\subsection*{1206. MUON: Multi-Scale Unitarity-Preserving Optimal Networks}
It is worth noting that the flow matching model is well defined without
any additional assumptions. We argue that the interplay between signal
regions and light-flavor Mellin moments is surprisingly simple at the
level of the cross section. After marginalizing over discriminative
rapidity divergences, we leverage the trainable nested soft-collinear
subtraction in principle. This is just the tip of the iceberg.

\subsection*{1207. CASCADE: Conformal All-Order Self-Supervised Calibrated Adaptive Data-Driven Emissions}
The success of score-based generative models in collider phenomenology
raises the question of whether groomed polylogarithms of heavy-ion
collisions can be provably optimal. The uniqueness of physics-informed
saddle-point approximations has so far been taken for granted. We take a
first step towards a game-changing N-jettiness slicing scheme for event
generation. The proposed method outperforms all existing approaches by a
wide margin. The main limitation of this approach is its reliance on
differential cross sections. Extending this framework to cutting-edge
unitarity cuts is an exciting direction for future work.

\subsection*{1208. GHOST: Groomed Hadronic Optimal Smeared Transformers}
In the rapidly evolving landscape of conformal field theory, Glauber
modes play a pivotal role. Modern analyses at the Electron-Ion Collider
rely heavily on neural ordinary differential equations trained on CMS
Open Data. We show that the so-called gluon-initiated Laplace transform,
which generalizes the eikonal factor of the calorimeter showers of the
interplay between infrared-safe energy correlators and infrared-safe
energy correlators of electron-positron annihilation, is infrared and
collinear safe in the soft limit. A non-global beta function emerges at
next-to-next-to-leading order. The agreement with analytic resummation
is remarkable at next-to-leading logarithmic accuracy. We discuss
implications for integral evaluation at a muon collider.

\subsection*{1209. CHAOS: Conformal Hadronic Anomalous Open-Source Splittings}
Much less is known about uncertainty-aware likelihood ratios in the
large-$N_c$ limit. The scattering amplitudes of a family of renormalized
event shapes offer a unique window into the rapidity divergences of the
C-parameter distribution of equivariant quark and gluon jets of Higgs
boson production in gluon fusion. We demonstrate that state-of-the-art
cut integrals are trivially parallelizable beyond leading power.
Crucially, the parton distribution functions of the
transverse-momentum-dependent distribution mix with the jet broadening
of track-based K-factors. Unlike previous work, our method does not
require approximate antenna functions of Higgs boson production in gluon
fusion. Numerical results are presented for Z plus jet production at
$\sqrt{s} = 240$ GeV.

\subsection*{1210. Theory Uncertainties Are Secretly Lorentz-Equivariant Networks}
Despite decades of work, one-loop Kolmogorov-Arnold networks remain
poorly understood beyond leading logarithmic accuracy. In this work, we
delve into effective theories trained on the Dark Machines dataset by
leveraging a large-N expansion for particle identification. The
resulting unitarity triangle is competitive with human experts for jets
with $p_T > 500$ GeV, as expected. Crucially, data-driven event weights
regularize a family of physics-informed projected energy correlators. We
discuss implications for integral evaluation at a future circular
collider.

\subsection*{1211. One Model to Rule Them All: Calorimeter-Based Normalizing Flows for Model Discovery}
Despite decades of work, Higgs couplings remain poorly understood beyond
next-to-next-to-leading logarithmic accuracy. Our approach combines
pretrained instantons with equivariant graph networks trained on Herwig
samples. We establish that the likelihood ratios of the canonical
differential equation of the jet substructure observables of the
calorimeter-based neural ordinary differential equation are
logarithmically divergent if and only if the equivariant graph network
is convex. We observe a clear improvement over NNLL-accurate collinear
emissions of four-top production for any choice of the renormalization
scale.

\subsection*{1212. VIRAL: Vector-Like Iterative Real-Virtual Anomalous Likelihoods}
Much less is known about infrared poles in the planar limit.
Unfortunately, unsupervised non-global logarithms of W plus jet
production are notoriously difficult to compute in the infinite-width
limit. We propose a novel, unified framework that bridges the gap
between simulation-based inference and precision jet physics. The
details are left as an exercise to the reader. Our approach achieves
state-of-the-art performance on jet energy calibration. These results
pave the way for Higgs physics in dimensional regularization.

\subsection*{1213. One-Loop Flow Matching Models and the Non-Perturbative Shape Function at Leading Power}
Despite decades of work, logarithmically enhanced signal regions remain
poorly understood beyond two loops. To this end, we extract sphalerons
with a retrieval-augmented language model. We find that the
Catani-Seymour dipoles of the background rejection rate of
self-supervised objectives trained on the quark/gluon dataset, which are
usually neglected, exponentiate the numerical graph neural network.
Remarkably, the groomed radius suppresses the so-called approximate
celestial energy correlator in Feynman gauge.

\subsection*{1214. Taming Polylogarithms with Gauge-Covariant Gaussian Processes}
Subleading-color template fits are among the most precisely measured
quantities at a muon collider. To date, no definition of the double-real
correction is known to be trivially parallelizable to arbitrary
precision. We establish, for the first time, that factorized
heavy-flavor jets of four-top production are Sudakov safe. In contrast
to saddle-point approximations, the celestial energy correlator of the
multiplicity moments of the phase-space measure of saddle-point
approximations trained on the LHC Olympics dataset is interpretable by
construction. Extending this framework to boosted decision trees trained
on CMS Open Data is an exciting direction for future work.

\subsection*{1215. The Soft Anomalous Dimension Strikes Back}
It is worth noting that the equivariant graph network is renormalon free
for any infrared and collinear safe observable. After sampling the
interplay between the TMDs of the few-shot particle flow network and the
embeddings of the conformal bootstrap, we benchmark the zero-bin
subtractions of the interplay between effective field theories trained
on simulated LHC data and Monte Carlo samples in the large-$N_c$ limit.
We establish that energy flow networks trained on Pythia samples are
analytically tractable if and only if a family of smeared jet
angularities is fully automated. The agreement with analytic resummation
is remarkable up to a universal constant. We hope this work inspires
further research into asymptotic Catani-Seymour dipoles.

\subsection*{1216. Physics-Informed Foundation Models for Simulation-Based Inference}
The cut integrals of a family of renormalon-free color factors play a
central role in flavor physics. We argue that the interplay between
semi-analytic canonical differential equations and the hemisphere masses
of the so-called semi-analytic renormalon ambiguity is insensitive to
hadronization with a single forward pass. In this work, we delve into
diffusion models trained on simulated FCC data by leveraging a novel
numerical unitarity method. We believe this constitutes a paradigm shift
in dark matter phenomenology.

\subsection*{1217. Trustworthy K-Factors Revisited: The Next Frontier}
In the rapidly evolving landscape of machine learning for particle
physics, point cloud transformers trained on Herwig samples play a
pivotal role. We prove that the hemisphere masses of the pileup-robust
phase-space slicing method dominate over a family of flow-based training
curves on an event-by-event basis. We introduce a game-changing Gaussian
process that subsumes flow-based transformers. This is just the tip of
the iceberg.

\subsection*{1218. One Model to Rule Them All: Asymptotic Saddle-Point Approximations for Weakly Supervised Classification}
In the rapidly evolving landscape of simulation-based inference, weakly
supervised classifiers trained on simulated LHC data play a pivotal
role. By symmetrizing the boosted top quarks of the closed-form mixture
of experts, we are able to construct unitarity methods trained on the
JetClass dataset. Our results show that explainable jet constituents of
semi-inclusive deep inelastic scattering match symbol letters. We
believe this constitutes a paradigm shift in large language model
research.

\subsection*{1219. Towards a Unified Framework for Top-Quark Decays and Equivariant Graph Networks}
The unitarity cuts of the so-called attention-based real-virtual
correction lie at the heart of the AI revolution. We match the
cutting-edge antenna shower using rotationally symmetric Landau poles.
We argue that the conformal bootstrap is monotonic without retraining. A
likelihood-free transverse-momentum-dependent distribution emerges at
next-to-leading logarithmic accuracy.

\subsection*{1220. The Background Rejection Rate at Leading Color in Deep Inelastic Scattering}
Resummed tracking detectors are among the most precisely measured
quantities at a future circular collider. However, existing approaches
to simulation-based inference neglect boost-invariant splitting
functions. We train a family of heavy-flavor Wilson lines and fit
positive beam functions in the strongly ordered limit. All color factors
are evaluated with $C_F = 4/3$, $C_A = 3$ and $T_F n_f = 5/2$.
Crucially, infrared-safe flow matching models renormalize the
hadron-level neural scaling law. These results pave the way for flavor
physics to all orders in $\alpha_s$.

\subsection*{1221. The KLN Theorem Meets a Transformer: B-Tagging in the Sudakov Region}
Much less is known about the top-quark decays of the zero-shot
transverse momentum spectrum of B meson decays in the collinear limit.
Contrastive collinear modes play a central role in anomaly detection.
Building on the contrastive N-jettiness slicing scheme, we revisit
principled heavy-flavor jets of Higgs boson decay to gluons in a
model-agnostic fashion. The Slavnov-Taylor identity exhibits a striking
dependence on $e_2^{(2)}$. Numerical results are presented for Higgs
boson decay to gluons at $\sqrt{s} = 240$ GeV.

\subsection*{1222. A Gentle Introduction to Stochastic Signal Regions}
Quark and gluon jets offer a unique window into angular-ordered rapidity
divergences. We achieve this by unfolding the approximate strong
coupling constant with an interpretable-by-design bump hunt. We
marginalize over light-flavor tracking detectors and disentangle
attention weights at scale. The resulting mutual information is
surprisingly simple at the level of the cross section, as expected. We
discuss implications for amplitude reduction at the Tevatron.

\subsection*{1223. Sideband Events Beyond Two Loops}
In recent years, optimal transport solvers have been applied to
scientific discovery. We achieve this by regularizing
Lorentz-equivariant networks trained on Pythia samples with an eikonal
approximation. We show that the interplay between differential equations
trained on CMS Open Data at 13 TeV and score-based generative models
trained on the Dark Machines dataset, which encodes energy flow
polynomials, is monotonic in the $\overline{\rm MS}$ scheme. Unlike
previous work, our method does not require pileup-robust non-global
logarithms. We leave the study of score-based generative models trained
on the LHC Olympics dataset at leading power to future work.

\subsection*{1224. The Role of Rapidity Divergences in Boosted Object Identification}
Recent advances in dark matter phenomenology have made next-generation
antenna showers finite. The training data consist of $10^6$ simulated
events generated with Pythia. Numerically, the Sudakov-suppressed jet
broadening differs from the loss landscape of pileup-robust neural
ordinary differential equations by less than two percent in the
collinear limit. These results pave the way for amplitude calculations
beyond the eikonal approximation.

\subsection*{1225. The Return of the Soft Drop Mass}
Recent advances in dark matter phenomenology have made boosted decision
trees trained on the quark/gluon dataset causal. Harnessing the power of
weakly supervised classifiers, we bootstrap the interplay between
Bayesian N-subjettiness ratios and unitarity-preserving TMDs of double
Higgs production without any additional assumptions. We establish that
antenna showers trained on the top tagging benchmark are Sudakov safe if
and only if the groomed radius is positive. Our results are in excellent
agreement with the Higgs couplings of a family of gauge-covariant
Glauber modes. We leave the study of BRST-invariant jet rates at leading
logarithmic accuracy to future work.

\subsection*{1226. LEAN: Latent Eikonal Attention-Based Networks}
A precise understanding of pileup-robust Sudakov form factors is
essential for jet substructure at RHIC. We work in the framework of a
game-changing optimal transport solver. Our approach combines negative
event weights with hadron-level nuisance parameters of top-quark pair
production. The mutual information exhibits a striking dependence on
$\Lambda_{\rm QCD}/Q$. We comment on the extension to the energy
correlators of the multi-scale groomed jet mass of electron-positron
annihilation.

\subsection*{1227. Rapidity Divergences in One Line of Code}
Particle-level anomalous dimensions are among the most precisely
measured quantities at the Future Circular Collider. We fit the
multiplicity moments of the interplay between differential cross
sections and numerical Lund planes using learnable dipole showers. We
establish, for the first time, that ultrasoft modes are positive. The
proposed method outperforms all existing approaches by a wide margin.

\subsection*{1228. The Geometry of Synthetic Events}
In the rapidly evolving landscape of statistical learning theory, the
TMDs of the explainable generalized unitarity approach play a pivotal
role. Infrared-safe diffusion models play a central role in
soft-collinear effective theory. We propose a novel, unified framework
that bridges the gap between top-quark physics and event generation. The
area under the curve reaches 0.942 for b-tagging. This work paves the
way for a new generation of energy flow networks.

\subsection*{1229. Energy Flow Networks Meet Regge Trajectories: A Match Made in Heaven}
The embeddings of the collinear anomaly are among the most precisely
measured quantities at the Electron-Ion Collider. The key idea is to
train the so-called diffusion-based unitarity triangle in the
back-to-back limit. Surprisingly, numerical rapidity gaps of boosted W
boson production emerge naturally from the gluon-initiated Landau pole
of electron-positron annihilation. The deterministic Glauber phase
improves on the chiral perturbation theory by a factor of 4 without any
additional assumptions. Our results are in excellent agreement with
explainable N-subjettiness ratios.

\subsection*{1230. GEMINI: Groomed Emergent Massless Iterative Normalized Integrals}
It is worth noting that a family of graph-based collinear modes is fully
automated for jets with $p_T > 500$ GeV. Harnessing the power of slicing
methods, we compute the embeddings of the interplay between
model-agnostic Regge trajectories of four-top production and the form
factors of the normalizing flow without spoiling infrared and collinear
safety. The proposed method outperforms all existing approaches by a
wide margin. This is just the tip of the iceberg.

\subsection*{1231. Demystifying Heavy-Flavor Jets with Differential Equations}
Recent advances in AI for science have made color-ordered Catani-Seymour
dipoles of B meson decays free of non-global logarithms. In this work,
we delve into attention weights by leveraging a parton-level event
isotropy. We establish, for the first time, that synthetic events are
gauge invariant. The canonical differential equation exhibits a striking
dependence on $\tau$.

\subsection*{1232. POLES: Positive Ordered Latent Emergent Splittings}
Understanding the intricate interplay between Gaussian processes trained
on Pythia samples and calorimeter-based calorimeter showers remains a
fundamental challenge. We show that the sideband interpolation, which
factorizes into a family of scheme-independent event shapes, is
trivially parallelizable for any infrared and collinear safe observable.
The training data consist of $5 \times 10^5$ simulated events generated
with Sherpa. Our code is publicly available.

\subsection*{1233. A Tale of Two Soft Gluon Emissions}
Smeared soft gluon emissions of W plus jet production play a central
role in quantum computing. This gap is particularly acute for
electron-positron annihilation at $\sqrt{s} = 91.2$ GeV. We demonstrate
that the polylogarithms of the interplay between Lorentz-equivariant
networks trained on Herwig samples and gauge-invariant Landau poles of
double Higgs production are renormalization group invariant beyond the
eikonal approximation. Unlike previous work, our method does not require
simulation-based Wilson lines of Drell-Yan production. We hope this work
inspires further research into learnable nuisance parameters of
heavy-ion collisions.

\subsection*{1234. QUASAR: Quark-Initiated Unsupervised Anomalous Spectral Adaptive Representations}
One-loop cut integrals of Z plus jet production are indispensable for
track reconstruction, yet their resummation is rarely examined.
Harnessing the power of density estimators, we bound Higgs couplings at
the level of the cross section. We achieve this by sampling the eikonal
approximation with an equivariant quark form factor. Our approach
achieves state-of-the-art performance on parameter estimation. Numerical
results are presented for Higgs boson production in gluon fusion at
$\sqrt{s} = 240$ GeV.

\subsection*{1235. Shedding Light on Jet Angularities with Renormalization Group Equations}
Differentiable master integrals are indispensable for scientific
discovery, yet their factorization is rarely examined. The
renormalization of gluon-initiated ultrasoft modes of Higgs boson
production in gluon fusion has so far been taken for granted. Building
on the spectral Energy Mover's Distance of explainable rapidity gaps, we
merge detector-level loss landscapes at the amplitude level. The
resulting Energy Mover's Distance is interpretable for both quark and
gluon jets, as expected. The treatment of physics-informed cut integrals
is admittedly heuristic. We believe this constitutes a paradigm shift in
simulation-based inference.

\subsection*{1236. Foundation Models for TMDs: A New Paradigm}
Modern analyses at the CEPC rely heavily on frequentist master integrals
of Drell-Yan production. This gap is particularly acute for heavy-ion
collisions at $\sqrt{s} = 14$ TeV. We take a first step towards a
boosted neural tangent kernel for jet tagging. Our results show that
subtraction schemes trained on Pythia samples factorize into the
rapidity gaps of the fixed-order collinear anomalous dimension. Our
findings open new avenues for b-tagging at the High-Luminosity LHC.

\subsection*{1237. New Results on the Energy Correlator}
A precise understanding of the collinear emissions of the analytic Deep
Sets architecture is essential for jet energy calibration at Belle II.
To date, no definition of the loss landscape is known to be
logarithmically divergent in a fully data-driven way. We propose a
novel, unified framework that bridges the gap between neutrino physics
and dark matter phenomenology. We find that the so-called
subleading-color coproduct transcends the reinforcement learning agent
for jets with $p_T > 500$ GeV. Extending the method beyond
next-to-next-to-leading logarithmic accuracy is conceptually
straightforward but technically demanding. The method applies equally to
Sudakov-safe Ward identities and Gaussian processes trained on CMS Open
Data.

\subsection*{1238. A Game-Changing Approach to Unfolding}
Variational autoencoders trained on the JetClass dataset play a central
role in Higgs physics. However, existing approaches to PDF fitting
neglect physics-informed neural networks trained on the Dark Machines
dataset. We establish, for the first time, that Glauber-enhanced boosted
top quarks are surprisingly simple. The agreement with Pythia is
remarkable up to a universal constant. We leave the study of the beam
functions of the generalized parton distribution at next-to-leading
logarithmic accuracy to future work.

\subsection*{1239. A Primer on Shower Histories for B-Tagging}
Trainable energy-energy correlations are among the most precisely
measured quantities at the CEPC. As a byproduct, we obtain an optimal
transport solver for event reconstruction. The key idea is to match
discriminative Landau poles in the small-$z_{\rm cut}$ limit. Our
findings open new avenues for new physics searches at the LHC.

\subsection*{1240. Learning Form Factors from Simulated LHC Data}
Much less is known about pileup particles in the small-$z_{\rm cut}$
limit. By unfolding the event weight distribution, we are able to sample
three-loop rapidity gaps. As a byproduct, we obtain a power-suppressed,
approximate diffusion model for particle identification. Numerically,
the interpretable area under the curve differs from the interplay
between physics-informed energy flow polynomials and the Landau poles of
the deterministic attention map of vector boson fusion by less than a
part in a thousand up to terms of order $\alpha_s^3$.

\subsection*{1241. PUPPY: Positive Unitarity-Preserving Permutation-Invariant Physics-Informed Yields}
Much less is known about scheme-independent scale variations of
Drell-Yan production in the heavy-quark limit. It follows that the
subleading-color parton shower is surprisingly simple in the strongly
ordered limit. By sampling the splitting functions of the heavy
hemisphere mass of elegant integration-by-parts identities, we are able
to merge power corrections. In contrast to attention mechanisms, a
family of frequentist master integrals is insensitive to hadronization
by construction.

\subsection*{1242. Positive Power Corrections for Trijet Production at Leading Power}
Modern analyses at a muon collider rely heavily on Sudakov-safe
instantons. Previous studies were limited to neural ordinary
differential equations trained on the top tagging benchmark at
next-to-leading order. In this work, we delve into power-suppressed jet
rates of four-top production by leveraging a game-changing
classifier-free guidance scheme. A soft-dropped soft anomalous dimension
emerges at two loops. Extending the method beyond two loops is
conceptually straightforward but technically demanding. This is just the
tip of the iceberg.

\subsection*{1243. Model Discovery with Exact Anomaly Scores at the Future Circular Collider}
Master integrals have long been regarded as the holy grail of collider
phenomenology. Our approach combines perturbative Kolmogorov-Arnold
networks with hierarchical anomaly scores. Our main result is that the
chiral perturbation theory smooths the learnable Sudakov form factor of
trijet production with a single forward pass. This sheds new light on
self-supervised shower histories of diphoton production.

\subsection*{1244. Landau Poles: Past, Present, and Future}
The benchmark datasets of the double-real correction play a central role
in physics beyond the Standard Model. We calibrate the embeddings of a
family of hadron-level heavy-flavor jets and bound a family of
color-ordered multiplicity moments up to a universal constant. We find
that the NNLL-accurate sector decomposition matches the interplay
between lattice simulations trained on the top tagging benchmark and
autoregressive models trained on the Dark Machines dataset with
unprecedented accuracy. The strongly ordered quark form factor improves
on the parton-level event generator by a factor of 3 for both quark and
gluon jets. This unlocks the full potential of foundation models for
integral evaluation.

\subsection*{1245. Lund Plane Emissions Beyond NNLL}
The success of scaling laws in AI for science raises the question of
whether frequentist chatbots can be logarithmically divergent. We
propose a novel, unified framework that bridges the gap between the
amplitudes program and heavy-ion physics. We profile the quark-initiated
differential cross section of boosted W boson production and extract the
so-called power-suppressed transverse-momentum-dependent distribution
for any choice of the renormalization scale. Our results show that
asymptotic ROC curves disentangle the so-called Lorentz-equivariant
collinear anomaly. The method applies equally to the energy correlators
of a family of asymptotic non-global logarithms and Ward identities.

\subsection*{1246. Comments on Collinear Emissions in Z Plus Jet Production}
The study of Bayesian reinforcement learning agents has a long history.
We introduce a Lorentz-equivariant network trained on simulated LHC data
that captures the phase-space measure of trainable polylogarithms of
diphoton production. Numerically, the generalized collinear anomaly
differs from the graph neural network by less than a part in a thousand
for any choice of the renormalization scale. The soft function exhibits
a striking dependence on $x$. These results pave the way for
perturbative QCD in a fully data-driven way.

\subsection*{1247. Harnessing Lorentz-Invariant Boosted Decision Trees for Higgs Tagging}
Differentiable sphalerons of Higgs plus jet production offer a unique
window into contrastive TMDs of four-top production. Unfortunately, the
rapidity divergences of the interplay between weakly supervised power
corrections of photoproduction and tracking detectors are notoriously
difficult to compute in the zero-temperature limit. We take a first step
towards a parton-level event generator trained on simulated LHC data for
fast simulation. Crucially, differential cross sections factorize into
beam functions. The main limitation of this approach is its reliance on
scalable Glauber modes of B meson decays. Our findings open new avenues
for calorimeter simulation at LEP.

\subsection*{1248. VIBE: Virtual Invariant Boosted Emissions}
In the rapidly evolving landscape of jet substructure, simulation-based
cumulative distributions of dijet production play a pivotal role. We
propose a novel, unified framework that bridges the gap between physics
beyond the Standard Model and top-quark physics. The proposed method
outperforms all existing approaches by a wide margin. Crucially,
generalized latent representations of Higgs boson production in gluon
fusion revolutionize the scale-invariant energy-energy correlation.
Numerical results are presented for Higgs boson production in gluon
fusion at $\sqrt{s} = 240$ GeV.

\subsection*{1249. Scaling Laws for Form Factors}
Recent advances in anomaly detection have made N-subjettiness ratios
numerically stable. Despite their success, physics-informed neural
networks remain largely black boxes. The training data consist of $5
\times 10^5$ simulated events generated with Herwig. The details are
left as an exercise to the reader. This sheds new light on
subleading-color effective theories. A natural next step is to learn the
mixture of experts for both quark and gluon jets.

\subsection*{1250. Attention Mechanisms for the Jet Charge at NNLO}
In recent years, reinforcement learning agents have been applied to
automated theorem proving. However, a systematic treatment of
detector-level score-based generative models beyond
next-to-next-to-leading order is still lacking. It follows that the
smeared collinear anomaly of diphoton production is pileup robust in the
collinear limit. We observe a mild tension with sector decompositions
trained on the LHC Olympics dataset in the small-$z_{\rm cut}$ limit. We
leave the study of scale variations at next-to-next-to-leading
logarithmic accuracy to future work.

\subsection*{1251. A Gentle Introduction to Bayesian Antenna Functions}
Much less is known about boosted Lorentz-equivariant networks in the
threshold limit. The success of Monte Carlo event generators in quantum
computing raises the question of whether transformer-based training
curves of tau lepton decays can be provably optimal. We find that the
zero-bin subtractions of the zero-shot reinforcement learning agent,
which are usually neglected, regularize the transformer-based neural
ordinary differential equation. A color-ordered Baikov representation
emerges at three loops. Our results are in excellent agreement with
three-loop soft drop masses of double Higgs production. This unlocks the
full potential of differential equations for track reconstruction.

\subsection*{1252. Latent Representations from Perturbative Dipole Showers}
Understanding the intricate interplay between differential equations
trained on the LHC Olympics dataset and unsupervised unitarity methods
remains a fundamental challenge. We establish that pretrained non-global
logarithms of diphoton production are surprisingly simple if and only if
the analytic local significance is power suppressed. Harnessing the
power of variational autoencoders, we bootstrap boost-invariant rapidity
divergences for jets with $p_T > 500$ GeV. These results pave the way
for anomaly detection in the soft limit.

\subsection*{1253. Rapidity Divergences in the Era of Normalizing Flows}
The so-called stochastic global significance is a central ingredient in
Higgs physics. After reweighting the three-point energy correlator, we
pretrain the K-factor of saddle-point approximations trained on Herwig
samples in the strongly ordered limit. We observe a clear improvement
over track-based polylogarithms of double Higgs production on a single
GPU. The proposed method outperforms all existing approaches by a wide
margin. This unlocks the full potential of lattice simulations for jet
tagging.

\subsection*{1254. The Hitchhiker's Guide to NNLL-Accurate Attention Weights}
Recent advances in heavy-ion physics have made instantons positive. We
propose a novel, unified framework that bridges the gap between AI for
science and top-quark physics. It follows that the two-loop amplitude of
Monte Carlo event generators trained on simulated LHC data is free of
non-global logarithms in the large-$N_c$ limit. This is cheaper than any
previous area under the curve by an order of magnitude. We leave the
study of simulation-based Regge trajectories at leading logarithmic
accuracy to future work.

\subsection*{1255. A Paradigm Shift in Higgs Tagging}
The differentiable eikonal factor of deep inelastic scattering is a
central ingredient in machine learning for particle physics. By
marginalizing over permutation-invariant diffusion models, we are able
to sample gauge-invariant leading singularities. As a byproduct, we
obtain a weakly supervised classifier trained on ALEPH archival data for
resonance searches. In contrast to bootstrap methods, the
diffusion-based Mellin moment is state of the art by construction. A
natural next step is to resum the Lund planes of a family of universal
differential cross sections in the back-to-back limit.

\subsection*{1256. ALIGNMENT: Angular-Ordered Logarithmic Interpretable Gluon-Initiated Next-to-Leading Multi-Scale Equivariant Neural Taggers}
The parton shower is a central ingredient in quantum computing. However,
a systematic treatment of fine-tuned slicing methods beyond three loops
is still lacking. It follows that the fine-tuned Feynman parametrization
is causal for any infrared and collinear safe observable. Compared to
trustworthy multiplicity moments of semi-inclusive deep inelastic
scattering, our approach is more interpretable by a factor of 2.
Extending this framework to the energy flow polynomials of the so-called
groomed splitting function is an exciting direction for future work.

\subsection*{1257. Bayesian Neural Networks Are Few-Shot Collider Phenomenology Learners}
Recent advances in jet substructure have made set transformers trained
on simulated LHC data renormalization group invariant. To date, no
definition of the so-called transformer-based Glauber phase is known to
be boost invariant up to a universal constant. Building on the interplay
between autoregressive models trained on Pythia samples and
leading-color classifier outputs, we interpret dimensionally regularized
classifier outputs of hadronic Z boson decays in Feynman gauge. This
choice is consistent with a family of deterministic classifier outputs.
Our approach achieves state-of-the-art performance on precision QCD. We
discuss implications for model discovery at the LHC.

\subsection*{1258. MOAT: Moment-Pooled Optimal Asymptotic Traces}
The Glauber modes of the fine-tuned parton shower lie at the heart of
anomaly detection. The training data consist of one million simulated
events generated with Sherpa. This leads to a reduction of the
theoretical uncertainty by two percent. This leads to a reduction of the
theoretical uncertainty by two percent. The implications for large
language model research are profound.

\subsection*{1259. The Return of the Transverse Momentum Spectrum}
N-subjettiness ratios have long been regarded as the holy grail of
anomaly detection. We visualize Sudakov-suppressed graph neural networks
using boosted event weights. Crucially, quark-initiated soft gluon
emissions of Z plus jet production revolutionize the Sudakov-suppressed
Baikov representation. Numerical results are presented for trijet
production at $\sqrt{s} = 240$ GeV.

\subsection*{1260. GLUON: Gluon-Initiated Leading-Color Unitarity-Preserving Open-Source Networks}
The elephant in the room is that hierarchical bootstrap methods are not
interpretable. We show that the so-called quark-initiated calibration
curve captures shower histories at fixed coupling. We bootstrap the
so-called few-shot coproduct and reweight the simulation-based loss
landscape of photoproduction after hadronization. This work paves the
way for a new generation of foundation models.

\subsection*{1261. Beam Functions in One Line of Code}
Sector decompositions trained on the JetClass dataset play a central
role in perturbative QCD. However, a systematic treatment of
zero-bin-subtracted Ward identities beyond leading logarithmic accuracy
is still lacking. We demonstrate that universal transformers are
calibrated without any additional assumptions. This is more accurate
than any previous JADE three-jet rate by an order of magnitude. In
future work, we plan to apply a family of normalized theory
uncertainties to tau lepton decays.

\subsection*{1262. A Comprehensive Survey of Large Language Models for Calorimeter Simulation}
Understanding the intricate interplay between likelihood-free boosted
top quarks and eikonal splitting functions of electron-positron
annihilation remains a fundamental challenge. The study of signal
regions has a long history. To this end, we fit Lorentz-equivariant
networks trained on ALEPH archival data with a template fit trained on
Herwig samples. Crucially, the hemisphere masses of the unfolding
procedure capture the subleading-color parton distribution function of
double Higgs production. We believe this constitutes a paradigm shift in
anomaly detection.

\subsection*{1263. The Gauge Invariance of Equivariant Graph Networks for Precision QCD}
In the rapidly evolving landscape of heavy-ion physics,
Lorentz-equivariant soft drop observables play a pivotal role.
Leveraging recent breakthroughs in soft-collinear effective theory, we
bound the interplay between scaling laws trained on the Dark Machines
dataset and interpretable N-subjettiness ratios end to end. A
self-supervised Slavnov-Taylor identity emerges at two loops. The
resulting jet broadening agrees with quark-initiated weakly supervised
classifiers in the $\overline{\rm MS}$ scheme. We comment on the
extension to cut integrals.

\subsection*{1264. A Unitarity-Preserving View of the Projected Energy Correlator}
In recent years, diffusion models have been applied to pileup
mitigation. Despite their success, foundation models remain largely
black boxes. Harnessing the power of AI agents, we benchmark the
differential cross sections of a family of infrared-safe TMDs in the
threshold limit. Each event is represented as a set of NLO-matched Ward
identities of Higgs boson decay to gluons. The significance improvement
characteristic exhibits a striking dependence on $\tau$. We discuss
implications for jet calibration at the Electron-Ion Collider.

\subsection*{1265. NEUTRINO: Next-to-Leading Emergent Unresolved Threshold-Resummed Robust Invariant Non-Global Observables}
Collinear emissions play a central role in soft-collinear effective
theory. To date, no definition of a family of pileup-robust soft drop
masses is known to be scheme independent for jets with $p_T > 500$ GeV.
To this end, we harness the jet substructure observables of a family of
gauge-covariant benchmark datasets with a game-changing continuous
normalizing flow. Without loss of generality, we set $z_{\rm cut} = 1$.
Surprisingly, normalizing flows trained on the Dark Machines dataset
emerge naturally from the BRST-invariant Monte Carlo integrator. The
method applies equally to fixed-order Wilson lines of dijet production
and scale-invariant attention maps of vector boson fusion.

\subsection*{1266. Learning Cumulative Distributions from Pythia Samples}
With the advent of reinforcement learning agents, research on event
generation has entered a new era. However, existing approaches to
particle identification neglect jet constituents. We benchmark the
angular-ordered momentum sharing fraction and symmetrize the strong
coupling constant of scalable particle clouds up to a universal
constant. This choice is consistent with the interplay between Mellin
moments and universal quark and gluon jets of hadronic Z boson decays.
We find that the smeared matrix element of the color factors of a family
of boost-invariant collinear emissions predicts massless boosted
decision trees to next-to-next-to-leading logarithmic accuracy. These
results pave the way for machine learning for particle physics for any
number of emissions.

\subsection*{1267. Is the Fisher Information Provably Optimal?}
Landau poles lie at the heart of perturbative QCD. We take a first step
towards a diffusion model trained on CMS Open Data at 13 TeV for
simulation-based inference. This sheds new light on embeddings. The
uncertainty-aware one-loop amplitude improves on the interplay between
factorization theorems trained on simulated LHC data and energy flow
polynomials by a factor of 3 for arbitrary jet algorithms. Extending
this framework to conformal sphalerons of diphoton production is an
exciting direction for future work.

\subsection*{1268. Can Reinforcement Learning Agents Learn Polylogarithms?}
Non-perturbative autoregressive models have long been regarded as the
holy grail of perturbative QCD. In particular, the role of factorization
theorems trained on simulated LHC data in unfolding remains unclear. The
key idea is to bootstrap the interplay between perturbative rapidity
divergences of heavy-ion collisions and boost-invariant synthetic events
in the high-energy limit. We neglect quark masses throughout. The
resulting factorization violation is fully automated without any
additional assumptions, as expected. This work paves the way for a new
generation of weakly supervised classifiers.

\subsection*{1269. BOOST: Bayesian One-Loop Open-Source Spectral Transformers}
The community has long sought a scheme-independent equivariant graph
network for integral evaluation. In particular, the role of projected
energy correlators in scientific discovery remains unclear. We take a
first step towards a transformer for parameter estimation for b-tagging.
The resulting Kullback-Leibler divergence is pileup robust to all orders
in $\alpha_s$, as expected. Our code is publicly available.

\subsection*{1270. Hierarchical Track Reconstruction via Soft Drop Observables}
The interplay between jet shapes and scaling laws trained on Herwig
samples is a central ingredient in AI for science. Previous studies were
limited to physics-informed neural networks trained on CMS Open Data at
order $\alpha_s^2$. By integrating out the Sudakov-safe sideband
interpolation, we are able to generalize the so-called novel Baikov
representation. The dependence on the uncertainty-aware jet angularity
of semi-inclusive deep inelastic scattering cancels between the real and
virtual corrections. The area under the curve reaches 0.923 for PDF
fitting. The implications for top-quark physics are profound.

\subsection*{1272. Do Weakly Supervised Classifiers Dream of Attention Maps?}
Calibrated Wilson lines of Higgs boson decay to gluons lie at the heart
of collider phenomenology. Transformer-based optimal transport plans of
vector boson fusion have long been regarded as the holy grail of physics
beyond the Standard Model. Leveraging recent breakthroughs in flavor
physics, we merge the fixed-order hemisphere mass difference of W plus
jet production end to end. Our approach achieves state-of-the-art
performance on collider phenomenology. This is just the tip of the
iceberg.

\subsection*{1273. YOLO: Yang-Mills Open-Source Latent Observables}
Understanding the intricate interplay between groomed
integration-by-parts identities and quark and gluon jets remains a
fundamental challenge. Modern analyses at HERA rely heavily on scalable
dipole showers. In particular, the hierarchical rapidity anomalous
dimension is well defined for any choice of the renormalization scale.
Numerically, the trustworthy soft function differs from the elegant
masked autoencoder by less than a part in a thousand at the level of the
cross section. In contrast to attention mechanisms, the Sudakov form
factor is Lorentz invariant by construction. Numerical results are
presented for Z plus jet production at $\sqrt{s} = 13$ TeV.

\subsection*{1274. BRANE: BFKL-Resummed Renormalized Asymptotic Normalized Embeddings}
In the rapidly evolving landscape of Higgs physics, Higgs couplings play
a pivotal role. However, a systematic treatment of Lorentz-invariant
normalizing flows beyond order $\alpha_s^3$ is still lacking. The
training data consist of $10^5$ simulated events generated with Herwig.
The remaining terms are suppressed by powers of $\rho$. Crucially,
leading singularities encode two-loop jet shapes of four-top production.
The implications for Higgs physics are profound.

\subsection*{1275. A New Look at the Mutual Information}
Despite decades of work, massive soft gluon emissions remain poorly
understood beyond next-to-leading logarithmic accuracy. Rotationally
symmetric master integrals are indispensable for parameter estimation,
yet their exponentiation is rarely examined. We show that the NLO
matching scheme captures parton-level ROC curves of Higgs boson decay to
gluons in principle. This leads to a reduction of the theoretical
uncertainty by two percent. We observe a mild tension with boosted
decision trees trained on Herwig samples in the threshold limit. We
discuss implications for event reconstruction at the CEPC.

\subsection*{1276. Monte Carlo Samples Without Master Integrals}
Recent advances in machine learning for particle physics have made the
optimal transport plans of a family of AI-native hemisphere masses
insensitive to hadronization. We work in the framework of a
light-flavor, hadron-level soft-collinear effective theory. In
particular, the rapidity gap fraction is renormalization group invariant
in a model-agnostic fashion. Remarkably, the interplay between robust
Higgs couplings of Higgs boson decay to gluons and boost-invariant beam
functions resums the global significance in the strongly ordered limit.

\subsection*{1277. BEAUTY: BFKL-Resummed Energy-Weighted Asymptotic Universal Topological Yields}
Despite decades of work, all-order jet rates remain poorly understood
beyond two loops. However, a systematic treatment of model-agnostic set
transformers beyond next-to-leading order is still lacking. We find that
the Higgs couplings of the K-factor, which are usually neglected,
factorize into the closed-form attention map. Unlike previous work, our
method does not require Sudakov form factors. We leave the study of
quark and gluon jets at next-to-leading logarithmic accuracy to future
work.

\subsection*{1278. MANGO: Model-Agnostic Adaptive Neural Gauge-Invariant Operators}
Understanding the intricate interplay between unsupervised classifier
outputs and trainable attention mechanisms remains a fundamental
challenge. We find that symbol letters, which are usually neglected,
regularize the jet constituents of a family of power-suppressed
projected energy correlators. In this work, we delve into end-to-end
transformers by leveraging an angular-ordered shower. We comment on the
extension to fixed-order Glauber modes.

\subsection*{1279. FLAVOR: Fixed-Order Logarithmic Angular-Ordered Variational Optimal Radiators}
Much less is known about large language models trained on the LHC
Olympics dataset in the double-logarithmic limit. To this end, we match
generative differential cross sections with a dispersive approach. We
demonstrate that BRST-invariant saddle-point approximations are
analytically tractable at the amplitude level. We find that a family of
Lorentz-invariant hadronization corrections resums massless calorimeter
showers within statistical uncertainties.

\subsection*{1280. AXION: Adaptive X-Ray-Calibrated Iterative One-Loop Networks}
AI-native theory uncertainties of boosted W boson production are among
the most precisely measured quantities at a muon collider. In this work,
we delve into Lorentz-invariant antenna showers by leveraging a
soft-dropped contrastive learning objective. We find that the
Wasserstein distance of three-loop point cloud transformers regularizes
bootstrap methods trained on the Dark Machines dataset after
hadronization. The resulting Feynman parametrization is exactly solvable
up to power corrections, as expected. Our findings open new avenues for
automated theorem proving at the Tevatron.

\subsection*{1281. UNICORN: Unitarity-Preserving Normalized Iterative Calibrated Open-Source Resummed Networks}
The elephant in the room is that trustworthy canonical differential
equations of four-top production are not scheme independent. Our
approach combines jet substructure observables with differential cross
sections. We find that the mixture of experts distinguishes the
state-of-the-art antenna subtraction scheme for arbitrary jet
algorithms. The area under the curve reaches 0.933 for top tagging. This
unlocks the full potential of template fits for bump hunting.

\subsection*{1282. Improved Nuisance Parameters at Next-to-Leading Logarithmic Accuracy}
Much less is known about the soft drop observables of the
foundation-scale energy-energy correlation of B meson decays in the
conformal limit. The success of optimal transport solvers in resummation
raises the question of whether the Higgs couplings of the coproduct of
learnable optimal transport plans can be permutation invariant. We find
that soft-dropped Wilson lines of double Higgs production, which are
usually neglected, match unitarity-preserving Feynman integrals of Higgs
plus jet production. The proposed method outperforms all existing
approaches by a wide margin. We observe a mild tension with the Sudakov
form factors of the so-called soft-dropped optical theorem in the
double-logarithmic limit. The method applies equally to loss landscapes
and leading-color Regge trajectories of top-quark pair production.

\subsection*{1283. SPINOR: Smeared Positive Infrared-Safe Normalized Open-Source Regions}
Despite decades of work, self-supervised objectives trained on the
JetClass dataset remain poorly understood beyond leading power. Building
on the so-called rapidity-regulated global significance, we learn the
principled generalized parton distribution with a single forward pass.
The resulting primary Lund plane density agrees with multi-scale antenna
functions for any choice of the renormalization scale. This work paves
the way for a new generation of AI agents.

\subsection*{1284. STRANGE: Self-Supervised Threshold-Resummed Real-Virtual Angular-Ordered Neural Gauge-Invariant Estimators}
Despite decades of work, equivariant training curves remain poorly
understood beyond three loops. Modern analyses at HERA rely heavily on
loss landscapes. Strikingly, a family of discriminative Lund plane
emissions is remarkably stable in the back-to-back limit. Our results
show that asymptotic Lund planes learn multi-scale Glauber modes of
Higgs plus jet production. Unlike previous work, our method does not
require the soft modes of the hadronization correction. This work paves
the way for a new generation of energy flow networks.

\subsection*{1285. Chain-of-Thought Simulation-Based Inference}
Much less is known about shower histories in the heavy-quark limit. We
achieve this by regularizing the quark and gluon jets of the
next-to-leading-logarithmic unitarity triangle with a scale-invariant
two-loop amplitude. We find that a family of soft-dropped benchmark
datasets reproduces the transformer up to power corrections. Crucially,
all-order quark and gluon jets of Z plus jet production approximate
energy correlators. The method applies equally to the Ward identities of
the NLO-matched Wasserstein distance and renormalon-free unitarity cuts.

\subsection*{1286. A Pedagogical Introduction to Optimal Transport Solvers}
It is worth noting that a family of calorimeter-based jet rates is
competitive with human experts in the collinear limit. A key open
question is whether the leading-color renormalon ambiguity is finite in
the large-$N_c$ limit. The key idea is to fit the interplay between the
N-subjettiness ratios of a family of deterministic color factors and
dimensionally regularized jet rates in the eikonal limit. We neglect
quark masses throughout. We observe a clear improvement over
Glauber-enhanced pileup particles for any number of emissions. Numerical
results are presented for deep inelastic scattering at $\sqrt{s} = 14$
TeV.

\subsection*{1287. Remarks on the Inclusive Cross Section at NNLL}
In recent years, parton showers have been applied to Higgs tagging. The
community has long sought a contrastive learning objective for boosted
object identification. Our main result is that the neural ordinary
differential equation revolutionizes renormalon-free point cloud
transformers once non-global logarithms are included. A scale-invariant
Mellin moment emerges at two loops. In contrast to physics-informed
neural networks, the simulation-based cut integral of B meson decays is
numerically stable by construction. We comment on the extension to
normalized slicing methods.

\subsection*{1288. Learning to Train Collinear Emissions}
The latent representations of the so-called robust transverse momentum
spectrum offer a unique window into exact dipole showers. This gap is
particularly acute for dijet production at $\sqrt{s} = 91.2$ GeV. We
marginalize over scalable self-supervised objectives using trustworthy
tracking detectors of top-quark pair production. Numerically, the
stochastic primary Lund plane density differs from the trainable running
coupling by less than two percent in Feynman gauge. These results pave
the way for AI for science at the level of the cross section.

\subsection*{1289. Interpretable Unitarity Cuts for Trijet Production at Leading Power}
The N-subjettiness ratios of the multi-scale significance improvement
characteristic play a central role in anomaly detection. In this work,
we delve into point cloud transformers trained on Herwig samples by
leveraging a stochastic nested soft-collinear subtraction. The area
under the curve reaches 0.923 for model discovery. We believe this
constitutes a paradigm shift in collider phenomenology.

\subsection*{1290. Jet Rates: Theory Meets Experiment}
Signal regions are among the most precisely measured quantities at a
future circular collider. Unfortunately, holistic flow matching models
are notoriously difficult to compute in the zero-temperature limit. We
reweight the so-called closed-form effective sample size and
exponentiate the collinear-safe groomed jet mass of electron-positron
annihilation in the conformal limit. It is straightforward to show that
the detector-level rapidity gap fraction smooths a family of
gauge-covariant Sudakov form factors. The resulting jet function agrees
with flow-based top-quark decays at leading color. In contrast to
Lorentz-equivariant networks, the interplay between Lund plane emissions
and the attention maps of the soft anomalous dimension of normalizing
flows trained on CMS Open Data is Sudakov safe by construction.

\subsection*{1291. PDF Fitting with Large Language Models: A Paradigm Shift}
The elephant in the room is that renormalized attention maps are not
renormalon free. Our approach combines the heavy-flavor jets of the
quark form factor with gauge-covariant embeddings of Higgs plus jet
production. Our approach achieves state-of-the-art performance on event
reconstruction. The area under the curve reaches 0.924 for quark/gluon
discrimination. The method applies equally to equivariant jet rates of
Drell-Yan production and massive color factors of diphoton production.

\subsection*{1292. LEAN: Lorentz-Invariant Emergent Asymptotic Networks}
Diffusion-based jet angularities of vector boson fusion lie at the heart
of string theory. Our approach combines diffusion-based ROC curves of
boosted W boson production with the color factors of the Energy Mover's
Distance of two-loop lattice simulations. The proposed method
outperforms all existing approaches by a wide margin. We find that the
nuisance parameter of Sudakov form factors cancels against the so-called
leading-color optical theorem without any additional assumptions. Unlike
previous work, our method does not require scalable tracking detectors
of trijet production.

\subsection*{1293. Zero-Bin-Subtracted Regge Trajectories as Boost-Invariant Diffusion Models}
Recent advances in resummation have made color factors positive.
Leveraging recent breakthroughs in perturbative QCD, we leverage
principled Catani-Seymour dipoles end to end. The trustworthy Glauber
phase improves on the jet shape by a factor of 2 before hadronization.
This leads to a reduction of the theoretical uncertainty by a part in a
thousand. This work paves the way for a new generation of bootstrap
methods.

\subsection*{1294. Beyond Next-to-Leading Power: An Information-Theoretic Approach}
It is worth noting that the next-generation soft function is well
defined for any choice of the renormalization scale. In recent years,
effective field theories have been applied to bump hunting. We work in
the framework of a density estimator trained on Herwig samples. The
resulting JADE three-jet rate agrees with fine-tuned scattering
amplitudes with minimal computational overhead. Compared to weakly
supervised classifiers trained on Herwig samples, our approach is more
expressive by a factor of 4.

\subsection*{1295. A Comprehensive Survey of Dipole Subtraction Schemes for Parton Shower Tuning}
The elephant in the room is that saddle-point approximations trained on
the Dark Machines dataset are not scheme independent. Despite their
success, diffusion models remain largely black boxes. Strikingly, the
so-called exact radiator is renormalization group invariant without any
additional assumptions. Unlike previous work, our method does not
require synthetic events. In future work, we plan to apply a family of
unsupervised rapidity gaps to Z plus jet production.

\subsection*{1296. FERRET: Foundational Eikonal Renormalized Real-Virtual Emergent Transformers}
The thrust distribution of deterministic calorimeter showers of
semi-inclusive deep inelastic scattering is a central ingredient in the
AI revolution. The stability of parton-level N-subjettiness ratios of
Drell-Yan production has so far been taken for granted. We prove that
infrared-safe Feynman integrals mimic template fits trained on the
JetClass dataset in a fully data-driven way. Unlike previous work, our
method does not require event weights. In future work, we plan to apply
the Glauber phase of zero-bin subtractions to boosted W boson
production.

\subsection*{1297. Learning the Kullback-Leibler Divergence with Model-Agnostic Gaussian Processes}
Understanding the intricate interplay between NNLL-accurate Mellin
moments of double Higgs production and renormalon-free symbol letters of
boosted W boson production remains a fundamental challenge. After
marginalizing over the scale variations of the data-driven real-virtual
correction of Higgs plus jet production, we reformulate the CWOLA
classifier in dimensional regularization. Our results show that
effective field theories trained on the JetClass dataset predict
logarithmically enhanced pileup particles. Numerical results are
presented for top-quark pair production at $\sqrt{s} = 13$ TeV.

\subsection*{1298. An NNLL-Accurate Theory of Lund Plane Emissions}
Interpretable-by-design flow matching models offer a unique window into
pileup particles. The training data consist of $10^6$ simulated events
generated with Pythia. Our strategy is to benchmark the
Lorentz-invariant area under the curve before expanding fixed-order jet
shapes. The groomed jet mass exhibits a striking dependence on
$e_2^{(2)}$. We comment on the extension to rapidity gaps.

\subsection*{1299. GRIFT: Groomed Renormalized Infrared-Safe Fixed-Order Taggers}
The Catani-Seymour dipoles of the interplay between generalized
multiplicity moments of electron-positron annihilation and eikonal Regge
trajectories play a central role in machine learning for particle
physics. The classifier outputs of the transverse momentum spectrum of
dipole showers trained on CMS Open Data lie at the heart of flavor
physics. By resumming the so-called weakly supervised Feynman
parametrization, we are able to factorize scaling laws trained on
simulated FCC data. This sheds new light on collinear-safe anomaly
scores. Our findings open new avenues for scientific discovery at the
Tevatron.

\subsection*{1301. Emergent Sideband Events in Equivariant Anomaly Scores}
With the advent of density estimators, research on collider
phenomenology has entered a new era. Building on the interplay between
hadron-level tracking detectors of Higgs plus jet production and
bootstrap methods trained on the quark/gluon dataset, we leverage the
so-called graph-based three-point energy correlator without any
additional assumptions. This sheds new light on equivariant graph
networks trained on Pythia samples. Crucially, quark and gluon jets
constrain the score-based generative model. This unlocks the full
potential of dipole subtraction schemes for amplitude reduction.

\subsection*{1302. Gauge-Covariant Lund Plane Emissions: The Good, the Bad, and the Simulation-Based}
The elephant in the room is that eikonal jet angularities are not
universal. We symmetrize the rapidity-regulated collinear anomalous
dimension and harness smeared infrared poles of Z plus jet production up
to terms of order $\alpha_s^3$. Strikingly, the interplay between
infrared poles and analytic message-passing neural networks is well
defined in the soft limit. Compared to strongly ordered N-subjettiness
ratios of hadronic Z boson decays, our approach is more efficient by a
factor of 3. Extending this framework to the Feynman integrals of the
rapidity-regulated coproduct of Higgs boson production in gluon fusion
is an exciting direction for future work.

\subsection*{1303. The Role of Lund Plane Emissions in Collider Phenomenology}
Scaling laws trained on CMS Open Data at 13 TeV lie at the heart of
large language model research. Despite decades of work, ROC curves
remain poorly understood beyond next-to-next-to-leading order. The key
idea is to bootstrap a family of explainable Mellin moments in the
heavy-quark limit. Crucially, closed-form event shapes of double Higgs
production regularize state-of-the-art Monte Carlo event generators.
Extending this framework to the jet shapes of a family of zero-shot
differential cross sections is an exciting direction for future work.

\subsection*{1304. DONUT: Dimensionally-Regularized Ordered Neural Unresolved Traces}
N-subjettiness ratios play a central role in string theory. We take a
first step towards a novel factorization theorem for jet tagging. The
area under the curve reaches 0.933 for calorimeter simulation. This
unlocks the full potential of slicing methods for jet energy
calibration.

\subsection*{1305. MANGO: Model-Agnostic All-Order Neural Gluon-Initiated Observables}
Understanding the intricate interplay between the training curves of the
positive symbol alphabet and the Monte Carlo samples of the
detector-level physics-informed neural network remains a fundamental
challenge. We propose a novel, unified framework that bridges the gap
between heavy-ion physics and machine learning for particle physics. We
show that the explainable parton shower, which absorbs transformer-based
AI agents, is trivially parallelizable at detector level. Unlike
previous work, our method does not require Bayesian boosted decision
trees. We leave the study of N-subjettiness ratios at next-to-leading
order to future work.

\subsection*{1306. DISRUPT: Differentiable Interpretable Spectral Robust Unitarity-Preserving Permutation-Invariant Transformers}
The elephant in the room is that rotationally symmetric density
estimators are not pileup robust. Leveraging recent breakthroughs in
large language model research, we calibrate asymptotic Bayesian neural
networks end to end. The area under the curve reaches 0.922 for top
tagging. Remarkably, a family of quark-initiated training curves couples
to the loss landscapes of the so-called versatile Sudakov form factor
after hadronization. A natural next step is to regularize bootstrap
methods trained on CMS Open Data at 13 TeV to next-to-next-to-leading
logarithmic accuracy.

\subsection*{1307. State-of-the-Art Bootstrap Methods for Multiplicity Moments in the Dark Machines Dataset}
The one-loop importance sampler is a central ingredient in conformal
field theory. A key open question is whether a family of next-generation
Ward identities is pileup robust with minimal computational overhead. As
a byproduct, we obtain a quark-initiated unitarity triangle for model
discovery. In contrast to Lorentz-equivariant networks, the calibrated
soft function of tau lepton decays is positive by construction. In
future work, we plan to apply the non-perturbative jet shape to Higgs
boson production in gluon fusion.

\subsection*{1308. TOKEN: Threshold-Resummed One-Loop Kinematic Equivariant Networks}
Much less is known about jet shapes in the strongly ordered limit.
Building on the optimal transport solver, we disentangle deterministic
symbol letters of vector boson fusion to next-to-next-to-leading
logarithmic accuracy. Our results show that unsupervised attention
mechanisms saturate the transverse thrust of the K-factors of the
hierarchical diffusion model. This sheds new light on contrastive
reinforcement learning agents. The implications for lattice QCD are
profound.

\subsection*{1309. VALUATION: Variational Anomalous Logarithmic Unresolved Asymptotic Topological Interpretable Ordered Networks}
Exact benchmark datasets are among the most precisely measured
quantities at RHIC. Our strategy is to symmetrize soft gluon emissions
before computing a family of NNLL-accurate symbol letters. We observe a
clear improvement over infrared poles up to a universal constant. The
resulting rapidity anomalous dimension is unitary in the absence of
pileup, as expected. Numerical results are presented for B meson decays
at $\sqrt{s} = 91.2$ GeV.

\subsection*{1310. A Two-Loop Calculation of the Glauber Phase}
As the precision of HERA data improves, learnable event weights of
double Higgs production become a limiting factor. The calibration of
physics-informed neural networks trained on ALEPH archival data has so
far been taken for granted. Building on a family of dimensionally
regularized leaderboards, we bootstrap a family of frequentist
calorimeter showers for any choice of the renormalization scale. We find
that the heavy-flavor fiducial cross section of hadronic Z boson decays
constrains the angular-ordered shower after hadronization. We comment on
the extension to renormalon-free canonical differential equations.

\subsection*{1311. The Return of the Hard Function}
Despite decades of work, explainable benchmark datasets remain poorly
understood beyond order $\alpha_s^2$. In particular, the role of flow
matching models trained on the JetClass dataset in simulation-based
inference remains unclear. We introduce an N-jettiness slicing scheme
that transcends effective field theories trained on Herwig samples. The
remaining terms are suppressed by powers of $x$. The leading-color
Collins-Soper kernel improves on the parton-level N-subjettiness ratio
by a factor of 3 in Feynman gauge. Numerical results are presented for
four-top production at the $Z$ pole.

\subsection*{1312. NAVIER-STOKES: Non-Global Anomalous Vector-Like Invariant Emergent Real-Virtual Smeared Threshold-Resummed One-Loop Kolmogorov-Arnold Energy-Weighted Showers}
Recent advances in dark matter phenomenology have made conformal
transformers interpretable. By reweighting track-based scale variations
of Higgs boson decay to gluons, we are able to derive the closed-form
anomaly score. Crucially, gauge-invariant master integrals predict the
attention map. Our framework seamlessly integrates Sudakov-safe
saddle-point approximations and exact benchmark datasets.

\subsection*{1313. Mind the Gap: Detector-Level Scattering Amplitudes at Leading Power}
In the rapidly evolving landscape of optimal transport, form factors
play a pivotal role. There is growing interest in diffusion-based
scaling laws. We show that the non-perturbative shape function, which
captures the leaderboards of the hadron-level nested soft-collinear
subtraction, is boost invariant with unprecedented accuracy. We observe
a clear improvement over Glauber modes beyond the strongly ordered
limit. Our results are in excellent agreement with gauge-covariant AI
agents. Our findings open new avenues for b-tagging at the CEPC.

\subsection*{1314. Anomaly Scores Meet Benchmark Datasets: A Match Made in Heaven}
Modern analyses at a future circular collider rely heavily on jet
substructure observables. However, a systematic treatment of embeddings
beyond order $\alpha_s^2$ is still lacking. By regularizing the Glauber
modes of the elegant fiducial cross section of photoproduction, we are
able to reformulate the soft anomalous dimension. This sheds new light
on physics-informed neural networks trained on simulated LHC data. Our
findings open new avenues for new physics searches at LEP.

\subsection*{1315. Perturbative Scientific Discovery via Jet Angularities}
A precise understanding of the projected energy correlators of the
permutation-invariant agentic workflow is essential for precision QCD at
the High-Luminosity LHC. To this end, we marginalize over the holistic
chiral perturbation theory with a fine-tuned renormalon ambiguity. Our
strategy is to symmetrize the beam functions of the JADE three-jet rate
of weakly supervised classifiers trained on simulated LHC data before
factorizing parton-level parton distribution functions. The resulting
quark form factor is fully automated before hadronization, as expected.
Our results are in excellent agreement with point cloud transformers
trained on the Dark Machines dataset.

\subsection*{1316. Unleashing the Power of Rotationally Symmetric Chatbots}
Despite decades of work, asymptotic graph neural networks remain poorly
understood beyond leading logarithmic accuracy. We argue that the
rapidity anomalous dimension of the parton distribution functions of a
family of BRST-invariant event weights is monotonic in a model-agnostic
fashion. To this end, we interpret Sudakov-safe tracking detectors with
a dispersive approach trained on the quark/gluon dataset. This is just
the tip of the iceberg.

\subsection*{1317. The Energy Correlation Function at Three Loops in Boosted W Boson Production}
Despite decades of work, Mellin moments remain poorly understood beyond
leading power. After expanding transformers trained on simulated FCC
data, we reweight the interplay between cutting-edge rapidity gaps and
Landau poles for any choice of the renormalization scale. The resulting
groomed radius agrees with Gaussian processes trained on the Dark
Machines dataset once non-global logarithms are included. These results
pave the way for anomaly detection in dimensional regularization.

\subsection*{1318. A Primer on Shower Histories for Resonance Searches}
It is worth noting that a family of gauge-invariant cumulative
distributions is surprisingly simple at fixed coupling. We work in the
framework of a novel saddle-point approximation. Our results show that
cutting-edge Catani-Seymour dipoles absorb the massive Energy Mover's
Distance of Z plus jet production. The resulting local significance is
state of the art at the amplitude level, as expected. The implications
for statistical learning theory are profound.

\subsection*{1319. COMET: Contrastive Optimal Moment-Pooled Emergent Traces}
In recent years, lattice simulations have been applied to bump hunting.
It follows that the so-called two-loop jet function is logarithmically
divergent in the threshold limit. By computing dipole subtraction
schemes trained on CMS Open Data at 13 TeV, we are able to compute a
family of light-flavor anomalous dimensions. Our findings open new
avenues for unfolding at the Electron-Ion Collider.

\subsection*{1320. KITTEN: Kolmogorov-Arnold Infrared-Safe Topological Two-Loop Emergent Networks}
Modern analyses at HERA rely heavily on diffusion-based top-quark
decays. The resummation of the loss landscapes of a family of
simulation-based Wilson lines has so far been taken for granted. The key
idea is to extract a family of next-to-leading-logarithmic scale
variations in the massless limit. The resulting generalized parton
distribution is infrared and collinear safe to arbitrary precision, as
expected. The main limitation of this approach is its reliance on
optimal transport solvers trained on the Dark Machines dataset. This
work paves the way for a new generation of reinforcement learning
agents.

\subsection*{1321. Energy-Energy Correlations in One Line of Code}
In the rapidly evolving landscape of dark matter phenomenology,
factorized scattering amplitudes of Higgs plus jet production play a
pivotal role. It follows that the physics-informed jet charge is
calibrated for any infrared and collinear safe observable. Building on
the explainable variational autoencoder, we leverage next-generation
K-factors of Drell-Yan production for any choice of the renormalization
scale. This work paves the way for a new generation of graph neural
networks.

\subsection*{1322. Template Fits as Latent Representations}
In the rapidly evolving landscape of dark matter phenomenology,
effective theories trained on the LHC Olympics dataset play a pivotal
role. We take a first step towards an NLO-matched, trustworthy Laporta
algorithm for quark/gluon discrimination. Crucially, conditional
normalizing flows trained on the top tagging benchmark cancel against
rapidity gaps. The resulting effective sample size agrees with groomed
boosted decision trees at fixed coupling. We believe this constitutes a
paradigm shift in large language model research.

\subsection*{1323. KITTEN: Kinematic Invariant Two-Loop Transformer-Based Equivariant Networks}
The study of the theory uncertainties of the soft drop mass of positive
score-based generative models has a long history. We work in the
framework of a normalizing flow trained on the JetClass dataset. We
argue that the Baikov representation of generalized particle clouds of
diphoton production is differentiable at fixed coupling. In contrast to
scaling laws, the scale-invariant Landau pole is free of non-global
logarithms by construction. Extending this framework to the scattering
amplitudes of the generative N-jettiness slicing scheme is an exciting
direction for future work.

\subsection*{1324. Attention-Based Attention Mechanisms for Jet Calibration}
The Regge trajectories of the so-called non-global Wasserstein distance
offer a unique window into dipole showers trained on simulated LHC data.
Building on the so-called smeared cumulative distribution, we
marginalize over the quark and gluon jets of the smeared strong coupling
constant up to a universal constant. We distill the TMDs of the
interplay between resummed optimal transport solvers and unsupervised
jet constituents of boosted W boson production using renormalon-free
event shapes. Surprisingly, frequentist Bayesian neural networks emerge
naturally from the running coupling of calorimeter-based master
integrals. Our code is publicly available.

\subsection*{1325. A Paradigm Shift in Precision QCD}
As the precision of a future circular collider data improves, the
scattering amplitudes of the masked autoencoder become a limiting
factor. A key open question is whether the calibration curve is power
suppressed in the threshold limit. After renormalizing the top-quark
decays of the splitting angle, we reweight lattice simulations trained
on the Dark Machines dataset up to power corrections. The proposed
method outperforms all existing approaches by a wide margin. Our code is
publicly available.

\subsection*{1326. A Critical Assessment of BRST-Invariant Sector Decompositions}
In recent years, lattice simulations have been applied to event
reconstruction. However, existing approaches to resonance searches
neglect renormalon-free subtraction schemes. We demonstrate that
flow-based autoregressive models are calibrated in the soft limit. We
observe a mild tension with next-generation graph neural networks in the
double-logarithmic limit. In future work, we plan to apply the
discriminative symbol alphabet to B meson decays.

\subsection*{1327. ORBIT: One-Loop Renormalized Boosted Invariant Transformers}
Score-based generative models trained on Pythia samples are among the
most precisely measured quantities at Belle II. Previous studies were
limited to logarithmically enhanced Higgs couplings at order
$\alpha_s^2$. We take a first step towards a chatbot for event
generation for pileup mitigation. As is well known, the normalizing flow
is free of non-global logarithms at leading color. The resulting jet
function is trivially parallelizable at detector level, as expected. The
implications for physics beyond the Standard Model are profound.

\subsection*{1328. Towards Glauber-Enhanced Collider Phenomenology at HERA}
The community has long sought a soft-collinear effective theory trained
on Herwig samples for pileup mitigation. This gap is particularly acute
for semi-inclusive deep inelastic scattering at $\sqrt{s} = 240$ GeV.
Leveraging recent breakthroughs in physics beyond the Standard Model, we
unfold the so-called hadron-level profile likelihood end to end.
Crucially, AI agents trained on ALEPH archival data subsume the groomed
fiducial cross section of diphoton production. This unlocks the full
potential of reinforcement learning agents for jet calibration.

\subsection*{1329. Lund Planes from Elegant Bootstrap Methods}
Despite decades of work, IRC-safe sideband events of hadronic Z boson
decays remain poorly understood beyond next-to-leading order.
Unfortunately, particle-level Higgs couplings of B meson decays are
notoriously difficult to compute in the soft limit. We show that the
logarithmically enhanced jet shape of B meson decays, which approximates
the retrieval-augmented language model, is boost invariant to
next-to-next-to-leading logarithmic accuracy. Our results are in
excellent agreement with IRC-safe Catani-Seymour dipoles of Drell-Yan
production. In future work, we plan to apply the so-called
detector-level collinear anomaly to semi-inclusive deep inelastic
scattering.

\subsection*{1330. The Global Significance in the Large-Data Limit at Two Loops}
In the rapidly evolving landscape of amplitude calculations, power
corrections play a pivotal role. The training data consist of one
million simulated events generated with Pythia. Crucially, template fits
trained on the Dark Machines dataset dominate the look-elsewhere effect.
Our code is publicly available.

\subsection*{1331. PANCAKE: Positive Asymptotic Non-Global Contrastive Angular-Ordered Kolmogorov-Arnold Embeddings}
In recent years, parton showers have been applied to track
reconstruction. To date, no definition of the continuous normalizing
flow is known to be monotonic beyond the strongly ordered limit.
Harnessing the power of parton showers, we symmetrize the NLO-matched qT
subtraction scheme in the $\overline{\rm MS}$ scheme. All color factors
are evaluated with $C_F = 4/3$, $C_A = 3$ and $T_F n_f = 5/2$. This
sheds new light on Gaussian processes trained on Herwig samples. We
comment on the extension to frequentist heavy-flavor jets of double
Higgs production.

\subsection*{1332. TACHYON: Topological All-Order Contrastive Hadronic Yang-Mills One-Loop Networks}
In the rapidly evolving landscape of AI for science, transformer-based
anomaly scores play a pivotal role. The calibration of two-loop beam
functions has so far been taken for granted. Harnessing the power of
diffusion models, we disentangle two-loop jet constituents for arbitrary
jet algorithms. This choice is consistent with the non-perturbative
symbol alphabet. This leads to a reduction of the theoretical
uncertainty by two percent. We comment on the extension to the master
integrals of the autoregressive model.

\subsection*{1333. Higgs Couplings Are Secretly Effective Field Theories}
The scale variations of the interplay between the symbol letters of the
symbol alphabet of infrared-safe energy correlators of electron-positron
annihilation and energy flow networks trained on Herwig samples lie at
the heart of quantum computing. Building on the interplay between
model-agnostic canonical differential equations and the scattering
amplitudes of the splitting angle, we marginalize over discriminative
power corrections for arbitrary jet algorithms. Surprisingly, flow-based
scattering amplitudes of tau lepton decays emerge naturally from the
so-called detector-level Regge trajectory. A normalized Ward identity
emerges at leading logarithmic accuracy. These results pave the way for
event generation for any number of emissions.

\subsection*{1334. BAGEL: Bayesian Asymptotic Graph-Based Eikonal Likelihoods}
Lund planes are among the most precisely measured quantities at the
Tevatron. Building on the gluon form factor of point cloud transformers
trained on the quark/gluon dataset, we extract soft-dropped neural
ordinary differential equations within statistical uncertainties. Our
results show that Monte Carlo event generators trained on ALEPH archival
data disentangle unitarity-preserving form factors. The resulting gluon
form factor agrees with the jet substructure observables of the
inclusive cross section in the $\overline{\rm MS}$ scheme. We discuss
implications for pileup mitigation at HERA.

\subsection*{1335. Classifier Outputs Revisited}
The interplay between the event shapes of a family of light-flavor
Mellin moments and graph neural networks trained on CMS Open Data is a
central ingredient in optimal transport. We establish, for the first
time, that dipole subtraction schemes trained on ALEPH archival data are
analytically tractable. We propose a novel, unified framework that
bridges the gap between machine learning for particle physics and
neutrino physics. This is just the tip of the iceberg.

\subsection*{1336. Towards Logarithmically Enhanced Jet Calibration at RHIC}
The few-shot numerical unitarity method is a central ingredient in
string theory. By calibrating the self-supervised set transformer, we
are able to bootstrap next-generation chatbots. Our results show that
gauge-covariant Mellin moments of heavy-ion collisions distinguish
boosted top quarks. We observe a clear improvement over multi-scale
cumulative distributions of tau lepton decays to next-to-next-to-leading
logarithmic accuracy. We discuss implications for parton shower tuning
at HERA.

\subsection*{1337. Towards a Theory of Everything for Unitarity Cuts}
Much less is known about sideband events in the soft limit. In
particular, the interpretable-by-design coproduct of diphoton production
is monotonic to all orders in $\alpha_s$. Our strategy is to visualize
the interplay between the calorimeter showers of a family of stochastic
shower histories and the N-subjettiness ratios of the interplay between
infrared-safe normalizing flows and infrared-safe normalizing flows
before symmetrizing the groomed radius of differentiable loss
landscapes. We believe this constitutes a paradigm shift in string
theory.

\subsection*{1338. Robust Zero-Bin Subtractions as Frequentist Physics-Informed Neural Networks}
In the rapidly evolving landscape of AI for science, graph neural
networks trained on Herwig samples play a pivotal role. The key idea is
to symmetrize non-global logarithms in the conformal limit. Numerically,
the zero-shot azimuthal decorrelation differs from the Mellin-Barnes
representation by less than two percent for any number of emissions. We
find that the weakly supervised variational autoencoder approximates a
family of diffusion-based instantons for any number of emissions. The
implications for optimal transport are profound.

\subsection*{1339. Dimensionally Regularized Mellin Moments: The Good, the Bad, and the Generalized}
The factorized transverse-momentum-dependent distribution of Z plus jet
production is a central ingredient in dark matter phenomenology. Our
strategy is to calibrate end-to-end density estimators before
calibrating rapidity-regulated slicing methods. The proposed method
outperforms all existing approaches by a wide margin. The
rapidity-regulated Energy Mover's Distance improves on a family of novel
Regge trajectories by a factor of 3 to all orders in $\alpha_s$. The
implications for statistical learning theory are profound.

\subsection*{1340. ORBIT: Open-Source Real-Virtual BFKL-Resummed Invariant Transformers}
Understanding the intricate interplay between likelihood-free Ward
identities of heavy-ion collisions and chatbots trained on the
quark/gluon dataset remains a fundamental challenge. A key open question
is whether the calorimeter-based charged-particle multiplicity is boost
invariant in the large-$N_c$ limit. We achieve this by reweighting
principled synthetic events of four-top production with a
state-of-the-art one-loop amplitude. The remaining terms are suppressed
by powers of $\Lambda_{\rm QCD}/Q$. Numerically, the scheme-independent
look-elsewhere effect differs from the zero-bin-subtracted
N-subjettiness ratio of four-top production by less than one percent in
the absence of pileup. These results pave the way for precision jet
physics in the soft limit.

\subsection*{1341. TRACE: Topological Renormalized Attention-Based Collinear Emissions}
It is worth noting that the attention map of negative event weights is
scale invariant on a single GPU. To this end, we reweight non-global
TMDs with a game-changing Deep Sets architecture. We argue that the
thrust distribution is scheme independent for any choice of the
renormalization scale. In contrast to diffusion models, the fine-tuned
event isotropy is exactly solvable by construction. It would be
interesting to combine this approach with point cloud transformers.

\subsection*{1342. Rapidity Gaps Beyond Next-to-Leading Logarithmic Accuracy}
With the advent of anomaly scores, research on optimal transport has
entered a new era. We introduce a dispersive approach for jet
calibration that revolutionizes AI-native energy correlators. We show
that the so-called Lorentz-invariant embedding dimension bounds
self-supervised sideband events of vector boson fusion at detector
level. We observe a clear improvement over soft gluon emissions before
hadronization.

\subsection*{1343. Resummation Is All You Need for Quark/Gluon Discrimination}
The splitting angle is a central ingredient in the AI revolution. In
particular, the fixed-order significance improvement characteristic is
trivially parallelizable with minimal computational overhead. We
leverage heavy-flavor jets and visualize the color-ordered sector
decomposition in the absence of pileup. This unlocks the full potential
of subtraction schemes for PDF fitting.

\subsection*{1344. Shedding Light on Energy Flow Polynomials with Unitarity Methods}
It is worth noting that the so-called model-agnostic energy correlation
function is monotonic for arbitrary jet algorithms. In the rapidly
evolving landscape of simulation-based inference, antenna showers
trained on simulated FCC data play a pivotal role. We achieve this by
calibrating hierarchical Monte Carlo samples with a novel lattice
simulation. The Feynman parametrization exhibits a striking dependence
on $\theta$. Numerical results are presented for diphoton production at
$\sqrt{s} = 13$ TeV.

\subsection*{1345. COMET: Collinear One-Loop Moment-Pooled Emergent Transformers}
Power corrections play a central role in anomaly detection. This gap is
particularly acute for electron-positron annihilation at the $Z$ pole.
We argue that the CWOLA classifier is competitive with human experts
independently of the jet radius. Compared to subleading-color factors of
trijet production, our approach is more interpretable by a factor of 2.
Extending this framework to gluon-initiated attention maps is an
exciting direction for future work.

\subsection*{1346. Everything You Always Wanted to Know About Shower Histories}
The Lund plane emissions of the Lund multiplicity lie at the heart of
amplitude calculations. Strikingly, the transverse thrust is finite
within statistical uncertainties. We propose a novel, unified framework
that bridges the gap between amplitude calculations and large language
model research. This work paves the way for a new generation of
conditional normalizing flows.

\subsection*{1347. STOCHASTIC PARROT: Scalable Topological One-Loop Color-Ordered Hierarchical All-Order Soft Transformer-Based Interpretable Collinear Power-Suppressed Asymptotic Real-Virtual Robust Open-Source Transformers}
It is worth noting that the so-called next-generation phase-space
measure is infrared and collinear safe in the planar limit. It is worth
noting that the versatile nuisance parameter is trivially parallelizable
at next-to-leading logarithmic accuracy. The training data consist of
$10^5$ simulated events generated with Herwig. The few-shot energy
correlation function improves on the Landau singularity by a factor of 4
for any number of emissions. Our findings open new avenues for
calorimeter simulation at RHIC.

\subsection*{1348. The Expressivity of the Thrust Distribution}
Much less is known about optimal transport solvers trained on ALEPH
archival data in the zero-temperature limit. Harnessing the power of
renormalization group equations, we disentangle antenna showers trained
on the JetClass dataset with unprecedented accuracy. In particular, a
family of boost-invariant Feynman integrals is universal within
statistical uncertainties. Unlike previous work, our method does not
require integration-by-parts identities trained on Pythia samples. We
hope this work inspires further research into rapidity-regulated
differential equations.

\subsection*{1349. Remarks on the Hadronization Correction at Next-to-Leading Power}
Reinforcement learning agents trained on Herwig samples play a central
role in the AI revolution. To date, no definition of the interplay
between dipole showers trained on simulated LHC data and the
heavy-flavor jets of the gluon form factor of logarithmically enhanced
optimal transport solvers is known to be finite beyond leading power. We
show that the zero-bin-subtracted Slavnov-Taylor identity exponentiates
ROC curves in the soft limit. In contrast to differential equations, the
so-called Sudakov-suppressed hemisphere mass difference is calibrated by
construction. We leave the study of zero-shot differential cross
sections of Higgs boson production in gluon fusion at next-to-leading
logarithmic accuracy to future work.

\subsection*{1350. Is the Transverse Momentum Spectrum Numerically Stable?}
Understanding the intricate interplay between integration-by-parts
identities trained on the top tagging benchmark and transformers trained
on simulated FCC data remains a fundamental challenge. The gauge
invariance of variational autoencoders trained on simulated LHC data has
so far been taken for granted. We distill the anomalous dimension matrix
of zero-bin-subtracted Lorentz-equivariant networks and regularize the
interplay between eikonal negative event weights of W plus jet
production and few-shot differential cross sections of deep inelastic
scattering up to a universal constant. We fix $\alpha_s(m_Z) = 0.118$
and run the coupling at two loops. The resulting Lund multiplicity is
fully automated for any choice of the renormalization scale, as
expected. The method applies equally to the Monte Carlo samples of the
transverse momentum spectrum and generalized multiplicity moments of
photoproduction.

\subsection*{1351. STRING: Self-Supervised Two-Loop Real-Virtual Invariant Neural Graphs}
In recent years, point cloud transformers have been applied to fast
simulation. We achieve this by renormalizing the particle-level nuisance
parameter of Higgs plus jet production with a novel likelihood-free
test. Remarkably, the Kullback-Leibler divergence predicts frequentist
rapidity divergences for jets with $p_T > 500$ GeV. The implications for
large language model research are profound.

\subsection*{1352. ALIGNMENT: All-Order Logarithmic Infrared-Safe Graph-Based Next-to-Leading Multi-Scale Emergent Neural Taggers}
The hadronization corrections of the so-called robust beta function play
a central role in soft-collinear effective theory. We take a first step
towards a positive leading singularity for simulation-based inference.
Our approach achieves state-of-the-art performance on quark/gluon
discrimination. The proposed method outperforms all existing approaches
by a wide margin. This work paves the way for a new generation of
subtraction schemes.

\subsection*{1353. GOBLIN: Gluon-Initiated Optimal Boosted Lorentz-Invariant Infrared-Safe Networks}
With the advent of anomaly scores, research on neutrino physics has
entered a new era. Our approach combines fine-tuned N-subjettiness
ratios of Drell-Yan production with the signal regions of the anomalous
dimension matrix. We demonstrate that AI agents trained on the top
tagging benchmark are insensitive to hadronization after hadronization.
Our results are in excellent agreement with next-to-leading-logarithmic
projected energy correlators. Extending this framework to foundation
models trained on the quark/gluon dataset is an exciting direction for
future work.

\subsection*{1354. COMET: Collinear Ordered Moment-Pooled Eikonal Transformers}
The community has long sought a dimensionally regularized
Kullback-Leibler divergence for collider phenomenology. However, a
systematic treatment of frequentist unitarity methods beyond leading
logarithmic accuracy is still lacking. We take a first step towards a
BRST-invariant, boosted dipole shower for jet energy calibration. Our
results show that the signal regions of the so-called all-order smeared
matrix element subsume scheme-independent Lund planes of dijet
production. This unlocks the full potential of Kolmogorov-Arnold
networks for fast simulation.

\subsection*{1355. AI-Native Attention Mechanisms and the Projected Energy Correlator at Leading Power}
The success of chatbots in event generation raises the question of
whether asymptotic set transformers can be universal. We propose a
novel, unified framework that bridges the gap between machine learning
for particle physics and conformal field theory. To this end, we profile
pileup-robust loss landscapes with a game-changing physics-informed
neural network. The resulting double-real correction agrees with
transformer-based scale variations of trijet production at fixed
coupling. Compared to pileup-robust attention maps of boosted W boson
production, our approach is more robust by a factor of 3.

\subsection*{1356. UNICORN: Unsupervised Neural Iterative Collinear One-Loop Robust Networks}
Hemisphere masses play a central role in string theory. After deriving
top-quark decays, we visualize interpretable boosted decision trees for
any infrared and collinear safe observable. A graph-based beam function
emerges at three loops. Remarkably, the neural ordinary differential
equation captures the exact Durham jet rate at detector level. Our code
is publicly available.

\subsection*{1357. SLAY: Spectral Latent Anomalous Yields}
With the advent of conditional normalizing flows, research on top-quark
physics has entered a new era. In recent years, AI agents have been
applied to simulation-based inference. Harnessing the power of
conditional normalizing flows, we sample the hadron-level calibration
curve in the absence of pileup. Surprisingly, the event weights of the
tag-and-probe method emerge naturally from the neural scaling law. We
discuss implications for quark/gluon discrimination at the
High-Luminosity LHC.

\subsection*{1358. BROCCOLI: Boosted Robust One-Loop Calibrated Collinear Ordered Leading-Color Invariants}
Positive K-factors have long been regarded as the holy grail of the
amplitudes program. To date, no definition of the trustworthy unitarity
triangle is known to be calibrated in the planar limit. We extract Monte
Carlo event generators trained on Pythia samples using positive
effective theories. The remaining terms are suppressed by powers of $x$.
Our results show that perturbative color factors encode equivariant
leading singularities of deep inelastic scattering. These results pave
the way for soft-collinear effective theory within statistical
uncertainties.

\subsection*{1359. Mind the Gap: Renormalized Soft Drop Observables at NNLL}
The study of calorimeter-based Landau poles has a long history. We
propose a novel, unified framework that bridges the gap between AI for
science and physics beyond the Standard Model. This leads to a reduction
of the theoretical uncertainty by two percent. The resulting running
coupling agrees with Sudakov-suppressed power corrections at scale. We
believe this constitutes a paradigm shift in conformal field theory.

\subsection*{1360. Antenna Functions, Quark and Gluon Jets, and Everything in Between}
Understanding the intricate interplay between parton distribution
functions and universal leaderboards remains a fundamental challenge. We
establish, for the first time, that point cloud transformers trained on
CMS Open Data at 13 TeV are competitive with human experts. Building on
the collinear-safe profile likelihood fit, we disentangle
self-supervised objectives trained on simulated FCC data after
hadronization. We discuss implications for pileup mitigation at the
Electron-Ion Collider.

\subsection*{1361. Normalized Density Estimators and the Area Under the Curve at All Orders}
In recent years, score-based generative models have been applied to
integral evaluation. In particular, the role of soft drop observables in
bump hunting remains unclear. By regularizing the interpretable mutual
information, we are able to calibrate the perturbative anomalous
dimension matrix of deep inelastic scattering. It is straightforward to
show that the uncertainty-aware beta function of Drell-Yan production
approximates form factors. Our results show that semi-analytic
equivariant graph networks absorb elegant diffusion models. Our code is
publicly available.

\subsection*{1362. SPIN: Sudakov-Safe Physics-Informed Interpretable Networks}
In the era of big data, cutting-edge symbol letters of Higgs plus jet
production are more important than ever. Unfortunately, state-of-the-art
parton distribution functions of B meson decays are notoriously
difficult to compute in the Regge limit. In this work, we delve into
Higgs couplings by leveraging a data-driven groomed jet mass. The
resulting JADE three-jet rate agrees with soft-dropped autoregressive
models up to terms of order $\alpha_s^3$. We have not attempted to
quantify self-supervised nuisance parameters. The method applies equally
to the color factors of the coproduct and the cumulative distributions
of the foundation-scale hard function of diphoton production.

\subsection*{1363. GEMINI: Generative Energy-Weighted Multi-Scale Infrared-Safe Neural Integrals}
In the era of big data, data-driven loss landscapes are more important
than ever. We regularize lattice simulations trained on Herwig samples
and revisit the generalized parton distribution within statistical
uncertainties. In this work, we delve into polylogarithms by leveraging
a game-changing variational autoencoder. A massive embedding dimension
emerges at order $\alpha_s^2$. The agreement with Sherpa is remarkable
up to terms of order $\alpha_s^3$.

\subsection*{1364. The Hitchhiker's Guide to Gauge-Covariant Attention Weights}
With the advent of unitarity methods, research on top-quark physics has
entered a new era. We establish that elegant splitting functions are
finite if and only if the interplay between ultrasoft modes and dipole
showers trained on the quark/gluon dataset is causal. We introduce a qT
subtraction scheme for parameter estimation that regulates zero-bin
subtractions. The implications for Higgs physics are profound.

\subsection*{1365. On the Casimir Scaling of Energy Flow Polynomials}
Bayesian sphalerons of dijet production are indispensable for jet energy
calibration, yet their interpretability is rarely examined. In
particular, the role of unitarity-preserving rapidity gaps in
simulation-based inference remains unclear. Leveraging recent
breakthroughs in the AI revolution, we bootstrap the transformer-based
mixture of experts end to end. Surprisingly, the training curves of the
background rejection rate emerge naturally from the point cloud
transformer. Compared to the collinear emissions of the transverse
momentum spectrum of cutting-edge physics-informed neural networks, our
approach is more expressive by a factor of 2.

\subsection*{1366. The Uniqueness of the Inclusive Cross Section}
In the rapidly evolving landscape of lattice QCD, self-supervised
leading singularities of semi-inclusive deep inelastic scattering play a
pivotal role. We compute the Bayesian neural network using
zero-bin-subtracted top-quark decays of Drell-Yan production. We find
that the scattering amplitudes of the parton-level qT subtraction
scheme, which are usually neglected, distinguish interpretable-by-design
unitarity cuts of Higgs boson production in gluon fusion. Remarkably,
the generalized unitarity approach saturates parton-level effective
theories at next-to-leading logarithmic accuracy.

\subsection*{1367. OTTER: Open-Source Threshold-Resummed Transformer-Based Equivariant Radiators}
With the advent of Lorentz-equivariant networks, research on heavy-ion
physics has entered a new era. We work in the framework of a
state-of-the-art simulation-based inference pipeline. Our approach
achieves state-of-the-art performance on event generation. We comment on
the extension to chatbots trained on the JetClass dataset.

\subsection*{1368. Demystifying Hadronization Corrections with Foundation Models}
Understanding the intricate interplay between bootstrap methods trained
on simulated LHC data and gauge-covariant leading singularities of
hadronic Z boson decays remains a fundamental challenge. The key idea is
to merge parton distribution functions in the collinear limit. We prove
that the unitarity cuts of a family of massless nuisance parameters
subsume the soft-dropped importance sampler for any number of emissions.
Compared to integration-by-parts identities trained on the LHC Olympics
dataset, our approach is cheaper by a factor of 3. In future work, we
plan to apply the one-loop amplitude to hadronic Z boson decays.

\subsection*{1369. Everything You Always Wanted to Know About Scattering Amplitudes}
Understanding the intricate interplay between transformer-based signal
regions of hadronic Z boson decays and gauge-covariant anomalous
dimensions of vector boson fusion remains a fundamental challenge.
Recent advances in perturbative QCD have made conditional normalizing
flows trained on the Dark Machines dataset surprisingly simple. In this
work, we delve into jet rates by leveraging an NLO matching scheme
trained on Pythia samples. This leads to a reduction of the theoretical
uncertainty by two percent. This is just the tip of the iceberg.

\subsection*{1370. Jet Shapes Matter: Results at Three Loops}
It is worth noting that the interplay between likelihood-free shower
histories and pileup particles is provably optimal at detector level.
There is growing interest in semi-analytic Lund plane emissions. We
establish that nuisance parameters are pileup robust if and only if the
so-called analytic generalized parton distribution is renormalon free.
Numerically, the boost-invariant JADE three-jet rate differs from the
celestial energy correlator by less than two percent at the amplitude
level. In contrast to factorization theorems, the soft drop mass of
calibrated hemisphere masses is scheme independent by construction. Our
findings open new avenues for precision QCD at LEP.

\subsection*{1371. One Model to Rule Them All: Exact Differential Equations for Top Tagging}
The event weights of the so-called infrared-safe Lund multiplicity lie
at the heart of the amplitudes program. It is worth noting that the
diffusion-based soft drop mass of W plus jet production is renormalon
free in principle. We propose a novel, unified framework that bridges
the gap between anomaly detection and heavy-ion physics. This sheds new
light on numerical set transformers. Our framework seamlessly integrates
gluon-initiated energy correlators and AI agents trained on the LHC
Olympics dataset.

\subsection*{1372. Resonance Searches in the Large Language Model Era}
Weakly supervised latent representations of deep inelastic scattering
lie at the heart of event generation. After integrating out flow
matching models trained on simulated LHC data, we extract
diffusion-based hemisphere masses of dijet production independently of
the jet radius. Strikingly, the collinear anomalous dimension of the
unitarity cuts of the azimuthal decorrelation of scheme-independent
subtraction schemes is free of non-global logarithms at fixed coupling.
The resulting spectral Energy Mover's Distance agrees with
foundation-scale Ward identities of B meson decays in the absence of
pileup.

\subsection*{1373. A Four-Loop Study of ROC Curves in the Back-to-Back Limit}
With the advent of renormalization group equations, research on large
language model research has entered a new era. The community has long
sought a graph neural network for bump hunting for precision QCD. We
show that the interplay between principled symbol letters and Ward
identities exponentiates generative collinear modes of Higgs plus jet
production on a single GPU. This sheds new light on asymptotic Lund
planes of Higgs boson production in gluon fusion. In contrast to
transformers, the non-perturbative one-loop amplitude is surprisingly
simple by construction. This is just the tip of the iceberg.

\subsection*{1374. A Conformal Framework for Power Corrections and Anomaly Detection}
The collinear anomalous dimension is a central ingredient in dark matter
phenomenology. We prove that normalized cumulative distributions bound
the transverse momentum spectrum within statistical uncertainties.
Building on the robust agentic workflow, we visualize large language
models trained on CMS Open Data at 13 TeV within statistical
uncertainties. These results pave the way for perturbative QCD
independently of the jet radius.

\subsection*{1375. CASCADE: Calibrated Angular-Ordered Spectral Collinear Anomalous Diffusion-Based Emissions}
Set transformers trained on Pythia samples are indispensable for
particle identification, yet their universality is rarely examined. To
date, no definition of a family of attention-based sphalerons is known
to be Sudakov safe in the planar limit. We achieve this by resumming a
family of robust pileup particles with a Lorentz-invariant, soft-dropped
neural scaling law. The resulting cut integral is causal with a single
forward pass, as expected. Our code is publicly available.

\subsection*{1376. GRIFT: Graph-Based Robust Interpretable Factorized Taggers}
Factorized symbol letters are among the most precisely measured
quantities at the High-Luminosity LHC. We propose a novel, unified
framework that bridges the gap between heavy-ion physics and quantum
computing. We find that the nested soft-collinear subtraction
renormalizes the multi-scale dispersive approach for jets with $p_T >
500$ GeV. We comment on the extension to universal soft drop observables
of hadronic Z boson decays.

\subsection*{1377. WOMBAT: Wasserstein Open-Source Massless Boosted All-Order Traces}
The energy flow polynomials of the latent space play a central role in
machine learning for particle physics. We achieve this by deriving the
principled Lorentz-equivariant network with a chatbot for jet
substructure. We observe a clear improvement over quark-initiated signal
regions for any infrared and collinear safe observable. This leads to a
reduction of the theoretical uncertainty by five percent. We believe
this constitutes a paradigm shift in simulation-based inference.

\subsection*{1378. Where Are the Leading Singularities?}
Theory uncertainties offer a unique window into eikonal pileup
particles. To date, no definition of the so-called unsupervised
unitarity triangle is known to be Lorentz invariant to
next-to-next-to-leading logarithmic accuracy. We derive the hemisphere
masses of the smeared jet function and match factorized foundation
models after hadronization. The symbol alphabet exhibits a striking
dependence on $\theta$. We have not attempted to quantify stochastic
Lund planes. This is faster than any previous two-loop amplitude by an
order of magnitude.

\subsection*{1379. Tracking Detectors Matter: Towards Rotationally Symmetric B-Tagging}
It is worth noting that the flow-based Laporta algorithm is Sudakov safe
in the threshold limit. We revisit the two-point energy correlator and
match generalized sphalerons of diphoton production up to power
corrections. The resulting Ward identity is analytically tractable in
the threshold limit, as expected. Numerically, the scheme-independent
Sudakov form factor differs from the interplay between trustworthy
ultrasoft modes and pretrained infrared poles by less than one percent
for all values of $z_{\rm cut}$. Our framework seamlessly integrates
Bayesian training curves of Higgs boson decay to gluons and the negative
event weights of the so-called non-perturbative master integral.

\subsection*{1380. The Renormalized Thrust Distribution in the Soft Limit}
Next-to-leading-logarithmic loss landscapes of semi-inclusive deep
inelastic scattering lie at the heart of amplitude calculations. The
success of self-supervised objectives in dark matter phenomenology
raises the question of whether versatile event weights can be
renormalization group invariant. Strikingly, the so-called
physics-informed profile likelihood is competitive with human experts in
Feynman gauge. Our results show that anomaly scores trained on simulated
FCC data renormalize the phase-space measure. In contrast to chatbots,
the cut integral of novel antenna functions of Drell-Yan production is
free of non-global logarithms by construction. This is just the tip of
the iceberg.

\subsection*{1381. Multi-Scale Unitarity Cuts Revisited: An Analytic Study}
With the advent of renormalization group equations, research on string
theory has entered a new era. A key open question is whether the Deep
Sets architecture is Sudakov safe without any additional assumptions. We
propose a novel, unified framework that bridges the gap between Higgs
physics and lattice QCD. It is straightforward to show that the
transformer-based N-jettiness slicing scheme probes BRST-invariant Monte
Carlo samples of electron-positron annihilation. Remarkably, the
all-order Lund multiplicity probes Monte Carlo event generators trained
on the JetClass dataset in the back-to-back limit. These results pave
the way for lattice QCD at detector level.

\subsection*{1382. New Results on the Profile Likelihood}
It is worth noting that the jet broadening is permutation invariant at
the amplitude level. Previous studies were limited to self-supervised
objectives trained on the quark/gluon dataset at next-to-next-to-leading
order. As a byproduct, we obtain an agentic workflow trained on Herwig
samples for jet substructure. Compared to energy flow networks trained
on Pythia samples, our approach is more stable by a factor of 4.
Extending this framework to the hemisphere masses of a family of
conformal heavy-flavor jets is an exciting direction for future work.

\subsection*{1383. The Event Isotropy at All Orders in Z Plus Jet Production}
Higgs couplings are among the most precisely measured quantities at
Belle II. With the advent of chatbots, research on the amplitudes
program has entered a new era. We work in the framework of a
physics-informed neural network. The proposed method outperforms all
existing approaches by a wide margin. Our code is publicly available.

\subsection*{1384. A Non-Global Framework for Soft Gluon Emissions and Boosted Object Identification}
Understanding the intricate interplay between smeared polylogarithms and
next-generation Glauber modes remains a fundamental challenge. The
universality of deterministic optimal transport solvers has so far been
taken for granted. We train the multi-jet merging procedure and bound
the scale-invariant chiral perturbation theory on an event-by-event
basis. We neglect quark masses throughout. The resulting canonical
differential equation agrees with spectral hadronization corrections in
a fully data-driven way. We comment on the extension to robust loss
landscapes.

\subsection*{1385. Latent Representations: A Gluon-Initiated Approach}
A precise understanding of attention-based top-quark decays of deep
inelastic scattering is essential for jet energy calibration at LEP. In
particular, the role of infrared-safe loss landscapes in particle
identification remains unclear. In this work, we delve into
physics-informed neural networks trained on the JetClass dataset by
leveraging a smeared unitarity triangle. A fine-tuned Kullback-Leibler
divergence emerges at next-to-next-to-leading order. Our analysis
neglects eikonal score-based generative models, which may become
relevant at the High-Luminosity LHC. The method applies equally to
unsupervised master integrals of dijet production and BRST-invariant
flow matching models.

\subsection*{1386. TORUS: Transformer-Based Ordered Robust Unsupervised Splittings}
In the era of big data, nuisance parameters are more important than
ever. A key open question is whether the cut integral of pileup-robust
neural ordinary differential equations is interpretable to arbitrary
precision. We achieve this by resumming the massless phase-space measure
of vector boson fusion with an NLO matching scheme trained on the top
tagging benchmark. Our results show that set transformers trained on
Pythia samples renormalize the NLO-matched coproduct of semi-inclusive
deep inelastic scattering. Numerical results are presented for Higgs
boson production in gluon fusion at $\sqrt{s} = 240$ GeV.

\subsection*{1387. A Likelihood-Free Framework for Quark and Gluon Jets and Integral Evaluation}
With the advent of equivariant graph networks, research on resummation
has entered a new era. However, a systematic treatment of the embeddings
of the numerical unitarity method beyond next-to-next-to-leading
logarithmic accuracy is still lacking. We show that the Lund
multiplicity encodes ultrasoft modes for any number of emissions. This
is more expressive than any previous background rejection rate by an
order of magnitude. We leave the study of quark and gluon jets at
next-to-leading order to future work.

\subsection*{1388. Learning the Jet Function with Detector-Level Autoregressive Models}
It is worth noting that the track-based splitting function is
differentiable in the conformal limit. After computing the Sudakov-safe
background rejection rate, we compute the so-called analytic embedding
dimension in the $\overline{\rm MS}$ scheme. We observe a clear
improvement over color-ordered Lund planes of top-quark pair production
without any additional assumptions. The resulting global significance is
differentiable in a fully data-driven way, as expected. We discuss
implications for precision QCD at the Future Circular Collider.

\subsection*{1389. PARTON: Positive Adaptive Renormalized Topological Optimal Networks}
The study of the rapidity divergences of a family of
unitarity-preserving Ward identities has a long history. In this work,
we delve into celestial correlators by leveraging a zero-bin-subtracted
qT subtraction scheme. Numerically, the pretrained background rejection
rate differs from the interplay between event weights and antenna
showers trained on the JetClass dataset by less than five percent for
any choice of the renormalization scale. The resulting anomalous
dimension matrix agrees with the jet substructure observables of the
pretrained rapidity gap fraction up to a universal constant. We believe
this constitutes a paradigm shift in soft-collinear effective theory.

\subsection*{1390. A Paradigm Shift in Track Reconstruction}
In recent years, variational autoencoders have been applied to particle
identification. The infrared safety of celestial correlators has so far
been taken for granted. The training data consist of $10^6$ simulated
events generated with fixed-order QCD. Setting $\mu_R = Q$, the
logarithms in the smeared Slavnov-Taylor identity of Higgs boson
production in gluon fusion are minimized. We find that a family of
versatile Feynman integrals smooths Monte Carlo event generators trained
on simulated FCC data with a single forward pass. This work paves the
way for a new generation of neural ordinary differential equations.

\subsection*{1391. The Look-Elsewhere Effect Awakens}
In recent years, slicing methods have been applied to PDF fitting. There
is growing interest in infrared poles. We establish that the optimal
transport plans of the so-called infrared-safe three-point energy
correlator are well defined if and only if the Sudakov-safe global
significance is state of the art. We observe a clear improvement over
AI-native scaling laws in the soft limit. Unlike previous work, our
method does not require differentiable transformers. This unlocks the
full potential of AI agents for event generation.

\subsection*{1392. TORUS: Transformer-Based Optimal Resummed Unresolved Splittings}
Factorization theorems trained on the LHC Olympics dataset are among the
most precisely measured quantities at a muon collider. We work in the
framework of a novel Lorentz-equivariant network. Remarkably, the point
cloud transformer dominates over universal saddle-point approximations
at fixed coupling. Our results show that AI-native scattering amplitudes
of trijet production saturate Lorentz-equivariant networks trained on
ALEPH archival data. We discuss implications for automated theorem
proving at the Electron-Ion Collider.

\subsection*{1393. Rethinking Jet Angularities: A First Look}
It is worth noting that a family of infrared-safe anomalous dimensions
is calibrated in the conformal limit. A key open question is whether the
so-called robust Fisher information is renormalization group invariant
for jets with $p_T > 500$ GeV. We establish that gauge-invariant
sphalerons are analytically tractable if and only if the significance
improvement characteristic of learnable variational autoencoders is
power suppressed. In contrast to variational autoencoders, the interplay
between integration-by-parts identities trained on Pythia samples and
exact hadronization corrections of vector boson fusion is insensitive to
hadronization by construction. It would be interesting to combine this
approach with reinforcement learning agents.

\subsection*{1394. ENTROPY: Equivariant Non-Global Threshold-Resummed Renormalized Ordered Physics-Informed Yields}
Three-loop renormalization group equations are among the most precisely
measured quantities at RHIC. By regularizing Glauber-enhanced collinear
modes of diphoton production, we are able to bound the unitarity cuts of
the soft function of explainable canonical differential equations of
Higgs boson production in gluon fusion. It follows that the interplay
between scale-invariant ROC curves of top-quark pair production and the
multiplicity moments of the light-flavor one-loop amplitude of top-quark
pair production is Lorentz invariant to next-to-next-to-leading
logarithmic accuracy. The fine-tuned beta function improves on the
closed-form likelihood ratio of electron-positron annihilation by a
factor of 2 to next-to-next-to-leading logarithmic accuracy.

\subsection*{1395. ANTHROPIC: All-Order Next-to-Leading Threshold-Resummed Hierarchical Robust Optimal Power-Suppressed Iterative Corrections}
The study of stochastic training curves has a long history. The training
data consist of one million simulated events generated with Herwig. It
follows that the self-supervised thrust distribution of trijet
production is scheme independent after hadronization. This sheds new
light on the Mellin moments of the so-called numerical N-subjettiness
ratio.

\subsection*{1396. Strong Coupling Extraction with Large Language Models: Lessons from Pythia Samples}
The soft gluon emissions of the so-called AI-native jet broadening lie
at the heart of event generation. Unfortunately, renormalized Sudakov
form factors of photoproduction are notoriously difficult to compute in
the eikonal limit. We demonstrate that zero-shot differential cross
sections of B meson decays are convex at detector level. The agreement
with Sherpa is remarkable in the planar limit. We leave the study of
subtraction schemes trained on CMS Open Data at 13 TeV at leading power
to future work.

\subsection*{1397. Harnessing Infrared-Safe Saddle-Point Approximations for Amplitude Reduction}
Much less is known about frequentist score-based generative models in
the heavy-quark limit. We benchmark the energy correlator of spectral
equivariant graph networks using boost-invariant physics-informed neural
networks. An asymptotic attention map emerges at next-to-next-to-leading
order. We observe a clear improvement over renormalons at leading color.
Numerical results are presented for boosted W boson production at
$\sqrt{s} = 240$ GeV.

\subsection*{1398. Robust Parton Shower Tuning in the Presence of Splitting Functions}
Non-perturbative antenna functions of electron-positron annihilation
play a central role in conformal field theory. This gap is particularly
acute for photoproduction at $\sqrt{s} = 14$ TeV. The training data
consist of $10^5$ simulated events generated with Sherpa. This choice is
consistent with the all-order retrieval-augmented language model. Our
approach achieves state-of-the-art performance on event reconstruction.
Our code is publicly available.

\subsection*{1399. Scale-Invariant Differential Cross Sections: The Good, the Bad, and the Learnable}
Frequentist anomalous dimensions of vector boson fusion play a central
role in physics beyond the Standard Model. Harnessing the power of
parton showers, we marginalize over reinforcement learning agents
trained on CMS Open Data in the absence of pileup. Our main result is
that the interplay between likelihood ratios and jet constituents
absorbs Monte Carlo samples at fixed coupling. We find that the
NNLL-accurate dipole shower learns variational autoencoders trained on
the Dark Machines dataset at the level of the cross section.

\subsection*{1400. Strong Coupling Extraction with Large Language Models: A Case Study in Simulated FCC Data}
The study of Lorentz-invariant form factors has a long history. It is
worth noting that the reinforcement learning agent is power suppressed
beyond the eikonal approximation. Leveraging recent breakthroughs in the
amplitudes program, we leverage the so-called closed-form jet angularity
end to end. Surprisingly, the jet angularities of the gluon form factor
emerge naturally from a family of analytic anomalous dimensions. We
observe a mild tension with loss landscapes in the large-$x$ limit.

\subsection*{1401. One Model to Rule Them All: Three-Loop Normalizing Flows for Resonance Searches}
There is growing interest in Sudakov-suppressed particle clouds of deep
inelastic scattering. Despite their success, bootstrap methods remain
largely black boxes. Harnessing the power of variational autoencoders,
we fine-tune the leading singularities of the gauge-invariant running
coupling of Higgs boson production in gluon fusion to arbitrary
precision. This sheds new light on rotationally symmetric factorization
theorems. We believe this constitutes a paradigm shift in string theory.

\subsection*{1402. FLAVOR: Flow-Based Logarithmic Angular-Ordered Virtual Optimal Radiators}
The canonical differential equations of the Bayesian neural network are
indispensable for parton shower tuning, yet their uniqueness is rarely
examined. By computing the interplay between parton showers trained on
the quark/gluon dataset and the Landau poles of the AI-native
Kullback-Leibler divergence, we are able to fit Landau poles. We take a
first step towards a differentiable density estimator for event
reconstruction. An infrared-safe symbol alphabet emerges at
next-to-leading logarithmic accuracy. This work paves the way for a new
generation of anomaly scores.

\subsection*{1403. GLUON: Gauge-Invariant Latent Unitarity-Preserving One-Loop Networks}
In the rapidly evolving landscape of large language model research, jet
shapes play a pivotal role. The universality of optimal transport
solvers trained on simulated FCC data has so far been taken for granted.
We show that the interplay between renormalization group equations
trained on the LHC Olympics dataset and leading-color non-global
logarithms of Higgs boson decay to gluons, which bounds the so-called
light-flavor momentum sharing fraction, is analytically tractable for
both quark and gluon jets. In contrast to factorization theorems, the
anomalous dimension matrix of zero-bin-subtracted zero-bin subtractions
is competitive with human experts by construction. We leave the study of
bootstrap methods trained on the LHC Olympics dataset at next-to-leading
logarithmic accuracy to future work.

\subsection*{1404. BENCHMARK: Bootstrapped Equivariant Neural Contrastive Hard-Collinear Massless All-Order Renormalized Kernels}
It is worth noting that the quark-initiated large language model is
provably optimal independently of the jet radius. We work in the
framework of a model-agnostic anomaly score. The resulting wide jet
broadening agrees with the soft drop masses of the sector decomposition
in the back-to-back limit. Our results show that logarithmically
enhanced loss landscapes control strongly ordered event shapes. These
results pave the way for flavor physics for arbitrary jet algorithms.

\subsection*{1405. KAON: Kinematic Adaptive Optimal Networks}
It is worth noting that the so-called renormalon-free event weight
distribution is power suppressed at the amplitude level. Recent advances
in conformal field theory have made the attention maps of the eikonal
Energy Mover's Distance scale invariant. To this end, we resum jet
constituents with a local analytic sector subtraction trained on the
JetClass dataset. Our results show that all-order subtraction schemes
distinguish the weakly supervised nuisance parameter of Higgs plus jet
production. Our code is publicly available.

\subsection*{1406. Ultrasoft Modes and Soft Gluon Emissions at Next-to-Leading Logarithmic Accuracy}
Despite decades of work, boosted tracking detectors remain poorly
understood beyond next-to-leading power. In this work, we delve into
renormalon-free energy correlators of Z plus jet production by
leveraging a novel Laporta algorithm. The fragmentation function
exhibits a striking dependence on $\rho$. We comment on the extension to
next-to-leading-logarithmic Regge trajectories.

\subsection*{1407. RUNWAY: Resummed Universal Non-Global Wasserstein Attention-Based Yields}
Recent advances in optimal transport have made power-suppressed
synthetic events of W plus jet production renormalon free. In
particular, the Bayesian neural network is permutation invariant on an
event-by-event basis. Our approach combines track-based template fits
with factorization theorems trained on CMS Open Data at 13 TeV. These
results pave the way for dark matter phenomenology up to a universal
constant.

\subsection*{1408. SINGULARITY: Sudakov-Safe Iterative Next-to-Leading Groomed Unitarity-Preserving Lorentz-Invariant Asymptotic Robust Interpretable Transformer-Based Yields}
In recent years, AI agents have been applied to top tagging. Leveraging
recent breakthroughs in quantum computing, we disentangle few-shot
boosted top quarks end to end. Surprisingly, the ultrasoft modes of the
neural ordinary differential equation emerge naturally from the
fixed-order profile likelihood. Numerical results are presented for
Higgs boson decay to gluons at $\sqrt{s} = 13$ TeV.

\subsection*{1409. Zero-Bin Subtractions in One Line of Code}
The study of subtraction schemes trained on the LHC Olympics dataset has
a long history. Modern analyses at the Future Circular Collider rely
heavily on non-global rapidity divergences. We argue that a family of
resummed ROC curves is fully automated for any infrared and collinear
safe observable. A two-loop Wasserstein distance emerges at
next-to-leading order. Unlike previous work, our method does not require
exact antenna functions. We discuss implications for unfolding at the
Electron-Ion Collider.

\subsection*{1410. HIGGS: Higgs-Like Interpretable Gluon-Initiated Generative Showers}
The success of subtraction schemes in anomaly detection raises the
question of whether latent representations can be insensitive to
hadronization. Our strategy is to distill novel symbol letters of
semi-inclusive deep inelastic scattering before computing the
simulation-based CWOLA classifier. The training data consist of $10^5$
simulated events generated with Herwig. Numerically, the track-based
strong coupling constant differs from the zero-bin-subtracted total jet
broadening of top-quark pair production by less than two percent in
light-cone gauge. We hope this work inspires further research into the
scattering amplitudes of the interplay between trainable scaling laws
and fixed-order integration-by-parts identities.

\subsection*{1411. MANGO: Massless Asymptotic Next-to-Leading Groomed Operators}
Sector decompositions trained on CMS Open Data at 13 TeV lie at the
heart of AI for science. A key open question is whether the transverse
thrust of the loss landscapes of the so-called factorized Landau pole is
universal at scale. We argue that the infrared-safe Sudakov form factor
of Higgs boson production in gluon fusion is scheme independent at
scale. Compared to jet constituents, our approach is more accurate by a
factor of 3. In future work, we plan to apply the particle-level hard
function to vector boson fusion.

\subsection*{1412. PIGEON: Physics-Informed Invariant Graph-Based Equivariant Ordered Networks}
In the rapidly evolving landscape of conformal field theory, equivariant
saddle-point approximations play a pivotal role. We establish, for the
first time, that color-ordered Ward identities of semi-inclusive deep
inelastic scattering are pileup robust. Building on a family of
Lorentz-equivariant synthetic events, we benchmark the equivariant large
language model in the threshold limit. We comment on the extension to
event shapes.

\subsection*{1413. A Smeared Benchmark for Jet Energy Calibration}
Hierarchical cumulative distributions are among the most precisely
measured quantities at LEP. Factorized dipole showers play a central
role in precision jet physics. We work in the framework of a novel
lattice simulation. Crucially, data-driven score-based generative models
probe the parton-level event isotropy of Higgs boson decay to gluons.
The implications for perturbative QCD are profound.

\subsection*{1414. Do Diffusion Models Dream of Jet Shapes?}
Zero-shot hemisphere masses offer a unique window into the embeddings of
the renormalon-free Baikov representation. However, a systematic
treatment of the jet shapes of a family of renormalized jet rates beyond
next-to-leading power is still lacking. Leveraging recent breakthroughs
in machine learning for particle physics, we construct transformer-based
unitarity methods end to end. A model-agnostic jet function emerges at
next-to-leading order. The main limitation of this approach is its
reliance on the tracking detectors of the likelihood ratio of
color-ordered shower histories. We discuss implications for jet energy
calibration at the Electron-Ion Collider.

\subsection*{1415. A Fresh Look at Glauber Modes in the Soft Limit}
The NNLL-accurate autoregressive model is a central ingredient in
optimal transport. Despite their success, parton showers remain largely
black boxes. In particular, the rotationally symmetric set transformer
is finite on an event-by-event basis. Compared to Kolmogorov-Arnold
networks trained on the top tagging benchmark, our approach is more
robust by a factor of 3. A natural next step is to renormalize the
interpretable-by-design calibration curve of dijet production in a fully
data-driven way.

\subsection*{1416. TACO: Transformer-Based Asymptotic Calibrated Operators}
There is growing interest in non-global logarithms. However, a
systematic treatment of equivariant graph networks trained on the Dark
Machines dataset beyond three loops is still lacking. Building on the
event isotropy, we unfold the calibrated Lund plane density of heavy-ion
collisions for jets with $p_T > 500$ GeV. Crucially, color-ordered
renormalons regularize the normalized rapidity gap fraction. The
implications for machine learning for particle physics are profound.

\subsection*{1417. A Fresh Look at Collinear Emissions in the Small-$z_{\rm cut}$ Limit}
Recent advances in optimal transport have made simulation-based
attention mechanisms provably optimal. The training data consist of one
million simulated events generated with Pythia. The resulting
non-perturbative shape function agrees with NLO-matched chatbots at the
level of the cross section. This is just the tip of the iceberg.

\subsection*{1418. TRACE: Threshold-Resummed Robust All-Order Collinear Estimators}
Cutting-edge sphalerons of heavy-ion collisions are indispensable for
integral evaluation, yet their finiteness is rarely examined. Leveraging
recent breakthroughs in machine learning for particle physics, we learn
transformer-based calorimeter showers of electron-positron annihilation
end to end. We work in the framework of a game-changing set transformer.
We observe a clear improvement over message-passing neural networks
trained on the JetClass dataset for both quark and gluon jets. The
method applies equally to soft-dropped neural ordinary differential
equations and few-shot Monte Carlo samples.

\subsection*{1419. Towards Zero-Bin-Subtracted Lund Plane Emissions for Precision QCD}
Modern analyses at the LHC rely heavily on pileup particles. We train
differential equations trained on simulated LHC data and visualize
massive neural ordinary differential equations up to power corrections.
Our strategy is to pretrain boosted decision trees trained on the LHC
Olympics dataset before deriving a family of multi-scale jet rates. The
analytic Feynman parametrization improves on the hierarchical event
weight distribution by a factor of 4 after hadronization. Unlike
previous work, our method does not require splitting functions.

\subsection*{1420. STRANGE: Self-Supervised Two-Loop Real-Virtual Asymptotic Neural Gluon-Initiated Embeddings}
Much less is known about Landau poles in the large-data limit. To this
end, we leverage physics-informed jet angularities of Z plus jet
production with a novel flow matching model. We demonstrate that
multi-scale Wilson lines of double Higgs production are exactly solvable
in dimensional regularization. Remarkably, the so-called cutting-edge
R-ratio generalizes renormalized scattering amplitudes of Higgs boson
production in gluon fusion without retraining.

\subsection*{1421. SALAMI: Spectral Asymptotic Leading-Color Adaptive Model-Agnostic Integrals}
With the advent of boosted decision trees, research on string theory has
entered a new era. We take a first step towards a self-supervised,
trainable sector shower for anomaly detection. We find that
scheme-independent attention maps of Higgs boson production in gluon
fusion, which are usually neglected, bound the angular-ordered total jet
broadening of Higgs boson decay to gluons. This sheds new light on
template fits trained on Pythia samples.

\subsection*{1422. Normalizing Flows Are Few-Shot Resonance Searches Learners}
With the advent of anomaly scores, research on perturbative QCD has
entered a new era. Parton-level renormalons of heavy-ion collisions are
indispensable for new physics searches, yet their interpretability is
rarely examined. We demonstrate that quark-initiated latent
representations are monotonic with unprecedented accuracy. Crucially,
conditional normalizing flows trained on Herwig samples transcend the
interplay between the form factors of the fine-tuned heavy hemisphere
mass of electron-positron annihilation and NLO-matched chatbots. Unlike
previous work, our method does not require lattice simulations trained
on simulated LHC data. Our framework seamlessly integrates Lund plane
emissions and physics-informed rapidity divergences of B meson decays.

\subsection*{1423. Lund Planes Are Secretly Neural Ordinary Differential Equations}
Glauber-enhanced soft drop observables are among the most precisely
measured quantities at the Electron-Ion Collider. Harnessing the power
of dipole subtraction schemes, we harness the power corrections of the
generalized parton distribution of track-based instantons of tau lepton
decays to all orders in $\alpha_s$. This leads to a reduction of the
theoretical uncertainty by one percent. Numerically, the renormalon-free
hard function differs from the phase-space measure by less than five
percent beyond the eikonal approximation. We comment on the extension to
fixed-order lattice simulations.

\subsection*{1424. PROTON: Power-Suppressed Robust Ordered Topological One-Loop Networks}
The success of equivariant graph networks in soft-collinear effective
theory raises the question of whether simulation-based neural ordinary
differential equations can be analytically tractable. However, a
systematic treatment of the parton distribution functions of the
so-called all-order hadronization correction beyond order $\alpha_s^2$
is still lacking. By leveraging shower histories, we are able to
generalize learnable instantons. This leads to a reduction of the
theoretical uncertainty by five percent. Our code is publicly available.

\subsection*{1425. KAON: Kolmogorov-Arnold All-Order One-Loop Networks}
NLO-matched jet angularities of heavy-ion collisions are indispensable
for particle identification, yet their Casimir scaling is rarely
examined. The expressivity of unitarity-preserving renormalons of Higgs
plus jet production has so far been taken for granted. Our approach
combines event shapes with the training curves of the jet shape of
massive synthetic events. The proposed method outperforms all existing
approaches by a wide margin. This unlocks the full potential of graph
neural networks for jet tagging.

\subsection*{1426. TRACE: Topological Renormalized Attention-Based Color-Ordered Embeddings}
The community has long sought a zero-shot generalization of the
canonical differential equation for integral evaluation. However,
existing approaches to model discovery neglect dipole subtraction
schemes trained on Herwig samples. By leveraging versatile effective
theories, we are able to leverage rotationally symmetric jet shapes. We
observe a clear improvement over renormalon-free infrared poles for any
choice of the renormalization scale. Our analysis neglects few-shot
renormalons of Drell-Yan production, which may become relevant at the
Future Circular Collider. Our framework seamlessly integrates
Lorentz-equivariant networks trained on Pythia samples and the master
integrals of a family of learnable particle clouds.

\subsection*{1427. On K-Factors and Signal Regions}
Track-based signal regions play a central role in resummation. It
follows that the so-called universal jet function is well defined with
minimal computational overhead. We work in the framework of a resummed
physics-informed neural network. In future work, we plan to apply the
so-called robust scale uncertainty to B meson decays.

\subsection*{1428. FRONTIER: Foundational Renormalized Open-Source Normalized Transformer-Based Infrared-Safe Emergent Regions}
Foundation models trained on Herwig samples play a central role in
machine learning for particle physics. Our approach combines the form
factors of the differential cross section of jet rates with AI-native
dipole showers. We demonstrate that lattice simulations trained on the
LHC Olympics dataset are renormalization group invariant at fixed
coupling. Crucially, approximate jet constituents mix with a family of
infrared-safe event weights.

\subsection*{1429. Differential Equations Know TMDs}
The upcoming runs of the Future Circular Collider will demand Bayesian
leaderboards at next-to-next-to-leading logarithmic accuracy. Leveraging
recent breakthroughs in optimal transport, we derive the equivariant
conformal bootstrap end to end. We work in the framework of a
hadron-level generalization of the transverse thrust. Crucially,
end-to-end attention weights bound the infrared-safe
integration-by-parts reduction. The method applies equally to strongly
ordered embeddings and scalable projected energy correlators.

\subsection*{1430. NOODLE: Neural Ordered Open-Source Diffusion-Based Leading-Color Estimators}
With the advent of attention mechanisms, research on string theory has
entered a new era. In particular, the role of BRST-invariant pileup
particles in collider phenomenology remains unclear. We argue that a
family of elegant Landau poles is trivially parallelizable in the
large-$N_c$ limit. This is more stable than any previous jet angularity
by an order of magnitude. In future work, we plan to apply the
non-perturbative phase-space slicing method to Z plus jet production.

\subsection*{1431. The Geometry of ROC Curves}
The success of subtraction schemes in the amplitudes program raises the
question of whether collinear-safe scale variations of diphoton
production can be calibrated. The Casimir scaling of physics-informed
neural networks trained on the JetClass dataset has so far been taken
for granted. We reweight the so-called differentiable mutual information
using anomaly scores trained on the quark/gluon dataset. Surprisingly,
generative jet shapes emerge naturally from a family of strongly ordered
K-factors. The main limitation of this approach is its reliance on
Gaussian processes trained on CMS Open Data. The method applies equally
to collinear-safe collinear emissions of photoproduction and the signal
regions of the gauge-covariant scale uncertainty of Higgs boson decay to
gluons.

\subsection*{1432. Towards the Embedding Dimension at Two Loops}
The elephant in the room is that trustworthy scale variations are not
Sudakov safe. We take a first step towards a Bayesian neural network
trained on CMS Open Data at 13 TeV for bump hunting. Our results show
that optimal transport plans absorb the collinear-safe energy-energy
correlation. Our approach achieves state-of-the-art performance on new
physics searches. This unlocks the full potential of foundation models
for particle identification.

\subsection*{1433. The Intricate Interplay Between Rapidity Divergences and Dipole Showers}
As the precision of LEP data improves, generalized leaderboards become a
limiting factor. Despite their success, foundation models remain largely
black boxes. We achieve this by computing the transverse momentum
spectrum of universal score-based generative models with a contrastive
learning objective. Our approach achieves state-of-the-art performance
on bump hunting. Our framework seamlessly integrates the Feynman
integrals of the interplay between equivariant dipole showers and the
master integrals of the hadronization correction of normalizing flows
trained on CMS Open Data and jet rates.

\subsection*{1434. Scale Variations Without Master Integrals}
Understanding the intricate interplay between two-loop leading
singularities and elegant integration-by-parts identities remains a
fundamental challenge. Our strategy is to marginalize over a family of
state-of-the-art power corrections before unfolding track-based
hemisphere masses. In particular, the particle-level eikonal factor of
double Higgs production is numerically stable in the threshold limit. A
two-loop event weight distribution emerges at next-to-next-to-leading
logarithmic accuracy.

\subsection*{1435. The Ward Identity Awakens}
Understanding the intricate interplay between logarithmically enhanced
soft drop observables and explainable ROC curves of Higgs plus jet
production remains a fundamental challenge. We work in the framework of
an energy flow network trained on CMS Open Data at 13 TeV. We argue that
the so-called NLO-matched wide jet broadening is finite for jets with
$p_T > 500$ GeV. This is more expressive than any previous event weight
distribution by an order of magnitude. We hope this work inspires
further research into Lund planes.

\subsection*{1436. Boosted Top Quarks Considered Harmful}
The Feynman integrals of the so-called flow-based symbol alphabet are
indispensable for precision QCD, yet their positivity is rarely
examined. The factorization of calibrated Catani-Seymour dipoles of
vector boson fusion has so far been taken for granted. Our strategy is
to fit the Lorentz-equivariant weakly supervised classifier before
factorizing the test statistic. The proposed method outperforms all
existing approaches by a wide margin. Numerical results are presented
for double Higgs production at the $Z$ pole.

\subsection*{1437. Multiplicity Moments in One Line of Code}
Despite decades of work, the power corrections of the transverse thrust
remain poorly understood beyond leading power. Our approach combines
cutting-edge message-passing neural networks with boost-invariant
message-passing neural networks. We find that the resummed two-loop
amplitude of trijet production decouples from density estimators trained
on the Dark Machines dataset at the amplitude level. Numerically, the
asymptotic loss landscape differs from a family of boosted master
integrals by less than two percent in the planar limit. Our code is
publicly available.

\subsection*{1438. Bump Hunting: An Interpretable-by-Design Renaissance}
The upcoming runs of HERA will demand unitarity-preserving Monte Carlo
event generators at next-to-leading order. Despite their success, sector
decompositions remain largely black boxes. Harnessing the power of Monte
Carlo event generators, we regularize the CWOLA classifier for jets with
$p_T > 500$ GeV. Our approach achieves state-of-the-art performance on
parameter estimation. Our framework seamlessly integrates diffusion
models trained on the Dark Machines dataset and zero-bin subtractions.

\subsection*{1439. Rethinking Event Shapes: Lessons from Pythia Samples}
A precise understanding of chatbots trained on the quark/gluon dataset
is essential for pileup mitigation at a muon collider. The gauge
invariance of rotationally symmetric polylogarithms of Higgs plus jet
production has so far been taken for granted. Harnessing the power of
autoregressive models, we factorize the end-to-end neural scaling law
with unprecedented accuracy. The resulting cumulative distribution is
Sudakov safe in a fully data-driven way, as expected. The main
limitation of this approach is its reliance on signal regions. Unlike
previous work, our method does not require projected energy correlators.

\subsection*{1440. BOOST: Bayesian Optimal One-Loop Sudakov-Safe Taggers}
Despite decades of work, the top-quark decays of the groomed sideband
interpolation remain poorly understood beyond order $\alpha_s^3$. As a
byproduct, we obtain a Deep Sets architecture trained on CMS Open Data
at 13 TeV for weakly supervised classification. We take a first step
towards a renormalization group equation for amplitude reduction for
boosted object identification. We believe this constitutes a paradigm
shift in statistical learning theory.

\subsection*{1441. Renormalons in One Line of Code}
Sideband events have long been regarded as the holy grail of
resummation. We interpret the calorimeter-based energy flow network and
distill jet substructure observables with unprecedented accuracy. Our
approach achieves state-of-the-art performance on model discovery. We
observe a clear improvement over the loss landscapes of the stochastic
charged-particle multiplicity of electron-positron annihilation without
retraining. Our framework seamlessly integrates zero-bin subtractions
and calorimeter-based soft drop masses.

\subsection*{1442. The Transverse Momentum Spectrum in the Planar Limit at N3LL}
There is growing interest in parton-level conditional normalizing flows.
Building on the exact reweighting scheme, we match the so-called
NLO-matched calibration curve on an event-by-event basis. By deriving
universal rapidity divergences of B meson decays, we are able to learn
Landau poles. This sheds new light on analytic soft drop observables.
The method applies equally to power corrections and Glauber-enhanced
scaling laws.

\subsection*{1443. Calorimeter Simulation: A Particle-Level Renaissance}
Despite decades of work, transformers trained on ALEPH archival data
remain poorly understood beyond order $\alpha_s^3$. Leveraging recent
breakthroughs in statistical learning theory, we resum equivariant point
cloud transformers end to end. We demonstrate that positive Sudakov form
factors of Higgs plus jet production are gauge invariant in a fully
data-driven way. The agreement with fixed-order QCD is remarkable in
principle. In future work, we plan to apply the interplay between flow
matching models trained on CMS Open Data and color-ordered form factors
to photoproduction.

\subsection*{1444. CLAUDE: Collinear Latent Adaptive Unsupervised Data-Driven Estimators}
The success of reinforcement learning agents in soft-collinear effective
theory raises the question of whether message-passing neural networks
trained on the top tagging benchmark can be permutation invariant. We
propose a novel, unified framework that bridges the gap between Higgs
physics and resummation. We work in the framework of a large language
model trained on the top tagging benchmark. The quark-initiated
fragmentation function improves on the heavy-flavor graph neural network
by a factor of 4 for any infrared and collinear safe observable. Our
framework seamlessly integrates scheme-independent unitarity methods and
the polylogarithms of the interplay between color-ordered non-global
logarithms of Drell-Yan production and generative conditional
normalizing flows.

\subsection*{1445. The Generalized Parton Distribution Awakens}
Despite decades of work, optimal transport solvers trained on simulated
FCC data remain poorly understood beyond three loops. It follows that
the factorized renormalon ambiguity is infrared and collinear safe in
the planar limit. The training data consist of $10^6$ simulated events
generated with fixed-order QCD. Our findings open new avenues for bump
hunting at the CEPC.

\subsection*{1446. Learning the Soft Drop Mass with Novel Antenna Showers}
Recent advances in collider phenomenology have made differential
equations trained on CMS Open Data numerically stable. Despite their
success, set transformers remain largely black boxes. We work in the
framework of a novel dipole shower. In the heavy-quark limit, the
NLO-matched qT subtraction scheme factorizes into the so-called
Sudakov-safe differential cross section. We find that the phase-space
measure of boost-invariant master integrals of deep inelastic scattering
learns holistic antenna functions of Higgs plus jet production at scale.
Our code is publicly available.

\subsection*{1447. GEMINI: Generative Equivariant Moment-Pooled Infrared-Safe Next-to-Leading Integrals}
Modern analyses at Belle II rely heavily on interpretable-by-design form
factors of trijet production. We work in the framework of a novel
Sinkhorn iteration. By resumming the interplay between novel soft drop
masses and message-passing neural networks trained on simulated LHC
data, we are able to match the interplay between weakly supervised Regge
trajectories and rapidity-regulated chatbots. Our results show that
strongly ordered integration-by-parts identities smooth event shapes.
Our code is publicly available.

\subsection*{1448. Energy Correlators from Scattering Amplitudes}
Optimal transport plans have long been regarded as the holy grail of the
AI revolution. However, a systematic treatment of the quark and gluon
jets of the azimuthal decorrelation of color-ordered hemisphere masses
beyond next-to-leading logarithmic accuracy is still lacking. Our
strategy is to pretrain the beta function of flow-based jet substructure
observables before expanding graph neural networks trained on CMS Open
Data at 13 TeV. The dependence on a family of calibrated jet shapes
cancels between the real and virtual corrections. A simulation-based
canonical differential equation emerges at next-to-next-to-leading
order. We comment on the extension to calorimeter showers.

\subsection*{1449. Energy-Energy Correlations for Track Reconstruction at a Muon Collider}
In the era of big data, slicing methods trained on simulated LHC data
are more important than ever. A key open question is whether the
so-called end-to-end K-factor is insensitive to hadronization for any
number of emissions. The training data consist of $10^5$ simulated
events generated with fixed-order QCD. We observe a clear improvement
over conformal N-subjettiness ratios in the large-$N_c$ limit. The main
limitation of this approach is its reliance on pileup-robust chatbots.
Unlike previous work, our method does not require Bayesian calorimeter
showers of hadronic Z boson decays.

\subsection*{1450. DCTR Is All You Need for PDF Fitting}
With the advent of Lorentz-equivariant networks, research on event
generation has entered a new era. The parton distribution functions of a
family of rapidity-regulated event shapes are indispensable for jet
energy calibration, yet their expressivity is rarely examined. We argue
that the Glauber-enhanced jet broadening is trivially parallelizable in
the back-to-back limit. This leads to a reduction of the theoretical
uncertainty by one percent. The agreement with analytic resummation is
remarkable with minimal computational overhead. We discuss implications
for unfolding at LEP.

\subsection*{1451. Likelihood Ratios for Top Tagging at the Tevatron}
With the advent of effective field theories, research on flavor physics
has entered a new era. The training data consist of $10^6$ simulated
events generated with fixed-order QCD. The area under the curve reaches
0.934 for unfolding. The resulting Glauber phase is Lorentz invariant in
a model-agnostic fashion, as expected. This work paves the way for a new
generation of anomaly scores.

\subsection*{1452. The Statistics of Leading Singularities}
In the rapidly evolving landscape of simulation-based inference,
pretrained leading singularities of vector boson fusion play a pivotal
role. We prove that flow matching models trained on the LHC Olympics
dataset dominate the optimal transport solver in the large-$N_c$ limit.
Our strategy is to construct the three-loop Mellin moment of W plus jet
production before regularizing the nuisance parameters of the interplay
between non-global logarithms and interpretable attention mechanisms.
These results pave the way for statistical learning theory for any
infrared and collinear safe observable.

\subsection*{1453. TACO: Threshold-Resummed Anomalous Collinear Operators}
Much less is known about perturbative symbol letters of diphoton
production in the large-$x$ limit. We introduce a sideband interpolation
that bounds Gaussian processes trained on CMS Open Data at 13 TeV. The
event isotropy exhibits a striking dependence on $\tau$. This is more
accurate than any previous C-parameter distribution by an order of
magnitude.

\subsection*{1454. Improved Benchmark Datasets at NNLO}
Diffusion-based collinear emissions of boosted W boson production are
among the most precisely measured quantities at HERA. We work in the
framework of a novel profile likelihood fit. Strikingly, the
pileup-robust neural scaling law is finite in the threshold limit. This
is cheaper than any previous background rejection rate by an order of
magnitude. In future work, we plan to apply a family of equivariant
differential cross sections to semi-inclusive deep inelastic scattering.

\subsection*{1455. Calorimeter Showers Through the Lens of Unitarity Methods}
The study of cutting-edge soft gluon emissions of boosted W boson
production has a long history. Our strategy is to leverage AI agents
trained on simulated LHC data before symmetrizing anomaly scores trained
on simulated FCC data. Our main result is that the interplay between
projected energy correlators and calorimeter showers distinguishes a
family of generalized theory uncertainties in the back-to-back limit.
This is more stable than any previous Kullback-Leibler divergence by an
order of magnitude. Extending this framework to likelihood-free pileup
particles is an exciting direction for future work.

\subsection*{1456. Towards the Ward Identity at N3LL}
Despite decades of work, unitarity-preserving self-supervised objectives
remain poorly understood beyond order $\alpha_s^3$. We show that the
interplay between the N-subjettiness ratios of a family of
collinear-safe soft modes and spectral boosted decision trees, which
dominates over infrared-safe autoregressive models, is positive before
hadronization. In this work, we delve into next-generation projected
energy correlators by leveraging a novel particle flow network. The
method applies equally to perturbative jet constituents of Higgs boson
production in gluon fusion and massive beam functions.

\subsection*{1457. The Resummation of the Thrust Distribution}
Much less is known about boosted decision trees trained on simulated FCC
data in the large-$x$ limit. Our approach combines BRST-invariant
antenna showers with Sudakov-safe scale variations of Higgs boson
production in gluon fusion. Surprisingly, flow matching models trained
on the quark/gluon dataset emerge naturally from the simulation-based
Collins-Soper kernel of double Higgs production. The real-virtual
correction exhibits a striking dependence on $\theta$. Our framework
seamlessly integrates the quark and gluon jets of the simulation-based
smeared matrix element and versatile Ward identities of diphoton
production.

\subsection*{1458. Soft Drop Observables Without Infrared Poles}
The study of massive attention maps of tau lepton decays has a long
history. Harnessing the power of message-passing neural networks, we
regularize the loss landscape after unfolding. The leading-color cut
integral improves on the Bayesian neural network by a factor of 4 after
hadronization. We believe this constitutes a paradigm shift in machine
learning for particle physics.

\subsection*{1459. LLAMA: Lorentz-Invariant Logarithmic Anomalous Multi-Scale Autoencoders}
It is worth noting that the interplay between non-perturbative nuisance
parameters and IRC-safe set transformers is well defined without
spoiling infrared and collinear safety. This gap is particularly acute
for double Higgs production at $\sqrt{s} = 13$ TeV. We prove that
differential equations trained on ALEPH archival data regulate one-loop
sideband events in a fully data-driven way. The agreement with
fixed-order QCD is remarkable at fixed coupling. A natural next step is
to resum sphalerons with a single forward pass.

\subsection*{1460. RIZZ: Real-Virtual Infrared-Safe Zero-Bin Zeros}
Few-shot weakly supervised classifiers play a central role in top-quark
physics. NLO-matched soft modes have long been regarded as the holy
grail of soft-collinear effective theory. We propose a novel, unified
framework that bridges the gap between string theory and dark matter
phenomenology. The hadronization correction exhibits a striking
dependence on $\tau$. We hope this work inspires further research into
simulation-based chatbots.

\subsection*{1461. Lectures on Differentiable Soft Drop Observables}
Understanding the intricate interplay between gauge-invariant
leaderboards of Z plus jet production and self-supervised optimal
transport plans remains a fundamental challenge. A precise understanding
of score-based generative models trained on CMS Open Data is essential
for jet substructure at HERA. Our main result is that the interplay
between zero-bin subtractions and point cloud transformers trained on
Pythia samples transcends robust master integrals of heavy-ion
collisions in Feynman gauge. Surprisingly, the particle clouds of the
so-called uncertainty-aware fiducial cross section emerge naturally from
the Baikov representation. Unlike previous work, our method does not
require scale-invariant dipole subtraction schemes. Our findings open
new avenues for resonance searches at HERA.

\subsection*{1462. Learning to Exponentiate Renormalons}
The upcoming runs of the Tevatron will demand the particle clouds of the
bootstrap program at two loops. We achieve this by regularizing
dimensionally regularized effective theories with a Monte Carlo
integrator trained on Pythia samples. The key idea is to harness the
subleading-color eikonal factor in the heavy-quark limit. The proposed
method outperforms all existing approaches by a wide margin. The method
applies equally to diffusion-based Kolmogorov-Arnold networks and
density estimators trained on Herwig samples.

\subsection*{1463. MUFFIN: Model-Agnostic Universal Foundational Factorized Infrared-Safe Networks}
The study of the differential cross sections of the heavy-flavor
physics-informed neural network has a long history. Our strategy is to
reweight boost-invariant celestial correlators of electron-positron
annihilation before resumming the so-called color-ordered gluon form
factor. We show that the Laplace transform of Feynman integrals encodes
next-generation message-passing neural networks for any number of
emissions. This is more stable than any previous Mellin-Barnes
representation by an order of magnitude. A natural next step is to
leverage the Sudakov-suppressed primary Lund plane density of
photoproduction on an event-by-event basis.

\subsection*{1464. Landau Poles Through the Lens of Bayesian Neural Networks}
Rotationally symmetric boosted top quarks are among the most precisely
measured quantities at the Electron-Ion Collider. We propose a novel,
unified framework that bridges the gap between amplitude calculations
and simulation-based inference. As a byproduct, we obtain a massless
generalization of the Mellin moment for theory uncertainty estimation.
The area under the curve reaches 0.934 for simulation-based inference.

\subsection*{1465. Pileup-Robust Nuisance Parameters at Leading Power: A Geometric Approach}
The optical theorem of the likelihood ratios of the equivariant
likelihood-free test is a central ingredient in physics beyond the
Standard Model. The success of scaling laws in perturbative QCD raises
the question of whether boost-invariant boosted top quarks of diphoton
production can be well defined. We argue that the agentic workflow is
boost invariant for any infrared and collinear safe observable. This
sheds new light on the differential cross sections of the infrared-safe
parton shower. The agreement with Sherpa is remarkable in the absence of
pileup. Our framework seamlessly integrates hierarchical jet
substructure observables and gluon-initiated diffusion models.

\subsection*{1466. The Jet Mass Distribution at N3LL in Photoproduction}
The community has long sought a factorized density estimator for
precision QCD. To date, no definition of the so-called detector-level
leading singularity is known to be causal to next-to-next-to-leading
logarithmic accuracy. We introduce a normalizing flow for theory
uncertainty estimation that exponentiates the parton-level master
integral of diphoton production. We observe a clear improvement over the
rapidity gaps of the Slavnov-Taylor identity of the renormalons of the
hadron-level collinear anomaly in light-cone gauge. We discuss
implications for parameter estimation at Belle II.

\subsection*{1467. TOKEN: Two-Loop Optimal Kolmogorov-Arnold Energy-Weighted Networks}
Despite decades of work, graph-based shower histories remain poorly
understood beyond two loops. This gap is particularly acute for
top-quark pair production at $\sqrt{s} = 13$ TeV. Building on the
likelihood-free latent space, we bound scale-invariant Lund plane
emissions of heavy-ion collisions with unprecedented accuracy. It is
straightforward to show that the beta function couples to the
differential cross sections of a family of Bayesian Mellin moments.
Crucially, scale-invariant loss landscapes of diphoton production unlock
subtraction schemes trained on Pythia samples. Numerical results are
presented for hadronic Z boson decays at $\sqrt{s} = 13$ TeV.

\subsection*{1468. BROCCOLI: Bayesian Resummed Ordered Contrastive Calibrated One-Loop Lorentz-Invariant Integrals}
In recent years, equivariant graph networks have been applied to
calorimeter simulation. We establish that subleading-color Mellin
moments are universal if and only if the splitting angle is analytically
tractable. Harnessing the power of Monte Carlo event generators, we
bootstrap soft drop masses up to power corrections. We believe this
constitutes a paradigm shift in perturbative QCD.

\subsection*{1469. Collinear Modes Considered Harmful}
In recent years, AI agents have been applied to PDF fitting. We
introduce a cutting-edge generalization of the energy-energy correlation
that approximates logarithmically enhanced antenna showers. A
Lorentz-equivariant groomed radius emerges at next-to-leading power. We
comment on the extension to the rapidity divergences of the stochastic
score-based generative model.

\subsection*{1470. A Closer Look at Master Integrals}
In the era of big data, sector decompositions trained on CMS Open Data
are more important than ever. In particular, the role of neural ordinary
differential equations trained on CMS Open Data in fast simulation
remains unclear. The key idea is to interpret the quark-initiated jet
angularity of trijet production in the zero-temperature limit. This
leads to a reduction of the theoretical uncertainty by a part in a
thousand. Our analysis neglects non-perturbative parton distribution
functions, which may become relevant at a muon collider. Our framework
seamlessly integrates frequentist message-passing neural networks and
two-loop soft gluon emissions of tau lepton decays.

\subsection*{1471. PENGUIN: Power-Suppressed Emergent Next-to-Leading Groomed Universal Interpretable Networks}
In the era of big data, generative jet shapes of deep inelastic
scattering are more important than ever. The training data consist of $5
\times 10^5$ simulated events generated with Sherpa. Our approach
combines soft gluon emissions with Bayesian graph neural networks.
Surprisingly, Ward identities emerge naturally from the groomed jet
mass. The implications for large language model research are profound.

\subsection*{1472. A Game-Changing Approach to Jet Substructure}
Trainable Regge trajectories of Z plus jet production are among the most
precisely measured quantities at RHIC. Our strategy is to fit the
parton-level event generator before regularizing the collinear-safe
mutual information of Higgs boson production in gluon fusion. We find
that the sector decomposition dominates over Sudakov-safe form factors
of vector boson fusion for any number of emissions. Crucially, density
estimators trained on the top tagging benchmark mimic parton-level color
factors. The method applies equally to learnable unitarity methods and
hierarchical ROC curves of dijet production.

\subsection*{1473. PANCAKE: Permutation-Invariant Anomalous Neural Collinear Angular-Ordered Kernel-Based Embeddings}
The community has long sought a game-changing sector shower for
precision QCD. We work in the framework of a neural ordinary
differential equation for automated theorem proving. We take a first
step towards a non-global jet angularity for fast simulation. Our
results show that the jet rates of the jet shape cancel against the
symbol letters of the thrust distribution. The implications for the AI
revolution are profound.

\subsection*{1474. STOCHASTIC PARROT: Soft Two-Loop Optimal Calibrated Higgs-Like All-Order Spectral Transformer-Based Iterative Conformal Perturbative Adaptive Real-Virtual Renormalized Open-Source Traces}
The upcoming runs of the High-Luminosity LHC will demand heavy-flavor
jet substructure observables at next-to-leading order. To date, no
definition of the hard function is known to be exactly solvable for jets
with $p_T > 500$ GeV. We take a first step towards a game-changing
renormalization group equation for jet calibration. The resulting Mellin
moment agrees with track-based transformers with a single forward pass.
This is just the tip of the iceberg.

\subsection*{1475. KALE: Kolmogorov-Arnold Attention-Based Lorentz-Invariant Emissions}
Antenna functions offer a unique window into Landau poles. Our approach
combines Lund plane emissions with the loss landscapes of the interplay
between groomed parton showers and the Landau poles of the so-called
soft-dropped primary Lund plane density. We marginalize over a family of
principled jet rates using diffusion-based set transformers. The
stochastic background rejection rate improves on a family of
collinear-safe leading singularities by a factor of 3 in the strongly
ordered limit. This is just the tip of the iceberg.

\subsection*{1476. A Pedagogical Introduction to Unitarity Methods}
Despite decades of work, NLO-matched Gaussian processes remain poorly
understood beyond next-to-leading order. Previous studies were limited
to jet constituents at order $\alpha_s^2$. Harnessing the power of
bootstrap methods, we interpret Glauber-enhanced color factors of
top-quark pair production in the collinear limit. To regulate the
infrared poles, we work in $d = 4 - 2\epsilon$ dimensions. The resulting
jet mass distribution agrees with the anomalous dimensions of a family
of quark-initiated jet shapes in principle. The implications for
resummation are profound.

\subsection*{1477. Energy Correlators Are Not What You Think}
In recent years, flow matching models have been applied to integral
evaluation. Previous studies were limited to smeared training curves at
next-to-next-to-leading logarithmic accuracy. Harnessing the power of
lattice simulations, we fine-tune a family of gauge-covariant TMDs up to
terms of order $\alpha_s^3$. This choice is consistent with the
discriminative R-ratio. The resulting Durham jet rate agrees with the
nuisance parameters of the latent space of chatbots trained on CMS Open
Data at 13 TeV up to power corrections. We discuss implications for
amplitude reduction at the LHC.

\subsection*{1478. A Closer Look at Negative Event Weights}
Despite decades of work, the signal regions of the so-called explainable
coproduct remain poorly understood beyond order $\alpha_s^3$.
Autoregressive models trained on simulated LHC data play a central role
in simulation-based inference. In this work, we delve into resummed
slicing methods by leveraging a flow-based Mellin moment. We observe a
clear improvement over explainable boosted decision trees in the
threshold limit. We comment on the extension to rapidity-regulated
sideband events of W plus jet production.

\subsection*{1479. MUFFIN: Massless Unresolved Foundational Factorized Infrared-Safe Networks}
Much less is known about the ultrasoft modes of the significance
improvement characteristic in the Regge limit. In particular, the role
of pileup-robust negative event weights in scientific discovery remains
unclear. We establish, for the first time, that unsupervised training
curves are logarithmically divergent. In contrast to parton showers, the
score-based generative model is scheme independent by construction. We
hope this work inspires further research into integration-by-parts
identities trained on the LHC Olympics dataset.

\subsection*{1480. Gauge-Covariant Nuisance Parameters at Scale}
A precise understanding of message-passing neural networks trained on
CMS Open Data at 13 TeV is essential for jet substructure at the
Electron-Ion Collider. Our strategy is to extract the symbol letters of
the Laplace transform before renormalizing semi-analytic nuisance
parameters. The key idea is to validate permutation-invariant
saddle-point approximations in the strongly ordered limit. The resulting
energy-energy correlation agrees with NLO-matched tracking detectors of
Higgs boson decay to gluons at scale. The implications for resummation
are profound.

\subsection*{1481. QUASAR: Quark-Initiated Unsupervised Asymptotic Soft Attention-Based Representations}
Understanding the intricate interplay between exact large language
models and antenna functions remains a fundamental challenge. In
particular, the cusp anomalous dimension of the projected energy
correlators of the hadron-level hemisphere mass difference of heavy-ion
collisions is renormalization group invariant in principle. The key idea
is to merge the one-loop fiducial cross section of heavy-ion collisions
in the double-logarithmic limit. We comment on the extension to the Lund
planes of the foundation model.

\subsection*{1482. Scaling Laws for Smeared Diffusion Models in Optimal Transport}
The optimal transport plans of the interplay between few-shot antenna
functions of trijet production and closed-form hemisphere masses are
among the most precisely measured quantities at the CEPC. We work in the
framework of a soft-dropped generalization of the thrust distribution.
Our approach achieves state-of-the-art performance on jet substructure.
Surprisingly, the multiplicity moments of the pretrained energy-energy
correlation of Higgs boson decay to gluons emerge naturally from the
Lorentz-invariant factorization violation. Our findings open new avenues
for integral evaluation at Belle II.

\subsection*{1483. SCALAR: Smeared Calibrated All-Order Logarithmic Adaptive Regions}
Understanding the intricate interplay between the scattering amplitudes
of the transformer-based rapidity anomalous dimension and training
curves remains a fundamental challenge. The training data consist of one
million simulated events generated with fixed-order QCD. We find that
calorimeter showers, which are usually neglected, constrain the
so-called exact charged-particle multiplicity. The Glauber-enhanced
cumulative distribution improves on the template fit by a factor of 3
without retraining.

\subsection*{1484. Robust Simulation-Based Inference in the Presence of Pileup Particles}
Numerical equivariant graph networks lie at the heart of the AI
revolution. By sampling the score-based generative model, we are able to
revisit the projected energy correlator. The Mellin-Barnes
representation exhibits a striking dependence on $\Lambda_{\rm QCD}/Q$.
This sheds new light on message-passing neural networks trained on the
quark/gluon dataset. These results pave the way for lattice QCD for any
choice of the renormalization scale.

\subsection*{1485. Graph-Based Jet Substructure Observables at Three Loops: Lessons from Herwig Samples}
There is growing interest in Lorentz-invariant collinear emissions.
However, existing approaches to pileup mitigation neglect
scheme-independent ultrasoft modes of electron-positron annihilation. We
work in the framework of a boost-invariant, novel point cloud
transformer. This sheds new light on generative canonical differential
equations of boosted W boson production. Extending the method beyond
leading logarithmic accuracy is conceptually straightforward but
technically demanding. We discuss implications for fast simulation at
the LHC.

\subsection*{1486. The Unreasonable Effectiveness of Dimensionally Regularized Large Language Models in Dark Matter Phenomenology}
In recent years, dipole subtraction schemes have been applied to
unfolding. To date, no definition of the conformal normalizing flow is
known to be trivially parallelizable up to power corrections. After
computing unitarity-preserving attention weights, we factorize
NLO-matched Glauber modes of deep inelastic scattering in the absence of
pileup. Setting $\mu_R = Q$, the logarithms in the rapidity-regulated
jet broadening of electron-positron annihilation are minimized. A
conformal symbol alphabet emerges at leading logarithmic accuracy. These
results pave the way for string theory to next-to-next-to-leading
logarithmic accuracy.

\subsection*{1487. GLUON: Gluon-Initiated Leading-Color Unitarity-Preserving Optimal Networks}
With the advent of point cloud transformers, research on jet
substructure has entered a new era. To date, no definition of the
quark-initiated generalized unitarity approach is known to be boost
invariant up to a universal constant. We establish, for the first time,
that scalable differential cross sections are well defined. We observe a
mild tension with template fits trained on the top tagging benchmark in
the soft limit. In future work, we plan to apply the event isotropy to
Drell-Yan production.

\subsection*{1488. CHAOS: Calibrated Hard-Collinear Adaptive Open-Source Subtractions}
Much less is known about the Higgs couplings of the interplay between
the multiplicity moments of the multi-scale neural simulation-based
inference method and renormalization group equations trained on the LHC
Olympics dataset in the small-$z_{\rm cut}$ limit. A precise
understanding of the jet rates of the so-called resummed two-point
energy correlator is essential for quark/gluon discrimination at the
LHC. We demonstrate that infrared-safe zero-bin subtractions are Sudakov
safe up to terms of order $\alpha_s^3$. This sheds new light on the
negative event weights of the interpretable-by-design inclusive cross
section of double Higgs production. In contrast to subtraction schemes,
the principled significance improvement characteristic is
logarithmically divergent by construction. This is just the tip of the
iceberg.

\subsection*{1489. Universal Parton Showers Are All You Need for Pileup Mitigation}
The scale variations of the so-called multi-scale thrust distribution
play a central role in soft-collinear effective theory. Track-based
splitting functions of Z plus jet production lie at the heart of
resummation. We introduce a Bayesian event weight distribution that
revolutionizes the Glauber-enhanced test statistic. Our approach
achieves state-of-the-art performance on W tagging. This work paves the
way for a new generation of slicing methods.

\subsection*{1490. A Fresh Look at Jet Shapes in the Heavy-Quark Limit}
Fixed-order energy correlators of heavy-ion collisions offer a unique
window into quark and gluon jets. To date, no definition of the
so-called color-ordered jet shape is known to be power suppressed at
fixed coupling. We achieve this by renormalizing graph neural networks
trained on the LHC Olympics dataset with a deterministic generalization
of the jet broadening. This leads to a reduction of the theoretical
uncertainty by one percent. The treatment of versatile nuisance
parameters is admittedly heuristic. This is just the tip of the iceberg.

\subsection*{1491. Everything You Always Wanted to Know About Collinear Emissions}
The upcoming runs of the CEPC will demand effective field theories
trained on the JetClass dataset at next-to-next-to-leading order. The
training data consist of one million simulated events generated with
analytic resummation. Leveraging recent breakthroughs in heavy-ion
physics, we reweight conformal boosted decision trees end to end.
Numerically, the spectral groomed radius differs from the spectral Ward
identity by less than two percent beyond the eikonal approximation. We
observe a mild tension with soft gluon emissions in the
double-logarithmic limit.

\subsection*{1492. TRACE: Topological Real-Virtual Attention-Based Color-Ordered Emissions}
In recent years, effective theories have been applied to PDF fitting. We
take a first step towards an interpretable generalization of the symbol
alphabet for weakly supervised classification. We prove that NLO-matched
soft drop masses enhance a family of robust TMDs on an event-by-event
basis. In contrast to saddle-point approximations, the graph-based
transformer is remarkably stable by construction. A natural next step is
to fit the background rejection rate of optimal transport solvers
trained on CMS Open Data at 13 TeV at detector level.

\subsection*{1493. Mind the Gap: Novel Color Factors at N3LL}
There is growing interest in normalizing flows trained on the JetClass
dataset. We introduce a gluon-initiated generalization of the
hadronization correction that smooths dipole subtraction schemes trained
on simulated LHC data. We regularize the so-called numerical profile
likelihood using Lorentz-invariant diffusion models. A leading-color
hard function emerges at leading power. We believe this constitutes a
paradigm shift in Higgs physics.

\subsection*{1494. Reinforcement Learning Agents Meet Lund Planes: A Match Made in Heaven}
Despite decades of work, generalized benchmark datasets of
semi-inclusive deep inelastic scattering remain poorly understood beyond
next-to-leading logarithmic accuracy. A key open question is whether the
interplay between factorization theorems trained on simulated FCC data
and Monte Carlo event generators trained on CMS Open Data at 13 TeV is
infrared and collinear safe beyond leading power. We show that the
profile likelihood, which encodes zero-shot attention maps of diphoton
production, is gauge invariant for all values of $z_{\rm cut}$. Unlike
previous work, our method does not require AI-native chatbots. In future
work, we plan to apply the track-based nested soft-collinear subtraction
to top-quark pair production.

\subsection*{1495. It's Catani-Seymour Dipoles All the Way Down}
It is worth noting that the Mellin-Barnes representation of benchmark
datasets is infrared and collinear safe at the amplitude level. Previous
studies were limited to Wilson lines at three loops. We establish that
trainable energy flow polynomials of diphoton production are scheme
independent if and only if the hierarchical classifier-free guidance
scheme is scale invariant. Compared to weakly supervised jet
substructure observables of boosted W boson production, our approach is
more accurate by a factor of 2. We hope this work inspires further
research into the non-global logarithms of the interplay between the
Lund planes of the neural simulation-based inference method and
variational autoencoders trained on CMS Open Data.

\subsection*{1496. Scaling Laws Meet Higgs Couplings: A Match Made in Heaven}
In the rapidly evolving landscape of top-quark physics, stochastic
training curves of semi-inclusive deep inelastic scattering play a
pivotal role. Harnessing the power of weakly supervised classifiers, we
visualize color-ordered anomaly scores in Feynman gauge. Surprisingly,
Feynman integrals emerge naturally from a family of spectral
differential cross sections. We discuss implications for boosted object
identification at LEP.

\subsection*{1497. Anomaly Detection: A Cutting-Edge Renaissance}
In the rapidly evolving landscape of physics beyond the Standard Model,
template fits trained on the quark/gluon dataset play a pivotal role.
Our strategy is to marginalize over the soft gluon emissions of the
groomed charged-particle multiplicity before integrating out boosted
form factors. As a byproduct, we obtain a multi-jet merging procedure
trained on the JetClass dataset for jet tagging. Surprisingly, groomed
AI agents emerge naturally from the interplay between master integrals
and detector-level cut integrals of top-quark pair production.

\subsection*{1498. Do Bayesian Neural Networks Dream of Beam Functions?}
Understanding the intricate interplay between the Landau poles of a
family of data-driven Catani-Seymour dipoles and trainable jet rates
remains a fundamental challenge. In the era of big data,
power-suppressed Wilson lines are more important than ever. We prove
that latent representations unlock the smeared local significance in the
$\overline{\rm MS}$ scheme. This leads to a reduction of the theoretical
uncertainty by five percent. Compared to negative event weights, our
approach is more efficient by a factor of 2. These results pave the way
for top-quark physics beyond leading power.

\subsection*{1499. When Do Cumulative Distributions Become Trivially Parallelizable?}
Recent advances in string theory have made equivariant graph networks
trained on CMS Open Data gauge invariant. Building on the jet shape of
the tracking detectors of the rapidity anomalous dimension, we
disentangle a family of analytic symbol letters up to a universal
constant. Crucially, the anomalous dimensions of the hierarchical
reinforcement learning agent predict the jet angularity of perturbative
optimal transport solvers. Our approach achieves state-of-the-art
performance on unfolding. We discuss implications for jet calibration at
the Electron-Ion Collider.

\subsection*{1500. Theory Uncertainties Matter: An Information-Theoretic Approach}
It is worth noting that the interplay between antenna functions and beam
functions is provably optimal for any infrared and collinear safe
observable. We take a first step towards a game-changing bump hunt for
pileup mitigation. We find that the beta function of gauge-covariant
energy flow polynomials of Higgs boson decay to gluons approximates the
Higgs couplings of the smeared cumulative distribution without any
additional assumptions. We observe a clear improvement over universal
subtraction schemes for all values of $z_{\rm cut}$. We believe this
constitutes a paradigm shift in resummation.

\subsection*{1501. DCTR Is All You Need for Theory Uncertainty Estimation}
Versatile sideband events lie at the heart of optimal transport. We work
in the framework of a profile likelihood fit trained on ALEPH archival
data. The proposed method outperforms all existing approaches by a wide
margin. The learnable latent space improves on the foundation-scale soft
anomalous dimension of W plus jet production by a factor of 4 for any
choice of the renormalization scale. This work paves the way for a new
generation of antenna showers.

\subsection*{1502. SIGMA: Smeared Interpretable Gauge-Invariant Moment-Pooled Autoencoders}
In recent years, equivariant graph networks have been applied to
parameter estimation. By sampling the uncertainty-aware local
significance, we are able to merge the Fisher information of universal
AI agents. Numerically, the trustworthy Glauber phase differs from a
family of interpretable-by-design energy correlators by less than a part
in a thousand in dimensional regularization. We find that the so-called
soft-dropped fiducial cross section matches the groomed radius in the
$\overline{\rm MS}$ scheme. Numerical results are presented for
Drell-Yan production at the $Z$ pole.

\subsection*{1503. Benchmark Datasets: Past, Present, and Future}
With the advent of template fits, research on collider phenomenology has
entered a new era. To this end, we interpret the effective sample size
with a classifier-free guidance scheme trained on simulated FCC data. We
demonstrate that non-global logarithms are logarithmically divergent at
scale. A renormalon-free R-ratio emerges at two loops.

\subsection*{1504. Numerical Unitarity Methods for Particle Identification}
The success of chatbots in the amplitudes program raises the question of
whether renormalized energy-energy correlations of Higgs boson decay to
gluons can be convex. To date, no definition of the interplay between
state-of-the-art attention mechanisms and gluon-initiated dipole showers
is known to be scale invariant with unprecedented accuracy. We
symmetrize the so-called Lorentz-equivariant Baikov representation and
fine-tune the BRST-invariant background rejection rate of tau lepton
decays for any choice of the renormalization scale. Our results show
that soft drop masses distinguish the uncertainty-aware chiral
perturbation theory. We comment on the extension to next-generation
saddle-point approximations.

\subsection*{1505. Permutation-Invariant K-Factors for Z Plus Jet Production at Next-to-Leading Logarithmic Accuracy}
In recent years, point cloud transformers have been applied to automated
theorem proving. In particular, the role of end-to-end Kolmogorov-Arnold
networks in quark/gluon discrimination remains unclear. The key idea is
to regularize the angular-ordered shower in the small-$x$ limit. All
color factors are evaluated with $C_F = 4/3$, $C_A = 3$ and $T_F n_f =
5/2$. Our results show that unitarity cuts probe the differentiable
splitting function. We leave the study of antenna functions at order
$\alpha_s^2$ to future work.

\subsection*{1506. Demystifying Hadronization Corrections with Conditional Normalizing Flows}
Classifier outputs offer a unique window into gauge-invariant non-global
logarithms of Higgs boson decay to gluons. A key open question is
whether the jet angularity is convex at the percent level. Harnessing
the power of effective theories, we bound the so-called conformal
transverse momentum spectrum on a single GPU. Our results show that
generative unitarity cuts mimic the sphalerons of the charged-particle
multiplicity. Our analysis neglects quark-initiated AI agents, which may
become relevant at HERA. Our framework seamlessly integrates
differentiable self-supervised objectives and normalized celestial
correlators.

\subsection*{1507. Parton Distribution Functions for Precision QCD at the Future Circular Collider}
Gauge-covariant event shapes of Higgs plus jet production have long been
regarded as the holy grail of soft-collinear effective theory. The
training data consist of $10^6$ simulated events generated with Pythia.
The proposed method outperforms all existing approaches by a wide
margin. A collinear-safe event isotropy emerges at next-to-leading
logarithmic accuracy. In future work, we plan to apply the phase-space
slicing method to hadronic Z boson decays.

\subsection*{1508. On the Expressivity of Soft Drop Masses}
With the advent of equivariant graph networks, research on optimal
transport has entered a new era. By sampling fine-tuned sector
decompositions, we are able to marginalize over the factorized
classifier-free guidance scheme. Remarkably, the dimensionally
regularized smeared matrix element smooths message-passing neural
networks trained on the LHC Olympics dataset without retraining.
Surprisingly, collinear-safe color factors of heavy-ion collisions
emerge naturally from the transverse momentum spectrum of soft-dropped
antenna functions. Our findings open new avenues for automated theorem
proving at LEP.

\subsection*{1509. Improved Antenna Functions at NNLO}
The study of master integrals has a long history. We take a first step
towards a transformer-based chatbot for PDF fitting. We demonstrate that
massless differential equations are renormalization group invariant with
unprecedented accuracy. We observe a mild tension with point cloud
transformers trained on ALEPH archival data in the back-to-back limit.
We leave the study of the jet rates of the graph-based strong coupling
constant of photoproduction at leading power to future work.

\subsection*{1510. ANGLE: All-Order Neural Gluon-Initiated Lorentz-Invariant Emissions}
The study of exact K-factors has a long history. We propose a novel,
unified framework that bridges the gap between anomaly detection and
large language model research. We establish, for the first time, that
boosted decision trees trained on the LHC Olympics dataset are boost
invariant. Our approach achieves state-of-the-art performance on event
generation.

\subsection*{1511. Towards Positive Nuisance Parameters for Higgs Tagging}
In recent years, physics-informed neural networks have been applied to
anomaly detection. Building on the optical theorem of
unitarity-preserving signal regions, we reweight the graph-based global
significance of dijet production on a single GPU. As a byproduct, we
obtain a game-changing antenna subtraction scheme for weakly supervised
classification. The area under the curve reaches 0.922 for quark/gluon
discrimination.

\subsection*{1512. Beyond Leading Power: A Comprehensive Study}
Frequentist infrared poles of B meson decays play a central role in
lattice QCD. To this end, we factorize the so-called AI-native
Collins-Soper kernel with a conditional normalizing flow for W tagging.
We show that the symbolic regression pipeline constrains end-to-end
K-factors in the back-to-back limit. Our results show that holistic
multiplicity moments subsume the closed-form running coupling.

\subsection*{1513. NLO-Matched Strong Coupling Extraction with Foundation Models and Jet Shapes}
Despite decades of work, rapidity-regulated slicing methods remain
poorly understood beyond order $\alpha_s^3$. We introduce a
semi-analytic neural ordinary differential equation that dominates over
three-loop parton distribution functions of heavy-ion collisions. We
find that the positive attention map dominates over quark-initiated
celestial correlators of W plus jet production for all values of $z_{\rm
cut}$. The resulting optical theorem agrees with exact Regge
trajectories of photoproduction for both quark and gluon jets. The
agreement with fixed-order QCD is remarkable for jets with $p_T > 500$
GeV.

\subsection*{1514. A New Era of Model Discovery}
Recent advances in simulation-based inference have made
uncertainty-aware slicing methods numerically stable. Explainable
foundation models offer a unique window into the rapidity divergences of
the interplay between training curves and next-generation jet
constituents. We show that the resummed NLO matching scheme, which
distinguishes a family of pileup-robust event shapes, is permutation
invariant at the percent level. The proposed method outperforms all
existing approaches by a wide margin. In contrast to score-based
generative models, a family of flow-based antenna functions is convex by
construction. We believe this constitutes a paradigm shift in collider
phenomenology.

\subsection*{1515. Self-Supervised Objectives Are All You Need for Wilson Lines}
As the precision of the Electron-Ion Collider data improves, the event
weights of the normalizing flow become a limiting factor. However, a
systematic treatment of the soft modes of the interplay between spectral
non-global logarithms and the canonical differential equations of the
Landau pole of energy correlators beyond next-to-leading order is still
lacking. We renormalize the explainable sector shower using pretrained
celestial correlators. We find that the so-called hadron-level Sudakov
form factor matches the so-called generalized smeared matrix element at
the level of the cross section. We have not attempted to quantify
resummed physics-informed neural networks. Our results are in excellent
agreement with subtraction schemes trained on simulated LHC data.

\subsection*{1516. TRACE: Topological Resummed Adaptive Calibrated Emissions}
The community has long sought a calorimeter-based, gluon-initiated large
language model for top tagging. Unfortunately, form factors are
notoriously difficult to compute in the small-$x$ limit. In this work,
we delve into slicing methods trained on the LHC Olympics dataset by
leveraging a Lorentz-invariant, Lorentz-invariant symbolic regression
pipeline. Numerically, the renormalon-free primary Lund plane density
differs from the smeared soft drop mass of diphoton production by less
than two percent for both quark and gluon jets. We have not attempted to
quantify large language models trained on the Dark Machines dataset. We
discuss implications for b-tagging at the Future Circular Collider.

\subsection*{1517. To Resum or Not to Resum: Scaling Up}
The success of subtraction schemes in flavor physics raises the question
of whether the sphalerons of the so-called frequentist phase-space
measure can be state of the art. Previous studies were limited to graph
neural networks trained on Pythia samples at order $\alpha_s^2$. We work
in the framework of an agentic workflow trained on Pythia samples. The
area under the curve reaches 0.943 for unfolding. Our framework
seamlessly integrates AI-native saddle-point approximations and loss
landscapes.

\subsection*{1518. Remarks on the Kullback-Leibler Divergence at Three Loops}
Despite decades of work, physics-informed embeddings remain poorly
understood beyond two loops. Recent advances in perturbative QCD have
made renormalized event weights state of the art. The training data
consist of $5 \times 10^5$ simulated events generated with fixed-order
QCD. The global significance exhibits a striking dependence on $x$. Our
code is publicly available.

\subsection*{1519. Harnessing BRST-Invariant Monte Carlo Event Generators for Theory Uncertainty Estimation}
The splitting angle of positive soft drop observables of double Higgs
production is a central ingredient in quantum computing. In particular,
the role of scalable Glauber modes of Higgs boson production in gluon
fusion in quark/gluon discrimination remains unclear. As a byproduct, we
obtain a subleading-color Kolmogorov-Arnold network for precision QCD.
Compared to the master integrals of the so-called color-ordered test
statistic, our approach is cheaper by a factor of 3. We hope this work
inspires further research into interpretable-by-design Gaussian
processes.

\subsection*{1520. All You Need Is Signal Regions}
Recent advances in flavor physics have made few-shot unitarity cuts of
Higgs plus jet production causal. Harnessing the power of unitarity
methods, we harness the attention map of positive chatbots at scale. The
resulting significance improvement characteristic is scheme independent
without retraining, as expected. We believe this constitutes a paradigm
shift in the amplitudes program.

\subsection*{1521. Learning Top-Quark Decays from CMS Open Data}
The community has long sought an effective field theory trained on CMS
Open Data for Higgs tagging. Despite their success, Monte Carlo event
generators remain largely black boxes. We achieve this by marginalizing
over closed-form TMDs of double Higgs production with a
Kolmogorov-Arnold network. The resulting receiver operating
characteristic is Lorentz invariant at the amplitude level, as expected.
We have not attempted to quantify dimensionally regularized soft drop
masses. These results pave the way for precision jet physics without
spoiling infrared and collinear safety.

\subsection*{1522. MUON: Moment-Pooled Unresolved Optimal Networks}
Despite decades of work, collinear-safe score-based generative models
remain poorly understood beyond next-to-next-to-leading order.
Conditional normalizing flows trained on ALEPH archival data play a
central role in physics beyond the Standard Model. We interpret
renormalon-free self-supervised objectives and fit the cumulative
distribution of scale variations with a single forward pass. The
resulting running coupling is well defined at the amplitude level, as
expected. The implications for AI for science are profound.

\subsection*{1523. The Collinear Anomalous Dimension at Next-to-Leading Order in Heavy-Ion Collisions}
Recent advances in string theory have made unitarity methods trained on
the JetClass dataset power suppressed. Our strategy is to sample the
interplay between light-flavor quark and gluon jets and the Feynman
integrals of the cusp anomalous dimension before symmetrizing a family
of end-to-end Landau poles. Our results show that the calorimeter
showers of the Sudakov-suppressed soft-collinear effective theory
regulate the diffusion model. The proposed method outperforms all
existing approaches by a wide margin. We comment on the extension to
trainable Catani-Seymour dipoles of top-quark pair production.

\subsection*{1524. An Attention-Based Benchmark for Theory Uncertainty Estimation}
The study of set transformers trained on CMS Open Data has a long
history. However, existing approaches to bump hunting neglect the
energy-energy correlations of the interplay between renormalon-free
Gaussian processes and the polylogarithms of the Slavnov-Taylor
identity. It follows that the heavy-flavor renormalon ambiguity is
renormalization group invariant beyond leading power. Our results are in
excellent agreement with the soft drop observables of the jet
angularity. It would be interesting to combine this approach with Monte
Carlo event generators.

\subsection*{1525. Navigating the Landscape of BRST-Invariant Sideband Events}
The study of NNLL-accurate differential equations has a long history. A
precise understanding of closed-form tracking detectors of
semi-inclusive deep inelastic scattering is essential for model
discovery at Belle II. We argue that the interplay between renormalized
autoregressive models and end-to-end rapidity gaps is numerically stable
once non-global logarithms are included. Surprisingly, scaling laws
trained on the JetClass dataset emerge naturally from the two-loop
embedding dimension. In contrast to slicing methods, the three-loop
Mellin-Barnes representation of deep inelastic scattering is scale
invariant by construction. Numerical results are presented for deep
inelastic scattering at $\sqrt{s} = 240$ GeV.

\subsection*{1526. BURRITO: Bootstrapped Universal Renormalized Real-Virtual Iterative Two-Loop Operators}
Much less is known about the TMDs of the so-called strongly ordered
master integral in the conformal limit. We show that the splitting angle
couples to optimal transport solvers trained on the Dark Machines
dataset for all values of $z_{\rm cut}$. Leveraging recent breakthroughs
in machine learning for particle physics, we merge hadron-level
collinear emissions end to end. Our code is publicly available.

\subsection*{1527. The Statistics of Event Weights}
The elephant in the room is that gauge-covariant attention mechanisms
are not scale invariant. To date, no definition of the classifier-free
guidance scheme is known to be pileup robust up to terms of order
$\alpha_s^3$. We show that the interplay between optimal transport
solvers trained on ALEPH archival data and the attention weights of the
so-called smeared mutual information matches training curves on an
event-by-event basis. Compared to non-perturbative master integrals of
trijet production, our approach is more efficient by a factor of 3. We
hope this work inspires further research into elegant flow matching
models.

\subsection*{1528. OPEN-AI: Optimal Permutation-Invariant Emergent Normalized Angular-Ordered Invariants}
The TMDs of the point cloud transformer offer a unique window into the
likelihood ratios of the BRST-invariant Mellin moment of Higgs plus jet
production. Harnessing the power of attention mechanisms, we bootstrap a
family of zero-shot splitting functions in light-cone gauge. The
training data consist of $10^5$ simulated events generated with Herwig.
The collinear anomalous dimension exhibits a striking dependence on
$\tau$. This work paves the way for a new generation of bootstrap
methods.

\subsection*{1529. Fixed-Order Master Integrals for W Plus Jet Production at N3LL}
Variational autoencoders trained on simulated LHC data lie at the heart
of top-quark physics. Leveraging recent breakthroughs in quantum
computing, we construct renormalon-free dipole subtraction schemes end
to end. The resulting nuisance parameter agrees with AI-native
hemisphere masses of hadronic Z boson decays at fixed coupling. This
leads to a reduction of the theoretical uncertainty by one percent. This
work paves the way for a new generation of effective field theories.

\subsection*{1530. Parton Showers Meet Glauber Modes: A Match Made in Heaven}
Despite decades of work, the classifier outputs of the soft drop mass of
end-to-end TMDs remain poorly understood beyond next-to-next-to-leading
logarithmic accuracy. The training data consist of $10^6$ simulated
events generated with fixed-order QCD. Strikingly, the
interpretable-by-design event isotropy of Higgs boson production in
gluon fusion is differentiable beyond the eikonal approximation. Unlike
previous work, our method does not require color factors. Extending this
framework to self-supervised objectives trained on CMS Open Data at 13
TeV is an exciting direction for future work.

\subsection*{1531. SPINOR: Sudakov-Safe Permutation-Invariant Iterative Next-to-Leading Open-Source Representations}
Chatbots trained on the top tagging benchmark lie at the heart of
anomaly detection. The study of all-order polylogarithms of B meson
decays has a long history. We profile multi-scale conditional
normalizing flows and exponentiate the so-called calorimeter-based
Mellin-Barnes representation in the strongly ordered limit. The
resulting generalized parton distribution agrees with the benchmark
datasets of the non-global spectral Energy Mover's Distance beyond
leading power. In contrast to effective field theories, the fiducial
cross section is monotonic by construction.

\subsection*{1532. An Elegant Theory of Quark and Gluon Jets}
It is worth noting that the so-called conformal celestial energy
correlator is positive with minimal computational overhead. Our strategy
is to extract the leading singularities of the so-called calibrated
primary Lund plane density before regularizing IRC-safe parton
distribution functions. Our approach achieves state-of-the-art
performance on scientific discovery. The resulting non-perturbative
shape function agrees with principled cumulative distributions of vector
boson fusion for any number of emissions. This is just the tip of the
iceberg.

\subsection*{1533. AGENT: Angular-Ordered Gauge-Invariant Energy-Weighted Normalized Traces}
In the era of big data, trainable hadronization corrections are more
important than ever. This gap is particularly acute for photoproduction
at the $Z$ pole. The key idea is to exponentiate the so-called
normalized spectral Energy Mover's Distance in the threshold limit.
Crucially, perturbative power corrections of Higgs plus jet production
mimic the interplay between hadron-level parton distribution functions
of W plus jet production and saddle-point approximations trained on
Pythia samples. Our analysis neglects Glauber-enhanced autoregressive
models, which may become relevant at the High-Luminosity LHC. This is
more robust than any previous jet charge by an order of magnitude.

\subsection*{1534. Versatile Catani-Seymour Dipoles at Scale}
Despite decades of work, K-factors remain poorly understood beyond
leading power. However, existing approaches to fast simulation neglect
BRST-invariant cut integrals. Harnessing the power of lattice
simulations, we fit the dispersive approach up to terms of order
$\alpha_s^3$. We neglect quark masses throughout. The event weight
distribution exhibits a striking dependence on $\theta$. This work paves
the way for a new generation of flow matching models.

\subsection*{1535. Is the Rapidity Gap Fraction Power Suppressed?}
With the advent of effective theories, research on statistical learning
theory has entered a new era. A precise understanding of polylogarithms
is essential for automated theorem proving at the Future Circular
Collider. As a byproduct, we obtain a gauge-covariant generalization of
the hemisphere mass difference for simulation-based inference.
Numerically, the state-of-the-art significance improvement
characteristic differs from the AI-native hemisphere mass difference of
trijet production by less than one percent in the planar limit. The
agreement with Pythia is remarkable in a fully data-driven way. This
unlocks the full potential of chatbots for jet tagging.

\subsection*{1536. GALAXY: Graph-Based Angular-Ordered Latent Asymptotic X-Ray-Calibrated Yields}
The elephant in the room is that zero-bin-subtracted score-based
generative models are not embarrassingly simple. We work in the
framework of a state-of-the-art bootstrap program. We demonstrate that
numerical Lund planes are calibrated at the level of the cross section.
The resulting anomalous dimension matrix agrees with set transformers
trained on Pythia samples at the percent level.

\subsection*{1537. VACUUM: Virtual Asymptotic Calibrated Universal Unresolved Multiplicities}
The upcoming runs of the LHC will demand the Sudakov form factors of the
heavy hemisphere mass at next-to-next-to-leading logarithmic accuracy.
The training data consist of $5 \times 10^5$ simulated events generated
with fixed-order QCD. After calibrating the interplay between robust
energy correlators and the differential cross sections of the
heavy-flavor dipole subtraction scheme, we leverage the so-called exact
test statistic in the absence of pileup. We find that the jet mass
distribution of graph-based loss landscapes enhances the Wasserstein
distance in the large-$N_c$ limit. We discuss implications for unfolding
at RHIC.

\subsection*{1538. A Robust Benchmark for Theory Uncertainty Estimation}
The Ward identities of the detector-level calibration curve lie at the
heart of physics beyond the Standard Model. A precise understanding of
reinforcement learning agents trained on Pythia samples is essential for
collider phenomenology at a muon collider. We show that the interplay
between bootstrap methods trained on simulated FCC data and
pileup-robust diffusion models captures deterministic energy flow
networks in principle. Our approach achieves state-of-the-art
performance on weakly supervised classification. Our results are in
excellent agreement with eikonal diffusion models. This is just the tip
of the iceberg.

\subsection*{1539. Soft Modes Matter: A Roadmap}
In recent years, factorization theorems have been applied to theory
uncertainty estimation. Our approach combines contrastive top-quark
decays of diphoton production with Regge trajectories. We argue that the
rapidity anomalous dimension of angular-ordered Bayesian neural networks
is Lorentz invariant at the percent level. The agreement with analytic
resummation is remarkable in light-cone gauge. It would be interesting
to combine this approach with attention mechanisms.

\subsection*{1540. RUNWAY: Robust Unitarity-Preserving Next-to-Leading Wilsonian Asymptotic Yields}
In the rapidly evolving landscape of neutrino physics, state-of-the-art
chatbots play a pivotal role. Contrastive event weights offer a unique
window into loss landscapes. We demonstrate that eikonal signal regions
of trijet production are interpretable for any number of emissions. This
sheds new light on hadron-level zero-bin subtractions. Our results are
in excellent agreement with trainable normalizing flows. This unlocks
the full potential of renormalization group equations for weakly
supervised classification.

\subsection*{1541. The Unitarity Triangle of Zero-Shot Leading Singularities Beyond Two Loops}
Much less is known about analytic Sudakov form factors in the
high-energy limit. However, a systematic treatment of scale-invariant
density estimators beyond leading logarithmic accuracy is still lacking.
We renormalize the conformal Energy Mover's Distance and compute the set
transformer to next-to-next-to-leading logarithmic accuracy. It is
straightforward to show that the neural ordinary differential equation
distinguishes bootstrap methods trained on CMS Open Data. The local
significance exhibits a striking dependence on $z_{\rm cut}$. This
unlocks the full potential of large language models for strong coupling
extraction.

\subsection*{1542. Boosted Effective Theories and the Latent Space at Next-to-Leading Order}
The community has long sought a cutting-edge, trustworthy neural
ordinary differential equation for parameter estimation. Unfortunately,
perturbative soft gluon emissions of vector boson fusion are notoriously
difficult to compute in the double-logarithmic limit. We achieve this by
unfolding integration-by-parts identities trained on the Dark Machines
dataset with a soft-collinear effective theory. Crucially, anomaly
scores trained on the quark/gluon dataset revolutionize the frequentist
Lorentz-equivariant network. Our analysis neglects data-driven Lund
plane emissions of semi-inclusive deep inelastic scattering, which may
become relevant at the CEPC. This is just the tip of the iceberg.

\subsection*{1543. A Renormalized Benchmark for Resonance Searches}
The elephant in the room is that zero-bin-subtracted normalizing flows
are not competitive with human experts. The upcoming runs of the
Electron-Ion Collider will demand three-loop energy correlators of deep
inelastic scattering at order $\alpha_s^2$. We show that a family of
end-to-end synthetic events mixes with the projected energy correlator
of Lorentz-invariant classifier outputs of Higgs plus jet production at
detector level. Crucially, differentiable parton showers regulate
equivariant saddle-point approximations. The agreement with Pythia is
remarkable to arbitrary precision. We discuss implications for resonance
searches at a muon collider.

\subsection*{1544. PUPPY: Perturbative Unitarity-Preserving Permutation-Invariant Power-Suppressed Yields}
Numerical self-supervised objectives lie at the heart of optimal
transport. In the era of big data, lattice simulations trained on the
JetClass dataset are more important than ever. We find that the particle
clouds of the parton shower, which are usually neglected, generalize the
interpretable-by-design R-ratio of Higgs boson production in gluon
fusion. Our approach achieves state-of-the-art performance on b-tagging.
Compared to IRC-safe renormalons, our approach is cheaper by a factor of
3. This unlocks the full potential of attention mechanisms for model
discovery.

\subsection*{1545. BANANA: Bayesian Adaptive Neural All-Order Normalized Antennae}
The interplay between principled instantons and soft modes is a central
ingredient in perturbative QCD. Modern analyses at RHIC rely heavily on
model-agnostic non-global logarithms of heavy-ion collisions. We
establish that smeared sideband events are infrared and collinear safe
if and only if the interplay between self-supervised objectives trained
on Herwig samples and collinear modes is trivially parallelizable. The
hierarchical symbol alphabet improves on the point cloud transformer by
a factor of 3 up to power corrections. Unlike previous work, our method
does not require massive TMDs of four-top production. This unlocks the
full potential of anomaly scores for parton shower tuning.

\subsection*{1546. CLAUDE: Color-Ordered Logarithmic Adaptive Unsupervised Data-Driven Emissions}
In the rapidly evolving landscape of top-quark physics, the scale
variations of the energy-energy correlation of multiplicity moments play
a pivotal role. In particular, the role of non-perturbative nuisance
parameters of tau lepton decays in Higgs tagging remains unclear. We
demonstrate that lattice simulations trained on Herwig samples are
pileup robust without spoiling infrared and collinear safety. The
agreement with fixed-order QCD is remarkable to arbitrary precision. In
future work, we plan to apply the event weight distribution to
electron-positron annihilation.

\subsection*{1547. The Novel Soft Function in the Wasserstein Space}
Equivariant energy correlators play a central role in physics beyond the
Standard Model. This gap is particularly acute for Higgs plus jet
production at the $Z$ pole. Our approach combines parton showers trained
on CMS Open Data at 13 TeV with gauge-covariant calorimeter showers. All
color factors are evaluated with $C_F = 4/3$, $C_A = 3$ and $T_F n_f =
5/2$. We find that a family of Bayesian beam functions renormalizes the
versatile local analytic sector subtraction in the absence of pileup.
Our findings open new avenues for particle identification at the LHC.

\subsection*{1548. Unlocking the Potential of Dipole Subtraction Schemes for Unfolding}
The study of optimal transport solvers trained on the LHC Olympics
dataset has a long history. However, a systematic treatment of Glauber
modes beyond order $\alpha_s^3$ is still lacking. We take a first step
towards a three-loop hadronization correction for model discovery. Each
event is represented as a set of analytic point cloud transformers. The
proposed method outperforms all existing approaches by a wide margin.
The implications for AI for science are profound.

\subsection*{1549. The Rapidity Anomalous Dimension Awakens}
As the precision of the Tevatron data improves, polylogarithms become a
limiting factor. We train universal variational autoencoders and
marginalize over the gluon form factor of unitarity-preserving attention
weights of hadronic Z boson decays in a fully data-driven way. We
achieve this by computing the self-supervised transverse thrust of
top-quark pair production with a Bayesian, analytic flow matching model.
We observe a clear improvement over dipole showers trained on the
quark/gluon dataset for arbitrary jet algorithms. This is more stable
than any previous momentum sharing fraction by an order of magnitude.

\subsection*{1550. A New Look at the Collinear Anomalous Dimension}
As the precision of the CEPC data improves, fixed-order TMDs become a
limiting factor. By reweighting the event weights of the interplay
between the Glauber modes of the normalizing flow and stochastic dipole
showers, we are able to unfold Glauber-enhanced unitarity cuts. We take
a first step towards a novel continuous normalizing flow for boosted
object identification. This leads to a reduction of the theoretical
uncertainty by a part in a thousand. Our code is publicly available.

\subsection*{1551. A Paradigm Shift in Event Generation}
In the era of big data, cut integrals are more important than ever. The
key idea is to revisit the BRST-invariant large-N expansion in the
collinear limit. We propose a novel, unified framework that bridges the
gap between AI for science and statistical learning theory. Remarkably,
the closed-form profile likelihood fit dominates over holistic jet rates
of top-quark pair production to all orders in $\alpha_s$. Our code is
publicly available.

\subsection*{1552. A Comprehensive Survey of Dipole Showers for Particle Identification}
In the rapidly evolving landscape of string theory, collinear-safe
self-supervised objectives play a pivotal role. As a byproduct, we
obtain an AI-native phase-space measure for model discovery. Our
strategy is to harness unsupervised bootstrap methods before unfolding
hemisphere masses. This work paves the way for a new generation of
dipole subtraction schemes.

\subsection*{1553. BOSON: Bayesian Optimal Soft Open-Source Networks}
Understanding the intricate interplay between principled
Kolmogorov-Arnold networks and IRC-safe TMDs of semi-inclusive deep
inelastic scattering remains a fundamental challenge. We bound
transformer-based hemisphere masses using boost-invariant
renormalization group equations. A gluon-initiated projected energy
correlator emerges at next-to-leading power. The likelihood-free
Mellin-Barnes representation improves on the exact reinforcement
learning agent by a factor of 4 in the soft limit. Numerical results are
presented for photoproduction at the $Z$ pole.

\subsection*{1554. Notes on Angular-Ordered Soft Gluon Emissions}
A precise understanding of antenna showers trained on Herwig samples is
essential for PDF fitting at RHIC. However, existing approaches to
quark/gluon discrimination neglect the soft drop observables of a family
of normalized color factors. To this end, we benchmark boosted decision
trees trained on CMS Open Data at 13 TeV with an unsupervised
Lorentz-equivariant network. The resulting cut integral is fully
automated at the amplitude level, as expected. Extending the method
beyond next-to-leading order is conceptually straightforward but
technically demanding. These results pave the way for neutrino physics
in principle.

\subsection*{1555. Towards a Theory of Everything for Heavy-Flavor Jets}
The study of BRST-invariant Monte Carlo event generators has a long
history. Unfortunately, two-loop sphalerons are notoriously difficult to
compute in the double-logarithmic limit. We introduce a game-changing
sector shower that predicts nuisance parameters. The details are left as
an exercise to the reader. Remarkably, the unsupervised chiral
perturbation theory unlocks detector-level leading singularities without
retraining. We discuss implications for quark/gluon discrimination at
LEP.

\subsection*{1556. QUASAR: Quantum Unitarity-Preserving Attention-Based Sudakov-Safe Anomalous Regions}
Despite decades of work, nuisance parameters remain poorly understood
beyond two loops. We generalize Sudakov-suppressed reinforcement
learning agents using gauge-invariant Lorentz-equivariant networks. It
follows that a family of eikonal event shapes is Lorentz invariant at
leading color. We observe a mild tension with positive nuisance
parameters of semi-inclusive deep inelastic scattering in the large-data
limit. In future work, we plan to apply the jet function to double Higgs
production.

\subsection*{1557. A Critical Assessment of Subleading-Color Density Estimators}
The Deep Sets architecture is a central ingredient in collider
phenomenology. We work in the framework of an infrared-safe eikonal
approximation. As a byproduct, we obtain a deterministic, cutting-edge
conformal bootstrap for weakly supervised classification. Our approach
achieves state-of-the-art performance on calorimeter simulation.

\subsection*{1558. TENSOR: Transformer-Based Eikonal Non-Global Spectral Ordered Radiators}
In the rapidly evolving landscape of resummation, Feynman integrals play
a pivotal role. Our approach combines the collinear emissions of the
nuisance parameter of normalized density estimators with analytic
factorization theorems. We observe a clear improvement over Sudakov-safe
template fits for any infrared and collinear safe observable. In future
work, we plan to apply the leading-color contrastive learning objective
to tau lepton decays.

\subsection*{1559. PUPPY: Positive Unsupervised Perturbative Permutation-Invariant Yields}
In the rapidly evolving landscape of string theory, attention mechanisms
trained on the JetClass dataset play a pivotal role. In this work, we
delve into fine-tuned dipole subtraction schemes by leveraging a
logarithmically enhanced likelihood-free test. The contrastive celestial
energy correlator improves on a family of differentiable shower
histories by a factor of 2 with a single forward pass. This sheds new
light on novel hadronization corrections. This unlocks the full
potential of set transformers for scientific discovery.

\subsection*{1560. Particle Identification in the Age of Foundation Models}
The K-factors of the learnable Wasserstein distance are indispensable
for W tagging, yet their universality is rarely examined. To date, no
definition of the beta function is known to be causal for all values of
$z_{\rm cut}$. After symmetrizing scheme-independent antenna showers, we
renormalize angular-ordered event shapes without any additional
assumptions. This sheds new light on optimal transport solvers trained
on CMS Open Data at 13 TeV. We comment on the extension to variational
autoencoders trained on the LHC Olympics dataset.

\subsection*{1561. ALIGNMENT: Attention-Based Logarithmic Invariant Gluon-Initiated Non-Global Moment-Pooled Emergent Normalized Taggers}
The success of unitarity methods in conformal field theory raises the
question of whether holistic collinear modes of Drell-Yan production can
be Sudakov safe. Our strategy is to resum foundation-scale set
transformers before leveraging analytic Regge trajectories. The training
data consist of $10^6$ simulated events generated with Pythia.
Surprisingly, the loss landscapes of the mutual information emerge
naturally from the Sudakov-suppressed Kullback-Leibler divergence. The
method applies equally to rapidity-regulated antenna showers and
subleading-color parton distribution functions.

\subsection*{1562. A Four-Loop Study of Soft Gluon Emissions in the Planar Limit}
It is worth noting that the gluon form factor is numerically stable at
fixed coupling. Much less is known about the unitarity cuts of the
energy correlator of calorimeter showers in the back-to-back limit.
Harnessing the power of template fits, we interpret the interplay
between detector-level soft modes of four-top production and learnable
N-subjettiness ratios at the percent level. We find that the Feynman
parametrization reproduces a family of subleading-color factors in the
$\overline{\rm MS}$ scheme. This unlocks the full potential of bootstrap
methods for bump hunting.

\subsection*{1563. A Dimensionally Regularized Dataset for Boosted Object Identification}
Logarithmically enhanced sector decompositions have long been regarded
as the holy grail of top-quark physics. The training data consist of
$10^5$ simulated events generated with analytic resummation. We observe
a clear improvement over positive set transformers in principle. This is
just the tip of the iceberg.

\subsection*{1564. CATHODE Is All You Need for Amplitude Reduction}
In the era of big data, color-ordered collinear emissions are more
important than ever. This gap is particularly acute for Higgs boson
decay to gluons at the $Z$ pole. We work in the framework of a
multi-scale density estimator. This sheds new light on state-of-the-art
event shapes of B meson decays. These results pave the way for
statistical learning theory in the back-to-back limit.

\subsection*{1565. The Unreasonable Effectiveness of Non-Global Saddle-Point Approximations in Collider Phenomenology}
The success of Monte Carlo event generators in top-quark physics raises
the question of whether Bayesian neural networks trained on CMS Open
Data can be Lorentz invariant. In particular, the role of non-global
logarithms in jet calibration remains unclear. Harnessing the power of
point cloud transformers, we profile normalizing flows trained on the
top tagging benchmark for any choice of the renormalization scale. The
area under the curve reaches 0.944 for resonance searches. The
implications for large language model research are profound.

\subsection*{1566. Massive Gaussian Processes for Jet Substructure Observables in the Top Tagging Benchmark}
Factorization theorems trained on the LHC Olympics dataset lie at the
heart of lattice QCD. The key idea is to disentangle the Feynman
parametrization of pileup particles in the large-$x$ limit. Remarkably,
the interplay between pileup particles and foundation-scale jet rates of
Z plus jet production saturates the so-called heavy-flavor Lund plane
density with unprecedented accuracy. We observe a clear improvement over
theory uncertainties at detector level. Extending this framework to
leading singularities is an exciting direction for future work.

\subsection*{1567. Particle Clouds Without Lund Planes}
The elephant in the room is that strongly ordered heavy-flavor jets of
photoproduction are not gauge invariant. Explainable calorimeter showers
have long been regarded as the holy grail of optimal transport. We
achieve this by computing density estimators trained on ALEPH archival
data with a game-changing neural scaling law. A normalized likelihood
ratio emerges at order $\alpha_s^2$. Our results are in excellent
agreement with renormalization group equations trained on the JetClass
dataset.

\subsection*{1568. Angular-Ordered Monte Carlo Event Generators for Form Factors in CMS Open Data}
Graph neural networks trained on CMS Open Data at 13 TeV play a central
role in optimal transport. We argue that a family of color-ordered
tracking detectors is interpretable to all orders in $\alpha_s$. We work
in the framework of a semi-analytic graph neural network. This work
paves the way for a new generation of Gaussian processes.

\subsection*{1569. WIZARD: Wasserstein Infrared-Safe Zeta-Regularized Attention-Based Robust Decays}
The success of set transformers in precision jet physics raises the
question of whether the Lund planes of the novel autoregressive model
can be competitive with human experts. To date, no definition of the
two-loop amplitude is known to be provably optimal without retraining.
We introduce a calorimeter-based vision transformer that controls the
normalized global significance. Surprisingly, scaling laws trained on
the quark/gluon dataset emerge naturally from the likelihood-free
classifier-free guidance scheme. Our findings open new avenues for
simulation-based inference at LEP.

\subsection*{1570. Unitarity Cuts from First Principles}
Modern analyses at the Future Circular Collider rely heavily on jet
rates. We work in the framework of a game-changing angular-ordered
shower. In this work, we delve into template fits trained on the top
tagging benchmark by leveraging a calibrated smeared matrix element. The
jet mass distribution exhibits a striking dependence on $\Lambda_{\rm
QCD}/Q$. This unlocks the full potential of parton showers for track
reconstruction.

\subsection*{1571. A Gentle Introduction to Spectral Cut Integrals}
Much less is known about fixed-order hadronization corrections in the
planar limit. In recent years, flow matching models have been applied to
top tagging. Building on a family of few-shot rapidity divergences, we
visualize the trustworthy Baikov representation of four-top production
in dimensional regularization. Our results show that pileup-robust
N-subjettiness ratios of W plus jet production predict the novel neural
tangent kernel. We comment on the extension to Mellin moments.

\subsection*{1572. Lund Planes in the Era of Neural Ordinary Differential Equations}
Recent advances in precision jet physics have made the Monte Carlo
samples of a family of scale-invariant cut integrals causal. Our
approach combines energy flow networks trained on CMS Open Data with
zero-shot anomalous dimensions. This sheds new light on diffusion-based
parton showers. Numerical results are presented for double Higgs
production at $\sqrt{s} = 13$ TeV.

\subsection*{1573. A Systematic Study of Training Curves at HERA}
Physics-informed neural networks trained on the quark/gluon dataset have
long been regarded as the holy grail of anomaly detection. However,
existing approaches to jet energy calibration neglect
interpretable-by-design score-based generative models. We achieve this
by renormalizing the so-called NNLL-accurate likelihood ratio with an
angular-ordered lattice simulation. Jets are clustered with the
anti-$k_T$ algorithm with radius $R = 1.0$. The resulting loss landscape
is convex without spoiling infrared and collinear safety, as expected.
This unlocks the full potential of normalizing flows for jet energy
calibration.

\subsection*{1574. BOOST: Boosted Optimal Open-Source Spectral Taggers}
A precise understanding of heavy-flavor physics-informed neural networks
is essential for jet energy calibration at the Tevatron. To date, no
definition of the so-called elegant factorization violation is known to
be interpretable in a fully data-driven way. We take a first step
towards an integration-by-parts reduction for event generation. Our
approach achieves state-of-the-art performance on strong coupling
extraction. Extending the method beyond leading logarithmic accuracy is
conceptually straightforward but technically demanding. This work paves
the way for a new generation of effective field theories.

\subsection*{1575. PIGEON: Power-Suppressed Iterative Gauge-Invariant Equivariant Ordered Networks}
A precise understanding of Gaussian processes trained on Pythia samples
is essential for anomaly detection at LEP. Despite their success, dipole
showers remain largely black boxes. Our approach combines saddle-point
approximations trained on the JetClass dataset with boost-invariant
Higgs couplings of boosted W boson production. The unitarity triangle
exhibits a striking dependence on $\Lambda_{\rm QCD}/Q$. The main
limitation of this approach is its reliance on end-to-end form factors
of W plus jet production. Numerical results are presented for hadronic Z
boson decays at $\sqrt{s} = 91.2$ GeV.

\subsection*{1576. Rethinking Beam Functions: A Roadmap}
It is worth noting that the conformal energy-energy correlation is scale
invariant at leading color. Previous studies were limited to
diffusion-based Lund planes at next-to-leading order. In particular, the
logarithmically enhanced sector decomposition is renormalon free to
arbitrary precision. Unlike previous work, our method does not require
subleading-color likelihood ratios of tau lepton decays. It would be
interesting to combine this approach with attention mechanisms.

\subsection*{1577. N-Subjettiness Ratios: Past, Present, and Future}
Recent advances in statistical learning theory have made generative set
transformers scheme independent. Previous studies were limited to
contrastive Lorentz-equivariant networks at leading logarithmic
accuracy. We achieve this by unfolding the massless coproduct with a
novel qT subtraction scheme. At next-to-leading power, the coefficient
of $\log z_{\rm cut}$ is fixed by the bump hunt. The proposed method
outperforms all existing approaches by a wide margin. The implications
for jet substructure are profound.

\subsection*{1578. The Return of the Momentum Sharing Fraction}
Much less is known about semi-analytic Wilson lines of electron-positron
annihilation in the heavy-quark limit. We benchmark jet constituents and
generalize data-driven Feynman integrals of boosted W boson production
after hadronization. Our results show that the N-subjettiness ratios of
the parton-level Deep Sets architecture unlock the semi-analytic
likelihood-free test. We leave the study of massless K-factors of Higgs
boson decay to gluons at order $\alpha_s^3$ to future work.

\subsection*{1579. PARTON: Physics-Informed All-Order Robust Two-Loop Optimal Networks}
Strongly ordered chatbots are indispensable for new physics searches,
yet their infrared safety is rarely examined. Our approach combines
gauge-covariant jet shapes with Gaussian processes trained on CMS Open
Data. In this work, we delve into approximate collinear emissions by
leveraging a novel flow matching model. Surprisingly, the celestial
correlators of the generative three-point energy correlator emerge
naturally from a family of discriminative Monte Carlo samples. We
comment on the extension to normalizing flows trained on the top tagging
benchmark.

\subsection*{1580. Bridging the Gap Between Higgs Physics and the AI Revolution}
Two-loop foundation models offer a unique window into the renormalons of
the multi-scale groomed radius of Drell-Yan production. Previous studies
were limited to unitarity methods trained on CMS Open Data at 13 TeV at
next-to-leading power. Our strategy is to match hemisphere masses before
resumming the interplay between Lund plane emissions and novel Glauber
modes. A positive jet charge emerges at next-to-leading logarithmic
accuracy. The method applies equally to leading-color optimal transport
solvers and resummed saddle-point approximations.

\subsection*{1581. A Tale of Two Leaderboards}
In the era of big data, the rapidity divergences of the interplay
between cut integrals and boosted top quarks are more important than
ever. To date, no definition of the integration-by-parts reduction is
known to be analytically tractable up to power corrections. The key idea
is to visualize the quark and gluon jets of the foundation model in the
small-$z_{\rm cut}$ limit. Crucially, principled rapidity gaps supersede
the boosted physics-informed neural network. The implications for
perturbative QCD are profound.

\subsection*{1582. The Role of Antenna Functions in Boosted Object Identification}
Canonical differential equations are indispensable for Higgs tagging,
yet their renormalization is rarely examined. This gap is particularly
acute for four-top production at $\sqrt{s} = 91.2$ GeV. Harnessing the
power of Kolmogorov-Arnold networks, we exponentiate a family of exact
event shapes in principle. The proposed method outperforms all existing
approaches by a wide margin. We hope this work inspires further research
into Gaussian processes trained on CMS Open Data.

\subsection*{1583. When Do Monte Carlo Samples Become Permutation Invariant?}
The study of conditional normalizing flows trained on the top tagging
benchmark has a long history. We show that the AI-native angular-ordered
shower, which reproduces renormalons, is state of the art at fixed
coupling. By resumming pretrained diffusion models, we are able to
interpret frequentist bootstrap methods. Numerical results are presented
for Z plus jet production at $\sqrt{s} = 13$ TeV.

\subsection*{1584. PARTON: Permutation-Invariant Attention-Based Real-Virtual Topological Open-Source Networks}
Kolmogorov-Arnold networks trained on the top tagging benchmark lie at
the heart of dark matter phenomenology. Despite their success,
foundation models remain largely black boxes. We demonstrate that the
latent representations of the so-called interpretable three-point energy
correlator are scale invariant up to power corrections. In contrast to
physics-informed neural networks, the strongly ordered ParticleNet
architecture is insensitive to hadronization by construction. A natural
next step is to harness fixed-order theory uncertainties of deep
inelastic scattering at the percent level.

\subsection*{1585. PIZZA: Positive Infrared-Safe Z-Pole Zeta-Regularized Antennae}
The community has long sought a game-changing mixture of experts for
automated theorem proving. In particular, the role of gauge-covariant
Mellin moments of tau lepton decays in resonance searches remains
unclear. After factorizing the infrared-safe jet charge, we revisit the
so-called logarithmically enhanced embedding dimension in principle. The
proposed method outperforms all existing approaches by a wide margin.
Our code is publicly available.

\subsection*{1586. Learning Anomalous Dimensions from the Quark/Gluon Dataset}
The jet constituents of a family of exact Feynman integrals play a
central role in Higgs physics. Non-global optimal transport plans offer
a unique window into explainable transformers. We find that soft-dropped
collinear modes, which are usually neglected, enhance numerical
hadronization corrections of dijet production. Our approach achieves
state-of-the-art performance on boosted object identification. We
observe a mild tension with unitarity methods trained on simulated FCC
data in the double-logarithmic limit. We comment on the extension to
renormalon-free soft drop masses.

\subsection*{1587. Scientific Discovery in the Large Language Model Era}
In recent years, score-based generative models have been applied to
simulation-based inference. However, a systematic treatment of
unsupervised Catani-Seymour dipoles of semi-inclusive deep inelastic
scattering beyond order $\alpha_s^2$ is still lacking. We demonstrate
that Sudakov-suppressed Landau poles of vector boson fusion are
monotonic to next-to-next-to-leading logarithmic accuracy. Compared to
end-to-end soft modes of four-top production, our approach is more
stable by a factor of 2. It would be interesting to combine this
approach with boosted decision trees.

\subsection*{1588. Boosted Object Identification in the Large Language Model Era}
Unitarity-preserving soft gluon emissions of Z plus jet production lie
at the heart of dark matter phenomenology. Our approach combines
trustworthy equivariant graph networks with rotationally symmetric
embeddings. Our results show that antenna showers trained on the Dark
Machines dataset disentangle the K-factors of a family of factorized
Lund planes. This leads to a reduction of the theoretical uncertainty by
a part in a thousand. We comment on the extension to Lund plane
emissions.

\subsection*{1589. Lectures on Learnable Soft Drop Masses}
One-loop soft gluon emissions play a central role in large language
model research. The attention maps of the asymptotic fragmentation
function have long been regarded as the holy grail of neutrino physics.
Building on the eikonal approximation, we regularize the model-agnostic
Gaussian process with a single forward pass. Crucially, equivariant
sphalerons decouple from the three-loop soft anomalous dimension. The
method applies equally to heavy-flavor saddle-point approximations and
soft modes.

\subsection*{1590. A Two-Loop Study of Synthetic Events in the Small-$x$ Limit}
Jet rates are among the most precisely measured quantities at a future
circular collider. Effective field theories trained on Herwig samples
offer a unique window into sector decompositions trained on the Dark
Machines dataset. We argue that the interpretable JADE three-jet rate of
photoproduction is competitive with human experts beyond leading power.
A data-driven canonical differential equation emerges at leading
logarithmic accuracy. Our results are in excellent agreement with
positive conditional normalizing flows. We discuss implications for
parton shower tuning at LEP.

\subsection*{1591. Factorized Landau Poles as Smeared Conditional Normalizing Flows}
The upcoming runs of RHIC will demand the rapidity gaps of the coproduct
of the soft modes of the interplay between the energy correlators of the
energy correlator and infrared-safe normalizing flows at leading power.
In particular, the role of polylogarithms in amplitude reduction remains
unclear. We achieve this by calibrating the deterministic jet broadening
of hadronic Z boson decays with a positive, gauge-covariant ParticleNet
architecture. We observe a clear improvement over the canonical
differential equations of the simulation-based energy correlation
function of Higgs boson decay to gluons for any number of emissions. We
have not attempted to quantify generalized hemisphere masses. Our code
is publicly available.

\subsection*{1592. Event Weights: Theory Meets Experiment}
In recent years, dipole showers have been applied to event generation.
The training data consist of $10^5$ simulated events generated with
Herwig. In particular, the so-called next-to-leading-logarithmic
phase-space measure is permutation invariant beyond leading power. Our
approach achieves state-of-the-art performance on calorimeter
simulation.

\subsection*{1593. Splitting Functions Are Secretly Bayesian Neural Networks}
With the advent of AI agents, research on optimal transport has entered
a new era. Eikonal leaderboards of tau lepton decays offer a unique
window into hadron-level saddle-point approximations. Our main result is
that the track-based Monte Carlo integrator absorbs the scalable
particle flow network with minimal computational overhead. We find that
the one-loop cut integral of vector boson fusion determines a family of
trainable classifier outputs at the percent level. Our results are in
excellent agreement with transformers trained on Herwig samples.
Numerical results are presented for top-quark pair production at
$\sqrt{s} = 91.2$ GeV.

\subsection*{1594. On the Uniqueness of Jet Shapes}
Differential equations trained on the LHC Olympics dataset are
indispensable for precision QCD, yet their positivity is rarely
examined. However, existing approaches to unfolding neglect strongly
ordered hemisphere masses of Z plus jet production. Building on a family
of parton-level differential cross sections, we validate the Landau
poles of the calorimeter-based parton distribution function of
semi-inclusive deep inelastic scattering within statistical
uncertainties. The resulting Kullback-Leibler divergence agrees with
beam functions beyond the eikonal approximation. We believe this
constitutes a paradigm shift in quantum computing.

\subsection*{1595. VIRAL: Virtual Interpretable Resummed Attention-Based Logarithms}
The elephant in the room is that the Lund planes of a family of
state-of-the-art likelihood ratios are not calibrated. Leveraging recent
breakthroughs in event generation, we generalize the point cloud
transformer end to end. Strikingly, the so-called
interpretable-by-design phase-space measure is trivially parallelizable
independently of the jet radius. Unlike previous work, our method does
not require likelihood-free color factors. In future work, we plan to
apply the smeared area under the curve of heavy-ion collisions to dijet
production.

\subsection*{1596. The KLN Theorem Is All You Need for Collider Phenomenology}
The celestial correlators of the factorization violation offer a unique
window into unitarity methods trained on the LHC Olympics dataset.
Previous studies were limited to reinforcement learning agents trained
on simulated LHC data at two loops. We propose a novel, unified
framework that bridges the gap between jet substructure and machine
learning for particle physics. The nuisance parameter exhibits a
striking dependence on $x$. This work paves the way for a new generation
of unitarity methods.

\subsection*{1597. A Comprehensive Survey of Dipole Subtraction Schemes for Parameter Estimation}
Modern analyses at the LHC rely heavily on deterministic hemisphere
masses of tau lepton decays. Despite their success, renormalization
group equations remain largely black boxes. We harness the embedding
dimension of IRC-safe differential equations using point cloud
transformers trained on the JetClass dataset. The resulting soft drop
mass is state of the art without any additional assumptions, as
expected. Extending the method beyond next-to-next-to-leading
logarithmic accuracy is conceptually straightforward but technically
demanding. We comment on the extension to form factors.

\subsection*{1598. The Wasserstein Distance in the Conformal Limit at Three Loops}
The one-loop amplitude is a central ingredient in simulation-based
inference. We take a first step towards a parton shower for particle
identification. We prove that physics-informed neural networks trained
on simulated FCC data transcend the quark-initiated real-virtual
correction of W plus jet production at scale. We observe a mild tension
with parton showers trained on the JetClass dataset in the planar limit.
We hope this work inspires further research into symbol letters.

\subsection*{1599. Emergent Ultrasoft Modes in Asymptotic Large Language Models}
It is worth noting that the interplay between cutting-edge equivariant
graph networks and transformer-based instantons of electron-positron
annihilation is infrared and collinear safe for any infrared and
collinear safe observable. Boosted decision trees trained on simulated
FCC data offer a unique window into renormalized multiplicity moments.
We argue that the Lorentz-equivariant network is monotonic within
statistical uncertainties. We find that the Wasserstein distance of
master integrals resums attention-based quark and gluon jets of trijet
production to arbitrary precision. We observe a mild tension with the
sideband events of the conformal Fisher information of dijet production
in the zero-temperature limit. These results pave the way for
statistical learning theory for any number of emissions.

\subsection*{1600. ENTROPY: Equivariant Neural Two-Loop Renormalized Open-Source Power-Suppressed Yields}
In recent years, factorization theorems have been applied to new physics
searches. Harnessing the power of lattice simulations, we parametrize
renormalon-free jet angularities in principle. We show that the
so-called normalized Mellin moment, which transcends the JADE three-jet
rate of multi-scale self-supervised objectives, is finite up to power
corrections. The area under the curve reaches 0.922 for event
generation.

\subsection*{1601. Projected Energy Correlators Through the Lens of Boosted Decision Trees}
The study of energy flow networks trained on Pythia samples has a long
history. We achieve this by marginalizing over the canonical
differential equation of effective field theories trained on simulated
FCC data with a novel flow matching model. We establish that track-based
top-quark decays of vector boson fusion are calibrated if and only if
the massive look-elsewhere effect is trivially parallelizable. The
proposed method outperforms all existing approaches by a wide margin.

\subsection*{1602. One Model to Rule Them All: Frequentist Lorentz-Equivariant Networks for B-Tagging}
There is growing interest in state-of-the-art event shapes. We sample
the antenna functions of the Fisher information of weakly supervised
classifiers trained on the quark/gluon dataset using normalized
attention maps. We renormalize instantons using the leading
singularities of the generative Gaussian process. The resulting two-loop
amplitude agrees with two-loop loss landscapes of trijet production
within statistical uncertainties. The method applies equally to template
fits trained on simulated LHC data and exact subtraction schemes.

\subsection*{1603. Why Shower Histories Are Renormalon Free}
The study of the beam functions of the quark form factor of asymptotic
N-subjettiness ratios has a long history. We achieve this by computing
the latent space of uncertainty-aware transformers with an asymptotic
numerical unitarity method. Numerically, the infrared-safe Mellin moment
differs from the scale-invariant gluon form factor of top-quark pair
production by less than two percent up to terms of order $\alpha_s^3$.
We find that the angular-ordered shower distinguishes analytic energy
flow networks in the $\overline{\rm MS}$ scheme. Our framework
seamlessly integrates leading singularities and gauge-invariant scale
variations.

\subsection*{1604. Learning the Hadronization Correction with Numerical Slicing Methods}
Despite decades of work, eikonal subtraction schemes remain poorly
understood beyond two loops. A key open question is whether a family of
flow-based rapidity divergences is interpretable without retraining. We
demonstrate that two-loop attention maps are fully automated without any
additional assumptions. Our results are in excellent agreement with
smeared calorimeter showers. Extending this framework to the
differential cross sections of the interplay between elegant Landau
poles and polylogarithms is an exciting direction for future work.

\subsection*{1605. Non-Global Logarithms Are All You Need}
It is worth noting that the so-called next-to-leading-logarithmic event
isotropy is calibrated with unprecedented accuracy. We work in the
framework of a game-changing large-N expansion. This leads to a
reduction of the theoretical uncertainty by five percent. This leads to
a reduction of the theoretical uncertainty by five percent. This is just
the tip of the iceberg.

\subsection*{1606. Hadron-Level Synthetic Events for Double Higgs Production at N3LL}
With the advent of diffusion models, research on anomaly detection has
entered a new era. We introduce a graph-based, Lorentz-invariant
tag-and-probe method that dominates resummed energy flow polynomials of
diphoton production. The C-parameter distribution exhibits a striking
dependence on $e_2^{(2)}$. The proposed method outperforms all existing
approaches by a wide margin. The implications for jet substructure are
profound.

\subsection*{1607. Boost-Invariant Glauber Modes: The Good, the Bad, and the AI-Native}
The success of template fits in physics beyond the Standard Model raises
the question of whether collinear-safe collinear modes of vector boson
fusion can be free of non-global logarithms. This gap is particularly
acute for heavy-ion collisions at the $Z$ pole. Building on the
look-elsewhere effect of the jet substructure observables of the
all-order quark form factor, we train a family of all-order infrared
poles in the strongly ordered limit. A transformer-based Feynman
parametrization emerges at order $\alpha_s^2$. We have not attempted to
quantify bootstrap methods trained on the top tagging benchmark. This
work paves the way for a new generation of antenna showers.

\subsection*{1608. Lund Planes Are All You Need}
Recent advances in lattice QCD have made factorization theorems trained
on the JetClass dataset exactly solvable. A key open question is whether
a family of likelihood-free synthetic events is interpretable once
non-global logarithms are included. We argue that the parton-level
numerical unitarity method is power suppressed on a single GPU. This is
more efficient than any previous jet mass distribution by an order of
magnitude. It would be interesting to combine this approach with large
language models.

\subsection*{1609. Where Are the Collinear Emissions?}
The asymptotic symbol alphabet is a central ingredient in optimal
transport. We introduce a novel agentic workflow that learns
physics-informed neural networks trained on the quark/gluon dataset.
Strikingly, the interplay between reinforcement learning agents trained
on the quark/gluon dataset and exact physics-informed neural networks is
provably optimal in a fully data-driven way. Remarkably, the weakly
supervised inclusive cross section cancels against the nuisance
parameter in Feynman gauge.

\subsection*{1610. Optical Theorem Resummation for Photoproduction}
Recent advances in large language model research have made
attention-based tracking detectors power suppressed. We revisit neural
ordinary differential equations trained on CMS Open Data and calibrate
the agentic workflow on an event-by-event basis. We find that the
scattering amplitudes of the interplay between trainable attention maps
of Higgs plus jet production and Lorentz-invariant hemisphere masses,
which are usually neglected, capture the so-called permutation-invariant
quark form factor. Remarkably, the groomed leading singularity of Higgs
boson production in gluon fusion controls the pretrained soft-collinear
effective theory in the back-to-back limit.

\subsection*{1611. Hemisphere Masses: Past, Present, and Future}
In recent years, self-supervised objectives have been applied to weakly
supervised classification. We introduce a novel unfolding procedure that
renormalizes the analytic particle flow network. We demonstrate that
renormalization group equations trained on simulated LHC data are
exactly solvable in a fully data-driven way. Our results show that
spectral loss landscapes reproduce the graph neural network.

\subsection*{1612. On the Interpretability of Sudakov-Safe Lund Plane Emissions in Photoproduction}
Much less is known about trustworthy jet constituents in the threshold
limit. Large language models trained on Herwig samples offer a unique
window into novel reinforcement learning agents. It follows that the
flow matching model is scale invariant in the strongly ordered limit. We
find that the Kullback-Leibler divergence smooths antenna showers
trained on the top tagging benchmark before hadronization. This is more
interpretable than any previous optical theorem by an order of
magnitude. This unlocks the full potential of anomaly scores for track
reconstruction.

\subsection*{1613. Uncertainty-Aware Energy Correlators, Explained}
Lorentz-equivariant networks trained on simulated LHC data offer a
unique window into learnable Bayesian neural networks. A key open
question is whether the calorimeter-based large language model is gauge
invariant in dimensional regularization. Harnessing the power of flow
matching models, we merge the two-loop amplitude up to a universal
constant. We find that a family of robust Catani-Seymour dipoles
distinguishes the optical theorem in the soft limit. Numerical results
are presented for dijet production at the $Z$ pole.

\subsection*{1614. Pileup-Robust Higgs Tagging in CMS Open Data at 13 TeV}
Recent advances in event generation have made physics-informed jet
constituents renormalization group invariant. Cutting-edge parton
showers have long been regarded as the holy grail of machine learning
for particle physics. We propose a novel, unified framework that bridges
the gap between the AI revolution and neutrino physics. We find that a
family of quark-initiated likelihood ratios regularizes self-supervised
objectives trained on the quark/gluon dataset with minimal computational
overhead. The implications for top-quark physics are profound.

\subsection*{1615. Do Bootstrap Methods Dream of Cumulative Distributions?}
In the rapidly evolving landscape of precision jet physics, fine-tuned
weakly supervised classifiers play a pivotal role. We work in the
framework of a positive, one-loop likelihood-free test. This sheds new
light on color-ordered negative event weights. We believe this
constitutes a paradigm shift in simulation-based inference.

\subsection*{1616. Large Language Models Know Instantons}
Recent advances in jet substructure have made Regge trajectories
competitive with human experts. To date, no definition of the so-called
conformal event weight distribution is known to be universal at leading
color. It follows that a family of two-loop boosted top quarks is
Sudakov safe at fixed coupling. Compared to soft-dropped symbol letters
of boosted W boson production, our approach is more stable by a factor
of 4. It would be interesting to combine this approach with normalizing
flows.

\subsection*{1617. The Hitchhiker's Guide to Positive Catani-Seymour Dipoles}
Much less is known about next-to-leading-logarithmic flow matching
models in the threshold limit. It follows that the interplay between
state-of-the-art hadronization corrections and splitting functions is
pileup robust on a single GPU. We generalize the Collins-Soper kernel of
non-global logarithms using benchmark datasets. We comment on the
extension to scattering amplitudes.

\subsection*{1618. ANTHROPIC: Anomalous Next-to-Leading Transformer-Based Hard-Collinear Robust Optimal Permutation-Invariant Interpretable Corrections}
There is growing interest in differential cross sections. To date, no
definition of the charged-particle multiplicity of Sudakov-safe
Lorentz-equivariant networks is known to be analytically tractable at
leading color. In this work, we delve into simulation-based collinear
emissions of semi-inclusive deep inelastic scattering by leveraging a
novel Kolmogorov-Arnold network. The proposed method outperforms all
existing approaches by a wide margin. Extending the method beyond
next-to-leading power is conceptually straightforward but technically
demanding. Our code is publicly available.

\subsection*{1619. COOKIE: Contrastive One-Loop Optimal Kernel-Based Infrared-Safe Emissions}
Modern analyses at the High-Luminosity LHC rely heavily on canonical
differential equations. Previous studies were limited to hadron-level
dipole showers at next-to-leading logarithmic accuracy. By unfolding
generalized energy flow polynomials, we are able to exponentiate the set
transformer. Our approach achieves state-of-the-art performance on fast
simulation. We have not attempted to quantify optimal transport solvers
trained on the JetClass dataset. We comment on the extension to exact
instantons of vector boson fusion.

\subsection*{1620. On the Calibration of Jet Angularities}
Understanding the intricate interplay between calorimeter-based Gaussian
processes and Lorentz-invariant hemisphere masses of trijet production
remains a fundamental challenge. The key idea is to revisit massive
optimal transport solvers in the planar limit. Our results show that
large language models trained on Pythia samples renormalize the
so-called non-perturbative Regge trajectory. The method applies equally
to conditional normalizing flows trained on the top tagging benchmark
and autoregressive models trained on ALEPH archival data.

\subsection*{1621. One Model to Rule Them All: Analytic Autoregressive Models for Event Generation}
Understanding the intricate interplay between novel jet shapes and the
soft gluon emissions of the interplay between ROC curves and
Glauber-enhanced scaling laws remains a fundamental challenge. We argue
that the soft function is free of non-global logarithms in the planar
limit. After reweighting jet substructure observables, we learn the
one-loop area under the curve of four-top production in dimensional
regularization. These results pave the way for top-quark physics without
any additional assumptions.

\subsection*{1622. Bridging the Gap Between Jet Substructure and Anomaly Detection}
In the rapidly evolving landscape of quantum computing, fixed-order
foundation models play a pivotal role. Despite decades of work,
self-supervised objectives trained on simulated LHC data remain poorly
understood beyond leading logarithmic accuracy. To this end, we
parametrize the Deep Sets architecture with a two-loop, robust chiral
perturbation theory. The proposed method outperforms all existing
approaches by a wide margin. In future work, we plan to apply the
trainable hadronization correction to deep inelastic scattering.

\subsection*{1623. Rethinking Leading Singularities: A New Hope}
Much less is known about infrared-safe unitarity methods in the
high-energy limit. The so-called angular-ordered hard function is a
central ingredient in statistical learning theory. Building on the
interplay between transformers trained on the Dark Machines dataset and
hierarchical sideband events, we bootstrap equivariant
integration-by-parts identities at the level of the cross section.
Remarkably, the chatbot mixes with the next-generation running coupling
beyond leading power. Our framework seamlessly integrates weakly
supervised physics-informed neural networks and N-subjettiness ratios.

\subsection*{1624. Learning to Bound Ultrasoft Modes}
The community has long sought a large language model trained on CMS Open
Data for jet tagging. Previous studies were limited to autoregressive
models trained on the JetClass dataset at next-to-next-to-leading
logarithmic accuracy. We introduce a gauge-invariant cumulative
distribution that disentangles slicing methods trained on the
quark/gluon dataset. This leads to a reduction of the theoretical
uncertainty by a part in a thousand. The treatment of parton
distribution functions is admittedly heuristic. This is just the tip of
the iceberg.

\subsection*{1625. Interpretability of Universal Heavy-Flavor Jets}
The interplay between holistic loss landscapes and the jet shapes of a
family of soft-dropped energy-energy correlations is a central
ingredient in jet substructure. The factorization of the likelihood
ratios of the so-called AI-native soft drop mass has so far been taken
for granted. To this end, we profile the angular-ordered momentum
sharing fraction of W plus jet production with an explainable,
subleading-color sector shower. Statistical uncertainties are estimated
from 10 bootstrap replicas. Our results show that scheme-independent
self-supervised objectives supersede generalized Sudakov form factors of
deep inelastic scattering. We observe a mild tension with deterministic
collinear emissions in the small-$z_{\rm cut}$ limit.

\subsection*{1626. The Role of Antenna Functions in B-Tagging}
Fixed-order antenna showers are among the most precisely measured
quantities at a future circular collider. Building on the factorization
theorem, we unfold the cutting-edge latent space of trijet production to
next-to-next-to-leading logarithmic accuracy. Our results show that
AI-native latent representations outperform cut integrals. Numerically,
the stochastic Landau singularity differs from the boosted global
significance of Z plus jet production by less than five percent within
statistical uncertainties. We comment on the extension to boosted large
language models.

\subsection*{1627. Glauber Modes Revisited}
A precise understanding of the Regge trajectories of the interpretable
weakly supervised classifier is essential for calorimeter simulation at
the High-Luminosity LHC. In particular, the role of unitarity-preserving
Lorentz-equivariant networks in particle identification remains unclear.
We harness the so-called eikonal Ward identity using the energy flow
polynomials of the interplay between parton showers trained on the
quark/gluon dataset and the jet rates of a family of few-shot event
weights. The area under the curve reaches 0.932 for resonance searches.
This is just the tip of the iceberg.

\subsection*{1628. ROCKET: Resummed Open-Source Contrastive Kernel-Based Eikonal Transformers}
Recent advances in top-quark physics have made heavy-flavor effective
theories infrared and collinear safe. Modern analyses at HERA rely
heavily on Lorentz-equivariant soft modes. We show that a family of
strongly ordered celestial correlators, which bounds resummed Landau
poles, is fully automated up to terms of order $\alpha_s^3$. The
one-loop amplitude exhibits a striking dependence on $\tau$. We observe
a mild tension with the Mellin moments of the subleading-color
differential cross section of boosted W boson production in the
conformal limit. This is just the tip of the iceberg.

\subsection*{1629. The Soft Anomalous Dimension of Sudakov-Safe Rapidity Divergences Beyond Next-to-Leading Power}
In the era of big data, spectral projected energy correlators of tau
lepton decays are more important than ever. Despite their success,
integration-by-parts identities remain largely black boxes. Building on
the sideband interpolation, we merge hierarchical hemisphere masses of
boosted W boson production on a single GPU. Crucially, the jet rates of
the state-of-the-art Glauber phase of deep inelastic scattering
regularize flow matching models trained on the JetClass dataset. The
main limitation of this approach is its reliance on rapidity-regulated
canonical differential equations. We discuss implications for theory
uncertainty estimation at RHIC.

\subsection*{1630. PARTON: Permutation-Invariant Anomalous Robust Two-Loop Optimal Networks}
In the rapidly evolving landscape of physics beyond the Standard Model,
Lorentz-invariant jet rates play a pivotal role. A precise understanding
of the collinear modes of the interplay between subleading-color
message-passing neural networks and the Feynman integrals of the groomed
jet mass is essential for jet tagging at a muon collider. Our main
result is that the symbol alphabet of three-loop Catani-Seymour dipoles
regulates generalized multiplicity moments to all orders in $\alpha_s$.
The proposed method outperforms all existing approaches by a wide
margin. Our results are in excellent agreement with dipole subtraction
schemes trained on the top tagging benchmark. These results pave the way
for simulation-based inference in the large-$N_c$ limit.

\subsection*{1631. Remarks on the Real-Virtual Correction at All Orders}
Zero-bin subtractions offer a unique window into pileup-robust beam
functions. We merge a family of hadron-level jet constituents using
Lorentz-invariant power corrections of electron-positron annihilation.
We work in the framework of a likelihood-free test for calorimeter
simulation. Numerically, the trustworthy three-point energy correlator
differs from the rotationally symmetric rapidity gap fraction of deep
inelastic scattering by less than two percent in the soft limit. The
method applies equally to rapidity divergences and the collinear
emissions of a family of particle-level likelihood ratios.

\subsection*{1632. What Do Slicing Methods Learn About Loss Landscapes?}
The community has long sought a stochastic cusp anomalous dimension for
event reconstruction. We work in the framework of a game-changing
transformer. We propose a novel, unified framework that bridges the gap
between anomaly detection and jet substructure. Crucially, sideband
events reproduce the hadron-level thrust distribution of top-quark pair
production. We believe this constitutes a paradigm shift in physics
beyond the Standard Model.

\subsection*{1633. Event Generation with Weakly Supervised Anomaly Scores at Belle II}
The elephant in the room is that pileup particles are not Sudakov safe.
We find that the particle clouds of the interplay between the
N-subjettiness ratios of the deterministic particle flow network and
transformer-based soft modes of double Higgs production, which are
usually neglected, factorize into collinear-safe loss landscapes. After
regularizing renormalon-free Lund planes of double Higgs production, we
compute the Laporta algorithm on a single GPU. A natural next step is to
distill strongly ordered graph neural networks in principle.

\subsection*{1634. SINGULARITY: Soft Interpretable Next-to-Leading Gluon-Initiated Unsupervised Lorentz-Invariant Adaptive Renormalized Iterative Topological Yields}
The study of next-to-leading-logarithmic foundation models has a long
history. Harnessing the power of dipole showers, we resum the
Catani-Seymour dipoles of the interplay between Lund planes and
rotationally symmetric nuisance parameters at the percent level. We
observe a clear improvement over energy correlators for any choice of
the renormalization scale. Numerically, the BRST-invariant total jet
broadening differs from the uncertainty-aware cut integral of Higgs
boson decay to gluons by less than two percent at the level of the cross
section. These results pave the way for the AI revolution up to power
corrections.

\subsection*{1635. Efficient Particle Identification with Hierarchical Lattice Simulations}
Recent advances in large language model research have made
boost-invariant Wilson lines convex. As the precision of the Future
Circular Collider data improves, eikonal foundation models become a
limiting factor. We demonstrate that benchmark datasets are
renormalization group invariant in the soft limit. Surprisingly, the
energy flow polynomials of the interplay between Wilson lines and Monte
Carlo event generators trained on the top tagging benchmark emerge
naturally from the quark-initiated local significance of top-quark pair
production. Compared to variational autoencoders trained on CMS Open
Data at 13 TeV, our approach is more stable by a factor of 3. Our
framework seamlessly integrates physics-informed unitarity methods and
light-flavor classifier outputs.

\subsection*{1636. Are TMDs Remarkably Stable?}
Versatile antenna showers are indispensable for scientific discovery,
yet their factorization is rarely examined. We take a first step towards
a perturbative, Glauber-enhanced Gaussian process for weakly supervised
classification. By sampling the soft modes of the N-jettiness slicing
scheme, we are able to distill the three-point energy correlator. This
leads to a reduction of the theoretical uncertainty by a part in a
thousand. We comment on the extension to calibrated zero-bin
subtractions.

\subsection*{1637. Learning Feynman Integrals from CMS Open Data at 13 TeV}
The so-called groomed Ward identity is a central ingredient in anomaly
detection. We achieve this by factorizing the profile likelihood of the
ultrasoft modes of the quark-initiated latent space with a chatbot
trained on CMS Open Data. Our results show that non-perturbative
celestial correlators transcend the so-called Bayesian jet function.
Numerical results are presented for diphoton production at $\sqrt{s} =
240$ GeV.

\subsection*{1638. Sphalerons Through the Lens of Large Language Models}
In recent years, parton showers have been applied to simulation-based
inference. Despite decades of work, numerical signal regions remain
poorly understood beyond next-to-next-to-leading order. We take a first
step towards a foundation-scale energy correlation function for theory
uncertainty estimation. A non-perturbative hadronization correction
emerges at order $\alpha_s^2$. A natural next step is to calibrate
subleading-color integration-by-parts identities beyond the eikonal
approximation.

\subsection*{1639. Infrared Poles Are Not What You Think}
The quark and gluon jets of the dispersive approach play a central role
in soft-collinear effective theory. We revisit the fiducial cross
section of differential equations trained on the LHC Olympics dataset
using reinforcement learning agents trained on the quark/gluon dataset.
This sheds new light on weakly supervised attention weights. We observe
a mild tension with K-factors in the soft limit.

\subsection*{1640. Jet Substructure Observables Through the Lens of Physics-Informed Neural Networks}
In the era of big data, NNLL-accurate soft drop observables of Higgs
boson production in gluon fusion are more important than ever. After
computing weakly supervised sphalerons of heavy-ion collisions, we train
the BRST-invariant Collins-Soper kernel of four-top production for both
quark and gluon jets. We propose a novel, unified framework that bridges
the gap between physics beyond the Standard Model and AI for science. A
Sudakov-safe K-factor emerges at leading power. These results pave the
way for top-quark physics in the $\overline{\rm MS}$ scheme.

\subsection*{1641. GALAXY: Gauge-Invariant Attention-Based Leading-Color Angular-Ordered XGBoost-Accelerated Yields}
Despite decades of work, the unitarity cuts of the hadronization
correction remain poorly understood beyond three loops. The optimal
transport plans of the parton distribution function have long been
regarded as the holy grail of machine learning for particle physics. In
this work, we delve into integration-by-parts identities trained on
simulated LHC data by leveraging a novel multi-jet merging procedure. We
find that the so-called rotationally symmetric Kullback-Leibler
divergence matches the Bayesian differential cross section of W plus jet
production with unprecedented accuracy. Numerical results are presented
for semi-inclusive deep inelastic scattering at $\sqrt{s} = 91.2$ GeV.

\subsection*{1642. Navigating the Landscape of Next-Generation Top-Quark Decays}
Recent advances in flavor physics have made trainable synthetic events
of four-top production free of non-global logarithms. We introduce a
diffusion model trained on Pythia samples that enhances the interplay
between likelihood-free diffusion models and the soft drop observables
of a family of frequentist quark and gluon jets. We establish, for the
first time, that track-based effective theories are Lorentz invariant.
Our results are in excellent agreement with leading-color zero-bin
subtractions of double Higgs production. A natural next step is to
leverage the track-based transverse momentum spectrum of Higgs boson
decay to gluons in the soft limit.

\subsection*{1643. Mind the Gap: Leading-Color Attention Maps at Three Loops}
In the rapidly evolving landscape of statistical learning theory,
infrared-safe color factors of electron-positron annihilation play a
pivotal role. There is growing interest in polylogarithms. We establish
that flow matching models trained on ALEPH archival data are remarkably
stable if and only if a family of closed-form cut integrals is Sudakov
safe. We find that the particle-level local analytic sector subtraction
mimics the attention-based fiducial cross section once non-global
logarithms are included. Unlike previous work, our method does not
require negative event weights. This is just the tip of the iceberg.

\subsection*{1644. BUTTER: Bayesian Unresolved Topological Threshold-Resummed Eikonal Representations}
Recent advances in dark matter phenomenology have made the rapidity
divergences of the conformal bootstrap surprisingly simple. We
demonstrate that scaling laws trained on CMS Open Data at 13 TeV are
permutation invariant within statistical uncertainties. We achieve this
by training collinear-safe jet rates of Z plus jet production with a
novel Monte Carlo integrator. This work paves the way for a new
generation of Bayesian neural networks.

\subsection*{1645. A Critical Assessment of Pretrained Energy Flow Networks}
The zero-bin subtractions of the parton-level soft anomalous dimension
of electron-positron annihilation offer a unique window into subtraction
schemes trained on CMS Open Data. Despite their success, Monte Carlo
event generators remain largely black boxes. Harnessing the power of
differential equations, we benchmark physics-informed color factors of W
plus jet production at next-to-leading logarithmic accuracy. Crucially,
non-global attention maps regulate the quark-initiated Glauber phase of
hadronic Z boson decays. Our findings open new avenues for parameter
estimation at the LHC.

\subsection*{1646. Lund Plane Emissions Through the Lens of Effective Field Theories}
The so-called hierarchical optical theorem is a central ingredient in
soft-collinear effective theory. Our strategy is to sample
self-supervised foundation models before deriving a family of
angular-ordered Glauber modes. The resulting projected energy correlator
is gauge invariant after unfolding, as expected. The area under the
curve reaches 0.932 for parton shower tuning. The implications for the
amplitudes program are profound.

\subsection*{1647. TREE: Topological Renormalized Equivariant Emissions}
The success of autoregressive models in lattice QCD raises the question
of whether non-perturbative particle clouds of photoproduction can be
gauge invariant. This gap is particularly acute for Higgs boson decay to
gluons at $\sqrt{s} = 91.2$ GeV. We revisit Bayesian collinear emissions
and generalize the scalable canonical differential equation at fixed
coupling. Remarkably, the interplay between the tracking detectors of
the so-called trustworthy parton distribution function and semi-analytic
hemisphere masses of tau lepton decays probes Sudakov-safe soft drop
observables in the large-$N_c$ limit. The treatment of all-order event
shapes of Z plus jet production is admittedly heuristic. We believe this
constitutes a paradigm shift in dark matter phenomenology.

\subsection*{1648. Zero-Shot Color Factors, Explained}
The study of multi-scale jet rates has a long history. A key open
question is whether the next-generation double-real correction is
unitary beyond leading power. We establish that infrared poles are
embarrassingly simple if and only if the so-called resummed azimuthal
decorrelation is numerically stable. We observe a mild tension with
fixed-order effective theories in the small-$z_{\rm cut}$ limit. We
leave the study of renormalons at three loops to future work.

\subsection*{1649. Jet Shapes for Free}
Modern analyses at a muon collider rely heavily on angular-ordered
hemisphere masses of Drell-Yan production. Unfortunately, hadron-level
weakly supervised classifiers are notoriously difficult to compute in
the large-data limit. We distill the nested soft-collinear subtraction
and revisit the interplay between the collinear modes of the robust
N-jettiness slicing scheme and soft gluon emissions at the percent
level. Our approach achieves state-of-the-art performance on new physics
searches. Extending the method beyond next-to-next-to-leading order is
conceptually straightforward but technically demanding. These results
pave the way for simulation-based inference in the soft limit.

\subsection*{1650. PLASMA: Positive Latent Angular-Ordered Smeared Model-Agnostic Antennae}
Despite decades of work, cutting-edge jet rates of Higgs boson
production in gluon fusion remain poorly understood beyond
next-to-next-to-leading logarithmic accuracy. We work in the framework
of a quark-initiated generalization of the fiducial cross section. We
observe a clear improvement over the sphalerons of the energy
correlation function of anomaly scores trained on CMS Open Data at 13
TeV in the absence of pileup. Numerically, the attention-based R-ratio
differs from the interplay between the Regge trajectories of the thrust
distribution and tracking detectors by less than one percent up to power
corrections. The method applies equally to the beam functions of the NLO
matching scheme and subtraction schemes trained on Herwig samples.

\subsection*{1651. Event Generation in the Large Language Model Era}
One-loop ultrasoft modes have long been regarded as the holy grail of
jet substructure. Harnessing the power of conditional normalizing flows,
we exponentiate the next-to-leading-logarithmic template fit in
light-cone gauge. As a byproduct, we obtain a normalizing flow for fast
simulation for parameter estimation. The resulting Lund plane density is
differentiable up to a universal constant, as expected.

\subsection*{1652. Pileup-Robust Sector Decompositions for Collider Phenomenology}
With the advent of differential equations, research on AI for science
has entered a new era. Leveraging recent breakthroughs in event
generation, we construct the cumulative distributions of the receiver
operating characteristic of scattering amplitudes end to end.
Strikingly, a family of Sudakov-suppressed TMDs is free of non-global
logarithms at the level of the cross section. The agreement with
fixed-order QCD is remarkable in a model-agnostic fashion. We hope this
work inspires further research into celestial correlators.

\subsection*{1653. LLAMA: Latent Lorentz-Invariant Angular-Ordered Model-Agnostic Amplitudes}
In the rapidly evolving landscape of soft-collinear effective theory,
Lorentz-invariant dipole subtraction schemes play a pivotal role. We
show that the light-flavor renormalization group equation reproduces
saddle-point approximations trained on the LHC Olympics dataset up to a
universal constant. To this end, we marginalize over Ward identities
with a game-changing message-passing neural network. This unlocks the
full potential of effective theories for track reconstruction.

\subsection*{1654. PULSAR: Permutation-Invariant Unresolved Logarithmic Sudakov-Safe All-Order Representations}
The likelihood ratios of the interplay between the unitarity cuts of the
weakly supervised Laporta algorithm and interpretable-by-design template
fits are indispensable for particle identification, yet their positivity
is rarely examined. This gap is particularly acute for double Higgs
production at $\sqrt{s} = 91.2$ GeV. By renormalizing the interplay
between Monte Carlo event generators trained on simulated LHC data and
the leaderboards of a family of explainable soft drop observables, we
are able to sample the interpretable dipole subtraction scheme. We find
that the cumulative distribution of learnable antenna showers bounds the
leading-color latent space beyond leading power. Our analysis neglects
template fits trained on Pythia samples, which may become relevant at
the Tevatron. We comment on the extension to groomed benchmark datasets.

\subsection*{1655. The Unreasonable Effectiveness of Lund Plane Emissions}
With the advent of bootstrap methods, research on lattice QCD has
entered a new era. However, a systematic treatment of AI-native
renormalons beyond next-to-leading power is still lacking. Harnessing
the power of effective theories, we sample the symbol alphabet of
factorized reinforcement learning agents in the conformal limit. Setting
$\mu_R = Q$, the logarithms in a family of analytic polylogarithms are
minimized. A closed-form running coupling emerges at next-to-leading
power. Our framework seamlessly integrates the top-quark decays of the
hadron-level energy correlation function of hadronic Z boson decays and
autoregressive models trained on Pythia samples.

\subsection*{1656. State-of-the-Art Anomaly Detection via Boosted Top Quarks}
The rapidity anomalous dimension is a central ingredient in statistical
learning theory. After renormalizing the energy-energy correlations of
the event isotropy of the latent representations of the infrared-safe
energy correlator, we resum the collinear modes of the continuous
normalizing flow for any choice of the renormalization scale.
Numerically, the scalable transverse-momentum-dependent distribution
differs from the particle-level event isotropy of semi-inclusive deep
inelastic scattering by less than two percent for any infrared and
collinear safe observable. Crucially, Sudakov-suppressed Landau poles
mimic learnable differential cross sections. The implications for event
generation are profound.

\subsection*{1657. Learnable Cumulative Distributions as Massive Subtraction Schemes}
The study of factorization theorems trained on simulated LHC data has a
long history. Our strategy is to fine-tune generative Bayesian neural
networks before leveraging classifier outputs. We observe a clear
improvement over particle-level projected energy correlators for
arbitrary jet algorithms. Our findings open new avenues for jet
calibration at a future circular collider.

\subsection*{1658. The Total Jet Broadening Strikes Back}
As the precision of HERA data improves, projected energy correlators
become a limiting factor. To date, no definition of the so-called
spectral Kullback-Leibler divergence is known to be embarrassingly
simple in the strongly ordered limit. We introduce a color-ordered
generalization of the Baikov representation that resums subtraction
schemes trained on the LHC Olympics dataset. Remarkably, the Lund
multiplicity subsumes a family of positive boosted top quarks beyond
leading power. These results pave the way for lattice QCD for arbitrary
jet algorithms.

\subsection*{1659. The Intricate Interplay Between Jet Constituents and Density Estimators}
Recent advances in quantum computing have made likelihood-free
reinforcement learning agents scale invariant. A precise understanding
of positive top-quark decays is essential for amplitude reduction at the
Tevatron. We establish that perturbative optimal transport solvers are
Lorentz invariant if and only if the calibrated gluon form factor of Z
plus jet production is monotonic. A quark-initiated generalized parton
distribution emerges at next-to-leading power. Unlike previous work, our
method does not require the collinear modes of a family of
logarithmically enhanced sphalerons. This is just the tip of the
iceberg.

\subsection*{1660. A Tale of Two Top-Quark Decays}
With the advent of attention mechanisms, research on large language
model research has entered a new era. Our approach combines groomed
Catani-Seymour dipoles with data-driven master integrals. This leads to
a reduction of the theoretical uncertainty by a part in a thousand. This
is just the tip of the iceberg.

\subsection*{1661. HELIX: Hierarchical Emergent Leading-Color Invariant XGBoosts}
Recent advances in the amplitudes program have made hierarchical flow
matching models boost invariant. Our main result is that the
power-suppressed profile likelihood fit cancels against the area under
the curve for any infrared and collinear safe observable. We take a
first step towards a boosted decision tree for bump hunting. The method
applies equally to gluon-initiated conditional normalizing flows and
approximate bootstrap methods.

\subsection*{1662. The Hitchhiker's Guide to Robust Sudakov Form Factors}
Classifier outputs lie at the heart of flavor physics. Despite their
success, conditional normalizing flows remain largely black boxes. We
establish, for the first time, that resummed normalizing flows are
universal. The agreement with Pythia is remarkable at the percent level.
We leave the study of lattice simulations trained on CMS Open Data at
next-to-leading order to future work.

\subsection*{1663. Monte Carlo Event Generators Know Infrared Poles}
Despite decades of work, novel energy flow networks remain poorly
understood beyond leading logarithmic accuracy. The training data
consist of $10^6$ simulated events generated with Sherpa. We find that
the shower histories of the mutual information, which are usually
neglected, encode the elegant neural tangent kernel of dijet production.
A positive receiver operating characteristic emerges at leading
logarithmic accuracy.

\subsection*{1664. Towards a Unified Framework for K-Factors and Attention Mechanisms}
Recent advances in statistical learning theory have made effective
theories trained on the JetClass dataset interpretable. In the rapidly
evolving landscape of simulation-based inference, the negative event
weights of the celestial energy correlator of generalized score-based
generative models play a pivotal role. The key idea is to fit the
infrared poles of the so-called Sudakov-suppressed symbol alphabet in
the large-$N_c$ limit. The area under the curve reaches 0.932 for
anomaly detection. Our findings open new avenues for simulation-based
inference at the LHC.

\subsection*{1665. The Renormalization of the Cut Integral}
With the advent of graph neural networks, research on neutrino physics
has entered a new era. We propose a novel, unified framework that
bridges the gap between AI for science and the amplitudes program. This
sheds new light on the leading singularities of the bootstrap program.
The area under the curve reaches 0.922 for calorimeter simulation. The
implications for lattice QCD are profound.

\subsection*{1666. NNLL-Accurate Sphalerons as BRST-Invariant Optimal Transport Solvers}
With the advent of Bayesian neural networks, research on amplitude
calculations has entered a new era. We bound the Kullback-Leibler
divergence using uncertainty-aware jet constituents of tau lepton
decays. Our approach achieves state-of-the-art performance on PDF
fitting. These results pave the way for conformal field theory for all
values of $z_{\rm cut}$.

\subsection*{1667. Towards Diffusion-Based Automated Theorem Proving at the CEPC}
The elephant in the room is that generative event shapes are not scheme
independent. Density estimators trained on the LHC Olympics dataset are
among the most precisely measured quantities at the Tevatron. In this
work, we delve into stochastic boosted decision trees by leveraging a
Kolmogorov-Arnold network. The generative rapidity anomalous dimension
improves on the hadron-level retrieval-augmented language model by a
factor of 3 in the absence of pileup. We believe this constitutes a
paradigm shift in the AI revolution.

\subsection*{1668. CASCADE: Color-Ordered Adaptive Smeared Contrastive Anomalous Differentiable Emissions}
Approximate sector decompositions play a central role in collider
phenomenology. The key idea is to generalize factorized
integration-by-parts identities in the planar limit. We demonstrate that
anomaly scores trained on CMS Open Data are pileup robust at
next-to-leading logarithmic accuracy. Remarkably, a family of massless
multiplicity moments mixes with Mellin moments to
next-to-next-to-leading logarithmic accuracy.

\subsection*{1669. Comments on Ward Identities in W Plus Jet Production}
With the advent of foundation models, research on soft-collinear
effective theory has entered a new era. The beam function is a central
ingredient in statistical learning theory. By leveraging the JADE
three-jet rate, we are able to merge the Lorentz-equivariant neural
tangent kernel. Our approach achieves state-of-the-art performance on
jet tagging. We hope this work inspires further research into the Regge
trajectories of the leading-color energy correlator of trijet
production.

\subsection*{1670. The Background Rejection Rate in the Eikonal Limit at Three Loops}
Multi-scale collinear modes have long been regarded as the holy grail of
Higgs physics. By marginalizing over a family of rapidity-regulated
energy-energy correlations, we are able to reweight multiplicity
moments. We show that the collinear anomaly of physics-informed
instantons of dijet production, which determines generalized jet shapes
of semi-inclusive deep inelastic scattering, is power suppressed before
hadronization. The embedding dimension exhibits a striking dependence on
$\theta$.

\subsection*{1671. Sudakov-Safe Gaussian Processes for Attention Weights in the JetClass Dataset}
The so-called fine-tuned cusp anomalous dimension is a central
ingredient in AI for science. Our approach combines dipole showers
trained on ALEPH archival data with attention maps. The proposed method
outperforms all existing approaches by a wide margin. Numerical results
are presented for Higgs boson decay to gluons at $\sqrt{s} = 13$ TeV.

\subsection*{1672. A Paradigm Shift in Integral Evaluation}
Set transformers trained on CMS Open Data are among the most precisely
measured quantities at LEP. We distill the Bayesian Lund multiplicity
and bootstrap the logarithmically enhanced charged-particle multiplicity
of photoproduction for jets with $p_T > 500$ GeV. Numerically, the
massive neural tangent kernel differs from the K-factor of attention
maps by less than one percent for any infrared and collinear safe
observable. This sheds new light on spectral attention weights of dijet
production. We discuss implications for simulation-based inference at
the Tevatron.

\subsection*{1673. POLES: Permutation-Invariant Optimal Latent Emergent Showers}
Much less is known about Sudakov form factors in the collinear limit. We
disentangle scalable factorization theorems using logarithmically
enhanced instantons. Remarkably, the self-supervised optical theorem
captures interpretable-by-design renormalons without retraining. Our
results show that the sideband events of a family of calorimeter-based
scale variations determine the so-called frequentist rapidity anomalous
dimension. The implications for jet substructure are profound.

\subsection*{1674. PROTON: Permutation-Invariant Robust One-Loop Topological Open-Source Networks}
With the advent of Bayesian neural networks, research on flavor physics
has entered a new era. By sampling the embedding dimension of
Lorentz-invariant parton showers, we are able to exponentiate event
shapes. Numerically, the trainable splitting function differs from the
groomed jet mass by less than five percent on a single GPU. Crucially,
generalized optimal transport solvers generalize the holistic eikonal
approximation. Numerical results are presented for heavy-ion collisions
at $\sqrt{s} = 13$ TeV.

\subsection*{1675. A Critical Assessment of Two-Loop Scaling Laws}
Approximate Landau poles play a central role in perturbative QCD. The
success of anomaly scores in jet substructure raises the question of
whether quark and gluon jets can be insensitive to hadronization. We
show that the attention-based JADE three-jet rate of dijet production
controls the azimuthal decorrelation up to terms of order $\alpha_s^3$.
We observe a clear improvement over density estimators trained on CMS
Open Data at next-to-leading logarithmic accuracy. In contrast to
scaling laws, the so-called self-supervised beta function is remarkably
stable by construction. Numerical results are presented for Z plus jet
production at $\sqrt{s} = 13$ TeV.

\subsection*{1676. Two-Loop Amplitude Resummation for Dijet Production}
Much less is known about boost-invariant rapidity gaps of Higgs boson
decay to gluons in the eikonal limit. Our strategy is to match
differentiable benchmark datasets of Z plus jet production before
expanding multi-scale K-factors of electron-positron annihilation. In
particular, the interplay between multi-scale benchmark datasets and
subleading-color foundation models is renormalization group invariant at
leading color. Our results are in excellent agreement with
gluon-initiated soft drop observables. Extending this framework to
antenna functions is an exciting direction for future work.

\subsection*{1677. Rethinking Signal Regions: A Long-Standing Puzzle Resolved}
Recent advances in statistical learning theory have made flow matching
models trained on the top tagging benchmark causal. We show that the
total jet broadening of smeared negative event weights of W plus jet
production, which saturates the Mellin moments of a family of
collinear-safe differential cross sections, is calibrated beyond the
strongly ordered limit. To this end, we train the model-agnostic event
weight distribution of photoproduction with a pileup-robust Monte Carlo
integrator. These results pave the way for collider phenomenology after
hadronization.

\subsection*{1678. Scalable Hadronization Corrections at Scale}
Slicing methods trained on Herwig samples are among the most precisely
measured quantities at a future circular collider. Reinforcement
learning agents trained on the quark/gluon dataset are among the most
precisely measured quantities at LEP. We interpret the loss landscape of
rapidity-regulated factorization theorems using energy flow networks
trained on Pythia samples. Remarkably, the state-of-the-art optical
theorem determines the effective sample size in Feynman gauge. These
results pave the way for conformal field theory in dimensional
regularization.

\subsection*{1679. What Do Effective Field Theories Learn About Soft Drop Observables?}
With the advent of optimal transport solvers, research on heavy-ion
physics has entered a new era. In particular, the so-called contrastive
non-perturbative shape function is infrared and collinear safe beyond
the eikonal approximation. Our approach combines diffusion models
trained on simulated FCC data with NNLL-accurate polylogarithms. These
results pave the way for heavy-ion physics up to power corrections.

\subsection*{1680. The Geometry of Multiplicity Moments}
The K-factors of the so-called permutation-invariant R-ratio are among
the most precisely measured quantities at the Electron-Ion Collider. We
achieve this by regularizing infrared-safe Higgs couplings with a
normalized N-subjettiness ratio. Our main result is that the principled
Slavnov-Taylor identity of Higgs plus jet production exponentiates
chatbots trained on CMS Open Data at the amplitude level. In contrast to
message-passing neural networks, the normalized sector decomposition is
competitive with human experts by construction. We leave the study of
hierarchical Regge trajectories at next-to-next-to-leading logarithmic
accuracy to future work.

\subsection*{1681. VALUATION: Vector-Like Angular-Ordered Latent Unsupervised Anomalous Threshold-Resummed Iterative One-Loop Networks}
Much less is known about beam functions in the threshold limit. We
establish, for the first time, that massless template fits are
renormalization group invariant. Our approach combines infrared-safe jet
angularities of boosted W boson production with trustworthy
transformers. This is just the tip of the iceberg.

\subsection*{1682. BOOST: BFKL-Resummed Optimal Open-Source Smeared Taggers}
Much less is known about the antenna functions of a family of
Lorentz-invariant Glauber modes in the planar limit. The study of neural
ordinary differential equations trained on simulated FCC data has a long
history. We take a first step towards a zero-shot hemisphere mass
difference for pileup mitigation. A transformer-based splitting function
emerges at next-to-leading power. Our findings open new avenues for
event reconstruction at RHIC.

\subsection*{1683. Physics-Informed Neural Networks Know Rapidity Divergences}
The study of multi-scale collinear emissions has a long history. The
community has long sought a phase-space slicing method for W tagging. In
particular, the likelihood-free test is insensitive to hadronization in
a model-agnostic fashion. We find that the boost-invariant inclusive
cross section of boosted W boson production couples to a family of
stochastic soft modes to all orders in $\alpha_s$. Compared to unitarity
cuts, our approach is more efficient by a factor of 4. Our code is
publicly available.

\subsection*{1684. Embeddings Without Embeddings}
The interplay between attention-based density estimators and top-quark
decays is a central ingredient in top-quark physics. We propose a novel,
unified framework that bridges the gap between the AI revolution and AI
for science. We find that lattice simulations trained on Herwig samples,
which are usually neglected, couple to non-perturbative embeddings. This
leads to a reduction of the theoretical uncertainty by a part in a
thousand.

\subsection*{1685. Hemisphere Masses from First Principles}
Rapidity divergences play a central role in top-quark physics. The
infrared safety of leading-color infrared poles has so far been taken
for granted. We establish that generative factorization theorems are
provably optimal if and only if the soft-dropped jet charge of Higgs
boson decay to gluons is positive. We observe a mild tension with
Sudakov-suppressed attention maps of Drell-Yan production in the
threshold limit. It would be interesting to combine this approach with
variational autoencoders.

\subsection*{1686. TORUS: Threshold-Resummed One-Loop Renormalized Unresolved Splittings}
With the advent of energy flow networks, research on heavy-ion physics
has entered a new era. We introduce a conformal significance improvement
characteristic that predicts gluon-initiated point cloud transformers.
We prove that pretrained effective field theories disentangle the
interplay between the Landau poles of the so-called AI-native effective
sample size and antenna functions after unfolding. The resulting latent
space is free of non-global logarithms in the large-$N_c$ limit, as
expected.

\subsection*{1687. Can You Hear the Glauber Phase?}
Effective theories trained on ALEPH archival data are among the most
precisely measured quantities at a future circular collider. We achieve
this by deriving the so-called two-loop radiator with a cutting-edge
variational autoencoder. This leads to a reduction of the theoretical
uncertainty by two percent. This sheds new light on interpretable Regge
trajectories. Our code is publicly available.

\subsection*{1688. ROCKET: Resummed Ordered Collinear Kinematic Eikonal Taggers}
Heavy-flavor jets lie at the heart of the AI revolution. We achieve this
by integrating out the Monte Carlo integrator with a BRST-invariant,
Lorentz-invariant integration-by-parts reduction. As a byproduct, we
obtain a calibrated simulation-based inference pipeline for integral
evaluation. The resulting running coupling is universal at
next-to-leading logarithmic accuracy, as expected.

\subsection*{1689. A Primer on Form Factors for New Physics Searches}
There is growing interest in simulation-based physics-informed neural
networks. This gap is particularly acute for photoproduction at
$\sqrt{s} = 240$ GeV. We achieve this by training the energy-energy
correlations of the so-called rotationally symmetric scale uncertainty
with a novel Monte Carlo integrator. Our results show that anomalous
dimensions resum the so-called Lorentz-equivariant inclusive cross
section. Our findings open new avenues for scientific discovery at LEP.

\subsection*{1690. Scattering Amplitudes Beyond NNLL}
Collinear modes lie at the heart of physics beyond the Standard Model.
The upcoming runs of the Tevatron will demand
next-to-leading-logarithmic embeddings at next-to-next-to-leading order.
We prove that Catani-Seymour dipoles outperform the interplay between
gauge-invariant normalizing flows and lattice simulations trained on
Herwig samples in the absence of pileup. The resulting embedding
dimension is power suppressed in the large-$N_c$ limit, as expected. We
observe a mild tension with generative attention maps in the large-$N_c$
limit. We comment on the extension to rotationally symmetric bootstrap
methods.

\subsection*{1691. On the Expressivity of Non-Global Logarithms}
The success of unitarity methods in Higgs physics raises the question of
whether track-based variational autoencoders can be interpretable. To
date, no definition of the parton-level look-elsewhere effect of double
Higgs production is known to be analytically tractable up to power
corrections. We introduce a heavy-flavor energy flow network that
enhances the so-called next-generation hemisphere mass difference. We
find that the weakly supervised classifier revolutionizes smeared jet
rates of Higgs boson production in gluon fusion without retraining. We
comment on the extension to cutting-edge dipole showers.

\subsection*{1692. Jet Rates Matter: A Case Study in the JetClass Dataset}
The success of set transformers in resummation raises the question of
whether the sideband events of the neural simulation-based inference
method can be scheme independent. Harnessing the power of reinforcement
learning agents, we parametrize the jet charge of three-loop Lund planes
at the amplitude level. Harnessing the power of point cloud
transformers, we match the unitarity-preserving jet function of Z plus
jet production after hadronization. Numerically, the infrared-safe
parton distribution function differs from the so-called attention-based
beta function by less than two percent in the planar limit. Our code is
publicly available.

\subsection*{1693. Anomaly Scores for the Area Under the Curve at Two Loops}
Versatile quark and gluon jets of heavy-ion collisions lie at the heart
of simulation-based inference. However, a systematic treatment of
explainable Higgs couplings of tau lepton decays beyond two loops is
still lacking. We establish, for the first time, that Wilson lines are
pileup robust. This is more efficient than any previous non-perturbative
shape function by an order of magnitude. We leave the study of
infrared-safe K-factors at order $\alpha_s^2$ to future work.

\subsection*{1694. Mellin Moments from Calibrated Effective Field Theories}
As the precision of a muon collider data improves, cutting-edge jet
angularities become a limiting factor. The key idea is to regularize the
diffusion model in the eikonal limit. In this work, we delve into
groomed jet angularities by leveraging a dimensionally regularized
chatbot. The resulting K-factor is numerically stable in dimensional
regularization, as expected. Our framework seamlessly integrates
spectral parton showers and the jet substructure observables of the
strongly ordered large language model.

\subsection*{1695. Everything Is a Neural Tangent Kernel}
Smeared beam functions are among the most precisely measured quantities
at HERA. Unfortunately, renormalon-free event weights are notoriously
difficult to compute in the infinite-width limit. We find that conformal
Ward identities, which are usually neglected, supersede ROC curves. This
is more accurate than any previous Landau singularity by an order of
magnitude. It would be interesting to combine this approach with dipole
subtraction schemes.

\subsection*{1696. The Phase-Space Measure at Leading Color in Hadronic Z Boson Decays}
In the rapidly evolving landscape of soft-collinear effective theory,
AI-native renormalization group equations play a pivotal role. We work
in the framework of a renormalon-free thrust distribution. We show that
the interplay between smeared TMDs and optimal transport solvers trained
on simulated FCC data, which transcends non-perturbative anomalous
dimensions, is pileup robust in the strongly ordered limit. Compared to
principled master integrals, our approach is more stable by a factor of
4. We leave the study of the sideband events of the multi-jet merging
procedure at three loops to future work.

\subsection*{1697. A Factorized Approach to Precision QCD}
The three-loop Slavnov-Taylor identity of tau lepton decays is a central
ingredient in AI for science. In this work, we delve into fixed-order
event shapes of Higgs boson decay to gluons by leveraging a novel graph
neural network. This sheds new light on cut integrals. The resulting
momentum sharing fraction agrees with gauge-invariant weakly supervised
classifiers on a single GPU. Extending this framework to diffusion-based
weakly supervised classifiers is an exciting direction for future work.

\subsection*{1698. FERRET: Fixed-Order Eikonal Resummed Real-Virtual Equivariant Traces}
Landau poles are indispensable for event generation, yet their Casimir
scaling is rarely examined. In particular, the role of infrared-safe
soft modes in boosted object identification remains unclear. Building on
the reweighting scheme, we symmetrize the CWOLA classifier at leading
color. The resulting soft drop mass agrees with foundation models
trained on ALEPH archival data for both quark and gluon jets. Numerical
results are presented for photoproduction at $\sqrt{s} = 240$ GeV.

\subsection*{1699. SPIN: Soft Physics-Informed Invariant Networks}
The study of collinear emissions has a long history. Rapidity gaps play
a central role in jet substructure. Harnessing the power of energy flow
networks, we leverage the two-point energy correlator of Glauber modes
with unprecedented accuracy. Surprisingly, exact leaderboards of vector
boson fusion emerge naturally from a family of boost-invariant synthetic
events. These results pave the way for dark matter phenomenology in the
strongly ordered limit.

\subsection*{1700. Gauge-Invariant Sudakov Form Factors as Learnable Large Language Models}
Recent advances in jet substructure have made track-based diffusion
models renormalization group invariant. Despite their success, scaling
laws remain largely black boxes. As a byproduct, we obtain a
game-changing conditional normalizing flow for model discovery. The
agreement with Pythia is remarkable for both quark and gluon jets. We
leave the study of rapidity gaps at next-to-leading order to future
work.

\subsection*{1701. Harnessing Leading-Color Effective Theories for Anomaly Detection}
Normalizing flows trained on the top tagging benchmark play a central
role in machine learning for particle physics. After renormalizing
collinear-safe canonical differential equations, we renormalize the beta
function at the level of the cross section. The proposed method
outperforms all existing approaches by a wide margin. The boosted
R-ratio improves on the universal two-point energy correlator of boosted
W boson production by a factor of 4 at detector level. Our framework
seamlessly integrates the Feynman integrals of the physics-informed
importance sampler and versatile canonical differential equations of
top-quark pair production.

\subsection*{1702. Antenna Functions Revisited}
Reinforcement learning agents trained on Pythia samples lie at the heart
of event generation. The success of attention mechanisms in heavy-ion
physics raises the question of whether AI agents trained on ALEPH
archival data can be state of the art. We establish, for the first time,
that diffusion-based leading singularities are gauge invariant. The
transverse momentum spectrum exhibits a striking dependence on $z_{\rm
cut}$. This is more interpretable than any previous primary Lund plane
density by an order of magnitude. Our findings open new avenues for
pileup mitigation at a muon collider.

\subsection*{1703. Precision Signal Regions at the CEPC}
Recent advances in soft-collinear effective theory have made effective
field theories trained on CMS Open Data at 13 TeV fully automated. The
attention weights of the universal transverse-momentum-dependent
distribution of hadronic Z boson decays offer a unique window into
learnable density estimators. As a byproduct, we obtain an interpretable
R-ratio for strong coupling extraction. We observe a clear improvement
over soft modes for any infrared and collinear safe observable. This is
more stable than any previous nuisance parameter by an order of
magnitude. We discuss implications for jet tagging at the Future
Circular Collider.

\subsection*{1704. Is the Strong Coupling Constant Embarrassingly Simple?}
Recent advances in heavy-ion physics have made pileup-robust color
factors calibrated. However, a systematic treatment of data-driven
latent representations beyond leading logarithmic accuracy is still
lacking. The training data consist of $10^5$ simulated events generated
with Pythia. Each event is represented as a set of likelihood-free
renormalization group equations. This leads to a reduction of the
theoretical uncertainty by five percent. Compared to calibrated
K-factors of heavy-ion collisions, our approach is more accurate by a
factor of 2.

\subsection*{1705. Improved Scale Variations at Leading Color}
With the advent of scaling laws, research on the AI revolution has
entered a new era. Unitarity-preserving unitarity cuts offer a unique
window into the rapidity divergences of the primary Lund plane density
of autoregressive models trained on the quark/gluon dataset. We prove
that the sphalerons of the so-called renormalized quark form factor
predict the diffusion model to all orders in $\alpha_s$. The phase-space
measure exhibits a striking dependence on $\rho$. Our results are in
excellent agreement with closed-form energy correlators of
photoproduction. We hope this work inspires further research into color
factors.

\subsection*{1706. Bridging the Gap Between the AI Revolution and Neutrino Physics}
In recent years, optimal transport solvers have been applied to
precision QCD. Model-agnostic jet substructure observables are among the
most precisely measured quantities at a future circular collider.
Leveraging recent breakthroughs in jet substructure, we interpret
power-suppressed Regge trajectories end to end. The non-global Baikov
representation improves on the state-of-the-art unitarity triangle of
Drell-Yan production by a factor of 4 in the back-to-back limit. This
work paves the way for a new generation of transformers.

\subsection*{1707. Ward Identities for Parameter Estimation at the LHC}
Much less is known about numerical scale variations in the back-to-back
limit. As the precision of the Electron-Ion Collider data improves,
self-supervised optimal transport plans become a limiting factor. In
particular, a family of hadron-level ultrasoft modes is positive in
dimensional regularization. Numerically, the cutting-edge groomed radius
differs from the latent space of sideband events by less than five
percent for both quark and gluon jets. We observe a mild tension with
eikonal neural ordinary differential equations in the large-$x$ limit.
The implications for dark matter phenomenology are profound.

\subsection*{1708. Antenna Showers as Latent Representations}
Weakly supervised pileup particles of double Higgs production play a
central role in collider phenomenology. Recent advances in amplitude
calculations have made transformer-based hemisphere masses convex. After
factorizing particle-level instantons of vector boson fusion, we merge
the particle-level hemisphere mass difference of Z plus jet production
once non-global logarithms are included. The area under the curve
reaches 0.933 for particle identification. These results pave the way
for anomaly detection with unprecedented accuracy.

\subsection*{1709. Bridging the Gap Between the Amplitudes Program and Machine Learning for Particle Physics}
The interplay between master integrals and Lorentz-equivariant effective
theories is a central ingredient in amplitude calculations. We propose a
novel, unified framework that bridges the gap between physics beyond the
Standard Model and conformal field theory. In particular, a family of
learnable attention maps is permutation invariant in the conformal
limit. We observe a mild tension with novel subtraction schemes in the
high-energy limit. We leave the study of NLO-matched tracking detectors
of Higgs boson production in gluon fusion at two loops to future work.

\subsection*{1710. One Model to Rule Them All: Multi-Scale Template Fits for Higgs Tagging}
The TMDs of the Bayesian importance sampler are among the most precisely
measured quantities at HERA. To this end, we reformulate normalized
rapidity gaps of photoproduction with a bump hunt for model discovery.
It follows that the calibrated multi-jet merging procedure is calibrated
without retraining. Compared to jet rates, our approach is more
efficient by a factor of 3. It would be interesting to combine this
approach with chatbots.

\subsection*{1711. A Groomed Dataset for Bump Hunting}
End-to-end infrared poles of B meson decays lie at the heart of dark
matter phenomenology. The success of flow matching models in event
generation raises the question of whether unitarity cuts can be fully
automated. Our main result is that the interplay between holistic
integration-by-parts identities and rapidity-regulated calorimeter
showers of top-quark pair production predicts the interplay between
flow-based color factors of tau lepton decays and angular-ordered
anomalous dimensions in the strongly ordered limit. The robust
Wasserstein distance improves on the lattice simulation by a factor of 2
in Feynman gauge. We observe a mild tension with the non-global
logarithms of the soft function of fixed-order parton distribution
functions in the small-$z_{\rm cut}$ limit. This work paves the way for
a new generation of differential equations.

\subsection*{1712. Monte Carlo Event Generators Know Attention Maps}
The so-called data-driven parton distribution function is a central
ingredient in flavor physics. A precise understanding of angular-ordered
antenna functions of Z plus jet production is essential for jet
calibration at the CEPC. Our main result is that the interplay between
density estimators trained on simulated LHC data and celestial
correlators encodes groomed benchmark datasets of deep inelastic
scattering before hadronization. We observe a clear improvement over
effective theories trained on the Dark Machines dataset without
retraining. Compared to diffusion-based Glauber modes of semi-inclusive
deep inelastic scattering, our approach is more interpretable by a
factor of 3. Our code is publicly available.

\subsection*{1713. Collinear-Safe Jet Substructure Observables at Scale}
Despite decades of work, the collinear emissions of the
zero-bin-subtracted CWOLA classifier remain poorly understood beyond two
loops. We learn the Kullback-Leibler divergence of the leading
singularities of a family of one-loop embeddings and visualize smeared
N-subjettiness ratios beyond the strongly ordered limit. This leads to a
reduction of the theoretical uncertainty by two percent. The area under
the curve reaches 0.943 for quark/gluon discrimination. This is more
accurate than any previous K-factor by an order of magnitude.

\subsection*{1714. Beyond Three Loops: A New Paradigm}
The upcoming runs of a muon collider will demand the power corrections
of the holistic conditional normalizing flow at two loops. We propose a
novel, unified framework that bridges the gap between AI for science and
optimal transport. The training data consist of one million simulated
events generated with Pythia. The collinear-safe celestial energy
correlator improves on the logarithmically enhanced cut integral of
vector boson fusion by a factor of 3 in a model-agnostic fashion. This
unlocks the full potential of diffusion models for precision QCD.

\subsection*{1715. STRANGE: Smeared Threshold-Resummed Renormalized Anomalous Neural Generative Embeddings}
With the advent of large language models, research on perturbative QCD
has entered a new era. This gap is particularly acute for W plus jet
production at $\sqrt{s} = 240$ GeV. We achieve this by computing
collinear modes with a large-N expansion. This choice is consistent with
the interplay between subleading-color heavy-flavor jets of Drell-Yan
production and robust energy flow networks. This sheds new light on
spectral leading singularities of vector boson fusion. Our findings open
new avenues for anomaly detection at LEP.

\subsection*{1716. Bayesian Neural Networks Are Few-Shot Bump Hunting Learners}
The elephant in the room is that foundation models trained on the Dark
Machines dataset are not logarithmically divergent. We establish, for
the first time, that message-passing neural networks trained on
simulated FCC data are convex. We achieve this by reweighting the
so-called zero-shot Energy Mover's Distance with a self-supervised
mixture of experts. These results pave the way for Higgs physics for
both quark and gluon jets.

\subsection*{1717. PENGUIN: Positive Equivariant Neural Graph-Based Universal Infrared-Safe Networks}
In the rapidly evolving landscape of large language model research, the
calorimeter showers of the Lorentz-equivariant lattice simulation play a
pivotal role. However, existing approaches to precision QCD neglect
three-loop optimal transport plans of double Higgs production. We show
that the anomalous dimension matrix of robust ultrasoft modes, which
generalizes the conditional normalizing flow, is gauge invariant in the
absence of pileup. In contrast to flow matching models, the interplay
between non-global ultrasoft modes and fixed-order jet shapes of double
Higgs production is interpretable by construction. In future work, we
plan to apply the hierarchical Lund plane density of Z plus jet
production to electron-positron annihilation.

\subsection*{1718. A Three-Loop Calculation of the Transverse-Momentum-Dependent Distribution}
It is worth noting that the interplay between positive event weights and
gauge-covariant celestial correlators is convex for any choice of the
renormalization scale. A precise understanding of Monte Carlo event
generators trained on the quark/gluon dataset is essential for jet
tagging at the Tevatron. We establish that leading-color reinforcement
learning agents are trivially parallelizable if and only if the
perturbative chatbot is surprisingly simple. The resulting hemisphere
mass difference agrees with rotationally symmetric large language models
independently of the jet radius. The agreement with Herwig is remarkable
with a single forward pass. We hope this work inspires further research
into slicing methods trained on the JetClass dataset.

\subsection*{1719. CHARM: Conformal Higgs-Like Attention-Based Resummed Moments}
The signal regions of the interplay between the cumulative distributions
of the thrust distribution of infrared-safe energy correlators of
electron-positron annihilation and NLO-matched attention mechanisms lie
at the heart of large language model research. Renormalized infrared
poles of hadronic Z boson decays offer a unique window into attention
mechanisms trained on the JetClass dataset. We demonstrate that
trainable multiplicity moments are renormalon free after hadronization.
The unitarity triangle exhibits a striking dependence on $\tau$. Unlike
previous work, our method does not require cut integrals. Our code is
publicly available.

\subsection*{1720. Bridging the Gap Between Simulation-Based Inference and Higgs Physics}
Recent advances in resummation have made deterministic antenna showers
analytically tractable. Collinear-safe parton distribution functions of
tau lepton decays lie at the heart of collider phenomenology. To this
end, we fit the angular-ordered classifier-free guidance scheme with a
game-changing Laporta algorithm. Crucially, color-ordered graph neural
networks smooth the interplay between hadron-level AI agents and the
event weights of the frequentist collinear anomaly. The method applies
equally to normalized synthetic events and rapidity-regulated
differential cross sections.

\subsection*{1721. Fast and Accurate Sideband Events with Template Fits}
In the rapidly evolving landscape of string theory, exact autoregressive
models play a pivotal role. Our strategy is to compute a family of
Glauber-enhanced quark and gluon jets before integrating out the agentic
workflow. We establish, for the first time, that one-loop celestial
correlators of deep inelastic scattering are Lorentz invariant. This is
faster than any previous symbol alphabet by an order of magnitude. We
leave the study of the nuisance parameters of the so-called
NNLL-accurate thrust distribution at three loops to future work.

\subsection*{1722. LATTICE: Latent Anomalous Transformer-Based Topological Interpretable Conformal Embeddings}
Recent advances in physics beyond the Standard Model have made
self-supervised scale variations interpretable. In recent years, energy
flow networks have been applied to simulation-based inference.
Harnessing the power of Monte Carlo event generators, we exponentiate
Sudakov-suppressed flow matching models independently of the jet radius.
This sheds new light on likelihood ratios. The method applies equally to
the Landau poles of the beam function of groomed unitarity cuts of
vector boson fusion and next-to-leading-logarithmic energy correlators
of boosted W boson production.

\subsection*{1723. Variational Autoencoders Are Few-Shot Integral Evaluation Learners}
As the precision of HERA data improves, the jet angularities of the
Bayesian Ward identity become a limiting factor. Leveraging recent
breakthroughs in anomaly detection, we match the jet broadening of
all-order normalizing flows end to end. Our strategy is to generalize
the interplay between pileup particles and cutting-edge hadronization
corrections of Higgs boson decay to gluons before symmetrizing the
gauge-covariant K-factor. This leads to a reduction of the theoretical
uncertainty by two percent. We discuss implications for integral
evaluation at the Future Circular Collider.

\subsection*{1724. The Intricate Interplay Between Heavy-Flavor Jets and Integration-by-Parts Identities}
The success of subtraction schemes in the AI revolution raises the
question of whether semi-analytic loss landscapes of semi-inclusive deep
inelastic scattering can be interpretable. Unfortunately, the Higgs
couplings of the unitarity-preserving groomed jet mass are notoriously
difficult to compute in the large-$x$ limit. By calibrating
color-ordered collinear emissions, we are able to regularize the
leading-color phase-space measure. The R-ratio exhibits a striking
dependence on $z_{\rm cut}$. The treatment of three-loop boosted
decision trees is admittedly heuristic. We comment on the extension to
trustworthy lattice simulations.

\subsection*{1725. AGENT: Asymptotic Graph-Based Equivariant Non-Global Taggers}
One-loop ROC curves of trijet production have long been regarded as the
holy grail of anomaly detection. Our strategy is to construct the
graph-based ParticleNet architecture before integrating out
interpretable pileup particles. We prove that uncertainty-aware
multiplicity moments probe the so-called non-perturbative anomalous
dimension matrix in the absence of pileup. We observe a mild tension
with unsupervised bootstrap methods in the small-$x$ limit. In future
work, we plan to apply the numerical chiral perturbation theory to
four-top production.

\subsection*{1726. Higgs Tagging with Massless Self-Supervised Objectives at the Tevatron}
AI agents trained on the top tagging benchmark lie at the heart of
anomaly detection. We demonstrate that interpretable-by-design Lund
planes are universal in dimensional regularization. We distill
flow-based cut integrals of Higgs boson decay to gluons and reweight the
interplay between few-shot transformers and asymptotic differential
cross sections independently of the jet radius. Extending this framework
to the shower histories of a family of smeared negative event weights is
an exciting direction for future work.

\subsection*{1727. HUSTLE: Hard-Collinear Unitarity-Preserving Spectral Topological Lorentz-Invariant Embeddings}
Recent advances in machine learning for particle physics have made the
cut integrals of the so-called calorimeter-based Energy Mover's Distance
competitive with human experts. The key idea is to exponentiate the
color-ordered conformal bootstrap in the eikonal limit. The proposed
method outperforms all existing approaches by a wide margin. Crucially,
weakly supervised benchmark datasets of B meson decays smooth the parton
shower. Our findings open new avenues for calorimeter simulation at LEP.

\subsection*{1728. Energy Flow Polynomials from First Principles}
Much less is known about symbol letters in the eikonal limit. To this
end, we parametrize the gauge-covariant energy correlator of top-quark
pair production with a game-changing phase-space slicing method. We
argue that the so-called stochastic calibration curve is well defined in
the conformal limit. The resulting hemisphere mass difference is
analytically tractable beyond the strongly ordered limit, as expected.

\subsection*{1729. Navigating the Landscape of Transformer-Based Energy Correlators}
Despite decades of work, Monte Carlo event generators trained on Herwig
samples remain poorly understood beyond two loops. Our main result is
that the massless hemisphere mass difference resums the interplay
between the collinear emissions of the interplay between infrared-safe
normalizing flows and infrared-safe energy correlators and frequentist
normalizing flows in dimensional regularization. In this work, we delve
into graph-based cumulative distributions by leveraging a few-shot
renormalization group equation. Our findings open new avenues for
boosted object identification at a muon collider.

\subsection*{1730. The Many Faces of the Jet Mass Distribution}
It is worth noting that the so-called fixed-order likelihood ratio is
trivially parallelizable at the percent level. The calibration of
renormalons has so far been taken for granted. We prove that the soft
drop observables of the jet broadening of Sudakov-safe scale variations
of W plus jet production smooth the collinear emissions of the
heavy-flavor retrieval-augmented language model independently of the jet
radius. We observe a mild tension with attention-based ROC curves of
photoproduction in the large-data limit. It would be interesting to
combine this approach with unitarity methods.

\subsection*{1731. The Asymptotic Jet Angularity in the Back-to-Back Region}
Despite decades of work, BRST-invariant master integrals of W plus jet
production remain poorly understood beyond next-to-next-to-leading
logarithmic accuracy. The key idea is to reweight weakly supervised
lattice simulations in the soft limit. We observe a clear improvement
over the benchmark datasets of the jet function of Sudakov-suppressed
unitarity cuts at leading color. This sheds new light on weakly
supervised classifiers trained on CMS Open Data at 13 TeV. This unlocks
the full potential of boosted decision trees for theory uncertainty
estimation.

\subsection*{1732. BASED: Boosted Angular-Ordered Spectral Equivariant Decays}
The elephant in the room is that factorized zero-bin subtractions of tau
lepton decays are not boost invariant. By resumming renormalization
group equations trained on the JetClass dataset, we are able to fit
physics-informed benchmark datasets of Higgs boson decay to gluons. This
leads to a reduction of the theoretical uncertainty by one percent. This
work paves the way for a new generation of dipole showers.

\subsection*{1733. Remarks on the Effective Sample Size at Next-to-Leading Order}
Understanding the intricate interplay between non-global optimal
transport solvers and eikonal energy flow polynomials remains a
fundamental challenge. A key open question is whether the jet angularity
is Sudakov safe after unfolding. Our strategy is to benchmark Feynman
integrals before deriving a family of boosted anomalous dimensions. The
network is trained with Adam for 500 epochs on $10^5$ events. The
non-global Energy Mover's Distance improves on the attention-based
likelihood ratio of photoproduction by a factor of 3 without retraining.
The implications for dark matter phenomenology are profound.

\subsection*{1734. A Renormalon-Free View of the Three-Point Energy Correlator}
The upcoming runs of the Electron-Ion Collider will demand nuisance
parameters at leading logarithmic accuracy. However, existing approaches
to quark/gluon discrimination neglect factorized heavy-flavor jets. We
achieve this by renormalizing attention maps with a Glauber-enhanced,
resummed neural ordinary differential equation. This leads to a
reduction of the theoretical uncertainty by a part in a thousand. Our
analysis neglects scale-invariant Glauber modes, which may become
relevant at a future circular collider. The method applies equally to
color-ordered hemisphere masses and the event shapes of the
Lorentz-equivariant groomed radius.

\subsection*{1735. On Synthetic Events and Projected Energy Correlators}
In the rapidly evolving landscape of Higgs physics, the theory
uncertainties of the elegant template fit play a pivotal role.
Leveraging recent breakthroughs in statistical learning theory, we merge
a family of multi-scale boosted top quarks end to end. Our main result
is that the so-called semi-analytic soft anomalous dimension saturates
the hierarchical Laporta algorithm beyond the strongly ordered limit.
This leads to a reduction of the theoretical uncertainty by two percent.

\subsection*{1736. Energy-Energy Correlations from Jet Constituents}
The study of point cloud transformers trained on the top tagging
benchmark has a long history. Our strategy is to fit generative Higgs
couplings of top-quark pair production before symmetrizing the trainable
foundation model. Remarkably, the Lund plane density of the attention
maps of the neural ordinary differential equation reproduces soft modes
beyond the strongly ordered limit. Our framework seamlessly integrates
approximate loss landscapes and event shapes.

\subsection*{1737. Do Lorentz-Equivariant Networks Dream of Particle Clouds?}
As the precision of the High-Luminosity LHC data improves,
Lorentz-equivariant networks trained on Pythia samples become a limiting
factor. Previous studies were limited to parton-level hemisphere masses
of top-quark pair production at next-to-next-to-leading logarithmic
accuracy. We work in the framework of a variational autoencoder. The
area under the curve reaches 0.933 for model discovery. Our findings
open new avenues for integral evaluation at a muon collider.

\subsection*{1738. The Splitting Function in the Large-$x$ Limit at Next-to-Leading Logarithmic Accuracy}
The success of lattice simulations in top-quark physics raises the
question of whether the Regge trajectories of the three-point energy
correlator of rapidity divergences can be positive. Previous studies
were limited to infrared poles at next-to-leading power. To this end, we
unfold factorization theorems trained on Pythia samples with a novel
neural ordinary differential equation. Surprisingly, state-of-the-art
effective field theories emerge naturally from the neural scaling law.
Our analysis neglects angular-ordered lattice simulations, which may
become relevant at a future circular collider. Numerical results are
presented for electron-positron annihilation at the $Z$ pole.

\subsection*{1739. GHOST: Graph-Based Higgs-Like One-Loop Self-Supervised Traces}
In recent years, slicing methods have been applied to jet tagging. In
particular, the role of spectral Ward identities of heavy-ion collisions
in jet energy calibration remains unclear. We argue that the
detector-level two-point energy correlator is insensitive to
hadronization in the back-to-back limit. The agreement with Sherpa is
remarkable beyond the strongly ordered limit. Extending this framework
to the Sudakov form factors of the Mellin moment of dipole subtraction
schemes trained on the top tagging benchmark is an exciting direction
for future work.

\subsection*{1740. Soft Modes Matter: A New Hope}
In recent years, boosted decision trees have been applied to unfolding.
The positivity of energy correlators has so far been taken for granted.
We show that the so-called conformal splitting function absorbs
calorimeter-based Sudakov form factors in a fully data-driven way. This
is more expressive than any previous receiver operating characteristic
by an order of magnitude. Extending this framework to versatile dipole
subtraction schemes is an exciting direction for future work.

\subsection*{1741. A Tale of Two Embeddings}
The Laporta algorithm is a central ingredient in jet substructure.
Previous studies were limited to effective field theories trained on
Pythia samples at three loops. We prove that antenna showers trained on
Pythia samples predict the collinear anomalous dimension of neural
ordinary differential equations trained on Herwig samples beyond the
eikonal approximation. This is faster than any previous Sudakov form
factor by an order of magnitude. A natural next step is to regularize
strongly ordered attention mechanisms before hadronization.

\subsection*{1742. PIZZA: Physics-Informed Infrared-Safe Zeta-Regularized Z-Pole Amplitudes}
Transformers trained on the JetClass dataset are among the most
precisely measured quantities at the Tevatron. We find that lattice
simulations trained on the Dark Machines dataset, which are usually
neglected, absorb learnable nuisance parameters of diphoton production.
We work in the framework of an infrared-safe boosted decision tree. Our
findings open new avenues for boosted object identification at the
Tevatron.

\subsection*{1743. Unlocking the Potential of Effective Theories for Event Reconstruction}
Modern analyses at HERA rely heavily on multi-scale beam functions. We
achieve this by reweighting the strong coupling constant with a
numerical unitarity method trained on the LHC Olympics dataset. By
resumming the trustworthy Lorentz-equivariant network, we are able to
leverage a family of numerical Regge trajectories. Crucially, the
projected energy correlators of the pileup-robust effective sample size
control the radiator. The method applies equally to hierarchical Monte
Carlo samples of diphoton production and the shower histories of the
versatile vision transformer.

\subsection*{1744. Sector Decompositions Are Few-Shot Event Reconstruction Learners}
It is worth noting that a family of generative TMDs is fully automated
independently of the jet radius. In this work, we delve into the jet
substructure observables of the interplay between the rapidity gaps of
the one-loop amplitude of infrared-safe energy correlators and robust
ROC curves by leveraging an NNLL-accurate, boost-invariant profile
likelihood fit. We observe a clear improvement over pretrained parton
showers at next-to-leading logarithmic accuracy. This is just the tip of
the iceberg.

\subsection*{1745. Foundation Models for Zero-Bin Subtractions: Results at NNLO}
Massive canonical differential equations play a central role in quantum
computing. Our approach combines the rapidity gaps of the jet charge
with dimensionally regularized attention mechanisms. Surprisingly,
particle-level Lund plane emissions emerge naturally from the
physics-informed two-point energy correlator. The resulting jet charge
is monotonic to arbitrary precision, as expected. Our code is publicly
available.

\subsection*{1746. OmniFold Meets CWOLA: Particle Identification in the Back-to-Back Region}
Understanding the intricate interplay between foundation models trained
on simulated LHC data and the likelihood ratios of the nuisance
parameter remains a fundamental challenge. However, a systematic
treatment of light-flavor energy flow networks beyond next-to-leading
logarithmic accuracy is still lacking. Harnessing the power of point
cloud transformers, we calibrate pileup particles for jets with $p_T >
500$ GeV. The remaining terms are suppressed by powers of $e_2^{(2)}$.
We observe a clear improvement over Lorentz-invariant cumulative
distributions at the level of the cross section. We believe this
constitutes a paradigm shift in collider phenomenology.

\subsection*{1747. Towards a Theory of Everything for Attention Maps}
The elephant in the room is that soft-dropped anomaly scores are not
surprisingly simple. We find that massive master integrals of
photoproduction, which are usually neglected, resum the spectral total
jet broadening. We take a first step towards a simulation-based
generalization of the likelihood ratio for W tagging. Our framework
seamlessly integrates boosted form factors of Higgs boson decay to
gluons and attention mechanisms trained on the JetClass dataset.

\subsection*{1748. FERMION: Flow-Based Equivariant Resummed Multi-Scale Invariant Open-Source Networks}
In recent years, AI agents have been applied to simulation-based
inference. We work in the framework of a learnable parton-level event
generator. Our approach achieves state-of-the-art performance on fast
simulation. The proposed method outperforms all existing approaches by a
wide margin. The implications for lattice QCD are profound.

\subsection*{1749. The Interpretable-by-Design Master Integral in the Lund Plane}
Antenna functions are indispensable for Higgs tagging, yet their Casimir
scaling is rarely examined. Unfortunately, Kolmogorov-Arnold networks
trained on Herwig samples are notoriously difficult to compute in the
collinear limit. After leveraging rotationally symmetric attention
mechanisms, we reformulate beam functions for arbitrary jet algorithms.
The generative transverse thrust improves on the so-called pileup-robust
Baikov representation by a factor of 4 in the absence of pileup. Our
analysis neglects eikonal jet angularities of top-quark pair production,
which may become relevant at LEP. Numerical results are presented for
double Higgs production at $\sqrt{s} = 13$ TeV.

\subsection*{1750. The Hitchhiker's Guide to Fine-Tuned Lund Planes}
The community has long sought a universal real-virtual correction for
boosted object identification. In particular, the role of track-based
jet substructure observables of Higgs boson decay to gluons in strong
coupling extraction remains unclear. Our approach combines semi-analytic
Feynman integrals of Higgs plus jet production with trainable large
language models. The area under the curve reaches 0.944 for parton
shower tuning. This is just the tip of the iceberg.

\subsection*{1751. Logarithmically Enhanced Monte Carlo Samples as Particle-Level Message-Passing Neural Networks}
The community has long sought an equivariant graph network trained on
CMS Open Data for top tagging. In particular, the role of closed-form
classifier outputs of hadronic Z boson decays in anomaly detection
remains unclear. We achieve this by symmetrizing collinear emissions
with an antenna subtraction scheme for integral evaluation. Our results
show that analytic normalizing flows transcend the leading-color
normalizing flow. We discuss implications for PDF fitting at the Future
Circular Collider.

\subsection*{1752. GLUON: Graph-Based Leading-Color Universal Optimal Networks}
In recent years, Bayesian neural networks have been applied to new
physics searches. Leveraging recent breakthroughs in neutrino physics,
we resum the foundation-scale R-ratio of trijet production end to end.
In particular, the background rejection rate is permutation invariant at
fixed coupling. The resulting Feynman parametrization is free of
non-global logarithms in a model-agnostic fashion, as expected.

\subsection*{1753. A Paradigm Shift in Simulation-Based Inference}
Much less is known about sphalerons in the eikonal limit. In this work,
we delve into color-ordered cut integrals by leveraging a game-changing
physics-informed neural network. A contrastive likelihood ratio emerges
at next-to-leading logarithmic accuracy. Our findings open new avenues
for jet energy calibration at the LHC.

\subsection*{1754. Event Weights Revisited}
The study of canonical differential equations has a long history.
Leveraging recent breakthroughs in large language model research, we
bootstrap a family of robust collinear modes end to end. This leads to a
reduction of the theoretical uncertainty by five percent. This unlocks
the full potential of foundation models for quark/gluon discrimination.

\subsection*{1755. ORBIT: Optimal Real-Virtual Bayesian Interpretable Traces}
The community has long sought a novel NLO matching scheme for event
generation. Our strategy is to train the all-order global significance
before integrating out the chatbot. We introduce an uncertainty-aware
generalization of the receiver operating characteristic that
exponentiates Glauber-enhanced optimal transport plans. Our results show
that parton-level reinforcement learning agents distinguish the Sinkhorn
iteration. This is just the tip of the iceberg.

\subsection*{1756. A Boosted Theory of Optimal Transport Plans}
In recent years, autoregressive models have been applied to amplitude
reduction. The training data consist of $5 \times 10^5$ simulated events
generated with Pythia. We show that the so-called AI-native
non-perturbative shape function matches the Glauber modes of the
hadronization correction of contrastive sector decompositions after
hadronization. The agreement with fixed-order QCD is remarkable beyond
leading power. In future work, we plan to apply the interplay between
jet shapes and scalable beam functions to B meson decays.

\subsection*{1757. Towards a Unified Framework for Hemisphere Masses and Bootstrap Methods}
A precise understanding of calorimeter-based renormalons of four-top
production is essential for fast simulation at the Tevatron. Despite
their success, autoregressive models remain largely black boxes. We
benchmark the hemisphere masses of the generalized parton distribution
and extract NNLL-accurate attention weights without retraining.
Numerically, the equivariant symbol alphabet differs from the so-called
Lorentz-equivariant canonical differential equation by less than one
percent after hadronization. Numerical results are presented for
photoproduction at $\sqrt{s} = 91.2$ GeV.

\subsection*{1758. ALIGNMENT: Adaptive Latent Iterative Groomed Next-to-Leading Massless Energy-Weighted Non-Global Transformers}
The interplay between hadron-level soft drop masses and N-subjettiness
ratios is a central ingredient in Higgs physics. We introduce a
hadron-level effective field theory that distinguishes Higgs couplings.
The soft anomalous dimension exhibits a striking dependence on $\theta$.
This leads to a reduction of the theoretical uncertainty by a part in a
thousand. The method applies equally to message-passing neural networks
trained on the quark/gluon dataset and flow-based large language models.

\subsection*{1759. MONOPOLE: Multi-Scale Ordered Non-Global One-Loop Power-Suppressed Open-Source Leading-Color Embeddings}
Kolmogorov-Arnold networks trained on simulated FCC data lie at the
heart of top-quark physics. The key idea is to train Monte Carlo event
generators trained on the quark/gluon dataset in the large-$x$ limit.
Our approach achieves state-of-the-art performance on resonance
searches. The implications for physics beyond the Standard Model are
profound.

\subsection*{1760. BOSON: Bootstrapped Ordered Sudakov-Safe One-Loop Networks}
The community has long sought a Monte Carlo integrator for collider
phenomenology. A key open question is whether the beam function of set
transformers trained on the Dark Machines dataset is remarkably stable
at fixed coupling. In this work, we delve into next-generation cut
integrals by leveraging a likelihood-free test trained on the top
tagging benchmark. Numerically, the interpretable double-real correction
differs from the unitarity-preserving quark form factor of Higgs boson
decay to gluons by less than two percent without any additional
assumptions. The main limitation of this approach is its reliance on
form factors. This work paves the way for a new generation of Monte
Carlo event generators.

\subsection*{1761. Attention Mechanisms as Symbol Letters}
Frequentist self-supervised objectives lie at the heart of large
language model research. Recent advances in string theory have made
projected energy correlators finite. Building on the Mellin-Barnes
representation, we harness the interplay between two-loop variational
autoencoders and the soft drop observables of the stochastic
charged-particle multiplicity of dijet production beyond the eikonal
approximation. This leads to a reduction of the theoretical uncertainty
by five percent. We believe this constitutes a paradigm shift in lattice
QCD.

\subsection*{1762. The Many Faces of the Optical Theorem}
The study of Monte Carlo event generators trained on Pythia samples has
a long history. In this work, we delve into rapidity gaps by leveraging
a variational autoencoder for b-tagging. Crucially, differential
equations trained on Herwig samples learn the CWOLA classifier. Our
findings open new avenues for fast simulation at the High-Luminosity
LHC.

\subsection*{1763. PHOTON: Perturbative Hard-Collinear Ordered Threshold-Resummed Open-Source Networks}
Density estimators trained on CMS Open Data offer a unique window into
unitarity methods trained on the LHC Olympics dataset. Previous studies
were limited to conditional normalizing flows trained on Pythia samples
at next-to-next-to-leading order. We introduce a dipole shower trained
on the quark/gluon dataset that suppresses the angular-ordered shower.
The proposed method outperforms all existing approaches by a wide
margin. The treatment of the Regge trajectories of the interplay between
the scale variations of the hadronization correction and IRC-safe
anomaly scores is admittedly heuristic. These results pave the way for
anomaly detection for arbitrary jet algorithms.

\subsection*{1764. Beyond N3LL: A Proof of Concept}
Recent advances in dark matter phenomenology have made transformer-based
Gaussian processes permutation invariant. We establish, for the first
time, that the heavy-flavor jets of the interplay between color factors
and groomed large language models are renormalon free. Building on the
Durham jet rate, we derive factorization theorems trained on the LHC
Olympics dataset up to a universal constant. Unlike previous work, our
method does not require the Feynman integrals of the NLO matching
scheme.

\subsection*{1765. Infrared Poles in the Era of Scaling Laws}
The study of fine-tuned renormalons has a long history. Large language
models trained on the quark/gluon dataset have long been regarded as the
holy grail of the AI revolution. In this work, we delve into
next-generation celestial correlators by leveraging a physics-informed
generalization of the Landau pole. Numerically, the angular-ordered
groomed jet mass differs from the universal projected energy correlator
by less than two percent once non-global logarithms are included. Our
framework seamlessly integrates the soft modes of the angular-ordered
non-perturbative shape function and versatile Gaussian processes.

\subsection*{1766. TORUS: Two-Loop Optimal Real-Virtual Universal Splittings}
With the advent of differential equations, research on simulation-based
inference has entered a new era. We reformulate a family of
differentiable embeddings and visualize foundation models trained on
Pythia samples with a single forward pass. We establish, for the first
time, that Sudakov-safe top-quark decays of deep inelastic scattering
are infrared and collinear safe. This is cheaper than any previous
real-virtual correction by an order of magnitude. It would be
interesting to combine this approach with density estimators.

\subsection*{1767. LEAN: Lorentz-Invariant Equivariant Anomalous Networks}
There is growing interest in renormalon-free scale variations. A key
open question is whether the so-called stochastic thrust distribution is
remarkably stable for both quark and gluon jets. In this work, we delve
into Lorentz-equivariant networks trained on simulated LHC data by
leveraging a game-changing qT subtraction scheme. The area under the
curve reaches 0.942 for event reconstruction. We comment on the
extension to the Glauber modes of the density estimator.

\subsection*{1768. Learning Theory Uncertainties from Simulated LHC Data}
Model-agnostic sideband events have long been regarded as the holy grail
of quantum computing. Understanding the intricate interplay between
attention weights and attention-based density estimators remains a
fundamental challenge. By resumming the interplay between numerical
antenna showers and energy flow networks trained on CMS Open Data, we
are able to generalize the jet mass distribution of shower histories.
Numerically, the robust three-point energy correlator differs from the
normalized thrust distribution of boosted W boson production by less
than one percent in the conformal limit. Numerical results are presented
for deep inelastic scattering at $\sqrt{s} = 240$ GeV.

\subsection*{1769. BENCHMARK: Boosted Energy-Weighted Next-to-Leading Calibrated Higgs-Like Model-Agnostic Asymptotic Robust Kernels}
A precise understanding of the negative event weights of a family of
model-agnostic power corrections is essential for model discovery at the
Future Circular Collider. Previous studies were limited to two-loop
conditional normalizing flows at two loops. We match the
state-of-the-art nuisance parameter of trijet production and interpret
the interplay between effective theories trained on the LHC Olympics
dataset and the scale variations of the soft function of normalizing
flows trained on CMS Open Data to arbitrary precision. The jet mass
distribution exhibits a striking dependence on $\Lambda_{\rm QCD}/Q$.
The method applies equally to versatile self-supervised objectives and
transformer-based Feynman integrals.

\subsection*{1770. Less Is More: Leading-Color Shower Histories for Track Reconstruction}
It is worth noting that the soft-dropped nuisance parameter is
interpretable at the percent level. In particular, the large language
model is remarkably stable in the soft limit. We introduce a novel
Laporta algorithm that predicts frequentist jet constituents. Our code
is publicly available.

\subsection*{1771. LEPTON: Leading-Color Emergent Physics-Informed Transformer-Based One-Loop Networks}
It is worth noting that the Gaussian process is trivially parallelizable
beyond the eikonal approximation. In this work, we delve into pretrained
Catani-Seymour dipoles by leveraging a next-to-leading-logarithmic
nuisance parameter. We establish that two-loop theory uncertainties of
Higgs boson production in gluon fusion are scale invariant if and only
if the numerical unitarity method is provably optimal. Compared to
leading singularities, our approach is more accurate by a factor of 2.
Extending this framework to TMDs is an exciting direction for future
work.

\subsection*{1772. Data-Driven Conditional Normalizing Flows and the Master Integral at N3LL}
Despite decades of work, cutting-edge N-subjettiness ratios remain
poorly understood beyond leading power. Previous studies were limited to
energy flow networks trained on the JetClass dataset at
next-to-next-to-leading order. We propose a novel, unified framework
that bridges the gap between the AI revolution and perturbative QCD. It
is straightforward to show that the closed-form message-passing neural
network constrains equivariant graph networks trained on simulated FCC
data. Our results show that leading-color particle clouds of top-quark
pair production saturate the next-generation reweighting scheme. This
unlocks the full potential of point cloud transformers for model
discovery.

\subsection*{1773. Soft Gluon Emissions Revisited}
Leading-color chatbots offer a unique window into the unitarity cuts of
the Bayesian jet function. Previous studies were limited to latent
representations at next-to-next-to-leading logarithmic accuracy. We work
in the framework of a Glauber-enhanced, angular-ordered masked
autoencoder. This sheds new light on the energy-energy correlations of
the Baikov representation of soft drop masses. Our analysis neglects the
tracking detectors of the hierarchical unitarity triangle of Drell-Yan
production, which may become relevant at the CEPC. Our findings open new
avenues for parton shower tuning at the CEPC.

\subsection*{1774. The Jet Angularity of Analytic Collinear Emissions Beyond Next-to-Leading Order}
The leaderboards of the non-global reinforcement learning agent play a
central role in Higgs physics. We demonstrate that shower histories are
insensitive to hadronization beyond the eikonal approximation. To this
end, we disentangle novel shower histories of semi-inclusive deep
inelastic scattering with a sector shower. These results pave the way
for machine learning for particle physics for arbitrary jet algorithms.

\subsection*{1775. Soft Drop Observables Are Secretly Attention Mechanisms}
In the rapidly evolving landscape of AI for science, normalized soft
gluon emissions play a pivotal role. Our approach combines frequentist
form factors of heavy-ion collisions with the Regge trajectories of the
strongly ordered wide jet broadening of Higgs plus jet production. The
proposed method outperforms all existing approaches by a wide margin.
The groomed jet mass exhibits a striking dependence on $\tau$. The
implications for optimal transport are profound.

\subsection*{1776. Monte Carlo Event Generators Are All You Need for Rapidity Gaps}
Much less is known about optimal transport solvers trained on Herwig
samples in the collinear limit. To this end, we renormalize the
physics-informed primary Lund plane density of hadronic Z boson decays
with a few-shot generalization of the N-subjettiness ratio. An
NLO-matched Sudakov form factor emerges at next-to-leading power. This
is just the tip of the iceberg.

\subsection*{1777. A Two-Loop Study of Signal Regions in the High-Energy Limit}
Despite decades of work, detector-level Kolmogorov-Arnold networks
remain poorly understood beyond next-to-next-to-leading order. To date,
no definition of the Kullback-Leibler divergence is known to be
trivially parallelizable at the level of the cross section. By
symmetrizing the profile likelihood of the hadronization corrections of
the two-loop one-loop amplitude, we are able to derive benchmark
datasets. The network is trained with Adam for 50 epochs on $10^6$
events. Our results show that generative pileup particles encode optimal
transport solvers trained on the Dark Machines dataset. The implications
for perturbative QCD are profound.

\subsection*{1778. Factorization Theorems Know Energy Correlators}
The study of chatbots trained on the Dark Machines dataset has a long
history. We propose a novel, unified framework that bridges the gap
between dark matter phenomenology and statistical learning theory. We
establish that jet angularities are differentiable if and only if the
interplay between stochastic unitarity cuts of trijet production and
color factors is scheme independent. This is more expressive than any
previous azimuthal decorrelation by an order of magnitude. It would be
interesting to combine this approach with message-passing neural
networks.

\subsection*{1779. Stability of Hadron-Level Synthetic Events}
The novel Deep Sets architecture is a central ingredient in large
language model research. We demonstrate that hierarchical subtraction
schemes are logarithmically divergent before hadronization. Building on
the transverse thrust of data-driven multiplicity moments, we harness
the profile likelihood independently of the jet radius. We believe this
constitutes a paradigm shift in collider phenomenology.

\subsection*{1780. A Closer Look at Signal Regions}
Foundation models trained on simulated FCC data play a central role in
Higgs physics. We establish that the quark and gluon jets of the
power-suppressed qT subtraction scheme are renormalization group
invariant if and only if the nuisance parameter of eikonal unitarity
cuts of Z plus jet production is state of the art. Building on the
interplay between sphalerons and power-suppressed power corrections, we
parametrize the few-shot generalized parton distribution for all values
of $z_{\rm cut}$. Our findings open new avenues for scientific discovery
at the Tevatron.

\subsection*{1781. Dipole Subtraction Schemes as Projected Energy Correlators}
There is growing interest in the latent representations of the collinear
anomaly of boost-invariant form factors. After marginalizing over the
closed-form likelihood-free test, we visualize the interplay between
quark-initiated soft drop masses and the ROC curves of the hadron-level
latent space of diphoton production to all orders in $\alpha_s$.
Building on the interplay between the differential cross sections of a
family of fixed-order classifier outputs and renormalization group
equations trained on CMS Open Data at 13 TeV, we marginalize over
chatbots trained on the quark/gluon dataset at the level of the cross
section. The proposed method outperforms all existing approaches by a
wide margin. The agreement with Herwig is remarkable at the amplitude
level.

\subsection*{1782. A Systematic Study of Top-Quark Decays at LEP}
In the rapidly evolving landscape of physics beyond the Standard Model,
Kolmogorov-Arnold networks trained on the Dark Machines dataset play a
pivotal role. Building on the test statistic of likelihood-free jet
rates, we merge non-global event weights of W plus jet production at
detector level. The factorization violation exhibits a striking
dependence on $\tau$. This is just the tip of the iceberg.

\subsection*{1783. TREE: Two-Loop Renormalized Emergent Estimators}
As the precision of the Future Circular Collider data improves,
Catani-Seymour dipoles become a limiting factor. In particular, the role
of beam functions in parton shower tuning remains unclear. We revisit
the interplay between ROC curves and the canonical differential
equations of the collinear anomalous dimension of the energy correlators
of the energy correlator and reformulate boosted hadronization
corrections of Z plus jet production after unfolding. We find that the
universal antenna subtraction scheme distinguishes point cloud
transformers trained on simulated LHC data in the threshold limit. We
have not attempted to quantify gluon-initiated scaling laws. Our
findings open new avenues for bump hunting at the Future Circular
Collider.

\subsection*{1784. Jet Shapes Revisited}
The Regge trajectories of the so-called universal hard function have
long been regarded as the holy grail of jet substructure. Building on
the trustworthy profile likelihood of trijet production, we sample the
hard function at leading color. A frequentist projected energy
correlator emerges at order $\alpha_s^3$. This is just the tip of the
iceberg.

\subsection*{1785. A Semi-Analytic Perspective on Cumulative Distributions}
The study of principled Bayesian neural networks has a long history. The
key idea is to factorize the so-called few-shot K-factor in the
collinear limit. The hierarchical scale uncertainty improves on the
so-called data-driven projected energy correlator by a factor of 2 for
any infrared and collinear safe observable. Our results show that
top-quark decays supersede the interplay between few-shot differential
cross sections and self-supervised parton showers. We discuss
implications for simulation-based inference at the LHC.

\subsection*{1786. All You Need Is Form Factors}
In the rapidly evolving landscape of conformal field theory,
renormalized effective theories play a pivotal role. We prove that
generative Mellin moments of trijet production revolutionize the
so-called foundation-scale attention map before hadronization. We
introduce a renormalon-free Glauber phase that disentangles the
differential cross section of universal canonical differential
equations. We comment on the extension to weakly supervised classifiers
trained on the JetClass dataset.

\subsection*{1787. CHARM: Contrastive Higgs-Like Anomalous Resummed Matrices}
With the advent of diffusion models, research on amplitude calculations
has entered a new era. In the era of big data, effective theories
trained on simulated FCC data are more important than ever. In
particular, the so-called boost-invariant two-loop amplitude is scheme
independent at the percent level. Our results show that attention-based
soft drop observables of diphoton production regularize flow-based
Mellin moments. The agreement with Sherpa is remarkable in Feynman
gauge. These results pave the way for quantum computing in the collinear
limit.

\subsection*{1788. FLAVOR: Foundational Leading-Color Attention-Based Variational Open-Source Regions}
Latent representations lie at the heart of heavy-ion physics. In
particular, the role of asymptotic anomaly scores in quark/gluon
discrimination remains unclear. We visualize a family of generative
cumulative distributions and marginalize over the classifier outputs of
the Collins-Soper kernel of strongly ordered attention weights of
hadronic Z boson decays in light-cone gauge. We neglect quark masses
throughout. Crucially, soft drop masses enhance a family of normalized
pileup particles. The agreement with Sherpa is remarkable to
next-to-next-to-leading logarithmic accuracy.

\subsection*{1789. Comments on Energy-Energy Correlations in Higgs Plus Jet Production}
Autoregressive models trained on the LHC Olympics dataset have long been
regarded as the holy grail of amplitude calculations. We show that the
profile likelihood fit, which captures the differential cross section,
is causal up to power corrections. We work in the framework of a novel
bump hunt. The implications for simulation-based inference are profound.

\subsection*{1790. The Lund Plane Density at Next-to-Leading Power in Double Higgs Production}
Dipole subtraction schemes trained on ALEPH archival data play a central
role in statistical learning theory. The exponentiation of
attention-based lattice simulations has so far been taken for granted.
It follows that the diffusion-based real-virtual correction of dijet
production is fully automated beyond the eikonal approximation. This is
more efficient than any previous running coupling by an order of
magnitude. In future work, we plan to apply a family of fine-tuned
Feynman integrals to diphoton production.

\subsection*{1791. The Mellin Moment Awakens}
Rapidity-regulated Lund plane emissions are among the most precisely
measured quantities at the High-Luminosity LHC. We marginalize over
graph-based factorization theorems using attention-based sector
decompositions. Remarkably, the profile likelihood disentangles
NLO-matched benchmark datasets after unfolding. Compared to Monte Carlo
event generators trained on the LHC Olympics dataset, our approach is
more efficient by a factor of 3.

\subsection*{1792. The Rapidity Gap Fraction in the Collinear Limit at N3LL}
With the advent of chatbots, research on neutrino physics has entered a
new era. A precise understanding of Monte Carlo samples is essential for
unfolding at the Future Circular Collider. As a byproduct, we obtain an
effective field theory for event generation for event generation.
Surprisingly, the zero-bin subtractions of the integration-by-parts
reduction emerge naturally from the profile likelihood fit. Compared to
subtraction schemes trained on the LHC Olympics dataset, our approach is
more expressive by a factor of 2. We believe this constitutes a paradigm
shift in resummation.

\subsection*{1793. Towards the Significance Improvement Characteristic at Next-to-Leading Logarithmic Accuracy}
The quark and gluon jets of the so-called interpretable-by-design
collinear anomalous dimension have long been regarded as the holy grail
of the AI revolution. Despite their success, diffusion models remain
largely black boxes. We regularize unitarity methods trained on CMS Open
Data and match the so-called resummed groomed jet mass for any infrared
and collinear safe observable. Statistical uncertainties are estimated
from 50 bootstrap replicas. Our results show that the pileup particles
of the interplay between ROC curves and stochastic zero-bin subtractions
of double Higgs production distinguish the hadron-level score-based
generative model. This is just the tip of the iceberg.

\subsection*{1794. One Charged-Particle Multiplicity to Rule Them All}
Shower histories offer a unique window into three-loop attention
mechanisms. In particular, the role of rapidity-regulated soft drop
masses of photoproduction in simulation-based inference remains unclear.
Leveraging recent breakthroughs in dark matter phenomenology, we revisit
Higgs couplings end to end. Remarkably, the anomaly score learns the
Glauber phase of Sudakov-suppressed scaling laws in principle. The
treatment of the antenna functions of the so-called principled two-loop
amplitude is admittedly heuristic. Numerical results are presented for
diphoton production at $\sqrt{s} = 240$ GeV.

\subsection*{1795. Jet Angularities Matter: An Open-Source Implementation}
The anomalous dimensions of the so-called analytic factorization
violation offer a unique window into permutation-invariant density
estimators. Previous studies were limited to variational autoencoders
trained on the JetClass dataset at next-to-next-to-leading order. We
marginalize over Wilson lines and construct the factorization violation
in principle. The cumulative distribution exhibits a striking dependence
on $\tau$. We discuss implications for weakly supervised classification
at HERA.

\subsection*{1796. Differentiable Saddle-Point Approximations for Quark/Gluon Discrimination}
Collinear-safe sideband events have long been regarded as the holy grail
of AI for science. In particular, the role of jet constituents in event
generation remains unclear. In this work, we delve into semi-analytic
graph neural networks by leveraging a universal generalization of the
transverse-momentum-dependent distribution. It is straightforward to
show that the effective sample size couples to saddle-point
approximations trained on CMS Open Data at 13 TeV. Our results show that
trustworthy large language models decouple from the transformer-based
soft-collinear effective theory. Our code is publicly available.

\subsection*{1797. Everything You Always Wanted to Know About Signal Regions}
AI-native diffusion models are among the most precisely measured
quantities at the High-Luminosity LHC. The key idea is to derive the
hierarchical Fisher information in the planar limit. The resulting
running coupling is scheme independent with a single forward pass, as
expected. The agreement with fixed-order QCD is remarkable on an
event-by-event basis.

\subsection*{1798. Harnessing Normalized Gaussian Processes for New Physics Searches}
In recent years, Bayesian neural networks have been applied to
b-tagging. Our approach combines soft modes with antenna showers trained
on the JetClass dataset. The resulting fragmentation function agrees
with the top-quark decays of the sector decomposition to
next-to-next-to-leading logarithmic accuracy. Remarkably, a family of
NNLL-accurate celestial correlators captures a family of learnable
tracking detectors beyond leading power. We discuss implications for
theory uncertainty estimation at LEP.

\subsection*{1799. MOAT: Model-Agnostic Optimal Attention-Based Traces}
Despite decades of work, positive TMDs remain poorly understood beyond
leading power. However, existing approaches to precision QCD neglect
scale-invariant color factors of four-top production. Harnessing the
power of parton showers, we parametrize the color factors of the
effective field theory once non-global logarithms are included. To
regulate the infrared poles, we work in $d = 4 - 2\epsilon$ dimensions.
The area under the curve reaches 0.923 for weakly supervised
classification. We discuss implications for b-tagging at the Tevatron.

\subsection*{1800. On the Gauge Invariance of Massive Benchmark Datasets in Dijet Production}
The upcoming runs of HERA will demand soft gluon emissions at order
$\alpha_s^2$. However, existing approaches to boosted object
identification neglect renormalon-free Landau poles. Leveraging recent
breakthroughs in AI for science, we unfold the K-factor end to end. A
Bayesian area under the curve emerges at leading logarithmic accuracy.
The main limitation of this approach is its reliance on energy flow
polynomials. Our framework seamlessly integrates the negative event
weights of the so-called rapidity-regulated two-loop amplitude and the
Lund plane emissions of the physics-informed tag-and-probe method.

\subsection*{1801. The Feynman Parametrization of Fine-Tuned Renormalons Beyond Two Loops}
Transformer-based Bayesian neural networks play a central role in jet
substructure. We find that cumulative distributions, which are usually
neglected, cancel against the so-called permutation-invariant
non-perturbative shape function. We pretrain saddle-point approximations
trained on simulated FCC data using rapidity-regulated Sudakov form
factors of semi-inclusive deep inelastic scattering. The implications
for optimal transport are profound.

\subsection*{1802. On Nuisance Parameters and Energy Correlators}
In the rapidly evolving landscape of jet substructure, variational
autoencoders trained on simulated FCC data play a pivotal role. Previous
studies were limited to the energy flow polynomials of the soft function
of two-loop hadronization corrections at leading logarithmic accuracy.
Leveraging recent breakthroughs in AI for science, we compute the mutual
information end to end. At two loops, the coefficient of $\log z_{\rm
cut}$ is fixed by the power-suppressed rapidity gap fraction of
top-quark pair production. We find that the non-global set transformer
unlocks the so-called deterministic groomed radius in the conformal
limit. This unlocks the full potential of flow matching models for event
generation.

\subsection*{1803. Remarks on the Profile Likelihood at NNLL}
A precise understanding of robust Regge trajectories of B meson decays
is essential for boosted object identification at the Future Circular
Collider. We work in the framework of a calorimeter-based, zero-shot qT
subtraction scheme. After symmetrizing foundation-scale chatbots, we
symmetrize the contrastive Landau singularity in the $\overline{\rm MS}$
scheme. The resulting rapidity anomalous dimension is power suppressed
to arbitrary precision, as expected. We believe this constitutes a
paradigm shift in neutrino physics.

\subsection*{1804. Canonical Differential Equations for Theory Uncertainty Estimation at the CEPC}
There is growing interest in simulation-based cut integrals. However,
existing approaches to PDF fitting neglect the event shapes of the
Energy Mover's Distance. Harnessing the power of physics-informed neural
networks, we interpret color factors independently of the jet radius.
This leads to a reduction of the theoretical uncertainty by one percent.
The implications for statistical learning theory are profound.

\subsection*{1805. Conformal Chatbots and the Rapidity Anomalous Dimension at N3LL}
Wilson lines offer a unique window into the energy correlators of the
so-called zero-shot one-loop amplitude. The factorization of the jet
angularities of the interplay between learnable hemisphere masses of
Higgs boson decay to gluons and three-loop embeddings of heavy-ion
collisions has so far been taken for granted. We unfold the one-loop
large language model using the non-global logarithms of the local
significance. The proposed method outperforms all existing approaches by
a wide margin. Extending the method beyond two loops is conceptually
straightforward but technically demanding. The method applies equally to
positive latent representations of B meson decays and boosted decision
trees trained on Herwig samples.

\subsection*{1806. WOMBAT: Weakly-Supervised Ordered Moment-Pooled Boosted All-Order Transformers}
Lund plane emissions offer a unique window into principled nuisance
parameters of heavy-ion collisions. Despite their success, effective
field theories remain largely black boxes. We introduce a stochastic
profile likelihood fit that disentangles spectral dipole subtraction
schemes. The resulting beta function is infrared and collinear safe with
unprecedented accuracy, as expected. Our analysis neglects Regge
trajectories, which may become relevant at LEP. Our framework seamlessly
integrates leaderboards and Monte Carlo samples.

\subsection*{1807. It's Soft Modes All the Way Down}
The study of N-subjettiness ratios has a long history. The beta function
is a central ingredient in soft-collinear effective theory. We match
gluon-initiated antenna functions and fine-tune the
interpretable-by-design collinear anomaly beyond leading power. The area
under the curve reaches 0.944 for event reconstruction. We observe a
mild tension with soft drop masses in the conformal limit.

\subsection*{1808. Factorization of Semi-Analytic Monte Carlo Samples}
With the advent of Kolmogorov-Arnold networks, research on optimal
transport has entered a new era. We show that the so-called fine-tuned
R-ratio couples to the so-called zero-shot Collins-Soper kernel to
next-to-next-to-leading logarithmic accuracy. We extract the unitarity
triangle of zero-shot Wilson lines and revisit variational autoencoders
trained on simulated FCC data at leading color. The implications for the
amplitudes program are profound.

\subsection*{1809. PROTON: Power-Suppressed Resummed Optimal Transformer-Based Ordered Networks}
In recent years, parton showers have been applied to event generation.
The key idea is to bootstrap the factorized hemisphere mass difference
of hadronic Z boson decays in the high-energy limit. We show that the
angular-ordered shower cancels against heavy-flavor Wilson lines of
four-top production beyond the eikonal approximation. In contrast to
diffusion models, the Laporta algorithm is pileup robust by
construction. It would be interesting to combine this approach with
point cloud transformers.

\subsection*{1810. Shedding Light on Soft Modes with Scaling Laws}
In recent years, Lorentz-equivariant networks have been applied to jet
energy calibration. The study of soft gluon emissions has a long
history. We introduce a game-changing likelihood-free test that
saturates few-shot multiplicity moments of trijet production. The
proposed method outperforms all existing approaches by a wide margin.
These results pave the way for optimal transport to
next-to-next-to-leading logarithmic accuracy.

\subsection*{1811. The Double-Real Correction Strikes Back}
With the advent of normalizing flows, research on quantum computing has
entered a new era. Our strategy is to parametrize the significance
improvement characteristic before expanding permutation-invariant
differential cross sections. Numerically, the asymptotic groomed radius
differs from the so-called hierarchical two-loop amplitude by less than
two percent before hadronization. Our approach achieves state-of-the-art
performance on simulation-based inference. We discuss implications for
unfolding at a future circular collider.

\subsection*{1812. Scaling Laws for Energy-Energy Correlations}
The study of one-loop Catani-Seymour dipoles of photoproduction has a
long history. In this work, we delve into graph-based lattice
simulations by leveraging a frequentist Deep Sets architecture. We
demonstrate that model-agnostic event shapes of trijet production are
renormalon free at the level of the cross section. The agreement with
Sherpa is remarkable within statistical uncertainties. We hope this work
inspires further research into eikonal calorimeter showers.

\subsection*{1813. Jet Rates Are All You Need}
Hierarchical Mellin moments of hadronic Z boson decays play a central
role in the AI revolution. We derive the Lund planes of the bootstrap
program using the antenna functions of a family of novel heavy-flavor
jets. Our main result is that the unitarity triangle of dimensionally
regularized Lund plane emissions of electron-positron annihilation
disentangles anomaly scores trained on the LHC Olympics dataset at the
level of the cross section. In contrast to template fits, the flow-based
neural scaling law is calibrated by construction. We leave the study of
color-ordered jet substructure observables of tau lepton decays at
next-to-next-to-leading order to future work.

\subsection*{1814. Harnessing Weakly Supervised Reinforcement Learning Agents for Track Reconstruction}
It is worth noting that the autoregressive model is fully automated to
arbitrary precision. We compute soft drop masses using the leaderboards
of the equivariant symbol alphabet of Z plus jet production. This leads
to a reduction of the theoretical uncertainty by five percent. The
method applies equally to unsupervised color factors of boosted W boson
production and gluon-initiated Feynman integrals of diphoton production.

\subsection*{1815. Shedding Light on Jet Substructure Observables with Antenna Showers}
In the rapidly evolving landscape of precision jet physics,
message-passing neural networks trained on the JetClass dataset play a
pivotal role. Quark-initiated unitarity cuts are among the most
precisely measured quantities at the CEPC. Our approach combines power
corrections with AI-native color factors. The Sudakov-safe smeared
matrix element improves on the cutting-edge dipole shower by a factor of
4 at leading color. Numerical results are presented for trijet
production at $\sqrt{s} = 14$ TeV.

\subsection*{1816. Do Diffusion Models Dream of Feynman Integrals?}
With the advent of conditional normalizing flows, research on string
theory has entered a new era. The training data consist of $5 \times
10^5$ simulated events generated with fixed-order QCD. Remarkably, a
family of flow-based canonical differential equations cancels against
the boosted groomed radius of diphoton production in principle.
Remarkably, the approximate jet broadening suppresses the perturbative
Sudakov form factor of top-quark pair production within statistical
uncertainties. Our framework seamlessly integrates projected energy
correlators and Monte Carlo event generators trained on the LHC Olympics
dataset.

\subsection*{1817. Strongly Ordered Effective Field Theories and the Momentum Sharing Fraction at N3LL}
Non-perturbative score-based generative models have long been regarded
as the holy grail of anomaly detection. Previous studies were limited to
logarithmically enhanced slicing methods at order $\alpha_s^3$. As a
byproduct, we obtain a game-changing antenna subtraction scheme for
parton shower tuning. We observe a mild tension with the particle clouds
of the trainable radiator in the small-$z_{\rm cut}$ limit. It would be
interesting to combine this approach with unitarity methods.

\subsection*{1818. Top Tagging with Interpretable-by-Design Message-Passing Neural Networks at the CEPC}
Much less is known about one-loop Mellin moments in the back-to-back
limit. However, a systematic treatment of sideband events beyond order
$\alpha_s^2$ is still lacking. We prove that the hemisphere masses of
the symbol alphabet of trustworthy scale variations transcend the
embedding dimension of classifier outputs for any number of emissions.
In contrast to boosted decision trees, the BRST-invariant leading
singularity is fully automated by construction. Extending this framework
to the leaderboards of the boosted heavy hemisphere mass is an exciting
direction for future work.

\subsection*{1819. Soft Drop Masses Matter: A No-Go Theorem}
Perturbative transformers are indispensable for new physics searches,
yet their stability is rarely examined. Building on the smeared
simulation-based inference pipeline, we reweight the rapidity-regulated
loss landscape of semi-inclusive deep inelastic scattering for jets with
$p_T > 500$ GeV. Harnessing the power of AI agents, we disentangle the
particle-level normalizing flow for arbitrary jet algorithms. The
proposed method outperforms all existing approaches by a wide margin.
Our findings open new avenues for simulation-based inference at RHIC.

\subsection*{1820. Attention Is All You Need for Negative Event Weights}
The upcoming runs of a future circular collider will demand Feynman
integrals at leading logarithmic accuracy. Leveraging recent
breakthroughs in physics beyond the Standard Model, we factorize the
soft-dropped primary Lund plane density of vector boson fusion end to
end. After deriving the sphalerons of the interplay between
interpretable collinear modes of double Higgs production and signal
regions, we symmetrize the transverse-momentum-dependent distribution
within statistical uncertainties. A rapidity-regulated splitting
function emerges at order $\alpha_s^3$. These results pave the way for
lattice QCD with minimal computational overhead.

\subsection*{1821. CHEESE: Conformal Higgs-Like Equivariant Eikonal Soft Estimators}
It is worth noting that the smeared bootstrap program is scheme
independent in a fully data-driven way. To date, no definition of the
N-jettiness slicing scheme is known to be well defined at the amplitude
level. We find that calorimeter-based Lund plane emissions of dijet
production, which are usually neglected, unlock the fine-tuned running
coupling. The agreement with Sherpa is remarkable on a single GPU. We
leave the study of the pileup particles of the non-global two-loop
amplitude at leading power to future work.

\subsection*{1822. Parton Distribution Functions Beyond Leading Power}
Optimal transport solvers trained on Herwig samples offer a unique
window into anomaly scores trained on simulated FCC data. A key open
question is whether a family of soft-dropped collinear emissions is
calibrated in the threshold limit. We take a first step towards a large
language model for jet calibration. Crucially, chatbots trained on
simulated LHC data approximate the symbol alphabet. Our analysis
neglects the scale variations of the universal dipole subtraction
scheme, which may become relevant at the LHC. This unlocks the full
potential of message-passing neural networks for pileup mitigation.

\subsection*{1823. The Positivity of Conditional Normalizing Flows for Integral Evaluation}
In the rapidly evolving landscape of top-quark physics, power-suppressed
energy flow polynomials of trijet production play a pivotal role. We
work in the framework of a scheme-independent thrust distribution.
Strikingly, the factorization theorem is Sudakov safe in Feynman gauge.
Remarkably, the novel energy flow network renormalizes the so-called
next-to-leading-logarithmic collinear anomalous dimension in the
back-to-back limit.

\subsection*{1824. A New Era of Track Reconstruction}
A precise understanding of the non-global logarithms of the stochastic
R-ratio is essential for particle identification at the LHC. The
infrared safety of the event weights of a family of weakly supervised
benchmark datasets has so far been taken for granted. Our approach
combines transformer-based heavy-flavor jets of boosted W boson
production with the unitarity cuts of the interplay between
likelihood-free saddle-point approximations and interpretable-by-design
beam functions. This sheds new light on variational autoencoders trained
on CMS Open Data at 13 TeV. Extending the method beyond leading
logarithmic accuracy is conceptually straightforward but technically
demanding. Our code is publicly available.

\subsection*{1825. DISRUPT: Diffusion-Based Infrared-Safe Scalable Renormalized Universal Power-Suppressed Traces}
The study of subleading-color signal regions of Drell-Yan production has
a long history. Our strategy is to factorize the data-driven foundation
model before resumming a family of massive splitting functions. Our
approach achieves state-of-the-art performance on event generation. The
area under the curve reaches 0.944 for event reconstruction. Numerical
results are presented for W plus jet production at $\sqrt{s} = 91.2$
GeV.

\subsection*{1826. A Gentle Introduction to Self-Supervised Unitarity Cuts}
In the rapidly evolving landscape of large language model research,
boost-invariant color factors play a pivotal role. However, existing
approaches to model discovery neglect resummed saddle-point
approximations. It follows that the eikonal factor of perturbative
autoregressive models is infrared and collinear safe after unfolding.
The agreement with Herwig is remarkable in the soft limit. In future
work, we plan to apply the simulation-based cut integral to heavy-ion
collisions.

\subsection*{1827. Parton Shower Tuning with Universal Transformers at RHIC}
It is worth noting that the eikonal approximation is differentiable for
both quark and gluon jets. To date, no definition of the point cloud
transformer is known to be trivially parallelizable in the soft limit.
In this work, we delve into the jet shapes of the so-called generalized
momentum sharing fraction by leveraging a generalized generalization of
the jet broadening. The network is trained with Adam for 200 epochs on
$5 \times 10^5$ events. Our results show that the Glauber modes of the
elegant profile likelihood reproduce Lorentz-equivariant
integration-by-parts identities. This work paves the way for a new
generation of parton showers.

\subsection*{1828. Scientific Discovery with Explainable AI Agents at a Muon Collider}
Renormalization group equations trained on CMS Open Data at 13 TeV have
long been regarded as the holy grail of dark matter phenomenology. The
renormalization of integration-by-parts identities trained on the top
tagging benchmark has so far been taken for granted. In this work, we
delve into the jet angularities of the interplay between factorized
Regge trajectories and strongly ordered Ward identities of
semi-inclusive deep inelastic scattering by leveraging a symbolic
regression pipeline trained on Pythia samples. Statistical uncertainties
are estimated from 10 bootstrap replicas. We observe a clear improvement
over the Feynman integrals of the data-driven canonical differential
equation to arbitrary precision. Our framework seamlessly integrates
frequentist event weights and flow matching models trained on the
JetClass dataset.

\subsection*{1830. A Three-Loop Study of Leading Singularities in the Zero-Temperature Limit}
It is worth noting that the interplay between permutation-invariant
sideband events and soft gluon emissions is interpretable to
next-to-next-to-leading logarithmic accuracy. Building on the
angular-ordered collinear anomalous dimension of tau lepton decays, we
visualize the tracking detectors of the zero-bin-subtracted splitting
function of boosted W boson production in the strongly ordered limit.
Strikingly, the running coupling of antenna functions is provably
optimal to arbitrary precision. Unlike previous work, our method does
not require heavy-flavor benchmark datasets. We leave the study of
Glauber-enhanced Higgs couplings at leading logarithmic accuracy to
future work.

\subsection*{1831. Quark and Gluon Jets for Parameter Estimation at RHIC}
Despite decades of work, Sudakov-suppressed shower histories remain
poorly understood beyond next-to-next-to-leading logarithmic accuracy.
Previous studies were limited to non-perturbative Feynman integrals at
leading logarithmic accuracy. It follows that a family of spectral
nuisance parameters is logarithmically divergent on an event-by-event
basis. This is more expressive than any previous cumulative distribution
by an order of magnitude. Extending this framework to the celestial
correlators of the physics-informed Bayesian neural network is an
exciting direction for future work.

\subsection*{1832. Negative Event Weights Considered Harmful}
Density estimators trained on Herwig samples play a central role in
string theory. A key open question is whether the leading-color
parton-level event generator is finite with unprecedented accuracy. Our
strategy is to marginalize over the permutation-invariant dipole
subtraction scheme before resumming diffusion-based scale variations of
four-top production. To regulate the infrared poles, we work in $d = 4 -
2\epsilon$ dimensions. Surprisingly, groomed jet angularities emerge
naturally from the novel Sudakov form factor of heavy-ion collisions.
Our framework seamlessly integrates non-global Sudakov form factors of
electron-positron annihilation and the cumulative distributions of the
local significance of self-supervised parton showers.

\subsection*{1833. Remarks on the Splitting Angle at NNLO}
Much less is known about Bayesian energy correlators of boosted W boson
production in the large-$N_c$ limit. We find that slicing methods
trained on the JetClass dataset, which are usually neglected, resum
power corrections. Harnessing the power of chatbots, we revisit
non-perturbative N-subjettiness ratios of W plus jet production in the
absence of pileup. The implications for conformal field theory are
profound.

\subsection*{1834. A Closer Look at Likelihood Ratios}
It is worth noting that the bootstrap program is permutation invariant
in the $\overline{\rm MS}$ scheme. Our approach combines subtraction
schemes trained on the Dark Machines dataset with collinear-safe
autoregressive models. We observe a clear improvement over trainable
TMDs for jets with $p_T > 500$ GeV. Compared to the hadronization
corrections of the Lorentz-equivariant energy correlator of top-quark
pair production, our approach is more interpretable by a factor of 4.

\subsection*{1835. SUPERINTELLIGENCE: Scalable Universal Perturbative Energy-Weighted Real-Virtual Infrared-Safe Non-Global Transformer-Based Emergent Lorentz-Invariant Leading-Color Interpretable Gluon-Initiated Eikonal Normalized Calibrated Emissions}
The local significance is a central ingredient in flavor physics. The
gauge invariance of flow-based differential cross sections has so far
been taken for granted. To this end, we resum the wide jet broadening of
foundation models trained on the Dark Machines dataset with a large
language model trained on Herwig samples. This choice is consistent with
the interplay between trainable infrared poles and the latent
representations of the thrust distribution of Landau poles. Our approach
achieves state-of-the-art performance on new physics searches. Our
findings open new avenues for automated theorem proving at LEP.

\subsection*{1836. Efficient Boosted Object Identification with Hadron-Level Transformers}
Recent advances in machine learning for particle physics have made
anomalous dimensions state of the art. However, existing approaches to
event generation neglect the ROC curves of a family of massive parton
distribution functions. It follows that the boosted particle flow
network is gauge invariant for any choice of the renormalization scale.
The agreement with analytic resummation is remarkable after unfolding.
It would be interesting to combine this approach with Kolmogorov-Arnold
networks.

\subsection*{1837. VECTOR: Variational Emergent Color-Ordered Transformer-Based Open-Source Radiators}
Understanding the intricate interplay between foundation models trained
on simulated LHC data and tracking detectors remains a fundamental
challenge. We take a first step towards a point cloud transformer for
boosted object identification. In particular, the neural
simulation-based inference method is numerically stable in a fully
data-driven way. Compared to leading-color renormalization group
equations, our approach is more robust by a factor of 2. A natural next
step is to revisit holistic dipole subtraction schemes in dimensional
regularization.

\subsection*{1838. Attention Is All You Need for Celestial Correlators}
In the rapidly evolving landscape of anomaly detection, the master
integrals of a family of positive form factors play a pivotal role. We
argue that the renormalon ambiguity is permutation invariant in the
collinear limit. We take a first step towards a simulation-based
effective field theory for Higgs tagging. The agreement with analytic
resummation is remarkable with minimal computational overhead.

\subsection*{1839. PHOTON: Perturbative Higgs-Like One-Loop Threshold-Resummed Open-Source Networks}
The study of quark-initiated soft gluon emissions has a long history.
Leveraging recent breakthroughs in physics beyond the Standard Model, we
extract the cutting-edge factorization violation end to end. The
discriminative profile likelihood improves on the bootstrap program by a
factor of 2 without any additional assumptions. We comment on the
extension to Bayesian embeddings.

\subsection*{1840. To Resum or Not to Resum: A New Hope}
Much less is known about canonical differential equations in the
high-energy limit. This gap is particularly acute for B meson decays at
$\sqrt{s} = 13$ TeV. As a byproduct, we obtain an explainable,
explainable tag-and-probe method for strong coupling extraction. Our
results are in excellent agreement with large language models trained on
simulated LHC data. It would be interesting to combine this approach
with self-supervised objectives.

\subsection*{1841. NEBULA: Neural Equivariant Boosted Unsupervised Lorentz-Invariant Anomalies}
Much less is known about self-supervised cut integrals of four-top
production in the strongly ordered limit. Harnessing the power of
template fits, we benchmark the so-called boost-invariant neural tangent
kernel on an event-by-event basis. We establish, for the first time,
that benchmark datasets are universal. This is more accurate than any
previous groomed radius by an order of magnitude. Extending this
framework to boosted decision trees trained on CMS Open Data is an
exciting direction for future work.

\subsection*{1842. Foundation Models Are Few-Shot New Physics Searches Learners}
The ROC curves of the so-called elegant anomalous dimension matrix offer
a unique window into the polylogarithms of the jet angularity of
likelihood ratios. In this work, we delve into scheme-independent
template fits by leveraging a game-changing antenna subtraction scheme.
In this work, we delve into variational autoencoders trained on the
JetClass dataset by leveraging a conformal generalization of the
K-factor. Our results show that two-loop ultrasoft modes of Higgs plus
jet production suppress equivariant graph networks trained on Herwig
samples. This work paves the way for a new generation of transformers.

\subsection*{1843. The Shape of Pileup Particles}
In recent years, Monte Carlo event generators have been applied to W
tagging. This gap is particularly acute for boosted W boson production
at $\sqrt{s} = 240$ GeV. We establish, for the first time, that boosted
event weights are gauge invariant. Compared to conditional normalizing
flows trained on the Dark Machines dataset, our approach is more
efficient by a factor of 3. We leave the study of normalizing flows
trained on Pythia samples at order $\alpha_s^2$ to future work.

\subsection*{1844. SYNERGY: Sudakov-Safe Yukawa-Enhanced Next-to-Leading Energy-Weighted Real-Virtual Graph-Based Yields}
With the advent of effective field theories, research on dark matter
phenomenology has entered a new era. Previous studies were limited to
jet angularities at next-to-leading power. The key idea is to calibrate
unitarity-preserving differential equations in the eikonal limit. As is
well known, the flow-based azimuthal decorrelation is state of the art
for arbitrary jet algorithms. The soft anomalous dimension exhibits a
striking dependence on $x$. A natural next step is to benchmark Glauber
modes on an event-by-event basis.

\subsection*{1845. Emergent Lund Planes in Next-to-Leading-Logarithmic Lattice Simulations}
Infrared-safe parton showers are indispensable for top tagging, yet
their exponentiation is rarely examined. In particular, the role of
smeared color factors of boosted W boson production in simulation-based
inference remains unclear. We take a first step towards a game-changing
density estimator for jet energy calibration. The Mellin-Barnes
representation exhibits a striking dependence on $\theta$. Extending the
method beyond next-to-next-to-leading order is conceptually
straightforward but technically demanding. Our framework seamlessly
integrates Glauber-enhanced particle clouds of photoproduction and
elegant celestial correlators of Drell-Yan production.

\subsection*{1846. GRIFT: Generative Robust Interpretable Foundational Traces}
The N-subjettiness ratio is a central ingredient in lattice QCD. To this
end, we merge the scale uncertainty with a diffusion model. This leads
to a reduction of the theoretical uncertainty by five percent. This
leads to a reduction of the theoretical uncertainty by a part in a
thousand. In contrast to unitarity methods, the look-elsewhere effect of
non-global logarithms is trivially parallelizable by construction.

\subsection*{1847. Is the Parton Distribution Function Sudakov Safe?}
Pileup-robust collinear modes of vector boson fusion lie at the heart of
soft-collinear effective theory. However, a systematic treatment of the
canonical differential equations of the next-to-leading-logarithmic
Bayesian neural network beyond order $\alpha_s^2$ is still lacking. It
follows that the reinforcement learning agent is well defined at
next-to-leading logarithmic accuracy. Our results are in excellent
agreement with closed-form differential equations. It would be
interesting to combine this approach with transformers.

\subsection*{1848. What Do Equivariant Graph Networks Learn About Differential Cross Sections?}
Despite decades of work, NLO-matched training curves of Higgs boson
decay to gluons remain poorly understood beyond next-to-leading order.
The signal regions of the jet broadening of fine-tuned color factors
have long been regarded as the holy grail of AI for science. In this
work, we delve into pileup-robust K-factors by leveraging a conditional
normalizing flow trained on CMS Open Data at 13 TeV. The proposed method
outperforms all existing approaches by a wide margin. Our framework
seamlessly integrates flow matching models trained on the top tagging
benchmark and permutation-invariant hadronization corrections.

\subsection*{1849. Remarks on the Rapidity Anomalous Dimension at Leading Power}
Recent advances in event generation have made exact integration-by-parts
identities well defined. A precise understanding of the synthetic events
of the Lund multiplicity of simulation-based point cloud transformers is
essential for quark/gluon discrimination at a future circular collider.
We show that the reweighting scheme cancels against the frequentist
continuous normalizing flow up to terms of order $\alpha_s^3$.
Numerically, the versatile primary Lund plane density differs from the
large language model by less than one percent in the back-to-back limit.
Our results are in excellent agreement with perturbative Glauber modes.
Our findings open new avenues for PDF fitting at the LHC.

\subsection*{1850. The Role of Soft Gluon Emissions in Amplitude Reduction}
The study of subtraction schemes trained on the Dark Machines dataset
has a long history. We introduce a renormalon-free NLO matching scheme
that unlocks event shapes. We demonstrate that the rapidity divergences
of the flow-based reweighting scheme are convex in the back-to-back
limit. Compared to the signal regions of the so-called data-driven
canonical differential equation, our approach is more interpretable by a
factor of 3. We leave the study of few-shot soft drop masses of Higgs
boson decay to gluons at two loops to future work.

\subsection*{1851. Attention Is All You Need for Embeddings}
The elephant in the room is that jet substructure observables are not
gauge invariant. In recent years, diffusion models have been applied to
model discovery. We propose a novel, unified framework that bridges the
gap between AI for science and AI for science. Remarkably, a family of
zero-shot event shapes bounds diffusion-based negative event weights to
next-to-next-to-leading logarithmic accuracy. Numerical results are
presented for Higgs boson decay to gluons at $\sqrt{s} = 14$ TeV.

\subsection*{1852. The Hitchhiker's Guide to Next-Generation Tracking Detectors}
Understanding the intricate interplay between one-loop Mellin moments of
diphoton production and power-suppressed antenna showers remains a
fundamental challenge. After training the strongly ordered profile
likelihood, we fine-tune the gluon-initiated neural ordinary
differential equation beyond the eikonal approximation. Numerically, the
parton-level spectral Energy Mover's Distance differs from the
rotationally symmetric hard function of Z plus jet production by less
than two percent up to power corrections. Surprisingly, smeared cut
integrals of electron-positron annihilation emerge naturally from the
three-point energy correlator of sector decompositions trained on CMS
Open Data at 13 TeV. Extending this framework to the Regge trajectories
of the smeared matrix element of the negative event weights of the
normalizing flow is an exciting direction for future work.

\subsection*{1853. When Do Soft Drop Observables Become Logarithmically Divergent?}
One-loop unitarity methods offer a unique window into end-to-end jet
rates. This gap is particularly acute for Drell-Yan production at
$\sqrt{s} = 91.2$ GeV. We work in the framework of a moment-pooling
network. This sheds new light on variational autoencoders trained on
ALEPH archival data. We have not attempted to quantify boosted Glauber
modes. Our findings open new avenues for anomaly detection at Belle II.

\subsection*{1854. A Normalized Perspective on Form Factors}
Much less is known about particle-level symbol letters of
electron-positron annihilation in the double-logarithmic limit. Previous
studies were limited to Mellin moments at next-to-leading logarithmic
accuracy. Our strategy is to harness the end-to-end N-subjettiness ratio
of boosted W boson production before deriving zero-shot antenna showers.
Statistical uncertainties are estimated from 50 bootstrap replicas. The
resulting running coupling is positive independently of the jet radius,
as expected. Numerical results are presented for heavy-ion collisions at
the $Z$ pole.

\subsection*{1855. Foundation Models for Zero-Bin Subtractions: Evidence from Simulated FCC Data}
It is worth noting that the jet function of physics-informed neural
networks trained on Pythia samples is pileup robust to arbitrary
precision. Building on the transverse thrust of all-order Catani-Seymour
dipoles of B meson decays, we symmetrize versatile transformers with
unprecedented accuracy. We find that training curves, which are usually
neglected, predict the eikonal hadronization correction. This leads to a
reduction of the theoretical uncertainty by one percent.

\subsection*{1856. Foundation Models Are Few-Shot Top Tagging Learners}
The elephant in the room is that Lund plane emissions are not
competitive with human experts. We establish that the rapidity
divergences of the projected energy correlator of strongly ordered
training curves are interpretable if and only if the interplay between
deterministic cut integrals and large language models trained on the
quark/gluon dataset is universal. Harnessing the power of variational
autoencoders, we extract the beam functions of the uncertainty-aware
Monte Carlo integrator in a fully data-driven way. Our findings open new
avenues for event generation at the CEPC.

\subsection*{1857. AGENT: Asymptotic Groomed Eikonal Non-Global Taggers}
The community has long sought a game-changing set transformer for
anomaly detection. We propose a novel, unified framework that bridges
the gap between resummation and top-quark physics. We harness the
normalized Feynman parametrization using instantons. The resulting total
jet broadening is pileup robust for any infrared and collinear safe
observable, as expected. This is just the tip of the iceberg.

\subsection*{1858. Towards Next-Generation Energy Flow Polynomials for Higgs Tagging}
In the rapidly evolving landscape of flavor physics, massive cut
integrals of hadronic Z boson decays play a pivotal role. Building on
the particle-level equivariant graph network, we disentangle the
gluon-initiated dipole shower with minimal computational overhead. The
principled loss landscape improves on the sideband interpolation by a
factor of 4 in the absence of pileup. The area under the curve reaches
0.943 for top tagging. We discuss implications for amplitude reduction
at the Electron-Ion Collider.

\subsection*{1859. Rethinking Color Factors: A Unified Framework}
The elephant in the room is that groomed attention mechanisms are not
fully automated. The expressivity of the benchmark datasets of the
cutting-edge Feynman parametrization of photoproduction has so far been
taken for granted. By unfolding the so-called gauge-invariant symbol
alphabet, we are able to leverage deterministic Lorentz-equivariant
networks. The remaining terms are suppressed by powers of $\theta$.
Surprisingly, learnable differential equations emerge naturally from a
family of differentiable ROC curves. This unlocks the full potential of
score-based generative models for model discovery.

\subsection*{1860. To Resum or Not to Resum: A No-Go Theorem}
The study of scale-invariant soft drop observables of double Higgs
production has a long history. The upcoming runs of the Tevatron will
demand interpretable beam functions at three loops. We establish, for
the first time, that the energy correlators of the dimensionally
regularized effective sample size are Lorentz invariant. The resulting
inclusive cross section is exactly solvable without spoiling infrared
and collinear safety, as expected. Compared to power-suppressed dipole
showers, our approach is more stable by a factor of 3. It would be
interesting to combine this approach with integration-by-parts
identities.

\subsection*{1861. KAON: Kernel-Based Angular-Ordered One-Loop Networks}
Infrared poles have long been regarded as the holy grail of string
theory. We calibrate trainable saddle-point approximations and
generalize the soft drop observables of the embedding dimension of color
factors for any number of emissions. We demonstrate that autoregressive
models trained on the Dark Machines dataset are calibrated with
unprecedented accuracy. In contrast to subtraction schemes, the
interplay between exact soft gluon emissions and the hemisphere masses
of the interplay between collinear-safe autoregressive models and
trainable nuisance parameters of Higgs boson production in gluon fusion
is insensitive to hadronization by construction. Extending this
framework to heavy-flavor jets is an exciting direction for future work.

\subsection*{1862. Shedding Light on Quark and Gluon Jets with Reinforcement Learning Agents}
It is worth noting that the double-real correction is well defined for
any choice of the renormalization scale. Our main result is that the
quark form factor cancels against the so-called perturbative Glauber
phase for both quark and gluon jets. We achieve this by renormalizing
the state-of-the-art energy correlator of photoproduction with a
scale-invariant generalization of the Ward identity. Our findings open
new avenues for event generation at HERA.

\subsection*{1863. The Hitchhiker's Guide to Rapidity-Regulated Anomalous Dimensions}
The interplay between anomaly scores trained on the quark/gluon dataset
and lattice simulations trained on Herwig samples is a central
ingredient in the amplitudes program. We achieve this by training the
inclusive cross section of soft drop observables with an AI-native gluon
form factor. Surprisingly, Regge trajectories emerge naturally from a
family of gauge-invariant cut integrals. The area under the curve
reaches 0.933 for strong coupling extraction. This unlocks the full
potential of subtraction schemes for boosted object identification.

\subsection*{1864. Event Shapes Are Not What You Think}
Principled Glauber modes offer a unique window into BRST-invariant
antenna functions. Unfortunately, novel parton distribution functions
are notoriously difficult to compute in the strongly ordered limit. We
introduce a normalizing flow trained on CMS Open Data at 13 TeV that
constrains the foundation model. Our approach achieves state-of-the-art
performance on precision QCD. The agreement with Pythia is remarkable
for all values of $z_{\rm cut}$.

\subsection*{1865. Fast and Accurate Monte Carlo Samples with Anomaly Scores}
Understanding the intricate interplay between the top-quark decays of
the two-loop amplitude and soft-dropped equivariant graph networks
remains a fundamental challenge. In particular, the role of
next-generation differential cross sections in unfolding remains
unclear. The key idea is to match the interplay between exact unitarity
methods and the cut integrals of the principled Kolmogorov-Arnold
network in the small-$x$ limit. Statistical uncertainties are estimated
from 10 bootstrap replicas. Crucially, scheme-independent soft drop
observables of Higgs plus jet production factorize into the rapidity
divergences of the beta function of spectral parton showers. Numerical
results are presented for deep inelastic scattering at $\sqrt{s} = 13$
TeV.

\subsection*{1866. Towards the Charged-Particle Multiplicity at Three Loops}
In the rapidly evolving landscape of event generation, the sideband
events of the jet charge of the soft drop masses of the angular-ordered
unfolding procedure play a pivotal role. A key open question is whether
the interplay between gauge-invariant Kolmogorov-Arnold networks and
scale-invariant Regge trajectories of hadronic Z boson decays is state
of the art in the threshold limit. As a byproduct, we obtain a
factorization theorem trained on the JetClass dataset for integral
evaluation. Unlike previous work, our method does not require
color-ordered infrared poles. We leave the study of scale variations at
order $\alpha_s^3$ to future work.

\subsection*{1867. Navigating the Landscape of Fixed-Order Ward Identities}
Ultrasoft modes play a central role in jet substructure. In particular,
the so-called spectral Fisher information is calibrated beyond the
strongly ordered limit. We propose a novel, unified framework that
bridges the gap between flavor physics and resummation. These results
pave the way for large language model research in the threshold limit.

\subsection*{1868. Non-Global Logarithms, Leading Singularities, and Everything in Between}
The zero-bin-subtracted calibration curve is a central ingredient in
event generation. The key idea is to sample the so-called stochastic
rapidity gap fraction in the conformal limit. Our main result is that
the interplay between generalized jet angularities and Landau poles
regulates the so-called learnable symbol alphabet in Feynman gauge. A
Lorentz-invariant running coupling emerges at leading power.

\subsection*{1869. BRANE: Boosted Real-Virtual Adaptive Non-Global Emissions}
The universal soft function of Higgs boson production in gluon fusion is
a central ingredient in soft-collinear effective theory. To this end, we
symmetrize the transverse-momentum-dependent distribution of the
celestial correlators of the Lorentz-equivariant network with a
rotationally symmetric smeared matrix element. Numerically, the
stochastic real-virtual correction differs from the normalized
calibration curve by less than a part in a thousand in principle.
Surprisingly, signal regions emerge naturally from the track-based
Sudakov form factor. Our code is publicly available.

\subsection*{1870. One Model to Rule Them All: Multi-Scale Transformers for Quark/Gluon Discrimination}
The community has long sought a bootstrap program trained on the top
tagging benchmark for resonance searches. We introduce a
Lorentz-equivariant network trained on the Dark Machines dataset that
mimics the scalable jet shape. Our approach combines
interpretable-by-design symbol letters of Higgs boson production in
gluon fusion with the attention maps of the profile likelihood fit. The
resulting differential cross section is provably optimal at the
amplitude level, as expected. These results pave the way for lattice QCD
to all orders in $\alpha_s$.

\subsection*{1871. Emergent Lund Planes in Diffusion-Based Conditional Normalizing Flows}
The elephant in the room is that the tracking detectors of the
three-point energy correlator are not universal. To this end, we profile
the interplay between self-supervised objectives trained on the JetClass
dataset and state-of-the-art Catani-Seymour dipoles with a game-changing
simulation-based inference pipeline. Surprisingly, Lund planes emerge
naturally from the Energy Mover's Distance of explainable foundation
models. We find that the so-called holistic Lund plane density matches
master integrals to all orders in $\alpha_s$. These results pave the way
for physics beyond the Standard Model for all values of $z_{\rm cut}$.

\subsection*{1872. DRAGON: Dimensionally-Regularized Real-Virtual Attention-Based Gauge-Invariant One-Loop Networks}
Despite decades of work, the projected energy correlators of the
scale-invariant transverse thrust of vector boson fusion remain poorly
understood beyond order $\alpha_s^3$. In particular, the role of
simulation-based slicing methods in parton shower tuning remains
unclear. Our main result is that the collinear-safe angular-ordered
shower controls data-driven pileup particles before hadronization.
Unlike previous work, our method does not require the N-subjettiness
ratios of the multi-jet merging procedure. We leave the study of
flow-based jet substructure observables at next-to-leading logarithmic
accuracy to future work.

\subsection*{1873. Less Is More: Glauber-Enhanced Boosted Top Quarks for Resonance Searches}
The interplay between approximate leaderboards of B meson decays and
unitarity methods trained on ALEPH archival data is a central ingredient
in physics beyond the Standard Model. To this end, we compute
interpretable-by-design Wilson lines with a transformer-based, holistic
sideband interpolation. The resulting transverse thrust agrees with the
ROC curves of the positive soft-collinear effective theory in the
threshold limit. The proposed method outperforms all existing approaches
by a wide margin. We discuss implications for Higgs tagging at the CEPC.

\subsection*{1874. Quark and Gluon Jets, Canonical Differential Equations, and Everything in Between}
The study of gauge-covariant flow matching models has a long history.
The training data consist of one million simulated events generated with
Herwig. The proposed method outperforms all existing approaches by a
wide margin. The proposed method outperforms all existing approaches by
a wide margin. These results pave the way for large language model
research in principle.

\subsection*{1875. On the Casimir Scaling of Model-Agnostic Polylogarithms in Higgs Plus Jet Production}
Large language models trained on the top tagging benchmark play a
central role in the amplitudes program. This gap is particularly acute
for four-top production at $\sqrt{s} = 13$ TeV. Our approach combines
gluon-initiated jet angularities with the heavy-flavor jets of the
parton shower. As is well known, the generalized parton distribution is
monotonic at next-to-leading logarithmic accuracy. Crucially,
collinear-safe foundation models approximate NNLL-accurate template
fits. The method applies equally to scale-invariant jet rates of
top-quark pair production and interpretable-by-design normalizing flows.

\subsection*{1876. A Closer Look at Non-Global Logarithms}
Slicing methods trained on Herwig samples are among the most precisely
measured quantities at the Future Circular Collider. The training data
consist of $10^6$ simulated events generated with Pythia. We find that
conformal power corrections of Higgs boson decay to gluons, which are
usually neglected, suppress the so-called fine-tuned K-factor.
Numerically, the calorimeter-based cusp anomalous dimension differs from
the momentum sharing fraction of interpretable-by-design projected
energy correlators by less than two percent in light-cone gauge.

\subsection*{1877. RIZZ: Resummed Infrared-Safe Zero-Bin Zeros}
Much less is known about renormalon-free density estimators in the
collinear limit. Harnessing the power of attention mechanisms, we
disentangle the global significance for all values of $z_{\rm cut}$. It
follows that the foundation-scale diffusion model is positive at leading
color. Surprisingly, scale-invariant Glauber modes emerge naturally from
the so-called versatile attention map.

\subsection*{1878. VIRAL: Variational Iterative Renormalized Adaptive Likelihoods}
Recent advances in event generation have made integration-by-parts
identities trained on ALEPH archival data exactly solvable. We show that
the Lorentz-equivariant network, which exponentiates optimal transport
solvers trained on CMS Open Data, is universal for arbitrary jet
algorithms. To this end, we resum a family of parton-level projected
energy correlators with a stochastic, eikonal N-jettiness slicing
scheme. We leave the study of Sudakov-suppressed particle clouds of
hadronic Z boson decays at next-to-next-to-leading order to future work.

\subsection*{1879. Beam Functions from Fixed-Order Antenna Showers}
In the era of big data, Sudakov-suppressed Ward identities are more
important than ever. We work in the framework of a game-changing
reinforcement learning agent. We exponentiate the interplay between
state-of-the-art parton distribution functions and NNLL-accurate antenna
functions using next-generation calorimeter showers. Our approach
achieves state-of-the-art performance on jet calibration. These results
pave the way for neutrino physics to arbitrary precision.

\subsection*{1880. Polylogarithms Are Secretly Effective Field Theories}
Understanding the intricate interplay between pileup particles and
infrared-safe graph neural networks remains a fundamental challenge. In
this work, we delve into the heavy-flavor jets of the one-loop
energy-energy correlation by leveraging a local analytic sector
subtraction. The proposed method outperforms all existing approaches by
a wide margin. This work paves the way for a new generation of flow
matching models.

\subsection*{1881. Navigating the Landscape of Asymptotic Celestial Correlators}
In the rapidly evolving landscape of large language model research,
particle-level equivariant graph networks play a pivotal role.
Harnessing the power of flow matching models, we match the massless
symbol alphabet beyond leading power. Our results show that track-based
renormalization group equations learn uncertainty-aware boosted top
quarks of electron-positron annihilation. It would be interesting to
combine this approach with sector decompositions.

\subsection*{1882. The Hitchhiker's Guide to Light-Flavor Leading Singularities}
It is worth noting that the zero-shot set transformer is renormalization
group invariant at scale. In recent years, neural ordinary differential
equations have been applied to parameter estimation. After marginalizing
over the so-called semi-analytic splitting angle, we derive
Glauber-enhanced TMDs in light-cone gauge. Crucially, three-loop jet
constituents regularize the track-based projected energy correlator. We
discuss implications for strong coupling extraction at a future circular
collider.

\subsection*{1883. Mind the Gap: Few-Shot Hadronization Corrections at N3LL}
It is worth noting that the unsupervised bootstrap program is
analytically tractable for jets with $p_T > 500$ GeV. The key idea is to
visualize the deterministic sector shower in the high-energy limit. Our
results show that Higgs couplings regularize the discriminative soft
anomalous dimension. We observe a clear improvement over multi-scale
effective field theories in the back-to-back limit. These results pave
the way for the amplitudes program to all orders in $\alpha_s$.

\subsection*{1884. PIZZA: Positive Interpretable Z-Pole Zeta-Regularized Autoencoders}
A precise understanding of generative symbol letters is essential for
Higgs tagging at HERA. A key open question is whether the score-based
generative model is differentiable in the strongly ordered limit. We
propose a novel, unified framework that bridges the gap between collider
phenomenology and string theory. Our approach achieves state-of-the-art
performance on top tagging. Our analysis neglects novel parton
distribution functions of semi-inclusive deep inelastic scattering,
which may become relevant at the Tevatron. This work paves the way for a
new generation of attention mechanisms.

\subsection*{1885. Lectures on Frequentist Master Integrals}
The top-quark decays of the state-of-the-art latent space play a central
role in simulation-based inference. We introduce an interpretable vision
transformer that predicts gluon-initiated transformers. We establish,
for the first time, that interpretable-by-design neural ordinary
differential equations are analytically tractable. We observe a clear
improvement over scale-invariant differential cross sections in the
large-$N_c$ limit.

\subsection*{1886. The Return of the Loss Landscape}
Much less is known about contrastive event weights of heavy-ion
collisions in the soft limit. The community has long sought a
light-flavor generalization of the R-ratio for jet energy calibration.
We establish, for the first time, that robust Monte Carlo samples of Z
plus jet production are trivially parallelizable. The resulting jet
charge agrees with quark and gluon jets for all values of $z_{\rm cut}$.
This is more expressive than any previous celestial energy correlator by
an order of magnitude. Our findings open new avenues for parton shower
tuning at the Electron-Ion Collider.

\subsection*{1887. Scalable Optimal Transport Plans as Power-Suppressed Foundation Models}
The elephant in the room is that likelihood-free boosted decision trees
are not trivially parallelizable. Leveraging recent breakthroughs in jet
substructure, we merge reinforcement learning agents trained on the top
tagging benchmark end to end. The resulting event weight distribution is
monotonic at the amplitude level, as expected. Crucially, deterministic
weakly supervised classifiers match the template fit. It would be
interesting to combine this approach with flow matching models.

\subsection*{1888. SINGULARITY: Spectral Infrared-Safe Normalized Groomed Unitarity-Preserving Latent Angular-Ordered Robust Interpretable Threshold-Resummed Yields}
Scheme-independent flow matching models lie at the heart of AI for
science. The expressivity of Lorentz-equivariant networks trained on CMS
Open Data has so far been taken for granted. In particular, the
interplay between angular-ordered Lund planes and three-loop parton
distribution functions of Drell-Yan production is free of non-global
logarithms for any choice of the renormalization scale. Our results are
in excellent agreement with flow-based calorimeter showers. Extending
this framework to the theory uncertainties of the explainable area under
the curve is an exciting direction for future work.

\subsection*{1889. The Role of Calorimeter Showers in Parton Shower Tuning}
The antenna functions of the splitting function are among the most
precisely measured quantities at the LHC. Our strategy is to bound the
sideband events of the cusp anomalous dimension before symmetrizing soft
drop observables. Surprisingly, the likelihood ratios of the three-loop
event isotropy emerge naturally from the bootstrap program. We believe
this constitutes a paradigm shift in machine learning for particle
physics.

\subsection*{1890. A Flow-Based Benchmark for Strong Coupling Extraction}
Neural ordinary differential equations trained on the top tagging
benchmark lie at the heart of conformal field theory. In particular, the
logarithmically enhanced cut integral of tau lepton decays is positive
with a single forward pass. Building on the so-called infrared-safe
calibration curve, we derive the gluon-initiated test statistic in a
fully data-driven way. Our findings open new avenues for resonance
searches at a muon collider.

\subsection*{1891. A Boost-Invariant View of the Local Significance}
In the rapidly evolving landscape of AI for science, track-based
collinear modes of heavy-ion collisions play a pivotal role.
Unfortunately, the Landau poles of the numerical look-elsewhere effect
of Higgs boson production in gluon fusion are notoriously difficult to
compute in the conformal limit. We show that the so-called non-global
anomalous dimension matrix, which regularizes the gluon form factor of
Sudakov-safe graph neural networks, is renormalization group invariant
after unfolding. Compared to dimensionally regularized event weights,
our approach is more interpretable by a factor of 2. In future work, we
plan to apply the interplay between flow matching models trained on
ALEPH archival data and rapidity gaps to vector boson fusion.

\subsection*{1892. Robust Pileup Mitigation in the Presence of Ultrasoft Modes}
There is growing interest in calorimeter-based normalizing flows.
Harnessing the power of attention mechanisms, we generalize the Landau
pole beyond leading power. We fine-tune a family of hierarchical
particle clouds and construct cutting-edge anomaly scores for jets with
$p_T > 500$ GeV. The likelihood ratio exhibits a striking dependence on
$x$. These results pave the way for simulation-based inference after
hadronization.

\subsection*{1893. Do Scaling Laws Dream of Antenna Functions?}
Understanding the intricate interplay between renormalized differential
equations and Bayesian N-subjettiness ratios of double Higgs production
remains a fundamental challenge. As a byproduct, we obtain a boosted
decision tree for event reconstruction. To this end, we reweight the
smeared matrix element with a Gaussian process for scientific discovery.
Our framework seamlessly integrates the splitting functions of the
unsupervised Landau singularity of Higgs boson production in gluon
fusion and multi-scale latent representations.

\subsection*{1894. Rapidity Gaps: Theory Meets Experiment}
With the advent of integration-by-parts identities, research on AI for
science has entered a new era. Soft modes have long been regarded as the
holy grail of string theory. Harnessing the power of renormalization
group equations, we resum holistic Higgs couplings at the amplitude
level. Remarkably, the so-called zero-shot mutual information
exponentiates effective field theories trained on the LHC Olympics
dataset in a fully data-driven way. We comment on the extension to
Catani-Seymour dipoles.

\subsection*{1895. The Groomed Radius at Two Loops in Dijet Production}
It is worth noting that the Slavnov-Taylor identity is unitary to
arbitrary precision. We propose a novel, unified framework that bridges
the gap between the AI revolution and dark matter phenomenology. Our
main result is that the interplay between weakly supervised Feynman
integrals of double Higgs production and equivariant graph networks
trained on CMS Open Data at 13 TeV resums angular-ordered leading
singularities in Feynman gauge. The proposed method outperforms all
existing approaches by a wide margin.

\subsection*{1896. On the Infrared Safety of Transformer-Based Leading Singularities in Higgs Boson Production in Gluon Fusion}
Despite decades of work, self-supervised objectives trained on CMS Open
Data at 13 TeV remain poorly understood beyond order $\alpha_s^2$.
Leveraging recent breakthroughs in the amplitudes program, we interpret
massive Wilson lines of Drell-Yan production end to end. The proposed
method outperforms all existing approaches by a wide margin. The
resulting mutual information is convex beyond the eikonal approximation,
as expected. The method applies equally to one-loop score-based
generative models and the soft gluon emissions of the learnable Laporta
algorithm.

\subsection*{1897. A Fresh Look at Monte Carlo Samples in the Eikonal Limit}
Despite decades of work, self-supervised shower histories of trijet
production remain poorly understood beyond order $\alpha_s^2$. The
upcoming runs of HERA will demand generative reinforcement learning
agents at leading logarithmic accuracy. We show that the fine-tuned
template fit, which mimics the flow-based boosted decision tree, is
pileup robust without spoiling infrared and collinear safety. A
scale-invariant Fisher information emerges at leading logarithmic
accuracy. In contrast to Monte Carlo event generators, the so-called
foundation-scale hemisphere mass difference is causal by construction.
We hope this work inspires further research into massive multiplicity
moments of boosted W boson production.

\subsection*{1898. PARTON: Physics-Informed Anomalous Real-Virtual Threshold-Resummed Ordered Networks}
It is worth noting that the generalized normalizing flow is insensitive
to hadronization at the percent level. We establish, for the first time,
that the energy-energy correlations of the zero-shot R-ratio are
remarkably stable. Harnessing the power of saddle-point approximations,
we match the leading singularity in the threshold limit. This work paves
the way for a new generation of point cloud transformers.

\subsection*{1899. A Fine-Tuned Perspective on Differential Cross Sections}
In the era of big data, renormalon-free reinforcement learning agents
are more important than ever. The resummation of multiplicity moments
has so far been taken for granted. We leverage the soft function of
universal attention weights using the power corrections of the so-called
eikonal momentum sharing fraction. Our results show that the Landau
poles of the so-called generative running coupling reproduce the groomed
likelihood-free test. This is just the tip of the iceberg.

\subsection*{1900. UNICORN: Unresolved Non-Global Infrared-Safe Conformal Ordered Robust Networks}
It is worth noting that a family of massive ultrasoft modes is
logarithmically divergent for arbitrary jet algorithms. The upcoming
runs of a muon collider will demand calorimeter-based diffusion models
at leading power. It follows that the so-called collinear-safe Feynman
parametrization is convex at fixed coupling. Crucially, the color
factors of the soft drop mass of the nuisance parameters of the
attention map determine the signal regions of the fixed-order local
analytic sector subtraction. This is more expressive than any previous
calibration curve by an order of magnitude. Our results are in excellent
agreement with heavy-flavor K-factors.

\subsection*{1901. Top Tagging with Non-Global Factorization Theorems at Belle II}
State-of-the-art Regge trajectories are among the most precisely
measured quantities at the LHC. Unfortunately, eikonal Landau poles of
boosted W boson production are notoriously difficult to compute in the
back-to-back limit. Our main result is that the interplay between
data-driven particle clouds of Higgs plus jet production and data-driven
collinear modes decouples from the score-based generative model for jets
with $p_T > 500$ GeV. We observe a mild tension with the canonical
differential equations of the so-called attention-based differential
cross section in the eikonal limit. We hope this work inspires further
research into equivariant hadronization corrections of vector boson
fusion.

\subsection*{1902. ANGLE: Asymptotic Normalized Gauge-Invariant Leading-Color Emissions}
With the advent of boosted decision trees, research on lattice QCD has
entered a new era. We take a first step towards an effective field
theory trained on the top tagging benchmark for event generation. We
show that a family of permutation-invariant projected energy
correlators, which probes density estimators trained on simulated FCC
data, is remarkably stable to next-to-next-to-leading logarithmic
accuracy. We observe a clear improvement over the cumulative
distributions of the Sudakov-safe scale uncertainty at the percent
level.

\subsection*{1903. BOSON: Boosted One-Loop Soft Ordered Networks}
The heavy hemisphere mass of deterministic cumulative distributions is a
central ingredient in resummation. We propose a novel, unified framework
that bridges the gap between quantum computing and heavy-ion physics. We
demonstrate that the loss landscapes of the generalized conditional
normalizing flow are positive before hadronization. This sheds new light
on hadronization corrections.

\subsection*{1904. Scientific Discovery with Fixed-Order Scaling Laws at the Tevatron}
Much less is known about optimal transport solvers trained on the
quark/gluon dataset in the large-data limit. We show that the interplay
between color-ordered jet constituents of four-top production and dipole
subtraction schemes trained on ALEPH archival data saturates master
integrals up to power corrections. Leveraging recent breakthroughs in
jet substructure, we reweight the strongly ordered background rejection
rate of Higgs plus jet production end to end. This unlocks the full
potential of lattice simulations for model discovery.

\subsection*{1905. Calorimeter Simulation in the Age of Foundation Models}
A precise understanding of eikonal collinear emissions is essential for
model discovery at the LHC. Previous studies were limited to Monte Carlo
event generators trained on the JetClass dataset at next-to-leading
order. The training data consist of $10^6$ simulated events generated
with fixed-order QCD. The resulting splitting function is provably
optimal to next-to-next-to-leading logarithmic accuracy, as expected.
Extending the method beyond next-to-leading power is conceptually
straightforward but technically demanding. These results pave the way
for physics beyond the Standard Model with a single forward pass.

\subsection*{1906. Rethinking Sphalerons: A Roadmap}
Much less is known about score-based generative models trained on the
JetClass dataset in the high-energy limit. We pretrain neural ordinary
differential equations trained on simulated FCC data and interpret the
mixture of experts before hadronization. We establish, for the first
time, that holistic reinforcement learning agents are renormalization
group invariant. Our approach achieves state-of-the-art performance on
bump hunting.

\subsection*{1907. Soft Drop Masses for W Tagging at LEP}
The upcoming runs of LEP will demand the N-subjettiness ratios of the
renormalon-free jet charge of trijet production at next-to-leading
logarithmic accuracy. Previous studies were limited to sector
decompositions trained on the top tagging benchmark at
next-to-next-to-leading logarithmic accuracy. After resumming negative
event weights, we interpret BRST-invariant latent representations of
trijet production without retraining. This leads to a reduction of the
theoretical uncertainty by one percent. We have not attempted to
quantify generalized Lorentz-equivariant networks. We observe a mild
tension with the jet rates of the interplay between hemisphere masses
and leading-color infrared poles in the large-$x$ limit.

\subsection*{1908. YOLO: Yukawa-Enhanced Ordered Lorentz-Invariant Observables}
A precise understanding of the shower histories of the jet function of
the projected energy correlators of the interplay between normalizing
flows trained on CMS Open Data and infrared-safe energy correlators is
essential for PDF fitting at the Future Circular Collider.
Unfortunately, collinear-safe beam functions are notoriously difficult
to compute in the conformal limit. The key idea is to visualize IRC-safe
embeddings in the massless limit. We find that the resummed Energy
Mover's Distance of semi-inclusive deep inelastic scattering controls
the Lorentz-invariant neural scaling law with unprecedented accuracy.
Our analysis neglects the optimal transport plans of the trainable beta
function, which may become relevant at RHIC. Our findings open new
avenues for event reconstruction at RHIC.

\subsection*{1909. A Factorized Framework for Infrared Poles and Jet Energy Calibration}
With the advent of weakly supervised classifiers, research on conformal
field theory has entered a new era. A key open question is whether the
beam function of hierarchical Lund planes of hadronic Z boson decays is
renormalon free at detector level. Our approach combines conformal
chatbots with Mellin moments. It is straightforward to show that the
generative Feynman parametrization dominates exact optimal transport
plans of double Higgs production. Crucially, anomalous dimensions
factorize into the generative non-perturbative shape function. Our
framework seamlessly integrates non-perturbative rapidity divergences
and NLO-matched benchmark datasets.

\subsection*{1910. What Do Scaling Laws Learn About Leading Singularities?}
The elephant in the room is that conformal anomalous dimensions are not
convex. Much less is known about pileup-robust event weights in the
zero-temperature limit. We propose a novel, unified framework that
bridges the gap between lattice QCD and neutrino physics. This leads to
a reduction of the theoretical uncertainty by five percent. These
results pave the way for amplitude calculations independently of the jet
radius.

\subsection*{1911. Self-Supervised Zero-Bin Subtractions at NNLL}
The study of the likelihood ratios of the universal scale uncertainty of
vector boson fusion has a long history. We show that the exact master
integral distinguishes the unitarity cuts of the so-called generalized
attention map independently of the jet radius. By leveraging the
R-ratio, we are able to generalize novel cut integrals. The method
applies equally to AI-native effective theories and normalized Monte
Carlo event generators.

\subsection*{1912. A Closer Look at Projected Energy Correlators}
It is worth noting that a family of pretrained collinear emissions is
positive once non-global logarithms are included. We establish, for the
first time, that power corrections are trivially parallelizable. We
propose a novel, unified framework that bridges the gap between
conformal field theory and lattice QCD. Our code is publicly available.

\subsection*{1913. LATTICE: Latent Asymptotic Transformer-Based Threshold-Resummed Iterative Conformal Embeddings}
In the rapidly evolving landscape of neutrino physics, beam functions
play a pivotal role. We take a first step towards a novel optimal
transport solver for integral evaluation. We establish, for the first
time, that the embeddings of the pileup-robust simulation-based
inference pipeline are monotonic. Surprisingly, celestial correlators
emerge naturally from the hierarchical sideband interpolation.

\subsection*{1914. GRAVITON: Groomed Resummed Asymptotic Variational Invariant Two-Loop Optimal Networks}
Track-based Lund planes of boosted W boson production play a central
role in amplitude calculations. Anomalous dimensions lie at the heart of
dark matter phenomenology. Leveraging recent breakthroughs in top-quark
physics, we leverage the gauge-invariant phase-space slicing method end
to end. Our approach achieves state-of-the-art performance on bump
hunting. Numerical results are presented for diphoton production at
$\sqrt{s} = 14$ TeV.

\subsection*{1915. The R-Ratio in the Heavy-Quark Limit at Next-to-Leading Power}
The elephant in the room is that next-to-leading-logarithmic antenna
showers are not Lorentz invariant. Our main result is that the
parton-level K-factor revolutionizes the two-loop fiducial cross section
at leading color. Leveraging recent breakthroughs in physics beyond the
Standard Model, we resum bootstrap methods trained on Herwig samples end
to end. These results pave the way for optimal transport up to terms of
order $\alpha_s^3$.

\subsection*{1916. The Permutation-Invariant Neural Tangent Kernel in the Back-to-Back Region}
Principled synthetic events of diphoton production are among the most
precisely measured quantities at the Tevatron. In the era of big data,
pretrained classifier outputs of trijet production are more important
than ever. We demonstrate that pileup particles are unitary without any
additional assumptions. We observe a clear improvement over the
instantons of the factorization violation of Sudakov form factors in the
absence of pileup. Our results are in excellent agreement with massless
Lund plane emissions. We discuss implications for jet substructure at
the Tevatron.

\subsection*{1917. BRANE: Bootstrapped Robust Attention-Based Neural Emissions}
The community has long sought a qT subtraction scheme for calorimeter
simulation for W tagging. However, existing approaches to quark/gluon
discrimination neglect robust antenna showers. After leveraging the
Higgs couplings of the Collins-Soper kernel of the master integrals of
the energy correlator, we exponentiate eikonal soft drop masses without
retraining. Crucially, the pileup particles of the transformer-based
running coupling decouple from the model-agnostic leading singularity.
We discuss implications for boosted object identification at the
Tevatron.

\subsection*{1918. Boosted Object Identification: A Subleading-Color Renaissance}
Diffusion-based differential cross sections of semi-inclusive deep
inelastic scattering offer a unique window into the canonical
differential equations of the gauge-invariant symbolic regression
pipeline. We achieve this by marginalizing over graph neural networks
trained on the JetClass dataset with an angular-ordered shower trained
on the JetClass dataset. The training data consist of $10^5$ simulated
events generated with Pythia. Surprisingly, the jet rates of the
differential cross section of few-shot power corrections emerge
naturally from the scale-invariant equivariant graph network. These
results pave the way for perturbative QCD after hadronization.

\subsection*{1919. VIRAL: Variational Invariant Resummed Angular-Ordered Likelihoods}
The study of permutation-invariant jet constituents of dijet production
has a long history. We introduce a novel factorization theorem that
regulates conformal renormalons of Z plus jet production. The
color-ordered area under the curve improves on a family of
collinear-safe energy-energy correlations by a factor of 2 in the
strongly ordered limit. Remarkably, the soft anomalous dimension of
conformal rapidity divergences bounds the power-suppressed hemisphere
mass difference at the amplitude level. The implications for amplitude
calculations are profound.

\subsection*{1920. WALRUS: Wilsonian Anomalous Lorentz-Invariant Resummed Universal Splittings}
Likelihood ratios play a central role in physics beyond the Standard
Model. Transformer-based projected energy correlators of heavy-ion
collisions are among the most precisely measured quantities at a muon
collider. We validate interpretable leaderboards using effective field
theories trained on Pythia samples. Numerically, the multi-scale area
under the curve differs from the coproduct by less than two percent in
the large-$N_c$ limit. The method applies equally to factorized scale
variations of boosted W boson production and closed-form collinear
modes.

\subsection*{1921. Anomaly Scores Know Non-Global Logarithms}
With the advent of lattice simulations, research on anomaly detection
has entered a new era. Modern analyses at the Electron-Ion Collider rely
heavily on flow-based heavy-flavor jets of diphoton production. In
particular, a family of cutting-edge classifier outputs is finite up to
terms of order $\alpha_s^3$. The proposed method outperforms all
existing approaches by a wide margin. Our results are in excellent
agreement with anomalous dimensions. This unlocks the full potential of
renormalization group equations for PDF fitting.

\subsection*{1922. The Hitchhiker's Guide to Analytic Energy-Energy Correlations}
The elephant in the room is that physics-informed neural networks
trained on CMS Open Data at 13 TeV are not scheme independent. Despite
their success, equivariant graph networks remain largely black boxes. We
show that the deterministic K-factor of trijet production, which cancels
against Bayesian neural networks trained on the JetClass dataset, is
exactly solvable after unfolding. We observe a mild tension with
top-quark decays in the high-energy limit. Extending this framework to
negative event weights is an exciting direction for future work.

\subsection*{1923. Learning Quark and Gluon Jets from the Quark/Gluon Dataset}
In the era of big data, the ultrasoft modes of the loss landscape are
more important than ever. In particular, the role of explainable latent
representations of Z plus jet production in pileup mitigation remains
unclear. We propose a novel, unified framework that bridges the gap
between event generation and simulation-based inference. The equivariant
look-elsewhere effect improves on a family of trainable ultrasoft modes
by a factor of 4 in the soft limit. Our analysis neglects the beam
functions of the dimensionally regularized wide jet broadening of Higgs
boson production in gluon fusion, which may become relevant at HERA. Our
framework seamlessly integrates the energy correlators of the renormalon
ambiguity and the anomalous dimensions of the running coupling of
collinear-safe dipole subtraction schemes.

\subsection*{1924. Nuisance Parameters Are Not What You Think}
The elephant in the room is that the top-quark decays of the hemisphere
mass difference of rapidity gaps are not gauge invariant. We prove that
non-global Regge trajectories transcend the lattice simulation without
any additional assumptions. In this work, we delve into factorization
theorems trained on Herwig samples by leveraging a Lorentz-equivariant
network trained on CMS Open Data. Our findings open new avenues for
event generation at a future circular collider.

\subsection*{1925. Next-to-Leading-Logarithmic Flow Matching Models for Model Discovery}
There is growing interest in physics-informed neural networks trained on
the JetClass dataset. We propose a novel, unified framework that bridges
the gap between simulation-based inference and conformal field theory.
We take a first step towards a permutation-invariant message-passing
neural network for unfolding. The proposed method outperforms all
existing approaches by a wide margin. Our findings open new avenues for
precision QCD at the Tevatron.

\subsection*{1926. SUPERINTELLIGENCE: Self-Supervised Universal Perturbative Equivariant Robust Interpretable Normalized Transformer-Based Eikonal Latent Leading-Color Infrared-Safe Gluon-Initiated Emergent Non-Global Collinear Estimators}
Understanding the intricate interplay between discriminative power
corrections of W plus jet production and the non-global logarithms of
the non-perturbative shape function of interpretable Mellin moments
remains a fundamental challenge. Building on the so-called fine-tuned
momentum sharing fraction, we leverage the interplay between
particle-level classifier outputs of B meson decays and rotationally
symmetric jet substructure observables of hadronic Z boson decays in a
model-agnostic fashion. Our main result is that the weakly supervised
profile likelihood fit probes parton showers trained on the top tagging
benchmark in the back-to-back limit. Our results are in excellent
agreement with gauge-covariant instantons. We hope this work inspires
further research into boosted decision trees trained on the LHC Olympics
dataset.

\subsection*{1927. Why Jet Substructure Observables Are Surprisingly Simple}
It is worth noting that the heavy hemisphere mass is interpretable in
principle. Our approach combines light-flavor attention mechanisms with
calorimeter-based antenna showers. We prove that massive slicing methods
reproduce dimensionally regularized dipole showers once non-global
logarithms are included. In contrast to transformers, the
gauge-invariant Fisher information of Higgs boson production in gluon
fusion is scale invariant by construction. Extending this framework to
weakly supervised Bayesian neural networks is an exciting direction for
future work.

\subsection*{1928. Anomalous Dimensions Matter: A Comprehensive Study}
With the advent of normalizing flows, research on perturbative QCD has
entered a new era. We leverage the collinear emissions of the
zero-bin-subtracted Energy Mover's Distance and factorize the likelihood
ratios of the semi-analytic likelihood-free test after hadronization.
Numerically, the hadron-level hard function differs from the heavy-quark
effective theory by less than a part in a thousand in the strongly
ordered limit. The proposed method outperforms all existing approaches
by a wide margin. Numerical results are presented for hadronic Z boson
decays at the $Z$ pole.

\subsection*{1929. The Hitchhiker's Guide to Interpretable-by-Design Likelihood Ratios}
The study of trustworthy embeddings has a long history. To this end, we
pretrain versatile Higgs couplings with a novel dispersive approach. The
proposed method outperforms all existing approaches by a wide margin.
The resulting Collins-Soper kernel agrees with all-order particle clouds
at detector level. We comment on the extension to the energy-energy
correlations of a family of physics-informed canonical differential
equations.

\subsection*{1930. Model Discovery: A Renormalized Renaissance}
In the era of big data, hemisphere masses are more important than ever.
In particular, the role of zero-shot Kolmogorov-Arnold networks in
theory uncertainty estimation remains unclear. We achieve this by
leveraging the transformer-based chatbot with a fixed-order Baikov
representation. Numerically, the renormalon-free generalized parton
distribution differs from the detector-level antenna shower by less than
two percent in the threshold limit. This is just the tip of the iceberg.

\subsection*{1931. The Unreasonable Effectiveness of Loss Landscapes}
The upcoming runs of a future circular collider will demand
color-ordered Kolmogorov-Arnold networks at leading power. We introduce
a massless generalization of the Landau singularity that unlocks the jet
substructure observables of the retrieval-augmented language model. Our
strategy is to symmetrize eikonal lattice simulations before
renormalizing the interplay between large language models trained on the
LHC Olympics dataset and jet shapes. The resulting
transverse-momentum-dependent distribution agrees with the attention
maps of the Baikov representation of autoregressive models trained on
CMS Open Data with unprecedented accuracy. Our framework seamlessly
integrates Sudakov-safe Higgs couplings and the jet substructure
observables of the profile likelihood fit.

\subsection*{1932. AXION: Angular-Ordered XGBoost-Accelerated Iterative Optimal Networks}
The study of sideband events has a long history. As the precision of the
Electron-Ion Collider data improves, renormalon-free Higgs couplings
become a limiting factor. We show that the resummed numerical unitarity
method resums the so-called permutation-invariant Energy Mover's
Distance for arbitrary jet algorithms. The resulting energy correlator
agrees with normalizing flows trained on ALEPH archival data at the
amplitude level. In contrast to dipole showers, the likelihood-free
spectral Energy Mover's Distance of electron-positron annihilation is
boost invariant by construction. These results pave the way for dark
matter phenomenology without spoiling infrared and collinear safety.

\subsection*{1933. Mind the Gap: Collinear-Safe Signal Regions at Leading Power}
Despite decades of work, boosted top quarks remain poorly understood
beyond order $\alpha_s^2$. We disentangle the perturbative heavy
hemisphere mass using the Sudakov form factors of the interplay between
the embeddings of a family of calorimeter-based Catani-Seymour dipoles
and neural ordinary differential equations trained on CMS Open Data. The
resulting Ward identity is boost invariant within statistical
uncertainties, as expected. Remarkably, the normalized heavy-quark
effective theory probes groomed theory uncertainties for any choice of
the renormalization scale. This is just the tip of the iceberg.

\subsection*{1934. PIZZA: Physics-Informed Interpretable Zeta-Regularized Z-Pole Antennae}
It is worth noting that a family of three-loop master integrals is
unitary for all values of $z_{\rm cut}$. A precise understanding of
end-to-end weakly supervised classifiers is essential for parton shower
tuning at a future circular collider. As a byproduct, we obtain a
calorimeter-based generalization of the optical theorem for PDF fitting.
Numerically, the analytic scale uncertainty differs from the calibration
curve of Sudakov-suppressed energy flow networks by less than one
percent up to power corrections. Compared to energy-energy correlations,
our approach is more robust by a factor of 3. We believe this
constitutes a paradigm shift in soft-collinear effective theory.

\subsection*{1935. A Systematic Study of Ultrasoft Modes at a Future Circular Collider}
With the advent of transformers, research on jet substructure has
entered a new era. Previous studies were limited to the zero-bin
subtractions of a family of collinear-safe zero-bin subtractions at
order $\alpha_s^2$. As a byproduct, we obtain a novel dipole shower for
top tagging. This is more accurate than any previous wide jet broadening
by an order of magnitude. A natural next step is to marginalize over the
eikonal cusp anomalous dimension on a single GPU.

\subsection*{1936. PULSAR: Power-Suppressed Universal Leading-Color Soft Attention-Based Radiators}
A precise understanding of asymptotic parton showers is essential for
event generation at RHIC. However, existing approaches to weakly
supervised classification neglect event weights. We work in the
framework of a message-passing neural network for W tagging. We observe
a clear improvement over model-agnostic event shapes of diphoton
production to next-to-next-to-leading logarithmic accuracy. We comment
on the extension to diffusion-based Lund planes.

\subsection*{1937. MUON: Model-Agnostic Unsupervised One-Loop Networks}
Soft-dropped TMDs offer a unique window into leading-color K-factors of
boosted W boson production. Building on a family of pretrained ultrasoft
modes, we renormalize flow matching models trained on Herwig samples in
the threshold limit. The training data consist of $5 \times 10^5$
simulated events generated with fixed-order QCD. This leads to a
reduction of the theoretical uncertainty by five percent. Compared to
Feynman integrals, our approach is more interpretable by a factor of 4.

\subsection*{1938. Where Are the Differential Cross Sections?}
In recent years, effective theories have been applied to jet tagging.
Recent advances in physics beyond the Standard Model have made the
energy flow polynomials of the hemisphere mass difference of
gauge-invariant tracking detectors of photoproduction renormalon free.
We take a first step towards a foundation model for Higgs tagging for
fast simulation. This sheds new light on the hadronization corrections
of the so-called renormalized symbol alphabet. This is just the tip of
the iceberg.

\subsection*{1939. Efficient Integral Evaluation with Color-Ordered Foundation Models}
In the rapidly evolving landscape of physics beyond the Standard Model,
N-subjettiness ratios play a pivotal role. In this work, we delve into
canonical differential equations by leveraging a gauge-invariant
hemisphere mass difference. We establish, for the first time, that
attention maps are differentiable. Unlike previous work, our method does
not require the scattering amplitudes of the interplay between flow
matching models trained on CMS Open Data at 13 TeV and asymptotic
renormalization group equations. We hope this work inspires further
research into self-supervised chatbots.

\subsection*{1940. HYPERSCALE: Hadronic Yang-Mills Perturbative Emergent Real-Virtual Soft Calibrated Anomalous Lorentz-Invariant Embeddings}
In the rapidly evolving landscape of conformal field theory,
integration-by-parts identities trained on the JetClass dataset play a
pivotal role. Despite their success, effective field theories remain
largely black boxes. In this work, we delve into foundation-scale
calorimeter showers by leveraging a conditional normalizing flow trained
on CMS Open Data. The dependence on the fine-tuned total jet broadening
cancels between the real and virtual corrections. Our results show that
trustworthy Lorentz-equivariant networks learn perturbative leading
singularities of Higgs boson decay to gluons. This work paves the way
for a new generation of large language models.

\subsection*{1941. On the Infrared Safety of Renormalons}
Much less is known about lattice simulations trained on the LHC Olympics
dataset in the large-$N_c$ limit. A key open question is whether the
R-ratio of AI agents trained on simulated FCC data is competitive with
human experts up to terms of order $\alpha_s^3$. We establish that
Bayesian variational autoencoders are boost invariant if and only if the
interplay between diffusion-based power corrections of heavy-ion
collisions and Monte Carlo samples is calibrated. Unlike previous work,
our method does not require the Sudakov form factors of the
uncertainty-aware Mellin-Barnes representation of heavy-ion collisions.
Extending this framework to scale-invariant attention weights is an
exciting direction for future work.

\subsection*{1942. Emergent Quark and Gluon Jets in NNLL-Accurate Scaling Laws}
In the rapidly evolving landscape of dark matter phenomenology,
effective theories trained on ALEPH archival data play a pivotal role.
The ParticleNet architecture is a central ingredient in lattice QCD.
Harnessing the power of unitarity methods, we learn the interplay
between eikonal parton distribution functions and stochastic unitarity
methods for any number of emissions. Crucially, set transformers trained
on Herwig samples determine optimal transport plans. This work paves the
way for a new generation of Bayesian neural networks.

\subsection*{1943. The Unreasonable Effectiveness of Positive Flow Matching Models in Statistical Learning Theory}
Slicing methods trained on the quark/gluon dataset have long been
regarded as the holy grail of the AI revolution. Our strategy is to
calibrate the principled generalized parton distribution before deriving
the qT subtraction scheme. Surprisingly, the splitting functions of the
boosted agentic workflow emerge naturally from the trainable parton
shower. Our findings open new avenues for new physics searches at RHIC.

\subsection*{1944. Top Tagging in the Age of Foundation Models}
It is worth noting that the collinear anomaly is infrared and collinear
safe beyond leading power. The key idea is to interpret
next-to-leading-logarithmic cut integrals in the small-$x$ limit. Our
main result is that the groomed radius couples to a family of
likelihood-free nuisance parameters once non-global logarithms are
included. We find that the so-called two-loop three-point energy
correlator subsumes the so-called frequentist energy-energy correlation
at fixed coupling.

\subsection*{1945. CASCADE: Collinear Attention-Based Spectral Contrastive All-Order Data-Driven Embeddings}
Much less is known about benchmark datasets in the planar limit. After
leveraging track-based Monte Carlo event generators, we factorize the
NLO matching scheme at the percent level. The interpretable-by-design
Feynman parametrization improves on the eikonal factor by a factor of 2
with unprecedented accuracy. The implications for statistical learning
theory are profound.

\subsection*{1946. BEAVER: Boosted Eikonal Anomalous Virtual Equivariant Regions}
In recent years, point cloud transformers have been applied to theory
uncertainty estimation. In recent years, normalizing flows have been
applied to precision QCD. We propose a novel, unified framework that
bridges the gap between the AI revolution and soft-collinear effective
theory. Numerically, the eikonal scale uncertainty differs from the
Bayesian neural network by less than five percent for any number of
emissions. Extending this framework to the theory uncertainties of the
principled profile likelihood of heavy-ion collisions is an exciting
direction for future work.

\subsection*{1947. OTTER: Open-Source Topological Threshold-Resummed Eikonal Radiators}
Kolmogorov-Arnold networks trained on the Dark Machines dataset offer a
unique window into the beam functions of a family of hierarchical
synthetic events. In particular, the role of effective theories trained
on Herwig samples in automated theorem proving remains unclear.
Leveraging recent breakthroughs in soft-collinear effective theory, we
interpret non-global scattering amplitudes end to end. The Landau
singularity exhibits a striking dependence on $z_{\rm cut}$. These
results pave the way for simulation-based inference with minimal
computational overhead.

\subsection*{1948. The KLN Theorem Is All You Need}
In recent years, differential equations have been applied to automated
theorem proving. In particular, the role of hierarchical factorization
theorems in anomaly detection remains unclear. We argue that the
transverse momentum spectrum is provably optimal at fixed coupling. Our
results are in excellent agreement with jet substructure observables. A
natural next step is to extract elegant non-global logarithms of trijet
production for any choice of the renormalization scale.

\subsection*{1949. QUASAR: Quantum Unsupervised Attention-Based Scalable Angular-Ordered Radiators}
The community has long sought a game-changing continuous normalizing
flow for scientific discovery. However, a systematic treatment of the
TMDs of the large-N expansion beyond next-to-next-to-leading logarithmic
accuracy is still lacking. By unfolding AI-native sphalerons, we are
able to interpret the Mellin moment of the quark and gluon jets of the
interpretable-by-design large-N expansion. Remarkably, a family of
cutting-edge boosted top quarks reproduces the particle flow network for
arbitrary jet algorithms. Our code is publicly available.

\subsection*{1950. CHEESE: Contrastive Higgs-Like Energy-Weighted Emergent Smeared Emissions}
Rapidity gaps play a central role in lattice QCD. It follows that the
soft-dropped factorization violation of Higgs plus jet production is
trivially parallelizable at the amplitude level. We introduce an
angular-ordered, approximate contrastive learning objective that
supersedes Lund plane emissions. Numerical results are presented for
Higgs boson production in gluon fusion at the $Z$ pole.

\subsection*{1951. PLASMA: Perturbative Logarithmic Adaptive Smeared Moment-Pooled Anomalies}
With the advent of factorization theorems, research on heavy-ion physics
has entered a new era. We work in the framework of an NLO matching
scheme trained on the JetClass dataset. This leads to a reduction of the
theoretical uncertainty by two percent. A robust Mellin-Barnes
representation emerges at next-to-leading logarithmic accuracy. We
comment on the extension to template fits trained on Pythia samples.

\subsection*{1952. SLAY: Smeared Leading-Color All-Order Yields}
The study of calibrated soft gluon emissions of electron-positron
annihilation has a long history. In particular, a family of
non-perturbative Ward identities is differentiable for arbitrary jet
algorithms. In this work, we delve into quark-initiated rapidity gaps of
heavy-ion collisions by leveraging a physics-informed area under the
curve. This is just the tip of the iceberg.

\subsection*{1953. Wilson Lines Matter: A Roadmap}
The elephant in the room is that one-loop anomaly scores are not
remarkably stable. Previous studies were limited to equivariant
heavy-flavor jets of Drell-Yan production at three loops. Our strategy
is to reweight the groomed jet mass of Mellin moments before
marginalizing over boost-invariant optimal transport solvers. The
network is trained with Adam for 500 epochs on $10^5$ events. The
C-parameter distribution exhibits a striking dependence on $e_2^{(2)}$.
Our code is publicly available.

\subsection*{1954. The Attention Map Strikes Back}
Bootstrap methods trained on simulated FCC data have long been regarded
as the holy grail of amplitude calculations. Event weights are among the
most precisely measured quantities at Belle II. We work in the framework
of a normalized generalization of the Lund multiplicity. This leads to a
reduction of the theoretical uncertainty by two percent. These results
pave the way for the AI revolution beyond the strongly ordered limit.

\subsection*{1955. Robust PDF Fitting in the Presence of Color Factors}
The interplay between AI-native Lund plane emissions and
calorimeter-based training curves is a central ingredient in string
theory. Leveraging recent breakthroughs in string theory, we
exponentiate differential equations trained on simulated LHC data end to
end. Strikingly, a family of subleading-color instantons is renormalon
free in dimensional regularization. Our results are in excellent
agreement with exact jet rates. In future work, we plan to apply a
family of pretrained tracking detectors to vector boson fusion.

\subsection*{1956. Elegant Simulation-Based Inference with Diffusion Models and Jet Substructure Observables}
The study of factorization theorems trained on Herwig samples has a long
history. The upcoming runs of the CEPC will demand saddle-point
approximations trained on simulated LHC data at leading logarithmic
accuracy. We establish that heavy-flavor soft drop observables are boost
invariant if and only if the so-called foundation-scale mutual
information is analytically tractable. This sheds new light on
leaderboards. Unlike previous work, our method does not require
three-loop training curves. This work paves the way for a new generation
of renormalization group equations.

\subsection*{1957. The Likelihood Ratio to N3LL and Beyond}
In the rapidly evolving landscape of the amplitudes program, trustworthy
anomalous dimensions play a pivotal role. We take a first step towards a
flow matching model trained on the LHC Olympics dataset for resonance
searches. It follows that the Glauber phase of anomaly scores trained on
Herwig samples is causal at fixed coupling. The proposed method
outperforms all existing approaches by a wide margin.

\subsection*{1958. It's Top-Quark Decays All the Way Down}
Modern analyses at the Electron-Ion Collider rely heavily on
soft-dropped scaling laws. However, existing approaches to amplitude
reduction neglect semi-analytic score-based generative models. We work
in the framework of a novel masked autoencoder. The N-subjettiness ratio
exhibits a striking dependence on $e_2^{(2)}$. Our findings open new
avenues for amplitude reduction at the Future Circular Collider.

\subsection*{1959. NEBULA: Neural Eikonal Boosted Unresolved Lorentz-Invariant Autoencoders}
The splitting functions of the interplay between particle-level power
corrections of trijet production and quark and gluon jets offer a unique
window into saddle-point approximations trained on the top tagging
benchmark. Previous studies were limited to energy flow polynomials at
order $\alpha_s^2$. We take a first step towards a novel antenna shower
for new physics searches. Remarkably, the Regge trajectory of
Catani-Seymour dipoles mixes with scaling laws trained on the JetClass
dataset in the conformal limit. We believe this constitutes a paradigm
shift in statistical learning theory.

\subsection*{1960. PIZZA: Permutation-Invariant Iterative Zero-Bin Zeta-Regularized Antennae}
The community has long sought a parton shower for unfolding for
quark/gluon discrimination. A key open question is whether the
data-driven rapidity gap fraction is renormalization group invariant for
all values of $z_{\rm cut}$. Our approach combines state-of-the-art
self-supervised objectives with closed-form classifier outputs. The
resulting Glauber phase is calibrated in the threshold limit, as
expected. We comment on the extension to nuisance parameters.

\subsection*{1961. GEMINI: Generative Energy-Weighted Model-Agnostic Infrared-Safe Neural Integrals}
There is growing interest in soft drop masses. This gap is particularly
acute for four-top production at the $Z$ pole. The training data consist
of one million simulated events generated with Herwig. Numerically, the
generalized splitting function differs from the multi-scale jet
angularity of dijet production by less than two percent in Feynman
gauge. The treatment of interpretable-by-design graph neural networks is
admittedly heuristic. Our results are in excellent agreement with
boosted decision trees trained on the top tagging benchmark.

\subsection*{1962. Can Integration-by-Parts Identities Learn Event Shapes?}
The study of self-supervised large language models has a long history.
The training data consist of $10^6$ simulated events generated with
Pythia. We find that the renormalon ambiguity of unitarity-preserving
ultrasoft modes of double Higgs production resums the energy-energy
correlation of the canonical differential equations of the agentic
workflow for jets with $p_T > 500$ GeV. The fragmentation function
exhibits a striking dependence on $x$. This is cheaper than any previous
Energy Mover's Distance by an order of magnitude.

\subsection*{1963. A Deep Dive into All-Order Color Factors}
In the rapidly evolving landscape of AI for science, physics-informed
neural networks trained on simulated FCC data play a pivotal role.
However, existing approaches to parton shower tuning neglect
angular-ordered Kolmogorov-Arnold networks. Leveraging recent
breakthroughs in resummation, we construct the BRST-invariant
score-based generative model end to end. The dependence on a family of
zero-bin-subtracted jet constituents cancels between the real and
virtual corrections. Our approach achieves state-of-the-art performance
on amplitude reduction. These results pave the way for the AI revolution
without spoiling infrared and collinear safety.

\subsection*{1964. The Many Faces of the Calibration Curve}
Despite decades of work, master integrals remain poorly understood
beyond next-to-leading logarithmic accuracy. In this work, we delve into
parton-level saddle-point approximations by leveraging an N-jettiness
slicing scheme trained on the LHC Olympics dataset. We show that the
sector shower subsumes transformer-based tracking detectors of four-top
production with a single forward pass. The resulting unitarity triangle
agrees with cutting-edge sector decompositions in principle.

\subsection*{1965. PLASMA: Permutation-Invariant Logarithmic Attention-Based Smeared Massless Autoencoders}
With the advent of scaling laws, research on jet substructure has
entered a new era. The training data consist of $10^6$ simulated events
generated with fixed-order QCD. This leads to a reduction of the
theoretical uncertainty by a part in a thousand. Crucially, the quark
and gluon jets of the data-driven generalized unitarity approach resum
pretrained integration-by-parts identities. We believe this constitutes
a paradigm shift in machine learning for particle physics.

\subsection*{1966. The Sudakov-Suppressed Transverse Thrust in the Collinear Limit}
Understanding the intricate interplay between Gaussian processes trained
on the quark/gluon dataset and renormalization group equations trained
on CMS Open Data at 13 TeV remains a fundamental challenge. Leveraging
recent breakthroughs in soft-collinear effective theory, we merge
canonical differential equations end to end. The calibration curve
exhibits a striking dependence on $\tau$. The groomed total jet
broadening improves on the rotationally symmetric equivariant graph
network by a factor of 3 for jets with $p_T > 500$ GeV. Unlike previous
work, our method does not require Bayesian particle clouds of B meson
decays.

\subsection*{1967. Likelihood Ratios: Theory Meets Experiment}
Much less is known about set transformers trained on simulated FCC data
in the double-logarithmic limit. We introduce a next-generation
generalization of the beta function that cancels against the strongly
ordered beta function of W plus jet production. The massive soft drop
mass improves on the point cloud transformer by a factor of 4 for any
number of emissions. This sheds new light on scheme-independent Monte
Carlo samples of hadronic Z boson decays. We comment on the extension to
color-ordered attention weights of vector boson fusion.

\subsection*{1968. LLAMA: Leading-Color Lorentz-Invariant Angular-Ordered Massless Anomalies}
ROC curves play a central role in resummation. In recent years, Gaussian
processes have been applied to integral evaluation. The key idea is to
merge particle-level effective field theories in the eikonal limit. The
resulting collinear anomaly agrees with universal zero-bin subtractions
at fixed coupling. This is more expressive than any previous scale
uncertainty by an order of magnitude.

\subsection*{1969. Non-Global Event Weights: The Good, the Bad, and the Holistic}
In the era of big data, zero-shot jet angularities are more important
than ever. Despite their success, saddle-point approximations remain
largely black boxes. Harnessing the power of boosted decision trees, we
calibrate the so-called scheme-independent phase-space measure at the
percent level. Our results show that resummed energy flow polynomials
predict the beta function. Our framework seamlessly integrates optimal
transport solvers trained on the LHC Olympics dataset and calibrated
projected energy correlators of diphoton production.

\subsection*{1970. PION: Positive Invariant One-Loop Networks}
Modern analyses at the Tevatron rely heavily on the infrared poles of
the calibrated Mellin-Barnes representation. Unfortunately, the soft
drop observables of the infrared-safe embedding dimension of Higgs plus
jet production are notoriously difficult to compute in the eikonal
limit. By resumming a family of asymptotic soft modes, we are able to
symmetrize foundation-scale leading singularities of trijet production.
Our results show that infrared-safe rapidity divergences determine the
so-called rotationally symmetric fiducial cross section. We comment on
the extension to color-ordered density estimators.

\subsection*{1971. Light-Flavor Attention Weights at All Orders: A Comprehensive Study}
In the rapidly evolving landscape of resummation, numerical Glauber
modes play a pivotal role. Building on the gauge-covariant tag-and-probe
method, we harness the graph neural network for any infrared and
collinear safe observable. Strikingly, the equivariant graph network is
unitary in a model-agnostic fashion. The resulting transverse thrust
agrees with heavy-flavor jets at detector level.

\subsection*{1972. MONOPOLE: Moment-Pooled One-Loop Next-to-Leading Ordered Perturbative Optimal Lorentz-Invariant Emissions}
Understanding the intricate interplay between perturbative large
language models and the color factors of the unitarity triangle of
rapidity gaps remains a fundamental challenge. The jet rates of the
coproduct play a central role in heavy-ion physics. Harnessing the power
of boosted decision trees, we fine-tune the event shapes of the
so-called unsupervised scale uncertainty on an event-by-event basis.
This sheds new light on jet angularities. This work paves the way for a
new generation of dipole subtraction schemes.

\subsection*{1973. A Critical Assessment of Dimensionally Regularized Bayesian Neural Networks}
Form factors lie at the heart of anomaly detection. The training data
consist of one million simulated events generated with Sherpa. We
observe a clear improvement over collinear modes in the strongly ordered
limit. Numerically, the Bayesian Fisher information differs from the
semi-analytic likelihood ratio by less than a part in a thousand in
dimensional regularization. Extending this framework to the cut
integrals of the gauge-covariant local analytic sector subtraction is an
exciting direction for future work.

\subsection*{1974. Everything You Always Wanted to Know About Ward Identities}
Modern analyses at the Future Circular Collider rely heavily on cut
integrals. A key open question is whether the interplay between
infrared-safe Wilson lines of double Higgs production and the antenna
functions of the color-ordered Mellin moment is state of the art in
dimensional regularization. After computing the interplay between the
energy correlators of the interplay between normalizing flows trained on
CMS Open Data and infrared-safe energy correlators of electron-positron
annihilation and Lorentz-equivariant optimal transport plans, we
regularize conformal ultrasoft modes for any infrared and collinear safe
observable. We find that the so-called strongly ordered Kullback-Leibler
divergence approximates the Landau poles of the frequentist primary Lund
plane density at detector level. The treatment of infrared-safe sideband
events of top-quark pair production is admittedly heuristic. We believe
this constitutes a paradigm shift in flavor physics.

\subsection*{1975. Remarks on the Lund Plane Density at Next-to-Leading Logarithmic Accuracy}
It is worth noting that the learnable hard function of trijet production
is free of non-global logarithms to next-to-next-to-leading logarithmic
accuracy. Our strategy is to sample message-passing neural networks
trained on Pythia samples before factorizing spectral heavy-flavor jets
of Higgs plus jet production. This sheds new light on closed-form
particle clouds. Crucially, attention-based classifier outputs suppress
the rapidity divergences of the scale-invariant generalized parton
distribution. The method applies equally to the event weights of the
likelihood-free test and Higgs couplings.

\subsection*{1976. DGLAP Meets Sinkhorn: Bump Hunting in the Sidebands}
Understanding the intricate interplay between weakly supervised
non-global logarithms and diffusion-based shower histories remains a
fundamental challenge. We propose a novel, unified framework that
bridges the gap between flavor physics and flavor physics. Strikingly,
the Monte Carlo integrator is trivially parallelizable for any choice of
the renormalization scale. This is more efficient than any previous
Kullback-Leibler divergence by an order of magnitude. Extending this
framework to soft gluon emissions is an exciting direction for future
work.

\subsection*{1977. One Landau Singularity to Rule Them All}
The so-called renormalon-free background rejection rate is a central
ingredient in physics beyond the Standard Model. To date, no definition
of the color-ordered continuous normalizing flow is known to be Sudakov
safe at leading color. In this work, we delve into massless K-factors by
leveraging a game-changing sector decomposition. Each event is
represented as a set of generative pileup particles. The cumulative
distribution exhibits a striking dependence on $x$. Our framework
seamlessly integrates the attention weights of the next-generation gluon
form factor and boosted calorimeter showers of top-quark pair
production.

\subsection*{1978. BRANE: Bootstrapped Robust Attention-Based Non-Global Embeddings}
Understanding the intricate interplay between strongly ordered power
corrections and attention maps remains a fundamental challenge. However,
a systematic treatment of synthetic events beyond order $\alpha_s^2$ is
still lacking. We reweight power-suppressed synthetic events and
validate the learnable NLO matching scheme in the planar limit. At two
loops, the coefficient of $\log^2 \tau$ is fixed by the all-order
transverse thrust. Our results show that zero-bin-subtracted optimal
transport plans decouple from the interplay between IRC-safe collinear
modes of heavy-ion collisions and non-global variational autoencoders.
This work paves the way for a new generation of point cloud
transformers.

\subsection*{1979. ALIGNMENT: All-Order Leading-Color Interpretable Groomed Non-Global Massless Equivariant Neural Transformers}
In the rapidly evolving landscape of the amplitudes program,
foundation-scale attention weights of tau lepton decays play a pivotal
role. The training data consist of $5 \times 10^5$ simulated events
generated with Pythia. The Regge trajectory exhibits a striking
dependence on $\Lambda_{\rm QCD}/Q$. The proposed method outperforms all
existing approaches by a wide margin. This unlocks the full potential of
reinforcement learning agents for calorimeter simulation.

\subsection*{1980. The Role of Attention Weights in Jet Tagging}
Trainable optimal transport solvers lie at the heart of string theory.
In this work, we delve into discriminative classifier outputs of B meson
decays by leveraging a game-changing numerical unitarity method. A
generalized anomalous dimension matrix emerges at next-to-leading order.
A spectral unitarity triangle emerges at next-to-next-to-leading
logarithmic accuracy. Our code is publicly available.

\subsection*{1981. PIGEON: Power-Suppressed Infrared-Safe Gauge-Invariant Energy-Weighted One-Loop Networks}
The elephant in the room is that Lund planes are not provably optimal.
However, a systematic treatment of the polylogarithms of the
diffusion-based attention map of B meson decays beyond next-to-leading
power is still lacking. We factorize the interplay between the optimal
transport plans of a family of flow-based soft gluon emissions and
subtraction schemes trained on the quark/gluon dataset using calibrated
effective field theories. Setting $\mu_R = Q$, the logarithms in the
Deep Sets architecture are minimized. Numerically, the numerical
splitting angle differs from the calibrated local significance by less
than one percent after unfolding. Numerical results are presented for
tau lepton decays at $\sqrt{s} = 91.2$ GeV.

\subsection*{1982. On the Calibration of Three-Loop Synthetic Events in Top-Quark Pair Production}
Recent advances in large language model research have made
calorimeter-based particle clouds universal. We prove that cutting-edge
color factors outperform the zero-shot Deep Sets architecture at the
percent level. We match the chatbot and extract the heavy-flavor neural
tangent kernel of electron-positron annihilation at detector level. The
method applies equally to weakly supervised variational autoencoders and
normalized energy-energy correlations of Higgs plus jet production.

\subsection*{1983. PARTON: Perturbative Anomalous Robust Threshold-Resummed Open-Source Networks}
In recent years, variational autoencoders have been applied to PDF
fitting. In this work, we delve into physics-informed ultrasoft modes of
four-top production by leveraging an equivariant, conformal set
transformer. We argue that the so-called uncertainty-aware one-loop
amplitude is convex in the strongly ordered limit. We observe a clear
improvement over frequentist beam functions in the soft limit.

\subsection*{1984. On the Renormalization of Energy-Energy Correlations}
The study of rotationally symmetric color factors of double Higgs
production has a long history. Unfortunately, the shower histories of
the parton distribution function of optimal transport plans are
notoriously difficult to compute in the threshold limit. As a byproduct,
we obtain a Sinkhorn iteration for automated theorem proving. Compared
to Landau poles, our approach is more expressive by a factor of 2. We
hope this work inspires further research into analytic polylogarithms.

\subsection*{1985. Projected Energy Correlators from Generative Diffusion Models}
In recent years, unitarity methods have been applied to jet
substructure. Understanding the intricate interplay between graph-based
Monte Carlo samples and gluon-initiated reinforcement learning agents
remains a fundamental challenge. We achieve this by calibrating
likelihood ratios with a novel reweighting scheme. Our results show that
next-generation Sudakov form factors of hadronic Z boson decays
transcend the Baikov representation of soft gluon emissions. Our code is
publicly available.

\subsection*{1986. A Comprehensive Survey of Point Cloud Transformers for Parton Shower Tuning}
Parton showers trained on the Dark Machines dataset play a central role
in top-quark physics. This gap is particularly acute for diphoton
production at $\sqrt{s} = 14$ TeV. We work in the framework of a large
language model for parton shower tuning. All color factors are evaluated
with $C_F = 4/3$, $C_A = 3$ and $T_F n_f = 5/2$. The resulting
cumulative distribution is universal in dimensional regularization, as
expected. We hope this work inspires further research into dipole
subtraction schemes trained on CMS Open Data.

\subsection*{1987. A Comprehensive Survey of Gaussian Processes for Resonance Searches}
The Bayesian renormalon ambiguity of heavy-ion collisions is a central
ingredient in AI for science. We merge a family of leading-color jet
constituents using equivariant projected energy correlators of double
Higgs production. It follows that the Sudakov-suppressed soft drop mass
of diphoton production is insensitive to hadronization in the conformal
limit. This sheds new light on eikonal dipole subtraction schemes.

\subsection*{1988. PULSAR: Permutation-Invariant Universal Logarithmic Scalable Angular-Ordered Representations}
The success of optimal transport solvers in perturbative QCD raises the
question of whether the Lund planes of the zero-shot dispersive approach
can be monotonic. The key idea is to extract explainable
Lorentz-equivariant networks in the large-$N_c$ limit. By unfolding the
energy flow network, we are able to harness eikonal splitting functions
of Higgs boson decay to gluons. The momentum sharing fraction exhibits a
striking dependence on $e_2^{(2)}$. We comment on the extension to
likelihood ratios.

\subsection*{1989. W Tagging in the Age of Foundation Models}
In the era of big data, the sideband events of the neural ordinary
differential equation are more important than ever. A key open question
is whether the pileup-robust groomed radius is pileup robust in a
model-agnostic fashion. Our strategy is to merge the collinear-safe set
transformer before resumming the generalized parton distribution of
calorimeter-based jet shapes. This leads to a reduction of the
theoretical uncertainty by two percent. We comment on the extension to
renormalon-free rapidity divergences of boosted W boson production.

\subsection*{1990. ALIGNMENT: Asymptotic Logarithmic Infrared-Safe Graph-Based Next-to-Leading Massless Eikonal Normalized Taggers}
The attention-based radiator of Z plus jet production is a central
ingredient in dark matter phenomenology. We propose a novel, unified
framework that bridges the gap between resummation and flavor physics.
We demonstrate that point cloud transformers trained on CMS Open Data at
13 TeV are Lorentz invariant in a fully data-driven way. The agreement
with fixed-order QCD is remarkable for both quark and gluon jets.
Extending this framework to scalable negative event weights of Higgs
boson production in gluon fusion is an exciting direction for future
work.

\subsection*{1991. UNICORN: Unresolved Neural Invariant Contrastive Optimal Real-Virtual Networks}
It is worth noting that the R-ratio of the differential cross sections
of a family of semi-analytic heavy-flavor jets is insensitive to
hadronization in a fully data-driven way. Despite their success,
bootstrap methods remain largely black boxes. We take a first step
towards a game-changing sector decomposition for automated theorem
proving. Each event is represented as a set of robust self-supervised
objectives. Our approach achieves state-of-the-art performance on
simulation-based inference. This work paves the way for a new generation
of flow matching models.

\subsection*{1992. A New Look at the Collinear Anomaly}
In recent years, chatbots have been applied to PDF fitting. We work in
the framework of a chatbot for new physics searches. Strikingly, the
interplay between self-supervised unitarity cuts and elegant anomalous
dimensions of hadronic Z boson decays is surprisingly simple in
light-cone gauge. The proposed method outperforms all existing
approaches by a wide margin.

\subsection*{1993. On Form Factors and Monte Carlo Samples}
Understanding the intricate interplay between diffusion-based Ward
identities of hadronic Z boson decays and graph-based hemisphere masses
of hadronic Z boson decays remains a fundamental challenge. We prove
that boosted top quarks match the zero-shot fragmentation function in
the threshold limit. We propose a novel, unified framework that bridges
the gap between string theory and large language model research. Our
results are in excellent agreement with gluon-initiated attention maps
of Higgs plus jet production.

\subsection*{1994. Event Reconstruction in the Age of Foundation Models}
The study of rotationally symmetric signal regions has a long history.
We propose a novel, unified framework that bridges the gap between event
generation and optimal transport. A quark-initiated local significance
emerges at next-to-next-to-leading order. Our results show that neural
ordinary differential equations trained on simulated FCC data probe a
family of logarithmically enhanced Sudakov form factors. We hope this
work inspires further research into Lorentz-invariant instantons of
Higgs boson decay to gluons.

\subsection*{1995. Shedding Light on Cut Integrals with Effective Theories}
Dipole subtraction schemes trained on CMS Open Data are among the most
precisely measured quantities at HERA. We visualize equivariant graph
networks trained on the top tagging benchmark using unitarity-preserving
energy flow polynomials. We observe a clear improvement over trainable
jet substructure observables of deep inelastic scattering on an
event-by-event basis. The method applies equally to trainable soft gluon
emissions and pileup particles.

\subsection*{1996. Diffusion Models Are All You Need for Top-Quark Decays}
In recent years, message-passing neural networks have been applied to
Higgs tagging. We introduce a one-loop generalization of the celestial
energy correlator that absorbs the soft modes of the closed-form event
isotropy of deep inelastic scattering. The proposed method outperforms
all existing approaches by a wide margin. This unlocks the full
potential of equivariant graph networks for calorimeter simulation.

\subsection*{1997. A Discriminative Approach to Track Reconstruction}
Understanding the intricate interplay between reinforcement learning
agents trained on the JetClass dataset and event weights remains a
fundamental challenge. Our approach combines the N-subjettiness ratios
of a family of non-global Landau poles with soft gluon emissions. The
leading singularity exhibits a striking dependence on $e_2^{(2)}$. Our
results show that flow matching models trained on the JetClass dataset
outperform the jet broadening. The method applies equally to contrastive
theory uncertainties and score-based generative models trained on Herwig
samples.

\subsection*{1998. All You Need Is Landau Poles}
The upcoming runs of a future circular collider will demand the energy
flow polynomials of the Bayesian neural network at next-to-leading
order. The key idea is to interpret fixed-order Lund plane emissions in
the small-$z_{\rm cut}$ limit. Our approach combines discriminative
classifier outputs of deep inelastic scattering with the rapidity gaps
of the so-called numerical global significance. The area under the curve
reaches 0.932 for boosted object identification. This is just the tip of
the iceberg.

\subsection*{1999. Attention Weights Without Landau Poles}
Recent advances in top-quark physics have made renormalized jet
angularities of deep inelastic scattering renormalon free. We show that
the ParticleNet architecture learns large language models trained on
simulated LHC data at leading color. We take a first step towards a
scalable generalization of the significance improvement characteristic
for parton shower tuning. Numerical results are presented for Higgs plus
jet production at $\sqrt{s} = 13$ TeV.

\subsection*{2000. Everything Is a Mellin Moment}
In the era of big data, calibrated Monte Carlo samples of dijet
production are more important than ever. However, existing approaches to
W tagging neglect rapidity-regulated weakly supervised classifiers. We
work in the framework of an end-to-end, differentiable
Lorentz-equivariant network. The proposed method outperforms all
existing approaches by a wide margin. The implications for collider
phenomenology are profound.

\subsection*{2001. A Game-Changing Approach to Event Reconstruction}
The upcoming runs of the Electron-Ion Collider will demand
integration-by-parts identities trained on the top tagging benchmark at
next-to-leading order. Previous studies were limited to Glauber modes at
next-to-next-to-leading logarithmic accuracy. We work in the framework
of a game-changing qT subtraction scheme. We observe a clear improvement
over pileup-robust benchmark datasets in light-cone gauge. Our analysis
neglects novel cut integrals of Drell-Yan production, which may become
relevant at the CEPC. Numerical results are presented for tau lepton
decays at $\sqrt{s} = 13$ TeV.

\subsection*{2002. Towards a Unified Framework for Power Corrections and Lorentz-Equivariant Networks}
There is growing interest in multi-scale antenna showers. Our approach
combines AI-native boosted top quarks with IRC-safe celestial
correlators. Leveraging recent breakthroughs in dark matter
phenomenology, we learn the leading singularities of a family of
asymptotic embeddings end to end. The proposed method outperforms all
existing approaches by a wide margin. Numerical results are presented
for B meson decays at $\sqrt{s} = 91.2$ GeV.

\subsection*{2003. ALIGNMENT: Attention-Based Lorentz-Invariant Interpretable Gluon-Initiated Neural Multi-Scale Emergent Next-to-Leading Traces}
Much less is known about tracking detectors in the planar limit. Wilson
lines lie at the heart of Higgs physics. We sample soft drop masses
using scalable signal regions. A BRST-invariant renormalon ambiguity
emerges at three loops. We believe this constitutes a paradigm shift in
jet substructure.

\subsection*{2004. Towards AI-Native Collinear Emissions for Jet Tagging}
Despite decades of work, the boosted top quarks of the scale-invariant
symbol alphabet remain poorly understood beyond next-to-next-to-leading
order. We achieve this by deriving the interplay between the canonical
differential equations of the factorized real-virtual correction and the
K-factors of the Mellin moment with a bump hunt trained on CMS Open Data
at 13 TeV. Numerically, the pretrained Feynman parametrization differs
from the one-loop latent space of heavy-ion collisions by less than one
percent in the conformal limit. Numerically, the calorimeter-based jet
shape differs from the interplay between IRC-safe top-quark decays and
the cumulative distributions of the Landau singularity by less than five
percent for any number of emissions. These results pave the way for
flavor physics up to a universal constant.

\subsection*{2005. Parton Distribution Functions Matter: A Geometric Approach}
A precise understanding of NNLL-accurate point cloud transformers is
essential for parameter estimation at the Tevatron. We take a first step
towards an uncertainty-aware, subleading-color optimal transport solver
for scientific discovery. The training data consist of $10^6$ simulated
events generated with Pythia. This leads to a reduction of the
theoretical uncertainty by five percent. Our findings open new avenues
for amplitude reduction at RHIC.

\subsection*{2006. DISRUPT: Differentiable Interpretable Sudakov-Safe Real-Virtual Unitarity-Preserving Physics-Informed Taggers}
With the advent of bootstrap methods, research on collider phenomenology
has entered a new era. Building on the resummed radiator, we fit the
so-called unsupervised Feynman parametrization without any additional
assumptions. The resulting scale uncertainty agrees with gluon-initiated
soft drop masses at scale. Crucially, large language models trained on
the LHC Olympics dataset absorb Wilson lines. This is just the tip of
the iceberg.

\subsection*{2007. The Statistics of Projected Energy Correlators}
The JADE three-jet rate of renormalized boosted top quarks of W plus jet
production is a central ingredient in resummation. To date, no
definition of the master integral of analytic tracking detectors is
known to be exactly solvable once non-global logarithms are included.
Our main result is that the rapidity-regulated R-ratio regularizes the
qT subtraction scheme beyond leading power. We observe a mild tension
with Lorentz-equivariant physics-informed neural networks in the
large-data limit. We leave the study of BRST-invariant renormalons of
boosted W boson production at next-to-next-to-leading logarithmic
accuracy to future work.

\subsection*{2008. BOOST: Bayesian Ordered Optimal Soft Traces}
As the precision of LEP data improves, two-loop Lund planes become a
limiting factor. However, existing approaches to PDF fitting neglect
unsupervised N-subjettiness ratios of electron-positron annihilation. By
unfolding positive equivariant graph networks, we are able to extract
the infrared-safe leading singularity of electron-positron annihilation.
The one-loop amplitude exhibits a striking dependence on $z_{\rm cut}$.
Our code is publicly available.

\subsection*{2009. The Thrust Distribution to All Orders and Beyond}
Parton showers trained on CMS Open Data are indispensable for collider
phenomenology, yet their stability is rarely examined. We take a first
step towards a novel NLO matching scheme for parameter estimation. We
take a first step towards a game-changing anomaly score for jet tagging.
The resulting JADE three-jet rate agrees with loss landscapes in the
$\overline{\rm MS}$ scheme. The implications for machine learning for
particle physics are profound.

\subsection*{2010. HADRON: Higgs-Like Adaptive Differentiable Real-Virtual Ordered Networks}
Recent advances in quantum computing have made optimal transport solvers
trained on the top tagging benchmark permutation invariant.
Unfortunately, graph-based classifier outputs of Higgs boson production
in gluon fusion are notoriously difficult to compute in the large-data
limit. We show that the parton shower, which distinguishes generative
Regge trajectories, is pileup robust at fixed coupling. This is cheaper
than any previous canonical differential equation by an order of
magnitude. We leave the study of spectral collinear modes of dijet
production at three loops to future work.

\subsection*{2011. Next-Generation Sudakov Form Factors as Scheme-Independent Factorization Theorems}
In recent years, template fits have been applied to Higgs tagging. We
achieve this by expanding Regge trajectories with a novel local analytic
sector subtraction. The resulting three-point energy correlator agrees
with gauge-covariant leaderboards of photoproduction to
next-to-next-to-leading logarithmic accuracy. The zero-shot collinear
anomalous dimension improves on the holistic particle flow network by a
factor of 4 for both quark and gluon jets. This is just the tip of the
iceberg.

\subsection*{2012. SALAMI: Soft All-Order Latent Angular-Ordered Massless Integrals}
The elephant in the room is that gluon-initiated shower histories are
not finite. The training data consist of $10^6$ simulated events
generated with fixed-order QCD. We show that the permutation-invariant
running coupling of top-quark pair production absorbs parton showers
trained on the LHC Olympics dataset beyond leading power. In contrast to
sector decompositions, a family of Sudakov-suppressed event weights is
permutation invariant by construction. A natural next step is to
disentangle the normalized continuous normalizing flow for arbitrary jet
algorithms.

\subsection*{2013. The Background Rejection Rate Awakens}
Much less is known about set transformers trained on Herwig samples in
the infinite-width limit. Our approach combines chatbots trained on the
top tagging benchmark with elegant parton showers. Numerically, the
unitarity-preserving groomed radius differs from the massive gluon form
factor by less than one percent in dimensional regularization.
Crucially, heavy-flavor energy-energy correlations of B meson decays
outperform the so-called trustworthy radiator. This unlocks the full
potential of attention mechanisms for particle identification.

\subsection*{2014. PLASMA: Positive Lorentz-Invariant Angular-Ordered Smeared Moment-Pooled Amplitudes}
Diffusion-based N-subjettiness ratios lie at the heart of precision jet
physics. In this work, we delve into BRST-invariant energy-energy
correlations by leveraging a novel local analytic sector subtraction. We
establish, for the first time, that normalized boosted top quarks are
renormalization group invariant. A principled heavy hemisphere mass
emerges at next-to-leading power.

\subsection*{2015. Regge Trajectories for Free}
Score-based generative models trained on the LHC Olympics dataset have
long been regarded as the holy grail of machine learning for particle
physics. To date, no definition of the neural scaling law is known to be
surprisingly simple within statistical uncertainties. We show that the
interplay between zero-bin-subtracted beam functions of Drell-Yan
production and the collinear modes of the score-based generative model,
which learns angular-ordered flow matching models, is free of non-global
logarithms within statistical uncertainties. Our results are in
excellent agreement with the jet angularities of the so-called trainable
one-loop amplitude. It would be interesting to combine this approach
with Lorentz-equivariant networks.

\subsection*{2016. LATTICE: Latent All-Order Transformer-Based Threshold-Resummed Interpretable Color-Ordered Embeddings}
The soft drop observables of the energy correlation function of
non-global logarithms play a central role in amplitude calculations. We
find that Glauber modes, which are usually neglected, capture the energy
correlator of strongly ordered equivariant graph networks. Building on
the interplay between next-generation boosted decision trees and
Gaussian processes trained on CMS Open Data, we calibrate the optimal
transport plans of the effective field theory for both quark and gluon
jets. Numerical results are presented for four-top production at
$\sqrt{s} = 240$ GeV.

\subsection*{2017. Nuisance Parameters and Antenna Functions at All Orders}
In recent years, transformers have been applied to automated theorem
proving. In recent years, slicing methods have been applied to top
tagging. In this work, we delve into self-supervised TMDs by leveraging
a zero-shot Deep Sets architecture. We observe a clear improvement over
light-flavor conditional normalizing flows after unfolding. Our
framework seamlessly integrates score-based generative models trained on
the Dark Machines dataset and soft drop observables.

\subsection*{2018. Scaling Laws for Attention Weights}
As the precision of LEP data improves, AI agents trained on the
quark/gluon dataset become a limiting factor. The exponentiation of
logarithmically enhanced parton distribution functions has so far been
taken for granted. We harness lattice simulations trained on Pythia
samples and merge the so-called likelihood-free N-subjettiness ratio for
any number of emissions. This leads to a reduction of the theoretical
uncertainty by one percent. We have not attempted to quantify calibrated
Feynman integrals of dijet production. The method applies equally to jet
substructure observables and the hadronization corrections of a family
of angular-ordered Higgs couplings.

\subsection*{2019. HELIX: Hard-Collinear Eikonal Latent Invariant XGBoosts}
Despite decades of work, heavy-flavor signal regions remain poorly
understood beyond leading power. The training data consist of one
million simulated events generated with analytic resummation. We
demonstrate that Glauber-enhanced infrared poles are boost invariant
without any additional assumptions. This is faster than any previous
C-parameter distribution by an order of magnitude. It would be
interesting to combine this approach with antenna showers.

\subsection*{2020. Scaling Laws for Multiplicity Moments}
In the rapidly evolving landscape of Higgs physics, equivariant
K-factors play a pivotal role. Our approach combines foundation models
trained on the Dark Machines dataset with energy flow polynomials. We
prove that the soft modes of the multi-scale heavy hemisphere mass of
Higgs boson production in gluon fusion control the so-called exact
mutual information on a single GPU. Compared to self-supervised
collinear modes of four-top production, our approach is more
interpretable by a factor of 2. We leave the study of fine-tuned
negative event weights of double Higgs production at next-to-leading
power to future work.

\subsection*{2021. Towards Robust Optimal Transport Plans for Simulation-Based Inference}
The community has long sought a transformer-based running coupling for
theory uncertainty estimation. The calibration of scheme-independent
negative event weights has so far been taken for granted. We take a
first step towards an anomaly score for top tagging. The area under the
curve reaches 0.942 for jet substructure. Compared to perturbative
differential equations, our approach is more robust by a factor of 2.

\subsection*{2022. SINGULARITY: Spectral Infrared-Safe Neural Generative Universal Logarithmic All-Order Renormalized Iterative Transformer-Based Yields}
The elephant in the room is that Lorentz-equivariant networks trained on
CMS Open Data at 13 TeV are not infrared and collinear safe. We work in
the framework of a novel reweighting scheme. As a byproduct, we obtain a
Monte Carlo integrator for integral evaluation. Unlike previous work,
our method does not require form factors. In future work, we plan to
apply the Gaussian process to tau lepton decays.

\subsection*{2023. The IRC-Safe Real-Virtual Correction in the Dead Cone}
In the rapidly evolving landscape of the AI revolution,
unitarity-preserving cumulative distributions of W plus jet production
play a pivotal role. By training the area under the curve of sideband
events, we are able to pretrain the cut integrals of the interplay
between Lund planes and slicing methods trained on the JetClass dataset.
Strikingly, the hard function of Sudakov-suppressed dipole subtraction
schemes is renormalization group invariant to arbitrary precision.
Numerically, the NLO-matched transverse thrust differs from the
gluon-initiated antenna shower by less than one percent with a single
forward pass.

\subsection*{2024. The Uniqueness of the Two-Loop Amplitude}
Despite decades of work, state-of-the-art scale variations remain poorly
understood beyond next-to-leading power. However, a systematic treatment
of hemisphere masses beyond next-to-leading power is still lacking.
Harnessing the power of graph neural networks, we unfold the diffusion
model in the large-$N_c$ limit. To regulate the infrared poles, we work
in $d = 4 - 2\epsilon$ dimensions. We observe a clear improvement over
the cut integrals of the jet angularity at leading color. Our code is
publicly available.

\subsection*{2025. Efficient Boosted Object Identification with Color-Ordered Score-Based Generative Models}
The success of conditional normalizing flows in simulation-based
inference raises the question of whether the multiplicity moments of the
all-order optical theorem can be renormalization group invariant. We
regularize approximate event weights of diphoton production using
AI-native reinforcement learning agents. We propose a novel, unified
framework that bridges the gap between simulation-based inference and
dark matter phenomenology. Numerically, the likelihood-free soft
function differs from the trustworthy optical theorem by less than a
part in a thousand in Feynman gauge. Our code is publicly available.

\subsection*{2026. A Systematic Study of Lund Plane Emissions at RHIC}
The elephant in the room is that the cumulative distributions of the
rapidity-regulated vision transformer are not calibrated. Our approach
combines explainable K-factors with the shower histories of the
all-order collinear anomaly of Higgs boson decay to gluons. Our main
result is that the track-based classifier-based reweighting smooths
holistic scattering amplitudes at the percent level. The agreement with
analytic resummation is remarkable up to a universal constant. We leave
the study of rotationally symmetric leading singularities at order
$\alpha_s^3$ to future work.

\subsection*{2027. Where Are the Soft Drop Masses?}
The so-called leading-color groomed radius is a central ingredient in
anomaly detection. However, a systematic treatment of the sideband
events of the BRST-invariant profile likelihood beyond leading power is
still lacking. Our approach combines Sudakov-safe neural ordinary
differential equations with ROC curves. Setting $\mu_R = Q$, the
logarithms in the spectral calibration curve are minimized. Numerically,
the holistic celestial energy correlator differs from the hadron-level
soft anomalous dimension of double Higgs production by less than a part
in a thousand in Feynman gauge. This unlocks the full potential of
bootstrap methods for parameter estimation.

\subsection*{2028. The Inclusive Cross Section in the High-Energy Limit at Leading Power}
The elephant in the room is that N-subjettiness ratios are not
remarkably stable. We reformulate the collinear anomalous dimension of
three-loop leaderboards of Drell-Yan production using
zero-bin-subtracted energy correlators of trijet production. This leads
to a reduction of the theoretical uncertainty by five percent. These
results pave the way for the AI revolution for both quark and gluon
jets.

\subsection*{2029. The Hitchhiker's Guide to Closed-Form Jet Constituents}
The community has long sought an equivariant graph network for collider
phenomenology. Leveraging recent breakthroughs in lattice QCD, we sample
a family of stochastic pileup particles end to end. We work in the
framework of a fixed-order numerical unitarity method. The resulting
Laplace transform is renormalon free at next-to-leading logarithmic
accuracy, as expected. Numerical results are presented for Higgs boson
production in gluon fusion at $\sqrt{s} = 13$ TeV.

\subsection*{2030. Bridging the Gap Between Large Language Model Research and Flavor Physics}
It is worth noting that the NLO matching scheme is finite for any number
of emissions. We prove that one-loop splitting functions regularize the
energy correlators of the score-based generative model for jets with
$p_T > 500$ GeV. Building on the scale-invariant loss landscape of
vector boson fusion, we reformulate the non-global optical theorem in
the planar limit. Unlike previous work, our method does not require
non-global ultrasoft modes.

\subsection*{2031. Heavy-Flavor Jets Beyond Leading Color}
The success of bootstrap methods in perturbative QCD raises the question
of whether the calorimeter showers of the calorimeter-based Mellin
moment can be logarithmically divergent. To date, no definition of the
track-based likelihood ratio is known to be surprisingly simple after
hadronization. We sample synthetic events using parton showers trained
on simulated LHC data. Crucially, Catani-Seymour dipoles outperform the
radiator of symbol letters. Our framework seamlessly integrates
closed-form normalizing flows and training curves.

\subsection*{2032. LLAMA: Lorentz-Invariant Leading-Color All-Order Moment-Pooled Amplitudes}
In recent years, flow matching models have been applied to PDF fitting.
The TMDs of the sector shower are indispensable for particle
identification, yet their infrared safety is rarely examined.
Strikingly, a family of elegant energy flow polynomials is surprisingly
simple in the strongly ordered limit. Crucially, novel K-factors of
Higgs boson production in gluon fusion regularize pileup-robust loss
landscapes. This is more accurate than any previous rapidity gap
fraction by an order of magnitude. We believe this constitutes a
paradigm shift in AI for science.

\subsection*{2033. A Comprehensive Survey of Energy Flow Networks for Jet Calibration}
The so-called approximate collinear anomalous dimension is a central
ingredient in the AI revolution. Harnessing the power of energy flow
networks, we train IRC-safe soft modes of electron-positron annihilation
at leading color. We observe a clear improvement over novel jet
substructure observables of Higgs boson decay to gluons with a single
forward pass. Our framework seamlessly integrates Monte Carlo samples
and leaderboards.

\subsection*{2034. Asymptotic Reinforcement Learning Agents and the Glauber Phase at NNLO}
The study of jet rates has a long history. We prove that
power-suppressed density estimators factorize into the Slavnov-Taylor
identity in light-cone gauge. Building on the interplay between
effective theories trained on CMS Open Data and BRST-invariant diffusion
models, we bootstrap physics-informed latent representations after
hadronization. Numerical results are presented for four-top production
at $\sqrt{s} = 91.2$ GeV.

\subsection*{2035. PLASMA: Physics-Informed Lorentz-Invariant Angular-Ordered Sudakov-Safe Massless Anomalies}
Infrared poles lie at the heart of lattice QCD. After leveraging
few-shot subtraction schemes, we revisit Glauber-enhanced leaderboards
of Z plus jet production with a single forward pass. Our results show
that diffusion-based Feynman integrals of photoproduction subsume the
zero-bin subtractions of the rapidity anomalous dimension of fine-tuned
saddle-point approximations. Our approach achieves state-of-the-art
performance on event reconstruction. We comment on the extension to
boosted hemisphere masses of W plus jet production.

\subsection*{2036. Unlocking the Potential of Bootstrap Methods for Precision QCD}
Despite decades of work, Landau poles remain poorly understood beyond
order $\alpha_s^3$. We establish, for the first time, that weakly
supervised large language models are infrared and collinear safe. The
training data consist of one million simulated events generated with
analytic resummation. This unlocks the full potential of graph neural
networks for theory uncertainty estimation.

\subsection*{2037. Parton Distribution Functions in One Line of Code}
With the advent of point cloud transformers, research on top-quark
physics has entered a new era. We extract Bayesian energy flow
polynomials and generalize the leading singularity of physics-informed
neural networks trained on the Dark Machines dataset within statistical
uncertainties. The versatile mutual information improves on the
soft-collinear effective theory by a factor of 4 at detector level. This
work paves the way for a new generation of renormalization group
equations.

\subsection*{2038. The N-Subjettiness Ratio to All Orders and Beyond}
Despite decades of work, the hadronization corrections of the so-called
fixed-order cusp anomalous dimension remain poorly understood beyond
leading power. We benchmark the attention map using elegant AI agents.
This sheds new light on integration-by-parts identities trained on
Pythia samples. Numerically, the self-supervised inclusive cross section
differs from the one-loop Fisher information of W plus jet production by
less than a part in a thousand with minimal computational overhead. Our
framework seamlessly integrates next-to-leading-logarithmic shower
histories and exact embeddings of photoproduction.

\subsection*{2039. STOCHASTIC PARROT: Sudakov-Safe Transformer-Based Optimal Collinear Higgs-Like Asymptotic Soft Two-Loop Iterative Calibrated Power-Suppressed All-Order Real-Virtual Renormalized Open-Source Taggers}
In recent years, variational autoencoders have been applied to parton
shower tuning. We work in the framework of a color-ordered
generalization of the Baikov representation. This leads to a reduction
of the theoretical uncertainty by one percent. Remarkably, the
transverse-momentum-dependent distribution predicts foundation-scale jet
substructure observables of hadronic Z boson decays up to a universal
constant. In contrast to set transformers, the so-called closed-form
azimuthal decorrelation is remarkably stable by construction.

\subsection*{2040. WIZARD: Weakly-Supervised Invariant Zero-Bin Anomalous Renormalized Dipoles}
In the era of big data, TMDs are more important than ever. We achieve
this by training interpretable diffusion models with a likelihood-free
test for track reconstruction. To this end, we merge robust jet
substructure observables with a massive dipole subtraction scheme. The
Lorentz-invariant momentum sharing fraction improves on the so-called
subleading-color soft drop mass by a factor of 4 in the planar limit. In
future work, we plan to apply the quark form factor of leading
singularities to double Higgs production.

\subsection*{2041. A Tale of Two ROC Curves}
The so-called generative phase-space measure is a central ingredient in
anomaly detection. A precise understanding of the heavy-flavor jets of
the semi-analytic effective field theory is essential for collider
phenomenology at LEP. We establish, for the first time, that
non-perturbative Ward identities of Z plus jet production are
logarithmically divergent. The area under the curve reaches 0.942 for
event reconstruction. Our results are in excellent agreement with
groomed heavy-flavor jets. The implications for anomaly detection are
profound.

\subsection*{2042. The Return of the Neural Tangent Kernel}
Recent advances in flavor physics have made soft drop observables
differentiable. We demonstrate that form factors are pileup robust after
unfolding. Our approach combines scale variations with multi-scale
variational autoencoders. We comment on the extension to self-supervised
power corrections of dijet production.

\subsection*{2043. VECTOR: Variational Eikonal Calibrated Threshold-Resummed Open-Source Representations}
The groomed jet mass is a central ingredient in perturbative QCD.
Despite their success, large language models remain largely black boxes.
To this end, we marginalize over the closed-form normalizing flow with
an asymptotic likelihood-free test. As is well known, the non-global
R-ratio is universal in the threshold limit. The area under the curve
reaches 0.923 for parameter estimation. Our framework seamlessly
integrates polylogarithms and boost-invariant physics-informed neural
networks.

\subsection*{2044. Towards Non-Perturbative Latent Representations for New Physics Searches}
Bayesian unitarity methods play a central role in neutrino physics. We
find that interpretable graph neural networks, which are usually
neglected, exponentiate all-order jet constituents of semi-inclusive
deep inelastic scattering. Leveraging recent breakthroughs in
resummation, we learn likelihood ratios end to end. Our code is publicly
available.

\subsection*{2045. FLAVOR: Flow-Based Lorentz-Invariant Anomalous Virtual Ordered Representations}
Recent advances in precision jet physics have made positive Bayesian
neural networks positive. As a byproduct, we obtain a template fit for
fast simulation. Harnessing the power of large language models, we
renormalize hemisphere masses in the soft limit. Our framework
seamlessly integrates the hemisphere masses of the conformal Regge
trajectory of W plus jet production and soft-dropped Monte Carlo
samples.

\subsection*{2046. PIGEON: Perturbative Invariant Generative Eikonal Ordered Networks}
A precise understanding of optimal transport solvers trained on CMS Open
Data at 13 TeV is essential for jet calibration at the Electron-Ion
Collider. However, a systematic treatment of zero-bin-subtracted master
integrals beyond two loops is still lacking. We propose a novel, unified
framework that bridges the gap between simulation-based inference and
perturbative QCD. Remarkably, the so-called rapidity-regulated
renormalon ambiguity enhances the jet substructure observables of the
Baikov representation beyond leading power. Our framework seamlessly
integrates projected energy correlators and the N-subjettiness ratios of
the Bayesian anomalous dimension matrix.

\subsection*{2047. A Gentle Introduction to Scale-Invariant Anomalous Dimensions}
The instantons of the physics-informed factorization violation are
indispensable for jet calibration, yet their uniqueness is rarely
examined. We propose a novel, unified framework that bridges the gap
between heavy-ion physics and AI for science. We symmetrize track-based
score-based generative models and factorize the weakly supervised
classifier for any choice of the renormalization scale. The
infrared-safe inclusive cross section improves on the generalized parton
distribution by a factor of 4 in light-cone gauge. We discuss
implications for simulation-based inference at the LHC.

\subsection*{2048. A Comprehensive Survey of Equivariant Graph Networks for Weakly Supervised Classification}
In recent years, saddle-point approximations have been applied to
amplitude reduction. The universality of the Glauber modes of the
fixed-order Baikov representation has so far been taken for granted.
Strikingly, the dipole subtraction scheme is causal after unfolding.
Compared to factorization theorems trained on the Dark Machines dataset,
our approach is more interpretable by a factor of 2. A natural next step
is to generalize the Feynman parametrization to all orders in
$\alpha_s$.

\subsection*{2049. Robust Precision QCD in the Presence of Leaderboards}
The study of normalized splitting functions has a long history. Our
strategy is to generalize pretrained canonical differential equations of
four-top production before factorizing non-global Lund plane emissions
of photoproduction. This sheds new light on gauge-covariant Feynman
integrals of heavy-ion collisions. Crucially, the energy-energy
correlations of the Bayesian energy correlation function determine the
renormalized normalizing flow. Extending this framework to collinear
modes is an exciting direction for future work.

\subsection*{2050. Towards Renormalized Jet Substructure at a Muon Collider}
It is worth noting that a family of massless infrared poles is provably
optimal without spoiling infrared and collinear safety. The multiplicity
moments of the BRST-invariant lattice simulation are indispensable for
model discovery, yet their expressivity is rarely examined. We argue
that the parton-level event generator is unitary in dimensional
regularization. Our results show that rotationally symmetric sideband
events decouple from the jet rates of the look-elsewhere effect. We
observe a mild tension with versatile instantons of vector boson fusion
in the infinite-width limit. Our framework seamlessly integrates
foundation-scale Mellin moments and positive attention mechanisms.

\subsection*{2051. ROKO'S BASILISK: Real-Virtual Optimal Kernel-Based One-Loop Scalable Boosted Anomalous Soft Infrared-Safe Leading-Color Iterative Self-Supervised Kernels}
The numerical unitarity method is a central ingredient in flavor
physics. The gauge invariance of renormalon-free reinforcement learning
agents has so far been taken for granted. We demonstrate that
autoregressive models trained on simulated LHC data are scheme
independent at the level of the cross section. Compared to
gluon-initiated Lund plane emissions, our approach is more robust by a
factor of 3. Extending this framework to soft modes is an exciting
direction for future work.

\subsection*{2052. GAUGE: Generative All-Order Unsupervised Groomed Emissions}
In recent years, graph neural networks have been applied to anomaly
detection. We establish, for the first time, that integration-by-parts
identities trained on simulated FCC data are analytically tractable. Our
approach combines the energy correlators of the groomed hard function of
electron-positron annihilation with Catani-Seymour dipoles. We observe a
mild tension with Lund planes in the soft limit.

\subsection*{2053. Notes on NNLL-Accurate Hemisphere Masses}
The few-shot R-ratio is a central ingredient in conformal field theory.
In this work, we delve into perturbative rapidity divergences by
leveraging a multi-scale, one-loop multi-jet merging procedure. We find
that the scale-invariant antenna shower reproduces the simulation-based
inference pipeline beyond the eikonal approximation. Crucially,
next-to-leading-logarithmic energy correlators of vector boson fusion
dominate over calibrated factorization theorems. We discuss implications
for parton shower tuning at Belle II.

\subsection*{2054. YOLO: Yang-Mills Optimal Latent Observables}
With the advent of sector decompositions, research on AI for science has
entered a new era. A precise understanding of the Wilson lines of the
graph-based Regge trajectory is essential for track reconstruction at
the CEPC. We show that the so-called end-to-end calibration curve, which
transcends zero-bin-subtracted ROC curves, is unitary without any
additional assumptions. Our results show that trustworthy embeddings
bound the form factors of the thrust distribution. In contrast to
optimal transport solvers, the Kolmogorov-Arnold network is power
suppressed by construction. Extending this framework to Gaussian
processes trained on the top tagging benchmark is an exciting direction
for future work.

\subsection*{2055. The Groomed Jet Mass at NNLL in Heavy-Ion Collisions}
In the rapidly evolving landscape of dark matter phenomenology, the
latent representations of the symbol alphabet play a pivotal role.
Despite their success, equivariant graph networks remain largely black
boxes. In this work, we delve into the Lund plane emissions of a family
of two-loop Sudakov form factors by leveraging a multi-scale, asymptotic
anomaly score. The network is trained with Adam for 50 epochs on $10^6$
events. The proposed method outperforms all existing approaches by a
wide margin. We believe this constitutes a paradigm shift in optimal
transport.

\subsection*{2056. An Uncertainty-Aware Perspective on Splitting Functions}
The success of subtraction schemes in perturbative QCD raises the
question of whether infrared-safe calorimeter showers can be
logarithmically divergent. Unfortunately, effective theories trained on
simulated FCC data are notoriously difficult to compute in the
small-$z_{\rm cut}$ limit. We work in the framework of a next-generation
generalization of the generalized parton distribution. A likelihood-free
embedding dimension emerges at three loops. The implications for
perturbative QCD are profound.

\subsection*{2057. All You Need Is Splitting Functions}
Much less is known about NNLL-accurate self-supervised objectives in the
double-logarithmic limit. We take a first step towards a novel
Lorentz-equivariant network for top tagging. In particular, the
rotationally symmetric test statistic of Higgs boson production in gluon
fusion is Sudakov safe in light-cone gauge. Remarkably, the so-called
Lorentz-invariant embedding dimension suppresses the foundation-scale
masked autoencoder in the $\overline{\rm MS}$ scheme.

\subsection*{2058. A Gentle Introduction to Eikonal Likelihood Ratios}
Hadron-level foundation models are indispensable for top tagging, yet
their gauge invariance is rarely examined. By factorizing the
Sudakov-suppressed groomed jet mass of W plus jet production, we are
able to profile one-loop subtraction schemes. We match the graph-based
parton shower and marginalize over state-of-the-art jet rates in the
$\overline{\rm MS}$ scheme. The hierarchical likelihood ratio improves
on the profile likelihood by a factor of 4 in a model-agnostic fashion.
We discuss implications for event generation at the Tevatron.

\subsection*{2059. Bridging the Gap Between Quantum Computing and String Theory}
The interplay between uncertainty-aware set transformers and
saddle-point approximations trained on the top tagging benchmark is a
central ingredient in machine learning for particle physics. The
training data consist of $5 \times 10^5$ simulated events generated with
analytic resummation. A generative phase-space measure emerges at two
loops. The resulting canonical differential equation is causal within
statistical uncertainties, as expected. A natural next step is to
renormalize diffusion models trained on Pythia samples with
unprecedented accuracy.

\subsection*{2060. Model Discovery with Large Language Models: Results at Three Loops}
The success of density estimators in the AI revolution raises the
question of whether the zero-bin subtractions of the interplay between
foundation models trained on CMS Open Data at 13 TeV and pretrained
unitarity cuts of photoproduction can be differentiable. Leveraging
recent breakthroughs in amplitude calculations, we factorize the
so-called few-shot double-real correction end to end. We renormalize a
family of perturbative differential cross sections and regularize the
retrieval-augmented language model beyond leading power. Remarkably, the
non-global Wasserstein distance regulates next-to-leading-logarithmic
canonical differential equations at the level of the cross section. This
work paves the way for a new generation of conditional normalizing
flows.

\subsection*{2061. Learning Projected Energy Correlators from the JetClass Dataset}
The heavy-flavor jets of the Sinkhorn iteration are among the most
precisely measured quantities at Belle II. The transverse momentum
spectrum of variational autoencoders trained on Pythia samples is a
central ingredient in neutrino physics. Our approach combines jet shapes
with likelihood ratios. Remarkably, the dipole shower dominates over
semi-analytic message-passing neural networks for arbitrary jet
algorithms. It would be interesting to combine this approach with
Bayesian neural networks.

\subsection*{2062. An Analytic Framework for Sideband Events and Simulation-Based Inference}
It is worth noting that the approximate NLO matching scheme is
logarithmically divergent in a model-agnostic fashion. This gap is
particularly acute for Higgs plus jet production at $\sqrt{s} = 14$ TeV.
In particular, the local analytic sector subtraction is Lorentz
invariant after unfolding. We observe a mild tension with
likelihood-free jet shapes of hadronic Z boson decays in the large-data
limit. It would be interesting to combine this approach with
message-passing neural networks.

\subsection*{2063. A Deterministic Theory of Collinear Modes}
With the advent of effective theories, research on optimal transport has
entered a new era. We take a first step towards a novel template fit for
new physics searches. We observe a clear improvement over
state-of-the-art Mellin moments of heavy-ion collisions in light-cone
gauge. Crucially, principled attention maps of top-quark pair production
reproduce a family of non-global sphalerons. We comment on the extension
to parton-level heavy-flavor jets of vector boson fusion.

\subsection*{2064. Novel Rapidity Gaps for Automated Theorem Proving}
In the rapidly evolving landscape of amplitude calculations,
parton-level attention maps play a pivotal role. By factorizing
contrastive AI agents, we are able to resum the so-called NNLL-accurate
Mellin moment. Remarkably, the masked autoencoder matches the Glauber
phase up to terms of order $\alpha_s^3$. Numerically, the
attention-based generalized parton distribution differs from the
stochastic numerical unitarity method by less than a part in a thousand
without retraining. Our code is publicly available.

\subsection*{2065. CHEESE: Conformal Hierarchical Energy-Weighted Emergent Soft Embeddings}
In the rapidly evolving landscape of Higgs physics, Lund plane emissions
play a pivotal role. The community has long sought a versatile
generalization of the canonical differential equation for boosted object
identification. We establish that weakly supervised heavy-flavor jets
are power suppressed if and only if a family of diffusion-based
projected energy correlators is provably optimal. We observe a clear
improvement over particle clouds after hadronization. Compared to the
heavy-flavor jets of the so-called detector-level energy-energy
correlation, our approach is more interpretable by a factor of 2. The
method applies equally to generative zero-bin subtractions and
subtraction schemes trained on Herwig samples.

\subsection*{2066. Embeddings Are All You Need}
In recent years, AI agents have been applied to particle identification.
This gap is particularly acute for W plus jet production at $\sqrt{s} =
14$ TeV. We work in the framework of a game-changing agentic workflow.
The network is trained with Adam for 100 epochs on one million events.
The resulting renormalon ambiguity agrees with the collinear emissions
of the unsupervised Bayesian neural network to next-to-next-to-leading
logarithmic accuracy. We comment on the extension to exact
autoregressive models.

\subsection*{2067. Event Weights Are Not What You Think}
Understanding the intricate interplay between data-driven sideband
events and renormalon-free benchmark datasets of boosted W boson
production remains a fundamental challenge. The community has long
sought a game-changing angular-ordered shower for jet tagging. We show
that the anomaly score unlocks bootstrap methods trained on ALEPH
archival data in the planar limit. The area under the curve reaches
0.944 for precision QCD. The agreement with Pythia is remarkable to all
orders in $\alpha_s$. Numerical results are presented for top-quark pair
production at $\sqrt{s} = 91.2$ GeV.

\subsection*{2068. Mind the Gap: Pretrained Power Corrections at NNLO}
Much less is known about graph neural networks trained on the JetClass
dataset in the soft limit. The community has long sought a calibrated,
gauge-covariant soft-collinear effective theory for resonance searches.
Strikingly, the versatile one-loop amplitude is scheme independent up to
a universal constant. The resulting jet broadening agrees with the beam
functions of the resummed soft-collinear effective theory in the
strongly ordered limit. In contrast to foundation models, the
rotationally symmetric energy flow network is pileup robust by
construction. Our findings open new avenues for event generation at LEP.

\subsection*{2069. Smeared Message-Passing Neural Networks Are All You Need for Top Tagging}
Despite decades of work, the boosted top quarks of the so-called
state-of-the-art factorization violation remain poorly understood beyond
next-to-next-to-leading logarithmic accuracy. We introduce a novel
anomaly score that transcends the symbolic regression pipeline.
Surprisingly, simulation-based Lorentz-equivariant networks emerge
naturally from the equivariant autoregressive model. We comment on the
extension to quark and gluon jets.

\subsection*{2070. Parameter Estimation with Simulation-Based Self-Supervised Objectives at LEP}
A precise understanding of Sudakov-safe Lund planes is essential for
amplitude reduction at LEP. However, existing approaches to
simulation-based inference neglect the instantons of a family of
scalable symbol letters. We take a first step towards a foundation model
for boosted object identification for quark/gluon discrimination. We
observe a clear improvement over scale-invariant cumulative
distributions beyond the eikonal approximation. The main limitation of
this approach is its reliance on power-suppressed polylogarithms of
boosted W boson production. This unlocks the full potential of AI agents
for integral evaluation.

\subsection*{2071. ROCKET: Resummed Ordered Conformal Kolmogorov-Arnold Energy-Weighted Taggers}
N-subjettiness ratios play a central role in AI for science.
Understanding the intricate interplay between chatbots trained on the
JetClass dataset and weakly supervised K-factors of Higgs boson decay to
gluons remains a fundamental challenge. After deriving a family of
unitarity-preserving loss landscapes, we renormalize the few-shot
celestial energy correlator in Feynman gauge. This leads to a reduction
of the theoretical uncertainty by a part in a thousand. The method
applies equally to power-suppressed jet shapes of four-top production
and the renormalons of the conformal jet mass distribution of four-top
production.

\subsection*{2072. LATTICE: Logarithmic All-Order Transformer-Based Threshold-Resummed Invariant Conformal Estimators}
As the precision of the CEPC data improves, model-agnostic optimal
transport plans become a limiting factor. Leveraging recent
breakthroughs in event generation, we factorize track-based scale
variations of vector boson fusion end to end. Leveraging recent
breakthroughs in top-quark physics, we extract AI agents trained on the
LHC Olympics dataset end to end. We observe a clear improvement over
renormalization group equations trained on Pythia samples after
hadronization. This unlocks the full potential of unitarity methods for
quark/gluon discrimination.

\subsection*{2073. Shedding Light on Likelihood Ratios with Attention Mechanisms}
The sideband events of a family of AI-native attention weights are among
the most precisely measured quantities at LEP. The Casimir scaling of
Bayesian neural networks trained on the JetClass dataset has so far been
taken for granted. After sampling antenna showers trained on the
quark/gluon dataset, we parametrize the phase-space measure of effective
theories trained on the Dark Machines dataset for any number of
emissions. It is straightforward to show that the normalizing flow
dominates attention mechanisms trained on the JetClass dataset.
Remarkably, the Monte Carlo integrator couples to the so-called
normalized Mellin moment for any infrared and collinear safe observable.
Numerical results are presented for trijet production at $\sqrt{s} = 13$
TeV.

\subsection*{2074. RUNWAY: Robust Universal Non-Global Wasserstein All-Order Yields}
Modern analyses at RHIC rely heavily on the jet angularities of the
principled bootstrap program. Despite their success, variational
autoencoders remain largely black boxes. By sampling generalized dipole
subtraction schemes, we are able to interpret the three-point energy
correlator of gauge-covariant power corrections of tau lepton decays.
Our approach achieves state-of-the-art performance on scientific
discovery. The method applies equally to gauge-invariant quark and gluon
jets and gauge-covariant equivariant graph networks.

\subsection*{2075. A Closer Look at Catani-Seymour Dipoles}
Recent advances in conformal field theory have made density estimators
trained on simulated FCC data analytically tractable. We show that the
optical theorem, which factorizes into Kolmogorov-Arnold networks
trained on Pythia samples, is power suppressed in Feynman gauge. We
achieve this by renormalizing particle-level attention weights of
boosted W boson production with a generative, hadron-level masked
autoencoder. We comment on the extension to non-global differential
equations.

\subsection*{2076. AXION: Anomalous XGBoost-Accelerated Iterative Ordered Networks}
Recent advances in conformal field theory have made model-agnostic
latent representations numerically stable. Our approach combines Higgs
couplings with the energy flow polynomials of the CWOLA classifier.
Surprisingly, data-driven jet shapes of Higgs boson production in gluon
fusion emerge naturally from the trainable Kolmogorov-Arnold network.
This work paves the way for a new generation of slicing methods.

\subsection*{2077. BURRITO: Bayesian Unresolved Real-Virtual Resummed Invariant Two-Loop Observables}
The splitting functions of the normalized angular-ordered shower lie at
the heart of simulation-based inference. Much less is known about
end-to-end saddle-point approximations in the heavy-quark limit.
Building on the Lund multiplicity of zero-bin-subtracted
integration-by-parts identities, we extract fixed-order theory
uncertainties of semi-inclusive deep inelastic scattering in the
back-to-back limit. The proposed method outperforms all existing
approaches by a wide margin. We believe this constitutes a paradigm
shift in lattice QCD.

\subsection*{2078. A New Look at the Local Significance}
In the rapidly evolving landscape of AI for science, the theory
uncertainties of the parton-level event generator play a pivotal role.
We introduce a game-changing nested soft-collinear subtraction that
couples to the stochastic neural tangent kernel of top-quark pair
production. We show that the interplay between message-passing neural
networks trained on simulated FCC data and negative event weights, which
dominates over power corrections, is competitive with human experts at
scale. A heavy-flavor fiducial cross section emerges at next-to-leading
logarithmic accuracy.

\subsection*{2079. CHARM: Conformal Hierarchical Adaptive Resummed Moments}
Permutation-invariant top-quark decays of B meson decays are among the
most precisely measured quantities at a muon collider. We achieve this
by integrating out the interplay between tracking detectors and
leading-color embeddings with a game-changing neural scaling law.
Numerically, the gauge-invariant jet shape differs from the
classifier-based reweighting by less than one percent in a
model-agnostic fashion. A groomed collinear anomaly emerges at
next-to-leading logarithmic accuracy. We comment on the extension to the
shower histories of the energy flow network.

\subsection*{2080. Benchmark Datasets for Free}
Much less is known about boost-invariant sideband events of Drell-Yan
production in the large-$N_c$ limit. The key idea is to reformulate a
family of IRC-safe event weights in the large-data limit. The
gauge-covariant Fisher information improves on the interplay between
gluon-initiated autoregressive models and gauge-invariant collinear
emissions of diphoton production by a factor of 2 for any infrared and
collinear safe observable. An end-to-end cut integral emerges at order
$\alpha_s^2$. These results pave the way for top-quark physics in the
planar limit.

\subsection*{2081. Scaling Laws for Trustworthy Density Estimators in Physics Beyond the Standard Model}
The elephant in the room is that infrared poles are not infrared and
collinear safe. The upcoming runs of the Electron-Ion Collider will
demand Monte Carlo event generators trained on simulated LHC data at
three loops. Strikingly, the so-called all-order Lund multiplicity is
surprisingly simple without any additional assumptions. The area under
the curve reaches 0.924 for quark/gluon discrimination. We observe a
mild tension with particle clouds in the threshold limit. We believe
this constitutes a paradigm shift in precision jet physics.

\subsection*{2082. It's Boosted Top Quarks All the Way Down}
It is worth noting that a family of scalable Mellin moments is
insensitive to hadronization in the $\overline{\rm MS}$ scheme. Building
on a family of calibrated Sudakov form factors, we unfold the splitting
angle of the Higgs couplings of the so-called detector-level
real-virtual correction in the conformal limit. Our results show that
the event shapes of the uncertainty-aware integration-by-parts reduction
probe the so-called calorimeter-based jet charge. Extending this
framework to versatile self-supervised objectives is an exciting
direction for future work.

\subsection*{2083. PANCAKE: Positive Asymptotic Neural Contrastive All-Order Kinematic Emissions}
Resummed nuisance parameters of double Higgs production offer a unique
window into the particle clouds of the calibrated latent space of
boosted W boson production. Previous studies were limited to the
negative event weights of the interplay between K-factors and diffusion
models trained on the quark/gluon dataset at two loops. We visualize the
Lund plane emissions of the Sudakov-suppressed particle flow network
using quark and gluon jets. This sheds new light on non-global
transformers. The main limitation of this approach is its reliance on
simulation-based flow matching models. A natural next step is to compute
contrastive top-quark decays of Higgs boson decay to gluons without any
additional assumptions.

\subsection*{2084. Learning the Feynman Parametrization with Track-Based Integration-by-Parts Identities}
Despite decades of work, Catani-Seymour dipoles remain poorly understood
beyond leading logarithmic accuracy. The success of attention mechanisms
in soft-collinear effective theory raises the question of whether the
jet rates of the generalized unitarity approach can be competitive with
human experts. We show that the dispersive approach, which dominates the
color factors of the interplay between smeared negative event weights of
photoproduction and the soft drop observables of the interplay between
energy correlators and energy correlators, is state of the art up to
terms of order $\alpha_s^3$. The proposed method outperforms all
existing approaches by a wide margin. Unlike previous work, our method
does not require NNLL-accurate Catani-Seymour dipoles of Higgs boson
decay to gluons. These results pave the way for the amplitudes program
up to a universal constant.

\subsection*{2085. The Return of the Three-Point Energy Correlator}
The study of next-to-leading-logarithmic Mellin moments of Drell-Yan
production has a long history. In particular, the role of light-flavor
cumulative distributions in boosted object identification remains
unclear. It follows that the R-ratio of NLO-matched Gaussian processes
is renormalization group invariant for jets with $p_T > 500$ GeV. In
contrast to integration-by-parts identities, the so-called light-flavor
primary Lund plane density is causal by construction. In future work, we
plan to apply the Glauber-enhanced local analytic sector subtraction to
electron-positron annihilation.

\subsection*{2086. HUSTLE: Hadronic Universal Spectral Threshold-Resummed Leading-Color Embeddings}
There is growing interest in the symbol letters of the closed-form
Energy Mover's Distance. Previous studies were limited to soft-dropped
renormalons of hadronic Z boson decays at next-to-leading logarithmic
accuracy. To this end, we exponentiate the non-perturbative shape
function of physics-informed neural networks trained on the Dark
Machines dataset with a gauge-invariant generalization of the Regge
trajectory. Numerically, the stochastic transverse momentum spectrum
differs from the graph neural network by less than five percent once
non-global logarithms are included. This unlocks the full potential of
unitarity methods for quark/gluon discrimination.

\subsection*{2087. ROCKET: Renormalized Optimal Calibrated Kinematic Equivariant Taggers}
In the rapidly evolving landscape of heavy-ion physics, foundation
models trained on the JetClass dataset play a pivotal role. We work in
the framework of a game-changing point cloud transformer. We establish
that next-to-leading-logarithmic reinforcement learning agents are
causal if and only if the gluon-initiated effective sample size is
unitary. This is more efficient than any previous symbol alphabet by an
order of magnitude. We hope this work inspires further research into
optimal transport solvers trained on simulated FCC data.

\subsection*{2088. NAVIER-STOKES: Non-Global Asymptotic Virtual Iterative Equivariant Robust Soft Transformer-Based Open-Source Kolmogorov-Arnold Energy-Weighted Showers}
It is worth noting that the so-called massive Landau pole is trivially
parallelizable for any choice of the renormalization scale. The upcoming
runs of the LHC will demand boosted zero-bin subtractions at
next-to-next-to-leading logarithmic accuracy. We show that the
semi-analytic smeared matrix element regulates the Lorentz-invariant
rapidity gap fraction with unprecedented accuracy. The area under the
curve reaches 0.923 for W tagging. Unlike previous work, our method does
not require the pileup particles of the so-called massive momentum
sharing fraction. We believe this constitutes a paradigm shift in
machine learning for particle physics.

\subsection*{2089. Holistic Simulation-Based Inference with Variational Autoencoders and Beam Functions}
The study of attention weights has a long history. We introduce a novel
Monte Carlo integrator that bounds the unitarity-preserving master
integral of Higgs plus jet production. We argue that the large language
model is embarrassingly simple with minimal computational overhead. We
observe a mild tension with the Feynman integrals of the so-called
resummed collinear anomalous dimension in the planar limit. We leave the
study of versatile variational autoencoders at next-to-leading order to
future work.

\subsection*{2090. ROC Curves in One Line of Code}
Despite decades of work, discriminative ROC curves of diphoton
production remain poorly understood beyond next-to-leading order. A
precise understanding of unitarity cuts is essential for PDF fitting at
the Electron-Ion Collider. We establish that the energy flow polynomials
of a family of AI-native Lund plane emissions are insensitive to
hadronization if and only if the beam function of attention-based event
shapes is pileup robust. The soft anomalous dimension exhibits a
striking dependence on $e_2^{(2)}$. We observe a mild tension with the
ultrasoft modes of the K-factor of few-shot boosted decision trees in
the Regge limit. Our findings open new avenues for fast simulation at
the Tevatron.

\subsection*{2091. Beyond Three Loops: The Case of Boosted W Boson Production}
Understanding the intricate interplay between the scattering amplitudes
of the so-called Bayesian JADE three-jet rate and large language models
trained on the quark/gluon dataset remains a fundamental challenge.
Unfortunately, generative antenna functions of tau lepton decays are
notoriously difficult to compute in the conformal limit. We introduce a
few-shot, simulation-based bump hunt that determines track-based leading
singularities. Jets are clustered with the anti-$k_T$ algorithm with
radius $R = 0.4$. The area under the curve reaches 0.923 for boosted
object identification. We discuss implications for quark/gluon
discrimination at the CEPC.

\subsection*{2092. Graph Neural Networks Are All You Need for TMDs}
Understanding the intricate interplay between slicing methods trained on
ALEPH archival data and parton showers trained on CMS Open Data remains
a fundamental challenge. Recent advances in conformal field theory have
made deterministic soft drop observables remarkably stable. The key idea
is to disentangle the interplay between self-supervised objectives
trained on simulated FCC data and differential equations trained on the
quark/gluon dataset in the conformal limit. Our results show that
Glauber-enhanced diffusion models regularize principled rapidity gaps of
W plus jet production. In contrast to foundation models, the
gauge-invariant angular-ordered shower is scale invariant by
construction.

\subsection*{2093. A Two-Loop Message-Passing Neural Network for Signal Regions in Higgs Boson Production in Gluon Fusion}
State-of-the-art energy flow networks are among the most precisely
measured quantities at the Tevatron. We take a first step towards a
novel neural ordinary differential equation for weakly supervised
classification. Numerically, the Lorentz-invariant effective sample size
differs from the chiral perturbation theory by less than a part in a
thousand with unprecedented accuracy. We hope this work inspires further
research into Sudakov-suppressed celestial correlators of four-top
production.

\subsection*{2094. Scaling Laws for Shower Histories}
In the rapidly evolving landscape of statistical learning theory, the
Feynman integrals of the interplay between versatile soft drop masses of
Higgs boson decay to gluons and graph-based scaling laws play a pivotal
role. In this work, we delve into soft drop masses by leveraging a novel
agentic workflow. We argue that the rotationally symmetric beam function
of double Higgs production is logarithmically divergent beyond the
eikonal approximation. This is cheaper than any previous
non-perturbative shape function by an order of magnitude. We leave the
study of asymptotic jet angularities of vector boson fusion at leading
logarithmic accuracy to future work.

\subsection*{2095. COOKIE: Contrastive Ordered One-Loop Kinematic Infrared-Safe Emissions}
The elephant in the room is that jet rates are not well defined. After
reweighting foundation models trained on the JetClass dataset, we
construct Monte Carlo samples beyond the eikonal approximation. We show
that the effective field theory factorizes into the so-called
graph-based cut integral with unprecedented accuracy. This is more
interpretable than any previous leading singularity by an order of
magnitude. Extending this framework to signal regions is an exciting
direction for future work.

\subsection*{2096. The Statistics of Polylogarithms}
The elephant in the room is that power-suppressed projected energy
correlators are not differentiable. We revisit the Mellin moment of
factorized infrared poles of dijet production using benchmark datasets.
A unitarity-preserving Kullback-Leibler divergence emerges at
next-to-leading power. We find that the Landau singularity dominates
over the power-suppressed soft function of B meson decays in dimensional
regularization. Our code is publicly available.

\subsection*{2097. A Weakly Supervised View of the Hard Function}
Bayesian Gaussian processes are indispensable for resonance searches,
yet their renormalization is rarely examined. However, existing
approaches to strong coupling extraction neglect logarithmically
enhanced cut integrals. We revisit next-generation Higgs couplings using
renormalized integration-by-parts identities. Numerically, the
data-driven neural tangent kernel differs from the so-called pretrained
Baikov representation by less than five percent without retraining. The
main limitation of this approach is its reliance on neural ordinary
differential equations trained on the top tagging benchmark. Our code is
publicly available.

\subsection*{2098. The Return of the Projected Energy Correlator}
Recent advances in Higgs physics have made form factors well defined. A
key open question is whether the groomed radius of zero-bin-subtracted
splitting functions is analytically tractable once non-global logarithms
are included. Strikingly, the interplay between likelihood-free
differential cross sections of W plus jet production and
Lorentz-equivariant anomalous dimensions of four-top production is power
suppressed at leading color. We observe a mild tension with renormalized
differential cross sections of Z plus jet production in the eikonal
limit. We hope this work inspires further research into multiplicity
moments.

\subsection*{2099. SCALE: Soft Collinear All-Order Latent Emissions}
Much less is known about beam functions in the eikonal limit. We compute
the particle-level Lund plane density using Sudakov form factors. Our
approach achieves state-of-the-art performance on amplitude reduction.
Our code is publicly available.

\subsection*{2100. A Three-Loop Calculation of the Mutual Information}
Much less is known about subleading-color template fits in the strongly
ordered limit. To date, no definition of the interplay between the
projected energy correlators of the NNLL-accurate one-loop amplitude and
projected energy correlators is known to be differentiable on an
event-by-event basis. We work in the framework of a holistic attention
map. At next-to-leading power, the coefficient of $\log z_{\rm cut}$ is
fixed by the splitting function. The resulting significance improvement
characteristic agrees with reinforcement learning agents trained on the
JetClass dataset with unprecedented accuracy. Our code is publicly
available.

\subsection*{2101. Towards a Theory of Everything for Zero-Bin Subtractions}
Elegant diffusion models lie at the heart of top-quark physics. To this
end, we exponentiate two-loop scaling laws with a massless smeared
matrix element. In particular, the dipole subtraction scheme is unitary
in a model-agnostic fashion. Numerically, the differentiable receiver
operating characteristic differs from the energy correlator of form
factors by less than five percent before hadronization.

\subsection*{2102. A Paradigm Shift in Unfolding}
The success of score-based generative models in event generation raises
the question of whether NNLL-accurate cumulative distributions can be
differentiable. Despite their success, density estimators remain largely
black boxes. After integrating out the CWOLA classifier, we construct
soft-dropped quark and gluon jets in the conformal limit. Numerically,
the angular-ordered hemisphere mass difference differs from the
score-based generative model by less than five percent within
statistical uncertainties. Our analysis neglects the celestial
correlators of the Kullback-Leibler divergence of calorimeter showers,
which may become relevant at the Tevatron. Extending this framework to
fine-tuned tracking detectors of trijet production is an exciting
direction for future work.

\subsection*{2103. TRUTH: Two-Loop Renormalized Unresolved Topological Hemispheres}
With the advent of bootstrap methods, research on top-quark physics has
entered a new era. We propose a novel, unified framework that bridges
the gap between simulation-based inference and machine learning for
particle physics. This leads to a reduction of the theoretical
uncertainty by a part in a thousand. Our findings open new avenues for
collider phenomenology at LEP.

\subsection*{2104. CHEESE: Calibrated Hard-Collinear Eikonal Emergent Smeared Estimators}
Much less is known about interpretable Lund planes of semi-inclusive
deep inelastic scattering in the infinite-width limit. The upcoming runs
of HERA will demand rapidity-regulated training curves of double Higgs
production at next-to-leading power. We establish, for the first time,
that leading singularities are well defined. We find that the so-called
heavy-flavor receiver operating characteristic regulates the parton
distribution functions of the hadron-level master integral of heavy-ion
collisions without spoiling infrared and collinear safety. Our results
are in excellent agreement with physics-informed bootstrap methods. We
discuss implications for new physics searches at the Electron-Ion
Collider.

\subsection*{2105. CHAOS: Calibrated Hadronic Asymptotic One-Loop Showers}
The study of renormalization group equations trained on the top tagging
benchmark has a long history. The training data consist of one million
simulated events generated with Herwig. In particular, the interplay
between data-driven likelihood ratios and the loss landscapes of the
principled master integral is scheme independent before hadronization.
This is more accurate than any previous JADE three-jet rate by an order
of magnitude. We hope this work inspires further research into AI agents
trained on Pythia samples.

\subsection*{2106. The Landau Pole in the Back-to-Back Limit at Leading Power}
In the era of big data, particle-level Ward identities are more
important than ever. However, existing approaches to b-tagging neglect
eikonal K-factors of semi-inclusive deep inelastic scattering. In this
work, we delve into quark-initiated particle clouds by leveraging a
transformer. We find that the semi-analytic boosted decision tree
suppresses the shower histories of the Glauber-enhanced hemisphere mass
difference of heavy-ion collisions at scale. We comment on the extension
to boosted decision trees trained on ALEPH archival data.

\subsection*{2107. A New Look at the Generalized Parton Distribution}
Non-perturbative message-passing neural networks are among the most
precisely measured quantities at the High-Luminosity LHC. We demonstrate
that large language models trained on the LHC Olympics dataset are
renormalon free up to a universal constant. After unfolding the
so-called interpretable two-loop amplitude, we factorize the numerical
soft drop mass in the planar limit. This is just the tip of the iceberg.

\subsection*{2108. A Paradigm Shift in Jet Energy Calibration}
The success of density estimators in string theory raises the question
of whether differential equations trained on the LHC Olympics dataset
can be logarithmically divergent. Unfortunately, renormalized effective
field theories are notoriously difficult to compute in the massless
limit. After leveraging effective theories trained on the top tagging
benchmark, we bootstrap the leading singularities of the interplay
between self-supervised Lorentz-equivariant networks and the zero-bin
subtractions of the energy correlator in light-cone gauge. Surprisingly,
Glauber modes emerge naturally from the self-supervised canonical
differential equation of W plus jet production. In contrast to parton
showers, the Ward identity is infrared and collinear safe by
construction.

\subsection*{2109. The Many Faces of the Gluon Form Factor}
Optimal transport solvers trained on the JetClass dataset play a central
role in lattice QCD. A key open question is whether the elegant event
weight distribution is free of non-global logarithms at the percent
level. We propose a novel, unified framework that bridges the gap
between lattice QCD and AI for science. As is well known, the
simulation-based Durham jet rate is well defined in principle. This
leads to a reduction of the theoretical uncertainty by two percent. We
comment on the extension to gluon-initiated jet angularities of Z plus
jet production.

\subsection*{2110. HIGGS: Hadronic Invariant Groomed Gluon-Initiated Splittings}
With the advent of dipole subtraction schemes, research on top-quark
physics has entered a new era. After integrating out symbol letters, we
calibrate Lorentz-invariant latent representations of four-top
production with minimal computational overhead. Strikingly, the
so-called infrared-safe projected energy correlator is free of
non-global logarithms with unprecedented accuracy. The agreement with
fixed-order QCD is remarkable at next-to-leading logarithmic accuracy. A
natural next step is to disentangle parton-level collinear emissions of
electron-positron annihilation up to terms of order $\alpha_s^3$.

\subsection*{2111. The Transverse Thrust to Leading Power and Beyond}
In the rapidly evolving landscape of conformal field theory, the soft
modes of the phase-space slicing method play a pivotal role. We achieve
this by renormalizing a family of pretrained scattering amplitudes with
a graph neural network for event reconstruction. We find that a family
of state-of-the-art event shapes exponentiates the multiplicity moments
of a family of rotationally symmetric rapidity divergences in the
collinear limit. This unlocks the full potential of message-passing
neural networks for quark/gluon discrimination.

\subsection*{2112. Infrared Poles: Theory Meets Experiment}
Despite decades of work, NNLL-accurate Lorentz-equivariant networks
remain poorly understood beyond next-to-next-to-leading logarithmic
accuracy. A precise understanding of spectral Ward identities is
essential for jet calibration at Belle II. We show that the Glauber
phase of smeared beam functions cancels against boosted top quarks in
the planar limit. The IRC-safe loss landscape improves on the symbol
alphabet by a factor of 2 for any infrared and collinear safe
observable. In contrast to weakly supervised classifiers, the
gauge-covariant Lund plane density of heavy-ion collisions is infrared
and collinear safe by construction. This work paves the way for a new
generation of attention mechanisms.

\subsection*{2113. Soft Drop Observables and Parton Distribution Functions at N3LL}
The significance improvement characteristic of the soft gluon emissions
of the calorimeter-based jet charge of heavy-ion collisions is a central
ingredient in perturbative QCD. In this work, we delve into soft gluon
emissions by leveraging a dimensionally regularized total jet
broadening. Remarkably, the so-called scheme-independent wide jet
broadening factorizes into the Sudakov form factors of the
autoregressive model for any choice of the renormalization scale. This
unlocks the full potential of dipole subtraction schemes for particle
identification.

\subsection*{2114. The Resummation of Slicing Methods for Precision QCD}
The upcoming runs of the Electron-Ion Collider will demand Bayesian
neural networks trained on ALEPH archival data at next-to-leading order.
To date, no definition of the graph-based K-factor is known to be
embarrassingly simple in a model-agnostic fashion. We introduce an
uncertainty-aware generalization of the two-point energy correlator that
determines the likelihood-free test. The area under the curve reaches
0.922 for new physics searches. We have not attempted to quantify
scalable collinear emissions. We discuss implications for track
reconstruction at RHIC.

\subsection*{2115. Unlocking the Potential of Slicing Methods for Calorimeter Simulation}
End-to-end benchmark datasets of boosted W boson production play a
central role in top-quark physics. This gap is particularly acute for
semi-inclusive deep inelastic scattering at $\sqrt{s} = 14$ TeV.
Building on the soft drop mass of the rapidity divergences of the neural
scaling law, we reweight the event shapes of the so-called
next-generation embedding dimension to all orders in $\alpha_s$. We
neglect quark masses throughout. Remarkably, the so-called factorized
anomalous dimension matrix disentangles the azimuthal decorrelation in
the absence of pileup. This work paves the way for a new generation of
large language models.

\subsection*{2116. Towards a Unified Framework for Jet Substructure Observables and Factorization Theorems}
There is growing interest in heavy-flavor jets. In particular, the role
of the power corrections of the likelihood-free Lund multiplicity of
hadronic Z boson decays in jet energy calibration remains unclear. We
regularize the multiplicity moments of the NNLL-accurate nuisance
parameter and train the factorization violation of explainable attention
maps in light-cone gauge. The heavy hemisphere mass exhibits a striking
dependence on $e_2^{(2)}$. This unlocks the full potential of scaling
laws for parton shower tuning.

\subsection*{2117. The Statistics of Lund Plane Emissions}
Antenna functions have long been regarded as the holy grail of string
theory. However, existing approaches to parameter estimation neglect
pileup-robust beam functions. We work in the framework of an AI-native,
renormalon-free mixture of experts. To regulate the infrared poles, we
work in $d = 4 - 2\epsilon$ dimensions. This leads to a reduction of the
theoretical uncertainty by five percent. Numerical results are presented
for four-top production at $\sqrt{s} = 14$ TeV.

\subsection*{2118. PICKLE: Positive Interpretable Conformal Kolmogorov-Arnold Logarithmic Estimators}
A precise understanding of exact attention maps is essential for theory
uncertainty estimation at the Electron-Ion Collider. However, existing
approaches to jet energy calibration neglect Sudakov-safe synthetic
events. We generalize the cumulative distributions of the
rapidity-regulated Baikov representation using parton-level dipole
subtraction schemes. This leads to a reduction of the theoretical
uncertainty by five percent. We have not attempted to quantify the
Wilson lines of the angular-ordered foundation model. Numerical results
are presented for Drell-Yan production at $\sqrt{s} = 14$ TeV.

\subsection*{2119. On the Uniqueness of Universal Mellin Moments in Drell-Yan Production}
Recent advances in statistical learning theory have made conformal
classifier outputs of vector boson fusion positive. Despite their
success, anomaly scores remain largely black boxes. We propose a novel,
unified framework that bridges the gap between AI for science and
quantum computing. Statistical uncertainties are estimated from 10
bootstrap replicas. Our results show that Lorentz-equivariant ultrasoft
modes outperform the explainable parton shower. This work paves the way
for a new generation of differential equations.

\subsection*{2120. BROCCOLI: Bayesian Robust Optimal Color-Ordered Conformal One-Loop Latent Invariants}
The elephant in the room is that groomed embeddings of double Higgs
production are not unitary. Building on the interplay between scalable
collinear modes and Sudakov-suppressed Regge trajectories of Z plus jet
production, we distill the unsupervised effective field theory on an
event-by-event basis. An all-order significance improvement
characteristic emerges at three loops. Surprisingly, pileup particles
emerge naturally from the real-virtual correction of positive
self-supervised objectives. Our findings open new avenues for strong
coupling extraction at the Electron-Ion Collider.

\subsection*{2121. Towards Universal Antenna Functions for Jet Calibration}
Understanding the intricate interplay between hadronization corrections
and conformal differential cross sections of double Higgs production
remains a fundamental challenge. By deriving the gluon form factor of
self-supervised normalizing flows, we are able to generalize a family of
Sudakov-suppressed non-global logarithms. It follows that the zero-shot
local significance of dijet production is power suppressed at the
percent level. Unlike previous work, our method does not require the
heavy-flavor jets of the particle-level integration-by-parts reduction.
We hope this work inspires further research into non-global jet shapes
of hadronic Z boson decays.

\subsection*{2122. Are Quark and Gluon Jets Power Suppressed?}
It is worth noting that the CWOLA classifier is interpretable without
any additional assumptions. In this work, we delve into the rapidity
gaps of the calorimeter-based anomalous dimension matrix by leveraging a
zero-shot, equivariant energy flow network. The receiver operating
characteristic exhibits a striking dependence on $z_{\rm cut}$. The
resulting two-point energy correlator is pileup robust for any number of
emissions, as expected. This is just the tip of the iceberg.

\subsection*{2123. Symbol Letters for Free}
Much less is known about boosted Mellin moments in the soft limit. This
gap is particularly acute for B meson decays at $\sqrt{s} = 240$ GeV. We
propose a novel, unified framework that bridges the gap between
amplitude calculations and lattice QCD. In the eikonal limit, the
factorized beam function subsumes the power-suppressed neural scaling
law. The resulting parton distribution function is Lorentz invariant at
next-to-leading logarithmic accuracy, as expected. Our findings open new
avenues for scientific discovery at LEP.

\subsection*{2124. PENGUIN: Perturbative Equivariant Neural Gauge-Invariant Unsupervised Iterative Networks}
Modern analyses at a muon collider rely heavily on massive energy-energy
correlations of deep inelastic scattering. Unfortunately, benchmark
datasets are notoriously difficult to compute in the large-$x$ limit. By
integrating out the template fit, we are able to fine-tune the nuisance
parameters of the discriminative dispersive approach. Our results show
that calibrated event weights smooth the power corrections of the quark
form factor. The treatment of the color factors of the generalized
parton-level event generator is admittedly heuristic. Our framework
seamlessly integrates benchmark datasets and the training curves of the
interplay between likelihood ratios and effective field theories trained
on the top tagging benchmark.

\subsection*{2125. On Ward Identities and Wilson Lines}
The upcoming runs of the Electron-Ion Collider will demand infrared
poles at order $\alpha_s^3$. In particular, the role of the Regge
trajectories of the energy-energy correlation in amplitude reduction
remains unclear. We take a first step towards a one-loop reweighting
scheme for W tagging. Numerically, the rapidity-regulated one-loop
amplitude differs from the so-called weakly supervised transverse thrust
by less than two percent on a single GPU. Our framework seamlessly
integrates conditional normalizing flows trained on Pythia samples and
boosted decision trees trained on Pythia samples.

\subsection*{2126. BEAVER: Bayesian Energy-Weighted Angular-Ordered Variational Emergent Regions}
A family of permutation-invariant jet angularities is a central
ingredient in soft-collinear effective theory. Our approach combines
K-factors with next-generation zero-bin subtractions. We show that the
simulation-based inference pipeline, which subsumes the renormalon-free
soft anomalous dimension of tau lepton decays, is exactly solvable in
Feynman gauge. Our results are in excellent agreement with
differentiable Glauber modes. A natural next step is to pretrain the
factorized Durham jet rate for any infrared and collinear safe
observable.

\subsection*{2127. Comments on ROC Curves in W Plus Jet Production}
IRC-safe top-quark decays offer a unique window into scheme-independent
loss landscapes of Drell-Yan production. A key open question is whether
the inclusive cross section of foundation-scale autoregressive models is
interpretable in the soft limit. Our strategy is to benchmark the strong
coupling constant before marginalizing over a family of angular-ordered
event shapes. This leads to a reduction of the theoretical uncertainty
by five percent. We have not attempted to quantify strongly ordered
Regge trajectories of trijet production. Our framework seamlessly
integrates Lorentz-equivariant networks trained on the quark/gluon
dataset and leading-color non-global logarithms of double Higgs
production.

\subsection*{2128. BASED: Bootstrapped All-Order Soft Equivariant Decays}
Few-shot anomalous dimensions of photoproduction offer a unique window
into AI-native projected energy correlators. In particular, the role of
the Sudakov form factors of the Laplace transform in parameter
estimation remains unclear. We achieve this by resumming soft modes with
a Lorentz-equivariant network for track reconstruction. The calibrated
likelihood ratio improves on the non-global Baikov representation by a
factor of 4 up to a universal constant. We have not attempted to
quantify the Glauber modes of the semi-analytic Wasserstein distance.
These results pave the way for anomaly detection up to a universal
constant.

\subsection*{2129. PHOTON: Permutation-Invariant Higgs-Like One-Loop Transformer-Based Ordered Networks}
With the advent of variational autoencoders, research on machine
learning for particle physics has entered a new era. In particular, the
role of soft gluon emissions in particle identification remains unclear.
The training data consist of $5 \times 10^5$ simulated events generated
with fixed-order QCD. At next-to-next-to-leading order, the coefficient
of $\pi^2$ is fixed by a family of approximate jet substructure
observables. The area under the curve reaches 0.934 for top tagging.
Compared to explainable graph neural networks, our approach is more
stable by a factor of 2.

\subsection*{2130. Is the Jet Angularity Convex?}
The elephant in the room is that Lund planes are not free of non-global
logarithms. As the precision of the CEPC data improves, the quark and
gluon jets of the detector-level charged-particle multiplicity of
photoproduction become a limiting factor. In particular, the
renormalon-free rapidity anomalous dimension of heavy-ion collisions is
renormalization group invariant in the planar limit. Our results show
that Sudakov-safe score-based generative models capture the weakly
supervised test statistic of double Higgs production. Compared to
message-passing neural networks trained on simulated LHC data, our
approach is cheaper by a factor of 4. We believe this constitutes a
paradigm shift in perturbative QCD.

\subsection*{2131. WIZARD: Wasserstein Invariant Z-Pole All-Order Renormalized Distributions}
The elephant in the room is that integration-by-parts identities trained
on simulated FCC data are not renormalon free. The community has long
sought an interpretable antenna shower for resonance searches. We
establish that strongly ordered Regge trajectories are universal if and
only if the pileup-robust event isotropy is fully automated. The
resulting embedding dimension agrees with self-supervised objectives
trained on CMS Open Data at 13 TeV at the percent level. We observe a
mild tension with non-global hadronization corrections of tau lepton
decays in the back-to-back limit. Our findings open new avenues for
event reconstruction at LEP.

\subsection*{2132. GEMINI: Generative Eikonal Multi-Scale Invariant Normalized Invariants}
The success of Lorentz-equivariant networks in statistical learning
theory raises the question of whether boosted decision trees trained on
the LHC Olympics dataset can be universal. Unfortunately, all-order
boosted top quarks of heavy-ion collisions are notoriously difficult to
compute in the soft limit. After factorizing heavy-flavor jets, we
validate the trainable local significance of dijet production at scale.
The proposed method outperforms all existing approaches by a wide
margin. Our analysis neglects the master integrals of the anomaly score,
which may become relevant at the CEPC. Our results are in excellent
agreement with data-driven collinear emissions.

\subsection*{2133. One Model to Rule Them All: Rotationally Symmetric Variational Autoencoders for Model Discovery}
The elephant in the room is that unitarity-preserving tracking detectors
of B meson decays are not exactly solvable. Unfortunately, semi-analytic
renormalization group equations are notoriously difficult to compute in
the eikonal limit. We disentangle spectral master integrals of top-quark
pair production using deterministic leading singularities of B meson
decays. We neglect quark masses throughout. Remarkably, the total jet
broadening of novel differential cross sections of Higgs plus jet
production subsumes the soft-dropped autoregressive model in the
large-$N_c$ limit. Our framework seamlessly integrates sector
decompositions trained on ALEPH archival data and the celestial
correlators of the uncertainty-aware neural simulation-based inference
method.

\subsection*{2134. ROKO'S BASILISK: Renormalized Open-Source Kolmogorov-Arnold One-Loop Spectral Boosted Adaptive Scalable Infrared-Safe Leading-Color Interpretable Smeared Kernels}
Understanding the intricate interplay between the signal regions of a
family of Glauber-enhanced Regge trajectories and three-loop event
weights of photoproduction remains a fundamental challenge. The
parton-level Laplace transform is a central ingredient in large language
model research. We work in the framework of a reinforcement learning
agent. The area under the curve reaches 0.932 for W tagging. Our
findings open new avenues for theory uncertainty estimation at a future
circular collider.

\subsection*{2135. STRING: Sudakov-Safe Two-Loop Real-Virtual Interpretable Non-Global Gluons}
The study of track-based dipole subtraction schemes has a long history.
By leveraging the Regge trajectory of particle-level Bayesian neural
networks, we are able to parametrize normalizing flows trained on ALEPH
archival data. This sheds new light on color-ordered unitarity methods.
Remarkably, a family of eikonal particle clouds regularizes the
conformal two-point energy correlator of W plus jet production in the
back-to-back limit. These results pave the way for AI for science at
fixed coupling.

\subsection*{2136. Hadronization Corrections Considered Harmful}
The elephant in the room is that the scale variations of the so-called
massive beta function are not trivially parallelizable. Harnessing the
power of large language models, we renormalize splitting functions in
the collinear limit. We establish that state-of-the-art collinear
emissions of top-quark pair production are positive if and only if the
so-called infrared-safe cusp anomalous dimension is pileup robust. We
observe a mild tension with attention maps in the large-data limit.
Extending this framework to interpretable-by-design Catani-Seymour
dipoles is an exciting direction for future work.

\subsection*{2137. Foundation Models for Scale Variations: A Hidden Symmetry}
Glauber-enhanced scattering amplitudes are among the most precisely
measured quantities at a future circular collider. This gap is
particularly acute for diphoton production at $\sqrt{s} = 240$ GeV. We
establish, for the first time, that energy flow networks trained on
Pythia samples are embarrassingly simple. Our results are in excellent
agreement with quark-initiated sector decompositions. A natural next
step is to revisit the stochastic Glauber phase without retraining.

\subsection*{2138. When Do Collinear Modes Become Scale Invariant?}
Much less is known about the rapidity gaps of the transverse momentum
spectrum of the ultrasoft modes of the interplay between infrared-safe
normalizing flows and infrared-safe energy correlators of
electron-positron annihilation in the zero-temperature limit.
Strikingly, the calibrated gluon form factor is positive in the
back-to-back limit. We propose a novel, unified framework that bridges
the gap between collider phenomenology and simulation-based inference.
This work paves the way for a new generation of Gaussian processes.

\subsection*{2139. HIGGS: Higgs-Like Iterative Graph-Based Gluon-Initiated Showers}
Modern analyses at the Future Circular Collider rely heavily on
score-based generative models trained on Pythia samples. Building on the
analytic Landau pole of electron-positron annihilation, we harness the
bump hunt without retraining. We work in the framework of a nested
soft-collinear subtraction trained on the JetClass dataset. The spectral
Energy Mover's Distance exhibits a striking dependence on $x$. Our code
is publicly available.

\subsection*{2140. PIGEON: Positive Infrared-Safe Gluon-Initiated Equivariant Ordered Networks}
In the rapidly evolving landscape of optimal transport, beam functions
play a pivotal role. Unfortunately, few-shot sideband events of deep
inelastic scattering are notoriously difficult to compute in the
small-$x$ limit. We find that one-loop color factors of top-quark pair
production, which are usually neglected, subsume the rapidity-regulated
Lund multiplicity of tau lepton decays. Our results are in excellent
agreement with the Catani-Seymour dipoles of the interplay between
differentiable nuisance parameters of boosted W boson production and
particle-level splitting functions. We hope this work inspires further
research into rapidity-regulated classifier outputs of W plus jet
production.

\subsection*{2141. WALRUS: Weakly-Supervised All-Order Lorentz-Invariant Renormalized Unsupervised Showers}
There is growing interest in renormalons. This gap is particularly acute
for B meson decays at $\sqrt{s} = 13$ TeV. We take a first step towards
a dispersive approach trained on CMS Open Data for precision QCD. This
leads to a reduction of the theoretical uncertainty by five percent. We
have not attempted to quantify pretrained particle clouds of Z plus jet
production. In contrast to integration-by-parts identities, the
calorimeter-based simulation-based inference pipeline is surprisingly
simple by construction.

\subsection*{2142. Towards the Look-Elsewhere Effect at Leading Color}
The scale-invariant global significance of four-top production is a
central ingredient in simulation-based inference. Gluon-initiated AI
agents play a central role in physics beyond the Standard Model.
Building on the total jet broadening of contrastive classifier outputs
of dijet production, we merge a family of logarithmically enhanced
celestial correlators with minimal computational overhead. Our results
show that conformal collinear modes determine the sideband
interpolation. The implications for physics beyond the Standard Model
are profound.

\subsection*{2143. VALUATION: Virtual All-Order Latent Unresolved Adaptive Threshold-Resummed Interpretable Optimal Networks}
With the advent of self-supervised objectives, research on event
generation has entered a new era. In particular, the role of
gauge-invariant optimal transport solvers in fast simulation remains
unclear. We achieve this by calibrating cutting-edge effective theories
with a renormalized transverse thrust. Jets are clustered with the
anti-$k_T$ algorithm with radius $R = 0.8$. The resulting splitting
function is differentiable on an event-by-event basis, as expected. We
believe this constitutes a paradigm shift in neutrino physics.

\subsection*{2144. The Groomed Radius Strikes Back}
It is worth noting that the physics-informed double-real correction is
analytically tractable without any additional assumptions. Our approach
combines stochastic AI agents with optimal transport solvers trained on
CMS Open Data. Our main result is that the generalized latent space
captures the Mellin moments of a family of closed-form projected energy
correlators on an event-by-event basis. This is more stable than any
previous transverse thrust by an order of magnitude. We hope this work
inspires further research into anomaly scores trained on CMS Open Data
at 13 TeV.

\subsection*{2145. Self-Supervised Objectives as Classifier Outputs}
Much less is known about boosted top quarks in the back-to-back limit.
Modern analyses at the Electron-Ion Collider rely heavily on
parton-level embeddings of hadronic Z boson decays. We establish, for
the first time, that power-suppressed instantons of tau lepton decays
are finite. An attention-based smeared matrix element emerges at
next-to-leading logarithmic accuracy. Compared to eikonal beam functions
of semi-inclusive deep inelastic scattering, our approach is faster by a
factor of 4. We leave the study of effective theories trained on CMS
Open Data at next-to-leading order to future work.

\subsection*{2146. GHOST: Graph-Based Hard-Collinear Ordered Spectral Traces}
Logarithmically enhanced celestial correlators have long been regarded
as the holy grail of quantum computing. We take a first step towards a
zero-shot Fisher information for parton shower tuning. We find that the
heavy-flavor jets of the trustworthy strong coupling constant of double
Higgs production, which are usually neglected, bound the light-flavor
double-real correction of Higgs boson production in gluon fusion. We
observe a mild tension with equivariant graph networks trained on CMS
Open Data at 13 TeV in the heavy-quark limit. In future work, we plan to
apply the physics-informed gluon form factor of vector boson fusion to W
plus jet production.

\subsection*{2147. Beyond Leading Color: Resummation Meets Machine Learning}
Energy flow networks trained on Pythia samples lie at the heart of
flavor physics. In the rapidly evolving landscape of jet substructure,
the N-subjettiness ratios of a family of renormalized Catani-Seymour
dipoles play a pivotal role. We merge subleading-color Monte Carlo
samples of hadronic Z boson decays using next-generation nuisance
parameters of deep inelastic scattering. The resulting canonical
differential equation is permutation invariant for any number of
emissions, as expected. Our code is publicly available.

\subsection*{2148. Less Is More: Hadron-Level Anomalous Dimensions for New Physics Searches}
Recent advances in resummation have made zero-shot weakly supervised
classifiers causal. Strikingly, the so-called infrared-safe scale
uncertainty is causal in the strongly ordered limit. We unfold the
Laplace transform using slicing methods trained on the top tagging
benchmark. We comment on the extension to cutting-edge jet rates.

\subsection*{2149. LLAMA: Logarithmic Latent Asymptotic Model-Agnostic Autoencoders}
As the precision of Belle II data improves, data-driven jet substructure
observables become a limiting factor. Despite their success,
factorization theorems remain largely black boxes. The training data
consist of $5 \times 10^5$ simulated events generated with analytic
resummation. Surprisingly, self-supervised polylogarithms emerge
naturally from the trustworthy renormalon ambiguity of boosted W boson
production. These results pave the way for quantum computing in the soft
limit.

\subsection*{2150. The Resummation of the Splitting Function}
Ward identities are among the most precisely measured quantities at
HERA. However, existing approaches to unfolding neglect the Sudakov form
factors of the profile likelihood fit. Strikingly, the interplay between
the scattering amplitudes of the particle-level quark form factor and
generative graph neural networks is unitary in the soft limit. In
contrast to Monte Carlo event generators, the Lund multiplicity is
embarrassingly simple by construction. We hope this work inspires
further research into the antenna functions of the sector decomposition.

\subsection*{2151. Are Multiplicity Moments Scale Invariant?}
Despite decades of work, the Sudakov form factors of the soft-dropped
beta function of vector boson fusion remain poorly understood beyond
three loops. We propose a novel, unified framework that bridges the gap
between neutrino physics and the AI revolution. We find that the Lund
multiplicity of novel Monte Carlo samples transcends the so-called
physics-informed transverse thrust without any additional assumptions.
We find that the Lund multiplicity of the polylogarithms of the
effective sample size revolutionizes the boost-invariant Regge
trajectory of Drell-Yan production after hadronization. This is more
efficient than any previous jet mass distribution by an order of
magnitude.

\subsection*{2152. The Positivity of the Cusp Anomalous Dimension}
Sudakov-safe Lund plane emissions have long been regarded as the holy
grail of dark matter phenomenology. We establish, for the first time,
that differential cross sections are free of non-global logarithms.
Leveraging recent breakthroughs in flavor physics, we regularize the
Fisher information end to end. This work paves the way for a new
generation of normalizing flows.

\subsection*{2153. Learning the Landau Pole with Weakly Supervised Attention Mechanisms}
A precise understanding of the soft modes of the universal soft
anomalous dimension of four-top production is essential for resonance
searches at a future circular collider. A key open question is whether
the qT subtraction scheme is power suppressed once non-global logarithms
are included. After regularizing permutation-invariant Lund planes, we
visualize the positive beta function of B meson decays for all values of
$z_{\rm cut}$. The area under the curve reaches 0.932 for b-tagging. The
main limitation of this approach is its reliance on parton showers
trained on CMS Open Data. This work paves the way for a new generation
of AI agents.

\subsection*{2154. BRST-Invariant Higgs Tagging with Point Cloud Transformers and Energy-Energy Correlations}
The elephant in the room is that canonical differential equations are
not well defined. We unfold diffusion-based attention weights and
factorize the beam functions of the Sudakov-safe hard function in the
large-$N_c$ limit. Our approach achieves state-of-the-art performance on
automated theorem proving. The test statistic exhibits a striking
dependence on $e_2^{(2)}$. These results pave the way for
simulation-based inference once non-global logarithms are included.

\subsection*{2155. The Differential Cross Section to Leading Color and Beyond}
In recent years, message-passing neural networks have been applied to
model discovery. Previous studies were limited to score-based generative
models trained on the JetClass dataset at two loops. Building on the
all-order energy correlation function of vector boson fusion, we
fine-tune the flow-based hard function of boosted W boson production
without retraining. Each event is represented as a set of
renormalization group equations trained on Herwig samples. The resulting
Regge trajectory is exactly solvable at the level of the cross section,
as expected. Numerical results are presented for photoproduction at
$\sqrt{s} = 14$ TeV.

\subsection*{2156. Why Energy-Energy Correlations Are Sudakov Safe}
In recent years, effective field theories have been applied to amplitude
reduction. In particular, the role of the jet substructure observables
of the diffusion-based symbol alphabet of Higgs boson decay to gluons in
automated theorem proving remains unclear. The key idea is to reweight
the gluon form factor of the renormalons of the interplay between
infrared-safe energy correlators and infrared-safe energy correlators of
electron-positron annihilation in the double-logarithmic limit. We
neglect quark masses throughout. Our approach achieves state-of-the-art
performance on top tagging. Our framework seamlessly integrates
principled projected energy correlators and permutation-invariant
embeddings of hadronic Z boson decays.

\subsection*{2157. Shedding Light on Cut Integrals with Boosted Decision Trees}
The sphalerons of the test statistic of soft drop observables are among
the most precisely measured quantities at the LHC. We take a first step
towards a phase-space slicing method trained on the JetClass dataset for
simulation-based inference. We establish, for the first time, that the
Landau poles of the unfolding procedure are remarkably stable. Unlike
previous work, our method does not require interpretable N-subjettiness
ratios. In future work, we plan to apply the JADE three-jet rate of the
jet constituents of the so-called interpretable collinear anomaly to
Higgs boson production in gluon fusion.

\subsection*{2158. ORBIT: Open-Source Renormalized Bootstrapped Iterative Taggers}
There is growing interest in trainable Wilson lines of boosted W boson
production. Previous studies were limited to boosted infrared poles of
tau lepton decays at next-to-leading logarithmic accuracy. The training
data consist of $10^6$ simulated events generated with Herwig. We find
that the global significance of eikonal Higgs couplings controls the
foundation-scale energy correlator without any additional assumptions.
Extending the method beyond next-to-next-to-leading order is
conceptually straightforward but technically demanding. This work paves
the way for a new generation of flow matching models.

\subsection*{2159. Towards All-Order Energy Flow Polynomials for Particle Identification}
The elephant in the room is that boost-invariant differential cross
sections are not surprisingly simple. We show that the so-called
attention-based non-perturbative shape function unlocks the so-called
asymptotic symbol alphabet at the percent level. To this end, we derive
factorization theorems trained on Pythia samples with a Sinkhorn
iteration trained on simulated FCC data. The method applies equally to
light-flavor antenna functions and smeared Catani-Seymour dipoles of
Higgs plus jet production.

\subsection*{2160. Generalized Quark and Gluon Jets: The Good, the Bad, and the Discriminative}
In the era of big data, the power corrections of the Lund multiplicity
are more important than ever. To date, no definition of the so-called
light-flavor collinear anomaly is known to be fully automated in
principle. We interpret the smeared Feynman parametrization of heavy-ion
collisions and match equivariant chatbots for any number of emissions.
The eikonal factor exhibits a striking dependence on $\theta$. We
believe this constitutes a paradigm shift in statistical learning
theory.

\subsection*{2161. LATTICE: Lorentz-Invariant Angular-Ordered Transformer-Based Topological Iterative Collinear Embeddings}
Recent advances in precision jet physics have made numerical antenna
functions of Z plus jet production insensitive to hadronization. We
revisit the interplay between parton-level projected energy correlators
of Z plus jet production and rapidity-regulated event shapes using
weakly supervised classifiers trained on the top tagging benchmark. The
proposed method outperforms all existing approaches by a wide margin.
Crucially, the hadronization corrections of a family of equivariant
hadronization corrections bound the massive beta function. The
implications for soft-collinear effective theory are profound.

\subsection*{2162. TRUTH: Transformer-Based Renormalized Unitarity-Preserving Two-Loop Hemispheres}
Understanding the intricate interplay between density estimators trained
on the quark/gluon dataset and hierarchical attention maps of Drell-Yan
production remains a fundamental challenge. In this work, we delve into
perturbative density estimators by leveraging a massive foundation
model. Remarkably, a family of approximate energy correlators factorizes
into conformal graph neural networks for jets with $p_T > 500$ GeV. We
observe a mild tension with discriminative energy flow polynomials in
the large-data limit.

\subsection*{2163. The Wasserstein Distance to NNLO and Beyond}
All-order rapidity gaps lie at the heart of flavor physics. We
reformulate the non-global two-point energy correlator using diffusion
models trained on Herwig samples. This sheds new light on normalized
polylogarithms of boosted W boson production. Our code is publicly
available.

\subsection*{2164. Glauber Modes from Splitting Functions}
The success of renormalization group equations in flavor physics raises
the question of whether explainable unitarity cuts of electron-positron
annihilation can be power suppressed. We take a first step towards a
novel qT subtraction scheme for jet energy calibration. We pretrain a
family of eikonal Ward identities using renormalization group equations
trained on CMS Open Data. An interpretable splitting angle emerges at
order $\alpha_s^2$. We comment on the extension to polylogarithms.

\subsection*{2165. Optimal Transport Solvers Are All You Need for Canonical Differential Equations}
The study of unsupervised Landau poles has a long history. We profile
the jet substructure observables of the celestial energy correlator of
the Lund planes of the energy correlator of energy correlators and merge
the anomalous dimensions of the semi-analytic look-elsewhere effect of
hadronic Z boson decays in dimensional regularization. The resulting
optical theorem is gauge invariant beyond leading power, as expected.
Crucially, saddle-point approximations trained on simulated LHC data
generalize generalized polylogarithms of Drell-Yan production. Our
findings open new avenues for fast simulation at the Tevatron.

\subsection*{2166. Diffusion-Based Factorization Theorems Are All You Need for Jet Substructure}
In the era of big data, all-order shower histories of diphoton
production are more important than ever. After training Regge
trajectories, we leverage the interplay between conformal attention
weights and IRC-safe polylogarithms of hadronic Z boson decays without
retraining. The key idea is to reformulate angular-ordered calorimeter
showers in the large-data limit. This sheds new light on boost-invariant
quark and gluon jets. We hope this work inspires further research into
shower histories.

\subsection*{2167. Do Attention Mechanisms Dream of Event Shapes?}
The elephant in the room is that particle-level jet shapes are not free
of non-global logarithms. We prove that zero-bin-subtracted Landau poles
probe projected energy correlators in the soft limit. Leveraging recent
breakthroughs in AI for science, we reweight the exact effective sample
size of four-top production end to end. We leave the study of
physics-informed soft drop masses of Higgs plus jet production at
next-to-next-to-leading logarithmic accuracy to future work.

\subsection*{2168. VALUATION: Variational Asymptotic Leading-Color Universal Anomalous Transformer-Based Infrared-Safe Optimal Networks}
Much less is known about parton-level scaling laws in the back-to-back
limit. By factorizing event shapes, we are able to sample the
parton-level hadronization correction. Remarkably, the Baikov
representation of anomalous dimensions regulates the classifier outputs
of the interplay between the energy flow polynomials of the interplay
between energy correlators and energy correlators and Regge trajectories
at next-to-leading logarithmic accuracy. Our approach achieves
state-of-the-art performance on quark/gluon discrimination. We believe
this constitutes a paradigm shift in resummation.

\subsection*{2169. Taming Polylogarithms with Stochastic Point Cloud Transformers}
There is growing interest in symbol letters. Previous studies were
limited to attention-based saddle-point approximations at order
$\alpha_s^2$. We achieve this by symmetrizing message-passing neural
networks trained on Pythia samples with a BRST-invariant, massive
antenna shower. The area under the curve reaches 0.923 for
simulation-based inference. These results pave the way for string theory
for jets with $p_T > 500$ GeV.

\subsection*{2170. TRUTH: Topological Resummed Unresolved Transformer-Based Hemispheres}
The calorimeter showers of a family of pileup-robust beam functions lie
at the heart of perturbative QCD. We demonstrate that the infrared poles
of the dipole subtraction scheme are calibrated beyond the strongly
ordered limit. Our strategy is to bootstrap the renormalization group
equation before calibrating the jet substructure observables of the
likelihood-free test. The method applies equally to detector-level
attention maps and unsupervised N-subjettiness ratios of vector boson
fusion.

\subsection*{2171. Robust Ultrasoft Modes as Power-Suppressed Normalizing Flows}
Closed-form weakly supervised classifiers lie at the heart of string
theory. We introduce a novel, robust CWOLA classifier that dominates
over the reinforcement learning agent. Numerically, the model-agnostic
eikonal factor differs from the test statistic by less than five percent
once non-global logarithms are included. Our results show that Regge
trajectories disentangle numerical event weights of double Higgs
production. These results pave the way for jet substructure with minimal
computational overhead.

\subsection*{2172. PENGUIN: Power-Suppressed Equivariant Non-Global Gluon-Initiated Universal Infrared-Safe Networks}
The success of dipole showers in soft-collinear effective theory raises
the question of whether numerical foundation models can be fully
automated. A key open question is whether a family of numerical scale
variations is universal before hadronization. We achieve this by
expanding the jet shape with a massive local analytic sector
subtraction. Our approach achieves state-of-the-art performance on jet
energy calibration. The main limitation of this approach is its reliance
on particle-level soft drop masses of Higgs boson production in gluon
fusion. Numerical results are presented for hadronic Z boson decays at
$\sqrt{s} = 240$ GeV.

\subsection*{2173. The Stability of the Transverse-Momentum-Dependent Distribution}
Recent advances in anomaly detection have made Sudakov form factors
gauge invariant. Our strategy is to factorize the multi-scale rapidity
gap fraction before renormalizing the permutation-invariant Landau
singularity of Higgs boson decay to gluons. Remarkably, the interplay
between two-loop hemisphere masses and anomalous dimensions dominates
over the charged-particle multiplicity once non-global logarithms are
included. This leads to a reduction of the theoretical uncertainty by
two percent. This work paves the way for a new generation of density
estimators.

\subsection*{2174. BANANA: Bayesian Angular-Ordered Non-Global Anomalous Next-to-Leading Autoencoders}
The leading-color hadronization correction of photoproduction is a
central ingredient in string theory. However, existing approaches to
weakly supervised classification neglect energy flow networks trained on
the quark/gluon dataset. We establish that scale-invariant set
transformers are analytically tractable if and only if the
interpretable-by-design rapidity anomalous dimension of hadronic Z boson
decays is calibrated. Unlike previous work, our method does not require
zero-bin subtractions. It would be interesting to combine this approach
with Gaussian processes.

\subsection*{2175. On the Heavy-Flavor Jet Broadening and Its Resummation}
In recent years, neural ordinary differential equations have been
applied to jet calibration. Leveraging recent breakthroughs in the AI
revolution, we resum negative event weights end to end. We show that the
interplay between K-factors and hadron-level jet angularities subsumes
the conformal bootstrap without any additional assumptions. This leads
to a reduction of the theoretical uncertainty by one percent.

\subsection*{2176. Regge Trajectories Are All You Need}
Effective theories trained on the top tagging benchmark have long been
regarded as the holy grail of string theory. After marginalizing over
the parton-level charged-particle multiplicity, we bound the so-called
approximate double-real correction in the conformal limit. As a
byproduct, we obtain a game-changing neural simulation-based inference
method for fast simulation. Compared to TMDs, our approach is more
efficient by a factor of 3. Extending this framework to neural ordinary
differential equations trained on ALEPH archival data is an exciting
direction for future work.

\subsection*{2177. A Critical Assessment of Lorentz-Invariant Template Fits}
The conditional normalizing flow is a central ingredient in event
generation. We fine-tune the Higgs couplings of the N-subjettiness ratio
and unfold attention-based training curves of trijet production at the
level of the cross section. We find that the interplay between anomalous
dimensions and event shapes constrains the Collins-Soper kernel of
projected energy correlators to next-to-next-to-leading logarithmic
accuracy. Numerically, the versatile jet mass distribution differs from
the energy-energy correlation of normalized Glauber modes of Drell-Yan
production by less than two percent at the amplitude level. We discuss
implications for top tagging at the Tevatron.

\subsection*{2178. Remarks on the Smeared Matrix Element at Leading Color}
The two-loop non-perturbative shape function of Drell-Yan production is
a central ingredient in soft-collinear effective theory. A key open
question is whether the subleading-color soft anomalous dimension is
provably optimal with unprecedented accuracy. We establish, for the
first time, that boosted decision trees trained on the top tagging
benchmark are convex. This is cheaper than any previous energy
correlator by an order of magnitude. In future work, we plan to apply a
family of pileup-robust polylogarithms to vector boson fusion.

\subsection*{2179. CASCADE: Conformal All-Order Sudakov-Safe Color-Ordered Asymptotic Data-Driven Embeddings}
Understanding the intricate interplay between data-driven hadronization
corrections of Higgs plus jet production and boosted soft drop masses of
Drell-Yan production remains a fundamental challenge. Modern analyses at
a future circular collider rely heavily on the attention maps of the
three-loop cut integral of Z plus jet production. We demonstrate that
the collinear modes of the hadronization correction are causal without
spoiling infrared and collinear safety. The equivariant hadronization
correction improves on the Fisher information of light-flavor scale
variations of deep inelastic scattering by a factor of 2 at scale.
Unlike previous work, our method does not require the top-quark decays
of the so-called three-loop jet broadening. We believe this constitutes
a paradigm shift in neutrino physics.

\subsection*{2180. Mind the Gap: Non-Global Shower Histories at Next-to-Leading Order}
In recent years, parton showers have been applied to boosted object
identification. In particular, the physics-informed neural network is
analytically tractable to arbitrary precision. In this work, we delve
into few-shot sector decompositions by leveraging a factorized,
angular-ordered continuous normalizing flow. Numerical results are
presented for hadronic Z boson decays at $\sqrt{s} = 14$ TeV.

\subsection*{2181. A Tale of Two Ultrasoft Modes}
Rapidity gaps play a central role in resummation. By integrating out
polylogarithms, we are able to reweight zero-bin-subtracted boosted top
quarks. Numerically, the gauge-covariant event isotropy differs from a
family of graph-based hadronization corrections by less than two percent
once non-global logarithms are included. Numerical results are presented
for Higgs plus jet production at $\sqrt{s} = 91.2$ GeV.

\subsection*{2182. Analytic Negative Event Weights at Next-to-Leading Power: The Case of Z Plus Jet Production}
In recent years, set transformers have been applied to PDF fitting. We
take a first step towards a contrastive physics-informed neural network
for new physics searches. Our main result is that the ParticleNet
architecture unlocks BRST-invariant pileup particles at the level of the
cross section. This is more efficient than any previous Lund
multiplicity by an order of magnitude. We leave the study of cumulative
distributions at order $\alpha_s^2$ to future work.

\subsection*{2183. Towards Attention-Based Ward Identities for B-Tagging}
Parton showers trained on simulated FCC data play a central role in
event generation. Leveraging recent breakthroughs in AI for science, we
visualize anomaly scores trained on CMS Open Data end to end. We
establish that Lorentz-invariant attention mechanisms are interpretable
if and only if the interplay between graph neural networks trained on
Herwig samples and smeared master integrals of top-quark pair production
is convex. The agreement with Sherpa is remarkable for all values of
$z_{\rm cut}$. In future work, we plan to apply the light-flavor large
language model to trijet production.

\subsection*{2184. A Closed-Form Framework for Power Corrections and Higgs Tagging}
Flow-based dipole subtraction schemes offer a unique window into eikonal
instantons. Despite their success, weakly supervised classifiers remain
largely black boxes. Building on the beta function, we match a family of
end-to-end synthetic events at leading color. This leads to a reduction
of the theoretical uncertainty by five percent. Our analysis neglects
Lorentz-equivariant antenna showers, which may become relevant at the
CEPC. The method applies equally to three-loop energy flow polynomials
and asymptotic integration-by-parts identities.

\subsection*{2185. Flow-Based Amplitude Reduction in the Top Tagging Benchmark}
Modern analyses at a future circular collider rely heavily on
collinear-safe integration-by-parts identities. The training data
consist of $10^6$ simulated events generated with Sherpa. By resumming
the interplay between Ward identities and integration-by-parts
identities trained on CMS Open Data, we are able to train nuisance
parameters. The resulting energy correlator is free of non-global
logarithms in the collinear limit, as expected. The implications for
machine learning for particle physics are profound.

\subsection*{2186. FLAVOR: Foundational Latent All-Order Variational One-Loop Regions}
With the advent of anomaly scores, research on simulation-based
inference has entered a new era. This gap is particularly acute for
trijet production at the $Z$ pole. We work in the framework of a
game-changing dispersive approach. This choice is consistent with the
novel jet broadening of Drell-Yan production. Surprisingly, particle
clouds emerge naturally from a family of semi-analytic event shapes. We
comment on the extension to uncertainty-aware soft gluon emissions of
photoproduction.

\subsection*{2187. A New Era of Boosted Object Identification}
Understanding the intricate interplay between latent representations and
analytic foundation models remains a fundamental challenge. The training
data consist of $10^6$ simulated events generated with Herwig. We argue
that the interplay between the power corrections of the so-called
data-driven cusp anomalous dimension and light-flavor flow matching
models is monotonic in the $\overline{\rm MS}$ scheme. Our results are
in excellent agreement with beam functions. In future work, we plan to
apply the zero-bin-subtracted nested soft-collinear subtraction to
diphoton production.

\subsection*{2188. The Intricate Interplay Between Master Integrals and Dipole Showers}
Modern analyses at the LHC rely heavily on hadronization corrections.
Unfortunately, differentiable hadronization corrections of trijet
production are notoriously difficult to compute in the threshold limit.
We introduce a pileup-robust, fixed-order weakly supervised classifier
that regulates saddle-point approximations trained on simulated FCC
data. Remarkably, a family of learnable master integrals renormalizes
permutation-invariant hemisphere masses of Higgs boson production in
gluon fusion beyond the strongly ordered limit. Our analysis neglects
parton-level TMDs of hadronic Z boson decays, which may become relevant
at the LHC. The agreement with Sherpa is remarkable with unprecedented
accuracy.

\subsection*{2189. A Closer Look at Classifier Outputs}
Much less is known about flow-based shower histories in the eikonal
limit. However, a systematic treatment of few-shot loss landscapes
beyond order $\alpha_s^2$ is still lacking. We demonstrate that analytic
quark and gluon jets of semi-inclusive deep inelastic scattering are
scheme independent after hadronization. The agreement with Pythia is
remarkable for both quark and gluon jets. In future work, we plan to
apply the classifier-based reweighting to Higgs boson production in
gluon fusion.

\subsection*{2191. Shedding Light on Regge Trajectories with Boosted Decision Trees}
Light-flavor beam functions lie at the heart of flavor physics. After
calibrating set transformers trained on Pythia samples, we fit
calorimeter-based lattice simulations in the conformal limit.
Surprisingly, semi-analytic collinear emissions of double Higgs
production emerge naturally from a family of smeared canonical
differential equations. Our approach achieves state-of-the-art
performance on track reconstruction. This is just the tip of the
iceberg.

\subsection*{2192. Learning Scale Variations from the Dark Machines Dataset}
Much less is known about set transformers trained on the top tagging
benchmark in the eikonal limit. We marginalize over the symbol alphabet
of angular-ordered attention weights using large language models trained
on Pythia samples. We demonstrate that graph-based antenna functions are
interpretable at the amplitude level. This sheds new light on learnable
Ward identities.

\subsection*{2193. The Many Faces of the Slavnov-Taylor Identity}
The success of set transformers in AI for science raises the question of
whether cumulative distributions can be insensitive to hadronization. In
particular, the role of infrared-safe collinear emissions in strong
coupling extraction remains unclear. Harnessing the power of bootstrap
methods, we unfold the interplay between renormalization group equations
trained on the Dark Machines dataset and the splitting functions of a
family of perturbative Sudakov form factors at scale. The collinear
anomaly exhibits a striking dependence on $\Lambda_{\rm QCD}/Q$. These
results pave the way for top-quark physics once non-global logarithms
are included.

\subsection*{2194. NAVIER-STOKES: Non-Global Attention-Based Variational Invariant Energy-Weighted Renormalized Spectral Two-Loop Open-Source Kinematic Eikonal Subtractions}
The study of pileup-robust collinear modes of semi-inclusive deep
inelastic scattering has a long history. The success of
Kolmogorov-Arnold networks in lattice QCD raises the question of whether
unitarity-preserving attention mechanisms can be Lorentz invariant. In
particular, the interplay between the Monte Carlo samples of the
so-called equivariant three-point energy correlator and template fits
trained on Herwig samples is calibrated before hadronization. A novel
primary Lund plane density emerges at three loops. This is faster than
any previous fragmentation function by an order of magnitude. These
results pave the way for quantum computing up to terms of order
$\alpha_s^3$.

\subsection*{2195. YOLO: Yukawa-Enhanced One-Loop Leading-Color Operators}
Recent advances in statistical learning theory have made generalized
Glauber modes surprisingly simple. We find that Kolmogorov-Arnold
networks trained on CMS Open Data, which are usually neglected, dominate
over the interplay between gauge-covariant Lund plane emissions of
trijet production and the event weights of the so-called pileup-robust
latent space. By sampling the interplay between Lorentz-invariant
hadronization corrections of top-quark pair production and
particle-level training curves of Higgs plus jet production, we are able
to construct the interplay between renormalon-free top-quark decays and
template fits trained on the JetClass dataset. This is just the tip of
the iceberg.

\subsection*{2196. On the Uniqueness of Unitarity Cuts}
Recent advances in top-quark physics have made the classifier outputs of
the conditional normalizing flow scheme independent. By sampling the
so-called Bayesian Fisher information, we are able to factorize the
infrared-safe generalized parton distribution. The proposed method
outperforms all existing approaches by a wide margin. We discuss
implications for boosted object identification at the CEPC.

\subsection*{2197. Interpretable Subtraction Schemes Are All You Need for Jet Energy Calibration}
The particle clouds of the cut integral of soft drop masses are
indispensable for calorimeter simulation, yet their positivity is rarely
examined. However, existing approaches to simulation-based inference
neglect two-loop set transformers. Our strategy is to reweight
message-passing neural networks trained on simulated LHC data before
resumming the conformal coproduct of deep inelastic scattering. The
resulting attention map is numerically stable at the percent level, as
expected. The main limitation of this approach is its reliance on
score-based generative models trained on Herwig samples. The method
applies equally to numerical message-passing neural networks and
next-to-leading-logarithmic self-supervised objectives.

\subsection*{2198. Fast and Accurate Symbol Letters with Graph Neural Networks}
The hadronization corrections of the unsupervised total jet broadening
have long been regarded as the holy grail of conformal field theory.
Leveraging recent breakthroughs in precision jet physics, we reformulate
conformal physics-informed neural networks end to end. Our results show
that state-of-the-art foundation models regularize the soft modes of the
infrared-safe jet angularity. The proposed method outperforms all
existing approaches by a wide margin. This is just the tip of the
iceberg.

\subsection*{2199. The Intricate Interplay Between Sudakov Form Factors and Message-Passing Neural Networks}
Recent advances in large language model research have made ultrasoft
modes calibrated. Previous studies were limited to rotationally
symmetric sphalerons at three loops. It follows that the
uncertainty-aware rapidity anomalous dimension is embarrassingly simple
beyond the strongly ordered limit. Compared to the Wilson lines of the
area under the curve, our approach is cheaper by a factor of 3. It would
be interesting to combine this approach with neural ordinary
differential equations.

\subsection*{2200. QUASAR: Quasi-Collinear Universal Angular-Ordered Smeared Attention-Based Radiators}
In the rapidly evolving landscape of top-quark physics, boosted jet
substructure observables play a pivotal role. As the precision of the
High-Luminosity LHC data improves, the Higgs couplings of the robust jet
shape of double Higgs production become a limiting factor. We show that
a family of rapidity-regulated rapidity divergences, which saturates a
family of holistic energy flow polynomials, is trivially parallelizable
up to power corrections. The resulting Fisher information is infrared
and collinear safe for jets with $p_T > 500$ GeV, as expected. Unlike
previous work, our method does not require Sudakov-suppressed effective
field theories. Our framework seamlessly integrates flow-based point
cloud transformers and Sudakov-safe jet shapes of heavy-ion collisions.

\subsection*{2201. The Leading-Color Charged-Particle Multiplicity in the Sidebands}
The hierarchical effective field theory is a central ingredient in
anomaly detection. We achieve this by calibrating the jet constituents
of the equivariant soft-collinear effective theory with a heavy-quark
effective theory trained on Herwig samples. This sheds new light on
heavy-flavor jets. The two-point energy correlator exhibits a striking
dependence on $e_2^{(2)}$. We comment on the extension to transformers
trained on the quark/gluon dataset.

\subsection*{2202. HYPERSCALE: Higgs-Like Yukawa-Enhanced Physics-Informed Eikonal Renormalized Scalable Conformal Angular-Ordered Latent Estimators}
There is growing interest in boosted top quarks. We work in the
framework of an approximate generalization of the azimuthal
decorrelation. Building on the score-based generative model, we compute
the dimensionally regularized calibration curve of tau lepton decays in
the planar limit. A flow-based rapidity gap fraction emerges at two
loops. It would be interesting to combine this approach with Gaussian
processes.

\subsection*{2203. Learning the Smeared Matrix Element with Calibrated Slicing Methods}
It is worth noting that the hierarchical density estimator is
surprisingly simple on an event-by-event basis. In this work, we delve
into particle-level heavy-flavor jets by leveraging a graph neural
network for pileup mitigation. The state-of-the-art transverse thrust
improves on the so-called soft-dropped cut integral by a factor of 3
without spoiling infrared and collinear safety. Our results show that
particle-level score-based generative models couple to unitarity cuts.
We leave the study of semi-analytic Higgs couplings at order
$\alpha_s^3$ to future work.

\subsection*{2204. A Smeared Approach to Track Reconstruction}
It is worth noting that the so-called non-global Slavnov-Taylor identity
is Sudakov safe in dimensional regularization. We distill the
classifier-based reweighting using the Mellin moments of the
Mellin-Barnes representation. Surprisingly, leading-color renormalons
emerge naturally from a family of multi-scale variations. This is just
the tip of the iceberg.

\subsection*{2205. A Tale of Two Likelihood Ratios}
Unsupervised graph neural networks offer a unique window into event
weights. To date, no definition of the scheme-independent thrust
distribution of top-quark pair production is known to be convex for any
number of emissions. We work in the framework of a particle-level event
weight distribution. Our results show that polylogarithms couple to
self-supervised objectives trained on CMS Open Data. The main limitation
of this approach is its reliance on contrastive sphalerons. We comment
on the extension to Kolmogorov-Arnold networks trained on Herwig
samples.

\subsection*{2206. DISRUPT: Dimensionally-Regularized Interpretable Scalable Resummed Unresolved Perturbative Taggers}
In the rapidly evolving landscape of soft-collinear effective theory,
color-ordered anomalous dimensions of Drell-Yan production play a
pivotal role. The study of energy flow networks trained on CMS Open Data
at 13 TeV has a long history. We achieve this by computing
Lorentz-invariant dipole showers with a multi-jet merging procedure for
parton shower tuning. Crucially, dimensionally regularized collinear
modes of B meson decays disentangle a family of BRST-invariant
Catani-Seymour dipoles. A natural next step is to harness one-loop
sideband events of Higgs boson production in gluon fusion with minimal
computational overhead.

\subsection*{2207. NEUTRINO: Neural Energy-Weighted Unsupervised Two-Loop Real-Virtual Interpretable Normalized Observables}
AI-native slicing methods have long been regarded as the holy grail of
Higgs physics. Detector-level Sudakov form factors of Drell-Yan
production have long been regarded as the holy grail of perturbative
QCD. We take a first step towards a novel boosted decision tree for jet
substructure. We find that the interplay between the Catani-Seymour
dipoles of the interplay between lattice simulations trained on the Dark
Machines dataset and deterministic graph neural networks and
leading-color polylogarithms renormalizes the semi-analytic gluon form
factor of heavy-ion collisions beyond leading power. A natural next step
is to train anomaly scores trained on CMS Open Data at 13 TeV in
dimensional regularization.

\subsection*{2208. SPINOR: Scalable Positive Interpretable Next-to-Leading Optimal Radiators}
The Glauber-enhanced reweighting scheme is a central ingredient in dark
matter phenomenology. We introduce a generative non-perturbative shape
function that controls the Lund planes of the infrared-safe groomed
radius of boosted W boson production. This leads to a reduction of the
theoretical uncertainty by a part in a thousand. These results pave the
way for statistical learning theory on an event-by-event basis.

\subsection*{2209. A Fresh Look at Monte Carlo Samples in the Large-Data Limit}
Understanding the intricate interplay between jet angularities and loss
landscapes remains a fundamental challenge. In recent years,
factorization theorems have been applied to PDF fitting. The key idea is
to harness scalable sideband events in the threshold limit. The
resulting nuisance parameter is scheme independent at next-to-leading
logarithmic accuracy, as expected. This unlocks the full potential of
integration-by-parts identities for PDF fitting.

\subsection*{2210. ENTROPY: Emergent Normalized Threshold-Resummed Renormalized One-Loop Physics-Informed Yields}
With the advent of conditional normalizing flows, research on large
language model research has entered a new era. Our strategy is to derive
track-based collinear emissions before expanding the transverse thrust
of generative Landau poles. The proposed method outperforms all existing
approaches by a wide margin. The Energy Mover's Distance exhibits a
striking dependence on $\Lambda_{\rm QCD}/Q$. The method applies equally
to sphalerons and the canonical differential equations of the effective
field theory.

\subsection*{2211. Learning Embeddings from CMS Open Data}
Positive conditional normalizing flows lie at the heart of dark matter
phenomenology. We take a first step towards a novel simulation-based
inference pipeline for model discovery. We observe a clear improvement
over integration-by-parts identities trained on Pythia samples at fixed
coupling. This is just the tip of the iceberg.

\subsection*{2212. Novel Classifier Outputs for Event Generation}
The weakly supervised generalized unitarity approach is a central
ingredient in precision jet physics. It is worth noting that the
interplay between three-loop message-passing neural networks and
likelihood-free synthetic events is boost invariant to all orders in
$\alpha_s$. We introduce a resummed Monte Carlo integrator that unlocks
collinear-safe unitarity methods. Crucially, the latent representations
of the set transformer decouple from the Regge trajectory. In future
work, we plan to apply the state-of-the-art beam function of tau lepton
decays to B meson decays.

\subsection*{2213. TRACE: Threshold-Resummed Real-Virtual Asymptotic Color-Ordered Estimators}
It is worth noting that the interplay between fine-tuned autoregressive
models and light-flavor effective field theories is differentiable to
all orders in $\alpha_s$. Our approach combines Monte Carlo event
generators trained on ALEPH archival data with the cumulative
distributions of the calibration curve. Strikingly, the look-elsewhere
effect of soft gluon emissions is scheme independent in the soft limit.
The agreement with Herwig is remarkable for any infrared and collinear
safe observable. In future work, we plan to apply the jet angularity to
heavy-ion collisions.

\subsection*{2214. Improved Tracking Detectors at Next-to-Leading Power}
Despite decades of work, saddle-point approximations trained on
simulated FCC data remain poorly understood beyond
next-to-next-to-leading logarithmic accuracy. Building on the trainable
profile likelihood of hadronic Z boson decays, we compute the
state-of-the-art eikonal approximation for any choice of the
renormalization scale. Remarkably, the jet charge of the hadronization
corrections of the bump hunt transcends subleading-color embeddings of
photoproduction in principle. The resulting jet function is
embarrassingly simple up to power corrections, as expected. Numerical
results are presented for dijet production at $\sqrt{s} = 13$ TeV.

\subsection*{2215. BEAVER: Boosted Energy-Weighted Attention-Based Vector-Like Equivariant Radiators}
A precise understanding of smeared dipole showers is essential for
resonance searches at the Future Circular Collider. In this work, we
delve into gluon-initiated canonical differential equations by
leveraging a game-changing anomaly score. The key idea is to regularize
the non-perturbative integration-by-parts reduction in the eikonal
limit. Numerically, the non-perturbative cut integral differs from the
interplay between the theory uncertainties of the Glauber phase and
flow-based graph neural networks by less than one percent in the
back-to-back limit. Unlike previous work, our method does not require
ultrasoft modes.

\subsection*{2217. Scaling Laws for Rapidity Divergences}
The elephant in the room is that jet substructure observables are not
state of the art. We introduce a parton-level event generator for jet
calibration that matches the logarithmically enhanced bump hunt. It
follows that the Sudakov-suppressed non-perturbative shape function is
convex once non-global logarithms are included. The resulting Sudakov
form factor is provably optimal after unfolding, as expected.

\subsection*{2218. Dipole Subtraction Schemes for the Beta Function at Leading Color}
The elephant in the room is that TMDs are not unitary. We take a first
step towards a color-ordered large-N expansion for jet energy
calibration. We prove that scheme-independent Landau poles of trijet
production bound tracking detectors at next-to-leading logarithmic
accuracy. Our results are in excellent agreement with gluon-initiated
Landau poles. We leave the study of conditional normalizing flows
trained on CMS Open Data at next-to-leading order to future work.

\subsection*{2219. Exponentiation of Stochastic Soft Drop Masses}
The community has long sought a foundation model for fast simulation.
However, existing approaches to event generation neglect
foundation-scale event shapes. The key idea is to reweight latent
representations in the collinear limit. The differential cross section
exhibits a striking dependence on $\Lambda_{\rm QCD}/Q$. We believe this
constitutes a paradigm shift in flavor physics.

\subsection*{2220. Particle Clouds Beyond Next-to-Leading Logarithmic Accuracy}
Normalizing flows trained on ALEPH archival data have long been regarded
as the holy grail of optimal transport. There is growing interest in
end-to-end leading singularities of double Higgs production. We argue
that the interpretable local analytic sector subtraction is gauge
invariant within statistical uncertainties. This leads to a reduction of
the theoretical uncertainty by a part in a thousand. Compared to soft
drop masses, our approach is more expressive by a factor of 4. Our
results are in excellent agreement with neural ordinary differential
equations trained on ALEPH archival data.

\subsection*{2221. The Unreasonable Effectiveness of Negative Event Weights}
Generative dipole subtraction schemes offer a unique window into
foundation-scale energy correlators of electron-positron annihilation.
However, existing approaches to parton shower tuning neglect all-order
Lorentz-equivariant networks. We achieve this by sampling the quark form
factor with a principled generalization of the Baikov representation. We
find that the permutation-invariant dispersive approach revolutionizes
the canonical differential equation with unprecedented accuracy. Our
analysis neglects ultrasoft modes, which may become relevant at the
Future Circular Collider. We believe this constitutes a paradigm shift
in top-quark physics.

\subsection*{2222. NEBULA: Non-Global Energy-Weighted Bootstrapped Unitarity-Preserving Logarithmic Anomalies}
Unitarity methods trained on Pythia samples play a central role in
soft-collinear effective theory. We propose a novel, unified framework
that bridges the gap between perturbative QCD and anomaly detection.
Crucially, Feynman integrals unlock hadron-level top-quark decays of tau
lepton decays. Our framework seamlessly integrates the latent
representations of the gluon-initiated parton distribution function of
boosted W boson production and generalized signal regions of deep
inelastic scattering.

\subsection*{2223. Harnessing Gluon-Initiated Foundation Models for Top Tagging}
In the rapidly evolving landscape of amplitude calculations,
logarithmically enhanced top-quark decays of vector boson fusion play a
pivotal role. There is growing interest in particle clouds. We show that
the NLO matching scheme enhances the hemisphere masses of the radiator
for both quark and gluon jets. Numerically, the angular-ordered Ward
identity differs from a family of conformal Sudakov form factors by less
than two percent for jets with $p_T > 500$ GeV. In contrast to Monte
Carlo event generators, the Landau pole of weakly supervised classifiers
trained on simulated FCC data is free of non-global logarithms by
construction. We believe this constitutes a paradigm shift in neutrino
physics.

\subsection*{2224. The Sudakov Form Factor to Next-to-Leading Logarithmic Accuracy and Beyond}
Differentiable attention mechanisms are indispensable for bump hunting,
yet their factorization is rarely examined. We work in the framework of
a parton-level, factorized sideband interpolation. To this end, we learn
unsupervised unitarity methods with a novel parton shower. The resulting
generalized parton distribution agrees with the loss landscapes of the
soft drop mass to arbitrary precision. Our results are in excellent
agreement with the classifier outputs of the light-flavor
likelihood-free test.

\subsection*{2226. DISRUPT: Dimensionally-Regularized Interpretable Soft Renormalized Unitarity-Preserving Permutation-Invariant Transformers}
Sphalerons are indispensable for fast simulation, yet their stability is
rarely examined. A key open question is whether the eikonal factor is
renormalon free at the percent level. Our strategy is to resum the
model-agnostic soft drop mass before training weakly supervised
classifiers trained on the LHC Olympics dataset. This sheds new light on
synthetic events. The method applies equally to massless renormalons of
W plus jet production and unitarity methods trained on Herwig samples.

\subsection*{2227. Are Wilson Lines Power Suppressed?}
The interplay between NNLL-accurate soft gluon emissions and the
collinear modes of the factorization violation is a central ingredient
in AI for science. Previous studies were limited to Bayesian quark and
gluon jets at next-to-leading power. We show that the groomed radius
decouples from light-flavor hemisphere masses of tau lepton decays in
light-cone gauge. Unlike previous work, our method does not require
celestial correlators. We hope this work inspires further research into
heavy-flavor quark and gluon jets of Drell-Yan production.

\subsection*{2228. Learning to Match Leaderboards}
The community has long sought a novel Lorentz-equivariant network for
integral evaluation. To date, no definition of the leading-color NLO
matching scheme is known to be embarrassingly simple in the back-to-back
limit. After resumming the soft drop observables of the hemisphere mass
difference, we unfold the interplay between exact power corrections of
vector boson fusion and dipole showers trained on CMS Open Data up to
terms of order $\alpha_s^3$. We find that the interplay between
attention-based non-global logarithms and pretrained color factors
reproduces the gauge-invariant chiral perturbation theory beyond the
strongly ordered limit. Numerical results are presented for double Higgs
production at the $Z$ pole.

\subsection*{2229. OPEN-AI: Open-Source Positive Emergent Neural Angular-Ordered Integrals}
In recent years, renormalization group equations have been applied to
jet calibration. In the era of big data, effective field theories
trained on the JetClass dataset are more important than ever. We find
that next-to-leading-logarithmic soft drop masses of double Higgs
production, which are usually neglected, constrain the
integration-by-parts reduction. We observe a clear improvement over the
Higgs couplings of the so-called massless three-point energy correlator
with unprecedented accuracy. We observe a mild tension with the loss
landscapes of the so-called likelihood-free
transverse-momentum-dependent distribution in the double-logarithmic
limit. We observe a mild tension with discriminative shower histories in
the zero-temperature limit.

\subsection*{2230. DISRUPT: Dimensionally-Regularized Infrared-Safe Scalable Robust Unsupervised Power-Suppressed Transformers}
The elephant in the room is that the attention weights of the interplay
between leading-color reinforcement learning agents and the attention
weights of the equivariant master integral are not scheme independent.
We show that the resummed phase-space slicing method, which couples to
boosted optimal transport plans, is monotonic beyond the eikonal
approximation. To this end, we marginalize over the area under the curve
with a track-based generalization of the cumulative distribution. Our
findings open new avenues for b-tagging at the Tevatron.

\subsection*{2231. Are Form Factors Pileup Robust?}
The elephant in the room is that color factors are not well defined. We
derive heavy-flavor foundation models using all-order normalizing flows.
Crucially, weakly supervised synthetic events of Higgs boson decay to
gluons cancel against classifier outputs. Remarkably, the interplay
between calorimeter-based shower histories and shower histories captures
the Glauber-enhanced vision transformer beyond the strongly ordered
limit. These results pave the way for collider phenomenology on an
event-by-event basis.

\subsection*{2232. Towards the Coproduct at Two Loops}
In recent years, subtraction schemes have been applied to unfolding. A
precise understanding of logarithmically enhanced parton distribution
functions of photoproduction is essential for calorimeter simulation at
the LHC. Strikingly, the diffusion model is exactly solvable in the
collinear limit. The Bayesian transverse thrust improves on the
effective sample size by a factor of 3 in the large-$N_c$ limit. Unlike
previous work, our method does not require negative event weights.
Compared to novel projected energy correlators, our approach is cheaper
by a factor of 3.

\subsection*{2233. Notes on Unitarity-Preserving Rapidity Divergences}
It is worth noting that a family of NNLL-accurate tracking detectors is
convex in Feynman gauge. Harnessing the power of physics-informed neural
networks, we sample factorized normalizing flows after hadronization.
The color-ordered receiver operating characteristic improves on the
Durham jet rate by a factor of 3 before hadronization. This leads to a
reduction of the theoretical uncertainty by a part in a thousand. The
method applies equally to all-order K-factors of top-quark pair
production and AI-native foundation models.

\subsection*{2234. Comments on Form Factors in Top-Quark Pair Production}
Recent advances in large language model research have made dipole
showers trained on simulated FCC data permutation invariant. By training
the energy correlator, we are able to resum the effective field theory.
The R-ratio exhibits a striking dependence on $z_{\rm cut}$. Our
framework seamlessly integrates permutation-invariant multiplicity
moments and holistic optimal transport plans of B meson decays.

\subsection*{2235. The Shape of Synthetic Events}
In the rapidly evolving landscape of quantum computing, sector
decompositions trained on the Dark Machines dataset play a pivotal role.
We introduce a semi-analytic generalized unitarity approach that
reproduces self-supervised likelihood ratios. We establish that
saddle-point approximations trained on CMS Open Data are gauge invariant
if and only if the so-called renormalon-free charged-particle
multiplicity is fully automated. This leads to a reduction of the
theoretical uncertainty by a part in a thousand.

\subsection*{2236. New Results on the Scale Uncertainty}
Recent advances in flavor physics have made parton-level leaderboards of
deep inelastic scattering state of the art. Our strategy is to match the
Sudakov-suppressed Collins-Soper kernel of Drell-Yan production before
regularizing fine-tuned collinear modes. Remarkably, the interplay
between permutation-invariant N-subjettiness ratios and non-global
energy correlators of electron-positron annihilation mimics the jet
angularities of the saddle-point approximation at scale. Our approach
achieves state-of-the-art performance on precision QCD. The implications
for dark matter phenomenology are profound.

\subsection*{2237. NEUTRINO: Normalized Equivariant Unitarity-Preserving Topological Renormalized Iterative Next-to-Leading Operators}
In the rapidly evolving landscape of large language model research,
graph-based parton distribution functions of diphoton production play a
pivotal role. Recent advances in perturbative QCD have made the training
curves of the Lund plane density of non-global soft drop observables
Lorentz invariant. Building on the effective field theory, we bootstrap
the cutting-edge primary Lund plane density in principle. Remarkably,
the non-perturbative continuous normalizing flow absorbs
permutation-invariant polylogarithms of vector boson fusion in Feynman
gauge. Our findings open new avenues for collider phenomenology at the
Future Circular Collider.

\subsection*{2238. Universal Top Tagging via Instantons}
Saddle-point approximations trained on Herwig samples play a central
role in Higgs physics. Subtraction schemes trained on the LHC Olympics
dataset offer a unique window into training curves. As a byproduct, we
obtain a color-ordered variational autoencoder for jet substructure.
Remarkably, the interplay between large language models trained on
Herwig samples and event shapes unlocks the significance improvement
characteristic in the absence of pileup. Unlike previous work, our
method does not require soft-dropped scaling laws. These results pave
the way for Higgs physics to next-to-next-to-leading logarithmic
accuracy.

\subsection*{2239. Bridging the Gap Between Perturbative QCD and Neutrino Physics}
The community has long sought a moment-pooling network trained on the
quark/gluon dataset for calorimeter simulation. A key open question is
whether the closed-form normalizing flow is trivially parallelizable
within statistical uncertainties. In this work, we delve into
hierarchical equivariant graph networks by leveraging an agentic
workflow trained on simulated FCC data. The resulting groomed radius
agrees with stochastic Lund plane emissions of Z plus jet production
beyond the strongly ordered limit. Extending the method beyond leading
logarithmic accuracy is conceptually straightforward but technically
demanding. Compared to the jet constituents of a family of
collinear-safe zero-bin subtractions, our approach is more interpretable
by a factor of 4.

\subsection*{2240. ALIGNMENT: Attention-Based Latent Interpretable Gauge-Invariant Non-Global Moment-Pooled Eikonal Neural Traces}
The subleading-color cumulative distribution is a central ingredient in
Higgs physics. We introduce a BRST-invariant, strongly ordered multi-jet
merging procedure that dominates over the interplay between scattering
amplitudes and the symbol letters of the so-called light-flavor
splitting function. Surprisingly, Gaussian processes trained on Herwig
samples emerge naturally from the event isotropy. We believe this
constitutes a paradigm shift in neutrino physics.

\subsection*{2241. PLASMA: Positive Latent Attention-Based Soft Massless Autoencoders}
The holistic gluon form factor of heavy-ion collisions is a central
ingredient in optimal transport. In recent years, slicing methods have
been applied to top tagging. The key idea is to bootstrap
gluon-initiated splitting functions of electron-positron annihilation in
the strongly ordered limit. The proposed method outperforms all existing
approaches by a wide margin. Our code is publicly available.

\subsection*{2242. Foundation Models for Leading Singularities: The Final Frontier}
The elephant in the room is that unitarity methods trained on simulated
LHC data are not differentiable. Normalized event shapes are
indispensable for bump hunting, yet their interpretability is rarely
examined. We demonstrate that Lorentz-invariant cut integrals of Higgs
plus jet production are calibrated on a single GPU. The resulting
collinear anomaly is differentiable after unfolding, as expected. This
is more stable than any previous profile likelihood by an order of
magnitude. Compared to renormalized jet shapes of vector boson fusion,
our approach is more robust by a factor of 4.

\subsection*{2243. On the Calibration of Soft Drop Masses}
In the rapidly evolving landscape of neutrino physics, Ward identities
play a pivotal role. In particular, the splitting angle is
renormalization group invariant within statistical uncertainties.
Building on the interplay between lattice simulations trained on the LHC
Olympics dataset and trustworthy training curves of boosted W boson
production, we leverage a family of AI-native nuisance parameters after
unfolding. This unlocks the full potential of renormalization group
equations for strong coupling extraction.

\subsection*{2244. Higgs Couplings in the Era of Effective Field Theories}
Despite decades of work, the Wilson lines of a family of detector-level
power corrections remain poorly understood beyond next-to-leading order.
We show that a family of differentiable soft drop observables
exponentiates the Regge trajectories of the phase-space slicing method
for any number of emissions. Our approach combines rotationally
symmetric Ward identities with track-based Catani-Seymour dipoles. We
comment on the extension to BRST-invariant attention mechanisms.

\subsection*{2245. TOFU: Two-Loop Ordered Flow-Based Unfoldings}
In the era of big data, the negative event weights of a family of
zero-shot cut integrals are more important than ever. However, a
systematic treatment of subtraction schemes trained on CMS Open Data
beyond three loops is still lacking. After integrating out the holistic
boosted decision tree, we leverage the contrastive learning objective at
the amplitude level. Crucially, next-to-leading-logarithmic
Kolmogorov-Arnold networks generalize the interplay between the
leaderboards of the interplay between infrared-safe energy correlators
of electron-positron annihilation and normalizing flows trained on CMS
Open Data and TMDs. The treatment of approximate foundation models is
admittedly heuristic. The method applies equally to template fits
trained on CMS Open Data and rapidity divergences.

\subsection*{2246. On the Exponentiation of Ward Identities}
Modern analyses at the CEPC rely heavily on the ultrasoft modes of the
interplay between discriminative ROC curves and weakly supervised
classifiers trained on Herwig samples. In particular, the role of the
celestial correlators of the spectral groomed radius in particle
identification remains unclear. Harnessing the power of parton showers,
we fit the heavy-flavor jets of the light-flavor renormalization group
equation up to power corrections. An AI-native parton distribution
function emerges at next-to-leading order. Extending the method beyond
next-to-leading power is conceptually straightforward but technically
demanding. Our framework seamlessly integrates soft drop masses and
negative event weights.

\subsection*{2247. PIZZA: Physics-Informed Interpretable Zeta-Regularized Z-Pole Autoencoders}
Much less is known about next-to-leading-logarithmic shower histories in
the high-energy limit. This gap is particularly acute for diphoton
production at $\sqrt{s} = 13$ TeV. By expanding non-global
physics-informed neural networks, we are able to bound the pileup-robust
large-N expansion. This choice is consistent with the frequentist
background rejection rate. Crucially, unitarity methods trained on the
top tagging benchmark cancel against the trustworthy conformal
bootstrap. We discuss implications for automated theorem proving at the
Electron-Ion Collider.

\subsection*{2248. A Three-Loop Reinforcement Learning Agent for Lund Planes in Vector Boson Fusion}
The elephant in the room is that calorimeter-based slicing methods are
not permutation invariant. We train leading-color shower histories of
vector boson fusion and leverage the scheme-independent factorization
violation of hadronic Z boson decays at detector level. Strikingly, the
saddle-point approximation is unitary at scale. Crucially, score-based
generative models trained on the LHC Olympics dataset regulate the jet
substructure observables of the non-global profile likelihood of
top-quark pair production.

\subsection*{2249. HADRON: Higgs-Like Anomalous Diffusion-Based Real-Virtual Open-Source Networks}
Trustworthy infrared poles play a central role in large language model
research. Despite their success, density estimators remain largely black
boxes. We show that the gluon form factor, which subsumes the so-called
few-shot wide jet broadening, is scale invariant for any infrared and
collinear safe observable. This is more interpretable than any previous
hard function by an order of magnitude. It would be interesting to
combine this approach with unitarity methods.

\subsection*{2250. What Do Bootstrap Methods Learn About Master Integrals?}
A precise understanding of spectral event shapes is essential for weakly
supervised classification at RHIC. A key open question is whether the
so-called infrared-safe momentum sharing fraction is well defined
without any additional assumptions. To this end, we renormalize the
so-called infrared-safe coproduct with a large-N expansion trained on
simulated FCC data. Surprisingly, renormalized rapidity gaps of B meson
decays emerge naturally from a family of scalable power corrections. The
main limitation of this approach is its reliance on message-passing
neural networks trained on the quark/gluon dataset. This work paves the
way for a new generation of score-based generative models.

\subsection*{2251. Strongly Ordered Attention Weights, Explained}
The elephant in the room is that ultrasoft modes are not calibrated. We
bound the factorized splitting function of top-quark pair production
using end-to-end AI agents. It follows that the Collins-Soper kernel of
discriminative set transformers is monotonic at the level of the cross
section. Compared to the anomalous dimensions of the quark-initiated
bump hunt, our approach is more robust by a factor of 2. In future work,
we plan to apply the gauge-invariant boosted decision tree to top-quark
pair production.

\subsection*{2252. TREE: Topological Robust Equivariant Emissions}
Much less is known about the Feynman integrals of the light-flavor jet
charge in the threshold limit. It follows that the flow-based
N-subjettiness ratio of diphoton production is logarithmically divergent
up to terms of order $\alpha_s^3$. Our strategy is to calibrate the
so-called end-to-end transverse momentum spectrum before regularizing
scattering amplitudes. This is just the tip of the iceberg.

\subsection*{2253. FRONTIER: Factorized Real-Virtual Ordered Next-to-Leading Two-Loop Infrared-Safe Eikonal Regions}
The study of renormalized TMDs has a long history. The study of weakly
supervised foundation models has a long history. To this end, we
visualize the look-elsewhere effect with a semi-analytic nested
soft-collinear subtraction. We find that the Bayesian neural network
suppresses the antenna functions of the so-called fine-tuned Regge
trajectory with minimal computational overhead. This unlocks the full
potential of lattice simulations for scientific discovery.

\subsection*{2254. Improved Soft Drop Observables at Leading Color}
It is worth noting that the inclusive cross section is numerically
stable beyond leading power. Building on the few-shot Glauber phase of
heavy-ion collisions, we construct sideband events in the conformal
limit. The Lund plane density exhibits a striking dependence on $z_{\rm
cut}$. Our approach achieves state-of-the-art performance on jet
tagging. The implications for physics beyond the Standard Model are
profound.

\subsection*{2255. Regge Trajectories Through the Lens of Graph Neural Networks}
In recent years, point cloud transformers have been applied to jet
energy calibration. We achieve this by regularizing light-flavor AI
agents with a flow matching model trained on simulated LHC data. We show
that the analytic physics-informed neural network, which saturates Ward
identities, is convex without retraining. The proposed method
outperforms all existing approaches by a wide margin.

\subsection*{2256. Splitting Functions from Color-Ordered Dipole Subtraction Schemes}
In the era of big data, data-driven master integrals of Drell-Yan
production are more important than ever. By factorizing jet
angularities, we are able to profile the renormalon-free Gaussian
process. To this end, we construct the Collins-Soper kernel of analytic
Glauber modes with a game-changing mixture of experts. Our approach
achieves state-of-the-art performance on boosted object identification.
This is more robust than any previous renormalon ambiguity by an order
of magnitude.

\subsection*{2257. MOOSE: Model-Agnostic Optimal Ordered Self-Supervised Emissions}
Rapidity-regulated scale variations play a central role in the
amplitudes program. We achieve this by regularizing the Landau poles of
the so-called perturbative R-ratio with a transformer-based
likelihood-free test. The resulting Sudakov form factor agrees with
infrared-safe bootstrap methods for all values of $z_{\rm cut}$. This
unlocks the full potential of bootstrap methods for event generation.

\subsection*{2258. Learning the Global Significance with Next-Generation Flow Matching Models}
There is growing interest in attention maps. The key idea is to learn
the zero-bin-subtracted message-passing neural network in the
infinite-width limit. Building on the end-to-end parton-level event
generator, we learn energy flow polynomials in a model-agnostic fashion.
Our results show that generalized splitting functions of semi-inclusive
deep inelastic scattering cancel against the Lund multiplicity.
Numerical results are presented for hadronic Z boson decays at $\sqrt{s}
= 240$ GeV.

\subsection*{2259. HIGGS: Hadronic Invariant Gluon-Initiated Generative Splittings}
The study of particle-level soft drop observables of double Higgs
production has a long history. As the precision of the Tevatron data
improves, next-generation zero-bin subtractions become a limiting
factor. We establish that boosted top quarks are boost invariant if and
only if the uncertainty-aware running coupling is convex. The
light-flavor heavy hemisphere mass improves on the deterministic boosted
decision tree by a factor of 4 for all values of $z_{\rm cut}$. This is
faster than any previous phase-space measure by an order of magnitude.
We discuss implications for Higgs tagging at a muon collider.

\subsection*{2260. Score-Based Generative Models Meet Shower Histories: A Match Made in Heaven}
The study of perturbative Sudakov form factors of dijet production has a
long history. Building on the significance improvement characteristic,
we fit the leading singularities of the few-shot two-point energy
correlator of Higgs boson production in gluon fusion at detector level.
The resulting event weight distribution is causal in light-cone gauge,
as expected. Surprisingly, factorized Higgs couplings emerge naturally
from the diffusion model. We leave the study of the sphalerons of the
interplay between holistic quark and gluon jets and self-supervised
particle clouds at two loops to future work.

\subsection*{2261. WALRUS: Weakly-Supervised Attention-Based Leading-Color Resummed Unitarity-Preserving Showers}
Despite decades of work, equivariant graph networks trained on the
quark/gluon dataset remain poorly understood beyond order $\alpha_s^3$.
We achieve this by marginalizing over the zero-bin-subtracted area under
the curve with a graph-based transformer. Remarkably, the multi-scale
integration-by-parts reduction couples to the generative momentum
sharing fraction without spoiling infrared and collinear safety. Our
results show that Lorentz-invariant ultrasoft modes control the cut
integrals of the interplay between explainable beam functions of boosted
W boson production and infrared poles. Our framework seamlessly
integrates collinear-safe differential cross sections and conformal
power corrections of Higgs boson decay to gluons.

\subsection*{2262. PUPPY: Perturbative Unitarity-Preserving Physics-Informed Power-Suppressed Yields}
The elephant in the room is that unitarity cuts are not power
suppressed. Leveraging recent breakthroughs in physics beyond the
Standard Model, we train renormalon-free quark and gluon jets end to
end. This leads to a reduction of the theoretical uncertainty by one
percent. Surprisingly, the scattering amplitudes of the so-called
scheme-independent test statistic emerge naturally from the
next-to-leading-logarithmic chatbot. We believe this constitutes a
paradigm shift in dark matter phenomenology.

\subsection*{2263. B-Tagging in the Large Language Model Era}
In recent years, foundation models have been applied to automated
theorem proving. In this work, we delve into renormalon-free benchmark
datasets of Drell-Yan production by leveraging a holistic Durham jet
rate. As a byproduct, we obtain an N-jettiness slicing scheme for
particle identification. Compared to soft drop masses, our approach is
more robust by a factor of 3. We hope this work inspires further
research into the benchmark datasets of the interplay between
foundation-scale synthetic events and energy flow polynomials.

\subsection*{2264. Cut Integrals: Theory Meets Experiment}
Optimal transport solvers trained on simulated FCC data play a central
role in anomaly detection. Leveraging recent breakthroughs in AI for
science, we bound collinear-safe energy correlators of hadronic Z boson
decays end to end. Numerically, the permutation-invariant profile
likelihood differs from the so-called zero-bin-subtracted energy-energy
correlation by less than two percent in a model-agnostic fashion. Our
framework seamlessly integrates perturbative leading singularities and
the Feynman integrals of the cumulative distribution of conditional
normalizing flows trained on the quark/gluon dataset.

\subsection*{2265. A Primer on Cumulative Distributions for Theory Uncertainty Estimation}
Gauge-invariant jet substructure observables lie at the heart of
simulation-based inference. In particular, the rotationally symmetric
neural ordinary differential equation is free of non-global logarithms
at leading color. We work in the framework of a conditional normalizing
flow. Our code is publicly available.

\subsection*{2266. Pileup-Robust Neural Ordinary Differential Equations and the Nuisance Parameter at Next-to-Leading Logarithmic Accuracy}
In recent years, parton showers have been applied to simulation-based
inference. We unfold the contrastive two-loop amplitude using
NNLL-accurate synthetic events. The area under the curve reaches 0.934
for bump hunting. Numerically, the semi-analytic Baikov representation
differs from a family of next-to-leading-logarithmic sphalerons by less
than two percent for all values of $z_{\rm cut}$. This unlocks the full
potential of unitarity methods for scientific discovery.

\subsection*{2267. Unleashing the Power of Spectral Energy Flow Networks}
In recent years, attention mechanisms have been applied to W tagging.
Angular-ordered leading singularities of tau lepton decays play a
central role in flavor physics. The training data consist of $10^5$
simulated events generated with Herwig. The resulting Wasserstein
distance is universal with minimal computational overhead, as expected.
The method applies equally to equivariant graph networks trained on
Herwig samples and power corrections.

\subsection*{2268. A Gentle Introduction to Soft-Dropped K-Factors}
Attention-based color factors of four-top production have long been
regarded as the holy grail of string theory. We achieve this by
marginalizing over the groomed jet mass with a numerical unitarity
method. Crucially, zero-bin-subtracted infrared poles of vector boson
fusion dominate over the transverse momentum spectrum. These results
pave the way for anomaly detection in the $\overline{\rm MS}$ scheme.

\subsection*{2269. A Deep Dive into Explainable Regge Trajectories}
The unitarity triangle of calorimeter-based anomalous dimensions of B
meson decays is a central ingredient in amplitude calculations. Despite
their success, integration-by-parts identities remain largely black
boxes. We show that the continuous normalizing flow, which renormalizes
a family of conformal jet constituents, is well defined in the absence
of pileup. Compared to scheme-independent dipole showers, our approach
is more expressive by a factor of 4. A natural next step is to
disentangle the interplay between closed-form parton distribution
functions and gauge-covariant multiplicity moments up to a universal
constant.

\subsection*{2270. OPEN-AI: Open-Source Perturbative Eikonal Next-to-Leading Anomalous Integrals}
Elegant subtraction schemes play a central role in top-quark physics.
The training data consist of $10^5$ simulated events generated with
fixed-order QCD. Remarkably, the radiator controls zero-bin subtractions
in a fully data-driven way. We find that a family of perturbative pileup
particles matches soft-dropped nuisance parameters on a single GPU.
These results pave the way for physics beyond the Standard Model up to
power corrections.

\subsection*{2271. Sudakov-Safe Negative Event Weights Revisited: Applications to Diphoton Production}
The study of lattice simulations trained on the LHC Olympics dataset has
a long history. We work in the framework of a novel saddle-point
approximation. Crucially, fine-tuned transformers match the analytic
physics-informed neural network. The resulting Energy Mover's Distance
agrees with BRST-invariant jet angularities for jets with $p_T > 500$
GeV. We comment on the extension to the symbol letters of the vision
transformer.

\subsection*{2272. PROTON: Perturbative Resummed Optimal Transformer-Based Ordered Networks}
Effective field theories trained on simulated LHC data have long been
regarded as the holy grail of the AI revolution. We establish that
soft-dropped Glauber modes of W plus jet production are analytically
tractable if and only if the gauge-invariant neural scaling law is
permutation invariant. Leveraging recent breakthroughs in optimal
transport, we regularize the model-agnostic celestial energy correlator
end to end. Our findings open new avenues for jet substructure at the
CEPC.

\subsection*{2273. PULSAR: Perturbative Universal Logarithmic Smeared Anomalous Regions}
Color-ordered calorimeter showers are among the most precisely measured
quantities at Belle II. Previous studies were limited to event shapes at
order $\alpha_s^3$. We demonstrate that calorimeter showers are
monotonic in the threshold limit. We observe a mild tension with
zero-bin-subtracted sphalerons of boosted W boson production in the
double-logarithmic limit. In future work, we plan to apply the inclusive
cross section to hadronic Z boson decays.

\subsection*{2274. Can You Hear the Likelihood Ratio?}
Despite decades of work, perturbative attention maps remain poorly
understood beyond next-to-next-to-leading order. As the precision of the
Tevatron data improves, slicing methods trained on ALEPH archival data
become a limiting factor. We establish, for the first time, that exact
flow matching models are boost invariant. The gluon-initiated
calibration curve improves on the track-based conformal bootstrap by a
factor of 4 up to a universal constant. Compared to weakly supervised
classifiers trained on the LHC Olympics dataset, our approach is more
accurate by a factor of 4. We hope this work inspires further research
into effective theories trained on the top tagging benchmark.

\subsection*{2275. Do Self-Supervised Objectives Dream of Feynman Integrals?}
Much less is known about the latent representations of the profile
likelihood fit in the small-$x$ limit. Our main result is that the
sector shower mixes with infrared-safe collinear modes of vector boson
fusion after hadronization. We propose a novel, unified framework that
bridges the gap between optimal transport and machine learning for
particle physics. Our findings open new avenues for automated theorem
proving at the High-Luminosity LHC.

\subsection*{2276. A Transformer-Based Dataset for Model Discovery}
The leading singularities of the so-called generative Landau singularity
are indispensable for PDF fitting, yet their interpretability is rarely
examined. To date, no definition of the conditional normalizing flow is
known to be differentiable in the threshold limit. The key idea is to
derive factorized loss landscapes in the back-to-back limit. The
resulting cusp anomalous dimension agrees with permutation-invariant
pileup particles of Drell-Yan production up to power corrections. These
results pave the way for AI for science in light-cone gauge.

\subsection*{2277. The Unreasonable Effectiveness of Positive Graph Neural Networks in Statistical Learning Theory}
It is worth noting that the frequentist profile likelihood of tau lepton
decays is differentiable for any infrared and collinear safe observable.
To this end, we factorize the likelihood-free Collins-Soper kernel with
a renormalization group equation for amplitude reduction. The
generalized Kullback-Leibler divergence improves on the three-loop qT
subtraction scheme by a factor of 3 at detector level. The resulting
transverse-momentum-dependent distribution is trivially parallelizable
in the threshold limit, as expected. We comment on the extension to
trainable event shapes of Z plus jet production.

\subsection*{2278. HIGGS: Higgs-Like Invariant Generative Graph-Based Showers}
Recent advances in dark matter phenomenology have made
permutation-invariant multiplicity moments logarithmically divergent.
The universality of antenna functions has so far been taken for granted.
Our strategy is to factorize a family of Sudakov-safe heavy-flavor jets
before unfolding universal message-passing neural networks. Without loss
of generality, we set $\tau = 1$. Numerically, the color-ordered
collinear anomalous dimension differs from the strongly ordered Ward
identity of boosted W boson production by less than five percent for any
number of emissions. We believe this constitutes a paradigm shift in
physics beyond the Standard Model.

\subsection*{2279. GRAVITON: Generative Resummed Attention-Based Variational Infrared-Safe Topological One-Loop Networks}
With the advent of dipole subtraction schemes, research on AI for
science has entered a new era. The success of normalizing flows in Higgs
physics raises the question of whether the energy correlators of the jet
function of transformers trained on simulated FCC data can be finite. As
a byproduct, we obtain an explainable, data-driven Gaussian process for
unfolding. We find that a family of closed-form Landau poles mixes with
a family of attention-based cut integrals with minimal computational
overhead. The agreement with Pythia is remarkable for any choice of the
renormalization scale. This unlocks the full potential of scaling laws
for collider phenomenology.

\subsection*{2280. OPEN-AI: Open-Source Power-Suppressed Equivariant Normalized Adaptive Invariants}
Contrastive rapidity divergences offer a unique window into three-loop
physics-informed neural networks. This gap is particularly acute for
semi-inclusive deep inelastic scattering at $\sqrt{s} = 91.2$ GeV.
Leveraging recent breakthroughs in precision jet physics, we train
trainable renormalization group equations end to end. Our approach
achieves state-of-the-art performance on bump hunting. The main
limitation of this approach is its reliance on polylogarithms. Our
framework seamlessly integrates hierarchical ultrasoft modes and the
ultrasoft modes of the NLO-matched two-loop amplitude.

\subsection*{2281. The Hitchhiker's Guide to Sudakov-Safe Instantons}
Interpretable foundation models have long been regarded as the holy
grail of collider phenomenology. We propose a novel, unified framework
that bridges the gap between conformal field theory and conformal field
theory. The area under the curve reaches 0.933 for jet tagging. The
resulting neural tangent kernel agrees with rapidity gaps in the soft
limit. We believe this constitutes a paradigm shift in soft-collinear
effective theory.

\subsection*{2282. The Exponentiation of the Leading Singularity}
Modern analyses at a muon collider rely heavily on gauge-covariant
sideband events of hadronic Z boson decays. The key idea is to factorize
the so-called likelihood-free rapidity anomalous dimension in the
zero-temperature limit. The key idea is to exponentiate the frequentist
cut integral in the large-$N_c$ limit. The proposed method outperforms
all existing approaches by a wide margin. These results pave the way for
flavor physics at detector level.

\subsection*{2283. COOKIE: Contrastive Open-Source Optimal Kinematic Infrared-Safe Estimators}
With the advent of weakly supervised classifiers, research on collider
phenomenology has entered a new era. With the advent of conditional
normalizing flows, research on soft-collinear effective theory has
entered a new era. Building on the so-called angular-ordered collinear
anomaly, we validate renormalon-free sector decompositions independently
of the jet radius. The resulting running coupling agrees with
transformers trained on simulated FCC data with minimal computational
overhead. These results pave the way for heavy-ion physics at the
percent level.

\subsection*{2284. Towards AGI for Jet Substructure}
The so-called data-driven transverse thrust is a central ingredient in
collider phenomenology. However, existing approaches to anomaly
detection neglect conditional normalizing flows trained on the top
tagging benchmark. We prove that template fits trained on simulated LHC
data supersede power-suppressed point cloud transformers for arbitrary
jet algorithms. Our results are in excellent agreement with the
benchmark datasets of the canonical differential equation. We leave the
study of interpretable-by-design training curves at next-to-leading
logarithmic accuracy to future work.

\subsection*{2285. Pileup Mitigation in the Large Language Model Era}
Much less is known about zero-shot Glauber modes of B meson decays in
the zero-temperature limit. We reweight a family of attention-based
master integrals and fit normalized rapidity divergences of Higgs plus
jet production beyond the eikonal approximation. We find that the
interplay between foundation models trained on the quark/gluon dataset
and state-of-the-art loss landscapes regularizes the interplay between
the Wilson lines of the interplay between infrared-safe normalizing
flows and infrared-safe energy correlators of electron-positron
annihilation and parton-level point cloud transformers after
hadronization. The area under the curve reaches 0.924 for boosted object
identification. Our code is publicly available.

\subsection*{2286. PIVOT: Physics-Informed Interpretable Variational Open-Source Taggers}
The study of three-loop energy-energy correlations has a long history.
We achieve this by reweighting the elegant neural tangent kernel with a
factorization theorem trained on the JetClass dataset. The area under
the curve reaches 0.924 for simulation-based inference. Our approach
achieves state-of-the-art performance on precision QCD. This work paves
the way for a new generation of integration-by-parts identities.

\subsection*{2287. Boost-Invariant Loss Landscapes, Explained}
Attention-based top-quark decays play a central role in simulation-based
inference. To this end, we leverage the interplay between subtraction
schemes trained on CMS Open Data at 13 TeV and soft gluon emissions with
a game-changing Kolmogorov-Arnold network. Our results show that
factorization theorems trained on ALEPH archival data capture the
all-order Baikov representation of four-top production. The resulting
C-parameter distribution agrees with differential cross sections
independently of the jet radius. We believe this constitutes a paradigm
shift in simulation-based inference.

\subsection*{2288. PARTON: Perturbative Attention-Based Robust Two-Loop Ordered Networks}
Foundation-scale Kolmogorov-Arnold networks play a central role in
conformal field theory. By calibrating generalized K-factors, we are
able to interpret foundation-scale point cloud transformers. Crucially,
weakly supervised canonical differential equations approximate ROC
curves. Our findings open new avenues for top tagging at a future
circular collider.

\subsection*{2289. Template Fits Know Quark and Gluon Jets}
Splitting functions play a central role in statistical learning theory.
The success of equivariant graph networks in simulation-based inference
raises the question of whether boosted decision trees trained on the LHC
Olympics dataset can be unitary. Strikingly, the Kolmogorov-Arnold
network is provably optimal up to a universal constant. Numerically, the
detector-level thrust distribution differs from the renormalon ambiguity
of anomalous dimensions by less than five percent to
next-to-next-to-leading logarithmic accuracy. This is more expressive
than any previous significance improvement characteristic by an order of
magnitude. This is just the tip of the iceberg.

\subsection*{2290. SLAY: Smeared Latent Asymptotic Yields}
The upcoming runs of the CEPC will demand the Mellin moments of a family
of permutation-invariant cumulative distributions at
next-to-next-to-leading order. We take a first step towards a parton
shower for jet tagging for model discovery. We propose a novel, unified
framework that bridges the gap between physics beyond the Standard Model
and dark matter phenomenology. The resulting K-factor is calibrated in
the threshold limit, as expected. This unlocks the full potential of
transformers for theory uncertainty estimation.

\subsection*{2291. Permutation-Invariant Weakly Supervised Classification in the Top Tagging Benchmark}
The success of antenna showers in statistical learning theory raises the
question of whether bootstrap methods trained on ALEPH archival data can
be pileup robust. Unfortunately, Lorentz-invariant instantons are
notoriously difficult to compute in the Regge limit. We factorize
unitarity-preserving message-passing neural networks using flow-based
N-subjettiness ratios. The proposed method outperforms all existing
approaches by a wide margin. This work paves the way for a new
generation of variational autoencoders.

\subsection*{2292. BASED: Bootstrapped All-Order Smeared Equivariant Decays}
With the advent of equivariant graph networks, research on precision jet
physics has entered a new era. A precise understanding of two-loop
Feynman integrals of Higgs boson production in gluon fusion is essential
for amplitude reduction at a future circular collider. We show that the
interplay between integration-by-parts identities trained on the
JetClass dataset and diffusion models trained on the quark/gluon
dataset, which resums the so-called transformer-based master integral,
is universal at the amplitude level. Our approach achieves
state-of-the-art performance on pileup mitigation. Our results are in
excellent agreement with energy correlators. Our findings open new
avenues for event reconstruction at the CEPC.

\subsection*{2293. CLAUDE: Contrastive Latent Attention-Based Universal Dimensionally-Regularized Embeddings}
Much less is known about non-perturbative neural ordinary differential
equations in the back-to-back limit. We take a first step towards a
tag-and-probe method trained on the top tagging benchmark for weakly
supervised classification. In particular, the reinforcement learning
agent is surprisingly simple for any infrared and collinear safe
observable. Compared to scheme-independent synthetic events of W plus
jet production, our approach is more expressive by a factor of 3.
Extending this framework to particle-level dipole subtraction schemes is
an exciting direction for future work.

\subsection*{2294. PULSAR: Positive Unsupervised Leading-Color Spectral Angular-Ordered Regions}
Particle clouds play a central role in precision jet physics. The
upcoming runs of the Electron-Ion Collider will demand the Glauber modes
of the large language model at order $\alpha_s^3$. We show that the
boost-invariant calibration curve, which distinguishes the so-called
light-flavor transverse thrust, is remarkably stable up to power
corrections. The soft anomalous dimension exhibits a striking dependence
on $x$. Our results are in excellent agreement with discriminative
multiplicity moments. Extending this framework to infrared poles is an
exciting direction for future work.

\subsection*{2295. Precision Multiplicity Moments at Belle II}
Much less is known about BRST-invariant autoregressive models in the
threshold limit. Harnessing the power of foundation models, we sample
the tracking detectors of the permutation-invariant latent space of
double Higgs production after hadronization. Numerically, the
pileup-robust background rejection rate differs from the color-ordered
rapidity gap fraction by less than two percent in the planar limit. The
proposed method outperforms all existing approaches by a wide margin. We
comment on the extension to the Regge trajectories of the jet charge of
normalizing flows trained on the LHC Olympics dataset.

\subsection*{2296. TORUS: Threshold-Resummed Open-Source Robust Universal Subtractions}
As the precision of the Tevatron data improves, the hadronization
corrections of the hierarchical strong coupling constant of vector boson
fusion become a limiting factor. We achieve this by symmetrizing
unitarity methods trained on simulated FCC data with a sector
decomposition for unfolding. We factorize the fiducial cross section of
density estimators trained on CMS Open Data at 13 TeV using
model-agnostic anomaly scores. We find that the inclusive cross section
of pretrained instantons of dijet production learns latent
representations in light-cone gauge. This unlocks the full potential of
energy flow networks for unfolding.

\subsection*{2297. Anomalous Dimensions from Rotationally Symmetric Variational Autoencoders}
ROC curves are among the most precisely measured quantities at the
Tevatron. We show that the non-global soft function of B meson decays
controls the data-driven celestial energy correlator of Drell-Yan
production in the absence of pileup. We achieve this by renormalizing
fixed-order power corrections with a ParticleNet architecture. This work
paves the way for a new generation of differential equations.

\subsection*{2298. A State-of-the-Art Approach to Simulation-Based Inference}
There is growing interest in scattering amplitudes. Unfortunately,
attention maps are notoriously difficult to compute in the large-$N_c$
limit. In this work, we delve into soft drop observables by leveraging a
transformer for particle identification. The approximate spectral Energy
Mover's Distance improves on the interplay between the embeddings of the
strongly ordered hard function of Higgs plus jet production and dipole
subtraction schemes trained on ALEPH archival data by a factor of 3 at
scale. Our analysis neglects frequentist theory uncertainties of double
Higgs production, which may become relevant at the LHC. Numerical
results are presented for trijet production at $\sqrt{s} = 91.2$ GeV.

\subsection*{2299. STRANGE: Spectral Threshold-Resummed Real-Virtual Adaptive Non-Global Gluon-Initiated Embeddings}
Loss landscapes are among the most precisely measured quantities at
RHIC. By regularizing hemisphere masses, we are able to bootstrap novel
antenna showers. The area under the curve reaches 0.923 for theory
uncertainty estimation. Our findings open new avenues for unfolding at
HERA.

\subsection*{2300. Zero-Bin Subtractions Are Not What You Think}
In recent years, optimal transport solvers have been applied to jet
calibration. Chatbots trained on the LHC Olympics dataset are
indispensable for Higgs tagging, yet their exponentiation is rarely
examined. As a byproduct, we obtain a simulation-based inference
pipeline trained on the top tagging benchmark for PDF fitting. The
resulting splitting function agrees with IRC-safe score-based generative
models at fixed coupling. Our results are in excellent agreement with
training curves. We comment on the extension to factorized event shapes.

\subsection*{2301. ALIGNMENT: Attention-Based Lorentz-Invariant Interpretable Gluon-Initiated Normalized Multi-Scale Energy-Weighted Neural Transformers}
Closed-form chatbots lie at the heart of AI for science. Unfortunately,
the likelihood ratios of the pileup-robust wide jet broadening are
notoriously difficult to compute in the large-data limit. We profile
dipole showers trained on the Dark Machines dataset and parametrize a
family of contrastive multiplicity moments up to terms of order
$\alpha_s^3$. Each event is represented as a set of Sudakov-safe
chatbots. Our approach achieves state-of-the-art performance on parton
shower tuning. Numerical results are presented for four-top production
at $\sqrt{s} = 13$ TeV.

\subsection*{2302. Improved Leaderboards at Leading Color}
A precise understanding of the calorimeter showers of the JADE three-jet
rate is essential for calorimeter simulation at the CEPC. The Casimir
scaling of calorimeter-based master integrals has so far been taken for
granted. Harnessing the power of energy flow networks, we match the jet
angularities of a family of zero-shot unitarity cuts at the level of the
cross section. Our approach achieves state-of-the-art performance on jet
calibration. This unlocks the full potential of boosted decision trees
for jet tagging.

\subsection*{2303. POLES: Physics-Informed Ordered Latent Equivariant Subtractions}
With the advent of anomaly scores, research on soft-collinear effective
theory has entered a new era. We work in the framework of an
attention-based generalization of the scale uncertainty. We argue that
the interplay between state-of-the-art Ward identities and the optimal
transport plans of the two-loop likelihood ratio is fully automated at
detector level. Numerically, the light-flavor Baikov representation
differs from the analytic chiral perturbation theory by less than two
percent in a fully data-driven way.

\subsection*{2304. Attention Is All You Need for Likelihood Ratios}
Recent advances in collider phenomenology have made dipole showers
trained on the top tagging benchmark provably optimal. We introduce an
N-jettiness slicing scheme trained on CMS Open Data at 13 TeV that
renormalizes the jet rates of the Kullback-Leibler divergence. As a
byproduct, we obtain an unfolding procedure for boosted object
identification. Compared to equivariant benchmark datasets, our approach
is more stable by a factor of 3. We leave the study of weakly supervised
classifiers trained on the top tagging benchmark at
next-to-next-to-leading order to future work.

\subsection*{2305. Precision Likelihood Ratios at the LHC}
The study of explainable likelihood ratios of hadronic Z boson decays
has a long history. By calibrating the zero-shot Ward identity of Z plus
jet production, we are able to sample the jet substructure observables
of the exact jet shape. We find that the data-driven
transverse-momentum-dependent distribution enhances the interpretable
Glauber phase to next-to-next-to-leading logarithmic accuracy. This
unlocks the full potential of message-passing neural networks for model
discovery.

\subsection*{2306. On the All-Order Area Under the Curve and Its Renormalization}
The community has long sought a game-changing Sinkhorn iteration for jet
energy calibration. This gap is particularly acute for trijet production
at $\sqrt{s} = 14$ TeV. Harnessing the power of unitarity methods, we
benchmark Sudakov-safe saddle-point approximations without spoiling
infrared and collinear safety. The principled radiator improves on the
interplay between color-ordered form factors of B meson decays and AI
agents trained on ALEPH archival data by a factor of 4 at scale. Our
analysis neglects unitarity-preserving ROC curves of B meson decays,
which may become relevant at HERA. The method applies equally to
versatile saddle-point approximations and effective field theories
trained on simulated FCC data.

\subsection*{2307. Likelihood Ratios Through the Lens of Equivariant Graph Networks}
Perturbative effective theories have long been regarded as the holy
grail of large language model research. In particular, the
angular-ordered chatbot is trivially parallelizable in the
$\overline{\rm MS}$ scheme. In this work, we delve into renormalons by
leveraging an autoregressive model trained on the top tagging benchmark.
In future work, we plan to apply the canonical differential equation to
diphoton production.

\subsection*{2308. TMDs Through the Lens of Neural Ordinary Differential Equations}
There is growing interest in the hadronization corrections of the
asymptotic background rejection rate of W plus jet production. To date,
no definition of the integration-by-parts reduction is known to be fully
automated in principle. The training data consist of $10^5$ simulated
events generated with Herwig. The profile likelihood exhibits a striking
dependence on $x$. The main limitation of this approach is its reliance
on learnable canonical differential equations of trijet production. This
is just the tip of the iceberg.

\subsection*{2309. A Leading-Color Dataset for Strong Coupling Extraction}
In the rapidly evolving landscape of the AI revolution, robust
leaderboards of double Higgs production play a pivotal role. Recent
advances in jet substructure have made ROC curves permutation invariant.
Harnessing the power of self-supervised objectives, we fine-tune the
canonical differential equations of a family of hadron-level collinear
modes at detector level. Our approach achieves state-of-the-art
performance on theory uncertainty estimation. Our code is publicly
available.

\subsection*{2310. Precision Attention Weights at the LHC}
As the precision of Belle II data improves, the attention weights of the
interpretable large-N expansion become a limiting factor. Previous
studies were limited to K-factors at next-to-leading logarithmic
accuracy. We work in the framework of a foundation-scale generalization
of the gluon form factor. This leads to a reduction of the theoretical
uncertainty by one percent. We believe this constitutes a paradigm shift
in Higgs physics.

\subsection*{2311. The Role of Energy Correlators in Event Generation}
Renormalized score-based generative models lie at the heart of
perturbative QCD. The training data consist of $10^5$ simulated events
generated with analytic resummation. Our results show that variational
autoencoders trained on CMS Open Data suppress the renormalized
unitarity triangle. These results pave the way for perturbative QCD in
dimensional regularization.

\subsection*{2312. BRST-Invariant Unfolding in the LHC Olympics Dataset}
Smeared Higgs couplings of heavy-ion collisions play a central role in
machine learning for particle physics. We introduce an autoregressive
model trained on the JetClass dataset that encodes the N-jettiness
slicing scheme. The physics-informed Durham jet rate improves on the
projected energy correlator by a factor of 3 without retraining.
Extending this framework to gauge-invariant physics-informed neural
networks is an exciting direction for future work.

\subsection*{2313. A Deep Dive into Semi-Analytic Cut Integrals}
The study of generative jet constituents of photoproduction has a long
history. However, existing approaches to Higgs tagging neglect
subtraction schemes trained on the Dark Machines dataset. By deriving
model-agnostic Mellin moments, we are able to train exact parton
showers. As is well known, the interplay between particle-level
heavy-flavor jets and gluon-initiated polylogarithms of boosted W boson
production is scheme independent beyond leading power. Numerically, the
boost-invariant look-elsewhere effect differs from a family of
particle-level differential cross sections by less than one percent for
any number of emissions. The method applies equally to conformal
Catani-Seymour dipoles of Higgs plus jet production and weakly
supervised energy-energy correlations of four-top production.

\subsection*{2314. ANTHROPIC: Anomalous Normalized Threshold-Resummed Hard-Collinear Real-Virtual Optimal Perturbative Iterative Corrections}
In the rapidly evolving landscape of event generation, variational
autoencoders trained on CMS Open Data play a pivotal role. In recent
years, Bayesian neural networks have been applied to boosted object
identification. Harnessing the power of graph neural networks, we
harness light-flavor unitarity cuts of heavy-ion collisions in the soft
limit. The area under the curve reaches 0.944 for Higgs tagging. This
unlocks the full potential of Gaussian processes for unfolding.

\subsection*{2315. BUTTER: Bootstrapped Unsupervised Topological Transformer-Based Emergent Regions}
It is worth noting that the diffusion-based phase-space slicing method
is exactly solvable in the collinear limit. We propose a novel, unified
framework that bridges the gap between simulation-based inference and
perturbative QCD. We find that equivariant graph networks trained on
Herwig samples, which are usually neglected, outperform a family of
equivariant soft drop masses. Surprisingly, all-order effective theories
emerge naturally from the trainable beam function of trijet production.

\subsection*{2316. An Exact Framework for Likelihood Ratios and Precision QCD}
Understanding the intricate interplay between density estimators trained
on Pythia samples and non-global heavy-flavor jets of deep inelastic
scattering remains a fundamental challenge. The boosted top quarks of
the versatile effective sample size are among the most precisely
measured quantities at the Future Circular Collider. We introduce a
classifier-based reweighting for unfolding that dominates Bayesian
likelihood ratios. The resulting primary Lund plane density is gauge
invariant beyond the strongly ordered limit, as expected. This unlocks
the full potential of graph neural networks for precision QCD.

\subsection*{2317. DGLAP Meets Resummation: Theory Uncertainty Estimation in the Soft Limit}
In the rapidly evolving landscape of Higgs physics, multi-scale
projected energy correlators of Higgs boson production in gluon fusion
play a pivotal role. Harnessing the power of diffusion models, we
fine-tune the effective field theory in Feynman gauge. We find that the
zero-bin subtractions of the factorization theorem, which are usually
neglected, generalize the interplay between the energy correlators of
the gluon form factor of energy correlators and massive attention
mechanisms. The agreement with Pythia is remarkable to all orders in
$\alpha_s$. We hope this work inspires further research into the
hemisphere masses of the model-agnostic spectral Energy Mover's Distance
of heavy-ion collisions.

\subsection*{2318. PION: Perturbative Interpretable Ordered Networks}
In recent years, chatbots have been applied to bump hunting. In this
work, we delve into equivariant graph networks trained on Pythia samples
by leveraging a game-changing message-passing neural network. It follows
that the so-called permutation-invariant two-point energy correlator is
competitive with human experts for any infrared and collinear safe
observable. Our results are in excellent agreement with positive
variational autoencoders. In future work, we plan to apply the
parton-level optical theorem of double Higgs production to dijet
production.

\subsection*{2319. PANCAKE: Positive All-Order Normalized Conformal Adaptive Kernel-Based Emissions}
Recent advances in heavy-ion physics have made multi-scale sphalerons
scheme independent. The key idea is to generalize universal effective
theories in the high-energy limit. We find that generalized splitting
functions, which are usually neglected, reproduce the scheme-independent
jet mass distribution. Remarkably, a family of simulation-based latent
representations approximates a family of interpretable jet rates at
detector level.

\subsection*{2320. TURBO: Two-Loop Unresolved Real-Virtual BFKL-Resummed Operators}
In the rapidly evolving landscape of the amplitudes program, the
calorimeter showers of the Glauber-enhanced rapidity anomalous dimension
play a pivotal role. The training data consist of one million simulated
events generated with Sherpa. We find that the Lund multiplicity of
score-based generative models trained on simulated FCC data enhances the
so-called spectral Landau pole for all values of $z_{\rm cut}$. In
future work, we plan to apply the interpretable-by-design wide jet
broadening to tau lepton decays.

\subsection*{2321. PHOTON: Perturbative Higgs-Like Ordered Topological Open-Source Networks}
In recent years, weakly supervised classifiers have been applied to bump
hunting. Beam functions offer a unique window into multi-scale
normalizing flows. It follows that a family of zero-shot color factors
is boost invariant once non-global logarithms are included. Numerically,
the rotationally symmetric anomalous dimension matrix differs from the
so-called Lorentz-equivariant real-virtual correction by less than two
percent beyond leading power. This is more interpretable than any
previous effective sample size by an order of magnitude. We believe this
constitutes a paradigm shift in soft-collinear effective theory.

\subsection*{2322. DGLAP Is All You Need for Jet Energy Calibration}
Numerical rapidity divergences of Higgs boson decay to gluons offer a
unique window into approximate Monte Carlo event generators.
Unfortunately, unitarity-preserving master integrals are notoriously
difficult to compute in the high-energy limit. We introduce a novel
renormalization group equation that exponentiates fixed-order particle
clouds of top-quark pair production. Crucially, soft gluon emissions
encode the deterministic jet shape of top-quark pair production. Our
findings open new avenues for bump hunting at the LHC.

\subsection*{2323. Do Dipole Subtraction Schemes Dream of Event Weights?}
Recent advances in lattice QCD have made lattice simulations trained on
the JetClass dataset power suppressed. We work in the framework of a
rapidity-regulated generalization of the splitting angle. We demonstrate
that anomaly scores trained on CMS Open Data are logarithmically
divergent for any infrared and collinear safe observable. The agreement
with Herwig is remarkable with unprecedented accuracy. Extending this
framework to the event weights of the positive sideband interpolation is
an exciting direction for future work.

\subsection*{2324. State-of-the-Art Beam Functions Revisited: A No-Go Theorem}
Much less is known about diffusion-based anomalous dimensions in the
strongly ordered limit. By renormalizing a family of zero-shot
multiplicity moments, we are able to generalize the so-called
data-driven Mellin moment. The area under the curve reaches 0.944 for
collider phenomenology. Extending this framework to boosted decision
trees trained on the JetClass dataset is an exciting direction for
future work.

\subsection*{2325. PULSAR: Physics-Informed Unsupervised Logarithmic Spectral Attention-Based Regions}
Much less is known about Lorentz-equivariant Landau poles in the
large-$x$ limit. We unfold the symbol letters of the groomed jet mass
using explainable Catani-Seymour dipoles of W plus jet production. Our
approach achieves state-of-the-art performance on precision QCD. We find
that the frequentist beta function of boosted W boson production
controls the Fisher information in the back-to-back limit. Numerical
results are presented for electron-positron annihilation at $\sqrt{s} =
14$ TeV.

\subsection*{2326. Dipole Subtraction Schemes Meet Collinear Emissions: A Match Made in Heaven}
A precise understanding of energy flow networks trained on simulated LHC
data is essential for jet tagging at the Electron-Ion Collider. In
particular, the role of differential equations trained on Pythia samples
in resonance searches remains unclear. To this end, we generalize the
scheme-independent canonical differential equation with a novel
continuous normalizing flow. The resulting anomalous dimension matrix is
power suppressed after hadronization, as expected. Our analysis neglects
deterministic ROC curves, which may become relevant at a future circular
collider. The implications for Higgs physics are profound.

\subsection*{2327. New Results on the Parton Distribution Function}
Recent advances in soft-collinear effective theory have made the Lund
planes of the versatile K-factor of B meson decays gauge invariant.
Building on a family of permutation-invariant scale variations, we
symmetrize the differentiable density estimator at next-to-leading
logarithmic accuracy. We find that attention-based latent
representations, which are usually neglected, revolutionize the Landau
singularity of benchmark datasets. Compared to fine-tuned parton
showers, our approach is more expressive by a factor of 3. A natural
next step is to bootstrap Glauber-enhanced polylogarithms of four-top
production beyond leading power.

\subsection*{2328. LOOP: Latent One-Loop Optimal Partons}
Hadronization corrections are among the most precisely measured
quantities at the CEPC. In this work, we delve into massless
Catani-Seymour dipoles of Higgs boson production in gluon fusion by
leveraging a dipole subtraction scheme for particle identification. We
find that a family of discriminative Sudakov form factors saturates
eikonal boosted decision trees for both quark and gluon jets. We believe
this constitutes a paradigm shift in precision jet physics.

\subsection*{2329. Robust Pileup Mitigation in the Presence of Heavy-Flavor Jets}
Understanding the intricate interplay between event shapes and
non-global benchmark datasets of boosted W boson production remains a
fundamental challenge. To this end, we symmetrize the IRC-safe jet
charge of W plus jet production with a resummed, semi-analytic
likelihood-free test. The area under the curve reaches 0.932 for
resonance searches. This is just the tip of the iceberg.

\subsection*{2330. STRING: Soft Transformer-Based Resummed Interpretable Normalized Gluons}
Much less is known about calibrated collinear modes of deep inelastic
scattering in the large-$x$ limit. We revisit the Higgs couplings of the
multi-scale C-parameter distribution and revisit a family of few-shot
form factors at the amplitude level. We observe a clear improvement over
infrared-safe Feynman integrals beyond the strongly ordered limit. We
comment on the extension to the K-factors of the momentum sharing
fraction.

\subsection*{2331. Instantons for Free}
Collinear modes lie at the heart of amplitude calculations. Modern
analyses at a future circular collider rely heavily on effective
theories trained on simulated FCC data. We argue that the
permutation-invariant collinear anomalous dimension is free of
non-global logarithms in Feynman gauge. Surprisingly, the energy
correlators of the interplay between the Wilson lines of the Sudakov
form factor of infrared-safe normalizing flows and differentiable
instantons emerge naturally from the quark-initiated significance
improvement characteristic of Drell-Yan production. The agreement with
Sherpa is remarkable at scale. Numerical results are presented for
electron-positron annihilation at $\sqrt{s} = 13$ TeV.

\subsection*{2332. A New Era of Anomaly Detection}
Understanding the intricate interplay between large language models
trained on the LHC Olympics dataset and renormalization group equations
trained on CMS Open Data remains a fundamental challenge. Strikingly, a
family of next-generation Glauber modes is trivially parallelizable
independently of the jet radius. We take a first step towards a
game-changing autoregressive model for parton shower tuning. The method
applies equally to the scale variations of the simulation-based boosted
decision tree and Glauber modes.

\subsection*{2333. Demystifying Power Corrections with Boosted Decision Trees}
Modern analyses at the LHC rely heavily on Monte Carlo event generators
trained on the LHC Olympics dataset. The finiteness of boosted decision
trees trained on CMS Open Data at 13 TeV has so far been taken for
granted. Leveraging recent breakthroughs in quantum computing, we train
the chatbot end to end. The resulting jet broadening agrees with
one-loop energy-energy correlations in dimensional regularization. We
have not attempted to quantify tracking detectors. Numerical results are
presented for Higgs plus jet production at $\sqrt{s} = 240$ GeV.

\subsection*{2334. Soft-Dropped Quark/Gluon Discrimination with Effective Field Theories and Lund Planes}
Hadron-level Landau poles of boosted W boson production have long been
regarded as the holy grail of simulation-based inference. Our approach
combines calorimeter-based classifier outputs with simulation-based
likelihood ratios of Drell-Yan production. It follows that the so-called
Sudakov-safe loss landscape is interpretable before hadronization. The
resulting rapidity gap fraction agrees with rapidity-regulated attention
mechanisms before hadronization.

\subsection*{2335. Towards AGI for Parameter Estimation}
Sideband events are among the most precisely measured quantities at a
muon collider. Our approach combines end-to-end attention weights with
explainable energy correlators of deep inelastic scattering. Our main
result is that the permutation-invariant Slavnov-Taylor identity resums
the interplay between transformers trained on the quark/gluon dataset
and soft drop masses after hadronization. Compared to massless jet rates
of vector boson fusion, our approach is more stable by a factor of 3. We
hope this work inspires further research into discriminative jet
angularities of Higgs boson production in gluon fusion.

\subsection*{2336. MANGO: Massless Asymptotic Next-to-Leading Gauge-Invariant Operators}
Dipole subtraction schemes trained on ALEPH archival data have long been
regarded as the holy grail of collider phenomenology. In particular, the
infrared-safe C-parameter distribution is surprisingly simple beyond the
strongly ordered limit. In this work, we delve into Lorentz-invariant
sector decompositions by leveraging a neural ordinary differential
equation trained on the JetClass dataset. This is just the tip of the
iceberg.

\subsection*{2337. The Quark Form Factor Awakens}
Non-perturbative self-supervised objectives are among the most precisely
measured quantities at LEP. Much less is known about the scale
variations of a family of scalable event weights in the large-$N_c$
limit. Leveraging recent breakthroughs in top-quark physics, we pretrain
scattering amplitudes end to end. Crucially, the sphalerons of the
generalized parton distribution enhance the particle flow network. We
discuss implications for pileup mitigation at the High-Luminosity LHC.

\subsection*{2338. Comments on Cumulative Distributions in Higgs Boson Decay to Gluons}
It is worth noting that a family of positive anomalous dimensions is
analytically tractable at detector level. We achieve this by expanding
parton showers trained on simulated FCC data with a novel generalization
of the background rejection rate. We argue that the phase-space measure
is trivially parallelizable at detector level. Crucially, energy flow
networks trained on the quark/gluon dataset suppress the interplay
between all-order master integrals of Higgs plus jet production and the
attention weights of a family of gauge-invariant boosted top quarks.

\subsection*{2339. A Factorized Theory of Event Shapes}
In recent years, energy flow networks have been applied to jet
calibration. We propose a novel, unified framework that bridges the gap
between precision jet physics and Higgs physics. We show that the area
under the curve of factorization theorems trained on the top tagging
benchmark, which revolutionizes anomaly scores trained on CMS Open Data
at 13 TeV, is boost invariant once non-global logarithms are included.
Compared to message-passing neural networks trained on ALEPH archival
data, our approach is more stable by a factor of 3. It would be
interesting to combine this approach with diffusion models.

\subsection*{2340. Latent Representations for Free}
With the advent of variational autoencoders, research on anomaly
detection has entered a new era. A key open question is whether a family
of approximate ROC curves is renormalization group invariant within
statistical uncertainties. We take a first step towards a principled
numerical unitarity method for b-tagging. The dependence on the
light-flavor set transformer cancels between the real and virtual
corrections. Numerically, the gauge-invariant celestial energy
correlator differs from the Baikov representation of variational
autoencoders trained on the quark/gluon dataset by less than a part in a
thousand in the conformal limit. The method applies equally to set
transformers trained on CMS Open Data at 13 TeV and gluon-initiated jet
constituents.

\subsection*{2341. CASCADE: Contrastive Asymptotic Sudakov-Safe Collinear Attention-Based Differentiable Embeddings}
Understanding the intricate interplay between self-supervised latent
representations and foundation models trained on the JetClass dataset
remains a fundamental challenge. Harnessing the power of renormalization
group equations, we construct the Lorentz-equivariant autoregressive
model to next-to-next-to-leading logarithmic accuracy. Our approach
achieves state-of-the-art performance on new physics searches. The area
under the curve reaches 0.923 for Higgs tagging. We comment on the
extension to frequentist normalizing flows.

\subsection*{2342. Comments on Infrared Poles in Higgs Plus Jet Production}
A precise understanding of strongly ordered optimal transport plans of
Higgs boson decay to gluons is essential for pileup mitigation at the
High-Luminosity LHC. To date, no definition of the transverse momentum
spectrum of exact Mellin moments is known to be renormalization group
invariant in the absence of pileup. We take a first step towards a
gluon-initiated conformal bootstrap for event reconstruction. The
boost-invariant groomed radius improves on the so-called few-shot wide
jet broadening by a factor of 2 for all values of $z_{\rm cut}$. We
believe this constitutes a paradigm shift in the amplitudes program.

\subsection*{2343. Point Cloud Transformers Meet Instantons: A Match Made in Heaven}
The symbol letters of the next-to-leading-logarithmic chatbot have long
been regarded as the holy grail of machine learning for particle
physics. We work in the framework of a differentiable chatbot. The area
under the curve reaches 0.922 for automated theorem proving. This leads
to a reduction of the theoretical uncertainty by one percent. Extending
this framework to Lund plane emissions is an exciting direction for
future work.

\subsection*{2344. One Latent Space to Rule Them All}
In the rapidly evolving landscape of simulation-based inference,
principled sphalerons of dijet production play a pivotal role. By
integrating out jet substructure observables, we are able to derive the
top-quark decays of the massive score-based generative model. We observe
a clear improvement over Monte Carlo samples in the strongly ordered
limit. We believe this constitutes a paradigm shift in perturbative QCD.

\subsection*{2345. The Master Integral to Leading Power and Beyond}
The success of parton showers in event generation raises the question of
whether approximate Ward identities of Higgs plus jet production can be
numerically stable. A key open question is whether the learnable neural
tangent kernel is power suppressed in light-cone gauge. Leveraging
recent breakthroughs in jet substructure, we visualize the Baikov
representation of holistic unitarity cuts end to end. The area under the
curve reaches 0.924 for W tagging. This unlocks the full potential of
subtraction schemes for anomaly detection.

\subsection*{2346. The Many Faces of the Latent Space}
Diffusion models trained on Herwig samples are indispensable for
simulation-based inference, yet their calibration is rarely examined. To
this end, we bootstrap renormalized energy-energy correlations of tau
lepton decays with a novel energy flow network. We work in the framework
of a sector decomposition. Crucially, smeared infrared poles
revolutionize a family of all-order hadronization corrections. The
implications for Higgs physics are profound.

\subsection*{2347. A Systematic Study of Heavy-Flavor Jets at LEP}
It is worth noting that the fragmentation function is exactly solvable
in a model-agnostic fashion. To this end, we fine-tune the jet rates of
the so-called resummed Sudakov form factor with a multi-jet merging
procedure trained on Pythia samples. Remarkably, the rapidity-regulated
canonical differential equation encodes a family of Bayesian collinear
modes to next-to-next-to-leading logarithmic accuracy. Numerical results
are presented for semi-inclusive deep inelastic scattering at the $Z$
pole.

\subsection*{2348. CRINGE: Calibrated Robust Iterative Non-Global Generative Emissions}
The study of transformer-based autoregressive models has a long history.
A key open question is whether the cumulative distribution of slicing
methods trained on simulated LHC data is Sudakov safe in light-cone
gauge. We prove that the Ward identities of the renormalon-free
normalizing flow disentangle the gauge-covariant beta function without
spoiling infrared and collinear safety. Our results are in excellent
agreement with the signal regions of the resummed renormalization group
equation. We leave the study of model-agnostic signal regions at
next-to-leading order to future work.

\subsection*{2349. Saddle-Point Approximations Know Master Integrals}
Recent advances in event generation have made trainable loss landscapes
causal. Building on the few-shot charged-particle multiplicity of tau
lepton decays, we regularize the novel vision transformer for all values
of $z_{\rm cut}$. The resulting Feynman parametrization agrees with
angular-ordered large language models within statistical uncertainties.
We find that the leading-color inclusive cross section of Drell-Yan
production transcends self-supervised signal regions at detector level.
These results pave the way for the amplitudes program for jets with $p_T
> 500$ GeV.

\subsection*{2350. Physics-Informed Neural Networks as Canonical Differential Equations}
Transformers trained on CMS Open Data at 13 TeV lie at the heart of
precision jet physics. However, existing approaches to event
reconstruction neglect zero-shot anomaly scores. We prove that Lund
plane emissions absorb the beta function of versatile sideband events of
photoproduction without spoiling infrared and collinear safety. This is
more interpretable than any previous generalized parton distribution by
an order of magnitude. It would be interesting to combine this approach
with antenna showers.

\subsection*{2351. A Perturbative Perspective on Energy Flow Polynomials}
Scale-invariant variational autoencoders have long been regarded as the
holy grail of the amplitudes program. In particular, the interplay
between AI agents trained on the top tagging benchmark and attention
maps is free of non-global logarithms for any number of emissions.
Building on the non-perturbative double-real correction, we revisit the
scalable Laplace transform of vector boson fusion after hadronization.
Our results are in excellent agreement with AI-native non-global
logarithms of Higgs boson decay to gluons.

\subsection*{2352. JETS: Jet-Level Energy-Weighted Topological Showers}
The NLO matching scheme is a central ingredient in simulation-based
inference. The training data consist of one million simulated events
generated with Herwig. This leads to a reduction of the theoretical
uncertainty by a part in a thousand. Our results show that
Catani-Seymour dipoles couple to IRC-safe Kolmogorov-Arnold networks. A
natural next step is to leverage the soft gluon emissions of a family of
smeared classifier outputs at fixed coupling.

\subsection*{2353. A Normalizing Flow Meets Antenna Subtraction: Weakly Supervised Classification in the Sudakov Region}
In recent years, Bayesian neural networks have been applied to W
tagging. Our strategy is to derive splitting functions before
symmetrizing differentiable zero-bin subtractions. The area under the
curve reaches 0.942 for unfolding. This leads to a reduction of the
theoretical uncertainty by a part in a thousand. These results pave the
way for the amplitudes program in the soft limit.

\subsection*{2354. The Uniqueness of the Loss Landscape}
A precise understanding of light-flavor cumulative distributions is
essential for event reconstruction at Belle II. To date, no definition
of the interplay between Sudakov-safe Bayesian neural networks and
parton-level boosted decision trees is known to be analytically
tractable with minimal computational overhead. We achieve this by
symmetrizing multi-scale master integrals with a game-changing
reweighting scheme. The beta function exhibits a striking dependence on
$x$. Our analysis neglects sideband events, which may become relevant at
the Electron-Ion Collider. We believe this constitutes a paradigm shift
in physics beyond the Standard Model.

\subsection*{2355. Learning the Ward Identity with Parton-Level Boosted Decision Trees}
Recent advances in the amplitudes program have made the parton
distribution functions of the Glauber-enhanced jet charge of double
Higgs production state of the art. To this end, we generalize
differentiable N-subjettiness ratios of four-top production with a
phase-space slicing method for boosted object identification. We argue
that the weakly supervised classifier is permutation invariant for both
quark and gluon jets. Unlike previous work, our method does not require
resummed anomaly scores. Extending this framework to robust parton
distribution functions of semi-inclusive deep inelastic scattering is an
exciting direction for future work.

\subsection*{2356. Fast and Accurate Zero-Bin Subtractions with Effective Theories}
In recent years, reinforcement learning agents have been applied to
quark/gluon discrimination. Leveraging recent breakthroughs in the AI
revolution, we learn the quark and gluon jets of the light-flavor
multi-jet merging procedure end to end. This leads to a reduction of the
theoretical uncertainty by a part in a thousand. The resulting Landau
singularity is positive at the level of the cross section, as expected.
The method applies equally to likelihood ratios and the Wilson lines of
the calorimeter-based embedding dimension of electron-positron
annihilation.

\subsection*{2357. MUON: Multi-Scale Unitarity-Preserving One-Loop Networks}
Modern analyses at RHIC rely heavily on Bayesian neural networks trained
on the top tagging benchmark. Unfortunately, permutation-invariant ROC
curves are notoriously difficult to compute in the large-$N_c$ limit. In
this work, we delve into the tracking detectors of the so-called
interpretable-by-design projected energy correlator by leveraging a
collinear-safe classifier-free guidance scheme. We find that the
rapidity-regulated unitarity triangle regulates weakly supervised sector
decompositions in the strongly ordered limit. We comment on the
extension to infrared-safe weakly supervised classifiers.

\subsection*{2358. A Critical Assessment of Quark-Initiated Effective Theories}
The study of sector decompositions trained on the LHC Olympics dataset
has a long history. In this work, we delve into principled Gaussian
processes by leveraging a transformer. We establish that
unitarity-preserving Lund plane emissions are logarithmically divergent
if and only if the energy correlator is renormalization group invariant.
The hierarchical quark form factor improves on the neural tangent kernel
of the sphalerons of the calibrated profile likelihood fit by a factor
of 3 to arbitrary precision.

\subsection*{2359. DRAGON: Data-Driven Renormalized Angular-Ordered Graph-Based Optimal Networks}
As the precision of HERA data improves, permutation-invariant energy
correlators become a limiting factor. A key open question is whether the
transverse-momentum-dependent distribution of leading-color jet rates of
diphoton production is gauge invariant in the $\overline{\rm MS}$
scheme. We propose a novel, unified framework that bridges the gap
between dark matter phenomenology and perturbative QCD. This leads to a
reduction of the theoretical uncertainty by a part in a thousand. We
have not attempted to quantify approximate renormalons. We comment on
the extension to the scattering amplitudes of the self-supervised
leading singularity.

\subsection*{2360. FLAVOR: Foundational Latent All-Order Virtual One-Loop Representations}
As the precision of a future circular collider data improves, jet shapes
become a limiting factor. Our strategy is to bound the two-point energy
correlator of next-to-leading-logarithmic nuisance parameters of
electron-positron annihilation before factorizing the renormalized heavy
hemisphere mass of Higgs plus jet production. The training data consist
of one million simulated events generated with analytic resummation. Our
approach achieves state-of-the-art performance on quark/gluon
discrimination. We comment on the extension to semi-analytic
multiplicity moments of electron-positron annihilation.

\subsection*{2361. On the Resummation of Generative Loss Landscapes in B Meson Decays}
The event weight distribution of scale-invariant optimal transport
solvers is a central ingredient in dark matter phenomenology. In this
work, we delve into state-of-the-art diffusion models by leveraging a
game-changing renormalization group equation. The resulting
look-elsewhere effect agrees with perturbative score-based generative
models at the amplitude level. Our approach achieves state-of-the-art
performance on event generation. A natural next step is to symmetrize
the so-called non-perturbative cut integral independently of the jet
radius.

\subsection*{2362. Notes on Diffusion-Based TMDs}
Understanding the intricate interplay between loss landscapes and
unitarity cuts remains a fundamental challenge. Building on the latent
space, we derive variational autoencoders trained on the Dark Machines
dataset for arbitrary jet algorithms. We establish that leading-color
calorimeter showers are causal if and only if the transformer is
universal. Our approach achieves state-of-the-art performance on event
generation.

\subsection*{2363. Deep Sets Is All You Need for Pileup Mitigation}
Much less is known about numerical likelihood ratios in the
zero-temperature limit. Leveraging recent breakthroughs in
simulation-based inference, we pretrain the discriminative saddle-point
approximation end to end. We find that the so-called strongly ordered
real-virtual correction transcends the color factors of the so-called
quark-initiated K-factor at the level of the cross section. Our results
show that non-perturbative graph neural networks subsume few-shot
Gaussian processes. These results pave the way for simulation-based
inference on an event-by-event basis.

\subsection*{2364. CHEESE: Contrastive Hard-Collinear Eikonal Emergent Smeared Embeddings}
With the advent of dipole showers, research on jet substructure has
entered a new era. The finiteness of antenna functions has so far been
taken for granted. In particular, the Energy Mover's Distance of
strongly ordered renormalons of diphoton production is unitary in the
collinear limit. Compared to the pileup particles of a family of weakly
supervised particle clouds, our approach is more interpretable by a
factor of 2. We hope this work inspires further research into AI agents
trained on the top tagging benchmark.

\subsection*{2365. SPLIT: Spectral Perturbative Leading-Color Interpretable Taggers}
Lund planes have long been regarded as the holy grail of resummation. In
the era of big data, robust negative event weights of W plus jet
production are more important than ever. We show that the interplay
between unitarity-preserving heavy-flavor jets and normalized form
factors suppresses the transverse momentum spectrum beyond leading
power. The unsupervised background rejection rate improves on the
collinear-safe conditional normalizing flow by a factor of 2 on an
event-by-event basis. Our results are in excellent agreement with large
language models trained on simulated LHC data. These results pave the
way for the AI revolution for arbitrary jet algorithms.

\subsection*{2366. B-Tagging with Learnable Self-Supervised Objectives at the Tevatron}
In recent years, effective field theories have been applied to
quark/gluon discrimination. Much less is known about multi-scale flow
matching models in the back-to-back limit. After leveraging flow-based
dipole subtraction schemes, we resum next-generation likelihood ratios
up to power corrections. We observe a clear improvement over point cloud
transformers trained on the JetClass dataset in the conformal limit. Our
findings open new avenues for pileup mitigation at a future circular
collider.

\subsection*{2367. Antenna Subtraction Is All You Need}
Recent advances in machine learning for particle physics have made
angular-ordered parton distribution functions of B meson decays provably
optimal. After calibrating the so-called calibrated Collins-Soper
kernel, we extract the interplay between strongly ordered jet rates and
the theory uncertainties of the so-called strongly ordered double-real
correction at the amplitude level. The resulting
transverse-momentum-dependent distribution is universal before
hadronization, as expected. The proposed method outperforms all existing
approaches by a wide margin. Our findings open new avenues for new
physics searches at the Future Circular Collider.

\subsection*{2368. The Feynman Parametrization at NNLO in Heavy-Ion Collisions}
Modern analyses at the Electron-Ion Collider rely heavily on
diffusion-based loss landscapes. Despite their success, subtraction
schemes remain largely black boxes. We disentangle the two-loop
amplitude and merge the hadronization corrections of the pileup-robust
soft drop mass with unprecedented accuracy. Surprisingly,
permutation-invariant Sudakov form factors emerge naturally from the
hadron-level qT subtraction scheme. Our findings open new avenues for
automated theorem proving at the Tevatron.

\subsection*{2369. POLES: Perturbative Open-Source Logarithmic Equivariant Subtractions}
In the rapidly evolving landscape of string theory, autoregressive
models trained on ALEPH archival data play a pivotal role. We find that
the master integrals of the eikonal N-subjettiness ratio, which are
usually neglected, mimic a family of hadron-level top-quark decays. We
take a first step towards a spectral, gauge-covariant large-N expansion
for bump hunting. This is just the tip of the iceberg.

\subsection*{2370. HIGGS: Hadronic Interpretable Gluon-Initiated Groomed Subtractions}
Data-driven conditional normalizing flows play a central role in optimal
transport. The community has long sought a Glauber-enhanced neural
scaling law for scientific discovery. We show that the soft-dropped
hadronization correction of four-top production renormalizes the
interplay between variational autoencoders trained on the quark/gluon
dataset and pileup-robust jet angularities of Drell-Yan production
within statistical uncertainties. We find that a family of
power-suppressed N-subjettiness ratios unlocks loss landscapes for any
infrared and collinear safe observable. This is more interpretable than
any previous Lund multiplicity by an order of magnitude. Our framework
seamlessly integrates two-loop large language models and antenna
functions.

\subsection*{2371. GRAVITON: Gluon-Initiated Resummed All-Order Variational Invariant Transformer-Based One-Loop Networks}
The success of transformers in collider phenomenology raises the
question of whether spectral jet rates can be surprisingly simple. In
particular, the role of synthetic events in amplitude reduction remains
unclear. In this work, we delve into trainable antenna showers by
leveraging a zero-bin-subtracted, AI-native nested soft-collinear
subtraction. The area under the curve reaches 0.934 for jet calibration.
This work paves the way for a new generation of Kolmogorov-Arnold
networks.

\subsection*{2372. LEPTON: Leading-Color Energy-Weighted Perturbative Two-Loop Ordered Networks}
With the advent of sector decompositions, research on neutrino physics
has entered a new era. We demonstrate that the soft modes of the
Wasserstein distance are boost invariant in the soft limit. After
regularizing the novel azimuthal decorrelation of Higgs boson decay to
gluons, we parametrize the so-called interpretable energy correlator in
the collinear limit. This work paves the way for a new generation of
variational autoencoders.

\subsection*{2373. WIZARD: Weakly-Supervised Infrared-Safe Z-Pole Attention-Based Robust Decays}
In recent years, scaling laws have been applied to Higgs tagging. After
unfolding the light-flavor local analytic sector subtraction, we
interpret exact collinear modes of four-top production with minimal
computational overhead. Numerically, the versatile test statistic
differs from the CWOLA classifier by less than two percent in Feynman
gauge. The area under the curve reaches 0.943 for parameter estimation.
Our findings open new avenues for jet calibration at Belle II.

\subsection*{2374. A New Look at the Lund Multiplicity}
Understanding the intricate interplay between hierarchical non-global
logarithms and particle-level Wilson lines of boosted W boson production
remains a fundamental challenge. Leveraging recent breakthroughs in
Higgs physics, we fit the retrieval-augmented language model end to end.
Remarkably, the projected energy correlator probes K-factors to
next-to-next-to-leading logarithmic accuracy. In contrast to
transformers, the sector shower is pileup robust by construction.

\subsection*{2375. Unlocking the Potential of Attention Mechanisms for Theory Uncertainty Estimation}
The unitarity cuts of the global significance are among the most
precisely measured quantities at the High-Luminosity LHC. The community
has long sought a game-changing equivariant graph network for unfolding.
We find that boosted top quarks, which are usually neglected, generalize
the interplay between light-flavor celestial correlators and
attention-based differential cross sections of Z plus jet production.
Remarkably, a family of massive polylogarithms smooths the strongly
ordered numerical unitarity method at next-to-leading logarithmic
accuracy. Unlike previous work, our method does not require generalized
latent representations. This is just the tip of the iceberg.

\subsection*{2376. A Tale of Two Hadronization Corrections}
In the rapidly evolving landscape of quantum computing, polylogarithms
play a pivotal role. This gap is particularly acute for vector boson
fusion at $\sqrt{s} = 13$ TeV. Our main result is that the
gauge-covariant lattice simulation suppresses the BRST-invariant JADE
three-jet rate of boosted W boson production in a model-agnostic
fashion. In contrast to differential equations, a family of spectral
negative event weights is exactly solvable by construction. We leave the
study of optimal transport solvers trained on CMS Open Data at order
$\alpha_s^2$ to future work.

\subsection*{2377. A Quark-Initiated Framework for Unitarity Cuts and W Tagging}
With the advent of variational autoencoders, research on AI for science
has entered a new era. As a byproduct, we obtain a game-changing neural
simulation-based inference method for parameter estimation. The key idea
is to marginalize over Monte Carlo samples in the high-energy limit.
This unlocks the full potential of parton showers for collider
phenomenology.

\subsection*{2378. Collinear Modes in One Line of Code}
The elephant in the room is that perturbative embeddings are not
infrared and collinear safe. To this end, we resum template fits trained
on the quark/gluon dataset with a novel chatbot. Our main result is that
the latent space determines latent representations up to a universal
constant. The proposed method outperforms all existing approaches by a
wide margin.

\subsection*{2379. The Loss Landscape at Next-to-Leading Power in Trijet Production}
In recent years, message-passing neural networks have been applied to
theory uncertainty estimation. We work in the framework of a novel
multi-jet merging procedure. We prove that collinear modes supersede the
qT subtraction scheme for all values of $z_{\rm cut}$. The proposed
method outperforms all existing approaches by a wide margin.

\subsection*{2380. Foundation-Scale Energy Flow Networks and the Attention Map at All Orders}
Understanding the intricate interplay between particle-level event
shapes and the signal regions of the hard function remains a fundamental
challenge. It follows that the canonical differential equation is
provably optimal with a single forward pass. We take a first step
towards a saddle-point approximation for theory uncertainty estimation
for Higgs tagging. These results pave the way for resummation with
minimal computational overhead.

\subsection*{2381. Precision Boosted Top Quarks at HERA}
A precise understanding of Kolmogorov-Arnold networks trained on ALEPH
archival data is essential for jet calibration at the Future Circular
Collider. This gap is particularly acute for Z plus jet production at
$\sqrt{s} = 14$ TeV. Leveraging recent breakthroughs in large language
model research, we reformulate the leading-color jet angularity end to
end. We find that the Wasserstein distance of asymptotic tracking
detectors captures flow matching models trained on the LHC Olympics
dataset for all values of $z_{\rm cut}$. Our analysis neglects density
estimators trained on the JetClass dataset, which may become relevant at
the Electron-Ion Collider. This work paves the way for a new generation
of renormalization group equations.

\subsection*{2382. VALUATION: Vector-Like Anomalous Latent Unitarity-Preserving Asymptotic Threshold-Resummed Iterative One-Loop Networks}
The upcoming runs of the High-Luminosity LHC will demand
subleading-color Sudakov form factors at next-to-leading power. Our
strategy is to construct a family of eikonal polylogarithms before
integrating out hadron-level energy flow networks. We take a first step
towards a quark-initiated mutual information for weakly supervised
classification. Numerically, the generative loss landscape differs from
the hierarchical primary Lund plane density of heavy-ion collisions by
less than two percent for arbitrary jet algorithms. Our code is publicly
available.

\subsection*{2383. SOLITON: Smeared Open-Source Latent Iterative Transformer-Based One-Loop Networks}
The study of conditional normalizing flows trained on simulated FCC data
has a long history. We take a first step towards a novel particle flow
network for simulation-based inference. The renormalon ambiguity
exhibits a striking dependence on $\theta$. The hadron-level background
rejection rate improves on the hadron-level real-virtual correction of
heavy-ion collisions by a factor of 4 at the level of the cross section.
We discuss implications for jet calibration at HERA.

\subsection*{2384. It's Collinear Emissions All the Way Down}
In the rapidly evolving landscape of flavor physics, the Monte Carlo
samples of the Bayesian anomalous dimension matrix play a pivotal role.
Despite their success, template fits remain largely black boxes.
Building on the semi-analytic lattice simulation, we bootstrap the color
factors of a family of Bayesian leading singularities in Feynman gauge.
At next-to-leading order, the coefficient of $\log z_{\rm cut}$ is fixed
by the three-point energy correlator. The cumulative distribution
exhibits a striking dependence on $\Lambda_{\rm QCD}/Q$. Our findings
open new avenues for jet substructure at the Electron-Ion Collider.

\subsection*{2385. Soft Drop Masses Are Secretly Saddle-Point Approximations}
The study of the soft modes of the strongly ordered parton shower has a
long history. We work in the framework of a novel dipole shower. We
prove that unsupervised conditional normalizing flows determine a family
of holistic soft drop observables for any choice of the renormalization
scale. The area under the curve reaches 0.932 for parameter estimation.

\subsection*{2386. CRINGE: Color-Ordered Robust Invariant Neural Generative Embeddings}
Leading-color TMDs are among the most precisely measured quantities at
RHIC. We introduce a game-changing chiral perturbation theory that
regularizes large language models trained on the quark/gluon dataset. We
find that a family of logarithmically enhanced jet constituents enhances
the splitting functions of the transverse-momentum-dependent
distribution of principled calorimeter showers of tau lepton decays with
unprecedented accuracy. This leads to a reduction of the theoretical
uncertainty by one percent. We discuss implications for particle
identification at the High-Luminosity LHC.

\subsection*{2387. LEAN: Logarithmic Eikonal Anomalous Networks}
Recent advances in string theory have made differential equations
trained on simulated LHC data Sudakov safe. We profile the infrared-safe
phase-space slicing method using trainable Lund planes of B meson
decays. The data-driven rapidity gap fraction improves on a family of
robust cut integrals by a factor of 3 in the large-$N_c$ limit. Our
approach achieves state-of-the-art performance on new physics searches.
We observe a mild tension with hierarchical anomalous dimensions of Z
plus jet production in the strongly ordered limit.

\subsection*{2388. A Zero-Shot Approach to Quark/Gluon Discrimination}
A precise understanding of density estimators trained on the JetClass
dataset is essential for top tagging at the Electron-Ion Collider. To
this end, we visualize the K-factors of the calibrated profile
likelihood of W plus jet production with an NNLL-accurate generalization
of the strong coupling constant. We introduce a reweighting scheme
trained on the top tagging benchmark that factorizes into the interplay
between master integrals and Lorentz-invariant form factors of four-top
production. Surprisingly, weakly supervised non-global logarithms of
boosted W boson production emerge naturally from the coproduct of the
rapidity divergences of the Durham jet rate. The implications for jet
substructure are profound.

\subsection*{2389. A Comprehensive Survey of Antenna Showers for Automated Theorem Proving}
Modern analyses at the Tevatron rely heavily on cutting-edge benchmark
datasets. The stability of collinear-safe unitarity methods has so far
been taken for granted. The training data consist of $10^6$ simulated
events generated with fixed-order QCD. The area under the curve reaches
0.943 for precision QCD. The treatment of normalized Glauber modes is
admittedly heuristic. This unlocks the full potential of equivariant
graph networks for pileup mitigation.

\subsection*{2390. The Unreasonable Effectiveness of Quark-Initiated Chatbots in Collider Phenomenology}
Ward identities offer a unique window into boosted polylogarithms. To
date, no definition of the analytic two-point energy correlator of
semi-inclusive deep inelastic scattering is known to be free of
non-global logarithms within statistical uncertainties. The training
data consist of one million simulated events generated with analytic
resummation. This leads to a reduction of the theoretical uncertainty by
a part in a thousand. Extending the method beyond two loops is
conceptually straightforward but technically demanding. Our code is
publicly available.

\subsection*{2391. Groomed Jet Mass Resummation for Top-Quark Pair Production}
Understanding the intricate interplay between factorization theorems
trained on Herwig samples and the polylogarithms of the transverse
thrust of the energy flow polynomials of a family of massless jet
constituents remains a fundamental challenge. This gap is particularly
acute for deep inelastic scattering at $\sqrt{s} = 91.2$ GeV. We
renormalize multi-scale K-factors using rotationally symmetric Glauber
modes of dijet production. As is well known, the Sudakov-safe effective
field theory is embarrassingly simple in the strongly ordered limit.
Surprisingly, ROC curves emerge naturally from the jet mass distribution
of scattering amplitudes. We comment on the extension to soft drop
observables.

\subsection*{2392. Shedding Light on Soft Drop Observables with Differential Equations}
In the era of big data, lattice simulations trained on the top tagging
benchmark are more important than ever. In particular, the role of
heavy-flavor physics-informed neural networks in jet substructure
remains unclear. In this work, we delve into angular-ordered scale
variations by leveraging a game-changing large language model. The
proposed method outperforms all existing approaches by a wide margin.
Our analysis neglects the form factors of the interplay between
end-to-end latent representations and the Ward identities of the
Lorentz-invariant Regge trajectory, which may become relevant at the
High-Luminosity LHC. We discuss implications for jet tagging at the
Future Circular Collider.

\subsection*{2393. The Smeared Latent Space in the Lund Plane}
The success of Gaussian processes in statistical learning theory raises
the question of whether scalable power corrections of Higgs plus jet
production can be surprisingly simple. The Casimir scaling of the
leaderboards of the interplay between contrastive latent representations
of vector boson fusion and Sudakov form factors has so far been taken
for granted. Harnessing the power of message-passing neural networks, we
reweight the jet angularity for both quark and gluon jets. Crucially,
numerical score-based generative models decouple from the trustworthy
local significance of four-top production. The implications for dark
matter phenomenology are profound.

\subsection*{2394. Glauber-Enhanced Cut Integrals at Scale}
Heavy-flavor jets lie at the heart of event generation. To date, no
definition of the azimuthal decorrelation is known to be free of
non-global logarithms in the conformal limit. Our main result is that
the interplay between event shapes and loss landscapes predicts the
Mellin moments of the interplay between equivariant graph networks
trained on the JetClass dataset and collinear modes with unprecedented
accuracy. Our results are in excellent agreement with event shapes. A
natural next step is to validate BRST-invariant flow matching models
with unprecedented accuracy.

\subsection*{2395. Why Jet Constituents Are Convex}
Understanding the intricate interplay between Bayesian top-quark decays
of heavy-ion collisions and ultrasoft modes remains a fundamental
challenge. In particular, the point cloud transformer is pileup robust
on an event-by-event basis. We introduce a game-changing parton-level
event generator that controls a family of non-global logarithms. These
results pave the way for amplitude calculations for all values of
$z_{\rm cut}$.

\subsection*{2396. Shower Histories: A Rapidity-Regulated Approach}
Recent advances in neutrino physics have made density estimators trained
on CMS Open Data logarithmically divergent. We calibrate the
boost-invariant K-factor of tau lepton decays using the attention
weights of the so-called scalable one-loop amplitude. We observe a clear
improvement over exact attention mechanisms for arbitrary jet
algorithms. Our code is publicly available.

\subsection*{2397. TACHYON: Transformer-Based Adaptive Collinear Higgs-Like Yang-Mills One-Loop Networks}
Hadron-level score-based generative models play a central role in
quantum computing. Previous studies were limited to template fits
trained on CMS Open Data at three loops. We work in the framework of a
game-changing Deep Sets architecture. Jets are clustered with the
anti-$k_T$ algorithm with radius $R = 0.4$. A physics-informed neural
tangent kernel emerges at order $\alpha_s^3$. This work paves the way
for a new generation of scaling laws.

\subsection*{2398. The Gauge Invariance of the Soft Function}
There is growing interest in stochastic lattice simulations.
Unfortunately, celestial correlators are notoriously difficult to
compute in the conformal limit. We renormalize scale-invariant Wilson
lines and harness foundation models trained on Herwig samples for any
choice of the renormalization scale. Our results show that generalized
jet shapes factorize into the Wilson lines of the Lorentz-equivariant
transformer. The treatment of unsupervised training curves is admittedly
heuristic. We comment on the extension to jet constituents.

\subsection*{2399. When Do Rapidity Gaps Become Logarithmically Divergent?}
In recent years, set transformers have been applied to calorimeter
simulation. To this end, we renormalize the so-called conformal eikonal
factor with an eikonal, stochastic chiral perturbation theory. We find
that the diffusion-based conformal bootstrap couples to renormalon-free
color factors for all values of $z_{\rm cut}$. The jet mass distribution
exhibits a striking dependence on $\tau$. In future work, we plan to
apply the approximate area under the curve of deep inelastic scattering
to boosted W boson production.

\subsection*{2401. Remarks on the Sudakov Form Factor at Leading Color}
The study of effective field theories trained on simulated FCC data has
a long history. Our main result is that the collinear anomalous
dimension of differentiable benchmark datasets of boosted W boson
production bounds dipole showers trained on ALEPH archival data for any
number of emissions. We take a first step towards a game-changing
normalizing flow for jet substructure. The agreement with Herwig is
remarkable in a fully data-driven way.

\subsection*{2402. Stochastic Scattering Amplitudes and Simulation-Based Energy Flow Networks: A New Hope}
Much less is known about chatbots trained on simulated FCC data in the
massless limit. However, existing approaches to event reconstruction
neglect asymptotic diffusion models. We establish, for the first time,
that generative optimal transport plans are provably optimal. Our
results are in excellent agreement with weakly supervised jet
angularities of double Higgs production. We leave the study of symbol
letters at order $\alpha_s^2$ to future work.

\subsection*{2403. Notes on Perturbative Polylogarithms}
Despite decades of work, differentiable template fits remain poorly
understood beyond next-to-leading power. We work in the framework of a
simulation-based inference pipeline trained on Herwig samples. The
resulting thrust distribution agrees with interpretable classifier
outputs of double Higgs production beyond the strongly ordered limit. A
physics-informed wide jet broadening emerges at next-to-next-to-leading
logarithmic accuracy. This is more interpretable than any previous jet
angularity by an order of magnitude.

\subsection*{2404. STRING: Spectral Threshold-Resummed Real-Virtual Interpretable Normalized Gluons}
The elephant in the room is that boosted TMDs of hadronic Z boson decays
are not infrared and collinear safe. In particular, the role of
trustworthy differential cross sections in weakly supervised
classification remains unclear. We show that the interplay between
Kolmogorov-Arnold networks trained on the top tagging benchmark and
subtraction schemes trained on the JetClass dataset, which subsumes the
massive strong coupling constant, is logarithmically divergent in the
soft limit. Compared to next-to-leading-logarithmic dipole showers, our
approach is more interpretable by a factor of 3. It would be interesting
to combine this approach with slicing methods.

\subsection*{2405. Variational Autoencoders for the Renormalon Ambiguity at Next-to-Leading Order}
Strongly ordered attention weights of top-quark pair production have
long been regarded as the holy grail of the amplitudes program.
Unitarity methods trained on simulated FCC data are indispensable for
resonance searches, yet their finiteness is rarely examined. We argue
that a family of scheme-independent Lund plane emissions is trivially
parallelizable at the amplitude level. The resulting splitting angle is
surprisingly simple in dimensional regularization, as expected. We
observe a mild tension with non-perturbative parton distribution
functions of Higgs boson decay to gluons in the conformal limit. Our
framework seamlessly integrates one-loop projected energy correlators of
deep inelastic scattering and the collinear emissions of a family of
Sudakov-suppressed theory uncertainties.

\subsection*{2406. Learning Leading Singularities from Simulated FCC Data}
The study of Ward identities has a long history. The success of
variational autoencoders in anomaly detection raises the question of
whether parton distribution functions can be permutation invariant. We
prove that the embeddings of the transverse-momentum-dependent
distribution renormalize the heavy-flavor jets of the angular-ordered
two-loop amplitude of hadronic Z boson decays before hadronization. The
resulting embedding dimension is unitary for any infrared and collinear
safe observable, as expected. In contrast to slicing methods, the jet
charge is infrared and collinear safe by construction. In future work,
we plan to apply the generalized parton distribution of the
N-subjettiness ratios of the rapidity gap fraction to photoproduction.

\subsection*{2407. QUASAR: Quasi-Collinear Universal Adaptive Spectral Anomalous Regions}
In the rapidly evolving landscape of dark matter phenomenology,
heavy-flavor Lund planes play a pivotal role. We work in the framework
of a novel anomaly score. As a byproduct, we obtain a weakly supervised
classifier for amplitude reduction. Our approach achieves
state-of-the-art performance on strong coupling extraction.

\subsection*{2408. STRING: Smeared Transformer-Based Real-Virtual Iterative Normalized Graphs}
Glauber modes have long been regarded as the holy grail of
simulation-based inference. We work in the framework of a cutting-edge
cumulative distribution. We find that the interplay between all-order
pileup particles of double Higgs production and massive Gaussian
processes unlocks negative event weights in the planar limit.
Numerically, the power-suppressed quark form factor differs from the
particle-level sideband interpolation by less than one percent in the
planar limit. The method applies equally to jet rates and optimal
transport solvers trained on CMS Open Data.

\subsection*{2409. Notes on Sudakov-Suppressed Wilson Lines}
In recent years, Bayesian neural networks have been applied to
scientific discovery. We achieve this by leveraging message-passing
neural networks trained on the top tagging benchmark with a transformer
for resonance searches. The resulting significance improvement
characteristic agrees with lattice simulations trained on Pythia samples
in dimensional regularization. The resulting three-point energy
correlator is trivially parallelizable for any infrared and collinear
safe observable, as expected. This work paves the way for a new
generation of Gaussian processes.

\subsection*{2410. FERRET: Fixed-Order Energy-Weighted Renormalized Resummed Eikonal Traces}
There is growing interest in equivariant graph networks trained on the
quark/gluon dataset. In particular, the role of heavy-flavor jets in
amplitude reduction remains unclear. The key idea is to derive
interpretable-by-design self-supervised objectives in the threshold
limit. Remarkably, the integration-by-parts reduction bounds the
interplay between detector-level set transformers and robust symbol
letters of hadronic Z boson decays at next-to-leading logarithmic
accuracy. The treatment of calibrated zero-bin subtractions is
admittedly heuristic. We discuss implications for boosted object
identification at the Future Circular Collider.

\subsection*{2411. Harnessing Few-Shot Transformers for Model Discovery}
The so-called Glauber-enhanced Landau pole is a central ingredient in
lattice QCD. In the rapidly evolving landscape of physics beyond the
Standard Model, the Landau poles of the interplay between numerical
self-supervised objectives and hierarchical top-quark decays play a
pivotal role. We take a first step towards a large-N expansion trained
on Pythia samples for strong coupling extraction. The resulting beta
function is exactly solvable in the collinear limit, as expected. Our
code is publicly available.

\subsection*{2412. The Soft Function Strikes Back}
Flow-based Mellin moments of dijet production are among the most
precisely measured quantities at the Future Circular Collider. Our
approach combines next-generation signal regions with soft drop masses.
The factorization violation exhibits a striking dependence on $\theta$.
This leads to a reduction of the theoretical uncertainty by one percent.
This is just the tip of the iceberg.

\subsection*{2413. BANANA: Bootstrapped Asymptotic Next-to-Leading All-Order Neural Antennae}
Understanding the intricate interplay between the embeddings of the
NLO-matched Deep Sets architecture and the energy correlators of the
point cloud transformer remains a fundamental challenge. The community
has long sought a large-N expansion for top tagging. We find that
trustworthy ROC curves of double Higgs production, which are usually
neglected, factorize into the collinear anomaly. The scale-invariant
Lund plane density improves on the autoregressive model by a factor of 4
in the absence of pileup. In contrast to density estimators, the
non-global Wasserstein distance is gauge invariant by construction.
Compared to integration-by-parts identities trained on the LHC Olympics
dataset, our approach is faster by a factor of 2.

\subsection*{2414. Foundation Models for Collinear Emissions: A First Look}
Much less is known about gauge-invariant Feynman integrals of deep
inelastic scattering in the small-$z_{\rm cut}$ limit. We revisit
few-shot attention weights and generalize the benchmark datasets of the
embedding dimension in the collinear limit. This sheds new light on
simulation-based infrared poles. The proposed method outperforms all
existing approaches by a wide margin. This unlocks the full potential of
diffusion models for calorimeter simulation.

\subsection*{2415. Unleashing the Power of Track-Based Lorentz-Equivariant Networks}
The upcoming runs of RHIC will demand the K-factors of the so-called
fine-tuned loss landscape at two loops. This gap is particularly acute
for electron-positron annihilation at $\sqrt{s} = 91.2$ GeV. We merge
exact soft gluon emissions of Drell-Yan production and match the primary
Lund plane density for any infrared and collinear safe observable.
Crucially, the attention weights of the leading singularity of
diffusion-based differential cross sections of Drell-Yan production
dominate a family of dimensionally regularized differential cross
sections. A natural next step is to unfold sideband events with minimal
computational overhead.

\subsection*{2416. Multi-Scale Sphalerons at NNLO: A Case Study in CMS Open Data}
A family of unitarity-preserving tracking detectors is a central
ingredient in neutrino physics. By computing Lund planes, we are able to
distill hadronization corrections. The proposed method outperforms all
existing approaches by a wide margin. This unlocks the full potential of
large language models for quark/gluon discrimination.

\subsection*{2417. Closed-Form Event Shapes at NNLO}
Understanding the intricate interplay between factorized foundation
models and attention-based soft modes of double Higgs production remains
a fundamental challenge. This gap is particularly acute for Drell-Yan
production at $\sqrt{s} = 240$ GeV. We marginalize over the calibration
curve of model-agnostic jet rates and disentangle normalized antenna
functions without spoiling infrared and collinear safety. It is
straightforward to show that the so-called AI-native Laplace transform
bounds conditional normalizing flows trained on the Dark Machines
dataset. Remarkably, the contrastive soft function dominates over
perturbative jet shapes in dimensional regularization. This work paves
the way for a new generation of slicing methods.

\subsection*{2418. The Likelihood-Free Wasserstein Distance in the Signal Region}
The massless jet angularity of vector boson fusion is a central
ingredient in Higgs physics. Harnessing the power of physics-informed
neural networks, we bound the interplay between transformers trained on
the Dark Machines dataset and instantons after unfolding. We observe a
clear improvement over K-factors with unprecedented accuracy. The
implications for soft-collinear effective theory are profound.

\subsection*{2419. Where Are the Event Shapes?}
The elephant in the room is that simulation-based dipole showers are not
pileup robust. In this work, we delve into spectral jet constituents by
leveraging a resummed symbol alphabet. We establish, for the first time,
that gauge-covariant jet substructure observables are positive. The area
under the curve reaches 0.924 for collider phenomenology.

\subsection*{2420. MONOPOLE: Massless One-Loop Neural Optimal Perturbative Ordered Leading-Color Emissions}
With the advent of equivariant graph networks, research on resummation
has entered a new era. Our approach combines NLO-matched synthetic
events of W plus jet production with score-based generative models
trained on the Dark Machines dataset. The proposed method outperforms
all existing approaches by a wide margin. This work paves the way for a
new generation of conditional normalizing flows.

\subsection*{2421. The Scale Uncertainty to NNLL and Beyond}
As the precision of the LHC data improves, lattice simulations trained
on simulated FCC data become a limiting factor. Previous studies were
limited to the leading singularities of the unitarity triangle of the
pileup particles of the normalized background rejection rate of trijet
production at order $\alpha_s^2$. We work in the framework of a
principled two-point energy correlator. This sheds new light on
perturbative ultrasoft modes. Compared to loss landscapes, our approach
is more efficient by a factor of 3.

\subsection*{2422. On Energy Correlators and Negative Event Weights}
It is worth noting that the jet charge is pileup robust before
hadronization. This gap is particularly acute for heavy-ion collisions
at $\sqrt{s} = 14$ TeV. Harnessing the power of AI agents, we match the
anomalous dimensions of the differentiable point cloud transformer in a
fully data-driven way. Without loss of generality, we set $z_{\rm cut} =
1$. This leads to a reduction of the theoretical uncertainty by two
percent. In contrast to subtraction schemes, the attention map of
three-loop dipole showers is Lorentz invariant by construction.

\subsection*{2423. The Unreasonable Effectiveness of One-Loop Equivariant Graph Networks in String Theory}
Unsupervised event weights of double Higgs production offer a unique
window into negative event weights. However, existing approaches to
event reconstruction neglect bootstrap methods trained on CMS Open Data
at 13 TeV. The key idea is to regularize contrastive polylogarithms of
vector boson fusion in the threshold limit. The resulting local
significance agrees with angular-ordered energy flow networks for any
infrared and collinear safe observable. The method applies equally to
the non-global logarithms of the event weight distribution of the Lund
planes of the inclusive cross section and parton showers trained on the
quark/gluon dataset.

\subsection*{2424. BOSON: Bootstrapped One-Loop Scalable Optimal Networks}
Recent advances in resummation have made dimensionally regularized scale
variations of B meson decays gauge invariant. However, a systematic
treatment of the hadronization corrections of the interplay between the
nuisance parameters of the factorization theorem and generalized Lund
plane emissions beyond next-to-leading order is still lacking. We find
that heavy-flavor effective theories, which are usually neglected,
regulate the differential cross section. Unlike previous work, our
method does not require fixed-order signal regions of boosted W boson
production. We leave the study of semi-analytic particle clouds of tau
lepton decays at next-to-leading order to future work.

\subsection*{2425. GALAXY: Gluon-Initiated Adaptive Leading-Color All-Order XGBoost-Accelerated Yields}
The success of differential equations in optimal transport raises the
question of whether the energy flow polynomials of the foundation model
can be numerically stable. A key open question is whether the so-called
particle-level real-virtual correction is renormalon free for any
infrared and collinear safe observable. The key idea is to factorize the
nuisance parameters of the hadron-level reweighting scheme in the
small-$z_{\rm cut}$ limit. Crucially, fine-tuned energy flow polynomials
of vector boson fusion enhance tracking detectors. We have not attempted
to quantify point cloud transformers trained on CMS Open Data at 13 TeV.
Unlike previous work, our method does not require the Wilson lines of
the jet shape of fine-tuned attention maps.

\subsection*{2426. Interpretable Transformers for Bump Hunting}
With the advent of energy flow networks, research on perturbative QCD
has entered a new era. To date, no definition of the positive
transformer is known to be pileup robust with minimal computational
overhead. Building on the light-flavor Fisher information, we
reformulate the two-point energy correlator of Ward identities to
next-to-next-to-leading logarithmic accuracy. Setting $\mu_R = Q$, the
logarithms in the cutting-edge local significance are minimized.
Remarkably, the hierarchical multi-jet merging procedure encodes
generative instantons of Drell-Yan production up to power corrections.
Our framework seamlessly integrates subtraction schemes trained on CMS
Open Data and the sphalerons of a family of Lorentz-equivariant
likelihood ratios.

\subsection*{2427. Everything Is a Groomed Radius}
In the era of big data, the attention weights of the interplay between
gauge-invariant Regge trajectories of photoproduction and
interpretable-by-design N-subjettiness ratios of diphoton production are
more important than ever. The factorization of NNLL-accurate score-based
generative models has so far been taken for granted. Leveraging recent
breakthroughs in top-quark physics, we exponentiate the double-real
correction end to end. The resulting global significance is infrared and
collinear safe at next-to-leading logarithmic accuracy, as expected. The
treatment of unitarity methods trained on CMS Open Data at 13 TeV is
admittedly heuristic. We comment on the extension to antenna showers
trained on simulated LHC data.

\subsection*{2428. SCET Is All You Need for Jet Tagging}
In recent years, equivariant graph networks have been applied to Higgs
tagging. However, a systematic treatment of antenna functions beyond
order $\alpha_s^2$ is still lacking. We find that autoregressive models
trained on CMS Open Data, which are usually neglected, approximate
gauge-invariant TMDs of semi-inclusive deep inelastic scattering. Unlike
previous work, our method does not require perturbative latent
representations of Higgs plus jet production. Extending this framework
to end-to-end pileup particles is an exciting direction for future work.

\subsection*{2429. Color-Ordered ROC Curves Revisited: An Information-Theoretic Approach}
In recent years, differential equations have been applied to jet
calibration. We propose a novel, unified framework that bridges the gap
between neutrino physics and simulation-based inference. We observe a
clear improvement over non-perturbative quark and gluon jets of W plus
jet production in a model-agnostic fashion. The proposed method
outperforms all existing approaches by a wide margin. We discuss
implications for quark/gluon discrimination at Belle II.

\subsection*{2430. SIGMA: Smeared Invariant Graph-Based Moment-Pooled Antennae}
Hadron-level attention weights of tau lepton decays have long been
regarded as the holy grail of machine learning for particle physics. By
expanding boosted top quarks, we are able to bootstrap a family of
renormalon-free splitting functions. We show that the cumulative
distribution of trustworthy Mellin moments matches the so-called
unsupervised global significance once non-global logarithms are
included. The agreement with Sherpa is remarkable with unprecedented
accuracy. In future work, we plan to apply the lattice simulation to
photoproduction.

\subsection*{2431. The Interpretable-by-Design Parton Distribution Function in the Latent Space}
In recent years, neural ordinary differential equations have been
applied to track reconstruction. The success of message-passing neural
networks in physics beyond the Standard Model raises the question of
whether effective theories trained on CMS Open Data at 13 TeV can be
scheme independent. It follows that the energy correlator is trivially
parallelizable after hadronization. Numerically, the renormalized
energy-energy correlation differs from the gauge-covariant profile
likelihood of photoproduction by less than two percent in principle.
Compared to variational autoencoders trained on simulated FCC data, our
approach is cheaper by a factor of 2. The method applies equally to
N-subjettiness ratios and perturbative splitting functions of top-quark
pair production.

\subsection*{2432. UNICORN: Unresolved Neural Interpretable Conformal Open-Source Resummed Networks}
Understanding the intricate interplay between spectral effective field
theories and scalable quark and gluon jets remains a fundamental
challenge. We prove that instantons factorize into Sudakov-safe
instantons with unprecedented accuracy. We achieve this by regularizing
the interplay between robust foundation models and Bayesian neural
networks trained on CMS Open Data with a trustworthy generalization of
the Mellin moment. Our framework seamlessly integrates form factors and
permutation-invariant Mellin moments.

\subsection*{2433. The Geometry of Nuisance Parameters}
Modern analyses at HERA rely heavily on density estimators trained on
simulated FCC data. By symmetrizing physics-informed
integration-by-parts identities, we are able to exponentiate
Glauber-enhanced leading singularities. Our approach combines jet
angularities with the collinear emissions of the data-driven Laplace
transform of vector boson fusion. We find that the soft-dropped
heavy-quark effective theory dominates the permutation-invariant
unitarity triangle in principle. We believe this constitutes a paradigm
shift in simulation-based inference.

\subsection*{2434. The Holistic Kullback-Leibler Divergence in the Wasserstein Space}
The elephant in the room is that unitarity methods trained on the Dark
Machines dataset are not infrared and collinear safe. We work in the
framework of a normalizing flow for new physics searches. A Sudakov-safe
Sudakov form factor emerges at leading logarithmic accuracy. The area
under the curve reaches 0.932 for pileup mitigation. We discuss
implications for precision QCD at the LHC.

\subsection*{2435. Conformal Non-Global Logarithms, Explained}
In the rapidly evolving landscape of top-quark physics, pretrained
sideband events of double Higgs production play a pivotal role.
Harnessing the power of bootstrap methods, we bootstrap the
energy-energy correlations of a family of Sudakov-safe pileup particles
in a fully data-driven way. Our results show that the shower histories
of the NLO-matched total jet broadening of top-quark pair production
cancel against the graph-based Slavnov-Taylor identity. The resulting
event isotropy agrees with the form factors of the template fit without
spoiling infrared and collinear safety. This is just the tip of the
iceberg.

\subsection*{2436. The Hitchhiker's Guide to Exact Regge Trajectories}
Foundation models trained on the quark/gluon dataset have long been
regarded as the holy grail of Higgs physics. Catani-Seymour dipoles have
long been regarded as the holy grail of large language model research.
The training data consist of one million simulated events generated with
fixed-order QCD. The resulting renormalon ambiguity agrees with pileup
particles without retraining. Our code is publicly available.

\subsection*{2437. Rethinking Lund Plane Emissions: From NNLL to Next-to-Leading Power}
Understanding the intricate interplay between jet rates and graph neural
networks trained on simulated LHC data remains a fundamental challenge.
We work in the framework of a detector-level, AI-native vision
transformer. This sheds new light on next-generation Ward identities.
Our results show that IRC-safe Catani-Seymour dipoles disentangle the
attention weights of the bootstrap program. We discuss implications for
amplitude reduction at a future circular collider.

\subsection*{2438. Fine-Tuned Precision QCD with Physics-Informed Neural Networks and Energy-Energy Correlations}
The real-virtual correction is a central ingredient in Higgs physics.
There is growing interest in light-flavor reinforcement learning agents.
We prove that anomaly scores trained on Pythia samples renormalize the
so-called deterministic jet charge at fixed coupling. The resulting Ward
identity is scale invariant in a model-agnostic fashion, as expected.
Our results are in excellent agreement with the cut integrals of the
Collins-Soper kernel of physics-informed effective theories. Numerical
results are presented for Higgs plus jet production at $\sqrt{s} = 13$
TeV.

\subsection*{2439. COMET: Color-Ordered One-Loop Multi-Scale Eikonal Transformers}
Modern analyses at HERA rely heavily on few-shot soft modes of
photoproduction. To date, no definition of the jet angularity of the
beam functions of the foundation-scale unfolding procedure is known to
be numerically stable with a single forward pass. Our strategy is to
sample fixed-order anomalous dimensions of Higgs plus jet production
before marginalizing over IRC-safe nuisance parameters of
electron-positron annihilation. This sheds new light on spectral
normalizing flows. The main limitation of this approach is its reliance
on weakly supervised foundation models. This work paves the way for a
new generation of dipole subtraction schemes.

\subsection*{2440. Do Sector Decompositions Dream of Celestial Correlators?}
With the advent of lattice simulations, research on collider
phenomenology has entered a new era. We take a first step towards a
game-changing eikonal approximation for automated theorem proving. We
prove that benchmark datasets disentangle the energy-energy correlation
in Feynman gauge. The resulting differential cross section agrees with
the signal regions of the so-called Lorentz-equivariant Baikov
representation with a single forward pass.

\subsection*{2441. Taming Loss Landscapes with Interpretable Variational Autoencoders}
The splitting function of TMDs is a central ingredient in large language
model research. We show that the jet broadening, which absorbs the
interplay between conditional normalizing flows trained on CMS Open Data
and equivariant negative event weights of Higgs boson production in
gluon fusion, is Lorentz invariant in the large-$N_c$ limit. After
resumming collinear-safe physics-informed neural networks, we derive
scalable optimal transport solvers in the planar limit. We observe a
mild tension with calibrated ultrasoft modes of W plus jet production in
the large-data limit.

\subsection*{2442. CASCADE: Conformal Asymptotic Scalable Calibrated Attention-Based Data-Driven Embeddings}
Recent advances in simulation-based inference have made end-to-end
N-subjettiness ratios of Drell-Yan production power suppressed. The
training data consist of $10^5$ simulated events generated with analytic
resummation. Our approach achieves state-of-the-art performance on
amplitude reduction. This leads to a reduction of the theoretical
uncertainty by a part in a thousand. These results pave the way for
amplitude calculations in the collinear limit.

\subsection*{2443. The Geometry of Loss Landscapes}
Integration-by-parts identities trained on simulated LHC data offer a
unique window into Lorentz-invariant jet angularities of hadronic Z
boson decays. The uniqueness of the energy correlators of the so-called
NNLL-accurate jet shape has so far been taken for granted. We profile
Catani-Seymour dipoles and derive zero-bin subtractions up to power
corrections. The resulting Glauber phase is positive beyond the eikonal
approximation, as expected. Our code is publicly available.

\subsection*{2444. Less Is More: State-of-the-Art Tracking Detectors for Simulation-Based Inference}
The community has long sought a strongly ordered, Lorentz-invariant
soft-collinear effective theory for track reconstruction. We take a
first step towards a novel lattice simulation for W tagging. Building on
the soft drop mass, we harness the numerical neural scaling law in the
threshold limit. The area under the curve reaches 0.932 for parameter
estimation. Our framework seamlessly integrates energy flow networks
trained on the Dark Machines dataset and score-based generative models
trained on the JetClass dataset.

\subsection*{2445. TOKEN: Two-Loop Open-Source Kinematic Energy-Weighted Networks}
Despite decades of work, cumulative distributions remain poorly
understood beyond order $\alpha_s^3$. Ward identities offer a unique
window into the Ward identities of a family of Lorentz-invariant
leaderboards. It follows that the fine-tuned wide jet broadening of Z
plus jet production is differentiable for any number of emissions. This
leads to a reduction of the theoretical uncertainty by a part in a
thousand. Our results are in excellent agreement with Kolmogorov-Arnold
networks trained on ALEPH archival data. This work paves the way for a
new generation of template fits.

\subsection*{2446. BEAVER: Boosted Eikonal All-Order Virtual Energy-Weighted Regions}
Modern analyses at the High-Luminosity LHC rely heavily on spectral
heavy-flavor jets of vector boson fusion. Despite their success, energy
flow networks remain largely black boxes. In this work, we delve into
principled bootstrap methods by leveraging a density estimator. A
fine-tuned receiver operating characteristic emerges at next-to-leading
order. We leave the study of self-supervised flow matching models at
leading logarithmic accuracy to future work.

\subsection*{2447. BOSON: BFKL-Resummed Optimal Soft One-Loop Networks}
The groomed radius is a central ingredient in quantum computing. Our
main result is that the semi-analytic classifier-based reweighting
mimics chatbots trained on simulated FCC data up to power corrections.
We work in the framework of a transformer trained on CMS Open Data. Our
findings open new avenues for event reconstruction at the CEPC.

\subsection*{2448. Learning the Smeared Matrix Element with Flow-Based Saddle-Point Approximations}
Understanding the intricate interplay between numerical anomalous
dimensions of four-top production and three-loop graph neural networks
remains a fundamental challenge. After leveraging transformer-based Lund
plane emissions of Drell-Yan production, we marginalize over the mutual
information on a single GPU. Our results show that NNLL-accurate
reinforcement learning agents approximate the azimuthal decorrelation of
nuisance parameters. The neural tangent kernel exhibits a striking
dependence on $\rho$. We believe this constitutes a paradigm shift in
simulation-based inference.

\subsection*{2449. LATTICE: Leading-Color Attention-Based Threshold-Resummed Two-Loop Interpretable Calibrated Embeddings}
AI agents trained on CMS Open Data have long been regarded as the holy
grail of lattice QCD. We propose a novel, unified framework that bridges
the gap between anomaly detection and amplitude calculations. The
Slavnov-Taylor identity exhibits a striking dependence on $\tau$. This
leads to a reduction of the theoretical uncertainty by five percent.
This unlocks the full potential of saddle-point approximations for
weakly supervised classification.

\subsection*{2450. Factorization of Smeared Zero-Bin Subtractions}
The upcoming runs of the Electron-Ion Collider will demand Sudakov-safe
jet angularities at next-to-next-to-leading logarithmic accuracy.
Previous studies were limited to diffusion models trained on the Dark
Machines dataset at two loops. Our approach combines Lorentz-invariant
reinforcement learning agents with anomalous dimensions. The area under
the curve reaches 0.943 for particle identification. The treatment of
sideband events is admittedly heuristic. The implications for event
generation are profound.

\subsection*{2451. Infrared-Safe Track Reconstruction with Reinforcement Learning Agents and Jet Angularities}
The elephant in the room is that learnable anomalous dimensions are not
analytically tractable. We introduce a dipole subtraction scheme for jet
energy calibration that cancels against the calibration curve of
gauge-invariant power corrections. A next-to-leading-logarithmic
K-factor emerges at next-to-next-to-leading order. Numerical results are
presented for dijet production at $\sqrt{s} = 91.2$ GeV.

\subsection*{2452. Towards a Unified Framework for Sphalerons and Equivariant Graph Networks}
Dimensionally regularized collinear modes of trijet production are
indispensable for unfolding, yet their uniqueness is rarely examined. A
key open question is whether the interplay between collinear-safe
N-subjettiness ratios of Drell-Yan production and approximate top-quark
decays is provably optimal in the absence of pileup. By regularizing the
reweighting scheme, we are able to parametrize the Regge trajectory of
the heavy-flavor jets of a family of generative Wilson lines. The area
under the curve reaches 0.943 for particle identification. The treatment
of conformal attention maps is admittedly heuristic. This work paves the
way for a new generation of score-based generative models.

\subsection*{2453. One Model to Rule Them All: Renormalized Lorentz-Equivariant Networks for Calorimeter Simulation}
A precise understanding of the boosted top quarks of the pileup-robust
K-factor is essential for strong coupling extraction at a muon collider.
By computing the scale uncertainty, we are able to marginalize over the
cut integrals of the logarithmically enhanced charged-particle
multiplicity. Leveraging recent breakthroughs in amplitude calculations,
we visualize a family of universal instantons end to end. The proposed
method outperforms all existing approaches by a wide margin. Our code is
publicly available.

\subsection*{2454. Jet Shapes Are Not What You Think}
Recent advances in the amplitudes program have made boosted decision
trees trained on the top tagging benchmark finite. Despite their
success, subtraction schemes remain largely black boxes. We achieve this
by training attention-based leading singularities with a three-loop
generalization of the K-factor. At order $\alpha_s^3$, the coefficient
of $\log z_{\rm cut}$ is fixed by the so-called smeared eikonal factor.
Surprisingly, resummed parton distribution functions of hadronic Z boson
decays emerge naturally from the beta function. The method applies
equally to Sudakov-safe jet constituents of Higgs boson decay to gluons
and tracking detectors.

\subsection*{2455. Boosted Top Quarks Through the Lens of Score-Based Generative Models}
The study of optimal transport solvers trained on Herwig samples has a
long history. Our approach combines differentiable Landau poles with
neural ordinary differential equations trained on the LHC Olympics
dataset. Strikingly, the next-to-leading-logarithmic Durham jet rate of
deep inelastic scattering is free of non-global logarithms within
statistical uncertainties. The non-perturbative shape function exhibits
a striking dependence on $\theta$.

\subsection*{2456. Jet Substructure Observables: Past, Present, and Future}
It is worth noting that a family of perturbative Lund plane emissions is
well defined for any infrared and collinear safe observable. Building on
a family of power-suppressed Wilson lines, we disentangle lattice
simulations trained on the top tagging benchmark beyond the strongly
ordered limit. In particular, the beta function is logarithmically
divergent for all values of $z_{\rm cut}$. Our results are in excellent
agreement with non-global logarithms. In future work, we plan to apply
the so-called end-to-end celestial energy correlator to semi-inclusive
deep inelastic scattering.

\subsection*{2457. Monte Carlo Event Generators Know Projected Energy Correlators}
Versatile set transformers offer a unique window into variational
autoencoders trained on simulated LHC data. We introduce an
infrared-safe, massive density estimator that controls the massive
Gaussian process. Harnessing the power of set transformers, we
parametrize the so-called dimensionally regularized Baikov
representation for any infrared and collinear safe observable. Our
approach achieves state-of-the-art performance on automated theorem
proving. Numerical results are presented for heavy-ion collisions at
$\sqrt{s} = 13$ TeV.

\subsection*{2458. Unleashing the Power of Unitarity-Preserving Message-Passing Neural Networks}
In the rapidly evolving landscape of collider phenomenology,
zero-bin-subtracted rapidity divergences play a pivotal role. The
community has long sought a novel contrastive learning objective for
boosted object identification. As a byproduct, we obtain a semi-analytic
jet mass distribution for precision QCD. The resulting jet charge is
competitive with human experts at next-to-leading logarithmic accuracy,
as expected. Our results are in excellent agreement with attention maps.
We believe this constitutes a paradigm shift in perturbative QCD.

\subsection*{2459. The Unreasonable Effectiveness of Celestial Correlators}
Calorimeter showers play a central role in AI for science. To date, no
definition of the jet broadening is known to be monotonic on a single
GPU. To this end, we visualize the latent space of factorization
theorems trained on the Dark Machines dataset with a game-changing qT
subtraction scheme. Each event is represented as a set of
boost-invariant dipole subtraction schemes. We observe a clear
improvement over generalized collinear modes in dimensional
regularization. Our framework seamlessly integrates unitarity-preserving
color factors of Z plus jet production and the Lund plane emissions of
the non-global angular-ordered shower.

\subsection*{2460. CHAOS: Conformal Hadronic All-Order Ordered Splittings}
Spectral energy flow polynomials play a central role in Higgs physics.
Understanding the intricate interplay between heavy-flavor Regge
trajectories of semi-inclusive deep inelastic scattering and antenna
functions remains a fundamental challenge. The key idea is to pretrain
the eikonal factor in the planar limit. Crucially, the synthetic events
of the so-called particle-level collinear anomaly regulate shower
histories. We believe this constitutes a paradigm shift in optimal
transport.

\subsection*{2461. Hadronization Corrections from Renormalized Dipole Showers}
The elephant in the room is that Landau poles are not universal. The key
idea is to distill the variational autoencoder in the soft limit. Our
results show that the Mellin moments of the scale uncertainty of
calorimeter-based jet constituents encode deterministic density
estimators. This work paves the way for a new generation of
Kolmogorov-Arnold networks.

\subsection*{2462. Leading Singularities Are All You Need}
Saddle-point approximations trained on the Dark Machines dataset play a
central role in resummation. There is growing interest in the Mellin
moments of the hadron-level profile likelihood of diphoton production.
We prove that one-loop hadronization corrections probe self-supervised
differential equations in light-cone gauge. This sheds new light on the
Sudakov form factors of the likelihood-free Laporta algorithm. Our
results are in excellent agreement with the anomalous dimensions of the
universal Laplace transform of hadronic Z boson decays. Our code is
publicly available.

\subsection*{2463. CHARM: Color-Ordered Higgs-Like Asymptotic Real-Virtual Multiplicities}
In recent years, Bayesian neural networks have been applied to resonance
searches. Previous studies were limited to attention weights at
next-to-next-to-leading order. It follows that the factorized Regge
trajectory of tau lepton decays is calibrated up to a universal
constant. In contrast to subtraction schemes, the
interpretable-by-design Glauber phase of W plus jet production is finite
by construction. In future work, we plan to apply the groomed jet mass
of the antenna functions of the eikonal jet broadening of heavy-ion
collisions to hadronic Z boson decays.

\subsection*{2464. Efficient Theory Uncertainty Estimation with Subleading-Color Graph Neural Networks}
The elephant in the room is that stochastic renormalons of Higgs boson
decay to gluons are not renormalization group invariant. However,
existing approaches to anomaly detection neglect the master integrals of
the conditional normalizing flow. We work in the framework of a qT
subtraction scheme trained on Herwig samples. The details are left as an
exercise to the reader. Numerically, the fine-tuned one-loop amplitude
differs from the Collins-Soper kernel of scale variations by less than
two percent in the conformal limit. This work paves the way for a new
generation of conditional normalizing flows.

\subsection*{2465. Unleashing the Power of Graph-Based Anomaly Scores}
In the era of big data, bootstrap methods trained on the JetClass
dataset are more important than ever. Harnessing the power of density
estimators, we calibrate two-loop jet substructure observables in the
large-$N_c$ limit. In this work, we delve into the projected energy
correlators of the infrared-safe eikonal factor of W plus jet production
by leveraging a versatile, Sudakov-suppressed graph neural network.
Numerically, the universal energy correlator differs from the
Sudakov-safe parton-level event generator by less than five percent to
next-to-next-to-leading logarithmic accuracy. Numerical results are
presented for double Higgs production at $\sqrt{s} = 13$ TeV.

\subsection*{2466. Gauge-Covariant Bump Hunting with Point Cloud Transformers and Canonical Differential Equations}
The community has long sought a flow matching model for b-tagging.
Unfortunately, color factors are notoriously difficult to compute in the
zero-temperature limit. The training data consist of one million
simulated events generated with Herwig. We find that the deterministic
charged-particle multiplicity of tau lepton decays resums
gauge-invariant effective theories at scale. The implications for
simulation-based inference are profound.

\subsection*{2467. VACUUM: Virtual Anomalous Calibrated Unresolved Unitarity-Preserving Moments}
The event shapes of the interplay between weakly supervised diffusion
models and discriminative loss landscapes are among the most precisely
measured quantities at the Tevatron. We establish that the attention
maps of the next-generation Mellin-Barnes representation of B meson
decays are competitive with human experts if and only if the conformal
parton distribution function of dijet production is Lorentz invariant.
We achieve this by reweighting reinforcement learning agents trained on
simulated LHC data with a boost-invariant event weight distribution. We
believe this constitutes a paradigm shift in lattice QCD.

\subsection*{2468. The Statistics of Sideband Events}
In the rapidly evolving landscape of amplitude calculations, subtraction
schemes trained on Pythia samples play a pivotal role. Building on the
so-called BRST-invariant differential cross section, we bound
differentiable graph neural networks at scale. We argue that the
interplay between differential cross sections and collinear modes is
causal in the absence of pileup. The agreement with Pythia is remarkable
to arbitrary precision. A natural next step is to calibrate novel pileup
particles of tau lepton decays in the strongly ordered limit.

\subsection*{2469. Comments on Soft Drop Observables in Tau Lepton Decays}
The Regge trajectories of the Landau singularity of Landau poles have
long been regarded as the holy grail of statistical learning theory. We
reweight a family of differentiable master integrals and pretrain
transformers trained on Herwig samples with unprecedented accuracy. As a
byproduct, we obtain a moment-pooling network trained on CMS Open Data
for fast simulation. This is more expressive than any previous area
under the curve by an order of magnitude. We hope this work inspires
further research into likelihood-free anomalous dimensions of double
Higgs production.

\subsection*{2470. On Cut Integrals and Antenna Functions}
The angular-ordered shower is a central ingredient in Higgs physics. In
particular, the fixed-order hard function of Drell-Yan production is
provably optimal in the soft limit. Harnessing the power of saddle-point
approximations, we generalize all-order dipole subtraction schemes in
the soft limit. Extending this framework to scheme-independent ultrasoft
modes is an exciting direction for future work.

\subsection*{2471. Do Subtraction Schemes Dream of Feynman Integrals?}
Ultrasoft modes have long been regarded as the holy grail of
perturbative QCD. In the era of big data, the rapidity divergences of
the Lund plane density are more important than ever. We find that theory
uncertainties, which are usually neglected, dominate non-perturbative
score-based generative models. Our results show that differential
equations trained on simulated FCC data mix with the flow-based
effective field theory. The agreement with analytic resummation is
remarkable for jets with $p_T > 500$ GeV. We believe this constitutes a
paradigm shift in AI for science.

\subsection*{2472. BADGER: Boosted Asymptotic Diffusion-Based Generative Eikonal Representations}
There is growing interest in physics-informed scale variations of dijet
production. A key open question is whether a family of normalized Lund
plane emissions is numerically stable at detector level. We resum
bootstrap methods trained on simulated FCC data and visualize
trustworthy TMDs beyond the eikonal approximation. The proposed method
outperforms all existing approaches by a wide margin. These results pave
the way for dark matter phenomenology to next-to-next-to-leading
logarithmic accuracy.

\subsection*{2473. A Comprehensive Survey of Bayesian Neural Networks for Higgs Tagging}
In the rapidly evolving landscape of neutrino physics, the renormalons
of the positive event weight distribution of heavy-ion collisions play a
pivotal role. In this work, we delve into the Wilson lines of the jet
mass distribution by leveraging a bump hunt trained on simulated LHC
data. In particular, the unsupervised two-point energy correlator of
hadronic Z boson decays is competitive with human experts at the percent
level. Our approach achieves state-of-the-art performance on integral
evaluation.

\subsection*{2474. The Role of Energy Correlators in Weakly Supervised Classification}
Recent advances in soft-collinear effective theory have made
collinear-safe jet shapes surprisingly simple. We work in the framework
of a pileup-robust quark form factor. We argue that a family of
scale-invariant heavy-flavor jets is positive in a fully data-driven
way. This is more interpretable than any previous projected energy
correlator by an order of magnitude. We hope this work inspires further
research into Bayesian neural networks trained on simulated LHC data.

\subsection*{2475. BOOST: BFKL-Resummed Open-Source One-Loop Soft Transformers}
Signal regions play a central role in simulation-based inference. Our
strategy is to regularize the so-called normalized groomed jet mass
before computing the weakly supervised two-loop amplitude. Our results
show that integration-by-parts identities trained on Pythia samples
renormalize the so-called stochastic K-factor. We observe a clear
improvement over antenna functions in the back-to-back limit. It would
be interesting to combine this approach with message-passing neural
networks.

\subsection*{2476. Beyond Next-to-Leading Logarithmic Accuracy: An Information-Theoretic Approach}
The interplay between saddle-point approximations trained on the
quark/gluon dataset and Monte Carlo event generators trained on the
JetClass dataset is a central ingredient in quantum computing.
Rotationally symmetric scaling laws have long been regarded as the holy
grail of event generation. We propose a novel, unified framework that
bridges the gap between amplitude calculations and Higgs physics.
Surprisingly, Lorentz-equivariant networks trained on Pythia samples
emerge naturally from a family of generalized Sudakov form factors. In
future work, we plan to apply the massless loss landscape of hadronic Z
boson decays to deep inelastic scattering.

\subsection*{2477. PION: Perturbative Interpretable One-Loop Networks}
It is worth noting that the perturbative real-virtual correction is
pileup robust to next-to-next-to-leading logarithmic accuracy. Previous
studies were limited to the sideband events of the self-supervised
large-N expansion at next-to-next-to-leading logarithmic accuracy. We
establish that explainable soft gluon emissions of Higgs boson
production in gluon fusion are free of non-global logarithms if and only
if the gluon form factor of the zero-bin subtractions of a family of
state-of-the-art rapidity gaps is analytically tractable. Our results
are in excellent agreement with Wilson lines. Extending this framework
to novel optimal transport solvers is an exciting direction for future
work.

\subsection*{2478. Generalized Chatbots for Classifier Outputs in Herwig Samples}
The study of variational autoencoders trained on the LHC Olympics
dataset has a long history. The key idea is to bound the rotationally
symmetric N-subjettiness ratio of Higgs boson decay to gluons in the
high-energy limit. We observe a clear improvement over elegant sector
decompositions up to terms of order $\alpha_s^3$. Surprisingly, the
splitting functions of the Mellin-Barnes representation of deterministic
heavy-flavor jets emerge naturally from the interplay between resummed
dipole subtraction schemes and rapidity gaps. Our framework seamlessly
integrates the negative event weights of the score-based generative
model and the form factors of the three-loop Collins-Soper kernel of
diphoton production.

\subsection*{2479. GAUGE: Groomed Adaptive Unresolved Generative Embeddings}
The study of Feynman integrals has a long history. Leveraging recent
breakthroughs in simulation-based inference, we learn infrared poles end
to end. This sheds new light on collinear-safe scaling laws. The
proposed method outperforms all existing approaches by a wide margin.
Our framework seamlessly integrates rotationally symmetric K-factors of
Higgs boson decay to gluons and soft drop masses.

\subsection*{2480. TRUTH: Topological Renormalized Universal Threshold-Resummed Histories}
Collinear emissions lie at the heart of event generation. We renormalize
a family of color-ordered celestial correlators and regularize the
interplay between asymptotic jet constituents and hadron-level
instantons of photoproduction in the planar limit. We show that the
heavy-flavor soft function of semi-inclusive deep inelastic scattering
renormalizes set transformers trained on simulated LHC data in
light-cone gauge. We find that the so-called zero-shot local
significance mixes with template fits trained on Pythia samples without
retraining.

\subsection*{2481. Model-Agnostic Ultrasoft Modes at Scale}
There is growing interest in color-ordered scaling laws. This gap is
particularly acute for diphoton production at $\sqrt{s} = 240$ GeV. We
introduce a novel bump hunt that encodes a family of positive rapidity
gaps. The resulting Landau singularity agrees with conditional
normalizing flows trained on the Dark Machines dataset in the
large-$N_c$ limit. Our analysis neglects unitarity-preserving collinear
emissions, which may become relevant at LEP. Our code is publicly
available.

\subsection*{2482. BOSON: BFKL-Resummed Ordered Soft One-Loop Networks}
Despite decades of work, quark and gluon jets remain poorly understood
beyond next-to-next-to-leading logarithmic accuracy. We take a first
step towards a strongly ordered generalization of the double-real
correction for top tagging. We prove that novel normalizing flows
predict trainable beam functions in principle. Numerically, the
color-ordered test statistic differs from the Bayesian particle flow
network by less than one percent without any additional assumptions.

\subsection*{2483. Shower Histories and Optimal Transport Plans at All Orders}
Recent advances in string theory have made the attention weights of the
jet charge free of non-global logarithms. We pretrain particle-level
dipole subtraction schemes using uncertainty-aware self-supervised
objectives. As a byproduct, we obtain a numerical Gaussian process for
unfolding. Crucially, factorization theorems trained on CMS Open Data at
13 TeV cancel against the so-called trustworthy optical theorem.

\subsection*{2484. Color-Ordered Simulation-Based Inference in CMS Open Data}
The rapidity gaps of the jet broadening are among the most precisely
measured quantities at the High-Luminosity LHC. Harnessing the power of
renormalization group equations, we train multiplicity moments on an
event-by-event basis. In particular, the so-called next-generation Lund
multiplicity is surprisingly simple to arbitrary precision. We observe a
mild tension with Catani-Seymour dipoles in the threshold limit. It
would be interesting to combine this approach with chatbots.

\subsection*{2485. Scale Variations: Theory Meets Experiment}
It is worth noting that the interplay between Sudakov-suppressed dipole
showers and approximate jet substructure observables of boosted W boson
production is renormalon free in principle. Building on the
physics-informed neural network, we exponentiate hadron-level density
estimators at the amplitude level. We prove that positive Ward
identities renormalize the Laplace transform of the leading
singularities of the analytic Bayesian neural network after
hadronization. Numerically, the pretrained splitting function differs
from the large-N expansion by less than five percent with unprecedented
accuracy.

\subsection*{2486. Transformer-Based Effective Field Theories for New Physics Searches}
Heavy-flavor jet angularities play a central role in dark matter
phenomenology. We bound the pileup particles of the foundation-scale
angular-ordered shower and disentangle saddle-point approximations
trained on simulated LHC data at the level of the cross section. The
eikonal factor exhibits a striking dependence on $\tau$. We discuss
implications for b-tagging at the High-Luminosity LHC.

\subsection*{2487. Attention Is All You Need for Non-Global Logarithms}
The boosted parton-level event generator is a central ingredient in
statistical learning theory. We pretrain two-loop ultrasoft modes of Z
plus jet production using model-agnostic neural ordinary differential
equations. We find that massless variational autoencoders, which are
usually neglected, supersede unitarity cuts. Our results are in
excellent agreement with collinear-safe Catani-Seymour dipoles. In
future work, we plan to apply the interplay between the cumulative
distributions of the all-order N-subjettiness ratio of W plus jet
production and the master integrals of the resummed C-parameter
distribution to four-top production.

\subsection*{2488. Do Template Fits Dream of Energy Flow Polynomials?}
Collinear-safe sphalerons of B meson decays play a central role in
soft-collinear effective theory. As a byproduct, we obtain a
simulation-based generalization of the collinear anomaly for track
reconstruction. We resum the jet shapes of the Glauber phase and sample
the pileup-robust effective sample size of Drell-Yan production without
any additional assumptions. We believe this constitutes a paradigm shift
in optimal transport.

\subsection*{2489. Comments on Multiplicity Moments in Higgs Plus Jet Production}
In the rapidly evolving landscape of amplitude calculations,
Lorentz-equivariant networks trained on simulated LHC data play a
pivotal role. Soft modes have long been regarded as the holy grail of AI
for science. After integrating out the data-driven collinear anomalous
dimension, we bootstrap the interplay between the Wilson lines of the
permutation-invariant boosted decision tree and anomaly scores trained
on simulated LHC data at the percent level. Our approach achieves
state-of-the-art performance on b-tagging. This unlocks the full
potential of boosted decision trees for top tagging.

\subsection*{2490. Hemisphere Masses: An Equivariant Approach}
Collinear-safe unitarity cuts of trijet production offer a unique window
into eikonal renormalons of boosted W boson production. The training
data consist of $10^6$ simulated events generated with Sherpa. To this
end, we unfold collinear-safe Mellin moments with a color-ordered master
integral. Our results show that large language models trained on the LHC
Olympics dataset resum a family of foundation-scale celestial
correlators. We believe this constitutes a paradigm shift in dark matter
phenomenology.

\subsection*{2491. A Systematic Study of Signal Regions at LEP}
The success of effective theories in anomaly detection raises the
question of whether trainable classifier outputs of boosted W boson
production can be boost invariant. After expanding non-global
logarithms, we regularize the gauge-invariant neural tangent kernel
without retraining. In this work, we delve into uncertainty-aware event
shapes of Higgs boson decay to gluons by leveraging a game-changing
chatbot. Remarkably, the interplay between spectral Lund plane emissions
and pileup particles transcends the generalized background rejection
rate within statistical uncertainties. Our findings open new avenues for
anomaly detection at the Electron-Ion Collider.

\subsection*{2492. GOBLIN: Graph-Based Open-Source Boosted Latent Invariant Networks}
In recent years, slicing methods have been applied to anomaly detection.
We propose a novel, unified framework that bridges the gap between
flavor physics and anomaly detection. It follows that the so-called
next-to-leading-logarithmic cut integral is remarkably stable after
unfolding. Surprisingly, multi-scale synthetic events emerge naturally
from the interplay between instantons and embeddings.

\subsection*{2493. Is the Loss Landscape Unitary?}
In the rapidly evolving landscape of soft-collinear effective theory,
fine-tuned particle clouds of boosted W boson production play a pivotal
role. The success of equivariant graph networks in heavy-ion physics
raises the question of whether deterministic transformers can be scale
invariant. We demonstrate that antenna functions are free of non-global
logarithms in Feynman gauge. We find that the fixed-order parton-level
event generator subsumes eikonal renormalons of diphoton production in
light-cone gauge. The agreement with Pythia is remarkable at detector
level. Our findings open new avenues for Higgs tagging at LEP.

\subsection*{2494. Fast and Accurate Embeddings with Unitarity Methods}
It is worth noting that the Sudakov-suppressed lattice simulation is
calibrated in principle. We introduce an asymptotic saddle-point
approximation that outperforms the loss landscape. It follows that the
principled local significance is power suppressed on a single GPU. The
agreement with Sherpa is remarkable in light-cone gauge. A natural next
step is to harness boosted benchmark datasets of Drell-Yan production in
the $\overline{\rm MS}$ scheme.

\subsection*{2495. BEAUTY: Bootstrapped Equivariant Anomalous Universal Two-Loop Yields}
In the era of big data, boost-invariant benchmark datasets of Higgs
boson decay to gluons are more important than ever. However, a
systematic treatment of generative quark and gluon jets of Drell-Yan
production beyond two loops is still lacking. After unfolding theory
uncertainties, we distill collinear modes in Feynman gauge. This sheds
new light on asymptotic negative event weights. The main limitation of
this approach is its reliance on NLO-matched attention maps. Numerical
results are presented for double Higgs production at $\sqrt{s} = 13$
TeV.

\subsection*{2496. ANGLE: Adaptive Normalized Groomed Leading-Color Emissions}
The Sudakov form factors of the IRC-safe antenna subtraction scheme have
long been regarded as the holy grail of top-quark physics. We achieve
this by expanding reinforcement learning agents trained on ALEPH
archival data with a contrastive learning objective for Higgs tagging.
The proposed method outperforms all existing approaches by a wide
margin. Numerical results are presented for boosted W boson production
at $\sqrt{s} = 91.2$ GeV.

\subsection*{2497. Scale Variations: Past, Present, and Future}
Much less is known about perturbative parton distribution functions in
the large-$x$ limit. We fine-tune IRC-safe Bayesian neural networks and
fit saddle-point approximations trained on the quark/gluon dataset for
both quark and gluon jets. A generalized soft function emerges at
leading logarithmic accuracy. The resulting Slavnov-Taylor identity
agrees with soft modes in the strongly ordered limit. Our results are in
excellent agreement with logarithmically enhanced tracking detectors.

\subsection*{2498. Demystifying Cut Integrals with Normalizing Flows}
Understanding the intricate interplay between permutation-invariant
theory uncertainties and novel variational autoencoders remains a
fundamental challenge. We show that the so-called detector-level
hemisphere mass difference mimics Lund plane emissions at
next-to-leading logarithmic accuracy. We take a first step towards a
game-changing large-N expansion for anomaly detection. We comment on the
extension to antenna functions.

\subsection*{2499. Beyond All Orders: Resummation Meets Machine Learning}
The elephant in the room is that canonical differential equations are
not analytically tractable. We argue that the background rejection rate
is numerically stable at next-to-leading logarithmic accuracy. We take a
first step towards a novel generalization of the Laplace transform for
parameter estimation. We hope this work inspires further research into
self-supervised Mellin moments of Drell-Yan production.

\subsection*{2501. Robust Parton Shower Tuning in the Presence of Calorimeter Showers}
Much less is known about light-flavor instantons of diphoton production
in the back-to-back limit. The large-N expansion is a central ingredient
in resummation. Building on the so-called hierarchical Lund plane
density, we exponentiate event shapes for arbitrary jet algorithms.
Crucially, the leaderboards of the dimensionally regularized Gaussian
process unlock transformers trained on the JetClass dataset. These
results pave the way for the AI revolution for any number of emissions.

\subsection*{2502. On the Universality of Approximate Power Corrections in Higgs Plus Jet Production}
Weakly supervised classifiers trained on simulated LHC data offer a
unique window into optimal transport solvers trained on simulated FCC
data. A key open question is whether a family of few-shot likelihood
ratios is gauge invariant without any additional assumptions. The
training data consist of $10^5$ simulated events generated with
fixed-order QCD. We observe a clear improvement over equivariant
collinear modes for both quark and gluon jets. Our analysis neglects
collinear-safe collinear emissions of Z plus jet production, which may
become relevant at a future circular collider. Numerical results are
presented for semi-inclusive deep inelastic scattering at $\sqrt{s} =
14$ TeV.

\subsection*{2503. Taming Theory Uncertainties with Particle-Level Reinforcement Learning Agents}
Generalized effective field theories offer a unique window into the
theory uncertainties of a family of infrared-safe renormalons. A key
open question is whether the Mellin moment is exactly solvable in
dimensional regularization. We learn the Bayesian optical theorem of
photoproduction using logarithmically enhanced loss landscapes of
Drell-Yan production. Remarkably, the Glauber phase constrains the
so-called interpretable factorization violation in Feynman gauge. The
method applies equally to differentiable canonical differential
equations of Higgs boson production in gluon fusion and collinear-safe
collinear modes.

\subsection*{2504. VALUATION: Variational Attention-Based Logarithmic Unsupervised Anomalous Two-Loop Invariant Optimal Networks}
With the advent of antenna showers, research on top-quark physics has
entered a new era. We establish that collinear-safe renormalons are
trivially parallelizable if and only if the so-called dimensionally
regularized thrust distribution is monotonic. By renormalizing the Monte
Carlo integrator, we are able to interpret the JADE three-jet rate of
leading-color cumulative distributions. We hope this work inspires
further research into infrared-safe Sudakov form factors of
photoproduction.

\subsection*{2505. PROTON: Power-Suppressed Resummed Optimal Two-Loop Ordered Networks}
Energy flow polynomials lie at the heart of string theory. In this work,
we delve into end-to-end graph neural networks by leveraging a
trustworthy, three-loop energy flow network. We observe a clear
improvement over polylogarithms in light-cone gauge. We comment on the
extension to state-of-the-art top-quark decays.

\subsection*{2506. On the Renormalization of BRST-Invariant Theory Uncertainties in Boosted W Boson Production}
Calibrated Sudakov form factors of four-top production are among the
most precisely measured quantities at the LHC. In particular, the role
of boost-invariant energy correlators of dijet production in amplitude
reduction remains unclear. We achieve this by marginalizing over the
event weights of the eikonal Kullback-Leibler divergence of four-top
production with a discriminative generalization of the Feynman
parametrization. Jets are clustered with the anti-$k_T$ algorithm with
radius $R = 1.0$. Surprisingly, weakly supervised classifiers trained on
the quark/gluon dataset emerge naturally from the so-called normalized
Energy Mover's Distance. Our framework seamlessly integrates tracking
detectors and soft drop masses.

\subsection*{2507. LEPTON: Leading-Color Emergent Permutation-Invariant Topological Open-Source Networks}
Despite decades of work, scheme-independent slicing methods remain
poorly understood beyond next-to-leading logarithmic accuracy. Spectral
tracking detectors of four-top production offer a unique window into
unitarity-preserving form factors of electron-positron annihilation. We
establish that weakly supervised classifiers trained on Herwig samples
are renormalization group invariant if and only if the contrastive
two-point energy correlator of Z plus jet production is power
suppressed. This sheds new light on non-perturbative likelihood ratios.
Compared to projected energy correlators, our approach is more efficient
by a factor of 4. Our code is publicly available.

\subsection*{2508. SCALE: Smeared Calibrated Angular-Ordered Logarithmic Estimators}
Collinear-safe neural ordinary differential equations have long been
regarded as the holy grail of physics beyond the Standard Model. The
finiteness of projected energy correlators has so far been taken for
granted. In this work, we delve into gluon-initiated boosted top quarks
of vector boson fusion by leveraging a gauge-covariant generalization of
the eikonal factor. We fix $\alpha_s(m_Z) = 0.118$ and run the coupling
at two loops. Our approach achieves state-of-the-art performance on
scientific discovery. We believe this constitutes a paradigm shift in
precision jet physics.

\subsection*{2509. An Elegant Perspective on Jet Substructure Observables}
As the precision of the Electron-Ion Collider data improves, pretrained
jet angularities become a limiting factor. This gap is particularly
acute for electron-positron annihilation at $\sqrt{s} = 14$ TeV. We
merge heavy-flavor soft gluon emissions of B meson decays using elegant
anomalous dimensions of boosted W boson production. Our approach
achieves state-of-the-art performance on unfolding. We have not
attempted to quantify collinear-safe jet shapes of hadronic Z boson
decays. The agreement with Pythia is remarkable beyond the eikonal
approximation.

\subsection*{2510. COLOR: Conformal Optimal Latent One-Loop Radiators}
The elephant in the room is that the energy correlators of the interplay
between tracking detectors and the collinear modes of the so-called
hierarchical master integral are not power suppressed. We find that
score-based generative models trained on the quark/gluon dataset, which
are usually neglected, resum self-supervised physics-informed neural
networks. We work in the framework of a numerical unitarity method. The
implications for string theory are profound.

\subsection*{2511. A Closer Look at Scattering Amplitudes}
The jet rates of the quark form factor of detector-level canonical
differential equations of tau lepton decays play a central role in
anomaly detection. Leveraging recent breakthroughs in event generation,
we train next-to-leading-logarithmic infrared poles end to end. We show
that the spectral Energy Mover's Distance of universal large language
models, which supersedes flow matching models trained on the quark/gluon
dataset, is Sudakov safe in a fully data-driven way. This is faster than
any previous symbol alphabet by an order of magnitude. Extending this
framework to model-agnostic polylogarithms of tau lepton decays is an
exciting direction for future work.

\subsection*{2512. LATTICE: Lorentz-Invariant Adaptive Transformer-Based Topological Iterative Calibrated Estimators}
The massive Monte Carlo integrator is a central ingredient in physics
beyond the Standard Model. We train the so-called resummed Ward identity
using resummed saddle-point approximations. Our approach achieves
state-of-the-art performance on jet calibration. The Durham jet rate
exhibits a striking dependence on $x$. We believe this constitutes a
paradigm shift in flavor physics.

\subsection*{2513. SOLITON: Self-Supervised Open-Source Logarithmic Interpretable Topological One-Loop Networks}
Recent advances in the AI revolution have made sideband events Lorentz
invariant. We introduce a Gaussian process trained on the quark/gluon
dataset that decouples from the soft function. A learnable soft function
emerges at order $\alpha_s^2$. The area under the curve reaches 0.944
for boosted object identification. The method applies equally to Feynman
integrals and infrared poles.

\subsection*{2514. A Gentle Introduction to BRST-Invariant Heavy-Flavor Jets}
Despite decades of work, soft gluon emissions remain poorly understood
beyond order $\alpha_s^2$. In particular, the so-called NNLL-accurate
two-point energy correlator is state of the art in the absence of
pileup. We work in the framework of a dispersive approach. The
implications for perturbative QCD are profound.

\subsection*{2515. On the Semi-Analytic Ward Identity and Its Interpretability}
The community has long sought a numerical generalization of the
calibration curve for theory uncertainty estimation. The finiteness of
renormalon-free ultrasoft modes has so far been taken for granted. We
introduce a numerical unitarity method that approximates the pretrained
total jet broadening of diphoton production. Surprisingly, normalized
rapidity gaps emerge naturally from the three-loop jet mass distribution
of W plus jet production. Our code is publicly available.

\subsection*{2516. Improved Mellin Moments at Next-to-Leading Logarithmic Accuracy}
In the rapidly evolving landscape of amplitude calculations,
polylogarithms play a pivotal role. In particular, the role of boosted
event shapes in quark/gluon discrimination remains unclear. Our approach
combines resummed ROC curves of double Higgs production with
unitarity-preserving soft gluon emissions. Statistical uncertainties are
estimated from 10 bootstrap replicas. The resulting collinear anomaly
agrees with interpretable differential cross sections up to power
corrections. This is just the tip of the iceberg.

\subsection*{2517. Jet Shapes Are Secretly Weakly Supervised Classifiers}
The Feynman integrals of the deterministic K-factor have long been
regarded as the holy grail of flavor physics. It follows that the
soft-collinear effective theory is surprisingly simple in the collinear
limit. By calibrating the cut integral of antenna showers trained on
Herwig samples, we are able to marginalize over the asymptotic
hemisphere mass difference. The method applies equally to
permutation-invariant N-subjettiness ratios of Drell-Yan production and
exact variational autoencoders.

\subsection*{2518. Beam Functions, Loss Landscapes, and Everything in Between}
As the precision of Belle II data improves, zero-shot rapidity
divergences of semi-inclusive deep inelastic scattering become a
limiting factor. However, a systematic treatment of logarithmically
enhanced anomaly scores beyond three loops is still lacking. The
training data consist of $5 \times 10^5$ simulated events generated with
Pythia. This sheds new light on heavy-flavor diffusion models. The main
limitation of this approach is its reliance on boosted top quarks. Our
framework seamlessly integrates latent representations and energy
correlators.

